%============================================================
%  Bd. 07 - Die natürlichen Zahlen aus Mengenlehre und Peano-Axiomen
%============================================================

\documentclass[main.tex]{subfiles}

\ifSubfilesClassLoaded{
  \BandSetupStandalone{B07}
}{
}

\title{Bd. 07 - Die natürlichen Zahlen aus Mengenlehre und Peano-Axiomen}
\author{Martin Kunze}
\date{}
\setcounter{file}{7}

\begin{document}

\title{Bd. 07 - Die natürlichen Zahlen aus Mengenlehre und Peano-Axiomen}
\author{Martin Kunze}
\date{}
\hypersetup{pageanchor=false}
\maketitle
\hypersetup{pageanchor=true}
\setcounter{tocdepth}{3}
\tableofcontents
\setcounter{file}{7}
\renewcommand{\FormulaBandID}{7}

\chapter{Einleitung}

Der erste Teil dieses Bandes konstruiert die natürlichen Zahlen aus den
Mengenaxiomen von Band~03. Der zweite Teil nimmt die Peano-Axiome als
eigenständigen Ausgangspunkt und entwickelt daraus Rekursion, Arithmetik,
Ordnung und Stellenwertsysteme. Die allgemeine Funktionstheorie aus Band~04
und die allgemeine Ordnungstheorie aus Band~06 werden dabei nicht erneut
axiomatisiert.

\chapter{Konstruktion der natürlichen Zahlen}

Wir konstruieren die Menge der natürlichen Zahlen allein aus den Axiomen der Mengenlehre, insbesondere dem Unendlichkeitsaxiom.  Dieser Ansatz kommt ohne eine explizite Nachfolgerfunktion als primitives Symbol aus; diese wird vielmehr aus der Mengenoperation \(x \cup \{x\}\) gewonnen.


\section{Axiomatische Grundlage}

Das Unendlichkeitsaxiom ist bereits in Band~03 in funktionsfreier Form
registriert. Hier beginnt seine Anwendung auf induktive Mengen und die
von-Neumann-Konstruktion.

\section{Nachfolger und induktive Mengen}

\subsection{Induktive Mengen}

\FormulaDefDelta[induktive Menge]{\Induktiv(A) \coloneq \varnothing \in A \,\land\, \forall x\in A\,(x \cup \{x\} \in A)}{
\DeltaDecl{Mengen}{A}
}

% Existenz einer induktiven Menge (aus dem Unendlichkeitsaxiom)
\FormulaThmAuto{\exists A(\Induktiv(A))}
\begin{tabproof}
  \proofstep{}{ \exists A\bigl(\varnothing \in A \land \forall x \in A\,(x \cup \{x\} \in A)\bigr) }{\FormulaRefAuto{\exists A\bigl(\varnothing \in A \land \forall x \in A\,(x \cup \{x\} \in A)\bigr)}}
  \proofstep{}{\exists A(\Induktiv(A)) }{\rIE{\FormulaRefAuto{\Induktiv(A) \coloneq \varnothing \in A \,\land\, \forall x\in A\,(x \cup \{x\} \in A)},1}}
\end{tabproof}

% Sofortige Konsequenzen aus der Definition
\FormulaThmDelta{\Induktiv(A)\vdash \varnothing\in A}{
\DeltaDecl{Mengen}{A}
}
\begin{tabproof}
  \proofstep{1}{\Induktiv(A)}{\rA}
  \proofstep{1}{\varnothing \in A \land \forall x \in A\,(x \cup \{x\} \in A)}{\rIE{\FormulaRefAuto{\Induktiv(A) \coloneq \varnothing \in A \,\land\, \forall x\in A\,(x \cup \{x\} \in A)},1}}
  \proofstep{1}{\varnothing \in A}{\rAEa{2}}
\end{tabproof}

\FormulaThmDelta{\Induktiv(A)\vdash \forall x \in A\,(x \cup \{x\} \in A)}{
\DeltaDecl{Mengen}{A}
}
\begin{tabproof}
  \proofstep{1}{\Induktiv(A)}{\rA}
  \proofstep{1}{\varnothing \in A \land \forall x \in A\,(x \cup \{x\} \in A)}{\rIE{\FormulaRefAuto{\Induktiv(A) \coloneq \varnothing \in A \,\land\, \forall x\in A\,(x \cup \{x\} \in A)},1}}
  \proofstep{1}{\forall x \in A\,(x \cup \{x\} \in A)}{\rAEb{2}}
\end{tabproof}

\FormulaThmDelta{\Induktiv(A),\,x\in A\vdash x \cup \{x\} \in A}{
\DeltaDecl{Mengen}{A}
}
\begin{tabproof}
  \proofstep{1}{\Induktiv(A)}{\rA}
  \proofstep{2}{x\in A}{\rA}
  \proofstep{1}{\forall u \in A\,(u \cup \{u\} \in A)}{\FormulaRefAuto{\Induktiv(A)\vdash \forall x \in A\,(x \cup \{x\} \in A)}{1}}
  \proofstep{1,2}{x \cup \{x\} \in A}{\rRE{\rUE{3},2}}
\end{tabproof}

% Schnitt induktiver Mengen ist induktiv
\FormulaThmDelta{\Induktiv(A),\, \Induktiv(B)\vdash \Induktiv(A\cap B)}{
\DeltaDecl{Mengen}{A\dsep B}
}
\begin{tabproof}
  \proofstep{1}{\Induktiv(A)}{\rA}
  \proofstep{2}{\Induktiv(B)}{\rA}

  % Null in A∩B
  \proofstep{1}{\,\varnothing\in A}{\FormulaRefAuto{\Induktiv(A)\vdash \varnothing\in A}{1}}
  \proofstep{2}{\,\varnothing\in B}{\FormulaRefAuto{\Induktiv(A)\vdash \varnothing\in A}{2}}
  \proofstep{1,2}{\,\varnothing\in A\cap B}{\FormulaRefAuto{x\in A,\, x\in B \vdash x\in A\cap B}{3,4}}

  % Abschluss unter Nachfolger-Operation für A∩B
  \proofstep{1}{\forall x\in A\,(x\cup\{x\}\in A)}{\FormulaRefAuto{\Induktiv(A)\vdash \forall x \in A\,(x \cup \{x\} \in A)}{1}}
  \proofstep{2}{\forall x\in B\,(x\cup\{x\}\in B)}{\FormulaRefAuto{\Induktiv(A)\vdash \forall x \in A\,(x \cup \{x\} \in A)}{2}}
  \proofstep{8}{x\in A\cap B}{\rA}
  \proofstep{8}{x\in A}{\FormulaRefAuto{x\in A\cap B \vdash x\in A}{9}}
  \proofstep{8}{x\in B}{\FormulaRefAuto{x\in A\cap B \vdash x\in B}{9}}
  \proofstep{1,8}{x\cup\{x\}\in A}{\rRE{\rUE{6},9}}
  \proofstep{2,8}{x\cup\{x\}\in B}{\rRE{\rUE{7},10}}
  \proofstep{1,2,8}{x\cup\{x\}\in A\cap B}{\FormulaRefAuto{x\in A,\, x\in B \vdash x\in A\cap B}{11,12}}
  \proofstep{1,2}{\,\forall x\in A\cap B\,(x\cup\{x\}\in A\cap B)}{\rUI{\rRI{8,13}}}

  % Schluss: Induktivität von A∩B
  \proofstep{1,2}{\,\varnothing\in A\cap B \;\land\; \forall x\in A\cap B\,(x\cup\{x\}\in A\cap B)}{\rAI{5,14}}
  \proofstep{1,2}{\,\Induktiv(A\cap B)}%
    {\rIE{\FormulaRefAuto{\Induktiv(A) \coloneq \varnothing \in A \,\land\, \forall x\in A\,(x \cup \{x\} \in A)},15}}
\end{tabproof}

\subsection{Die Nachfolger-Funktion}

% ============================================================
% Nachfolger-Funktion (analog zu den Beispielen von Funktionen)
% Voraussetzung: Induktiv(A) wurde bereits definiert.
% ============================================================

\subsubsection{Definition der Nachfolgerfunktion}

\FormulaDefDelta[Funktionsvorschrift für die Nachfolgerfunktion auf \(A\)]%
{x\in A \vdash t_{\Succ_A}(x) \coloneqq x\cup\{x\}}%
{
  \DeltaRow{Mengen}{A\dsep x}
  \DeltaRow{Funktionensymbol}{t_{\Succ_A}}
}

\FormulaThmDelta[Typisierungsaxiom für \(t_{\Succ_A}\)]%
{\Induktiv(A)\dsep x\in A \vdash t_{\Succ_A}(x)\in A}%
{
  \DeltaRow{Mengen}{A\dsep x}
  \DeltaRow{Funktionensymbol}{t_{\Succ_A}}
}
\begin{tabproof}
  \proofstep{1}{\Induktiv(A)}{\rA}
  \proofstep{2}{x\in A}{\rA}
  \proofstep{1,2}{x\cup\{x\}\in A}{%
    \FormulaRefAuto{\Induktiv(A),\,x\in A\vdash x \cup \{x\} \in A}{1,2}}
  \proofstep{2}{t_{\Succ_A}(x)=x\cup\{x\}}{%
    \FormulaRefAuto{x\in A \vdash t_{\Succ_A}(x) \coloneqq x\cup\{x\}}{2}}
  \proofstep{1,2}{t_{\Succ_A}(x)\in A}{\rIE{4,3}}
\end{tabproof}

\FormulaThmDelta[Nachfolgerfunktion als Funktion]%
{\Induktiv(A)\vdash \exists! F\colon A\to A\,\forall x\in A\,F(x)=t_{\Succ_A}(x)}%
{
  \DeltaRow{Mengen}{A\dsep x}
  \DeltaRow{Funktionensymbol}{t_{\Succ_A}}
}
\begin{tabproof}
  \proofstep{1}{\Induktiv(A)}{\rA}
  \proofstep{}{%
    \forall x\in A\,t_{\Succ_A}(x)=x\cup\{x\}}{%
    \FormulaRefAuto{x\in A \vdash t_{\Succ_A}(x) \coloneqq x\cup\{x\}}}
  \proofstep{1}{%
    \forall x\in A\,t_{\Succ_A}(x)\in A}{%
    \rUI{\rRI{2,\FormulaRefAuto{\Induktiv(A)\dsep x\in A \vdash t_{\Succ_A}(x)\in A}{1,2}}}}
  \proofstep{1}{%
    \exists! F\colon A\to A\,\forall x\in A\,F(x)=t_{\Succ_A}(x)}{%
    \FormulaRefAuto{x\in A\vdash t(x)\coloneq y\dsep x\in A\vdash t(x)\in B\exists! F\colon A\to B\,\forall x\in A\,F(x) = t(x)}{2,3}}
\end{tabproof}

\FormulaDefDelta[Nachfolgerfunktion auf \(A\)]%
{\Induktiv(A)\vdash \Succ_A\coloneqq\iota F\Bigl(F\colon A\to A \land \forall x\in A\,F(x)=t_{\Succ_A}(x)\Bigr)}%
{
  \DeltaRow{Mengen}{A}
  \DeltaRow{Funktionensymbol}{t_{\Succ_A}}
}

\FormulaThmDelta[Nachfolgerfunktion auf \(A\)]%
{\Induktiv(A)\vdash \Succ_A\colon A\to A}%
{
  \DeltaRow{Mengen}{A}
}
\begin{tabproof}
  \proofstep{1}{\Induktiv(A)}{\rA}
  \proofstep{1}{\Succ_A\colon A\to A}{%
    \rAEa{\FormulaRefAuto{\Induktiv(A)\vdash \Succ_A\coloneqq\iota F\Bigl(F\colon A\to A \land \forall x\in A\,F(x)=t_{\Succ_A}(x)\Bigr)}{1}}}
\end{tabproof}

\FormulaThmDelta%
{\Induktiv(A)\vdash \TotRel{\Succ_A,A,A}}%
{
  \DeltaRow{Mengen}{A}
}
\begin{tabproof}
  \proofstep{1}{\Induktiv(A)}{\rA}
  \proofstep{1}{\Succ_A\colon A\to A}{%
    \FormulaRefAuto{\Induktiv(A)\vdash \Succ_A\colon A\to A}{1}}
  \proofstep{1}{\TotRel{\Succ_A,A,A}}{%
    \FormulaRefAuto{F\colon A\to B\vdash\TotRel{F,A,B}}{2}}
\end{tabproof}

\FormulaThmDelta[Grundgleichung der Nachfolgerfunktion]%
{\Induktiv(A)\dsep x\in A \vdash \Succ_A(x)=x\cup\{x\}}%
{
  \DeltaRow{Mengen}{A\dsep x}
  \DeltaRow{Funktionensymbol}{t_{\Succ_A}}
}
\begin{tabproof}
  \proofstep{1}{\Induktiv(A)}{\rA}
  \proofstep{2}{x\in A}{\rA}

  \proofstep{1}{%
    \forall x\in A\,\Succ_A(x)=t_{\Succ_A}(x)}{%
    \rAEb{\FormulaRefAuto{\Induktiv(A)\vdash \Succ_A\coloneqq\iota F\Bigl(F\colon A\to A \land \forall x\in A\,F(x)=t_{\Succ_A}(x)\Bigr)}{1}}}

  \proofstep{1,2}{\Succ_A(x)=t_{\Succ_A}(x)}{\rRE{2,\rUE{3}}}

  \proofstep{2}{t_{\Succ_A}(x)=x\cup\{x\}}{%
    \FormulaRefAuto{x\in A \vdash t_{\Succ_A}(x) \coloneqq x\cup\{x\}}{2}}

  \proofstep{1,2}{\Succ_A(x)=x\cup\{x\}}{%
    \FormulaRefAuto{a = b,\, b = c \vdash a = c}{4,5}}
\end{tabproof}

\FormulaThmDelta[Bild der Nachfolgerfunktion liegt in \(A\)]{%
  \Induktiv(A)\dsep x\in A \vdash \Succ_A(x)\in A%
}{
  \DeltaDecl{Mengen}{A\dsep x}
}
\begin{tabproof}
  \proofstep{1}{\Induktiv(A)}{\rA}
  \proofstep{2}{x\in A}{\rA}

  \proofstep{1,2}{x\cup\{x\}\in A}{%
    \FormulaRefAuto{\Induktiv(A),\,x\in A\vdash x \cup \{x\} \in A}{1,2}}

  \proofstep{1,2}{\Succ_A(x)=x\cup\{x\}}{%
    \FormulaRefAuto{\Induktiv(A)\dsep x\in A \vdash \Succ_A(x)=x\cup\{x\}}{1,2}}

  \proofstep{1,2}{\Succ_A(x)\in A}{\rIE{4,3}}
\end{tabproof}

\FormulaThmDelta[Nachfolgermenge einer natürlichen Zahl in Vereinigungsform]{%
  n\in\mathbb{N}\vdash (n\cup\{n\})\in\mathbb{N}
}{%
  \DeltaRow{Mengen}{n}
}
\begin{tabproof}
  \proofstep{1}{n\in\mathbb{N}}{\rA}
  \proofstep{1}{\Succ(n)=n\cup\{n\}}{%
    \FormulaRefAuto{n\in\mathbb{N}\vdash \Succ(n)=n\cup\{n\}}{1}}
  \proofstep{1}{\Succ(n)\in\mathbb{N}}{%
    \FormulaRefAuto{n\in\mathbb{N}\vdash \Succ(n)\in\mathbb{N}}{1}}
  \proofstep{1}{(n\cup\{n\})\in\mathbb{N}}{\rIE{2,3}}
\end{tabproof}



\section{Konstruktion der natürlichen Zahlen}

\FormulaThmAuto{
  \exists! C\; \forall B \Bigl(\Induktiv(B) \rightarrow
    C = \{ x \in B \mid \forall A (\Induktiv(A) \rightarrow x \in A) \}\Bigr)
}
\begin{tabproofwide}
  % Schritt 1: Existenz einer induktiven Menge
  \proofstepwidestar[]{\exists A(\Induktiv(A))}%
    {\FormulaRefAuto{\exists A(\Induktiv(A))}}

  % Schritt 2: Einzigartige Existenz des Schnitts aller induktiven Mengen
  \proofstepwide{}{}{\exists! C\; \forall B(\Induktiv(B)\rightarrow }%
    {\multirow{2}{*}{\FormulaRefAuto{P(D)\vdash \exists! C\forall B(P(B)\rightarrow C= \{ x \in B \mid \forall A (P(A) \rightarrow x \in A) \})}{1}}}
  \proofstepwide{}{}%
    {C=\{x\in B \mid \forall A (\Induktiv(A)\rightarrow x\in A)\})}%
    {}
\end{tabproofwide}


\FormulaDefAuto[Menge der natürlichen Zahlen]{\mathbb{N} \coloneq \bigcap \{ A \mid \Induktiv(A) \}}

\begin{remark}
  Die Definition sagt: Ein Element \(x\) gehört genau dann zu \(\mathbb{N}\), wenn es zu jeder induktiven Menge gehört.  Damit ist \(\mathbb{N}\) die \emph{kleinste} induktive Menge, denn für jede andere induktive Menge \(B\) gilt \(\mathbb{N} \subseteq B\).
\end{remark}

\FormulaThmAuto{x\in\mathbb{N}\eqvdash \forall \Induktiv(N)(x\in N)}
\begin{tabproofwide}
    \proofstepwide[]{x\in \mathbb{N}}{\leftrightarrow}{x\in \bigcap \{ A \mid \Induktiv(A) \}}{\FormulaRefAuto{\mathbb{N} \coloneq \bigcap \{ A \mid \Induktiv(A) \}}}
    \proofstepwide[]{}{\leftrightarrow}{\forall \Induktiv(N)(x\in N)}{\FormulaRefAuto{\exists A(P(A)) \vdash x \in \bigcap_{P(B)} B \leftrightarrow \forall C\, (P(C) \rightarrow x \in C)}{\FormulaRefAuto{\exists A(\Induktiv(A))}}}
\end{tabproofwide}

\section{Eigenschaften der Menge der natürlichen Zahlen}

Im Folgenden werden wir einige grundlegende Eigenschaften der so konstruierten Menge \(\mathbb{N}\) beweisen.

\subsection{Induktivität der Menge der natürlichen Zahlen}


\FormulaThmAuto[Null-Axiom]{\varnothing\in\mathbb{N}}
\begin{tabproof}
    \proofstep{1}{\Induktiv(N)}{\rA}
    \proofstep{1}{\varnothing\in N}{\FormulaRefAuto{\Induktiv(A)\vdash \varnothing\in A}{1}}
    \proofstep{}{\forall \Induktiv(N)(x\in N)}{\rUI{\rRI{1,2}}}
    \proofstep{}{\varnothing\in\mathbb{N}}{\FormulaRefAuto{x\in\mathbb{N}\eqvdash \forall \Induktiv(N)(x\in N)}{3}}
\end{tabproof}

\FormulaThmAuto[Nachfolger-Axiom]{\forall x \in \mathbb{N}\,(x\cup\{x\} \in \mathbb{N})}
\begin{tabproof}
  \proofstep{1}{x \in \mathbb{N}}{\rA}

  \proofstep{1}{\forall \Induktiv(N)(x\in N)}{%
    \FormulaRefAuto{x\in\mathbb{N}\eqvdash \forall \Induktiv(N)(x\in N)}{1}}

  \proofstep{3}{\Induktiv(N)}{\rA}

  \proofstep{1}{\Induktiv(N)\rightarrow x\in N}{\rUE{2}}
  \proofstep{1,3}{x\in N}{\rRE{4,3}}

  \proofstep{1,3}{x\cup\{x\}\in N}{%
    \FormulaRefAuto{\Induktiv(A),\,x\in A\vdash x \cup \{x\} \in A}{3,5}}

  \proofstep{1}{\Induktiv(N)\rightarrow x\cup\{x\}\in N}{\rRI{3,6}}
  \proofstep{1}{\forall \Induktiv(N)(x\cup\{x\}\in N)}{\rUI{7}}

  \proofstep{1}{x\cup\{x\}\in \mathbb{N}}{%
    \FormulaRefAuto{x\in\mathbb{N}\eqvdash \forall \Induktiv(N)(x\in N)}{8}}

  \proofstep{}{\,\forall x\in \mathbb{N}\,(x\cup\{x\}\in \mathbb{N})}{\rUI{\rRI{1,9}}}
\end{tabproof}


\FormulaThmAuto{\Induktiv(\mathbb{N})}
\begin{tabproof}
  \proofstep{}{\,\varnothing\in\mathbb{N}}{%
    \FormulaRefAuto{\varnothing\in\mathbb{N}}}

  \proofstep{}{\,\forall x\in\mathbb{N}\,(x\cup\{x\}\in\mathbb{N})}{%
    \FormulaRefAuto{\forall x \in \mathbb{N}\,(x\cup\{x\} \in \mathbb{N})}}

  \proofstep{}{\,\varnothing\in\mathbb{N} \;\land\; \forall x\in\mathbb{N}\,(x\cup\{x\}\in\mathbb{N})}{%
    \rAI{1,2}}

  \proofstep{}{\,\Induktiv(\mathbb{N})}{%
    \rIE{\FormulaRefAuto{\Induktiv(A) \coloneq \varnothing \in A \,\land\, \forall x\in A\,(x \cup \{x\} \in A)},3}}
\end{tabproof}


\subsection{Die Nachfolgerfunktion}

\FormulaDefAuto[Nachfolgerfunktion]{%
  \Succ \coloneqq \Succ_{\mathbb{N}}%
}

\FormulaThmAuto{%
  \Succ\colon \mathbb{N}\to \mathbb{N}%
}
\begin{tabproof}
  \proofstep{}{\,\Induktiv(\mathbb{N})}{%
    \FormulaRefAuto{\Induktiv(\mathbb{N})}}
  \proofstep{}{\,\Succ_{\mathbb{N}}\colon \mathbb{N}\to \mathbb{N}}{%
    \FormulaRefAuto{\Induktiv(A)\vdash \Succ_A\colon A\to A}}
  \proofstep{}{\,\Succ=\Succ_{\mathbb{N}}}{%
    \FormulaRefAuto{\Succ \coloneqq \Succ_{\mathbb{N}}}}
  \proofstep{}{\,\Succ\colon \mathbb{N}\to \mathbb{N}}{\rIE{3,2}}
\end{tabproof}

\FormulaThmDelta{%
  n\in\mathbb{N}\vdash \Succ(n)\in\mathbb{N}%
}{%
  \DeltaDecl{Mengen}{n}%
}
\begin{tabproof}
  \proofstep{1}{n\in\mathbb{N}}{\rA}
  \proofstep{}{\,\Succ\colon \mathbb{N}\to \mathbb{N}}{%
    \FormulaRefAuto{\Succ\colon \mathbb{N}\to \mathbb{N}}}
  \proofstep{1}{\,\Succ(n)\in \mathbb{N}}{%
    \FormulaRefAuto{F\colon A\to B\dsep x\in A\vdash F(x)\in B}{1,2}}
\end{tabproof}


\FormulaThmDelta{%
  x\in\mathbb{N}\vdash \Succ(x)=x\cup\{x\}%
}{
  \DeltaDecl{Mengen}{x}
}
\begin{tabproof}
  \proofstep{1}{x\in\mathbb{N}}{\rA}
  \proofstep{}{\,\Induktiv(\mathbb{N})}{%
    \FormulaRefAuto{\Induktiv(\mathbb{N})}}
  \proofstep{1}{\,\Succ_{\mathbb{N}}(x)=x\cup\{x\}}{%
    \FormulaRefAuto{\Induktiv(A)\dsep x\in A \vdash \Succ_A(x)=x\cup\{x\}}{2,1}}
  \proofstep{}{\,\Succ=\Succ_{\mathbb{N}}}{%
    \FormulaRefAuto{\Succ \coloneqq \Succ_{\mathbb{N}}}}
  \proofstep{1}{\,\Succ(x)=x\cup\{x\}}{\rIE{4,3}}
\end{tabproof}

\subsubsection{Mengen-theoretische Eigenschaften der Nachfolgermenge}

\paragraph{Zerlegung der Nachfolgermenge}


\FormulaThmDelta{%
  n\in\mathbb{N}\dsep x\in\Succ(n)\vdash x\in n\lor x=n%
}{%
  \DeltaDecl{Mengen}{n\dsep x}%
}
\begin{tabproof}
  \proofstep{1}{n\in\mathbb{N}}{\rA}
  \proofstep{2}{x\in\Succ(n)}{\rA}
  \proofstep{1}{\Succ(n)=n\cup\{n\}}{%
    \FormulaRefAuto{n\in\mathbb{N}\vdash \Succ(n) = n\cup\{n\}}{1}}
  \proofstep{1,2}{x\in n\cup\{n\}}{\rIE{3,2}}
  \proofstep{1,2}{x\in n\lor x\in\{n\}}{%
    \FormulaRefAuto{z \in A \cup B \eqvdash z \in A \lor z \in B}{2}}

  \proofstep{3}{x\in n}{\rA}
  \proofstep{3}{x\in n\lor x=n}{\rOIa{6}}

  \proofstep{4}{x\in\{n\}}{\rA}
  \proofstep{4}{x=n}{\FormulaRefAuto{x \in \{a\} \eqvdash x = a}{4}}
  \proofstep{4}{x\in n\lor x=n}{\rOIb{9}}

  \proofstep{1,2}{x\in n\lor x=n}{\rOE{5,6,7,8,10}}
\end{tabproof}

% ------------------------------------------------------------
% Hilfssatz: Element in Succ(n) und ungleich n liegt in n
% ------------------------------------------------------------
\FormulaThmDelta{%
  n\in\mathbb{N}\dsep z\in\Succ(n)\dsep z\neq n \vdash z\in n%
}{%
  \DeltaDecl{Mengen}{n\dsep z}%
}
\begin{tabproof}
  \proofstep{1}{n\in\mathbb{N}}{\rA}
  \proofstep{2}{z\in\Succ(n)}{\rA}
  \proofstep{3}{z\neq n}{\rA}
  \proofstep{1,2}{z\in n\lor z=n}{%
    \FormulaRefAuto{n\in\mathbb{N}\dsep x\in \Succ (n)\vdash x\in n\lor x=n}{1,2}}
  \proofstep{3,4}{z\in n}{%
    \FormulaRefAuto{P \lor Q,\neg Q \vdash P}{4,3}}
\end{tabproof}


\FormulaThmDelta[Teilmengen einer Nachfolgermenge ohne neues Element]{%
  n\in\mathbb{N}\dsep B\subseteq \Succ(n)\dsep n\notin B \vdash B\subseteq n
}{%
  \DeltaRow{Mengen}{B\dsep n\dsep x}
}
\begin{tabproof}
  \proofstep{1}{n\in\mathbb{N}}{\rA}
  \proofstep{2}{B\subseteq \Succ(n)}{\rA}
  \proofstep{3}{n\notin B}{\rA}

  \proofstep{4}{x\in B}{\rA}
  \proofstep{2,4}{x\in \Succ(n)}{\FormulaRefAuto{A\subseteq B,\, x\in A \vdash x\in B}{2,4}}
  \proofstep{1}{\Succ(n)=n\cup\{n\}}{\FormulaRefAuto{n\in\mathbb{N}\vdash \Succ(n)=n\cup\{n\}}{1}}
  \proofstep{1,2,4}{x\in n\cup\{n\}}{\FormulaRefAuto{A = B,\, x \in A \vdash x \in B}{6,5}}
  \proofstep{1,2,4}{x\in n\lor x=n}{\FormulaRefAuto{x\in A\cup\{a\} \vdash x\in A\lor x=a}{7}}
  \proofstep{3,4}{x\neq n}{\FormulaRefAuto{a \in A,\; b \not\in A \vdash a \neq b}{4,3}}

  \proofstep{10}{x\in n}{\rA}

  \proofstep{11}{x=n}{\rA}
  \proofstep{3,4,11}{\bot}{\rBI{11,9}}
  \proofstep{3,4,11}{x\in n}{\rCE{12}}

  \proofstep{1,2,3,4}{x\in n}{\rOE{8,10,10,11,13}}
  \proofstep{1,2,3}{\forall x\,\bigl(x\in B \rightarrow x\in n\bigr)}{\rUI{\rRI{4,14}}}
  \proofstep{1,2,3}{B\subseteq n}{%
    \FormulaRefAuto{A \subseteq B \coloneq \forall x\,(x\in A \rightarrow x\in B)}{15}}
\end{tabproof}




\FormulaThmDelta[Teilmengen einer Nachfolgermenge mit neuem Element]{%
  n\in\mathbb{N}\dsep B\subseteq \Succ(n)\dsep n\in B \vdash B=(B\cap n)\cup\{n\}
}{%
  \DeltaRow{Mengen}{B\dsep n\dsep x}
}
\begin{tabproofsplitwide}
  \proofpartwideindR[Hinrichtung]{%
    \begin{aligned}[t]
      &n\in\mathbb{N}\dsep B\subseteq \Succ(n)\dsep n\in B \vdash{}\\
      &B\subseteq (B\cap n)\cup\{n\}
    \end{aligned}
  }{%
    n\in\mathbb{N}\dsep B\subseteq \Succ(n)\dsep n\in B \vdash B\subseteq (B\cap n)\cup\{n\}
  }
    \proofstepwidestar[1]{n\in\mathbb{N}}{\rA}
    \proofstepwidestar[2]{B\subseteq \Succ(n)}{\rA}
    \proofstepwidestar[3]{n\in B}{\rA}

    \proofstepwidestar[4]{x\in B}{\rA}
    \proofstepwidestar[2,4]{x\in \Succ(n)}{\FormulaRefAuto{A\subseteq B\dsep x\in A\vdash x\in B}{2,4}}
    \proofstepwidestar[1]{\Succ(n)=n\cup\{n\}}{\FormulaRefAuto{n\in\mathbb{N}\vdash \Succ(n)=n\cup\{n\}}{1}}
    \proofstepwidestar[1,2,4]{x\in n\cup\{n\}}{\rIE{5,6}}
    \proofstepwidestar[1,2,4]{x\in n\lor x\in\{n\}}{\FormulaRefAuto{z\in A\cup B \eqvdash z\in A\lor z\in B}{7}}

    \proofstepwidestar[9]{x\in n}{\rA}
    \proofstepwidestar[4,9]{x\in B\cap n}{\FormulaRefAuto{x\in A\dsep x\in B\vdash x\in A\cap B}{4,9}}
    \proofstepwidestar[4,9]{x\in (B\cap n)\cup\{n\}}{\FormulaRefAuto{z\in A\vdash z\in A\cup B}{10}}

    \proofstepwidestar[12]{x\in\{n\}}{\rA}
    \proofstepwidestar[12]{x\in (B\cap n)\cup\{n\}}{\FormulaRefAuto{z\in B\vdash z\in A\cup B}{12}}

    \proofstepwidestar[1,2,4]{x\in (B\cap n)\cup\{n\}}{\rOE{8,9,11,12,13}}
    \proofstepwidestar[1,2,3]{\forall x\,\bigl(x\in B \rightarrow x\in (B\cap n)\cup\{n\}\bigr)}{\rUI{\rRI{4,14}}}
    \proofstepwidestar[1,2,3]{B\subseteq (B\cap n)\cup\{n\}}{%
      \FormulaRefAuto{A \subseteq B \coloneq \forall x\,(x\in A \rightarrow x\in B)}{15}}
  \closeproofpartwideindR

  \proofpartwideindR[Rückrichtung]{%
    \begin{aligned}[t]
      &n\in\mathbb{N}\dsep B\subseteq \Succ(n)\dsep n\in B \vdash{}\\
      &(B\cap n)\cup\{n\}\subseteq B
    \end{aligned}
  }{%
    n\in\mathbb{N}\dsep B\subseteq \Succ(n)\dsep n\in B \vdash (B\cap n)\cup\{n\}\subseteq B
  }
    \proofstepwidestar[1]{n\in\mathbb{N}}{\rA}
    \proofstepwidestar[2]{B\subseteq \Succ(n)}{\rA}
    \proofstepwidestar[3]{n\in B}{\rA}

    \proofstepwidestar[4]{x\in (B\cap n)\cup\{n\}}{\rA}
    \proofstepwidestar[4]{x\in B\cap n\lor x\in\{n\}}{\FormulaRefAuto{z\in A\cup B \eqvdash z\in A\lor z\in B}{4}}

    \proofstepwidestar[6]{x\in B\cap n}{\rA}
    \proofstepwidestar[6]{x\in B}{\FormulaRefAuto{x\in A\cap B\vdash x\in A}{6}}

    \proofstepwidestar[8]{x\in\{n\}}{\rA}
    \proofstepwidestar[8]{x=n}{\FormulaRefAuto{x\in\{a\}\eqvdash x=a}{8}}
    \proofstepwidestar[3,9]{x\in B}{\rIE{9,3}}

    \proofstepwidestar[1,2,3,4]{x\in B}{\rOE{5,6,7,8,10}}
    \proofstepwidestar[1,2,3]{\forall x\,\bigl(x\in (B\cap n)\cup\{n\} \rightarrow x\in B\bigr)}{\rUI{\rRI{4,11}}}
    \proofstepwidestar[1,2,3]{(B\cap n)\cup\{n\}\subseteq B}{%
      \FormulaRefAuto{A \subseteq B \coloneq \forall x\,(x\in A \rightarrow x\in B)}{12}}
  \closeproofpartwideindR

  \proofpartwide{Schluss}
    \proofstepwidestar[1]{n\in\mathbb{N}}{\rA}
    \proofstepwidestar[2]{B\subseteq \Succ(n)}{\rA}
    \proofstepwidestar[3]{n\in B}{\rA}

    \proofstepwidestar[1,2,3]{B\subseteq (B\cap n)\cup\{n\}}{%
      \FormulaRefAuto{n\in\mathbb{N}\dsep B\subseteq \Succ(n)\dsep n\in B \vdash B\subseteq (B\cap n)\cup\{n\}}{1,2,3}}
    \proofstepwidestar[1,2,3]{(B\cap n)\cup\{n\}\subseteq B}{%
      \FormulaRefAuto{n\in\mathbb{N}\dsep B\subseteq \Succ(n)\dsep n\in B \vdash (B\cap n)\cup\{n\}\subseteq B}{1,2,3}}
    \proofstepwidestar[1,2,3]{B=(B\cap n)\cup\{n\}}{%
      \FormulaRefAuto{A\subseteq B\dsep B\subseteq A \vdash A=B}{4,5}}
  \closeproofpartwide
\end{tabproofsplitwide}



\subsection{Minimalität der Menge der natürlichen Zahlen}

\FormulaThmDelta[Minimalität]{\Induktiv(A) \vdash \mathbb{N} \subseteq A}{\DeltaDecl{Mengen}{A}}
\begin{tabproof}
  \proofstep{1}{ \Induktiv(A) }{\rA}
  \proofstep{2}{ \mathbb{N} \subseteq A }{\FormulaRefAuto{P(C)\vdash \bigcap_{P(A)} A \subseteq C}{1}}
\end{tabproof}


\FormulaThmDelta{%
  n\in\mathbb{N}\vdash n\subseteq \Succ(n)%
}{
  \DeltaDecl{Mengen}{n}
}
\begin{tabproof}
  \proofstep{1}{n\in\mathbb{N}}{\rA}

  \proofstep{1}{\Succ(n)=n\cup\{n\}}{%
    \FormulaRefAuto{x\in\mathbb{N}\vdash \Succ(x)=x\cup\{x\}}{1}}

  \proofstep{}{\,n\subseteq n\cup\{n\}}{%
    \FormulaRefAuto{A\subseteq A\cup B}}

  \proofstep{1}{n\subseteq \Succ(n)}{\rIE{2,3}}
\end{tabproof}


\section{Herleitung der Peano-Axiome}

Die folgenden Abschnitte zeigen, wie sich die Peano-Axiome aus den zuvor eingeführten Definitionen und Sätzen ergeben.

\subsection{Das Null-Axiom}

Aus \FormulaRefAuto{\varnothing\in\mathbb{N}} folgt, dass die Null in den natürlichen Zahlen enthalten ist.

\subsection{Nachfolger‑Axiom}

\FormulaThmAuto[Nachfolger-Axiom]{%
  \forall x\in \mathbb{N}\,\Succ(x)\in \mathbb{N}%
}
\begin{tabproof}
  \proofstep{1}{x\in \mathbb{N}}{\rA}

  \proofstep{}{\,\Succ\colon \mathbb{N}\to \mathbb{N}}{%
    \FormulaRefAuto{\Succ\colon \mathbb{N}\to \mathbb{N}}}

  \proofstep{1}{\,\Succ(x)\in \mathbb{N}}{%
    \FormulaRefAuto{F\colon A\to B\dsep x\in A\vdash F(x)\in B}{1,2}}

  \proofstep{}{\,\forall x\in \mathbb{N}\,(\Succ(x)\in \mathbb{N})}{%
    \rUI{\rRI{1,3}}}
\end{tabproof}


\subsection{Kein Vorgänger von Null}

\FormulaThmDelta{x\cup\{x\}\neq \varnothing}{
  \DeltaRow{Mengen}{x}
}
\begin{tabproof}
  \proofstep{1}{x\in \{x\}}{%
    \FormulaRefAuto{x\in \{x\}}}

  \proofstep{1}{x\in x\cup \{x\}}{%
    \FormulaRefAuto{u\in B \vdash u\in A\cup B}{1}}

  \proofstep{}{x\cup\{x\}\neq \varnothing}{%
    \FormulaRefAuto{u\in A \vdash A\neq \varnothing}{2}}
\end{tabproof}


\FormulaThmDelta{%
  x\in\mathbb{N}\vdash \Succ(x) \neq \varnothing%
}{
  \DeltaRow{Mengen}{x}
}

\begin{tabproof}
  \proofstep{1}{x\in\mathbb{N}}{\rA}
  \proofstep{}{x\in \{x\}}{%
    \FormulaRefAuto{x\in \{x\}}}
  \proofstep{}{x\in x\cup \{x\}}{%
    \FormulaRefAuto{u\in B \vdash u\in A\cup B}{2}}
  \proofstep{}{\,\Succ(x)=x\cup \{x\}}{%
    \FormulaRefAuto{x\in\mathbb{N}\vdash \Succ(x)=x\cup\{x\}}}
  \proofstep{1}{x\in \Succ(x)}{\rIE{3,2}}
  \proofstep{}{\,\Succ(x)\neq \varnothing}{%
    \FormulaRefAuto{u\in A \vdash A\neq \varnothing}{4}}
\end{tabproof}

\FormulaThmDeltaK[Null liegt im Null-oder-Nachfolger-Träger]{%
  \varnothing\in
  \{z\in\mathbb{N}\mid z=\varnothing \lor
  \exists y\in\mathbb{N}\,(z=\Succ(y))\}%
}{NaturalZeroOrSuccessorCarrierContainsZero}{%
  \DeltaDecl{Mengen}{y\dsep z}%
}
\begin{tabproof}
  \proofstep{}{\varnothing\in\mathbb{N}}{%
    \FormulaRefAuto{\varnothing\in\mathbb{N}}}
  \proofstep{}{\varnothing=\varnothing}{\rII}
  \proofstep{}{%
    \varnothing=\varnothing \lor
    \exists y\in\mathbb{N}\,(\varnothing=\Succ(y))%
  }{\rOIa{2}}
  \proofstep{}{%
    \varnothing\in\mathbb{N}\land
    \bigl(\varnothing=\varnothing \lor
    \exists y\in\mathbb{N}\,(\varnothing=\Succ(y))\bigr)%
  }{\rAI{1,3}}
  \proofstep{}{%
    \varnothing\in
    \{z\in\mathbb{N}\mid z=\varnothing \lor
    \exists y\in\mathbb{N}\,(z=\Succ(y))\}%
  }{%
    \FormulaRefAuto{x \in \{u \in A \mid P(u)\} \eqvdash
      x \in A \land P(x)}{4}}
\end{tabproof}

\FormulaThmDeltaK[Nachfolgerabschluss des Null-oder-Nachfolger-Trägers]{%
  \begin{aligned}[t]
    &z\in\{u\in\mathbb{N}\mid u=\varnothing \lor
      \exists y\in\mathbb{N}\,(u=\Succ(y))\}\vdash{}\\
    &z\cup\{z\}\in\{u\in\mathbb{N}\mid u=\varnothing \lor
      \exists y\in\mathbb{N}\,(u=\Succ(y))\}
  \end{aligned}%
}{NaturalZeroOrSuccessorCarrierSuccessorClosed}{%
  \DeltaDecl{Mengen}{y\dsep z}%
}
\begin{tabproof}
  \proofstep{1}{%
    z\in\{u\in\mathbb{N}\mid u=\varnothing \lor
    \exists y\in\mathbb{N}\,(u=\Succ(y))\}%
  }{\rA}
  \proofstep{1}{z\in\mathbb{N}}{%
    \FormulaRefAuto{x \in \{x \in A \mid P(x)\}\vdash x\in A}{1}}
  \proofstep{1}{\Succ(z)=z\cup\{z\}}{%
    \FormulaRefAuto{x\in\mathbb{N}\vdash \Succ(x)=x\cup\{x\}}{2}}
  \proofstep{1}{z\cup\{z\}=\Succ(z)}{%
    \FormulaRefAuto{a=b\vdash b=a}{3}}
  \proofstep{1}{\Succ(z)\in\mathbb{N}}{%
    \FormulaRefAuto{n\in\mathbb{N}\vdash \Succ(n)\in\mathbb{N}}{2}}
  \proofstep{1}{z\cup\{z\}\in\mathbb{N}}{%
    \FormulaRefAuto{a=b \dsep b\in A \vdash a\in A}{4,5}}
  \proofstep{1}{z\in\mathbb{N}\land z\cup\{z\}=\Succ(z)}{\rAI{2,4}}
  \proofstep{1}{\exists y\in\mathbb{N}\,(z\cup\{z\}=\Succ(y))}{\rEI{7}}
  \proofstep{1}{%
    z\cup\{z\}=\varnothing \lor
    \exists y\in\mathbb{N}\,(z\cup\{z\}=\Succ(y))%
  }{\rOIb{8}}
  \proofstep{1}{%
    z\cup\{z\}\in\mathbb{N}\land
    \bigl(z\cup\{z\}=\varnothing \lor
    \exists y\in\mathbb{N}\,(z\cup\{z\}=\Succ(y))\bigr)%
  }{\rAI{6,9}}
  \proofstep{1}{%
    z\cup\{z\}\in\{u\in\mathbb{N}\mid u=\varnothing \lor
    \exists y\in\mathbb{N}\,(u=\Succ(y))\}%
  }{%
    \FormulaRefAuto{x \in \{u \in A \mid P(u)\} \eqvdash
      x \in A \land P(x)}{10}}
\end{tabproof}

\FormulaThmDeltaK[Der Null-oder-Nachfolger-Träger ist induktiv]{%
  \Induktiv\bigl(\{z\in\mathbb{N}\mid z=\varnothing \lor
  \exists y\in\mathbb{N}\,(z=\Succ(y))\}\bigr)%
}{NaturalZeroOrSuccessorCarrierInductive}{%
  \DeltaDecl{Mengen}{y\dsep z}%
}
\begin{tabproof}
  \proofstep{}{%
    \varnothing\in\{z\in\mathbb{N}\mid z=\varnothing \lor
    \exists y\in\mathbb{N}\,(z=\Succ(y))\}%
  }{\FormulaRefAuto{NaturalZeroOrSuccessorCarrierContainsZero}}
  \proofstep{2}{%
    z\in\{u\in\mathbb{N}\mid u=\varnothing \lor
    \exists y\in\mathbb{N}\,(u=\Succ(y))\}%
  }{\rA}
  \proofstep{2}{%
    z\cup\{z\}\in\{u\in\mathbb{N}\mid u=\varnothing \lor
    \exists y\in\mathbb{N}\,(u=\Succ(y))\}%
  }{\FormulaRefAuto{NaturalZeroOrSuccessorCarrierSuccessorClosed}{2}}
  \proofstep{}{%
    \forall z\in\{u\in\mathbb{N}\mid u=\varnothing \lor
    \exists y\in\mathbb{N}\,(u=\Succ(y))\}\,
    \bigl(z\cup\{z\}\in\{u\in\mathbb{N}\mid u=\varnothing \lor
    \exists y\in\mathbb{N}\,(u=\Succ(y))\}\bigr)%
  }{\rUI{\rRI{2,3}}}
  \proofstep{}{%
    \varnothing\in\{z\in\mathbb{N}\mid z=\varnothing \lor
    \exists y\in\mathbb{N}\,(z=\Succ(y))\}\land{}\allowbreak
    \forall z\in\{u\in\mathbb{N}\mid u=\varnothing \lor
    \exists y\in\mathbb{N}\,(u=\Succ(y))\}\,
    \bigl(z\cup\{z\}\in\{u\in\mathbb{N}\mid u=\varnothing \lor
    \exists y\in\mathbb{N}\,(u=\Succ(y))\}\bigr)%
  }{\rAI{1,4}}
  \proofstep{}{%
    \Induktiv\bigl(\{z\in\mathbb{N}\mid z=\varnothing \lor
    \exists y\in\mathbb{N}\,(z=\Succ(y))\}\bigr)%
  }{%
    \rIE{\FormulaRefAuto{\Induktiv(A) \coloneq \varnothing \in A \,\land\,
      \forall x\in A\,(x \cup \{x\} \in A)},5}}
\end{tabproof}

\FormulaThmDeltaR[Jede von \(\varnothing\) verschiedene natürliche Zahl ist ein Nachfolger]{%
  \begin{aligned}[t]
    &x\in\mathbb{N}\dsep x\neq \varnothing \vdash{}\\
    &\exists y\in\mathbb{N}\,\bigl(x=\Succ(y)\bigr)
  \end{aligned}
}{%
  x\in\mathbb{N}\dsep x\neq \varnothing
  \vdash \exists y\in\mathbb{N}\,\bigl(x=\Succ(y)\bigr)
}{%
  \DeltaDecl{Mengen}{x\dsep y}%
}
\begin{tabproof}
  \proofstep{1}{x\in\mathbb{N}}{\rA}
  \proofstep{2}{x\neq \varnothing}{\rA}
  \proofstep{}{%
    \Induktiv\bigl(\{z\in\mathbb{N}\mid z=\varnothing \lor
    \exists y\in\mathbb{N}\,(z=\Succ(y))\}\bigr)%
  }{\FormulaRefAuto{NaturalZeroOrSuccessorCarrierInductive}}
  \proofstep{}{%
    \mathbb{N}\subseteq
    \{z\in\mathbb{N}\mid z=\varnothing \lor
    \exists y\in\mathbb{N}\,(z=\Succ(y))\}%
  }{\FormulaRefAuto{\Induktiv(A)\vdash \mathbb{N} \subseteq A}{3}}
  \proofstep{1}{%
    x\in\{z\in\mathbb{N}\mid z=\varnothing \lor
    \exists y\in\mathbb{N}\,(z=\Succ(y))\}%
  }{\FormulaRefAuto{A\subseteq B\dsep x\in A\vdash x\in B}{4,1}}
  \proofstep{1}{%
    x=\varnothing \lor \exists y\in\mathbb{N}\,\bigl(x=\Succ(y)\bigr)
  }{\FormulaRefAuto{x \in \{x \in A \mid P(x)\}\vdash P(x)}{5}}
  \proofstep{1,2}{%
    \exists y\in\mathbb{N}\,\bigl(x=\Succ(y)\bigr)
  }{%
    \FormulaRefAuto{P \lor Q,\neg P \vdash Q}{6,2}}
\end{tabproof}

\subsection{Injektivität der Nachfolgerfunktion}

\FormulaThmAuto{%
  x\in\mathbb{N}\dsep y\in\mathbb{N}\dsep \Succ(x)=\Succ(y)\vdash x=y%
}
\begin{tabproof}
  \proofstep{1}{x\in\mathbb{N}}{\rA}
  \proofstep{2}{y\in\mathbb{N}}{\rA}
  \proofstep{3}{\Succ(x)=\Succ(y)}{\rA}

  \proofstep{1}{\Succ(x)=x\cup\{x\}}{%
    \FormulaRefAuto{x\in\mathbb{N}\vdash \Succ(x)=x\cup\{x\}}{1}}
  \proofstep{2}{\Succ(y)=y\cup\{y\}}{%
    \FormulaRefAuto{x\in\mathbb{N}\vdash \Succ(x)=x\cup\{x\}}{2}}

  \proofstep{1}{x\cup\{x\}=\Succ(x)}{%
    \FormulaRefAuto{a=b\vdash b=a}{4}}
  \proofstep{1,3}{x\cup\{x\}=\Succ(y)}{%
    \FormulaRefAuto{a=b,\,b=c\vdash a=c}{6,3}}
  \proofstep{1,2,3}{x\cup\{x\}=y\cup\{y\}}{%
    \FormulaRefAuto{a=b,\,b=c\vdash a=c}{7,5}}

  \proofstep{1,2,3}{x=y}{%
    \FormulaRefAuto{a\cup\{a\}=b\cup\{b\}\eqvdash a=b}{8}}
\end{tabproof}


\subsection{Induktionsprinzip}

\FormulaThmDeltaK[Nullaufnahme der Prädikataussonderung]{%
  P(\varnothing)\vdash
  \varnothing\in\{x\in\mathbb{N}\mid P(x)\}%
}{NaturalPredicateCarrierContainsZero}{%
  \DeltaRow{Einstelliges Prädikat}{P}%
}
\begin{tabproof}
  \proofstep{1}{P(\varnothing)}{\rA}
  \proofstep{}{\varnothing\in\mathbb{N}}{%
    \FormulaRefAuto{\varnothing\in\mathbb{N}}}
  \proofstep{1}{\varnothing\in\mathbb{N}\land P(\varnothing)}{\rAI{2,1}}
  \proofstep{1}{\varnothing\in\{x\in\mathbb{N}\mid P(x)\}}{%
    \FormulaRefAuto{x \in \{u \in A \mid P(u)\} \eqvdash
      x \in A \land P(x)}{3}}
\end{tabproof}

\FormulaThmDeltaK[Nachfolgerabschluss der Prädikataussonderung]{%
  \begin{aligned}[t]
    &\forall x\in\mathbb{N}\,(P(x)\rightarrow P(\Succ(x)))\vdash{}\\
    &\forall x\in\{u\in\mathbb{N}\mid P(u)\}\,
      \bigl(x\cup\{x\}\in\{u\in\mathbb{N}\mid P(u)\}\bigr)
  \end{aligned}%
}{NaturalPredicateCarrierSuccessorClosed}{%
  \DeltaRow{Einstelliges Prädikat}{P}%
  \DeltaRow{Mengen}{x}%
}
\begin{tabproof}
  \proofstep{1}{\forall x\in\mathbb{N}\,(P(x)\rightarrow P(\Succ(x)))}{\rA}
  \proofstep{2}{x\in\{u\in\mathbb{N}\mid P(u)\}}{\rA}
  \proofstep{2}{x\in\mathbb{N}}{%
    \FormulaRefAuto{x \in \{x \in A \mid P(x)\}\vdash x\in A}{2}}
  \proofstep{2}{P(x)}{%
    \FormulaRefAuto{x \in \{x \in A \mid P(x)\}\vdash P(x)}{2}}
  \proofstep{1,2}{P(x)\rightarrow P(\Succ(x))}{\rRE{\rUE{1},3}}
  \proofstep{1,2}{P(\Succ(x))}{\rRE{5,4}}
  \proofstep{2}{\Succ(x)=x\cup\{x\}}{%
    \FormulaRefAuto{x\in\mathbb{N}\vdash \Succ(x)=x\cup\{x\}}{3}}
  \proofstep{1,2}{P(x\cup\{x\})}{\rIE{7,6}}
  \proofstep{2}{\Succ(x)\in\mathbb{N}}{%
    \FormulaRefAuto{n\in\mathbb{N}\vdash \Succ(n)\in\mathbb{N}}{3}}
  \proofstep{2}{x\cup\{x\}\in\mathbb{N}}{\rIE{7,9}}
  \proofstep{1,2}{x\cup\{x\}\in\mathbb{N}\land P(x\cup\{x\})}{\rAI{10,8}}
  \proofstep{1,2}{x\cup\{x\}\in\{u\in\mathbb{N}\mid P(u)\}}{%
    \FormulaRefAuto{x \in \{u \in A \mid P(u)\} \eqvdash
      x \in A \land P(x)}{11}}
  \proofstep{1}{%
    x\in\{u\in\mathbb{N}\mid P(u)\}\rightarrow
    x\cup\{x\}\in\{u\in\mathbb{N}\mid P(u)\}%
  }{\rRI{2,12}}
  \proofstep{1}{%
    \forall x\in\{u\in\mathbb{N}\mid P(u)\}\,
    \bigl(x\cup\{x\}\in\{u\in\mathbb{N}\mid P(u)\}\bigr)%
  }{\rUI{13}}
\end{tabproof}

\FormulaThmAuto{%
  P(\varnothing),\, \forall x\in\mathbb{N}\,(P(x) \rightarrow P(\Succ(x)))
  \vdash \Induktiv(\{x \in \mathbb{N} \mid P(x)\})%
}
\begin{tabproof}
  \proofstep{1}{P(\varnothing)}{\rA}
  \proofstep{2}{\forall x\in\mathbb{N}\,(P(x)\rightarrow P(\Succ(x)))}{\rA}
  \proofstep{1}{\varnothing\in\{x\in\mathbb{N}\mid P(x)\}}{%
    \FormulaRefAuto{NaturalPredicateCarrierContainsZero}{1}}
  \proofstep{2}{%
    \forall x\in\{u\in\mathbb{N}\mid P(u)\}\,
    \bigl(x\cup\{x\}\in\{u\in\mathbb{N}\mid P(u)\}\bigr)%
  }{\FormulaRefAuto{NaturalPredicateCarrierSuccessorClosed}{2}}
  \proofstep{1,2}{%
    \begin{aligned}[t]
      \varnothing\in \{x\in\mathbb{N}\mid P(x)\}\land {}\\
      \forall x\in \{x\in\mathbb{N}\mid P(x)\}\,(x\cup\{x\}\in \{x\in\mathbb{N}\mid P(x)\})
    \end{aligned}%
  }{%
    \rAI{3,4}}

  \proofstep{1,2}{\Induktiv(\{x \in \mathbb{N} \mid P(x)\})}{%
    \rIE{\FormulaRefAuto{\Induktiv(A) \coloneq \varnothing \in A \,\land\, \forall x\in A\,(x \cup \{x\} \in A)},5}}
\end{tabproof}


\FormulaThmDelta[Induktion auf \(\mathbb{N}\)]{P(\varnothing)\dsep \forall x\in\mathbb{N}\,(P(x) \rightarrow P(\Succ(x))) \vdash \forall x\in\mathbb{N}\,(P(x))}{
  \DeltaRow{Einstelliges Prädikat}{P}
}
\begin{tabproof}
  \proofstep{1}{ P(\varnothing) }{\rA}
  \proofstep{2}{ \forall x \in \mathbb{N}\,(P(x) \rightarrow P(\Succ(x))) }{\rA}
  \proofstep{1,2}{ \Induktiv(\{x \in \mathbb{N} \mid P(x)\}) }{\FormulaRefAuto{P(\varnothing),\, \forall x\in\mathbb{N}\,(P(x) \rightarrow P(\Succ(x))) \vdash \Induktiv(\{x \in \mathbb{N} \mid P(x)\})}{1,2}}
  \proofstep{}{ \mathbb{N} \subseteq \{x \in \mathbb{N} \mid P(x)\} }{\FormulaRefAuto{\Induktiv(A) \vdash \mathbb{N} \subseteq A}{3}}
  \proofstep{1,2}{ \forall x\in\mathbb{N}\,(P(x)) }{\FormulaRefAuto{A\subseteq \{x\in B\mid P(x)\}\vdash \forall x\in A(P(x))}{4}}
\end{tabproof}

\FormulaThmDelta[Induktion auf \(\mathbb{N}\)]{%
  P(\varnothing)\dsep [x\in\mathbb{N},P(x)]\vdots P(\Succ(x))\dsep n\in\mathbb{N} \vdash P(n)%
}{
  \DeltaRow{Einstelliges Prädikat}{P}
}
\begin{tabproof}
  \proofstep{1}{P(\varnothing)}{\rA}
  \proofstep{2}{n\in\mathbb{N}}{\rA}

  \proofstep{3}{x\in\mathbb{N}}{\rA}
  \proofstep{4}{P(x)}{\rA}
  \proofstep{3,4}{P(\Succ(x))}{[x\in\mathbb{N},P(x)]\vdots P(\Succ(x))(3,4)}
  \proofstep{3}{P(x)\rightarrow P(\Succ(x))}{\rRI{4,5}}
  \proofstep{}{x\in\mathbb{N}\rightarrow \bigl(P(x)\rightarrow P(\Succ(x))\bigr)}{\rRI{3,6}}
  \proofstep{}{ \forall x\in\mathbb{N}\,\bigl(P(x)\rightarrow P(\Succ(x))\bigr) }{\rUI{7}}

  \proofstep{1}{\forall x\in\mathbb{N}\,(P(x))}{%
    \FormulaRefAuto{P(\varnothing)\dsep \forall x\in\mathbb{N}\,(P(x) \rightarrow P(\Succ(x))) \vdash \forall x\in\mathbb{N}\,(P(x))}{1,8}%
  }
  \proofstep{}{n\in\mathbb{N}\rightarrow P(n)}{\rUE{9}}
  \proofstep{2,9}{P(n)}{\rRE{2,10}}
\end{tabproof}


% Endlichkeitsunabhängiger Funktionenblock. Insbesondere wird hier der
% von-Neumann-Satz bewiesen, dass keine natürliche Zahl eine injektive Kopie
% von \(\mathbb{N}\) enthält. Band~08 verwendet diesen Satz in der
% Dedekind-Charakterisierung.
% Die folgenden weiterführenden Resultate über natürliche Zahlen werden in
% diesem Band in der dort verwendeten Peano-Sprache entwickelt. Die historische
% B03-Fassung bleibt als inaktive Arbeitsfassung im Quelltext erhalten.

\section{Übergang zur Peano-Schreibweise}

In der mengentheoretischen Konstruktion wurde die erste natürliche Zahl als
leere Menge geschrieben. Für die anschließende abstrakte Peano-Theorie führen
wir nun das übliche Zahlzeichen ein.

\FormulaDefAuto[Nullzeichen]{0\coloneqq\varnothing}

\chapter{Peano-Axiome}

Die Struktur \((\mathbb{N},0,\Succ)\) wird durch die folgenden Axiome
festgelegt. Zuerst stehen sie gesammelt wie die ZF-Axiome in Band~03; danach
werden sie einzeln registriert, damit spätere Beweise dieses Bandes direkt
auf diese Axiome verweisen können.

\FormulaDefDeltaK[Peano-Axiome mit Nachfolgerfunktion]{%
  (\mathbb{N},0,\Succ)%
}{Peano-Axiome}{%
  \DeltaRow{\textbf{Axiome}}{}
  \DeltaRow{Null}
           {0\in\mathbb{N}}
           [\FormulaRefAuto{PeanoZero}]
  \DeltaRow{\makecell[l]{Nachfolger-\\funktion}}
           {\Succ\colon \mathbb{N}\to \mathbb{N}}
           [\FormulaRefAuto{PeanoSuccFunction}]
  \DeltaRow{\makecell[l]{Nachfolger-\\Axiom}}
           {\begin{aligned}[t]
              &\forall x\in \mathbb{N}\\
              &\quad \Succ(x)\in \mathbb{N}
            \end{aligned}}
           [\FormulaRefAuto{PeanoSuccClosureUniversal}]
  \DeltaRow{\makecell[l]{Nachfolger-\\abschluss}}
           {n\in\mathbb{N}\vdash \Succ(n)\in\mathbb{N}}
           [\FormulaRefAuto{PeanoSuccClosure}]
  \DeltaRow{\makecell[l]{Kein Nachfolger\\ ist Null}}
           {x\in\mathbb{N}\vdash \Succ(x) \neq 0}
           [\FormulaRefAuto{PeanoSuccNonzero}]
  \DeltaRow{\makecell[l]{Nichtnullzahlen\\ sind Nachfolger}}
           {\begin{aligned}[t]
              &x\in\mathbb{N}\dsep x\neq 0 \vdash{}\\
              &\exists y\in\mathbb{N}\,(x=\Succ(y))
            \end{aligned}}
           [\FormulaRefAuto{PeanoNonzeroIsSucc}]
  \DeltaRow{Injektivität}
           {\begin{aligned}[t]
              &x\in\mathbb{N}\dsep y\in\mathbb{N}\dsep{}\\
              &\Succ(x)=\Succ(y)\vdash x=y
            \end{aligned}}
           [\FormulaRefAuto{PeanoSuccInjective}]
  \DeltaRow{Induktion}
           {\begin{aligned}[t]
              &P(0)\dsep \forall x\in\mathbb{N}\\
              &\quad (P(x)\rightarrow P(\Succ(x)))\vdash{}\\
              &\forall x\in\mathbb{N}\,(P(x))
            \end{aligned}}
           [\FormulaRefAuto{PeanoInductionAxiom}]
  \DeltaRow{\makecell[l]{Induktions-\\regel}}
           {\begin{aligned}[t]
              &P(0)\dsep [x\in\mathbb{N},P(x)]\vdots{}\\
              &P(\Succ(x))\dsep n\in\mathbb{N}\vdash P(n)
            \end{aligned}}
           [\FormulaRefAuto{PeanoInductionRule}]
  \DeltaRow{\textbf{Neue Symbole}}{}
  \DeltaRow{\makecell[l]{Natürliche\\ Zahlen}}{\mathbb{N}}
  \DeltaRow{Nullzeichen}{0}[\FormulaRefAuto{0\coloneqq\varnothing}]
  \DeltaRow{\makecell[l]{Nachfolger-\\funktion}}{\Succ}
}

\FormulaAxiomAutoR[Null-Axiom]{0\in\mathbb{N}}{PeanoZero}
\par\noindent\emph{Herkunft aus Band~03:} Dies ist die \(0\)-Schreibweise von
\FormulaRefAuto{\varnothing\in\mathbb{N}}[thm].\par

\FormulaAxiomAutoR[Nachfolgerfunktion]{%
  \Succ\colon \mathbb{N}\to \mathbb{N}%
}{PeanoSuccFunction}
\par\noindent\emph{Herkunft aus Band~03:} Dies entspricht
\FormulaRefAuto{\Succ\colon \mathbb{N}\to \mathbb{N}}[thm].\par

\FormulaAxiomAutoR[Nachfolger-Axiom]{%
  \forall x\in \mathbb{N}\,\Succ(x)\in \mathbb{N}%
}{PeanoSuccClosureUniversal}
\par\noindent\emph{Herkunft aus Band~03:} Dies entspricht
\FormulaRefAuto{\forall x\in \mathbb{N}\,\Succ(x)\in \mathbb{N}}[thm].\par

\FormulaAxiomAutoR[Nachfolgerabschluss]{n\in\mathbb{N}\vdash \Succ(n)\in\mathbb{N}}{PeanoSuccClosure}
\par\noindent\emph{Herkunft aus Band~03:} Dies entspricht
\FormulaRefAuto{n\in\mathbb{N}\vdash \Succ(n)\in\mathbb{N}}[thm].\par

\FormulaAxiomAutoR[Kein Nachfolger ist Null]{x\in\mathbb{N}\vdash \Succ(x) \neq 0}{PeanoSuccNonzero}
\par\noindent\emph{Herkunft aus Band~03:} Dies ist die \(0\)-Schreibweise von
\FormulaRefAuto{x\in\mathbb{N}\vdash \Succ(x) \neq \varnothing}[thm].\par

\FormulaAxiomAutoR[Jede von \(0\) verschiedene natürliche Zahl ist ein Nachfolger]{%
  x\in\mathbb{N}\dsep x\neq 0 \vdash \exists y\in\mathbb{N}\,(x=\Succ(y))%
}{PeanoNonzeroIsSucc}
\par\noindent\emph{Herkunft aus Band~03:} Dies ist die \(0\)-Schreibweise von
\FormulaRefAuto{x\in\mathbb{N}\dsep x\neq \varnothing
  \vdash \exists y\in\mathbb{N}\,\bigl(x=\Succ(y)\bigr)}[thm].\par

\FormulaAxiomAutoR[Injektivität der Nachfolgerfunktion]{%
  x\in\mathbb{N}\dsep y\in\mathbb{N}\dsep \Succ(x)=\Succ(y)\vdash x=y%
}{PeanoSuccInjective}
\par\noindent\emph{Herkunft aus Band~03:} Dies entspricht
\FormulaRefAuto{x\in\mathbb{N}\dsep y\in\mathbb{N}\dsep \Succ(x)=\Succ(y)\vdash x=y}[thm].\par

\FormulaAxiomAutoR[Induktionsaxiom]{%
  P(0)\dsep \forall x\in\mathbb{N}\,(P(x) \rightarrow P(\Succ(x)))
  \vdash \forall x\in\mathbb{N}\,(P(x))%
}{PeanoInductionAxiom}
\par\noindent\emph{Herkunft aus Band~03:} Dies ist die \(0\)-Schreibweise von
\FormulaRefAuto{P(\varnothing)\dsep \forall x\in\mathbb{N}\,(P(x) \rightarrow P(\Succ(x)))
  \vdash \forall x\in\mathbb{N}\,(P(x))}[thm].\par

\FormulaAxiomAutoR[Induktionsregel]{%
  P(0)\dsep [x\in\mathbb{N},P(x)]\vdots P(\Succ(x))\dsep n\in\mathbb{N} \vdash P(n)%
}{PeanoInductionRule}
\par\noindent\emph{Herkunft aus Band~03:} Dies ist die \(0\)-Schreibweise von
\FormulaRefAuto{P(\varnothing)\dsep [x\in\mathbb{N},P(x)]\vdots P(\Succ(x))\dsep n\in\mathbb{N} \vdash P(n)}[thm].\par

\chapter{Natürliche Zahlen aus den Peano-Axiomen}

\section{Grundlegende Peano-Eigenschaften der natürlichen Zahlen}

In diesem Band wird \(\mathbb{N}\) in den folgenden Abschnitten als
Peano-Struktur mit den primitiven Zeichen \(0\) und \(\Succ\) verwendet. Die
von-Neumann-spezifischen Aussagen aus Band~03, insbesondere
\(\Succ(n)=n\cup\{n\}\), \(x\in\Succ(n)\leftrightarrow x\in n\lor x=n\) und
Ordnungsdefinitionen über \(m\in n\) oder \(m\subseteq n\), werden hier nicht
als Folgerungen aus den Peano-Axiomen benutzt.

\iffalse
Von-Neumann-spezifischer Block aus Band 03. Diese Aussagen beruhen auf der
konkreten Konstruktion natürlicher Zahlen als Mengen und gehören nicht zur
abstrakten Entwicklung aus den Peano-Axiomen.
\section{Eigenschaften der Menge der natürlichen Zahlen}

\subsection{Die Nachfolgerfunktion}

\FormulaDefAuto[Nachfolgerfunktion]{%
  \Succ \coloneqq \Succ_{\mathbb{N}}%
}
\par\noindent\emph{Herkunft aus Band~03:} Die B03-Herleitung der Nachfolgerfunktion ist
\FormulaRefAuto{\Succ\colon \mathbb{N}\to \mathbb{N}}[thm]; die Schreibweise
\(\Succ=\Succ_{\mathbb{N}}\) entspricht der dortigen Konstruktion.\par

\FormulaThmDelta{%
  x\in\mathbb{N}\vdash \Succ(x)=x\cup\{x\}%
}{
  \DeltaDecl{Mengen}{x}
}
\begin{tabproof}
  \proofstep{1}{x\in\mathbb{N}}{\rA}
  \proofstep{}{\,\Induktiv(\mathbb{N})}{%
    \FormulaRefAuto{\Induktiv(\mathbb{N})}}
  \proofstep{1}{\,\Succ_{\mathbb{N}}(x)=x\cup\{x\}}{%
    \FormulaRefAuto{\Induktiv(A)\dsep x\in A \vdash \Succ_A(x)=x\cup\{x\}}{2,1}}
  \proofstep{}{\,\Succ=\Succ_{\mathbb{N}}}{%
    \FormulaRefAuto{\Succ \coloneqq \Succ_{\mathbb{N}}}}
  \proofstep{1}{\,\Succ(x)=x\cup\{x\}}{\rIE{4,3}}
\end{tabproof}

\subsubsection{Mengen-theoretische Eigenschaften der Nachfolgermenge}

\paragraph{Zerlegung der Nachfolgermenge}


\FormulaThmDelta{%
  n\in\mathbb{N}\dsep x\in\Succ(n)\vdash x\in n\lor x=n%
}{%
  \DeltaDecl{Mengen}{n\dsep x}%
}
\begin{tabproof}
  \proofstep{1}{n\in\mathbb{N}}{\rA}
  \proofstep{2}{x\in\Succ(n)}{\rA}
  \proofstep{1}{\Succ(n)=n\cup\{n\}}{%
    \FormulaRefAuto{n\in\mathbb{N}\vdash \Succ(n) = n\cup\{n\}}{1}}
  \proofstep{1,2}{x\in n\cup\{n\}}{\rIE{3,2}}
  \proofstep{1,2}{x\in n\lor x\in\{n\}}{%
    \FormulaRefAuto{z \in A \cup B \eqvdash z \in A \lor z \in B}{2}}

  \proofstep{3}{x\in n}{\rA}
  \proofstep{3}{x\in n\lor x=n}{\rOIa{6}}

  \proofstep{4}{x\in\{n\}}{\rA}
  \proofstep{4}{x=n}{\FormulaRefAuto{x \in \{a\} \eqvdash x = a}{4}}
  \proofstep{4}{x\in n\lor x=n}{\rOIb{9}}

  \proofstep{1,2}{x\in n\lor x=n}{\rOE{5,6,7,8,10}}
\end{tabproof}

% ------------------------------------------------------------
% Hilfssatz: Element in Succ(n) und ungleich n liegt in n
% ------------------------------------------------------------
\FormulaThmDelta{%
  n\in\mathbb{N}\dsep z\in\Succ(n)\dsep z\neq n \vdash z\in n%
}{%
  \DeltaDecl{Mengen}{n\dsep z}%
}
\begin{tabproof}
  \proofstep{1}{n\in\mathbb{N}}{\rA}
  \proofstep{2}{z\in\Succ(n)}{\rA}
  \proofstep{3}{z\neq n}{\rA}
  \proofstep{1,2}{z\in n\lor z=n}{%
    \FormulaRefAuto{n\in\mathbb{N}\dsep x\in \Succ (n)\vdash x\in n\lor x=n}{1,2}}
  \proofstep{3,4}{z\in n}{%
    \FormulaRefAuto{P \lor Q,\neg Q \vdash P}{4,3}}
\end{tabproof}


\FormulaThmDelta[Teilmengen einer Nachfolgermenge ohne neues Element]{%
  n\in\mathbb{N}\dsep B\subseteq \Succ(n)\dsep n\notin B \vdash B\subseteq n
}{%
  \DeltaRow{Mengen}{B\dsep n\dsep x}
}
\begin{tabproof}
  \proofstep{1}{n\in\mathbb{N}}{\rA}
  \proofstep{2}{B\subseteq \Succ(n)}{\rA}
  \proofstep{3}{n\notin B}{\rA}

  \proofstep{4}{x\in B}{\rA}
  \proofstep{2,4}{x\in \Succ(n)}{\FormulaRefAuto{A\subseteq B,\, x\in A \vdash x\in B}{2,4}}
  \proofstep{1}{\Succ(n)=n\cup\{n\}}{\FormulaRefAuto{n\in\mathbb{N}\vdash \Succ(n)=n\cup\{n\}}{1}}
  \proofstep{1,2,4}{x\in n\cup\{n\}}{\FormulaRefAuto{A = B,\, x \in A \vdash x \in B}{6,5}}
  \proofstep{1,2,4}{x\in n\lor x=n}{\FormulaRefAuto{x\in A\cup\{a\} \vdash x\in A\lor x=a}{7}}
  \proofstep{3,4}{x\neq n}{\FormulaRefAuto{a \in A,\; b \not\in A \vdash a \neq b}{4,3}}

  \proofstep{10}{x\in n}{\rA}

  \proofstep{11}{x=n}{\rA}
  \proofstep{3,4,11}{\bot}{\rBI{11,9}}
  \proofstep{3,4,11}{x\in n}{\rCE{12}}

  \proofstep{1,2,3,4}{x\in n}{\rOE{8,10,10,11,13}}
  \proofstep{1,2,3}{\forall x\,\bigl(x\in B \rightarrow x\in n\bigr)}{\rUI{\rRI{4,14}}}
  \proofstep{1,2,3}{B\subseteq n}{%
    \FormulaRefAuto{A \subseteq B \coloneq \forall x\,(x\in A \rightarrow x\in B)}{15}}
\end{tabproof}




\FormulaThmDelta[Teilmengen einer Nachfolgermenge mit neuem Element]{%
  n\in\mathbb{N}\dsep B\subseteq \Succ(n)\dsep n\in B \vdash B=(B\cap n)\cup\{n\}
}{%
  \DeltaRow{Mengen}{B\dsep n\dsep x}
}
\begin{tabproofsplitwide}
  \proofpartwideindR[Hinrichtung]{%
    \begin{aligned}[t]
      &n\in\mathbb{N}\dsep B\subseteq \Succ(n)\dsep n\in B \vdash{}\\
      &B\subseteq (B\cap n)\cup\{n\}
    \end{aligned}
  }{%
    n\in\mathbb{N}\dsep B\subseteq \Succ(n)\dsep n\in B \vdash B\subseteq (B\cap n)\cup\{n\}
  }
    \proofstepwidestar[1]{n\in\mathbb{N}}{\rA}
    \proofstepwidestar[2]{B\subseteq \Succ(n)}{\rA}
    \proofstepwidestar[3]{n\in B}{\rA}

    \proofstepwidestar[4]{x\in B}{\rA}
    \proofstepwidestar[2,4]{x\in \Succ(n)}{\FormulaRefAuto{A\subseteq B\dsep x\in A\vdash x\in B}{2,4}}
    \proofstepwidestar[1]{\Succ(n)=n\cup\{n\}}{\FormulaRefAuto{n\in\mathbb{N}\vdash \Succ(n)=n\cup\{n\}}{1}}
    \proofstepwidestar[1,2,4]{x\in n\cup\{n\}}{\rIE{5,6}}
    \proofstepwidestar[1,2,4]{x\in n\lor x\in\{n\}}{\FormulaRefAuto{z\in A\cup B \eqvdash z\in A\lor z\in B}{7}}

    \proofstepwidestar[9]{x\in n}{\rA}
    \proofstepwidestar[4,9]{x\in B\cap n}{\FormulaRefAuto{x\in A\dsep x\in B\vdash x\in A\cap B}{4,9}}
    \proofstepwidestar[4,9]{x\in (B\cap n)\cup\{n\}}{\FormulaRefAuto{z\in A\vdash z\in A\cup B}{10}}

    \proofstepwidestar[12]{x\in\{n\}}{\rA}
    \proofstepwidestar[12]{x\in (B\cap n)\cup\{n\}}{\FormulaRefAuto{z\in B\vdash z\in A\cup B}{12}}

    \proofstepwidestar[1,2,4]{x\in (B\cap n)\cup\{n\}}{\rOE{8,9,11,12,13}}
    \proofstepwidestar[1,2,3]{\forall x\,\bigl(x\in B \rightarrow x\in (B\cap n)\cup\{n\}\bigr)}{\rUI{\rRI{4,14}}}
    \proofstepwidestar[1,2,3]{B\subseteq (B\cap n)\cup\{n\}}{%
      \FormulaRefAuto{A \subseteq B \coloneq \forall x\,(x\in A \rightarrow x\in B)}{15}}
  \closeproofpartwideindR

  \proofpartwideindR[Rückrichtung]{%
    \begin{aligned}[t]
      &n\in\mathbb{N}\dsep B\subseteq \Succ(n)\dsep n\in B \vdash{}\\
      &(B\cap n)\cup\{n\}\subseteq B
    \end{aligned}
  }{%
    n\in\mathbb{N}\dsep B\subseteq \Succ(n)\dsep n\in B \vdash (B\cap n)\cup\{n\}\subseteq B
  }
    \proofstepwidestar[1]{n\in\mathbb{N}}{\rA}
    \proofstepwidestar[2]{B\subseteq \Succ(n)}{\rA}
    \proofstepwidestar[3]{n\in B}{\rA}

    \proofstepwidestar[4]{x\in (B\cap n)\cup\{n\}}{\rA}
    \proofstepwidestar[4]{x\in B\cap n\lor x\in\{n\}}{\FormulaRefAuto{z\in A\cup B \eqvdash z\in A\lor z\in B}{4}}

    \proofstepwidestar[6]{x\in B\cap n}{\rA}
    \proofstepwidestar[6]{x\in B}{\FormulaRefAuto{x\in A\cap B\vdash x\in A}{6}}

    \proofstepwidestar[8]{x\in\{n\}}{\rA}
    \proofstepwidestar[8]{x=n}{\FormulaRefAuto{x\in\{a\}\eqvdash x=a}{8}}
    \proofstepwidestar[3,9]{x\in B}{\rIE{9,3}}

    \proofstepwidestar[1,2,3,4]{x\in B}{\rOE{5,6,7,8,10}}
    \proofstepwidestar[1,2,3]{\forall x\,\bigl(x\in (B\cap n)\cup\{n\} \rightarrow x\in B\bigr)}{\rUI{\rRI{4,11}}}
    \proofstepwidestar[1,2,3]{(B\cap n)\cup\{n\}\subseteq B}{%
      \FormulaRefAuto{A \subseteq B \coloneq \forall x\,(x\in A \rightarrow x\in B)}{12}}
  \closeproofpartwideindR

  \proofpartwide{Schluss}
    \proofstepwidestar[1]{n\in\mathbb{N}}{\rA}
    \proofstepwidestar[2]{B\subseteq \Succ(n)}{\rA}
    \proofstepwidestar[3]{n\in B}{\rA}

    \proofstepwidestar[1,2,3]{B\subseteq (B\cap n)\cup\{n\}}{%
      \FormulaRefAuto{n\in\mathbb{N}\dsep B\subseteq \Succ(n)\dsep n\in B \vdash B\subseteq (B\cap n)\cup\{n\}}{1,2,3}}
    \proofstepwidestar[1,2,3]{(B\cap n)\cup\{n\}\subseteq B}{%
      \FormulaRefAuto{n\in\mathbb{N}\dsep B\subseteq \Succ(n)\dsep n\in B \vdash (B\cap n)\cup\{n\}\subseteq B}{1,2,3}}
    \proofstepwidestar[1,2,3]{B=(B\cap n)\cup\{n\}}{%
      \FormulaRefAuto{A\subseteq B\dsep B\subseteq A \vdash A=B}{4,5}}
  \closeproofpartwide
\end{tabproofsplitwide}



\subsection{Minimalität der Menge der natürlichen Zahlen}

\FormulaThmDelta[Minimalität]{\Induktiv(A) \vdash \mathbb{N} \subseteq A}{\DeltaDecl{Mengen}{A}}
\begin{tabproof}
  \proofstep{1}{ \Induktiv(A) }{\rA}
  \proofstep{2}{ \mathbb{N} \subseteq A }{\FormulaRefAuto{P(C)\vdash \bigcap_{P(A)} A \subseteq C}{1}}
\end{tabproof}


\FormulaThmDelta{%
  n\in\mathbb{N}\vdash n\subseteq \Succ(n)%
}{
  \DeltaDecl{Mengen}{n}
}
\begin{tabproof}
  \proofstep{1}{n\in\mathbb{N}}{\rA}

  \proofstep{1}{\Succ(n)=n\cup\{n\}}{%
    \FormulaRefAuto{x\in\mathbb{N}\vdash \Succ(x)=x\cup\{x\}}{1}}

  \proofstep{}{\,n\subseteq n\cup\{n\}}{%
    \FormulaRefAuto{A\subseteq A\cup B}}

  \proofstep{1}{n\subseteq \Succ(n)}{\rIE{2,3}}
\end{tabproof}
\fi
\section{Grundlegende Eigenschaften von Funktionen auf und in die natürlichen Zahlen}

\subsection{Funktionen von \(\mathbb{N}\)}

% Die beiden Folgenbegriffe sind nun im wieder aktiven, von \Finite
% unabhängigen Funktionsblock von Band~03 registriert. Dieser Band verwendet diese
% Definitionen, registriert sie aber nicht ein zweites Mal.
\iffalse
\FormulaDefDeltaK[Begriff der Folge in einer Menge]{%
  H\colon\mathbb{N}\to M%
}{Folge in \(M\)}{%
  \DeltaRow{Mengen}{M\dsep n}%
  \DeltaRow{Funktionen}{H}%
  \DeltaRow{Definitionsbereich}{\mathbb{N}}%
  \DeltaRow{Wertebereich}{M}%
  \DeltaRow{Folgenglied zum Index \(n\)}{H(n)}%
}

\FormulaDefDeltaK[Begriff der Folge natürlicher Zahlen]{%
  u\colon\mathbb{N}\to\mathbb{N}%
}{Folge natürlicher Zahlen}{%
  \DeltaRow{Funktionen}{u}%
  \DeltaRow{Definitionsbereich}{\mathbb{N}}%
  \DeltaRow{Wertebereich}{\mathbb{N}}%
}
\fi

\FormulaThmDelta[Anfangspunkt liegt im Bild einer Funktion von \(\mathbb{N}\)]{%
  H\colon \mathbb{N}\to M \vdash H(0)\in H[\mathbb{N}]
}{%
  \DeltaDecl{Mengen}{M}%
  \DeltaDecl{Funktionen}{H}%
}
\begin{tabproof}
  \proofstep{1}{H\colon \mathbb{N}\to M}{\rA}
  \proofstep{}{0\in\mathbb{N}}{%
    \FormulaRefAuto{PeanoZero}}
  \proofstep{1}{H(0)\in H[\mathbb{N}]}{%
    \FormulaRefAuto{C\subseteq A\dsep F\colon A\to B\dsep x\in C\vdash F(x)\in F[C]}{%
      \FormulaRefAuto{A\subseteq A},1,2}}
\end{tabproof}

\FormulaThmDelta[Folgenglieder bleiben im Bild einer Funktion von \(\mathbb{N}\)]{%
  H\colon \mathbb{N}\to M\dsep n\in\mathbb{N}
  \vdash H(\Succ(n))\in H[\mathbb{N}]
}{%
  \DeltaDecl{Mengen}{M\dsep n}%
  \DeltaDecl{Funktionen}{H}%
}
\begin{tabproof}
  \proofstep{1}{H\colon \mathbb{N}\to M}{\rA}
  \proofstep{2}{n\in\mathbb{N}}{\rA}
  \proofstep{2}{\Succ(n)\in\mathbb{N}}{%
    \FormulaRefAuto{PeanoSuccClosure}{2}}
  \proofstep{1,2}{H(\Succ(n))\in H[\mathbb{N}]}{%
    \FormulaRefAuto{C\subseteq A\dsep F\colon A\to B\dsep x\in C\vdash F(x)\in F[C]}{%
      \FormulaRefAuto{A\subseteq A},1,3}}
\end{tabproof}

\subsection{Injektive Funktionen von \(\mathbb{N}\)}

\FormulaThmDelta[Anfangspunkt liegt im Bild]{%
  H\colon \mathbb{N}\inj M \vdash H(0)\in H[\mathbb{N}]
}{%
  \DeltaDecl{Mengen}{M}%
  \DeltaDecl{Funktionen}{H}%
}
\begin{tabproof}
  \proofstep{1}{H\colon \mathbb{N}\inj M}{\rA}
  \proofstep{}{0\in\mathbb{N}}{%
    \FormulaRefAuto{PeanoZero}}
  \proofstep{1}{H(0)\in H[\mathbb{N}]}{%
    \FormulaRefAuto{C\subseteq A\dsep F\colon A\inj B\dsep x\in C\vdash F(x)\in F[C]}{%
      \FormulaRefAuto{A\subseteq A},1,2}}
\end{tabproof}

\FormulaThmDelta[Folgenglieder bleiben im Bild]{%
  H\colon \mathbb{N}\inj M\dsep n\in\mathbb{N}
  \vdash H(\Succ(n))\in H[\mathbb{N}]
}{%
  \DeltaDecl{Mengen}{M\dsep n}%
  \DeltaDecl{Funktionen}{H}%
}
\begin{tabproof}
  \proofstep{1}{H\colon \mathbb{N}\inj M}{\rA}
  \proofstep{2}{n\in\mathbb{N}}{\rA}
  \proofstep{2}{\Succ(n)\in\mathbb{N}}{%
    \FormulaRefAuto{PeanoSuccClosure}{2}}
  \proofstep{1,2}{H(\Succ(n))\in H[\mathbb{N}]}{%
    \FormulaRefAuto{C\subseteq A\dsep F\colon A\inj B\dsep x\in C\vdash F(x)\in F[C]}{%
      \FormulaRefAuto{A\subseteq A},1,3}}
\end{tabproof}

\FormulaThmDelta[Typisierung des inversen Index]{%
  H\colon \mathbb{N}\inj M\dsep x\in H[\mathbb{N}]
  \vdash H^{-1}(x)\in\mathbb{N}
}{%
  \DeltaDecl{Mengen}{M\dsep x}%
  \DeltaDecl{Funktionen}{H}%
}
\begin{tabproof}
  \proofstep{1}{H\colon \mathbb{N}\inj M}{\rA}
  \proofstep{2}{x\in H[\mathbb{N}]}{\rA}
  \proofstep{1,2}{H^{-1}(x)\in\mathbb{N}}{%
    \FormulaRefAuto{F\colon A\to B\dsep x\in A\vdash F(x)\in B}{%
      \FormulaRefAuto{F\colon A\bij B\vdash F^{-1}\colon B\to A}{%
        \FormulaRefAuto{F\colon A\inj B \vdash F\colon A\bij F[A]}{1}},
      2}}
\end{tabproof}

\FormulaThmDelta[Folgenglieder treffen den Anfangspunkt nicht]{%
  H\colon \mathbb{N}\inj M\dsep n\in\mathbb{N}
  \vdash H(\Succ(n))\neq H(0)
}{%
  \DeltaDecl{Mengen}{M\dsep n}%
  \DeltaDecl{Funktionen}{H}%
}
\begin{tabproof}
  \proofstep{1}{H\colon \mathbb{N}\inj M}{\rA}
  \proofstep{2}{n\in\mathbb{N}}{\rA}
  \proofstep{}{0\in\mathbb{N}}{%
    \FormulaRefAuto{PeanoZero}}
  \proofstep{2}{\Succ(n)\in\mathbb{N}}{%
    \FormulaRefAuto{PeanoSuccClosure}{2}}
  \proofstep{5}{H(\Succ(n))=H(0)}{\rA}
  \proofstep{1,2,5}{\Succ(n)=0}{%
    \FormulaRefAuto{F\colon A\inj B\dsep x\in A\dsep y\in A\dsep F(x)=F(y)\vdash x=y}{1,4,3,5}}
  \proofstep{2}{\Succ(n)\neq0}{%
    \FormulaRefAuto{PeanoSuccNonzero}{2}}
  \proofstep{1,2,5}{\bot}{\rBI{6,7}}
  \proofstep{1,2}{H(\Succ(n))\neq H(0)}{\rCI{5,8}}
\end{tabproof}

\subsection{Surjektive Funktionen von \(\mathbb{N}\)}

\FormulaThmDelta[Anfangspunkt liegt im Bild einer surjektiven Funktion von \(\mathbb{N}\)]{%
  H\colon \mathbb{N}\sur M \vdash H(0)\in H[\mathbb{N}]
}{%
  \DeltaDecl{Mengen}{M}%
  \DeltaDecl{Funktionen}{H}%
}
\begin{tabproof}
  \proofstep{1}{H\colon \mathbb{N}\sur M}{\rA}
  \proofstep{}{0\in\mathbb{N}}{%
    \FormulaRefAuto{PeanoZero}}
  \proofstep{1}{H(0)\in H[\mathbb{N}]}{%
    \FormulaRefAuto{C\subseteq A\dsep F\colon A\sur B\dsep x\in C\vdash F(x)\in F[C]}{%
      \FormulaRefAuto{A\subseteq A},1,2}}
\end{tabproof}

\FormulaThmDelta[Folgenglieder bleiben im Bild einer surjektiven Funktion von \(\mathbb{N}\)]{%
  H\colon \mathbb{N}\sur M\dsep n\in\mathbb{N}
  \vdash H(\Succ(n))\in H[\mathbb{N}]
}{%
  \DeltaDecl{Mengen}{M\dsep n}%
  \DeltaDecl{Funktionen}{H}%
}
\begin{tabproof}
  \proofstep{1}{H\colon \mathbb{N}\sur M}{\rA}
  \proofstep{2}{n\in\mathbb{N}}{\rA}
  \proofstep{2}{\Succ(n)\in\mathbb{N}}{%
    \FormulaRefAuto{PeanoSuccClosure}{2}}
  \proofstep{1,2}{H(\Succ(n))\in H[\mathbb{N}]}{%
    \FormulaRefAuto{C\subseteq A\dsep F\colon A\sur B\dsep x\in C\vdash F(x)\in F[C]}{%
      \FormulaRefAuto{A\subseteq A},1,3}}
\end{tabproof}

\subsection{Bijektive Funktionen von \(\mathbb{N}\)}

\FormulaThmDelta[Anfangspunkt liegt im Bild einer bijektiven Funktion von \(\mathbb{N}\)]{%
  H\colon \mathbb{N}\bij M \vdash H(0)\in H[\mathbb{N}]
}{%
  \DeltaDecl{Mengen}{M}%
  \DeltaDecl{Funktionen}{H}%
}
\begin{tabproof}
  \proofstep{1}{H\colon \mathbb{N}\bij M}{\rA}
  \proofstep{}{0\in\mathbb{N}}{%
    \FormulaRefAuto{PeanoZero}}
  \proofstep{1}{H(0)\in H[\mathbb{N}]}{%
    \FormulaRefAuto{C\subseteq A\dsep F\colon A\bij B\dsep x\in C\vdash F(x)\in F[C]}{%
      \FormulaRefAuto{A\subseteq A},1,2}}
\end{tabproof}

\FormulaThmDelta[Folgenglieder bleiben im Bild einer bijektiven Funktion von \(\mathbb{N}\)]{%
  H\colon \mathbb{N}\bij M\dsep n\in\mathbb{N}
  \vdash H(\Succ(n))\in H[\mathbb{N}]
}{%
  \DeltaDecl{Mengen}{M\dsep n}%
  \DeltaDecl{Funktionen}{H}%
}
\begin{tabproof}
  \proofstep{1}{H\colon \mathbb{N}\bij M}{\rA}
  \proofstep{2}{n\in\mathbb{N}}{\rA}
  \proofstep{2}{\Succ(n)\in\mathbb{N}}{%
    \FormulaRefAuto{PeanoSuccClosure}{2}}
  \proofstep{1,2}{H(\Succ(n))\in H[\mathbb{N}]}{%
    \FormulaRefAuto{C\subseteq A\dsep F\colon A\bij B\dsep x\in C\vdash F(x)\in F[C]}{%
      \FormulaRefAuto{A\subseteq A},1,3}}
\end{tabproof}

\iffalse
Von-Neumann-spezifische Anwendungen aus Band 03. Dieser Block enthält unter
anderem die alte Ordnung über \(m\in n\) und \(m\subseteq n\) sowie Folgerungen
aus \(\Succ(n)=n\cup\{n\}\). Die Peano-basierte Ordnung wird nach der Addition
neu eingeführt.
\section{Anwendungen des Induktionsprinzips}



\FormulaThmDelta{%
  n\in\mathbb{N}\vdash n\subseteq\mathbb{N}%
}{%
  \DeltaDecl{Mengen}{n\dsep x}%
}
\begin{tabproofsplitwide}
  \proofpartwideind[Induktionsanfang]{0 \subseteq \mathbb{N}}
    \proofstepwidestar{0 \subseteq \mathbb{N}}{%
      \FormulaRefAuto{\varnothing\subseteq A}}
  \closeproofpartwideind

  \proofpartwideind[Induktionsschritt]{%
    n\in\mathbb{N}\dsep n\subseteq\mathbb{N}\vdash \Succ(n)\subseteq\mathbb{N}}
    \proofstepwidestar[1]{n\in\mathbb{N}}{\rA}
    \proofstepwidestar[2]{n\subseteq\mathbb{N}}{\rA}

    \proofstepwidestar[3]{x\in\Succ(n)}{\rA}
    \proofstepwidestar[1,3]{x\in n\lor x=n}{%
      \FormulaRefAuto{n\in\mathbb{N}\dsep x\in\Succ(n)\vdash x\in n\lor x=n}{1,3}}

    \proofstepwidestar[5]{x\in n}{\rA}
    \proofstepwidestar[2,5]{x\in\mathbb{N}}{%
      \FormulaRefAuto{A\subseteq B,\, x\in A\vdash x\in B}{2,5}}

    \proofstepwidestar[7]{x=n}{\rA}
    \proofstepwidestar[1,7]{x\in\mathbb{N}}{\rIE{7,1}}

    \proofstepwidestar[1,2,3]{x\in\mathbb{N}}{\rOE{4,5,6,7,8}}

    \proofstepwidestar[1,2]{x\in\Succ(n)\rightarrow x\in\mathbb{N}}{\rRI{3,9}}
    \proofstepwidestar[1,2]{\forall x\,(x\in\Succ(n)\rightarrow x\in\mathbb{N})}{\rUI{10}}
    \proofstepwidestar[1,2]{\Succ(n)\subseteq\mathbb{N}}{%
      \FormulaRefAuto{A \subseteq B \coloneq \forall x\,(x \in A \rightarrow x \in B)}{11}}
  \closeproofpartwideind

  \proofpartwide{Induktionsschluss}
    \proofstepwidestar[1]{n\in\mathbb{N}}{\rA}
    \proofstepwidestar[1]{n\subseteq\mathbb{N}}{%
      \FormulaRefAuto{PeanoInductionRule}{%
        \FormulaRefAuto{\varnothing \subseteq \mathbb{N}},%
        \FormulaRefAuto{n\in\mathbb{N}\dsep n\subseteq\mathbb{N}\vdash \Succ(n)\subseteq\mathbb{N}},%
        1}}
  \closeproofpartwide
\end{tabproofsplitwide}




\subsection{Transitivitätskette}

\FormulaThmDelta[Transitivität der natürlichen Zahlen]{%
  n\in\mathbb{N}\vdash \bigcup n\subseteq n%
}{%
  \DeltaDecl{Mengen}{n}%
}
\begin{tabproofsplitwide}
  \proofpartwideind[Induktionsanfang]{\bigcup 0 \subseteq 0}
    \proofstepwidestar{\bigcup 0 = 0}{\FormulaRefAuto{\bigcup \varnothing = \varnothing}}
    \proofstepwidestar{\bigcup 0 \subseteq 0}{\FormulaRefAuto{A=B\vdash A\subseteq B}}
  \closeproofpartwideind

  \proofpartwideind[Induktionsschritt]{n\in\mathbb{N}\dsep \bigcup n\subseteq n\vdash \bigcup \Succ(n)\subseteq \Succ(n)}
    \proofstepwidestar[1]{n\in\mathbb{N}}{\rA}
    \proofstepwidestar[2]{\bigcup n\subseteq n}{\rA}
    \proofstepwidestar[1]{\Succ(n)=n\cup\{n\}}{%
      \FormulaRefAuto{n\in\mathbb{N}\vdash \Succ(n) = n\cup\{n\}}{1}}

    \proofstepwidestar[3]{x\in\bigcup\Succ(n)}{\rA}
    \proofstepwidestar[3]{\exists B\in\Succ(n)\,(x\in B)}{%
      \FormulaRefAuto{x \in \bigcup A \eqvdash \exists B\in A\,(x \in B)}{2}}

    \proofstepwidestar[4]{B\in\Succ(n)\land x\in B}{\rA}
    \proofstepwidestar[4]{B\in\Succ(n)}{\rAEa{4}}
    \proofstepwidestar[4]{x\in B}{\rAEb{4}}

    \proofstepwidestar[1,4]{B\in n\lor B=n}{%
      \FormulaRefAuto{n\in\mathbb{N}\dsep x\in\Succ(n)\vdash x\in n\lor x=n}{1,5}}

    \proofstepwidestar[5]{B\in n}{\rA}
    \proofstepwidestar[4,5]{x\in\bigcup n}{%
      \FormulaRefAuto{B\in A\dsep x\in B\vdash x \in \bigcup A}{5,8}}
    \proofstepwidestar[2,4,5]{x\in n}{%
      \FormulaRefAuto{A\subseteq B, x\in A\vdash x\in B}{2,11}}
    \proofstepwidestar[2,4,5]{x\in n\cup\{n\}}{%
      \FormulaRefAuto{z \in A \vdash z \in A \cup B}{13}}
    \proofstepwidestar[1,2,4,5]{x\in\Succ(n)}{\rIE{3,14}}

    \proofstepwidestar[6]{B=n}{\rA}
    \proofstepwidestar[4,6]{x\in n}{\rIE{6,8}}
    \proofstepwidestar[4,6]{x\in n\cup\{n\}}{%
      \FormulaRefAuto{z \in A \vdash z \in A \cup B}{16}}
    \proofstepwidestar[1,4,6]{x\in\Succ(n)}{\rIE{3,17}}

    \proofstepwidestar[1,2,4]{x\in\Succ(n)}{\rOE{9,10,15,16,18}}
    \proofstepwidestar[1,2,3]{x\in\Succ(n)}{\rEE{5,6,19}}

    \proofstepwidestar[1,2]{\bigcup\Succ(n)\subseteq \Succ(n)}{%
      \FormulaRefAuto{A \subseteq B \coloneq \forall x\,(x \in A \rightarrow x \in B)}{\rUI{\rRI{4,20}}}}
  \closeproofpartwideind
 \proofpartwide{Induktionsschluss}
    \proofstepwidestar[1]{n\in\mathbb{N}}{\rA}
    \proofstepwidestar[1]{\bigcup n\subseteq n}{%
      \FormulaRefAuto{PeanoInductionRule}{\FormulaRefAuto{\bigcup \varnothing \subseteq \varnothing},\FormulaRefAuto{n\in\mathbb{N}\dsep \bigcup n\subseteq n\vdash \bigcup \Succ(n)\subseteq \Succ(n)},1}
    }
  \closeproofpartwide
\end{tabproofsplitwide}



\FormulaThmDelta[Transitivitätskette]{%
  n\in\mathbb{N}\vdash \bigcup n\subseteq n = \bigcup\Succ(n) \subseteq \Succ(n)%
}{%
  \DeltaDecl{Mengen}{n}%
}
\begin{tabproofsplitwide}
  \proofpartwideind[Teil I]{n\in\mathbb{N}\vdash n=\bigcup\Succ(n)}
    \proofstepwidestar[1]{n\in\mathbb{N}}{\rA}

    \proofstepwidestar[2]{x\in n}{\rA}
    \proofstepwidestar[1]{\Succ(n)=n\cup\{n\}}{%
      \FormulaRefAuto{n\in\mathbb{N}\vdash \Succ(n) = n\cup\{n\}}{1}}
    \proofstepwidestar[]{n\in\{n\}}{\FormulaRefAuto{a\in\{a\}}}
    \proofstepwidestar[]{n\in n\cup\{n\}}{%
      \FormulaRefAuto{z\in B\vdash z\in A\cup B}{4}}
    \proofstepwidestar[1]{n\in\Succ(n)}{\rIE{3,5}}
    \proofstepwidestar[1,2]{x\in\bigcup\Succ(n)}{%
      \FormulaRefAuto{B\in A\dsep x\in B\vdash x \in \bigcup A}{6,2}}
    \proofstepwidestar[1]{n\subseteq\bigcup\Succ(n)}{%
      \FormulaRefAuto{A \subseteq B \coloneq \forall x\,(x \in A \rightarrow x \in B)}{\rUI{\rRI{2,7}}}}

    \proofstepwidestar[9]{y\in\bigcup\Succ(n)}{\rA}
    \proofstepwidestar[9]{\exists B\in\Succ(n)\,(y\in B)}{%
      \FormulaRefAuto{y \in \bigcup A \eqvdash \exists B\in A\,(y \in B)}{9}}

    \proofstepwidestar[11]{B\in\Succ(n)\land y\in B}{\rA}
    \proofstepwidestar[11]{B\in\Succ(n)}{\rAEa{11}}
    \proofstepwidestar[11]{y\in B}{\rAEb{11}}

    \proofstepwidestar[1,11]{B\in n\cup\{n\}}{\rIE{3,12}}
    \proofstepwidestar[1,11]{B\in n\lor B\in\{n\}}{%
      \FormulaRefAuto{z \in A \cup B \eqvdash z \in A \lor z \in B}{14}}

    \proofstepwidestar[16]{B\in n}{\rA}
    \proofstepwidestar[11,16]{y\in\bigcup n}{%
      \FormulaRefAuto{B\in A\dsep y\in B\vdash y \in \bigcup A}{16,13}}
    \proofstepwidestar[1]{\bigcup n\subseteq n}{%
      \FormulaRefAuto{n\in\mathbb{N}\vdash \bigcup n\subseteq n}{1}}
    \proofstepwidestar[1,11,16]{y\in n}{%
      \FormulaRefAuto{A\subseteq B, y\in A\vdash y\in B}{18,17}}

    \proofstepwidestar[20]{B\in\{n\}}{\rA}
    \proofstepwidestar[20]{B=n}{%
      \FormulaRefAuto{z \in \{a\} \eqvdash z = a}{20}}
    \proofstepwidestar[11,20]{y\in n}{\rIE{21,13}}

    \proofstepwidestar[1,11]{y\in n}{\rOE{15,16,19,20,22}}
    \proofstepwidestar[1,9]{y\in n}{\rEE{10,11,24}}

    \proofstepwidestar[1]{\bigcup\Succ(n)\subseteq n}{%
      \FormulaRefAuto{A \subseteq B \coloneq \forall y\,(y \in A \rightarrow y \in B)}{\rUI{\rRI{9,24}}}}

    \proofstepwidestar[1]{n=\bigcup\Succ(n)}{%
      \FormulaRefAuto{A\subseteq B, B\subseteq A\vdash A=B}{8,25}}
  \closeproofpartwide


  \proofpartwideind[Teil II]{n\in\mathbb{N}\vdash \bigcup\Succ(n)\subseteq \Succ(n)}
    \proofstepwidestar[1]{n\in\mathbb{N}}{\rA}
    \proofstepwidestar[]{\forall x\in \mathbb{N}\,\Succ(x)\in \mathbb{N}}{%
      \FormulaRefAuto{PeanoSuccClosureUniversal}}
    \proofstepwidestar[1]{\Succ(n)\in \mathbb{N}}{\rRE{\rUE{2},1}}
    \proofstepwidestar[1]{\bigcup\Succ(n)\subseteq \Succ(n)}{%
      \FormulaRefAuto{n\in\mathbb{N}\vdash \bigcup n\subseteq n}{3}}
  \closeproofpartwide


  \proofpartwide{Zusammenfassung}
    \proofstepwidestar[1]{n\in\mathbb{N}}{\rA}
    \proofstepwide[1]{\bigcup n}{\subseteq}{n}{
      \FormulaRefAuto{n\in\mathbb{N}\vdash \bigcup n\subseteq n}{1}}
    \proofstepwide[1]{}{=}{\bigcup\Succ(n)}{
      \FormulaRefAuto{n\in\mathbb{N}\vdash n=\bigcup\Succ(n)}{1}}
    \proofstepwide[1]{}{\subseteq}{\Succ(n)}{
      \FormulaRefAuto{n\in\mathbb{N}\vdash \bigcup\Succ(n)\subseteq \Succ(n)}{1}}
  \closeproofpartwide
\end{tabproofsplitwide}

\FormulaThmDelta{%
  u\in\mathbb{N}\vdash \bigcup u\in\mathbb{N}%
}{%
  \DeltaDecl{Mengen}{u\dsep n}%
}
\begin{tabproofsplitwide}

  \proofpartwideind[Induktionsanfang]{\bigcup0\in\mathbb{N}}
    \proofstepwidestar{\bigcup0=0}{%
      \FormulaRefAuto{\bigcup\varnothing=\varnothing}}
    \proofstepwidestar{0\in\mathbb{N}}{%
      \FormulaRefAuto{PeanoZero}}
    \proofstepwidestar{\bigcup0\in\mathbb{N}}{\rIE{1,2}}
  \closeproofpartwideind

  \proofpartwideind[Induktionsschritt]{%
    n\in\mathbb{N}\dsep \bigcup n\in\mathbb{N}\vdash \bigcup\Succ(n)\in\mathbb{N}}
    \proofstepwidestar[1]{n\in\mathbb{N}}{\rA}
    \proofstepwidestar[2]{\bigcup n\in\mathbb{N}}{\rA}

    \proofstepwidestar[1]{n=\bigcup\Succ(n)}{%
      \FormulaRefAuto{n\in\mathbb{N}\vdash n=\bigcup\Succ(n)}{1}}
    \proofstepwidestar[1]{\bigcup\Succ(n)=n}{%
      \FormulaRefAuto{a=b\vdash b=a}{3}}
    \proofstepwidestar[1]{\bigcup\Succ(n)\in\mathbb{N}}{\rIE{4,1}}
  \closeproofpartwideind

\end{tabproofsplitwide}


\FormulaThmDelta[Element liegt in seinem Nachfolger]{%
  n\in\mathbb{N}\vdash n\in\Succ(n)%
}{%
  \DeltaDecl{Mengen}{n}%
}
\begin{tabproof}
  \proofstep{1}{n\in\mathbb{N}}{\rA}

  \proofstep{}{n\in\{n\}}{%
    \FormulaRefAuto{a\in\{a\}}}

  \proofstep{}{n\in n\cup\{n\}}{%
    \FormulaRefAuto{z\in B\vdash z\in A\cup B}{2}}

  \proofstep{1}{\Succ(n)=n\cup\{n\}}{%
    \FormulaRefAuto{x\in\mathbb{N}\vdash \Succ(x)=x\cup\{x\}}{1}}

  \proofstep{1}{n\in\Succ(n)}{\rIE{4,3}}
\end{tabproof}

\subsection{Die Vorgängerfunktion}


\FormulaDefDelta[Funktionsvorschrift der Vorgängerfunktion]{%
  n\in\mathbb{N}\vdash t_{\Pred}(n) \coloneqq \bigcup n%
}{%
  \DeltaRow{Mengen}{n}%
  \DeltaRow{Funktionensymbol}{t_{\Pred}}%
}




\FormulaThmDelta[Typisierungsaxiom für $t_{\Pred}$]{%
  n\in\mathbb{N}\vdash t_{\Pred}(n)\in\mathbb{N}%
}{%
  \DeltaRow{Mengen}{n}%
  \DeltaRow{Funktionensymbol}{t_{\Pred}}%
}
\begin{tabproof}
  \proofstep{1}{n\in\mathbb{N}}{\rA}
  \proofstep{1}{t_{\Pred}(n)=\bigcup n}{%
    \FormulaRefAuto{u\in\mathbb{N}\vdash t_{\Pred}(u)\coloneqq \bigcup u}{1}}
  \proofstep{1}{\bigcup n\in\mathbb{N}}{%
    \FormulaRefAuto{u\in\mathbb{N}\vdash \bigcup u\in\mathbb{N}}{1}}
  \proofstep{1}{t_{\Pred}(n)\in\mathbb{N}}{\rIE{2,3}}
\end{tabproof}

\FormulaThmDelta[Vorgängerfunktion als Funktion]{%
  \exists! F\colon \mathbb{N}\to\mathbb{N}\,\forall n\in\mathbb{N}\,F(n)=t_{\Pred}(n)%
}{%
  \DeltaRow{Mengen}{n}%
  \DeltaRow{Funktionensymbol}{t_{\Pred}}%
}
\begin{tabproof}
  \proofstep{}{%
    \forall n\in\mathbb{N}\,t_{\Pred}(n)=\bigcup n}{%
    \FormulaRefAuto{u\in\mathbb{N}\vdash t_{\Pred}(u)\coloneqq \bigcup u}}
  \proofstep{}{%
    \forall n\in\mathbb{N}\,t_{\Pred}(n)\in\mathbb{N}}{%
    \FormulaRefAuto{u\in\mathbb{N}\vdash t_{\Pred}(u)\in\mathbb{N}}}
  \proofstep{}{%
    \exists! F\colon \mathbb{N}\to\mathbb{N}\,\forall n\in\mathbb{N}\,F(n)=t_{\Pred}(n)}{%
    \FormulaRefAuto{x\in A\vdash t(x)\coloneq y\dsep x\in A\vdash t(x)\in B\exists! F\colon A\to B\,\forall x\in A\,F(x) = t(x)}{1,2}}
\end{tabproof}

\FormulaDefDelta[Vorgängerfunktion]{%
  \Pred\coloneqq \iota F\Bigl(F\colon \mathbb{N}\to\mathbb{N}\land \forall n\in\mathbb{N}\,F(n)=t_{\Pred}(n)\Bigr)%
}{%
  \DeltaRow{Mengen}{n}%
  \DeltaRow{Funktionensymbol}{t_{\Pred}}%
}

\FormulaThmDelta[Vorgängerfunktion]{%
  \Pred\colon \mathbb{N}\to\mathbb{N}%
}{%
  \DeltaRow{Mengen}{n}%
}
\begin{tabproof}
  \proofstep{}{%
    \Pred\colon \mathbb{N}\to\mathbb{N}}{%
    \rAEa{\FormulaRefAuto{\Pred\coloneqq \iota F\Bigl(F\colon \mathbb{N}\to\mathbb{N}\land \forall n\in\mathbb{N}\,F(n)=t_{\Pred}(n)\Bigr)}}}
\end{tabproof}

\FormulaThmDelta[Vorgänger einer natürlichen Zahl ist natürlich]{%
  n\in\mathbb{N}\vdash \Pred(n)\in\mathbb{N}%
}{%
  \DeltaRow{Mengen}{n}%
}
\begin{tabproof}
  \proofstep{1}{n\in\mathbb{N}}{\rA}
  \proofstep{}{\Pred\colon \mathbb{N}\to\mathbb{N}}{%
    \FormulaRefAuto{\Pred\colon \mathbb{N}\to\mathbb{N}}}
  \proofstep{1}{\Pred(n)\in\mathbb{N}}{\FormulaRefAuto{F\colon A\to B\dsep a\in A\vdash F(a)\in B}{1,2}}
\end{tabproof}

\FormulaThmDelta[Grundgleichung der Vorgängerfunktion]{%
  n\in\mathbb{N}\vdash \Pred(n)=\bigcup n%
}{%
  \DeltaRow{Mengen}{n}%
}
\begin{tabproof}
  \proofstep{1}{n\in\mathbb{N}}{\rA}
  \proofstep{}{%
    \forall n\in\mathbb{N}\,\Pred(n)=t_{\Pred}(n)}{%
    \rAEb{\FormulaRefAuto{\Pred\coloneqq \iota F\Bigl(F\colon \mathbb{N}\to\mathbb{N}\land \forall n\in\mathbb{N}\,F(n)=t_{\Pred}(n)\Bigr)}}}
  \proofstep{1}{\Pred(n)=t_{\Pred}(n)}{\rRE{1,\rUE{2}}}
  \proofstep{1}{t_{\Pred}(n)=\bigcup n}{%
    \FormulaRefAuto{u\in\mathbb{N}\vdash t_{\Pred}(u)\coloneqq \bigcup u}{1}}
  \proofstep{1}{\Pred(n)=\bigcup n}{%
    \FormulaRefAuto{a=b,\,b=c\vdash a=c}{3,4}}
\end{tabproof}



\FormulaThmDelta[Vorgänger eines Nachfolgers]{%
  n\in\mathbb{N}\vdash \Pred(\Succ(n))=n%
}{%
  \DeltaDecl{Mengen}{n}%
}
\begin{tabproof}
  \proofstep{1}{n\in\mathbb{N}}{\rA}
  \proofstep{1}{\Succ(n)\in\mathbb{N}}{%
    \FormulaRefAuto{PeanoSuccClosure}{1}}
  \proofstep{1}{\Pred(\Succ(n))=\bigcup\Succ(n)}{%
    \FormulaRefAuto{u\in\mathbb{N}\vdash \Pred(u)=\bigcup u}{2}}
  \proofstep{1}{n=\bigcup\Succ(n)}{%
    \FormulaRefAuto{n\in\mathbb{N}\vdash n=\bigcup\Succ(n)}{1}}
  \proofstep{1}{\bigcup\Succ(n)=n}{%
    \FormulaRefAuto{a=b\vdash b=a}{4}}
  \proofstep{1}{\Pred(\Succ(n))=n}{%
    \FormulaRefAuto{a=b,\,b=c\vdash a=c}{3,5}}
\end{tabproof}

% =========================================================
% Im Abschnitt: Die Vorgängerfunktion
% Direkt nach: Vorgänger eines Nachfolgers
% =========================================================

\FormulaThmDelta{%
  n\in\mathbb{N}\dsep n\neq0\vdash n=\Succ(\Pred(n))
}{%
  \DeltaDecl{Mengen}{n\dsep y}%
}
\begin{tabproof}
  \proofstep{1}{n\in\mathbb{N}}{\rA}
  \proofstep{2}{n\neq0}{\rA}

  \proofstep{1,2}{\exists y\in\mathbb{N}\,\bigl(n=\Succ(y)\bigr)}{%
    \FormulaRefAuto{PeanoNonzeroIsSucc}{1,2}}

  \proofstep{4}{y\in\mathbb{N}}{\rA}
  \proofstep{5}{n=\Succ(y)}{\rA}

  \proofstep{4}{\Succ(y)\in\mathbb{N}}{%
    \FormulaRefAuto{PeanoSuccClosure}{4}}

  \proofstep{}{\,\Pred\colon \mathbb{N}\to\mathbb{N}}{%
    \FormulaRefAuto{\Pred\colon \mathbb{N}\to\mathbb{N}}}

  \proofstep{1,4,5}{\Pred(n)=\Pred(\Succ(y))}{%
    \FormulaRefAuto{x\in A\dsep y\in A\dsep x=y\vdash F(x)=F(y)}{1,6,5}}

  \proofstep{4}{\Pred(\Succ(y))=y}{%
    \FormulaRefAuto{n\in\mathbb{N}\vdash \Pred(\Succ(n))=n}{4}}

  \proofstep{1,4,5}{\Pred(n)=y}{%
    \FormulaRefAuto{a=b,\,b=c\vdash a=c}{8,9}}

  \proofstep{1,4,5}{y=\Pred(n)}{%
    \FormulaRefAuto{a=b\vdash b=a}{10}}

  \proofstep{1}{\Pred(n)\in\mathbb{N}}{%
    \FormulaRefAuto{F\colon A\to B\dsep x\in A\vdash F(x)\in B}{7,1}}

  \proofstep{1,4,5}{\Succ(y)=\Succ(\Pred(n))}{%
    \FormulaRefAuto{x\in A\dsep y\in A\dsep x=y\vdash F(x)=F(y)}{4,12,11}}

  \proofstep{1,4,5}{n=\Succ(\Pred(n))}{%
    \FormulaRefAuto{a=b,\,b=c\vdash a=c}{5,13}}

  \proofstep{1,2}{n=\Succ(\Pred(n))}{%
    \rEE{3,4,5,14}}
\end{tabproof}

\FormulaThmDelta{%
  k\in\mathbb{N}\dsep k\neq0\dsep \Pred(k)=n\vdash k=\Succ(n)
}{%
  \DeltaDecl{Mengen}{k\dsep n}%
}
\begin{tabproof}
  \proofstep{1}{k\in\mathbb{N}}{\rA}
  \proofstep{2}{k\neq0}{\rA}
  \proofstep{3}{\Pred(k)=n}{\rA}

  \proofstep{1,2}{k=\Succ(\Pred(k))}{%
    \FormulaRefAuto{n\in\mathbb{N}\dsep n\neq0\vdash n=\Succ(\Pred(n))}{1,2}}

  \proofstep{1}{\Pred(k)\in\mathbb{N}}{%
    \FormulaRefAuto{n\in\mathbb{N}\vdash \Pred(n)\in\mathbb{N}}{1}}

  \proofstep{1,3}{n\in\mathbb{N}}{\rIE{3,5}}

  \proofstep{1,3}{\Succ(\Pred(k))=\Succ(n)}{%
    \FormulaRefAuto{x\in A\dsep y\in A\dsep x=y\vdash F(x)=F(y)}{5,6,3}}

  \proofstep{1,2,3}{k=\Succ(n)}{%
    \FormulaRefAuto{a=b,\,b=c\vdash a=c}{4,7}}
\end{tabproof}

\subsection{Keine Selbstmitgliedschaft in den natürlichen Zahlen}

\FormulaThmDelta{%
  n\in\mathbb{N}\vdash n\notin n%
}{%
  \DeltaDecl{Mengen}{n}%
}
\begin{tabproofsplitwide}
  \proofpartwideind[Induktionsanfang]{0 \notin 0}
    \proofstepwidestar[1]{0\in0}{\rA}
    \proofstepwidestar[1]{0\neq0}{%
      \FormulaRefAuto{u\in A \vdash A\neq \varnothing}{1}}
    \proofstepwidestar{0=0}{\rII}
    \proofstepwidestar[1]{\bot}{\rBI{3,2}}
    \proofstepwidestar{0\notin0}{\rCI{1,4}}
  \closeproofpartwideind

  \proofpartwideind[Induktionsschritt]{n\in\mathbb{N}\dsep n\notin n\vdash \Succ(n)\notin\Succ(n)}
    \proofstepwidestar[1]{n\in\mathbb{N}}{\rA}
    \proofstepwidestar[2]{n\notin n}{\rA}

    \proofstepwidestar[1]{\Succ(n)=n\cup\{n\}}{%
      \FormulaRefAuto{x\in\mathbb{N}\vdash \Succ(x)=x\cup\{x\}}{1}}

    \proofstepwidestar{n\in\{n\}}{\FormulaRefAuto{a\in\{a\}}}
    \proofstepwidestar{n\in n\cup\{n\}}{%
      \FormulaRefAuto{u\in B \vdash u\in A\cup B}{4}}
    \proofstepwidestar[1]{n\in\Succ(n)}{\rIE{3,5}}

    \proofstepwidestar[7]{\Succ(n)\in\Succ(n)}{\rA}
    \proofstepwidestar[1,7]{\Succ(n)\in n\cup\{n\}}{\rIE{3,7}}
    \proofstepwidestar[1,7]{\Succ(n)\in n \lor \Succ(n)\in\{n\}}{%
      \FormulaRefAuto{z \in A \cup B \eqvdash z \in A \lor z \in B}{8}}

    \proofstepwidestar[10]{\Succ(n)\in n}{\rA}
    \proofstepwidestar[1,10]{n\in\bigcup n}{%
      \FormulaRefAuto{B\in A\dsep x\in B\vdash x \in \bigcup A}{10,6}}
    \proofstepwidestar[1]{\bigcup n\subseteq n}{%
      \FormulaRefAuto{n\in\mathbb{N}\vdash \bigcup n\subseteq n}{1}}
    \proofstepwidestar[1,10]{n\in n}{%
      \FormulaRefAuto{A\subseteq B, x\in A\vdash x\in B}{12,11}}
    \proofstepwidestar[1,2,10]{\bot}{\rBI{13,2}}

    \proofstepwidestar[15]{\Succ(n)\in\{n\}}{\rA}
    \proofstepwidestar[15]{\Succ(n)=n}{%
      \FormulaRefAuto{z \in \{a\} \eqvdash z = a}{15}}
    \proofstepwidestar[1,15]{n\in n}{\rIE{16,6}}
    \proofstepwidestar[1,2,15]{\bot}{\rBI{17,2}}

    \proofstepwidestar[1,2,7]{\bot}{\rOE{9,10,14,15,18}}
    \proofstepwidestar[1,2]{\Succ(n)\notin\Succ(n)}{\rCI{7,19}}
  \closeproofpartwideind

  \proofpartwide{Anwendung des Induktionsprinzips}
    \proofstepwidestar[1]{n\in\mathbb{N}}{\rA}
    \proofstepwidestar[1]{n\notin n}{%
      \FormulaRefAuto{PeanoInductionRule}{%
        \FormulaRefAuto{\varnothing \notin \varnothing},%
        \FormulaRefAuto{n\in\mathbb{N}\dsep n\notin n\vdash \Succ(n)\notin\Succ(n)},%
        1}}
  \closeproofpartwide
\end{tabproofsplitwide}

\subsection{Weitere Eigenschaften}

\FormulaThmDelta{%
  n\in\mathbb{N}\vdash 0 = n \lor 0 \in n%
}{%
  \DeltaDecl{Mengen}{n}%
}
\begin{tabproofsplitwide}
  \proofpartwideind[Induktionsanfang]{0 = 0 \lor 0 \in 0}
    \proofstepwidestar{0=0}{\rII}
    \proofstepwidestar{0 = 0 \lor 0 \in 0}{\rOIa{1}}
  \closeproofpartwideind

  \proofpartwideind[Induktionsschritt]{%
    n\in\mathbb{N}\dsep 0 = n \lor 0 \in n \vdash
    0 = \Succ(n) \lor 0 \in \Succ(n)}
    \proofstepwidestar[1]{n\in\mathbb{N}}{\rA}
    \proofstepwidestar[2]{0 = n \lor 0 \in n}{\rA}

    \proofstepwidestar[1]{\Succ(n)=n\cup\{n\}}{%
      \FormulaRefAuto{x\in\mathbb{N}\vdash \Succ(x)=x\cup\{x\}}{1}}

    \proofstepwidestar[4]{0 = n}{\rA}
    \proofstepwidestar[4]{n=0}{%
      \FormulaRefAuto{a=b\vdash b=a}{4}}
    \proofstepwidestar{n\in\{n\}}{\FormulaRefAuto{a\in\{a\}}}
    \proofstepwidestar[4]{0\in\{n\}}{\rIE{5,6}}
    \proofstepwidestar[4]{0\in n\cup\{n\}}{%
      \FormulaRefAuto{u\in B \vdash u\in A\cup B}{7}}
    \proofstepwidestar[1,4]{0\in\Succ(n)}{\rIE{3,8}}

    \proofstepwidestar[10]{0 \in n}{\rA}
    \proofstepwidestar[10]{0\in n\cup\{n\}}{%
      \FormulaRefAuto{z \in A \vdash z \in A \cup B}{10}}
    \proofstepwidestar[1,10]{0\in\Succ(n)}{\rIE{3,11}}

    \proofstepwidestar[1,2]{0\in\Succ(n)}{\rOE{2,4,9,10,12}}
    \proofstepwidestar[1,2]{0 = \Succ(n) \lor 0 \in \Succ(n)}{\rOIb{13}}
  \closeproofpartwideind

  \proofpartwide{Induktionsschluss}
    \proofstepwidestar[1]{n\in\mathbb{N}}{\rA}
    \proofstepwidestar[1]{0 = n \lor 0 \in n}{%
      \FormulaRefAuto{PeanoInductionRule}{%
        \FormulaRefAuto{\varnothing = \varnothing \lor \varnothing \in \varnothing},%
        \FormulaRefAuto{n\in\mathbb{N}\dsep \varnothing = n \lor \varnothing \in n \vdash \varnothing = \Succ(n) \lor \varnothing \in \Succ(n)},%
        1}}
  \closeproofpartwide
\end{tabproofsplitwide}

\FormulaThmDelta{%
  m\in\mathbb{N}\dsep n\in\mathbb{N}\dsep m\in n\vdash \Succ(m)=n\lor \Succ(m)\in n%
}{%
  \DeltaDecl{Mengen}{m\dsep n}%
}
\begin{tabproofsplitwide}
  \proofpartwideind[Induktionsanfang]{%
    m\in0 \rightarrow \bigl(\Succ(m)=0 \lor \Succ(m)\in0\bigr)}
    \proofstepwidestar[1]{m\in0}{\rA}
    \proofstepwidestar[1]{0\neq0}{%
      \FormulaRefAuto{u\in A \vdash A\neq \varnothing}{1}}
    \proofstepwidestar{0=0}{\rII}
    \proofstepwidestar[1]{0=0 \land 0\neq0}{\rAI{3,2}}
    \proofstepwidestar[1]{\Succ(m)=0 \lor \Succ(m)\in0}{%
      \FormulaRefAuto{P \land \neg P\vdash Q}{4}}
    \proofstepwidestar[]{m\in0 \rightarrow \bigl(\Succ(m)=0 \lor \Succ(m)\in0\bigr)}{\rRI{1,5}}
  \closeproofpartwideind

  \proofpartwideindDelta[Induktionsschritt]{x\in\mathbb{N}\dsep m\in\Succ(x)\vdash \Succ(m)=\Succ(x)\lor \Succ(m)\in\Succ(x)}{%
    \DeltaPrem{Induktionsvoraussetzung}{x\in\mathbb{N}\dsep m\in x\vdash \Succ(m)=x \lor \Succ(m)\in x}
  }[x\in\mathbb{N}\dsep \bigl(m\in x \rightarrow (\Succ(m)=x \lor \Succ(m)\in x)\bigr)\vdash m\in\Succ(x)\rightarrow\bigl(\Succ(m)=\Succ(x)\lor \Succ(m)\in\Succ(x)\bigr)]

    \proofstepwidestar[1]{x\in\mathbb{N}}{\rA}
    \proofstepwidestar[2]{m\in\Succ(x)}{\rA}

    \proofstepwidestar[1]{\Succ(x)=x\cup\{x\}}{%
      \FormulaRefAuto{x\in\mathbb{N}\vdash \Succ(x)=x\cup\{x\}}{1}}
    \proofstepwidestar[1,2]{m\in x\lor m=x}{%
      \FormulaRefAuto{n\in\mathbb{N}\dsep x\in\Succ(n)\vdash x\in n\lor x=n}{1,2}}

    \proofstepwidestar[5]{m\in x}{\rA}
    \proofstepwidestar[1,5]{\Succ(m)=x \lor \Succ(m)\in x}{IV(1,5)}

    \proofstepwidestar[7]{\Succ(m)=x}{\rA}
    \proofstepwidestar{x\in\{x\}}{\FormulaRefAuto{a\in\{a\}}}
    \proofstepwidestar{ x\in x\cup\{x\} }{%
      \FormulaRefAuto{u\in B \vdash u\in A\cup B}{8}}
    \proofstepwidestar[1]{x\in\Succ(x)}{\rIE{3,9}}
    \proofstepwidestar[7]{x=\Succ(m)}{\FormulaRefAuto{a=b\vdash b=a}{7}}
    \proofstepwidestar[1,7]{\Succ(m)\in\Succ(x)}{\rIE{11,10}}
    \proofstepwidestar[1,7]{\Succ(m)=\Succ(x)\lor \Succ(m)\in\Succ(x)}{\rOIb{12}}

    \proofstepwidestar[14]{\Succ(m)\in x}{\rA}
    \proofstepwidestar[14]{\Succ(m)\in x\cup\{x\}}{%
      \FormulaRefAuto{z \in A \vdash z \in A \cup B}{14}}
    \proofstepwidestar[1,14]{\Succ(m)\in\Succ(x)}{\rIE{3,15}}
    \proofstepwidestar[1,14]{\Succ(m)=\Succ(x)\lor \Succ(m)\in\Succ(x)}{\rOIb{16}}

    \proofstepwidestar[1,5]{\Succ(m)=\Succ(x)\lor \Succ(m)\in\Succ(x)}{\rOE{6,7,13,14,17}}

    \proofstepwidestar[19]{m=x}{\rA}
    \proofstepwidestar[19]{x=m}{\FormulaRefAuto{a=b\vdash b=a}{19}}
    \proofstepwidestar{\Succ(x)=\Succ(x)}{\rII}
    \proofstepwidestar[19]{\Succ(m)=\Succ(x)}{\rIE{20,21}}
    \proofstepwidestar[19]{\Succ(m)=\Succ(x)\lor \Succ(m)\in\Succ(x)}{\rOIa{22}}

    \proofstepwidestar[1,2]{\Succ(m)=\Succ(x)\lor \Succ(m)\in\Succ(x)}{\rOE{4,5,18,19,23}}
    \proofstepwidestar[1]{%
          \begin{aligned}[t]
            m\in\Succ(x)\rightarrow {}\\
            \bigl(\Succ(m)=\Succ(x)\lor \Succ(m)\in\Succ(x)\bigr)
          \end{aligned}%
        }{\rRI{2,24}}
  \closeproofpartwideindDelta

  \proofpartwide{Induktionsschluss}
    \proofstepwidestar[1]{m\in\mathbb{N}}{\rA}
    \proofstepwidestar[2]{n\in\mathbb{N}}{\rA}
    \proofstepwidestar[3]{m\in n}{\rA}

    \proofstepwidestar[2]{m\in n \rightarrow \bigl(\Succ(m)=n \lor \Succ(m)\in n\bigr)}{%
      \FormulaRefAuto{PeanoInductionRule}{IA,IS,2}}

    \proofstepwidestar[2,3]{\Succ(m)=n \lor \Succ(m)\in n}{\rRE{4,3}}
  \closeproofpartwide
\end{tabproofsplitwide}


\FormulaThmDelta{%
  m\in\mathbb{N}\dsep n\in\mathbb{N}\dsep m\in n\vdash \Succ(m)\subseteq n
}{%
  \DeltaDecl{Mengen}{m\dsep n}%
}
\begin{tabproof}
  \proofstep{1}{m\in\mathbb{N}}{\rA}
  \proofstep{2}{n\in\mathbb{N}}{\rA}
  \proofstep{3}{m\in n}{\rA}
  \proofstep{1,2,3}{\Succ(m)=n\lor \Succ(m)\in n}{%
    \FormulaRefAuto{m\in\mathbb{N}\dsep n\in\mathbb{N}\dsep m\in n\vdash \Succ(m)=n\lor \Succ(m)\in n}{1,2,3}}
  \proofstep{5}{\Succ(m)=n}{\rA}
  \proofstep{5}{\Succ(m)\subseteq n}{\FormulaRefAuto{A=B\vdash A\subseteq B}{5}}
  \proofstep{7}{\Succ(m)\in n}{\rA}
  \proofstep{1,2,7}{\Succ(m)\subseteq n}{%
    \FormulaRefAuto{m\in\mathbb{N}\dsep n\in\mathbb{N}\dsep m\in n \vdash m\subseteq n}{%
      \FormulaRefAuto{m\in\mathbb{N}\vdash \Succ(m)\in\mathbb{N}}{1},2,7}}
  \proofstep{1,2,3}{\Succ(m)\subseteq n}{\rOE{4,5,6,7,8}}
\end{tabproof}


\FormulaThmDelta{%
  m\in\mathbb{N}\dsep n\in\mathbb{N}\vdash m\in n\lor n\in m\lor m=n%
}{%
  \DeltaDecl{Mengen}{m\dsep n}%
}
\begin{tabproofsplitwide}
  \proofpartwideind[Induktionsanfang]{%
    m\in\mathbb{N}\vdash m\in0 \lor 0\in m\lor m=0}
    \proofstepwidestar[1]{m\in\mathbb{N}}{\rA}
    \proofstepwidestar[1]{0=m \lor 0\in m}{%
      \FormulaRefAuto{n\in\mathbb{N}\vdash \varnothing = n \lor \varnothing \in n}{1}}

    \proofstepwidestar[3]{0=m}{\rA}
    \proofstepwidestar[3]{m=0}{\FormulaRefAuto{a=b\vdash b=a}{3}}
    \proofstepwidestar[3]{m\in0 \lor 0\in m\lor m=0}{%
      \FormulaRefAuto{R\vdash P\lor Q\lor R}{4}}

    \proofstepwidestar[6]{0\in m}{\rA}
    \proofstepwidestar[6]{m\in0 \lor 0\in m\lor m=0}{%
      \FormulaRefAuto{Q\vdash P\lor Q\lor R}{6}}

    \proofstepwidestar[1]{m\in0 \lor 0\in m\lor m=0}{\rOE{2,3,5,6,7}}
  \closeproofpartwideind

  \proofpartwideindDelta[Induktionsschritt]{%
    m\in\mathbb{N}\dsep x\in\mathbb{N}
    \vdash
    \begin{aligned}[t]
      m\in\Succ(x)\lor {}\\
      \Succ(x)\in m\lor m=\Succ(x)
    \end{aligned}%
  }{%
    \DeltaPrem{Induktionsvoraussetzung}{m\in\mathbb{N}\dsep x\in\mathbb{N}\vdash m\in x\lor x\in m\lor m=x}
  }[m\in\mathbb{N}\dsep x\in\mathbb{N}\vdash m\in\Succ(x)\lor \Succ(x)\in m\lor m=\Succ(x)]

    \proofstepwidestar[1]{m\in\mathbb{N}}{\rA}
    \proofstepwidestar[2]{x\in\mathbb{N}}{\rA}
    \proofstepwidestar[1,2]{m\in x\lor x\in m\lor m=x}{IV(1,2)}

    \proofstepwidestar[2]{\Succ(x)=x\cup\{x\}}{%
      \FormulaRefAuto{x\in\mathbb{N}\vdash \Succ(x)=x\cup\{x\}}{2}}

    \proofstepwidestar[5]{m\in x}{\rA}
    \proofstepwidestar[5]{m\in x\cup\{x\}}{%
      \FormulaRefAuto{u\in A \vdash u\in A \cup B}{5}}
    \proofstepwidestar[2,5]{m\in\Succ(x)}{\rIE{4,6}}
    \proofstepwidestar[2,5]{%
      \begin{aligned}[t]
        m\in\Succ(x)\lor {}\\
        \Succ(x)\in m\lor m=\Succ(x)
      \end{aligned}%
    }{%
      \FormulaRefAuto{P\vdash P\lor Q\lor R}{7}}

    \proofstepwidestar[9]{x\in m}{\rA}
    \proofstepwidestar[1,2,9]{\Succ(x)=m \lor \Succ(x)\in m}{%
      \FormulaRefAuto{m\in\mathbb{N}\dsep n\in\mathbb{N}\dsep m\in n\vdash \Succ(m)=n\lor \Succ(m)\in n}{2,1,9}}

    \proofstepwidestar[11]{\Succ(x)=m}{\rA}
    \proofstepwidestar[11]{m=\Succ(x)}{\FormulaRefAuto{a=b\vdash b=a}{11}}
    \proofstepwidestar[11]{%
      \begin{aligned}[t]
        m\in\Succ(x)\lor {}\\
        \Succ(x)\in m\lor m=\Succ(x)
      \end{aligned}%
    }{%
      \FormulaRefAuto{R\vdash P\lor Q\lor R}{12}}

    \proofstepwidestar[14]{\Succ(x)\in m}{\rA}
    \proofstepwidestar[14]{%
      \begin{aligned}[t]
        m\in\Succ(x)\lor {}\\
        \Succ(x)\in m\lor m=\Succ(x)
      \end{aligned}%
    }{%
      \FormulaRefAuto{Q\vdash P\lor Q\lor R}{14}}

    \proofstepwidestar[1,2,9]{%
      \begin{aligned}[t]
        m\in\Succ(x)\lor {}\\
        \Succ(x)\in m\lor m=\Succ(x)
      \end{aligned}%
    }{\rOE{10,11,13,14,15}}

    \proofstepwidestar[17]{m=x}{\rA}
    \proofstepwidestar[17]{x=m}{\FormulaRefAuto{a=b\vdash b=a}{17}}
    \proofstepwidestar{x\in\{x\}}{\FormulaRefAuto{a\in\{a\}}}
    \proofstepwidestar{x\in x\cup\{x\}}{%
      \FormulaRefAuto{u\in B \vdash u\in A\cup B}{19}}
    \proofstepwidestar[2]{x\in\Succ(x)}{\rIE{4,20}}
    \proofstepwidestar[2,17]{m\in\Succ(x)}{\rIE{18,21}}
    \proofstepwidestar[2,17]{%
      \begin{aligned}[t]
        m\in\Succ(x)\lor {}\\
        \Succ(x)\in m\lor m=\Succ(x)
      \end{aligned}%
    }{%
      \FormulaRefAuto{P\vdash P\lor Q\lor R}{22}}

    \proofstepwidestar[1,2]{%
      \begin{aligned}[t]
        m\in\Succ(x)\lor {}\\
        \Succ(x)\in m\lor m=\Succ(x)
      \end{aligned}%
    }{%
      \FormulaRefAuto{P_1\lor P_2\lor P_3\dsep [P_1]\vdots R\dsep [P_2]\vdots R\dsep [P_3]\vdots R\vdash R}{%
        3,5,8,9,16,17,23}}
  \closeproofpartwideindDelta

  \proofpartwide{Induktionsschluss}
    \proofstepwidestar[1]{m\in\mathbb{N}}{\rA}
    \proofstepwidestar[2]{n\in\mathbb{N}}{\rA}
    \proofstepwidestar[1,2]{m\in n\lor n\in m\lor m=n}{%
      \FormulaRefAuto{PeanoInductionRule}{IA,IS,2,1}}
  \closeproofpartwide
\end{tabproofsplitwide}



\FormulaThmDelta{%
  m\in\mathbb{N}\dsep n\in\mathbb{N}\dsep \neg(m\subseteq n)\vdash n\in m%
}{%
  \DeltaDecl{Mengen}{m\dsep n}%
}
\begin{tabproof}
  \proofstep{1}{m\in\mathbb{N}}{\rA}
  \proofstep{2}{n\in\mathbb{N}}{\rA}
  \proofstep{3}{\neg(m\subseteq n)}{\rA}
  \proofstep{1,2}{m\in n\lor n\in m\lor m=n}{%
    \FormulaRefAuto{m\in\mathbb{N}\dsep n\in\mathbb{N}\vdash m\in n\lor n\in m\lor m=n}{1,2}}

  \proofstep{5}{m\in n}{\rA}
  \proofstep{1,2,5}{m\subseteq n}{%
    \FormulaRefAuto{m\in\mathbb{N}\dsep n\in\mathbb{N}\dsep m\in n \vdash m\subseteq n}{1,2,5}}
  \proofstep{1,2,3,5}{n\in m}{%
    \FormulaRefAuto{P,\,\neg P \vdash Q}{6,3}}

  \proofstep{8}{n\in m}{\rA}

  \proofstep{9}{m=n}{\rA}
  \proofstep{9}{m\subseteq n}{%
    \FormulaRefAuto{A=B\vdash A\subseteq B}{9}}
  \proofstep{3,9}{n\in m}{%
    \FormulaRefAuto{P,\,\neg P \vdash Q}{10,3}}

  \proofstep{1,2,3}{n\in m}{%
    \FormulaRefAuto{P_1\lor P_2\lor P_3\dsep [P_1]\vdots R\dsep [P_2]\vdots R\dsep [P_3]\vdots R\vdash R}{%
      4,5,7,8,8,9,11}}
\end{tabproof}


\section{Mitgliedschaft und (echte) Teilmengenrelation}


\FormulaThmDelta[Element ist Teilmenge]{%
  m\in\mathbb{N}\dsep n\in\mathbb{N}\dsep m\in n \vdash m\subseteq n%
}{%
  \DeltaDecl{Mengen}{m\dsep n\dsep x}%
}
\begin{tabproof}
  \proofstep{1}{m\in\mathbb{N}}{\rA}
  \proofstep{2}{n\in\mathbb{N}}{\rA}
  \proofstep{3}{m\in n}{\rA}

  \proofstep{4}{x\in m}{\rA}
  \proofstep{3,4}{x\in\bigcup n}{%
    \FormulaRefAuto{B\in A\dsep x\in B\vdash x \in \bigcup A}{3,4}}
  \proofstep{2}{\bigcup n\subseteq n}{%
    \FormulaRefAuto{n\in\mathbb{N}\vdash \bigcup n\subseteq n}{2}}
  \proofstep{2,3,4}{x\in n}{%
    \FormulaRefAuto{A\subseteq B, x\in A\vdash x\in B}{6,5}}

  \proofstep{2,3}{m\subseteq n}{%
    \FormulaRefAuto{A \subseteq B \coloneq \forall x\,(x \in A \rightarrow x \in B)}{\rUI{\rRI{4,7}}}}
\end{tabproof}


\FormulaThmDelta[Teilmenge oder Gleichheit]{%
  m\in\mathbb{N}\dsep n\in\mathbb{N}\dsep m\subseteq n \vdash m\in n\lor m=n%
}{%
  \DeltaDecl{Mengen}{m\dsep n}%
}
\begin{tabproof}
  \proofstep{1}{m\in\mathbb{N}}{\rA}
  \proofstep{2}{n\in\mathbb{N}}{\rA}
  \proofstep{3}{m\subseteq n}{\rA}

  \proofstep{1,2}{m\in n\lor n\in m\lor m=n}{%
    \FormulaRefAuto{m\in\mathbb{N}\dsep n\in\mathbb{N}\vdash m\in n\lor n\in m\lor m=n}{1,2}}

  \proofstep{5}{m\in n}{\rA}
  \proofstep{5}{m\in n\lor m=n}{\rOIa{5}}

  \proofstep{7}{n\in m}{\rA}
  \proofstep{1,2,7}{n\subseteq m}{%
    \FormulaRefAuto{m\in\mathbb{N}\dsep n\in\mathbb{N}\dsep m\in n \vdash m\subseteq n}{2,1,7}}
  \proofstep{1,2,3,7}{m=n}{%
    \FormulaRefAuto{A\subseteq B, B\subseteq A\vdash A=B}{3,8}}
  \proofstep{1,2,3,7}{m\in n\lor m=n}{\rOIb{9}}

  \proofstep{11}{m=n}{\rA}
  \proofstep{11}{m\in n\lor m=n}{\rOIb{11}}

  \proofstep{1,2,3}{m\in n\lor m=n}{%
    \FormulaRefAuto{P_1\lor P_2\lor P_3\dsep [P_1]\vdots R\dsep [P_2]\vdots R\dsep [P_3]\vdots R\vdash R}{4,5,6,7,10,11,12}}
\end{tabproof}

\FormulaThmDelta[Echte Teilmenge ist Element]{%
  m\in\mathbb{N}\dsep n\in\mathbb{N}\dsep m\subseteq n\dsep m\neq n \vdash m\in n%
}{%
  \DeltaDecl{Mengen}{m\dsep n}%
}
\begin{tabproof}
  \proofstep{1}{m\in\mathbb{N}}{\rA}
  \proofstep{2}{n\in\mathbb{N}}{\rA}
  \proofstep{3}{m\subseteq n}{\rA}
  \proofstep{4}{m\neq n}{\rA}

  \proofstep{1,2,3}{m\in n\lor m=n}{%
    \FormulaRefAuto{m\in\mathbb{N}\dsep n\in\mathbb{N}\dsep m\subseteq n \vdash m\in n\lor m=n}{1,2,3}}

  \proofstep{1,2,3,4}{m\in n}{%
    \FormulaRefAuto{P\lor Q,\,\neg Q \vdash P}{5,4}}
\end{tabproof}

% Hilfssatz: aus m∈n folgt m≠n (genauer: ¬(m=n))
\FormulaThmDelta{%
  m\in\mathbb{N}\dsep n\in\mathbb{N}\dsep m\in n \vdash m\neq n%
}{%
  \DeltaDecl{Mengen}{m\dsep n}%
}
\begin{tabproof}
  \proofstep{1}{m\in\mathbb{N}}{\rA}
  \proofstep{2}{n\in\mathbb{N}}{\rA}
  \proofstep{3}{m\in n}{\rA}

  \proofstep{4}{m=n}{\rA}
  \proofstep{3,4}{m\in m}{\rIE{4,3}}
  \proofstep{1}{m\notin m}{\FormulaRefAuto{n\in\mathbb{N}\vdash n\notin n}{1}}
  \proofstep{1,3,4}{\bot}{\rBI{5,6}}
  \proofstep{1,2,3}{m\neq n}{\rCI{4,7}}
\end{tabproof}

% Hauptsatz: m∈n ⇒ m⊂n
\FormulaThmDelta{%
  m\in\mathbb{N}\dsep n\in\mathbb{N}\dsep m\in n \vdash m\subset n%
}{%
  \DeltaDecl{Mengen}{m\dsep n}%
}
\begin{tabproof}
  \proofstep{1}{m\in\mathbb{N}}{\rA}
  \proofstep{2}{n\in\mathbb{N}}{\rA}
  \proofstep{3}{m\in n}{\rA}

  \proofstep{1,2,3}{m\subseteq n}{%
    \FormulaRefAuto{m\in\mathbb{N}\dsep n\in\mathbb{N}\dsep m\in n \vdash m\subseteq n}{1,2,3}}
  \proofstep{1,2,3}{m\neq n}{%
    \FormulaRefAuto{m\in\mathbb{N}\dsep n\in\mathbb{N}\dsep m\in n \vdash m\neq n}{1,2,3}}
  \proofstep{1,2,3}{m\subseteq n \land m\neq n}{\rAI{4,5}}
  \proofstep{1,2,3}{m\subset n}{%
    \FormulaRefAuto{A \subset B \coloneq A\subseteq B \land A\neq B}{6}}
\end{tabproof}

\FormulaThmDelta{%
  m\in\mathbb{N}\dsep n\in\mathbb{N}\dsep m\subset n \vdash m\in n%
}{%
  \DeltaDecl{Mengen}{m\dsep n}%
}
\begin{tabproof}
  \proofstep{1}{m\in\mathbb{N}}{\rA}
  \proofstep{2}{n\in\mathbb{N}}{\rA}
  \proofstep{3}{m\subset n}{\rA}

  \proofstep{3}{m\subseteq n \land m\neq n}{%
    \FormulaRefAuto{A \subset B \coloneq A\subseteq B \land A\neq B}{3}}
  \proofstep{3}{m\subseteq n}{\rAEa{4}}
  \proofstep{3}{m\neq n}{\rAEb{4}}

  \proofstep{1,2,3}{m\in n}{%
    \FormulaRefAuto{m\in\mathbb{N}\dsep n\in\mathbb{N}\dsep m\subseteq n\dsep m\neq n \vdash m\in n}{1,2,5,6}}
\end{tabproof}

\FormulaThmDelta[Mitgliedschaft genau echte Teilmenge]{%
  m\in\mathbb{N}\dsep n\in\mathbb{N}\vdash m\in n \leftrightarrow m\subset n%
}{%
  \DeltaDecl{Mengen}{m\dsep n}%
}
\begin{tabproof}
  \proofstep{1}{m\in\mathbb{N}}{\rA}
  \proofstep{2}{n\in\mathbb{N}}{\rA}

  \proofstep{3}{m\in n}{\rA}
  \proofstep{1,2,3}{m\subset n}{%
    \FormulaRefAuto{m\in\mathbb{N}\dsep n\in\mathbb{N}\dsep m\in n \vdash m\subset n}{1,2,3}}
  \proofstep{1,2}{m\in n \rightarrow m\subset n}{\rRI{3,4}}

  \proofstep{6}{m\subset n}{\rA}
  \proofstep{1,2,6}{m\in n}{%
    \FormulaRefAuto{m\in\mathbb{N}\dsep n\in\mathbb{N}\dsep m\subset n \vdash m\in n}{1,2,6}}
  \proofstep{1,2}{m\subset n \rightarrow m\in n}{\rRI{6,7}}

  \proofstep{1,2}{m\in n \leftrightarrow m\subset n}{\rLRI{5,8}}
\end{tabproof}


\section{Von-Neumann-Vergleich der natürlichen Zahlen}

Der folgende Block enthält Hilfsaussagen zur von-Neumann-Darstellung der
natürlichen Zahlen, also zu Mitgliedschaft und Teilmengenrelation zwischen
natürlichen Zahlen. Die additive natürliche Ordnung wird später in diesem
Band auf Grundlage der allgemeinen Ordnungsbegriffe aus Band~06 eingeführt.

\FormulaDefDelta[Strikte Ordnungsrelation der natürlichen Zahlen]{%
  <_{\mathbb{N}}\coloneqq
  \{(m,n)\in \mathbb{N}\times\mathbb{N}\mid m\in n\}%
}{%
  \DeltaRow{Mengen}{m\dsep n}%
  \DeltaRow{Relationen}{<_{\mathbb{N}}}%
}

\FormulaDefDelta[Notation der strikten natürlichen Ordnung]{%
  m\in\mathbb{N}\dsep n\in\mathbb{N}\vdash
  m < n \coloneqq m\in n%
}{%
  \DeltaRow{Mengen}{m\dsep n}%
  \DeltaRow{Relationen}{<_{\mathbb{N}}}%
  \DeltaRow{Neue Symbole}{<}%
}

\FormulaThmDelta[Mitgliedschaftscharakterisierung der strikten natürlichen Ordnung]{%
  m\in\mathbb{N}\dsep n\in\mathbb{N}\vdash
  m < n \leftrightarrow m\in n%
}{%
  \DeltaRow{Mengen}{m\dsep n}%
  \DeltaRow{Relationen}{<_{\mathbb{N}}}%
}
\begin{tabproof}
  \proofstep{1}{m\in\mathbb{N}}{\rA}
  \proofstep{2}{n\in\mathbb{N}}{\rA}
  \proofstep{1,2}{m < n \leftrightarrow m\in n}{%
    \FormulaRefAuto{m\in\mathbb{N}\dsep n\in\mathbb{N}\vdash m < n \coloneqq m\in n}{1,2}}
\end{tabproof}

\FormulaThmDelta[Charakterisierung der strikten natürlichen Ordnung]{%
  m\in\mathbb{N}\dsep n\in\mathbb{N}\vdash
  m < n \leftrightarrow m\subset n%
}{%
  \DeltaRow{Mengen}{m\dsep n}%
  \DeltaRow{Relationen}{<_{\mathbb{N}}}%
}
\begin{tabproof}
  \proofstep{1}{m\in\mathbb{N}}{\rA}
  \proofstep{2}{n\in\mathbb{N}}{\rA}
  \proofstep{1,2}{m < n \leftrightarrow m\in n}{%
    \FormulaRefAuto{m\in\mathbb{N}\dsep n\in\mathbb{N}\vdash m < n \leftrightarrow m\in n}{1,2}}
  \proofstep{1,2}{m\in n \leftrightarrow m\subset n}{%
    \FormulaRefAuto{m\in\mathbb{N}\dsep n\in\mathbb{N}\vdash m\in n \leftrightarrow m\subset n}{1,2}}
  \proofstep{1,2}{m < n \leftrightarrow m\subset n}{%
    \FormulaRefAuto{P\leftrightarrow Q\dsep Q\leftrightarrow R\vdash P\leftrightarrow R}{3,4}}
\end{tabproof}

\FormulaThmDelta[Einführung in die strikte natürliche Ordnungsrelation]{%
  m\in\mathbb{N}\dsep n\in\mathbb{N}\dsep m < n
  \vdash (m,n)\in <_{\mathbb{N}}%
}{%
  \DeltaRow{Mengen}{m\dsep n}%
  \DeltaRow{Relationen}{<_{\mathbb{N}}}%
}
\begin{tabproof}
  \proofstep{1}{m\in\mathbb{N}}{\rA}
  \proofstep{2}{n\in\mathbb{N}}{\rA}
  \proofstep{3}{m < n}{\rA}
  \proofstep{1,2,3}{m\in n}{%
    \FormulaRefAuto{m\in\mathbb{N}\dsep n\in\mathbb{N}\vdash m < n \leftrightarrow m\in n}{1,2,3}}
  \proofstep{1,2}{(m,n)\in \mathbb{N}\times\mathbb{N}}{%
    \FormulaRefAuto{a\in A\dsep b\in B\vdash (a,b)\in A\times B}{1,2}}
  \proofstep{1,2,3}{(m,n)\in \{(u,v)\in \mathbb{N}\times\mathbb{N}\mid u\in v\}}{%
    \FormulaRefAuto{x\in A \dsep P(x) \vdash x\in \{\,x\in A \mid P(x)\,\}}{5,4}}
  \proofstep{1,2,3}{(m,n)\in <_{\mathbb{N}}}{%
    \rIE{\FormulaRefAuto{<_{\mathbb{N}}\coloneqq \{(m,n)\in \mathbb{N}\times\mathbb{N}\mid m\in n\}},6}}
\end{tabproof}

\FormulaThmDelta[Elimination der strikten natürlichen Ordnungsrelation]{%
  (m,n)\in <_{\mathbb{N}}\vdash
  m\in\mathbb{N}\land n\in\mathbb{N}\land m < n%
}{%
  \DeltaRow{Mengen}{m\dsep n}%
  \DeltaRow{Relationen}{<_{\mathbb{N}}}%
}
\begin{tabproof}
  \proofstep{1}{(m,n)\in <_{\mathbb{N}}}{\rA}
  \proofstep{}{<_{\mathbb{N}}=\{(u,v)\in \mathbb{N}\times\mathbb{N}\mid u\in v\}}{%
    \FormulaRefAuto{<_{\mathbb{N}}\coloneqq \{(m,n)\in \mathbb{N}\times\mathbb{N}\mid m\in n\}}}
  \proofstep{1}{(m,n)\in \{(u,v)\in \mathbb{N}\times\mathbb{N}\mid u\in v\}}{\rIE{2,1}}
  \proofstep{1}{(m,n)\in \mathbb{N}\times\mathbb{N}\land m\in n}{%
    \FormulaRefAuto{x\in\{u\in A\mid P(u)\}\eqvdash x\in A\land P(x)}{3}}
  \proofstep{1}{(m,n)\in \mathbb{N}\times\mathbb{N}}{\rAEa{4}}
  \proofstep{1}{m\in n}{\rAEb{4}}
  \proofstep{1}{m\in\mathbb{N}\land n\in\mathbb{N}}{%
    \FormulaRefAuto{(a,b)\in A\times B \eqvdash a\in A \land b\in B}{5}}
  \proofstep{1}{m\in\mathbb{N}}{\rAEa{7}}
  \proofstep{1}{n\in\mathbb{N}}{\rAEb{7}}
  \proofstep{1}{m < n}{%
    \FormulaRefAuto{m\in\mathbb{N}\dsep n\in\mathbb{N}\vdash m < n \leftrightarrow m\in n}{8,9,6}}
  \proofstep{1}{m\in\mathbb{N}\land n\in\mathbb{N}\land m < n}{%
    \FormulaRefAuto{P\dsep Q\dsep R\vdash P\land Q\land R}{8,9,10}}
\end{tabproof}

\FormulaDefDelta[Ordnungsrelation der natürlichen Zahlen]{%
  \leq_{\mathbb{N}}\coloneqq
  \{(m,n)\in \mathbb{N}\times\mathbb{N}\mid m\subseteq n\}%
}{%
  \DeltaRow{Mengen}{m\dsep n}%
  \DeltaRow{Relationen}{\leq_{\mathbb{N}}}%
}

\FormulaDefDelta[Notation der natürlichen Ordnung]{%
  m\in\mathbb{N}\dsep n\in\mathbb{N}\vdash
  m\leq n \coloneqq m\subseteq n%
}{%
  \DeltaRow{Mengen}{m\dsep n}%
  \DeltaRow{Relationen}{\leq_{\mathbb{N}}}%
  \DeltaRow{Neue Symbole}{\leq}%
}
\begin{remark}
  Die Schreibweise \(m < n\) bezeichnet auf \(\mathbb{N}\) die strikte Ordnung. Nach dem vorherigen Satz stimmt sie dort mit der echten Teilmengenrelation \(m\subset n\) überein. Die Schreibweise \(m\leq n\) bezeichnet die zugehörige nicht-strikte Ordnung.
\end{remark}

\FormulaThmDelta[Charakterisierung der natürlichen Ordnung]{%
  m\in\mathbb{N}\dsep n\in\mathbb{N}\vdash
  m\leq n \leftrightarrow m\subseteq n%
}{%
  \DeltaRow{Mengen}{m\dsep n}%
  \DeltaRow{Relationen}{\leq_{\mathbb{N}}}%
}
\begin{tabproof}
  \proofstep{1}{m\in\mathbb{N}}{\rA}
  \proofstep{2}{n\in\mathbb{N}}{\rA}
  \proofstep{1,2}{m\leq n \leftrightarrow m\subseteq n}{%
    \FormulaRefAuto{m\in\mathbb{N}\dsep n\in\mathbb{N}\vdash m\leq n \coloneqq m\subseteq n}{1,2}}
\end{tabproof}

\FormulaThmDelta[Reflexivität der natürlichen Ordnung]{%
  n\in\mathbb{N}\vdash n\leq n%
}{%
  \DeltaRow{Mengen}{n}%
  \DeltaRow{Relationen}{\leq_{\mathbb{N}}}%
}
\begin{tabproof}
  \proofstep{1}{n\in\mathbb{N}}{\rA}
  \proofstep{}{\NatSeg{n}\subseteq \NatSeg{n}}{\FormulaRefAuto{A\subseteq A}}
  \proofstep{1}{n\leq n}{%
    \FormulaRefAuto{m\in\mathbb{N}\dsep n\in\mathbb{N}\vdash
    m\leq n \leftrightarrow m\subseteq n}{1,1,2}}
\end{tabproof}

\FormulaThmDelta[Einführung in die natürliche Ordnungsrelation]{%
  m\in\mathbb{N}\dsep n\in\mathbb{N}\dsep m\leq n
  \vdash (m,n)\in \leq_{\mathbb{N}}%
}{%
  \DeltaRow{Mengen}{m\dsep n}%
  \DeltaRow{Relationen}{\leq_{\mathbb{N}}}%
}
\begin{tabproof}
  \proofstep{1}{m\in\mathbb{N}}{\rA}
  \proofstep{2}{n\in\mathbb{N}}{\rA}
  \proofstep{3}{m\leq n}{\rA}
  \proofstep{1,2,3}{m\subseteq n}{%
    \FormulaRefAuto{m\in\mathbb{N}\dsep n\in\mathbb{N}\vdash m\leq n \leftrightarrow m\subseteq n}{1,2,3}}
  \proofstep{1,2}{(m,n)\in \mathbb{N}\times\mathbb{N}}{%
    \FormulaRefAuto{a\in A\dsep b\in B\vdash (a,b)\in A\times B}{1,2}}
  \proofstep{1,2,3}{(m,n)\in \{(u,v)\in \mathbb{N}\times\mathbb{N}\mid u\subseteq v\}}{%
    \FormulaRefAuto{x\in A \dsep P(x) \vdash x\in \{\,x\in A \mid P(x)\,\}}{5,4}}
  \proofstep{1,2,3}{(m,n)\in \leq_{\mathbb{N}}}{%
    \rIE{\FormulaRefAuto{\leq_{\mathbb{N}}\coloneqq \{(m,n)\in \mathbb{N}\times\mathbb{N}\mid m\subseteq n\}},6}}
\end{tabproof}

\FormulaThmDelta[Elimination der natürlichen Ordnungsrelation]{%
  (m,n)\in \leq_{\mathbb{N}}\vdash
  m\in\mathbb{N}\land n\in\mathbb{N}\land m\leq n%
}{%
  \DeltaRow{Mengen}{m\dsep n}%
  \DeltaRow{Relationen}{\leq_{\mathbb{N}}}%
}
\begin{tabproof}
  \proofstep{1}{(m,n)\in \leq_{\mathbb{N}}}{\rA}
  \proofstep{}{\leq_{\mathbb{N}}=\{(u,v)\in \mathbb{N}\times\mathbb{N}\mid u\subseteq v\}}{%
    \FormulaRefAuto{\leq_{\mathbb{N}}\coloneqq \{(m,n)\in \mathbb{N}\times\mathbb{N}\mid m\subseteq n\}}}
  \proofstep{1}{(m,n)\in \{(u,v)\in \mathbb{N}\times\mathbb{N}\mid u\subseteq v\}}{\rIE{2,1}}
  \proofstep{1}{(m,n)\in \mathbb{N}\times\mathbb{N}\land m\subseteq n}{%
    \FormulaRefAuto{x\in\{u\in A\mid P(u)\}\eqvdash x\in A\land P(x)}{3}}
  \proofstep{1}{(m,n)\in \mathbb{N}\times\mathbb{N}}{\rAEa{4}}
  \proofstep{1}{m\subseteq n}{\rAEb{4}}
  \proofstep{1}{m\in\mathbb{N}\land n\in\mathbb{N}}{%
    \FormulaRefAuto{(a,b)\in A\times B \eqvdash a\in A \land b\in B}{5}}
  \proofstep{1}{m\in\mathbb{N}}{\rAEa{7}}
  \proofstep{1}{n\in\mathbb{N}}{\rAEb{7}}
  \proofstep{1}{m\leq n}{%
    \FormulaRefAuto{m\in\mathbb{N}\dsep n\in\mathbb{N}\vdash m\leq n \leftrightarrow m\subseteq n}{8,9,6}}
  \proofstep{1}{m\in\mathbb{N}\land n\in\mathbb{N}\land m\leq n}{%
    \FormulaRefAuto{P\dsep Q\dsep R\vdash P\land Q\land R}{8,9,10}}
\end{tabproof}

\begin{remark}
  Die folgenden Sätze sind Ordnungsversionen bereits bewiesener Mitgliedschafts- und Teilmengenformulierungen. Die ursprünglichen Formulierungen bleiben als Begründungen erhalten.
\end{remark}

\FormulaThmDelta[Nachfolger eines kleineren Elements]{%
  m\in\mathbb{N}\dsep n\in\mathbb{N}\dsep m < n
  \vdash \Succ(m)=n\lor \Succ(m) < n%
}{%
  \DeltaRow{Mengen}{m\dsep n}%
  \DeltaRow{Relationen}{<_{\mathbb{N}}}%
}
\begin{tabproof}
  \proofstep{1}{m\in\mathbb{N}}{\rA}
  \proofstep{2}{n\in\mathbb{N}}{\rA}
  \proofstep{3}{m < n}{\rA}
  \proofstep{1,2,3}{m\in n}{%
    \FormulaRefAuto{m\in\mathbb{N}\dsep n\in\mathbb{N}\vdash m < n \leftrightarrow m\in n}{1,2,3}}
  \proofstep{1,2,3}{\Succ(m)=n\lor \Succ(m)\in n}{%
    \FormulaRefAuto{m\in\mathbb{N}\dsep n\in\mathbb{N}\dsep m\in n\vdash \Succ(m)=n\lor \Succ(m)\in n}{1,2,4}}
  \proofstep{1}{\Succ(m)\in\mathbb{N}}{%
    \FormulaRefAuto{m\in\mathbb{N}\vdash \Succ(m)\in\mathbb{N}}{1}}
  \proofstep{7}{\Succ(m)=n}{\rA}
  \proofstep{7}{\Succ(m)=n\lor \Succ(m) < n}{\rOIa{7}}
  \proofstep{9}{\Succ(m)\in n}{\rA}
  \proofstep{1,2,9}{\Succ(m) < n}{%
    \FormulaRefAuto{m\in\mathbb{N}\dsep n\in\mathbb{N}\vdash m < n \leftrightarrow m\in n}{6,2,9}}
  \proofstep{1,2,9}{\Succ(m)=n\lor \Succ(m) < n}{\rOIb{10}}
  \proofstep{1,2,3}{\Succ(m)=n\lor \Succ(m) < n}{\rOE{5,7,8,9,11}}
\end{tabproof}

\FormulaThmDelta[Nachfolger eines kleineren Elements ist kleiner oder gleich]{%
  m\in\mathbb{N}\dsep n\in\mathbb{N}\dsep m < n
  \vdash \Succ(m)\leq n%
}{%
  \DeltaRow{Mengen}{m\dsep n}%
  \DeltaRow{Relationen}{<_{\mathbb{N}}\dsep \leq_{\mathbb{N}}}%
}
\begin{tabproof}
  \proofstep{1}{m\in\mathbb{N}}{\rA}
  \proofstep{2}{n\in\mathbb{N}}{\rA}
  \proofstep{3}{m < n}{\rA}
  \proofstep{1,2,3}{m\in n}{%
    \FormulaRefAuto{m\in\mathbb{N}\dsep n\in\mathbb{N}\vdash m < n \leftrightarrow m\in n}{1,2,3}}
  \proofstep{1,2,3}{\Succ(m)\subseteq n}{%
    \FormulaRefAuto{m\in\mathbb{N}\dsep n\in\mathbb{N}\dsep m\in n\vdash \Succ(m)\subseteq n}{1,2,4}}
  \proofstep{1}{\Succ(m)\in\mathbb{N}}{%
    \FormulaRefAuto{m\in\mathbb{N}\vdash \Succ(m)\in\mathbb{N}}{1}}
  \proofstep{1,2,3}{\Succ(m)\leq n}{%
    \FormulaRefAuto{m\in\mathbb{N}\dsep n\in\mathbb{N}\vdash m\leq n \leftrightarrow m\subseteq n}{6,2,5}}
\end{tabproof}

\FormulaThmDelta[Element ist strikt kleiner als sein Nachfolger]{%
  n\in\mathbb{N}\vdash n < \Succ(n)%
}{%
  \DeltaRow{Mengen}{n}%
  \DeltaRow{Relationen}{<_{\mathbb{N}}}%
}
\begin{tabproof}
  \proofstep{1}{n\in\mathbb{N}}{\rA}
  \proofstep{1}{\Succ(n)\in\mathbb{N}}{%
    \FormulaRefAuto{PeanoSuccClosure}{1}}
  \proofstep{1}{n\in\Succ(n)}{%
    \FormulaRefAuto{n\in\mathbb{N}\vdash n\in\Succ(n)}{1}}
  \proofstep{1}{n < \Succ(n)}{%
    \FormulaRefAuto{m\in\mathbb{N}\dsep n\in\mathbb{N}\vdash m < n \leftrightarrow m\in n}{1,2,3}}
\end{tabproof}

\FormulaThmDelta[Element ist kleiner oder gleich seinem Nachfolger]{%
  n\in\mathbb{N}\vdash n\leq \Succ(n)%
}{%
  \DeltaRow{Mengen}{n}%
  \DeltaRow{Relationen}{\leq_{\mathbb{N}}}%
}
\begin{tabproof}
  \proofstep{1}{n\in\mathbb{N}}{\rA}
  \proofstep{1}{\Succ(n)\in\mathbb{N}}{%
    \FormulaRefAuto{PeanoSuccClosure}{1}}
  \proofstep{1}{n\subseteq \Succ(n)}{%
    \FormulaRefAuto{n\in\mathbb{N}\vdash n\subseteq \Succ(n)}{1}}
  \proofstep{1}{n\leq \Succ(n)}{%
    \FormulaRefAuto{m\in\mathbb{N}\dsep n\in\mathbb{N}\vdash m\leq n \leftrightarrow m\subseteq n}{1,2,3}}
\end{tabproof}

\FormulaThmDelta[Elemente einer Nachfolgermenge sind genau die kleineren oder gleichen]{%
  x\in\mathbb{N}\dsep n\in\mathbb{N}\vdash x\in\Succ(n)\leftrightarrow x\leq n%
}{%
  \DeltaRow{Mengen}{x\dsep n}%
  \DeltaRow{Relationen}{\leq_{\mathbb{N}}\dsep <_{\mathbb{N}}}%
}
\begin{tabproof}
  \proofstep{1}{x\in\mathbb{N}}{\rA}
  \proofstep{2}{n\in\mathbb{N}}{\rA}

  \proofstep{3}{x\in\Succ(n)}{\rA}
  \proofstep{2,3}{x\in n\lor x=n}{%
    \FormulaRefAuto{n\in\mathbb{N}\dsep x\in\Succ(n)\vdash x\in n\lor x=n}{2,3}}
  \proofstep{5}{x\in n}{\rA}
  \proofstep{1,2,5}{x < n}{%
    \FormulaRefAuto{m\in\mathbb{N}\dsep n\in\mathbb{N}\vdash m < n \leftrightarrow m\in n}{1,2,5}}
  \proofstep{1,2,5}{x\leq n}{%
    \FormulaRefAuto{m\in\mathbb{N}\dsep n\in\mathbb{N}\dsep m < n\vdash m\leq n}{1,2,6}}
  \proofstep{8}{x=n}{\rA}
  \proofstep{8}{x\subseteq n}{\FormulaRefAuto{A=B\vdash A\subseteq B}{8}}
  \proofstep{1,2,8}{x\leq n}{%
    \FormulaRefAuto{m\in\mathbb{N}\dsep n\in\mathbb{N}\vdash m\leq n \leftrightarrow m\subseteq n}{1,2,9}}
  \proofstep{1,2,3}{x\leq n}{\rOE{4,5,7,8,10}}
  \proofstep{1,2}{x\in\Succ(n)\rightarrow x\leq n}{\rRI{3,11}}

  \proofstep{13}{x\leq n}{\rA}
  \proofstep{1,2,13}{x\subseteq n}{%
    \FormulaRefAuto{m\in\mathbb{N}\dsep n\in\mathbb{N}\vdash m\leq n \leftrightarrow m\subseteq n}{1,2,13}}
  \proofstep{1,2,13}{x\in n\lor x=n}{%
    \FormulaRefAuto{m\in\mathbb{N}\dsep n\in\mathbb{N}\dsep m\subseteq n \vdash m\in n\lor m=n}{1,2,14}}
  \proofstep{2}{n\subseteq \Succ(n)}{%
    \FormulaRefAuto{n\in\mathbb{N}\vdash n\subseteq \Succ(n)}{2}}
  \proofstep{17}{x\in n}{\rA}
  \proofstep{2,17}{x\in\Succ(n)}{%
    \FormulaRefAuto{A\subseteq B,\, x\in A \vdash x\in B}{16,17}}
  \proofstep{19}{x=n}{\rA}
  \proofstep{2}{n\in\Succ(n)}{%
    \FormulaRefAuto{n\in\mathbb{N}\vdash n\in\Succ(n)}{2}}
  \proofstep{2,19}{x\in\Succ(n)}{%
    \FormulaRefAuto{a=b \dsep b\in A \vdash a\in A}{19,20}}
  \proofstep{1,2,13}{x\in\Succ(n)}{\rOE{15,17,18,19,21}}
  \proofstep{1,2}{x\leq n\rightarrow x\in\Succ(n)}{\rRI{13,22}}

  \proofstep{1,2}{x\in\Succ(n)\leftrightarrow x\leq n}{\rLRI{12,23}}
\end{tabproof}

\FormulaThmDelta[Natürliche Ordnung schliesst spätere Mitgliedschaft aus]{%
  m\in\mathbb{N}\dsep n\in\mathbb{N}\dsep m\leq n\vdash n\notin m%
}{%
  \DeltaRow{Mengen}{m\dsep n}%
  \DeltaRow{Relationen}{\leq_{\mathbb{N}}}%
}
\begin{tabproof}
  \proofstep{1}{m\in\mathbb{N}}{\rA}
  \proofstep{2}{n\in\mathbb{N}}{\rA}
  \proofstep{3}{m\leq n}{\rA}
  \proofstep{4}{n\in m}{\rA}
  \proofstep{1,2,3}{m\subseteq n}{%
    \FormulaRefAuto{m\in\mathbb{N}\dsep n\in\mathbb{N}\vdash m\leq n \leftrightarrow m\subseteq n}{1,2,3}}
  \proofstep{1,2,3,4}{n\in n}{%
    \FormulaRefAuto{A\subseteq B,\, x\in A \vdash x\in B}{5,4}}
  \proofstep{2}{n\notin n}{%
    \FormulaRefAuto{n\in\mathbb{N}\vdash n\notin n}{2}}
  \proofstep{1,2,3,4}{\bot}{\rBI{6,7}}
  \proofstep{1,2,3}{n\notin m}{\rCI{4,8}}
\end{tabproof}

\FormulaThmDelta[Trichotomie der natürlichen Ordnung]{%
  m\in\mathbb{N}\dsep n\in\mathbb{N}\vdash m < n\lor n < m\lor m=n%
}{%
  \DeltaRow{Mengen}{m\dsep n}%
  \DeltaRow{Relationen}{<_{\mathbb{N}}}%
}
\begin{tabproof}
  \proofstep{1}{m\in\mathbb{N}}{\rA}
  \proofstep{2}{n\in\mathbb{N}}{\rA}
  \proofstep{1,2}{m\in n\lor n\in m\lor m=n}{%
    \FormulaRefAuto{m\in\mathbb{N}\dsep n\in\mathbb{N}\vdash m\in n\lor n\in m\lor m=n}{1,2}}
  \proofstep{4}{m\in n}{\rA}
  \proofstep{1,2,4}{m < n}{%
    \FormulaRefAuto{m\in\mathbb{N}\dsep n\in\mathbb{N}\vdash m < n \leftrightarrow m\in n}{1,2,4}}
  \proofstep{1,2,4}{m < n\lor n < m\lor m=n}{%
    \FormulaRefAuto{P\vdash P\lor Q\lor R}{5}}
  \proofstep{7}{n\in m}{\rA}
  \proofstep{1,2,7}{n < m}{%
    \FormulaRefAuto{m\in\mathbb{N}\dsep n\in\mathbb{N}\vdash m < n \leftrightarrow m\in n}{2,1,7}}
  \proofstep{1,2,7}{m < n\lor n < m\lor m=n}{%
    \FormulaRefAuto{Q\vdash P\lor Q\lor R}{8}}
  \proofstep{10}{m=n}{\rA}
  \proofstep{10}{m < n\lor n < m\lor m=n}{%
    \FormulaRefAuto{R\vdash P\lor Q\lor R}{10}}
  \proofstep{1,2}{m < n\lor n < m\lor m=n}{%
    \FormulaRefAuto{P_1\lor P_2\lor P_3\dsep [P_1]\vdots R\dsep [P_2]\vdots R\dsep [P_3]\vdots R\vdash R}{3,4,6,7,9,10,11}}
\end{tabproof}

\emph{Beweisidee.}
Der folgende allgemeine Hilfssatz kapselt den Ordnungsanteil: Eine
Funktion auf \(\mathbb{N}\), die schon fuer jedes \(k<n\) verschiedene
Werte besitzt, ist injektiv. Die Trichotomie der Ordnung wird damit nur
einmal verwendet.

\FormulaThmDeltaR[Ordnungstrennung auf \(\mathbb{N}\) impliziert Injektivitaet]{%
  \begin{aligned}[t]
    &F\colon \mathbb{N}\to M\dsep
    \forall k,n\in\mathbb{N}\,\bigl(k<n\rightarrow F(k)\neq F(n)\bigr)
    \vdash{}\\
    &F\colon \mathbb{N}\inj M
  \end{aligned}
}{%
  F\colon \mathbb{N}\to M\dsep
  \forall k,n\in\mathbb{N}\,\bigl(k<n\rightarrow F(k)\neq F(n)\bigr)
  \vdash F\colon \mathbb{N}\inj M
}{%
  \DeltaDecl{Mengen}{M\dsep k\dsep n}%
  \DeltaDecl{Funktionen}{F}%
  \DeltaDecl{Relation}{<_{\mathbb{N}}}%
}
\begin{tabproof}
  \proofstep{1}{F\colon \mathbb{N}\to M}{\rA}
  \proofstep{2}{\forall k,n\in\mathbb{N}\,\bigl(k<n\rightarrow F(k)\neq F(n)\bigr)}{\rA}
  \proofstep{3}{k\in\mathbb{N}}{\rA}
  \proofstep{4}{n\in\mathbb{N}}{\rA}
  \proofstep{5}{F(k)=F(n)}{\rA}
  \proofstep{3,4}{k < n\lor n < k\lor k=n}{%
    \FormulaRefAuto{m\in\mathbb{N}\dsep n\in\mathbb{N}\vdash m < n\lor n < m\lor m=n}{3,4}}

  \proofstep{7}{k < n}{\rA}
  \proofstep{2,3,4,7}{F(k)\neq F(n)}{%
    \rRE{7,\rRE{4,\rRE{3,\rUE{2}}}}}
  \proofstep{2,3,4,5,7}{k=n}{%
    \FormulaRefAuto{P,\,\neg P \vdash Q}{5,8}}

  \proofstep{10}{n < k}{\rA}
  \proofstep{2,4,3,10}{F(n)\neq F(k)}{%
    \rRE{10,\rRE{3,\rRE{4,\rUE{2}}}}}
  \proofstep{5}{F(n)=F(k)}{\FormulaRefAuto{a=b\vdash b=a}{5}}
  \proofstep{2,3,4,5,10}{k=n}{%
    \FormulaRefAuto{P,\,\neg P \vdash Q}{12,11}}

  \proofstep{14}{k=n}{\rA}
  \proofstep{14}{k=n}{\FormulaRefAuto{P\vdash P}{14}}

  \proofstep{2,3,4,5}{k=n}{%
    \FormulaRefAuto{P_1\lor P_2\lor P_3\dsep [P_1]\vdots R\dsep [P_2]\vdots R\dsep [P_3]\vdots R\vdash R}{6,7,9,10,13,14,15}}
  \proofstep{1,2}{\forall k,n\in\mathbb{N}\,\bigl(F(k)=F(n)\rightarrow k=n\bigr)}{%
    \rUI{\rRI{3,\rRI{4,\rRI{5,16}}}}}
  \proofstep{1,2}{F\colon \mathbb{N}\inj M}{%
    \FormulaRefAuto{Injektive Funktion}{1,17}}
\end{tabproof}

\emph{Beweisidee.}
Der folgende allgemeine Hilfssatz kapselt den Induktionsanteil fuer
wachsende Folgen von Mengen: Wenn jede Stufe in ihrer Nachfolgerstufe
liegt, dann liegt die Nachfolgerstufe von \(k\) in jeder spaeteren
Stufe \(n\) mit \(k<n\).

\FormulaThmDeltaR[Schrittweise Inklusionen induzieren Monotonie]{%
  \begin{aligned}[t]
    &\forall u\in\mathbb{N}\,\bigl(F(u)\subseteq F(\Succ(u))\bigr)\dsep
    k\in\mathbb{N}\dsep n\in\mathbb{N}\dsep k<n
    \vdash{}\\
    &F(\Succ(k))\subseteq F(n)
  \end{aligned}
}{%
  \forall u\in\mathbb{N}\,\bigl(F(u)\subseteq F(\Succ(u))\bigr)\dsep
  k\in\mathbb{N}\dsep n\in\mathbb{N}\dsep k<n
  \vdash F(\Succ(k))\subseteq F(n)
}{%
  \DeltaDecl{Mengen}{k\dsep n\dsep u\dsep v}%
  \DeltaDecl{Funktionen}{F}%
  \DeltaDecl{Relation}{<_{\mathbb{N}}}%
}
\begin{tabproofsplitwide}
  \proofpartwideindR[Induktionsanfang]{%
    \forall u\in\mathbb{N}\,\bigl(F(u)\subseteq F(\Succ(u))\bigr)\vdash
    \forall k\in\mathbb{N}\,\bigl(k\in0\rightarrow
      F(\Succ(k))\subseteq F(0)\bigr)
  }{%
    \forall u\in\mathbb{N}\,\bigl(F(u)\subseteq F(\Succ(u))\bigr)\vdash
    \forall k\in\mathbb{N}\,\bigl(k\in0\rightarrow
      F(\Succ(k))\subseteq F(0)\bigr)
  }
    \proofstepwidestar[1]{\forall u\in\mathbb{N}\,\bigl(F(u)\subseteq F(\Succ(u))\bigr)}{\rA}
    \proofstepwidestar[2]{k\in\mathbb{N}}{\rA}
    \proofstepwidestar[3]{k\in0}{\rA}
    \proofstepwidestar[]{k\notin0}{\FormulaRefAuto{x\not\in\varnothing}}
    \proofstepwidestar[3]{F(\Succ(k))\subseteq F(0)}{%
      \FormulaRefAuto{P,\,\neg P \vdash Q}{3,4}}
    \proofstepwidestar[2]{k\in0\rightarrow F(\Succ(k))\subseteq F(0)}{\rRI{3,5}}
    \proofstepwidestar[]{%
      \forall k\in\mathbb{N}\,\bigl(k\in0\rightarrow
        F(\Succ(k))\subseteq F(0)\bigr)
    }{\rUI{\rRI{2,6}}}
  \closeproofpartwideindR

  \proofpartwideindR[Induktionsschritt]{%
    \begin{aligned}[t]
      &\forall v\in\mathbb{N}\,\bigl(F(v)\subseteq F(\Succ(v))\bigr)\dsep
      u\in\mathbb{N}\dsep{}\\
      &\forall k\in\mathbb{N}\,\bigl(k\in u\rightarrow
        F(\Succ(k))\subseteq F(u)\bigr)
      \vdash{}\\
      &\forall k\in\mathbb{N}\,\bigl(k\in\Succ(u)\rightarrow
        F(\Succ(k))\subseteq F(\Succ(u))\bigr)
    \end{aligned}
  }{%
    \forall v\in\mathbb{N}\,\bigl(F(v)\subseteq F(\Succ(v))\bigr)\dsep
    u\in\mathbb{N}\dsep
    \forall k\in\mathbb{N}\,\bigl(k\in u\rightarrow
      F(\Succ(k))\subseteq F(u)\bigr)\vdash
    \forall k\in\mathbb{N}\,\bigl(k\in\Succ(u)\rightarrow
      F(\Succ(k))\subseteq F(\Succ(u))\bigr)
  }
    \proofstepwidestar[1]{\forall v\in\mathbb{N}\,\bigl(F(v)\subseteq F(\Succ(v))\bigr)}{\rA}
    \proofstepwidestar[2]{u\in\mathbb{N}}{\rA}
    \proofstepwidestar[3]{%
      \forall k\in\mathbb{N}\,\bigl(k\in u\rightarrow
        F(\Succ(k))\subseteq F(u)\bigr)
    }{\rA}
    \proofstepwidestar[4]{k\in\mathbb{N}}{\rA}
    \proofstepwidestar[5]{k\in\Succ(u)}{\rA}
    \proofstepwidestar[2,5]{k\in u\lor k=u}{%
      \FormulaRefAuto{n\in\mathbb{N}\dsep x\in\Succ(n)\vdash x\in n\lor x=n}{2,5}}

    \proofstepwidestar[7]{k\in u}{\rA}
    \proofstepwidestar[3,4,7]{F(\Succ(k))\subseteq F(u)}{%
      \rRE{7,\rRE{4,\rUE{3}}}}
    \proofstepwidestar[1,2]{F(u)\subseteq F(\Succ(u))}{\rRE{2,\rUE{1}}}
    \proofstepwidestar[1,2,3,4,7]{F(\Succ(k))\subseteq F(\Succ(u))}{%
      \FormulaRefAuto{A\subseteq B, B\subseteq C \vdash A\subseteq C}{8,9}}

    \proofstepwidestar[11]{k=u}{\rA}
    \proofstepwidestar[11]{\Succ(k)=\Succ(u)}{\rIE{11}}
    \proofstepwidestar[11]{F(\Succ(k))=F(\Succ(u))}{\rIE{12}}
    \proofstepwidestar[11]{F(\Succ(k))\subseteq F(\Succ(u))}{%
      \FormulaRefAuto{A=B\vdash A\subseteq B}{13}}

    \proofstepwidestar[1,2,3,4,5]{F(\Succ(k))\subseteq F(\Succ(u))}{%
      \rOE{6,7,10,11,14}}
    \proofstepwidestar[1,2,3,4]{k\in\Succ(u)\rightarrow F(\Succ(k))\subseteq F(\Succ(u))}{%
      \rRI{5,15}}
    \proofstepwidestar[1,2,3]{%
      \forall k\in\mathbb{N}\,\bigl(k\in\Succ(u)\rightarrow
        F(\Succ(k))\subseteq F(\Succ(u))\bigr)
    }{\rUI{\rRI{4,16}}}
  \closeproofpartwideindR

  \proofpartwide{Induktionsschluss}
    \proofstepwidestar[1]{\forall u\in\mathbb{N}\,\bigl(F(u)\subseteq F(\Succ(u))\bigr)}{\rA}
    \proofstepwidestar[2]{k\in\mathbb{N}}{\rA}
    \proofstepwidestar[3]{n\in\mathbb{N}}{\rA}
    \proofstepwidestar[4]{k<n}{\rA}
    \proofstepwidestar[2,3,4]{k\in n}{%
      \FormulaRefAuto{m\in\mathbb{N}\dsep n\in\mathbb{N}\vdash m < n \leftrightarrow m\in n}{2,3,4}}
    \proofstepwidestar[1,3]{%
      \forall k\in\mathbb{N}\,\bigl(k\in n\rightarrow F(\Succ(k))\subseteq F(n)\bigr)
    }{%
      \FormulaRefAuto{PeanoInductionRule}{IA,IS,3}}
    \proofstepwidestar[1,2,3]{k\in n\rightarrow F(\Succ(k))\subseteq F(n)}{%
      \rRE{2,\rUE{6}}}
    \proofstepwidestar[1,2,3,4]{F(\Succ(k))\subseteq F(n)}{\rRE{5,7}}
  \closeproofpartwide
\end{tabproofsplitwide}

\FormulaThmDelta[Negation der natürlichen Ordnung]{%
  m\in\mathbb{N}\dsep n\in\mathbb{N}\dsep \neg(m\leq n)\vdash n < m%
}{%
  \DeltaRow{Mengen}{m\dsep n}%
  \DeltaRow{Relationen}{<_{\mathbb{N}}\dsep \leq_{\mathbb{N}}}%
}
\begin{tabproof}
  \proofstep{1}{m\in\mathbb{N}}{\rA}
  \proofstep{2}{n\in\mathbb{N}}{\rA}
  \proofstep{3}{\neg(m\leq n)}{\rA}
  \proofstep{4}{m\subseteq n}{\rA}
  \proofstep{1,2,4}{m\leq n}{%
    \FormulaRefAuto{m\in\mathbb{N}\dsep n\in\mathbb{N}\vdash m\leq n \leftrightarrow m\subseteq n}{1,2,4}}
  \proofstep{1,2,3,4}{\bot}{\rBI{5,3}}
  \proofstep{1,2,3}{\neg(m\subseteq n)}{\rCI{4,6}}
  \proofstep{1,2,3}{n\in m}{%
    \FormulaRefAuto{m\in\mathbb{N}\dsep n\in\mathbb{N}\dsep \neg(m\subseteq n)\vdash n\in m}{1,2,7}}
  \proofstep{1,2,3}{n < m}{%
    \FormulaRefAuto{m\in\mathbb{N}\dsep n\in\mathbb{N}\vdash m < n \leftrightarrow m\in n}{2,1,8}}
\end{tabproof}

\FormulaThmDelta[Strikte natürliche Ordnung impliziert natürliche Ordnung]{%
  m\in\mathbb{N}\dsep n\in\mathbb{N}\dsep m < n\vdash m\leq n%
}{%
  \DeltaRow{Mengen}{m\dsep n}%
  \DeltaRow{Relationen}{<_{\mathbb{N}}\dsep \leq_{\mathbb{N}}}%
}
\begin{tabproof}
  \proofstep{1}{m\in\mathbb{N}}{\rA}
  \proofstep{2}{n\in\mathbb{N}}{\rA}
  \proofstep{3}{m < n}{\rA}
  \proofstep{1,2,3}{m\in n}{%
    \FormulaRefAuto{m\in\mathbb{N}\dsep n\in\mathbb{N}\vdash m < n \leftrightarrow m\in n}{1,2,3}}
  \proofstep{1,2,3}{m\subseteq n}{%
    \FormulaRefAuto{m\in\mathbb{N}\dsep n\in\mathbb{N}\dsep m\in n \vdash m\subseteq n}{1,2,4}}
  \proofstep{1,2,3}{m\leq n}{%
    \FormulaRefAuto{m\in\mathbb{N}\dsep n\in\mathbb{N}\vdash m\leq n \leftrightarrow m\subseteq n}{1,2,5}}
\end{tabproof}

\FormulaThmDelta[Natürliche Ordnung ist strikte Ordnung oder Gleichheit]{%
  m\in\mathbb{N}\dsep n\in\mathbb{N}\dsep m\leq n\vdash m < n\lor m=n%
}{%
  \DeltaRow{Mengen}{m\dsep n}%
  \DeltaRow{Relationen}{<_{\mathbb{N}}\dsep \leq_{\mathbb{N}}}%
}
\begin{tabproof}
  \proofstep{1}{m\in\mathbb{N}}{\rA}
  \proofstep{2}{n\in\mathbb{N}}{\rA}
  \proofstep{3}{m\leq n}{\rA}
  \proofstep{1,2,3}{m\subseteq n}{%
    \FormulaRefAuto{m\in\mathbb{N}\dsep n\in\mathbb{N}\vdash m\leq n \leftrightarrow m\subseteq n}{1,2,3}}
  \proofstep{1,2,3}{m\in n\lor m=n}{%
    \FormulaRefAuto{m\in\mathbb{N}\dsep n\in\mathbb{N}\dsep m\subseteq n \vdash m\in n\lor m=n}{1,2,4}}
  \proofstep{6}{m\in n}{\rA}
  \proofstep{1,2,6}{m < n}{%
    \FormulaRefAuto{m\in\mathbb{N}\dsep n\in\mathbb{N}\vdash m < n \leftrightarrow m\in n}{1,2,6}}
  \proofstep{1,2,6}{m < n\lor m=n}{\rOIa{7}}
  \proofstep{9}{m=n}{\rA}
  \proofstep{9}{m < n\lor m=n}{\rOIb{9}}
  \proofstep{1,2,3}{m < n\lor m=n}{\rOE{5,6,8,9,10}}
\end{tabproof}

\FormulaThmDelta[Ungleiche natürliche Ordnung ist strikt]{%
  m\in\mathbb{N}\dsep n\in\mathbb{N}\dsep m\leq n\dsep m\neq n\vdash m < n%
}{%
  \DeltaRow{Mengen}{m\dsep n}%
  \DeltaRow{Relationen}{<_{\mathbb{N}}\dsep \leq_{\mathbb{N}}}%
}
\begin{tabproof}
  \proofstep{1}{m\in\mathbb{N}}{\rA}
  \proofstep{2}{n\in\mathbb{N}}{\rA}
  \proofstep{3}{m\leq n}{\rA}
  \proofstep{4}{m\neq n}{\rA}
  \proofstep{1,2,3}{m < n\lor m=n}{%
    \FormulaRefAuto{m\in\mathbb{N}\dsep n\in\mathbb{N}\dsep m\leq n\vdash m < n\lor m=n}{1,2,3}}
  \proofstep{1,2,3,4}{m < n}{%
    \FormulaRefAuto{P\lor Q,\,\neg Q \vdash P}{5,4}}
\end{tabproof}

\FormulaThmDelta[Strikte natürliche Ordnung impliziert Ungleichheit]{%
  m\in\mathbb{N}\dsep n\in\mathbb{N}\dsep m < n\vdash m\neq n%
}{%
  \DeltaRow{Mengen}{m\dsep n}%
  \DeltaRow{Relationen}{<_{\mathbb{N}}}%
}
\begin{tabproof}
  \proofstep{1}{m\in\mathbb{N}}{\rA}
  \proofstep{2}{n\in\mathbb{N}}{\rA}
  \proofstep{3}{m < n}{\rA}
  \proofstep{1,2,3}{m\in n}{%
    \FormulaRefAuto{m\in\mathbb{N}\dsep n\in\mathbb{N}\vdash m < n \leftrightarrow m\in n}{1,2,3}}
  \proofstep{1,2,3}{m\neq n}{%
    \FormulaRefAuto{m\in\mathbb{N}\dsep n\in\mathbb{N}\dsep m\in n \vdash m\neq n}{1,2,4}}
\end{tabproof}

\FormulaThmDelta[Vergleichbarkeit der natürlichen Ordnung]{%
  m\in\mathbb{N}\dsep n\in\mathbb{N}\vdash m\leq n\lor n\leq m%
}{%
  \DeltaRow{Mengen}{m\dsep n}%
  \DeltaRow{Relationen}{<_{\mathbb{N}}\dsep \leq_{\mathbb{N}}}%
}
\begin{tabproof}
  \proofstep{1}{m\in\mathbb{N}}{\rA}
  \proofstep{2}{n\in\mathbb{N}}{\rA}
  \proofstep{1,2}{m < n\lor n < m\lor m=n}{%
    \FormulaRefAuto{m\in\mathbb{N}\dsep n\in\mathbb{N}\vdash m < n\lor n < m\lor m=n}{1,2}}
  \proofstep{4}{m < n}{\rA}
  \proofstep{1,2,4}{m\leq n}{%
    \FormulaRefAuto{m\in\mathbb{N}\dsep n\in\mathbb{N}\dsep m < n\vdash m\leq n}{1,2,4}}
  \proofstep{1,2,4}{m\leq n\lor n\leq m}{\rOIa{5}}
  \proofstep{7}{n < m}{\rA}
  \proofstep{1,2,7}{n\leq m}{%
    \FormulaRefAuto{m\in\mathbb{N}\dsep n\in\mathbb{N}\dsep m < n\vdash m\leq n}{2,1,7}}
  \proofstep{1,2,7}{m\leq n\lor n\leq m}{\rOIb{8}}
  \proofstep{10}{m=n}{\rA}
  \proofstep{10}{m\subseteq n}{\FormulaRefAuto{A=B\vdash A\subseteq B}{10}}
  \proofstep{1,2,10}{m\leq n}{%
    \FormulaRefAuto{m\in\mathbb{N}\dsep n\in\mathbb{N}\vdash m\leq n \leftrightarrow m\subseteq n}{1,2,11}}
  \proofstep{1,2,10}{m\leq n\lor n\leq m}{\rOIa{12}}
  \proofstep{1,2}{m\leq n\lor n\leq m}{%
    \FormulaRefAuto{P_1\lor P_2\lor P_3\dsep [P_1]\vdots R\dsep [P_2]\vdots R\dsep [P_3]\vdots R\vdash R}{3,4,6,7,9,10,13}}
\end{tabproof}

\FormulaThmDelta[Transitivität der natürlichen Ordnung]{%
  m\in\mathbb{N}\dsep n\in\mathbb{N}\dsep k\in\mathbb{N}\dsep
  m\leq n\dsep n\leq k\vdash m\leq k%
}{%
  \DeltaRow{Mengen}{k\dsep m\dsep n}%
  \DeltaRow{Relationen}{\leq_{\mathbb{N}}}%
}
\begin{tabproof}
  \proofstep{1}{m\in\mathbb{N}}{\rA}
  \proofstep{2}{n\in\mathbb{N}}{\rA}
  \proofstep{3}{k\in\mathbb{N}}{\rA}
  \proofstep{4}{m\leq n}{\rA}
  \proofstep{5}{n\leq k}{\rA}
  \proofstep{1,2,4}{m\subseteq n}{%
    \FormulaRefAuto{m\in\mathbb{N}\dsep n\in\mathbb{N}\vdash
    m\leq n \leftrightarrow m\subseteq n}{1,2,4}}
  \proofstep{2,3,5}{n\subseteq k}{%
    \FormulaRefAuto{m\in\mathbb{N}\dsep n\in\mathbb{N}\vdash
    m\leq n \leftrightarrow m\subseteq n}{2,3,5}}
  \proofstep{1,2,3,4,5}{m\subseteq k}{%
    \FormulaRefAuto{A \subseteq B, B \subseteq C \vdash A \subseteq C}{6,7}}
  \proofstep{1,2,3,4,5}{m\leq k}{%
    \FormulaRefAuto{m\in\mathbb{N}\dsep n\in\mathbb{N}\vdash
    m\leq n \leftrightarrow m\subseteq n}{1,3,8}}
\end{tabproof}

\FormulaThmDelta[Totale Ordnung der natürlichen Zahlen]{%
  \TotOrd{\leq_{\mathbb{N}},\mathbb{N}}%
}{%
  \DeltaRow{Mengen}{k\dsep m\dsep n}
  \DeltaRow{Relationen}{\leq_{\mathbb{N}}}
  \DeltaRow{\textbf{Begründung}}{}
  \DeltaRow{Teilmenge als Ordnung}{m\leq n \leftrightarrow m\subseteq n}
    [\FormulaRefAuto{m\in\mathbb{N}\dsep n\in\mathbb{N}\vdash m\leq n \leftrightarrow m\subseteq n}]
  \DeltaRow{Vergleichbarkeit}{m\in\mathbb{N}\dsep n\in\mathbb{N}\vdash m\leq n\lor n\leq m}
    [\FormulaRefAuto{m\in\mathbb{N}\dsep n\in\mathbb{N}\vdash m\leq n\lor n\leq m}]
  \DeltaRow{Antisymmetrie}{m\leq n\dsep n\leq m\vdash m=n}
    [\FormulaRefAuto{A \subseteq B, B \subseteq A \vdash A = B}]
  \DeltaRow{Transitivität}{%
    m\in\mathbb{N}\dsep n\in\mathbb{N}\dsep k\in\mathbb{N}\dsep
    m\leq n\dsep n\leq k\vdash m\leq k}
    [\FormulaRefAuto{%
    m\in\mathbb{N}\dsep n\in\mathbb{N}\dsep k\in\mathbb{N}\dsep
    m\leq n\dsep n\leq k\vdash m\leq k}]
}{}


\section{Wohlordnungsprinzip}

\FormulaThmDelta[Wohlordnungsprinzip der natürlichen Zahlen]{%
  B\subseteq\mathbb{N}\dsep B\neq\emptyset
  \vdash \exists a\in B\,\forall b\in B\,a\subseteq b
}{%
  \DeltaDecl{Mengen}{B}%
}
\begin{tabproof}
  \proofstep{1}{B\subseteq\mathbb{N}}{\rA}
  \proofstep{2}{B\neq0}{\rA}

  \proofstep{2}{\exists x\in B\,\bigl(x\cap B=0\bigr)}{%
    \FormulaRefAuto{A \neq \varnothing \vdash \exists x \in A \,(x \cap A = \varnothing )}{2}}

  \proofstep{4}{a\in B\land a\cap B=0}{\rA}
  \proofstep{4}{a\in B}{\rAEa{4}}
  \proofstep{4}{a\cap B=0}{\rAEb{4}}

  \proofstep{7}{b\in B}{\rA}
  \proofstep{1,4}{a\in\mathbb{N}}{\FormulaRefAuto{A\subseteq B,\, x\in A \vdash x\in B}{1,5}}
  \proofstep{1,7}{b\in\mathbb{N}}{\FormulaRefAuto{A\subseteq B,\, x\in A \vdash x\in B}{1,7}}

  \proofstep{1,4,7}{a\in b\lor b\in a\lor a=b}{%
    \FormulaRefAuto{m\in\mathbb{N}\dsep n\in\mathbb{N}\vdash m\in n\lor n\in m\lor m=n}{8,9}}

  \proofstep{11}{a\in b}{\rA}
  \proofstep{1,4,7,11}{a\subseteq b}{%
    \FormulaRefAuto{m\in\mathbb{N}\dsep n\in\mathbb{N}\dsep m\in n\vdash m\subseteq n}{8,9,11}}

  \proofstep{13}{b\in a}{\rA}
  \proofstep{7,13}{b\in a\cap B}{\FormulaRefAuto{x \in A, x\in B \vdash x \in A\cap B}{13,7}}
  \proofstep{4,7,13}{b\in0}{\FormulaRefAuto{A = B,\, x \in A \vdash x \in B}{6,14}}
  \proofstep{}{b\not\in0}{\FormulaRefAuto{x \not\in \varnothing}}
  \proofstep{4,7,13}{b\in0 \land b\not\in0}{\rAI{15,16}}
  \proofstep{4,7,13}{a\subseteq b}{\FormulaRefAuto{P \land \neg P\vdash Q}{17}}

  \proofstep{19}{a=b}{\rA}
  \proofstep{19}{a\subseteq b}{\FormulaRefAuto{A=B \vdash A\subseteq B}{19}}

  \proofstep{1,4,7}{a\subseteq b}{%
    \FormulaRefAuto{P_1\lor P_2\lor P_3\dsep [P_1]\vdots R\dsep [P_2]\vdots R\dsep [P_3]\vdots R\vdash R}{10,11,12,13,18,19,20}}

  \proofstep{1,4}{\forall b\in B\,a\subseteq b}{\rUI{\rRI{7,21}}}
  \proofstep{1,4}{\exists a\in B\,\forall b\in B\,a\subseteq b}{\rEI{\rAI{5,22}}}

  \proofstep{1,2}{\exists a\in B\,\forall b\in B\,a\subseteq b}{\rEE{3,4,23}}
\end{tabproof}

\FormulaThmDelta[Wohlordnungsprinzip der natürlichen Zahlen mit \(\leq\)]{%
  B\subseteq\mathbb{N}\dsep B\neq0
  \vdash \exists a\in B\,\forall b\in B\,a\leq b
}{%
  \DeltaDecl{Mengen}{B\dsep a\dsep b}%
  \DeltaRow{Relationen}{\leq_{\mathbb{N}}}%
}
\begin{tabproof}
  \proofstep{1}{B\subseteq\mathbb{N}}{\rA}
  \proofstep{2}{B\neq0}{\rA}
  \proofstep{1,2}{\exists a\in B\,\forall b\in B\,a\subseteq b}{%
    \FormulaRefAuto{B\subseteq\mathbb{N}\dsep B\neq\emptyset \vdash \exists a\in B\,\forall b\in B\,a\subseteq b}{1,2}}
  \proofstep{4}{a\in B\land \forall b\in B\,a\subseteq b}{\rA}
  \proofstep{4}{a\in B}{\rAEa{4}}
  \proofstep{4}{\forall b\in B\,a\subseteq b}{\rAEb{4}}
  \proofstep{7}{b\in B}{\rA}
  \proofstep{1,4}{a\in\mathbb{N}}{%
    \FormulaRefAuto{A\subseteq B,\, x\in A \vdash x\in B}{1,5}}
  \proofstep{1,7}{b\in\mathbb{N}}{%
    \FormulaRefAuto{A\subseteq B,\, x\in A \vdash x\in B}{1,7}}
  \proofstep{4,7}{a\subseteq b}{\rRE{7,\rUE{6}}}
  \proofstep{1,4,7}{a\leq b}{%
    \FormulaRefAuto{m\in\mathbb{N}\dsep n\in\mathbb{N}\vdash m\leq n \leftrightarrow m\subseteq n}{8,9,10}}
  \proofstep{1,4}{\forall b\in B\,a\leq b}{\rUI{\rRI{7,11}}}
  \proofstep{1,4}{\exists a\in B\,\forall b\in B\,a\leq b}{\rEI{\rAI{5,12}}}
  \proofstep{1,2}{\exists a\in B\,\forall b\in B\,a\leq b}{\rEE{3,4,13}}
\end{tabproof}

\FormulaThmDelta[Wohlordnungsprinzip der natürlichen Zahlen in Relationsform]{%
  B\subseteq\mathbb{N}\dsep B\neq0
  \vdash \exists a\in B\,\forall b\in B\,\bigl((a,b)\in \leq_{\mathbb{N}}\bigr)
}{%
  \DeltaDecl{Mengen}{B\dsep a\dsep b}%
  \DeltaRow{Relationen}{\leq_{\mathbb{N}}}%
}
\begin{tabproof}
  \proofstep{1}{B\subseteq\mathbb{N}}{\rA}
  \proofstep{2}{B\neq0}{\rA}
  \proofstep{1,2}{\exists a\in B\,\forall b\in B\,a\leq b}{%
    \FormulaRefAuto{B\subseteq\mathbb{N}\dsep B\neq\varnothing \vdash \exists a\in B\,\forall b\in B\,a\leq b}{1,2}}
  \proofstep{4}{a\in B\land \forall b\in B\,a\leq b}{\rA}
  \proofstep{4}{a\in B}{\rAEa{4}}
  \proofstep{4}{\forall b\in B\,a\leq b}{\rAEb{4}}
  \proofstep{7}{b\in B}{\rA}
  \proofstep{1,4}{a\in\mathbb{N}}{%
    \FormulaRefAuto{A\subseteq B,\, x\in A \vdash x\in B}{1,5}}
  \proofstep{1,7}{b\in\mathbb{N}}{%
    \FormulaRefAuto{A\subseteq B,\, x\in A \vdash x\in B}{1,7}}
  \proofstep{4,7}{a\leq b}{\rRE{7,\rUE{6}}}
  \proofstep{1,4,7}{(a,b)\in \leq_{\mathbb{N}}}{%
    \FormulaRefAuto{m\in\mathbb{N}\dsep n\in\mathbb{N}\dsep m\leq n \vdash (m,n)\in \leq_{\mathbb{N}}}{8,9,10}}
  \proofstep{1,4}{\forall b\in B\,\bigl((a,b)\in \leq_{\mathbb{N}}\bigr)}{\rUI{\rRI{7,11}}}
  \proofstep{1,4}{\exists a\in B\,\forall b\in B\,\bigl((a,b)\in \leq_{\mathbb{N}}\bigr)}{\rEI{\rAI{5,12}}}
  \proofstep{1,2}{\exists a\in B\,\forall b\in B\,\bigl((a,b)\in \leq_{\mathbb{N}}\bigr)}{\rEE{3,4,13}}
\end{tabproof}

\FormulaThmDelta[Wohlordnung der natürlichen Zahlen]{%
  \WellOrd{\leq_{\mathbb{N}},\mathbb{N}}%
}{%
  \DeltaRow{Relationen}{\leq_{\mathbb{N}}}
  \DeltaRow{\textbf{Begründung}}{}
  \DeltaRow{Totale Ordnung}{\TotOrd{\leq_{\mathbb{N}},\mathbb{N}}}
    [\FormulaRefAuto{\TotOrd{\leq_{\mathbb{N}},\mathbb{N}}}]
  \DeltaRow{Wohlordnungsprinzip}{%
    B\subseteq\mathbb{N}\dsep B\neq0
    \vdash \exists a\in B\,\forall b\in B\,\bigl((a,b)\in \leq_{\mathbb{N}}\bigr)}
    [\FormulaRefAuto{B\subseteq\mathbb{N}\dsep B\neq\varnothing \vdash \exists a\in B\,\forall b\in B\,\bigl((a,b)\in \leq_{\mathbb{N}}\bigr)}]
}{}
\begin{tabproof}
  \proofstep{}{\TotOrd{\leq_{\mathbb{N}},\mathbb{N}}}{%
    \FormulaRefAuto{\TotOrd{\leq_{\mathbb{N}},\mathbb{N}}}}
  \proofstep{}{%
    B\subseteq\mathbb{N}\dsep B\neq0
    \vdash \exists a\in B\,\forall b\in B\,\bigl((a,b)\in \leq_{\mathbb{N}}\bigr)}{%
    \FormulaRefAuto{B\subseteq\mathbb{N}\dsep B\neq\varnothing \vdash \exists a\in B\,\forall b\in B\,\bigl((a,b)\in \leq_{\mathbb{N}}\bigr)}}
  \proofstep{}{\WellOrd{\leq_{\mathbb{N}},\mathbb{N}}}{%
    \FormulaRefAuto{Wohlordnung}{1,2}}
\end{tabproof}


% =========================================================
\section{Beispiele von Funktionen}



\subsection{Entfernen eines Elements aus den natürlichen Zahlen}

Im folgenden Konstruktionsschritt bezeichnet \(\tJomit{x_0}\) nur die
punktweise Funktionsvorschrift; die eigentliche Auslassungsabbildung ist
\(\Jomit{x_0}\).

\FormulaDefDelta[Ausgelassene Menge zu \(x_0\)]{%
  x_0\in\mathbb{N}\vdash
  \SOmit{x_0}\coloneqq \{x\in\mathbb{N}\mid x\neq x_0\}%
}{%
  \DeltaRow{Mengen}{x_0\dsep x}%
  \DeltaRow{Neue Symbole}{\SOmit{x_0}}%
}

\FormulaThmDelta{%
  x_0\in\mathbb{N}\vdash \{x\in\mathbb{N}\mid x\neq x_0\}=\SOmit{x_0}%
}{%
  \DeltaDecl{Mengen}{x_0}%
}
\begin{tabproof}
  \proofstep{1}{x_0\in\mathbb{N}}{\rA}
  \proofstep{1}{\SOmit{x_0}=\{x\in\mathbb{N}\mid x\neq x_0\}}{%
    \FormulaRefAuto{x_0\in\mathbb{N}\vdash
      \SOmit{x_0}\coloneqq \{x\in\mathbb{N}\mid x\neq x_0\}}{1}}
  \proofstep{2}{\{x\in\mathbb{N}\mid x\neq x_0\}=\SOmit{x_0}}{%
    \FormulaRefAuto{a=b\vdash b=a}{2}}
\end{tabproof}

\FormulaThmDelta{%
  x_0\in\mathbb{N}\vdash \SOmit{x_0}\subseteq\mathbb{N}%
}{%
  \DeltaDecl{Mengen}{x_0}%
}
\begin{tabproof}
  \proofstep{1}{x_0\in\mathbb{N}}{\rA}

  \proofstep{}{\{x\in\mathbb{N}\mid x\neq x_0\}\subseteq\mathbb{N}}{%
    \FormulaRefAuto{\{ x \in A \mid P(x) \} \subseteq A}}

  \proofstep{1}{\{x\in\mathbb{N}\mid x\neq x_0\}=\SOmit{x_0}}{%
    \FormulaRefAuto{x_0\in\mathbb{N}\vdash \{x\in\mathbb{N}\mid x\neq x_0\}=\SOmit{x_0}}{1}}

  \proofstep{2,3}{\SOmit{x_0}\subseteq\mathbb{N}}{\rIE{3,2}}
\end{tabproof}

\FormulaDefDelta[Funktionsvorschrift der Auslassungsabbildung]{%
  n\in\mathbb{N}\vdash
  \tJomit{x_0}(n)\coloneqq
  \IfThenElse{n\in x_0}{n}{\Succ(n)}%
}{%
  \DeltaRow{Mengen}{x_0\dsep n}%
  \DeltaRow{Funktionensymbol}{\tJomit{x_0}}%
}



\FormulaThmDelta[Funktionsvorschrift der Auslassungsabbildung im Elementfall]{%
  n\in\mathbb{N}\dsep n\in x_0 \vdash \tJomit{x_0}(n)=n%
}{%
  \DeltaDecl{Mengen}{x_0\dsep n}%
  \DeltaDecl{Funktionensymbol}{\tJomit{x_0}}%
}
\begin{tabproof}
  \proofstep{1}{n\in\mathbb{N}}{\rA}
  \proofstep{2}{n\in x_0}{\rA}
  \proofstep{1}{\tJomit{x_0}(n)=\IfThenElse{n\in x_0}{n}{\Succ(n)}}{%
    \FormulaRefAuto{n\in\mathbb{N}\vdash \tJomit{x_0}(n)\coloneqq \IfThenElse{n\in x_0}{n}{\Succ(n)}}{1}}
  \proofstep{2}{\IfThenElse{n\in x_0}{n}{\Succ(n)}=n}{%
    \FormulaRefAuto{P \vdash \IfThenElse{P}{a}{b}=a}{2}}
  \proofstep{1,2}{\tJomit{x_0}(n)=n}{%
    \FormulaRefAuto{a = b,\, b = c \vdash a = c}{3,4}}
\end{tabproof}

\FormulaThmDelta[Funktionsvorschrift der Auslassungsabbildung im Nicht-Elementfall]{%
  n\in\mathbb{N}\dsep n\notin x_0 \vdash \tJomit{x_0}(n)=\Succ(n)%
}{%
  \DeltaDecl{Mengen}{x_0\dsep n}%
  \DeltaDecl{Funktionensymbol}{\tJomit{x_0}}%
}
\begin{tabproof}
  \proofstep{1}{n\in\mathbb{N}}{\rA}
  \proofstep{2}{n\notin x_0}{\rA}
  \proofstep{1}{\tJomit{x_0}(n)=\IfThenElse{n\in x_0}{n}{\Succ(n)}}{%
    \FormulaRefAuto{n\in\mathbb{N}\vdash \tJomit{x_0}(n)\coloneqq \IfThenElse{n\in x_0}{n}{\Succ(n)}}{1}}
  \proofstep{2}{\IfThenElse{n\in x_0}{n}{\Succ(n)}=\Succ(n)}{%
    \FormulaRefAuto{\neg P \vdash \IfThenElse{P}{a}{b}=b}{2}}
  \proofstep{1,2}{\tJomit{x_0}(n)=\Succ(n)}{%
    \FormulaRefAuto{a = b,\, b = c \vdash a = c}{3,4}}
\end{tabproof}

\FormulaThmDelta[Typisierung der Funktionsvorschrift]{%
  x_0\in\mathbb{N}\dsep n\in\mathbb{N}\vdash
  \tJomit{x_0}(n)\in \SOmit{x_0}%
}{%
  \DeltaDecl{Mengen}{x_0\dsep n}%
  \DeltaDecl{Funktionensymbol}{\tJomit{x_0}}%
}
\begin{tabproofsplitwide}

  \proofpartwideind[Fall $n\in x_0$]{%
    x_0\in\mathbb{N}\dsep n\in\mathbb{N}\dsep n\in x_0 \vdash
    \tJomit{x_0}(n)\in \SOmit{x_0}%
  }
    \proofstepwidestar[1]{x_0\in\mathbb{N}}{\rA}
    \proofstepwidestar[2]{n\in\mathbb{N}}{\rA}
    \proofstepwidestar[3]{n\in x_0}{\rA}

    \proofstepwidestar[2,3]{\tJomit{x_0}(n)=n}{%
      \FormulaRefAuto{n\in\mathbb{N}\dsep n\in x_0 \vdash \tJomit{x_0}(n)=n}{2,3}}

    \proofstepwidestar[1,2]{n\in x_0 \leftrightarrow n\subset x_0}{%
      \FormulaRefAuto{m\in\mathbb{N}\dsep n\in\mathbb{N}\vdash m\in n \leftrightarrow m\subset n}{2,1}}
    \proofstepwidestar[1,2,3]{n\subset x_0}{%
      \FormulaRefAuto{P \leftrightarrow Q, P \vdash Q}{5,3}}
    \proofstepwidestar[1,2,3]{n\neq x_0}{%
      \FormulaRefAuto{A\subset B \vdash A\neq B}{6}}
    \proofstepwidestar[1,2,3]{n\in\mathbb{N}\land n\neq x_0}{\rAI{2,7}}
    \proofstepwidestar[1,2,3]{n\in \SOmit{x_0}}{%
      \FormulaRefAuto{x\in\{u\in A\mid P(u)\}\eqvdash x\in A\land P(x)}{8}}

    \proofstepwidestar[1,2,3]{\tJomit{x_0}(n)\in \SOmit{x_0}}{\rIE{4,9}}
  \closeproofpartwideind


  \proofpartwideind[Fall $n\notin x_0$]{%
    x_0\in\mathbb{N}\dsep n\in\mathbb{N}\dsep n\notin x_0 \vdash
    \tJomit{x_0}(n)\in \SOmit{x_0}%
  }
    \proofstepwidestar[1]{x_0\in\mathbb{N}}{\rA}
    \proofstepwidestar[2]{n\in\mathbb{N}}{\rA}
    \proofstepwidestar[3]{n\notin x_0}{\rA}

    \proofstepwidestar[2,3]{\tJomit{x_0}(n)=\Succ(n)}{%
      \FormulaRefAuto{n\in\mathbb{N}\dsep n\notin x_0 \vdash \tJomit{x_0}(n)=\Succ(n)}{2,3}}
    \proofstepwidestar[2]{\Succ(n)\in\mathbb{N}}{%
      \FormulaRefAuto{PeanoSuccClosure}{2}}

    \proofstepwidestar[6]{\Succ(n)=x_0}{\rA}
    \proofstepwidestar[2]{n\in\Succ(n)}{%
      \FormulaRefAuto{n\in\mathbb{N}\vdash n\in\Succ(n)}{2}}
    \proofstepwidestar[1,2,3,6]{n\in x_0}{\rIE{6,7}}
    \proofstepwidestar[1,2,3,6]{\bot}{\rBI{3,8}}
    \proofstepwidestar[1,2,3]{\Succ(n)\neq x_0}{\rCI{6,9}}

    \proofstepwidestar[1,2,3]{\Succ(n)\in\mathbb{N}\land \Succ(n)\neq x_0}{\rAI{5,10}}
    \proofstepwidestar[1,2,3]{\Succ(n)\in \SOmit{x_0}}{%
      \FormulaRefAuto{x\in\{u\in A\mid P(u)\}\eqvdash x\in A\land P(x)}{11}}

    \proofstepwidestar[1,2,3]{\tJomit{x_0}(n)\in \SOmit{x_0}}{\rIE{4,12}}
  \closeproofpartwideind


  \proofpartwide{Zusammenfassung}
    \proofstepwidestar{x_0\in\mathbb{N}}{\rA}
    \proofstepwidestar{n\in\mathbb{N}}{\rA}

    \proofstepwidestar[2]{n\in x_0\lor n\notin x_0}{\FormulaRefAuto{P\lor \neg P}}

    \proofstepwidestar[4]{n\in x_0}{\rA}
    \proofstepwidestar[1,2,4]{\tJomit{x_0}(n)\in \SOmit{x_0}}{%
      \FormulaRefAuto{x_0\in\mathbb{N}\dsep n\in\mathbb{N}\dsep n\in x_0 \vdash
        \tJomit{x_0}(n)\in \SOmit{x_0}}{1,2,4}}

    \proofstepwidestar[6]{n\notin x_0}{\rA}
    \proofstepwidestar[1,2,6]{\tJomit{x_0}(n)\in \SOmit{x_0}}{%
      \FormulaRefAuto{x_0\in\mathbb{N}\dsep n\in\mathbb{N}\dsep n\notin x_0 \vdash
        \tJomit{x_0}(n)\in \SOmit{x_0}}{1,2,6}}

    \proofstepwidestar[1,2]{\tJomit{x_0}(n)\in \SOmit{x_0}}{\rOE{3,4,5,6,7}}
  \closeproofpartwide

\end{tabproofsplitwide}

\FormulaThmDelta[Universelle Typisierung der Funktionsvorschrift]{%
  x_0\in\mathbb{N}\vdash
  \forall n\in\mathbb{N}\,\tJomit{x_0}(n)\in \SOmit{x_0}%
}{%
  \DeltaDecl{Mengen}{x_0\dsep n}%
  \DeltaDecl{Funktionensymbol}{\tJomit{x_0}}%
}
\begin{tabproof}
  \proofstep{1}{x_0\in\mathbb{N}}{\rA}

  \proofstep{2}{n\in\mathbb{N}}{\rA}
  \proofstep{1,2}{\tJomit{x_0}(n)\in \SOmit{x_0}}{%
    \FormulaRefAuto{x_0\in\mathbb{N}\dsep n\in\mathbb{N}\vdash
      \tJomit{x_0}(n)\in \SOmit{x_0}}{1,2}}

  \proofstep{1}{\forall n\in\mathbb{N}\,\tJomit{x_0}(n)\in \SOmit{x_0}}{%
    \rUI{\rRI{2,3}}}
\end{tabproof}

\FormulaThmDelta[Auslassungsabbildung als Funktion]{%
  x_0\in\mathbb{N}\vdash
  \exists! J\colon \mathbb{N}\to \SOmit{x_0}\,
  \forall n\in\mathbb{N}\,J(n)=\tJomit{x_0}(n)%
}{%
  \DeltaDecl{Mengen}{x_0}%
  \DeltaDecl{Funktionen}{J}%
  \DeltaDecl{Funktionensymbol}{\tJomit{x_0}}%
}
\begin{tabproof}
  \proofstep{1}{x_0\in\mathbb{N}}{\rA}

  \proofstep{1}{\forall n\in\mathbb{N}\,\tJomit{x_0}(n)\in \SOmit{x_0}}{%
    \FormulaRefAuto{%
      x_0\in\mathbb{N}\vdash
      \forall n\in\mathbb{N}\,\tJomit{x_0}(n)\in \SOmit{x_0}}{1}}

  \proofstep{1}{%
    \exists! J\colon \mathbb{N}\to \SOmit{x_0}\,
    \forall n\in\mathbb{N}\,J(n)=\tJomit{x_0}(n)}{%
    \FormulaRefAuto{%
      x\in A\vdash t(x)\coloneq y\dsep x\in A\vdash t(x)\in B
      \exists! F\colon A\to B\,\forall x\in A\,F(x) = t(x)}{2}}
\end{tabproof}

\FormulaDefDelta[Auslassungsabbildung]{%
  x_0\in\mathbb{N}\vdash
  \Jomit{x_0}\coloneqq \iota F\Bigl(
    F\colon \mathbb{N}\to \SOmit{x_0}\land
    \forall n\in\mathbb{N}\,F(n)=\tJomit{x_0}(n)
  \Bigr)%
}{%
  \DeltaRow{Mengen}{x_0\dsep n}%
  \DeltaRow{Funktionensymbole}{\Jomit{x_0}\dsep \tJomit{x_0}}%
}


\FormulaThmDelta[Auslassungsabbildung]{%
  x_0\in\mathbb{N}\vdash \Jomit{x_0}\colon \mathbb{N}\to \SOmit{x_0}%
}{%
  \DeltaRow{Mengen}{x_0}%
}
\begin{tabproof}
  \proofstep{1}{x_0\in\mathbb{N}}{\rA}
  \proofstep{1}{\Jomit{x_0}\colon \mathbb{N}\to \SOmit{x_0}}{%
    \rAEa{\FormulaRefAuto{%
      x_0\in\mathbb{N}\vdash
      \Jomit{x_0}\coloneqq \iota F\Bigl(
        F\colon \mathbb{N}\to \SOmit{x_0}\land
        \forall n\in\mathbb{N}\,F(n)=\tJomit{x_0}(n)
      \Bigr)}{1}}}
\end{tabproof}


\FormulaThmDelta[Grundgleichung der Auslassungsabbildung]{%
  x_0\in\mathbb{N}\dsep n\in\mathbb{N}\vdash \Jomit{x_0}(n)=\tJomit{x_0}(n)%
}{%
  \DeltaRow{Mengen}{x_0\dsep n}%
  \DeltaRow{Funktionensymbole}{\Jomit{x_0}\dsep \tJomit{x_0}}%
}
\begin{tabproof}
  \proofstep{1}{x_0\in\mathbb{N}}{\rA}
  \proofstep{2}{n\in\mathbb{N}}{\rA}

  \proofstep{1}{\forall n\in\mathbb{N}\,\Jomit{x_0}(n)=\tJomit{x_0}(n)}{%
    \rAEb{\FormulaRefAuto{%
      x_0\in\mathbb{N}\vdash
      \Jomit{x_0}\coloneqq \iota F\Bigl(
        F\colon \mathbb{N}\to \SOmit{x_0}\land
        \forall n\in\mathbb{N}\,F(n)=\tJomit{x_0}(n)
      \Bigr)}{1}}}

  \proofstep{1,2}{\Jomit{x_0}(n)=\tJomit{x_0}(n)}{\rRE{2,\rUE{3}}}
\end{tabproof}

\FormulaThmDelta[$\Jomit{x_0}$ im Elementfall]{%
  x_0\in\mathbb{N}\dsep n\in\mathbb{N}\dsep n\in x_0 \vdash \Jomit{x_0}(n)=n%
}{%
  \DeltaDecl{Mengen}{x_0\dsep n}%
  \DeltaDecl{Funktionensymbol}{\Jomit{x_0}}%
}
\begin{tabproof}
  \proofstep{1}{x_0\in\mathbb{N}}{\rA}
  \proofstep{2}{n\in\mathbb{N}}{\rA}
  \proofstep{3}{n\in x_0}{\rA}

  \proofstep{1,2}{\Jomit{x_0}(n)=\tJomit{x_0}(n)}{%
    \FormulaRefAuto{%
      x_0\in\mathbb{N}\dsep n\in\mathbb{N}\vdash \Jomit{x_0}(n)=\tJomit{x_0}(n)}{1,2}}

  \proofstep{2,3}{\tJomit{x_0}(n)=n}{%
    \FormulaRefAuto{n\in\mathbb{N}\dsep n\in x_0 \vdash \tJomit{x_0}(n)=n}{2,3}}

  \proofstep{1,2,3}{\Jomit{x_0}(n)=n}{%
    \FormulaRefAuto{a = b,\, b = c \vdash a = c}{4,5}}
\end{tabproof}

\FormulaThmDelta[$\Jomit{x_0}$ im strikten Ordnungsfall]{%
  x_0\in\mathbb{N}\dsep n\in\mathbb{N}\dsep n < x_0
  \vdash \Jomit{x_0}(n)=n%
}{%
  \DeltaDecl{Mengen}{x_0\dsep n}%
  \DeltaDecl{Funktionensymbol}{\Jomit{x_0}}%
  \DeltaRow{Relationen}{<_{\mathbb{N}}}%
}
\begin{tabproof}
  \proofstep{1}{x_0\in\mathbb{N}}{\rA}
  \proofstep{2}{n\in\mathbb{N}}{\rA}
  \proofstep{3}{n < x_0}{\rA}

  \proofstep{1,2,3}{n\in x_0}{%
    \FormulaRefAuto{m\in\mathbb{N}\dsep n\in\mathbb{N}\vdash m < n \leftrightarrow m\in n}{2,1,3}}
  \proofstep{1,2,3}{\Jomit{x_0}(n)=n}{%
    \FormulaRefAuto{x_0\in\mathbb{N}\dsep n\in\mathbb{N}\dsep n\in x_0 \vdash \Jomit{x_0}(n)=n}{1,2,4}}
\end{tabproof}


\FormulaThmDelta[$\Jomit{x_0}$ im Nicht-Elementfall]{%
  x_0\in\mathbb{N}\dsep n\in\mathbb{N}\dsep n\notin x_0 \vdash \Jomit{x_0}(n)=\Succ(n)%
}{%
  \DeltaDecl{Mengen}{x_0\dsep n}%
  \DeltaDecl{Funktionensymbol}{\Jomit{x_0}}%
}
\begin{tabproof}
  \proofstep{1}{x_0\in\mathbb{N}}{\rA}
  \proofstep{2}{n\in\mathbb{N}}{\rA}
  \proofstep{3}{n\notin x_0}{\rA}

  \proofstep{1,2}{\Jomit{x_0}(n)=\tJomit{x_0}(n)}{%
    \FormulaRefAuto{%
      x_0\in\mathbb{N}\dsep n\in\mathbb{N}\vdash \Jomit{x_0}(n)=\tJomit{x_0}(n)}{1,2}}

  \proofstep{2,3}{\tJomit{x_0}(n)=\Succ(n)}{%
    \FormulaRefAuto{n\in\mathbb{N}\dsep n\notin x_0 \vdash \tJomit{x_0}(n)=\Succ(n)}{2,3}}

  \proofstep{1,2,3}{\Jomit{x_0}(n)=\Succ(n)}{%
    \FormulaRefAuto{a = b,\, b = c \vdash a = c}{4,5}}
\end{tabproof}

\FormulaThmDelta[$\Jomit{x_0}$ im nicht-strikten Ordnungsfall]{%
  x_0\in\mathbb{N}\dsep n\in\mathbb{N}\dsep x_0\leq n
  \vdash \Jomit{x_0}(n)=\Succ(n)%
}{%
  \DeltaDecl{Mengen}{x_0\dsep n}%
  \DeltaDecl{Funktionensymbol}{\Jomit{x_0}}%
  \DeltaRow{Relationen}{\leq_{\mathbb{N}}}%
}
\begin{tabproof}
  \proofstep{1}{x_0\in\mathbb{N}}{\rA}
  \proofstep{2}{n\in\mathbb{N}}{\rA}
  \proofstep{3}{x_0\leq n}{\rA}

  \proofstep{4}{n\in x_0}{\rA}
  \proofstep{1,2,3}{x_0\subseteq n}{%
    \FormulaRefAuto{m\in\mathbb{N}\dsep n\in\mathbb{N}\vdash m\leq n \leftrightarrow m\subseteq n}{1,2,3}}
  \proofstep{1,2,3,4}{n\in n}{%
    \FormulaRefAuto{A\subseteq B,\, x\in A \vdash x\in B}{5,4}}
  \proofstep{2}{n\notin n}{%
    \FormulaRefAuto{n\in\mathbb{N}\vdash n\notin n}{2}}
  \proofstep{1,2,3,4}{\bot}{\rBI{6,7}}
  \proofstep{1,2,3}{n\notin x_0}{\rCI{4,8}}
  \proofstep{1,2,3}{\Jomit{x_0}(n)=\Succ(n)}{%
    \FormulaRefAuto{x_0\in\mathbb{N}\dsep n\in\mathbb{N}\dsep n\notin x_0 \vdash \Jomit{x_0}(n)=\Succ(n)}{1,2,9}}
\end{tabproof}

\FormulaThmDeltaK[Ordnungstrennung der Auslassungsabbildung]{%
  x_0\in\mathbb{N}\vdash
  \forall k,n\in\mathbb{N}\,
  \bigl(k<n\rightarrow
    \Jomit{x_0}(k)\neq\Jomit{x_0}(n)\bigr)%
}{JomitOrderSeparation}{%
  \DeltaRow{Mengen}{x_0\dsep k\dsep n}%
  \DeltaRow{Funktionensymbol}{\Jomit{x_0}}%
  \DeltaRow{Relationen}{<_{\mathbb{N}}\dsep\leq_{\mathbb{N}}}%
}
\begingroup
\small
\begin{tabproofsplitwide}
  \proofpartwideindR[Fall $k<n<x_0$]{%
    \begin{aligned}[t]
      &x_0\in\mathbb{N}\dsep k\in\mathbb{N}\dsep n\in\mathbb{N}\dsep
      k<n\dsep k<x_0\dsep n<x_0\vdash{}\\
      &\Jomit{x_0}(k)\neq\Jomit{x_0}(n)
    \end{aligned}
  }{%
    x_0\in\mathbb{N}\dsep k\in\mathbb{N}\dsep n\in\mathbb{N}\dsep
    k<n\dsep k<x_0\dsep n<x_0\vdash
    \Jomit{x_0}(k)\neq\Jomit{x_0}(n)%
  }
    \proofstepwidestar[1]{x_0\in\mathbb{N}}{\rA}
    \proofstepwidestar[2]{k\in\mathbb{N}}{\rA}
    \proofstepwidestar[3]{n\in\mathbb{N}}{\rA}
    \proofstepwidestar[4]{k<n}{\rA}
    \proofstepwidestar[5]{k<x_0}{\rA}
    \proofstepwidestar[6]{n<x_0}{\rA}
    \proofstepwidestar[1,2,5]{\Jomit{x_0}(k)=k}{%
      \FormulaRefAuto{x_0\in\mathbb{N}\dsep n\in\mathbb{N}\dsep n < x_0
        \vdash \Jomit{x_0}(n)=n}{1,2,5}}
    \proofstepwidestar[1,3,6]{\Jomit{x_0}(n)=n}{%
      \FormulaRefAuto{x_0\in\mathbb{N}\dsep n\in\mathbb{N}\dsep n < x_0
        \vdash \Jomit{x_0}(n)=n}{1,3,6}}
    \proofstepwidestar[2,3,4]{k\neq n}{%
      \FormulaRefAuto{m\in\mathbb{N}\dsep n\in\mathbb{N}\dsep
        m < n\vdash m\neq n}{2,3,4}}
    \proofstepwidestar[1,2,3,4,5,6]{%
      \Jomit{x_0}(k)\neq\Jomit{x_0}(n)}{%
      \FormulaRefAuto{a=b\dsep c=d\dsep P(b,d)\vdash P(a,c)}{7,8,9}}
  \closeproofpartwideindR

  \proofpartwideindR[Fall $k<x_0\leq n$]{%
    \begin{aligned}[t]
      &x_0\in\mathbb{N}\dsep k\in\mathbb{N}\dsep n\in\mathbb{N}\dsep
      k<n\dsep k<x_0\dsep x_0\leq n\vdash{}\\
      &\Jomit{x_0}(k)\neq\Jomit{x_0}(n)
    \end{aligned}
  }{%
    x_0\in\mathbb{N}\dsep k\in\mathbb{N}\dsep n\in\mathbb{N}\dsep
    k<n\dsep k<x_0\dsep x_0\leq n\vdash
    \Jomit{x_0}(k)\neq\Jomit{x_0}(n)%
  }
    \proofstepwidestar[1]{x_0\in\mathbb{N}}{\rA}
    \proofstepwidestar[2]{k\in\mathbb{N}}{\rA}
    \proofstepwidestar[3]{n\in\mathbb{N}}{\rA}
    \proofstepwidestar[4]{k<n}{\rA}
    \proofstepwidestar[5]{k<x_0}{\rA}
    \proofstepwidestar[6]{x_0\leq n}{\rA}
    \proofstepwidestar[1,2,5]{\Jomit{x_0}(k)=k}{%
      \FormulaRefAuto{x_0\in\mathbb{N}\dsep n\in\mathbb{N}\dsep n < x_0
        \vdash \Jomit{x_0}(n)=n}{1,2,5}}
    \proofstepwidestar[1,3,6]{\Jomit{x_0}(n)=\Succ(n)}{%
      \FormulaRefAuto{x_0\in\mathbb{N}\dsep n\in\mathbb{N}\dsep x_0\leq n
        \vdash \Jomit{x_0}(n)=\Succ(n)}{1,3,6}}
    \proofstepwidestar[2,3,4]{k\leq n}{%
      \FormulaRefAuto{m\in\mathbb{N}\dsep n\in\mathbb{N}\dsep
        m < n\vdash m\leq n}{2,3,4}}
    \proofstepwidestar[2,3,4]{k\in\Succ(n)}{%
      \FormulaRefAuto{x\in\mathbb{N}\dsep n\in\mathbb{N}\vdash
        x\in\Succ(n)\leftrightarrow x\leq n}{2,3,9}}
    \proofstepwidestar[3]{\Succ(n)\in\mathbb{N}}{%
      \FormulaRefAuto{PeanoSuccClosure}{3}}
    \proofstepwidestar[2,3,4]{k\neq\Succ(n)}{%
      \FormulaRefAuto{m\in\mathbb{N}\dsep n\in\mathbb{N}\dsep
        m\in n\vdash m\neq n}{2,11,10}}
    \proofstepwidestar[1,2,3,4,5,6]{%
      \Jomit{x_0}(k)\neq\Jomit{x_0}(n)}{%
      \FormulaRefAuto{a=b\dsep c=d\dsep P(b,d)\vdash P(a,c)}{7,8,12}}
  \closeproofpartwideindR

  \proofpartwideindR[Fall $x_0\leq k<n$]{%
    \begin{aligned}[t]
      &x_0\in\mathbb{N}\dsep k\in\mathbb{N}\dsep n\in\mathbb{N}\dsep
      k<n\dsep x_0\leq k\vdash{}\\
      &\Jomit{x_0}(k)\neq\Jomit{x_0}(n)
    \end{aligned}
  }{%
    x_0\in\mathbb{N}\dsep k\in\mathbb{N}\dsep n\in\mathbb{N}\dsep
    k<n\dsep x_0\leq k\vdash
    \Jomit{x_0}(k)\neq\Jomit{x_0}(n)%
  }
    \proofstepwidestar[1]{x_0\in\mathbb{N}}{\rA}
    \proofstepwidestar[2]{k\in\mathbb{N}}{\rA}
    \proofstepwidestar[3]{n\in\mathbb{N}}{\rA}
    \proofstepwidestar[4]{k<n}{\rA}
    \proofstepwidestar[5]{x_0\leq k}{\rA}
    \proofstepwidestar[2,3,4]{k\leq n}{%
      \FormulaRefAuto{m\in\mathbb{N}\dsep n\in\mathbb{N}\dsep
        m < n\vdash m\leq n}{2,3,4}}
    \proofstepwidestar[1,2,3,4,5]{x_0\leq n}{%
      \FormulaRefAuto{m\in\mathbb{N}\dsep n\in\mathbb{N}\dsep
        k\in\mathbb{N}\dsep m\leq n\dsep n\leq k\vdash m\leq k}{1,2,3,5,6}}
    \proofstepwidestar[1,2,5]{\Jomit{x_0}(k)=\Succ(k)}{%
      \FormulaRefAuto{x_0\in\mathbb{N}\dsep n\in\mathbb{N}\dsep x_0\leq n
        \vdash \Jomit{x_0}(n)=\Succ(n)}{1,2,5}}
    \proofstepwidestar[1,2,3,4,5]{\Jomit{x_0}(n)=\Succ(n)}{%
      \FormulaRefAuto{x_0\in\mathbb{N}\dsep n\in\mathbb{N}\dsep x_0\leq n
        \vdash \Jomit{x_0}(n)=\Succ(n)}{1,3,7}}
    \proofstepwidestar[2,3,4]{k\neq n}{%
      \FormulaRefAuto{m\in\mathbb{N}\dsep n\in\mathbb{N}\dsep
        m < n\vdash m\neq n}{2,3,4}}
    \proofstepwidestar[11]{\Succ(k)=\Succ(n)}{\rA}
    \proofstepwidestar[2,3,11]{k=n}{%
      \FormulaRefAuto{PeanoSuccInjective}{2,3,11}}
    \proofstepwidestar[2,3,4,11]{\bot}{\rBI{12,10}}
    \proofstepwidestar[2,3,4]{\Succ(k)\neq\Succ(n)}{\rCI{11,13}}
    \proofstepwidestar[1,2,3,4,5]{%
      \Jomit{x_0}(k)\neq\Jomit{x_0}(n)}{%
      \FormulaRefAuto{a=b\dsep c=d\dsep P(b,d)\vdash P(a,c)}{8,9,14}}
  \closeproofpartwideindR

  \proofpartwide{Zusammenfassung}
    \proofstepwidestar[1]{x_0\in\mathbb{N}}{\rA}
    \proofstepwidestar[2]{k\in\mathbb{N}}{\rA}
    \proofstepwidestar[3]{n\in\mathbb{N}}{\rA}
    \proofstepwidestar[4]{k<n}{\rA}

    \proofstepwidestar{x_0\leq k\lor\neg(x_0\leq k)}{%
      \FormulaRefAuto{P\lor\neg P}}

    \proofstepwidestar[6]{x_0\leq k}{\rA}
    \proofstepwidestar[1,2,3,4,6]{%
      \Jomit{x_0}(k)\neq\Jomit{x_0}(n)}{%
      \FormulaRefAuto{%
        x_0\in\mathbb{N}\dsep k\in\mathbb{N}\dsep n\in\mathbb{N}\dsep
        k<n\dsep x_0\leq k\vdash
        \Jomit{x_0}(k)\neq\Jomit{x_0}(n)%
      }{1,2,3,4,6}}

    \proofstepwidestar[8]{\neg(x_0\leq k)}{\rA}
    \proofstepwidestar[1,2,8]{k<x_0}{%
      \FormulaRefAuto{m\in\mathbb{N}\dsep n\in\mathbb{N}\dsep
        \neg(m\leq n)\vdash n < m}{1,2,8}}
    \proofstepwidestar{x_0\leq n\lor\neg(x_0\leq n)}{%
      \FormulaRefAuto{P\lor\neg P}}

    \proofstepwidestar[11]{x_0\leq n}{\rA}
    \proofstepwidestar[1,2,3,4,8,11]{%
      \Jomit{x_0}(k)\neq\Jomit{x_0}(n)}{%
      \FormulaRefAuto{%
        x_0\in\mathbb{N}\dsep k\in\mathbb{N}\dsep n\in\mathbb{N}\dsep
        k<n\dsep k<x_0\dsep x_0\leq n\vdash
        \Jomit{x_0}(k)\neq\Jomit{x_0}(n)%
      }{1,2,3,4,9,11}}

    \proofstepwidestar[13]{\neg(x_0\leq n)}{\rA}
    \proofstepwidestar[1,3,13]{n<x_0}{%
      \FormulaRefAuto{m\in\mathbb{N}\dsep n\in\mathbb{N}\dsep
        \neg(m\leq n)\vdash n < m}{1,3,13}}
    \proofstepwidestar[1,2,3,4,8,13]{%
      \Jomit{x_0}(k)\neq\Jomit{x_0}(n)}{%
      \FormulaRefAuto{%
        x_0\in\mathbb{N}\dsep k\in\mathbb{N}\dsep n\in\mathbb{N}\dsep
        k<n\dsep k<x_0\dsep n<x_0\vdash
        \Jomit{x_0}(k)\neq\Jomit{x_0}(n)%
      }{1,2,3,4,9,14}}
    \proofstepwidestar[1,2,3,4,8]{%
      \Jomit{x_0}(k)\neq\Jomit{x_0}(n)}{\rOE{10,11,12,13,15}}

    \proofstepwidestar[1,2,3,4]{%
      \Jomit{x_0}(k)\neq\Jomit{x_0}(n)}{\rOE{5,6,7,8,16}}

    \proofstepwidestar[1]{%
      \forall k,n\in\mathbb{N}\,
      \bigl(k<n\rightarrow
        \Jomit{x_0}(k)\neq\Jomit{x_0}(n)\bigr)}{%
      \rUI{\rRI{2,\rRI{3,\rRI{4,17}}}}}
  \closeproofpartwide
\end{tabproofsplitwide}
\endgroup

\FormulaThmDelta[Bild der Auslassungsabbildung liegt im Ziel]{%
  x_0\in\mathbb{N}\dsep n\in\mathbb{N}\vdash
  \Jomit{x_0}(n)\in \SOmit{x_0}%
}{%
  \DeltaRow{Mengen}{x_0\dsep n}%
}
\begin{tabproof}
  \proofstep{1}{x_0\in\mathbb{N}}{\rA}
  \proofstep{2}{n\in\mathbb{N}}{\rA}

  \proofstep{1,2}{\tJomit{x_0}(n)\in \SOmit{x_0}}{%
    \FormulaRefAuto{%
      x_0\in\mathbb{N}\dsep n\in\mathbb{N}\vdash
      \tJomit{x_0}(n)\in \SOmit{x_0}}{1,2}}

  \proofstep{1,2}{\Jomit{x_0}(n)=\tJomit{x_0}(n)}{%
    \FormulaRefAuto{%
      x_0\in\mathbb{N}\dsep n\in\mathbb{N}\vdash
      \Jomit{x_0}(n)=\tJomit{x_0}(n)}{1,2}}

  \proofstep{1,2}{\Jomit{x_0}(n)\in \SOmit{x_0}}{\rIE{4,3}}
\end{tabproof}

\FormulaThmDelta[Injektivität von $\Jomit{x_0}$]{%
  x_0\in\mathbb{N}\dsep x\in\mathbb{N}\dsep y\in\mathbb{N}\dsep
  \Jomit{x_0}(x)=\Jomit{x_0}(y)\vdash x=y%
}{%
  \DeltaRow{Mengen}{x_0\dsep x\dsep y}%
}
\begin{tabproof}
  \proofstep{1}{x_0\in\mathbb{N}}{\rA}
  \proofstep{2}{x\in\mathbb{N}}{\rA}
  \proofstep{3}{y\in\mathbb{N}}{\rA}
  \proofstep{4}{\Jomit{x_0}(x)=\Jomit{x_0}(y)}{\rA}
  \proofstep{1}{\Jomit{x_0}\colon\mathbb{N}\to\SOmit{x_0}}{%
    \FormulaRefAuto{%
      x_0\in\mathbb{N}\vdash
      \Jomit{x_0}\colon\mathbb{N}\to\SOmit{x_0}}{1}}
  \proofstep{1}{%
    \forall k,n\in\mathbb{N}\,
    \bigl(k<n\rightarrow
      \Jomit{x_0}(k)\neq\Jomit{x_0}(n)\bigr)}{%
    \FormulaRefAuto{JomitOrderSeparation}[thm]{1}}
  \proofstep{1}{\Jomit{x_0}\colon\mathbb{N}\inj\SOmit{x_0}}{%
    \FormulaRefAuto{%
      F\colon\mathbb{N}\to M\dsep
      \forall k,n\in\mathbb{N}\,\bigl(k<n\rightarrow F(k)\neq F(n)\bigr)
      \vdash F\colon\mathbb{N}\inj M}{5,6}}
  \proofstep{1,2,3,4}{x=y}{%
    \FormulaRefAuto{%
      F\colon A\inj B\dsep x\in A\dsep y\in A\dsep
      F(x)=F(y)\vdash x=y}{7,2,3,4}}
\end{tabproof}

\FormulaThmDeltaR[Fall $x<x_0,\; y<x_0$ der Injektivität von $\Jomit{x_0}$]{%
  \begin{aligned}[t]
    &x_0\in\mathbb{N}\dsep x\in\mathbb{N}\dsep y\in\mathbb{N}\dsep{}\\
    &\Jomit{x_0}(x)=\Jomit{x_0}(y)\dsep x<x_0\dsep y<x_0 \vdash x=y
  \end{aligned}
}{%
  x_0\in\mathbb{N}\dsep x\in\mathbb{N}\dsep y\in\mathbb{N}\dsep
  \Jomit{x_0}(x)=\Jomit{x_0}(y)\dsep x<x_0\dsep y<x_0 \vdash x=y
}{%
  \DeltaRow{Mengen}{x_0\dsep x\dsep y}%
  \DeltaRow{Relationen}{<_{\mathbb{N}}}%
}
\begin{tabproof}
  \proofstep{1}{x_0\in\mathbb{N}}{\rA}
  \proofstep{2}{x\in\mathbb{N}}{\rA}
  \proofstep{3}{y\in\mathbb{N}}{\rA}
  \proofstep{4}{\Jomit{x_0}(x)=\Jomit{x_0}(y)}{\rA}
  \proofstep{5}{x<x_0}{\rA}
  \proofstep{6}{y<x_0}{\rA}
  \proofstep{1,2,3,4,5,6}{x=y}{%
    \FormulaRefAuto{%
      x_0\in\mathbb{N}\dsep x\in\mathbb{N}\dsep y\in\mathbb{N}\dsep
      \Jomit{x_0}(x)=\Jomit{x_0}(y)\vdash x=y}{1,2,3,4}}
\end{tabproof}

\FormulaThmDeltaR[Fall $x<x_0,\; x_0\leq y$ der Injektivität von $\Jomit{x_0}$]{%
  \begin{aligned}[t]
    &x_0\in\mathbb{N}\dsep x\in\mathbb{N}\dsep y\in\mathbb{N}\dsep{}\\
    &\Jomit{x_0}(x)=\Jomit{x_0}(y)\dsep x<x_0\dsep x_0\leq y \vdash x=y
  \end{aligned}
}{%
  x_0\in\mathbb{N}\dsep x\in\mathbb{N}\dsep y\in\mathbb{N}\dsep
  \Jomit{x_0}(x)=\Jomit{x_0}(y)\dsep x<x_0\dsep x_0\leq y \vdash x=y
}{%
  \DeltaRow{Mengen}{x_0\dsep x\dsep y}%
  \DeltaRow{Relationen}{<_{\mathbb{N}}\dsep \leq_{\mathbb{N}}}%
}
\begin{tabproof}
  \proofstep{1}{x_0\in\mathbb{N}}{\rA}
  \proofstep{2}{x\in\mathbb{N}}{\rA}
  \proofstep{3}{y\in\mathbb{N}}{\rA}
  \proofstep{4}{\Jomit{x_0}(x)=\Jomit{x_0}(y)}{\rA}
  \proofstep{5}{x<x_0}{\rA}
  \proofstep{6}{x_0\leq y}{\rA}
  \proofstep{1,2,3,4,5,6}{x=y}{%
    \FormulaRefAuto{%
      x_0\in\mathbb{N}\dsep x\in\mathbb{N}\dsep y\in\mathbb{N}\dsep
      \Jomit{x_0}(x)=\Jomit{x_0}(y)\vdash x=y}{1,2,3,4}}
\end{tabproof}

\FormulaThmDeltaR[Fall $x_0\leq x,\; y<x_0$ der Injektivität von $\Jomit{x_0}$]{%
  \begin{aligned}[t]
    &x_0\in\mathbb{N}\dsep x\in\mathbb{N}\dsep y\in\mathbb{N}\dsep{}\\
    &\Jomit{x_0}(x)=\Jomit{x_0}(y)\dsep x_0\leq x\dsep y<x_0 \vdash x=y
  \end{aligned}
}{%
  x_0\in\mathbb{N}\dsep x\in\mathbb{N}\dsep y\in\mathbb{N}\dsep
  \Jomit{x_0}(x)=\Jomit{x_0}(y)\dsep x_0\leq x\dsep y<x_0 \vdash x=y
}{%
  \DeltaRow{Mengen}{x_0\dsep x\dsep y}%
  \DeltaRow{Relationen}{<_{\mathbb{N}}\dsep \leq_{\mathbb{N}}}%
}
\begin{tabproof}
  \proofstep{1}{x_0\in\mathbb{N}}{\rA}
  \proofstep{2}{x\in\mathbb{N}}{\rA}
  \proofstep{3}{y\in\mathbb{N}}{\rA}
  \proofstep{4}{\Jomit{x_0}(x)=\Jomit{x_0}(y)}{\rA}
  \proofstep{5}{x_0\leq x}{\rA}
  \proofstep{6}{y<x_0}{\rA}
  \proofstep{1,2,3,4,5,6}{x=y}{%
    \FormulaRefAuto{%
      x_0\in\mathbb{N}\dsep x\in\mathbb{N}\dsep y\in\mathbb{N}\dsep
      \Jomit{x_0}(x)=\Jomit{x_0}(y)\vdash x=y}{1,2,3,4}}
\end{tabproof}

\FormulaThmDeltaR[Fall $x_0\leq x,\; x_0\leq y$ der Injektivität von $\Jomit{x_0}$]{%
  \begin{aligned}[t]
    &x_0\in\mathbb{N}\dsep x\in\mathbb{N}\dsep y\in\mathbb{N}\dsep{}\\
    &\Jomit{x_0}(x)=\Jomit{x_0}(y)\dsep x_0\leq x\dsep x_0\leq y \vdash x=y
  \end{aligned}
}{%
  x_0\in\mathbb{N}\dsep x\in\mathbb{N}\dsep y\in\mathbb{N}\dsep
  \Jomit{x_0}(x)=\Jomit{x_0}(y)\dsep x_0\leq x\dsep x_0\leq y \vdash x=y
}{%
  \DeltaRow{Mengen}{x_0\dsep x\dsep y}%
  \DeltaRow{Relationen}{\leq_{\mathbb{N}}}%
}
\begin{tabproof}
  \proofstep{1}{x_0\in\mathbb{N}}{\rA}
  \proofstep{2}{x\in\mathbb{N}}{\rA}
  \proofstep{3}{y\in\mathbb{N}}{\rA}
  \proofstep{4}{\Jomit{x_0}(x)=\Jomit{x_0}(y)}{\rA}
  \proofstep{5}{x_0\leq x}{\rA}
  \proofstep{6}{x_0\leq y}{\rA}
  \proofstep{1,2,3,4,5,6}{x=y}{%
    \FormulaRefAuto{%
      x_0\in\mathbb{N}\dsep x\in\mathbb{N}\dsep y\in\mathbb{N}\dsep
      \Jomit{x_0}(x)=\Jomit{x_0}(y)\vdash x=y}{1,2,3,4}}
\end{tabproof}

% =========================================================
% Folgesätze: Aus der bereits konstruierten Auslassungsabbildung
%             folgt Injektivität und dann die Existenzbehauptung.
% Empfohlener Ort: direkt nach "Injektivität von \Jomit{x_0}".
% =========================================================

\FormulaThmDelta[$\Jomit{x_0}$ als injektive Funktion]{%
  x_0\in\mathbb{N}\vdash
  \Jomit{x_0}\colon \mathbb{N}\inj \SOmit{x_0}%
}{%
  \DeltaDecl{Mengen}{x_0}%
}
\begin{tabproof}
  \proofstep{1}{x_0\in\mathbb{N}}{\rA}
  \proofstep{1}{\Jomit{x_0}\colon \mathbb{N}\to \SOmit{x_0}}{%
    \FormulaRefAuto{%
      x_0\in\mathbb{N}\vdash
      \Jomit{x_0}\colon \mathbb{N}\to \SOmit{x_0}}{1}}
  \proofstep{1}{%
    \forall k,n\in\mathbb{N}\,
    \bigl(k<n\rightarrow
      \Jomit{x_0}(k)\neq\Jomit{x_0}(n)\bigr)}{%
    \FormulaRefAuto{JomitOrderSeparation}[thm]{1}}
  \proofstep{1}{\Jomit{x_0}\colon \mathbb{N}\inj \SOmit{x_0}}{%
    \FormulaRefAuto{%
      F\colon\mathbb{N}\to M\dsep
      \forall k,n\in\mathbb{N}\,\bigl(k<n\rightarrow F(k)\neq F(n)\bigr)
      \vdash F\colon\mathbb{N}\inj M}{2,3}}
\end{tabproof}


\FormulaThmDelta[Entfernen eines Elements aus $\mathbb{N}$]{%
  x_0\in\mathbb{N}\vdash \exists J\colon \mathbb{N}\inj \SOmit{x_0}%
}{%
  \DeltaDecl{Mengen}{x_0}%
  \DeltaDecl{Funktionen}{J}%
}
\begin{tabproof}
  \proofstep{1}{x_0\in\mathbb{N}}{\rA}
  \proofstep{1}{\Jomit{x_0}\colon \mathbb{N}\inj \SOmit{x_0}}{%
    \FormulaRefAuto{%
      x_0\in\mathbb{N}\vdash
      \Jomit{x_0}\colon \mathbb{N}\inj \SOmit{x_0}}{1}}
  \proofstep{1}{\exists J\colon \mathbb{N}\inj \SOmit{x_0}}{\rEI{2}}
\end{tabproof}


\subsection{Abstieg vom Nachfolger}
% ---------------------------------------------------------
% 1) Komposition ist Abbildung nach \Succ(n)
% ---------------------------------------------------------
\FormulaThmDelta[Komposition liefert Abbildung nach $\Succ(n)$]{%
  F\colon \mathbb{N}\to \Succ(n)\dsep
  x_0\in\mathbb{N}
  \vdash
  \bigl((F\circ \iota_{\SOmit{x_0},\mathbb{N}})\circ \Jomit{x_0}\bigr)\colon \mathbb{N}\to \Succ(n)%
}{%
  \DeltaDecl{Mengen}{n\dsep x_0}%
  \DeltaDecl{Funktionen}{F}%
}
\begin{tabproof}
  \proofstep{1}{F\colon \mathbb{N}\to \Succ(n)}{\rA}
  \proofstep{2}{x_0\in\mathbb{N}}{\rA}

  \proofstep{2}{\SOmit{x_0}\subseteq\mathbb{N}}{%
    \FormulaRefAuto{x_0\in\mathbb{N}\vdash \SOmit{x_0}\subseteq\mathbb{N}}{2}}

  \proofstep{2}{\Jomit{x_0}\colon \mathbb{N}\to \SOmit{x_0}}{%
    \FormulaRefAuto{x_0\in \mathbb{N}\vdash \Jomit {x_0}\colon \mathbb{N}\to \SOmit {x_0}}{2}}

  \proofstep{2}{ \iota_{\SOmit{x_0},\mathbb{N}}\colon \SOmit{x_0}\to \mathbb{N}}{%
    \FormulaRefAuto{A\subseteq B\vdash \iota _{A,B}\colon A\to B}{3}}

  \proofstep{1,2}{(F\circ \iota_{\SOmit{x_0},\mathbb{N}})\colon \SOmit{x_0}\to \Succ(n)}{%
    \FormulaRefAuto{G\circ F\colon A\to C}{1,5}}

  \proofstep{1,2}{\bigl((F\circ \iota_{\SOmit{x_0},\mathbb{N}})\circ \Jomit{x_0}\bigr)\colon \mathbb{N}\to \Succ(n)}{%
    \FormulaRefAuto{G\circ F\colon A\to C}{6,4}}
\end{tabproof}

% ---------------------------------------------------------
% Hilfslemma: Element von \SOmit{x_0} ist ungleich x_0
% ---------------------------------------------------------
\FormulaThmDelta{%
  x_0\in\mathbb{N}\dsep y\in \SOmit{x_0}\vdash y\neq x_0%
}{%
  \DeltaDecl{Mengen}{x_0\dsep y}%
}
\begin{tabproof}
  \proofstep{1}{x_0\in\mathbb{N}}{\rA}
  \proofstep{2}{y\in \SOmit{x_0}}{\rA}
  \proofstep{1}{\SOmit{x_0}=\{x\in\mathbb{N}\mid x\neq x_0\}}{%
    \FormulaRefAuto{x_0\in\mathbb{N}\vdash
      \SOmit{x_0}\coloneqq \{x\in\mathbb{N}\mid x\neq x_0\}}{1}}
  \proofstep{2,3}{y\in \{x\in\mathbb{N}\mid x\neq x_0\}}{\rIE{3,2}}
  \proofstep{4}{y\in\mathbb{N}\land y\neq x_0}{%
    \FormulaRefAuto{x\in\{u\in A\mid P(u)\}\eqvdash x\in A\land P(x)}{4}}
  \proofstep{5}{y\neq x_0}{\rAEb{5}}
\end{tabproof}


% ---------------------------------------------------------
% Typisierung der Abstiegs-Funktionsvorschrift nach n
% ---------------------------------------------------------
\FormulaThmDeltaR[Typisierungsaxiom]{%
  \begin{aligned}[t]
    &n\in\mathbb{N}\dsep
    F\colon \mathbb{N}\inj \Succ(n)\dsep
    x_0\in\mathbb{N}\dsep
    F(x_0)=n\dsep
    x\in\mathbb{N}\\
    &\vdash\ \bigl((F\circ \iota_{\SOmit{x_0},\mathbb{N}})\circ \Jomit{x_0}\bigr)(x)\in n
  \end{aligned}%
}{%
  n\in\mathbb{N}\dsep
  F\colon \mathbb{N}\inj \Succ(n)\dsep
  x_0\in\mathbb{N}\dsep
  F(x_0)=n\dsep
  x\in\mathbb{N}
  \vdash
  \bigl((F\circ \iota_{\SOmit{x_0},\mathbb{N}})\circ \Jomit{x_0}\bigr)(x)\in n%
}{%
  \DeltaDecl{Mengen}{n\dsep x_0\dsep x}%
  \DeltaDecl{Funktionen}{F}%
}
\begin{tabproof}
  \proofstep{1}{n\in\mathbb{N}}{\rA}
  \proofstep{2}{F\colon \mathbb{N}\inj \Succ(n)}{\rA}
  \proofstep{3}{x_0\in\mathbb{N}}{\rA}
  \proofstep{4}{F(x_0)=n}{\rA}
  \proofstep{5}{x\in\mathbb{N}}{\rA}

  \proofstep{2,3}{\bigl((F\circ \iota_{\SOmit{x_0},\mathbb{N}})\circ \Jomit{x_0}\bigr)\colon \mathbb{N}\to \Succ(n)}{%
    \FormulaRefAuto{F\colon \mathbb{N}\to \Succ(n)\dsep x_0\in\mathbb{N}\vdash
      \bigl((F\circ \iota_{\SOmit{x_0},\mathbb{N}})\circ \Jomit{x_0}\bigr)\colon \mathbb{N}\to \Succ(n)}{2,3}}

  \proofstep{2,3,5}{\bigl((F\circ \iota_{\SOmit{x_0},\mathbb{N}})\circ \Jomit{x_0}\bigr)(x)\in \Succ(n)}{%
    \FormulaRefAuto{F\colon A\to B\dsep x\in A\vdash F(x)\in B}{6,5}}

  \proofstep{3,5}{\Jomit{x_0}(x)\in \SOmit{x_0}}{%
    \FormulaRefAuto{x_0\in\mathbb{N}\dsep n\in\mathbb{N}\vdash \Jomit{x_0}(n)\in \SOmit{x_0}}{3,5}}
  \proofstep{3,5}{\Jomit{x_0}(x)\neq x_0}{%
    \FormulaRefAuto{x_0\in\mathbb{N}\dsep y\in \SOmit{x_0}\vdash y\neq x_0}{3,8}}

  \proofstep{10}{\bigl((F\circ \iota_{\SOmit{x_0},\mathbb{N}})\circ \Jomit{x_0}\bigr)(x)=n}{\rA}

  \proofstep{5}{\bigl((F\circ \iota_{\SOmit{x_0},\mathbb{N}})\circ \Jomit{x_0}\bigr)(x)=F\bigl(\iota_{\SOmit{x_0},\mathbb{N}}(\Jomit{x_0}(x))\bigr)}{%
    \FormulaRefAuto{x\in A \vdash ((F\circ G)\circ H)(x)=F(G(H(x)))}{5}}

  \proofstep{3}{\SOmit{x_0}\subseteq\mathbb{N}}{%
    \FormulaRefAuto{x_0\in\mathbb{N}\vdash \SOmit{x_0}\subseteq\mathbb{N}}{3}}
  \proofstep{3,5}{\iota_{\SOmit{x_0},\mathbb{N}}(\Jomit{x_0}(x))=\Jomit{x_0}(x)}{%
    \FormulaRefAuto{A\subseteq B \dsep x\in A \vdash \iota _{A,B}(x)=x}{12,8}}

  \proofstep{3,5}{F\bigl(\iota_{\SOmit{x_0},\mathbb{N}}(\Jomit{x_0}(x))\bigr)=F(\Jomit{x_0}(x))}{\rIE{13}}
  \proofstep{3,5}{\bigl((F\circ \iota_{\SOmit{x_0},\mathbb{N}})\circ \Jomit{x_0}\bigr)(x)=F(\Jomit{x_0}(x))}{%
    \FormulaRefAuto{a=b \dsep b=c \vdash a=c}{11,14}}

  \proofstep{3,5,10}{F(\Jomit{x_0}(x))=n}{%
    \FormulaRefAuto{a=b \dsep b=c \vdash a=c}{15,10}}

  \proofstep{4}{n=F(x_0)}{\FormulaRefAuto{a=b \vdash b=a}{4}}
  \proofstep{3,4,5,10}{F(\Jomit{x_0}(x))=F(x_0)}{%
    \FormulaRefAuto{a=b \dsep c=b \vdash a=c}{16,17}}

  \proofstep{3,5}{\Jomit{x_0}(x)\in \mathbb{N}}{%
    \FormulaRefAuto{A\subseteq B,\,z\in A\vdash z\in B}{12,8}}

  \proofstep{3,4,5,10}{\Jomit{x_0}(x)=x_0}{%
    \FormulaRefAuto{F\colon A\inj B\dsep x\in A\dsep y\in A\dsep F(x)=F(y)\vdash x=y}{2,19,3,18}}

  \proofstep{3,4,5,10}{\bot}{\rBI{9,20}}
  \proofstep{3,4,5}{\bigl((F\circ \iota_{\SOmit{x_0},\mathbb{N}})\circ \Jomit{x_0}\bigr)(x)\neq n}{\rCI{10,21}}

  \proofstep{1,2,3,4,5}{\bigl((F\circ \iota_{\SOmit{x_0},\mathbb{N}})\circ \Jomit{x_0}\bigr)(x)\in n}{%
    \FormulaRefAuto{n\in\mathbb{N}\dsep z\in\Succ(n)\dsep z\neq n\vdash z\in n}{1,7,22}}
\end{tabproof}

\FormulaThmDeltaR[Typisierungsaxiom in Ordnungssprache]{%
  \begin{aligned}[t]
    &n\in\mathbb{N}\dsep
    F\colon \mathbb{N}\inj \Succ(n)\dsep
    x_0\in\mathbb{N}\dsep
    F(x_0)=n\dsep
    x\in\mathbb{N}\\
    &\vdash\ \bigl((F\circ \iota_{\SOmit{x_0},\mathbb{N}})\circ \Jomit{x_0}\bigr)(x)<n
  \end{aligned}%
}{%
  n\in\mathbb{N}\dsep
  F\colon \mathbb{N}\inj \Succ(n)\dsep
  x_0\in\mathbb{N}\dsep
  F(x_0)=n\dsep
  x\in\mathbb{N}
  \vdash
  \bigl((F\circ \iota_{\SOmit{x_0},\mathbb{N}})\circ \Jomit{x_0}\bigr)(x)<n%
}{%
  \DeltaDecl{Mengen}{n\dsep x_0\dsep x}%
  \DeltaDecl{Funktionen}{F}%
  \DeltaDecl{Relation}{<_{\mathbb{N}}}%
}
\begin{tabproof}
  \proofstep{1}{n\in\mathbb{N}}{\rA}
  \proofstep{2}{F\colon \mathbb{N}\inj \Succ(n)}{\rA}
  \proofstep{3}{x_0\in\mathbb{N}}{\rA}
  \proofstep{4}{F(x_0)=n}{\rA}
  \proofstep{5}{x\in\mathbb{N}}{\rA}
  \proofstep{1,2,3,4,5}{\bigl((F\circ \iota_{\SOmit{x_0},\mathbb{N}})\circ \Jomit{x_0}\bigr)(x)\in n}{%
    \FormulaRefAuto{%
      n\in\mathbb{N}\dsep
      F\colon \mathbb{N}\inj \Succ(n)\dsep
      x_0\in\mathbb{N}\dsep
      F(x_0)=n\dsep
      x\in\mathbb{N}
      \vdash
      \bigl((F\circ \iota_{\SOmit{x_0},\mathbb{N}})\circ \Jomit{x_0}\bigr)(x)\in n%
    }{1,2,3,4,5}}
  \proofstep{1}{\NatSeg{n}\subseteq\mathbb{N}}{%
    \FormulaRefAuto{PeanoNatSegSubsetN}{1}}
  \proofstep{1,2,3,4,5}{\bigl((F\circ \iota_{\SOmit{x_0},\mathbb{N}})\circ \Jomit{x_0}\bigr)(x)\in\mathbb{N}}{%
    \FormulaRefAuto{A\subseteq B,\, x\in A\vdash x\in B}{7,6}}
  \proofstep{1,2,3,4,5}{\bigl((F\circ \iota_{\SOmit{x_0},\mathbb{N}})\circ \Jomit{x_0}\bigr)(x)<n}{%
    \FormulaRefAuto{m\in\mathbb{N}\dsep n\in\mathbb{N}\vdash m < n \leftrightarrow m\in n}{8,1,6}}
\end{tabproof}

% =========================================================
% (1) Universalisierung des Typisierungsaxioms (Werte liegen in n)
% =========================================================
\FormulaThmDeltaR{%
  \begin{aligned}[t]
    &n\in\mathbb{N}\dsep
    F\colon \mathbb{N}\inj \Succ(n)\dsep
    x_0\in\mathbb{N}\dsep
    F(x_0)=n\\
    &\vdash\ \forall x\in\mathbb{N}\,
      \bigl((F\circ \iota_{\SOmit{x_0},\mathbb{N}})\circ \Jomit{x_0}\bigr)(x)\in n
  \end{aligned}%
}{%
  n\in\mathbb{N}\dsep
  F\colon \mathbb{N}\inj \Succ(n)\dsep
  x_0\in\mathbb{N}\dsep
  F(x_0)=n
  \vdash
  \forall x\in\mathbb{N}\,
    \bigl((F\circ \iota_{\SOmit{x_0},\mathbb{N}})\circ \Jomit{x_0}\bigr)(x)\in n%
}{%
  \DeltaDecl{Mengen}{n\dsep x_0\dsep x}%
  \DeltaDecl{Funktionen}{F}%
}
\begin{tabproof}
  \proofstep{1}{n\in\mathbb{N}}{\rA}
  \proofstep{2}{F\colon \mathbb{N}\inj \Succ(n)}{\rA}
  \proofstep{3}{x_0\in\mathbb{N}}{\rA}
  \proofstep{4}{F(x_0)=n}{\rA}

  \proofstep{5}{x\in\mathbb{N}}{\rA}
  \proofstep{1,2,3,4,5}{\bigl((F\circ \iota_{\SOmit{x_0},\mathbb{N}})\circ \Jomit{x_0}\bigr)(x)\in n}{%
    \FormulaRefAuto{%
      n\in\mathbb{N}\dsep
      F\colon \mathbb{N}\inj \Succ(n)\dsep
      x_0\in\mathbb{N}\dsep
      F(x_0)=n\dsep
      x\in\mathbb{N}
      \vdash
      \bigl((F\circ \iota_{\SOmit{x_0},\mathbb{N}})\circ \Jomit{x_0}\bigr)(x)\in n%
    }{1,2,3,4,5}}

  \proofstep{1,2,3,4}{\forall x\in\mathbb{N}\,\bigl((F\circ \iota_{\SOmit{x_0},\mathbb{N}})\circ \Jomit{x_0}\bigr)(x)\in n}{%
    \rUI{\rRI{5,6}}}
\end{tabproof}

\FormulaThmDeltaR{%
  \begin{aligned}[t]
    &n\in\mathbb{N}\dsep
    F\colon \mathbb{N}\inj \Succ(n)\dsep
    x_0\in\mathbb{N}\dsep
    F(x_0)=n\\
    &\vdash\ \forall x\in\mathbb{N}\,
      \bigl((F\circ \iota_{\SOmit{x_0},\mathbb{N}})\circ \Jomit{x_0}\bigr)(x)<n
  \end{aligned}%
}{%
  n\in\mathbb{N}\dsep
  F\colon \mathbb{N}\inj \Succ(n)\dsep
  x_0\in\mathbb{N}\dsep
  F(x_0)=n
  \vdash
  \forall x\in\mathbb{N}\,
    \bigl((F\circ \iota_{\SOmit{x_0},\mathbb{N}})\circ \Jomit{x_0}\bigr)(x)<n%
}{%
  \DeltaDecl{Mengen}{n\dsep x_0\dsep x}%
  \DeltaDecl{Funktionen}{F}%
  \DeltaDecl{Relation}{<_{\mathbb{N}}}%
}
\begin{tabproof}
  \proofstep{1}{n\in\mathbb{N}}{\rA}
  \proofstep{2}{F\colon \mathbb{N}\inj \Succ(n)}{\rA}
  \proofstep{3}{x_0\in\mathbb{N}}{\rA}
  \proofstep{4}{F(x_0)=n}{\rA}

  \proofstep{5}{x\in\mathbb{N}}{\rA}
  \proofstep{1,2,3,4,5}{\bigl((F\circ \iota_{\SOmit{x_0},\mathbb{N}})\circ \Jomit{x_0}\bigr)(x)<n}{%
    \FormulaRefAuto{%
      n\in\mathbb{N}\dsep
      F\colon \mathbb{N}\inj \Succ(n)\dsep
      x_0\in\mathbb{N}\dsep
      F(x_0)=n\dsep
      x\in\mathbb{N}
      \vdash
      \bigl((F\circ \iota_{\SOmit{x_0},\mathbb{N}})\circ \Jomit{x_0}\bigr)(x)<n%
    }{1,2,3,4,5}}

  \proofstep{1,2,3,4}{\forall x\in\mathbb{N}\,\bigl((F\circ \iota_{\SOmit{x_0},\mathbb{N}})\circ \Jomit{x_0}\bigr)(x)<n}{%
    \rUI{\rRI{5,6}}}
\end{tabproof}


% =========================================================
% (2) Korestriktion: die Komposition ist eine Funktion nach n
% =========================================================
\FormulaThmDelta{%
  n\in\mathbb{N}\dsep
  F\colon \mathbb{N}\inj \Succ(n)\dsep
  x_0\in\mathbb{N}\dsep
  F(x_0)=n
  \vdash
  \bigl((F\circ \iota_{\SOmit{x_0},\mathbb{N}})\circ \Jomit{x_0}\bigr)\colon \mathbb{N}\to n%
}{%
  \DeltaDecl{Mengen}{n\dsep x_0}%
  \DeltaDecl{Funktionen}{F}%
}
\begin{tabproof}
  \proofstep{1}{n\in\mathbb{N}}{\rA}
  \proofstep{2}{F\colon \mathbb{N}\inj \Succ(n)}{\rA}
  \proofstep{3}{x_0\in\mathbb{N}}{\rA}
  \proofstep{4}{F(x_0)=n}{\rA}

  \proofstep{2}{F\colon \mathbb{N}\to \Succ(n)}{\FormulaRefAuto{Injektive Funktion}{2}}

  \proofstep{3,5}{\bigl((F\circ \iota_{\SOmit{x_0},\mathbb{N}})\circ \Jomit{x_0}\bigr)\colon \mathbb{N}\to \Succ(n)}{%
    \FormulaRefAuto{%
      F\colon \mathbb{N}\to \Succ(n)\dsep x_0\in\mathbb{N}
      \vdash
      \bigl((F\circ \iota_{\SOmit{x_0},\mathbb{N}})\circ \Jomit{x_0}\bigr)\colon \mathbb{N}\to \Succ(n)%
    }{5,3}}

  \proofstep{1,2,3,4}{\forall x\in\mathbb{N}\,\bigl((F\circ \iota_{\SOmit{x_0},\mathbb{N}})\circ \Jomit{x_0}\bigr)(x)\in n}{%
    \FormulaRefAuto{%
      n\in\mathbb{N}\dsep
      F\colon \mathbb{N}\inj \Succ(n)\dsep
      x_0\in\mathbb{N}\dsep
      F(x_0)=n
      \vdash
      \forall x\in\mathbb{N}\,
      \bigl((F\circ \iota_{\SOmit{x_0},\mathbb{N}})\circ \Jomit{x_0}\bigr)(x)\in n%
    }{1,2,3,4}}

  \proofstep{1,2,3,4}{\bigl((F\circ \iota_{\SOmit{x_0},\mathbb{N}})\circ \Jomit{x_0}\bigr)\colon \mathbb{N}\to n}{%
    \FormulaRefAuto{F\colon A\to B\dsep \forall x\in A\,F(x)\in C\vdash F\colon A\to C}{6,7}}
\end{tabproof}


% =========================================================
% (3) Inj-Korestriktion: aus Injektivität nach Succ(n) folgt Injektivität nach n
% =========================================================
\FormulaThmDelta{%
  n\in\mathbb{N}\dsep
  F\colon \mathbb{N}\inj \Succ(n)\dsep
  x_0\in\mathbb{N}\dsep
  F(x_0)=n
  \vdash
  \bigl((F\circ \iota_{\SOmit{x_0},\mathbb{N}})\circ \Jomit{x_0}\bigr)\colon \mathbb{N}\inj n%
}{%
  \DeltaDecl{Mengen}{n\dsep x_0}%
  \DeltaDecl{Funktionen}{F}%
}
\begin{tabproof}
  \proofstep{1}{n\in\mathbb{N}}{\rA}
  \proofstep{2}{F\colon \mathbb{N}\inj \Succ(n)}{\rA}
  \proofstep{3}{x_0\in\mathbb{N}}{\rA}
  \proofstep{4}{F(x_0)=n}{\rA}

  \proofstep{3}{\Jomit{x_0}\colon \mathbb{N}\inj \SOmit{x_0}}{%
    \FormulaRefAuto{x_0\in \mathbb{N}\vdash \Jomit{x_0}\colon \mathbb{N}\inj \SOmit{x_0}}{3}}
  \proofstep{3}{\SOmit{x_0}\subseteq \mathbb{N}}{%
    \FormulaRefAuto{x_0\in \mathbb{N}\vdash \SOmit{x_0}\subseteq \mathbb{N}}{3}}
  \proofstep{3}{\iota_{\SOmit{x_0},\mathbb{N}}\colon \SOmit{x_0}\inj \mathbb{N}}{%
    \FormulaRefAuto{A\subseteq B\vdash \iota_{A,B}\colon A\inj B}{6}}

  \proofstep{2,3}{F\circ \iota_{\SOmit{x_0},\mathbb{N}}\colon \SOmit{x_0}\inj \Succ(n)}{%
    \FormulaRefAuto{F\colon A\inj B \dsep G\colon B\inj C\vdash G\circ F\colon A\inj C}{7,2}}

  \proofstep{2,3}{\bigl((F\circ \iota_{\SOmit{x_0},\mathbb{N}})\circ \Jomit{x_0}\bigr)\colon \mathbb{N}\inj \Succ(n)}{%
    \FormulaRefAuto{F\colon A\inj B \dsep G\colon B\inj C\vdash G\circ F\colon A\inj C}{5,8}}

  \proofstep{1,2,3,4}{\forall x\in\mathbb{N}\,\bigl((F\circ \iota_{\SOmit{x_0},\mathbb{N}})\circ \Jomit{x_0}\bigr)(x)\in n}{%
    \FormulaRefAuto{%
      n\in\mathbb{N}\dsep
      F\colon \mathbb{N}\inj \Succ(n)\dsep
      x_0\in\mathbb{N}\dsep
      F(x_0)=n
      \vdash
      \forall x\in\mathbb{N}\,
      \bigl((F\circ \iota_{\SOmit{x_0},\mathbb{N}})\circ \Jomit{x_0}\bigr)(x)\in n%
    }{1,2,3,4}}

  \proofstep{1,2,3,4}{\bigl((F\circ \iota_{\SOmit{x_0},\mathbb{N}})\circ \Jomit{x_0}\bigr)\colon \mathbb{N}\inj n}{%
    \FormulaRefAuto{F\colon A\inj B\dsep \forall x\in A\,F(x)\in C\vdash F\colon A\inj C}{9,10}}
\end{tabproof}

% ---------------------------------------------------------
% (A) Abstieg einer Injektion nach n, wenn n nicht getroffen wird
% ---------------------------------------------------------
\FormulaThmDelta[Abstieg einer Injektion ohne Spitze]{%
  n\in\mathbb{N}\dsep
  H\colon A\inj \Succ(n)\dsep
  \forall x\in A\,\bigl(H(x)\neq n\bigr)
  \vdash
  H\colon A\inj n%
}{%
  \DeltaDecl{Mengen}{A\dsep n\dsep x}%
  \DeltaDecl{Funktionen}{H}%
}
\begin{tabproof}
  \proofstep{1}{n\in\mathbb{N}}{\rA}
  \proofstep{2}{H\colon A\inj \Succ(n)}{\rA}
  \proofstep{3}{\forall x\in A\,\bigl(H(x)\neq n\bigr)}{\rA}

  \proofstep{4}{x\in A}{\rA}
  \proofstep{2}{H\colon A\to \Succ(n)}{\FormulaRefAuto{Injektive Funktion}{2}}
  \proofstep{2,4}{H(x)\in \Succ(n)}{%
    \FormulaRefAuto{F\colon A\to B\dsep x\in A\vdash F(x)\in B}{5,4}}
  \proofstep{3,4}{H(x)\neq n}{\rRE{3,4}}
  \proofstep{1,2,3,4}{H(x)\in n}{%
    \FormulaRefAuto{n\in\mathbb{N}\dsep z\in\Succ(n)\dsep z\neq n\vdash z\in n}{1,6,7}}
  \proofstep{1,2,3}{\forall x\in A\,\bigl(H(x)\in n\bigr)}{\rUI{\rRI{4,8}}}

  \proofstep{1,2,3}{H\colon A\inj n}{%
    \FormulaRefAuto{F\colon A\inj B\dsep \forall x\in A\,F(x)\in C\vdash F\colon A\inj C}{2,9}}
\end{tabproof}


% ---------------------------------------------------------
% (B) Abstieg liefert Existenz einer Injektion nach n
% ---------------------------------------------------------
\FormulaThmDelta[Abstieg liefert Injektion nach $n$]{%
  n\in\mathbb{N}\dsep
  F\colon \mathbb{N}\inj \Succ(n)\dsep
  x_0\in\mathbb{N}\dsep
  F(x_0)=n
  \vdash
  \exists G\colon \mathbb{N}\inj n%
}{%
  \DeltaDecl{Mengen}{n\dsep x_0}%
  \DeltaDecl{Funktionen}{F\dsep G}%
}
\begin{tabproof}
  \proofstep{1}{n\in\mathbb{N}}{\rA}
  \proofstep{2}{F\colon \mathbb{N}\inj \Succ(n)}{\rA}
  \proofstep{3}{x_0\in\mathbb{N}}{\rA}
  \proofstep{4}{F(x_0)=n}{\rA}

  \proofstep{1,2,3,4}{%
    \bigl((F\circ \iota_{\SOmit{x_0},\mathbb{N}})\circ \Jomit{x_0}\bigr)\colon \mathbb{N}\inj n}{%
    \FormulaRefAuto{%
      n\in\mathbb{N}\dsep
      F\colon \mathbb{N}\inj \Succ(n)\dsep
      x_0\in\mathbb{N}\dsep
      F(x_0)=n
      \vdash
      \bigl((F\circ \iota_{\SOmit{x_0},\mathbb{N}})\circ \Jomit{x_0}\bigr)\colon \mathbb{N}\inj n%
    }{1,2,3,4}}

  \proofstep{1,2,3,4}{\exists G\colon \mathbb{N}\inj n}{\rEI{5}}
\end{tabproof}

\FormulaThmDelta[Keine Injektion von $\mathbb{N}$ nach $n$]{%
  n\in\mathbb{N}\vdash \neg\exists F\colon \mathbb{N}\inj n%
}{%
  \DeltaDecl{Mengen}{n\dsep x\dsep x_0}%
  \DeltaDecl{Funktionen}{F\dsep G}%
}
\begin{tabproofsplitwide}
  \proofpartwideind[Induktionsanfang]{\neg\exists F\colon \mathbb{N}\inj 0}
    \proofstepwidestar[1]{\exists F\colon \mathbb{N}\inj 0}{\rA}

    \proofstepwidestar[2]{F\colon \mathbb{N}\inj 0}{\rA}
    \proofstepwidestar[2]{F\colon \mathbb{N}\to 0}{\FormulaRefAuto{Injektive Funktion}{2}}

    \proofstepwidestar{0\in\mathbb{N}}{\FormulaRefAuto{PeanoZero}}
    \proofstepwidestar[2]{F(0)\in0}{%
      \FormulaRefAuto{F\colon A\to B\dsep x\in A\vdash F(x)\in B}{3,4}}

    \proofstepwidestar[2]{0\neq0}{%
      \FormulaRefAuto{u\in A\vdash A\neq \varnothing}{5}}
    \proofstepwidestar{0=0}{\rII}
    \proofstepwidestar[2]{\bot}{\rBI{6,7}}

    \proofstepwidestar[1]{\bot}{\rEE{1,2,8}}
    \proofstepwidestar{\neg\exists F\colon \mathbb{N}\inj 0}{\rCI{1,9}}
  \closeproofpartwideind

  \proofpartwideind[Induktionsschritt]{%
    n\in\mathbb{N}\dsep \neg\exists F\colon \mathbb{N}\inj n
    \vdash
    \neg\exists F\colon \mathbb{N}\inj \Succ(n)}
    \proofstepwidestar[1]{n\in\mathbb{N}}{\rA}
    \proofstepwidestar[2]{\neg\exists F\colon \mathbb{N}\inj n}{\rA}

    \proofstepwidestar[3]{\exists F\colon \mathbb{N}\inj \Succ(n)}{\rA}

    \proofstepwidestar[4]{F\colon \mathbb{N}\inj \Succ(n)}{\rA}

    \proofstepwidestar[5]{\neg\exists x_0\in\mathbb{N}\,\bigl(F(x_0)=n\bigr)}{\rA}

    \proofstepwidestar[6]{x\in\mathbb{N}}{\rA}
    \proofstepwidestar[7]{F(x)=n}{\rA}
    \proofstepwidestar[6,7]{x\in\mathbb{N}\land F(x)=n}{\rAI{6,7}}
    \proofstepwidestar[6,7]{\exists x_0\in\mathbb{N}\,\bigl(F(x_0)=n\bigr)}{\rEI{8}}
    \proofstepwidestar[5,6,7]{\bot}{\rBI{5,9}}
    \proofstepwidestar[5,6]{F(x)\neq n}{\rCI{7,10}}
    \proofstepwidestar[5]{\forall x\in\mathbb{N}\,\bigl(F(x)\neq n\bigr)}{\rUI{\rRI{6,11}}}

    \proofstepwidestar[1,4,5]{F\colon \mathbb{N}\inj n}{%
      \FormulaRefAuto{n\in\mathbb{N}\dsep H\colon A\inj \Succ(n)\dsep \forall x\in A\,\bigl(H(x)\neq n\bigr)\vdash H\colon A\inj n}{1,4,12}}
    \proofstepwidestar[1,4,5]{\exists G\colon \mathbb{N}\inj n}{\rEI{13}}
    \proofstepwidestar[1,2,4,5]{\bot}{\rBI{2,14}}

    \proofstepwidestar[1,2,4]{\exists x_0\in\mathbb{N}\,\bigl(F(x_0)=n\bigr)}{\rCI{5,15}}

    \proofstepwidestar[17]{x_0\in\mathbb{N}\land F(x_0)=n}{\rA}
    \proofstepwidestar[17]{x_0\in\mathbb{N}}{\rAEa{17}}
    \proofstepwidestar[17]{F(x_0)=n}{\rAEb{17}}

    \proofstepwidestar[1,4,17]{\exists G\colon \mathbb{N}\inj n}{%
      \FormulaRefAuto{n\in\mathbb{N}\dsep F\colon \mathbb{N}\inj \Succ(n)\dsep x_0\in\mathbb{N}\dsep F(x_0)=n\vdash \exists G\colon \mathbb{N}\inj n}{1,4,18,19}}
    \proofstepwidestar[1,2,4,17]{\bot}{\rBI{2,20}}
    \proofstepwidestar[1,2,4]{\bot}{\rEE{16,17,20}}

    \proofstepwidestar[1,2,3]{\bot}{\rEE{3,4,21}}
    \proofstepwidestar[1,2]{\neg\exists F\colon \mathbb{N}\inj \Succ(n)}{\rCI{3,22}}
  \closeproofpartwideind

  \proofpartwide{Induktionsschluss}
    \proofstepwidestar[1]{n\in\mathbb{N}}{\rA}
    \proofstepwidestar[1]{\neg\exists F\colon \mathbb{N}\inj n}{%
      \FormulaRefAuto{P(0)\dsep [x\in\mathbb{N},P(x)]\vdots P(\Succ(x))\dsep n\in\mathbb{N}\vdash P(n)}{%
        \FormulaRefAuto{\neg\exists F\colon \mathbb{N}\inj \varnothing},%
        \FormulaRefAuto{n\in\mathbb{N}\dsep \neg\exists F\colon \mathbb{N}\inj n\vdash \neg\exists F\colon \mathbb{N}\inj \Succ(n)},%
        1}}
  \closeproofpartwide
\end{tabproofsplitwide}

\fi

\subsection{Peano-Vorgaengerfunktion}

\FormulaDefDeltaK[Peano-Vorgaengerfunktion]{%
  n\in\mathbb{N}\vdash
  \Pred(n)\coloneqq
  \IfThenElse{n=0}{0}{\iota y\,\bigl(y\in\mathbb{N}\land n=\Succ(y)\bigr)}%
}{PeanoPred}{%
  \DeltaRow{Mengen}{n\dsep y}%
  \DeltaRow{Funktionensymbol}{\Pred}%
}

\FormulaThmDeltaK[Vorgaenger einer natuerlichen Zahl ist natuerlich]{%
  n\in\mathbb{N}\vdash \Pred(n)\in\mathbb{N}%
}{PeanoPredNatural}{%
  \DeltaRow{Mengen}{n}%
}
\begin{tabproof}
  \proofstep{1}{n\in\mathbb{N}}{\rA}
  \proofstep{1}{\Pred(n)\in\mathbb{N}}{\FormulaRefAuto{PeanoPred}{1}}
\end{tabproof}

\FormulaThmDeltaK[Vorgaenger eines Nachfolgers]{%
  n\in\mathbb{N}\vdash \Pred(\Succ(n))=n%
}{PeanoPredSucc}{%
  \DeltaRow{Mengen}{n}%
}
\begin{tabproof}
  \proofstep{1}{n\in\mathbb{N}}{\rA}
  \proofstep{1}{\Succ(n)\in\mathbb{N}}{\FormulaRefAuto{PeanoSuccClosure}{1}}
  \proofstep{1}{\Succ(n)\neq0}{\FormulaRefAuto{PeanoSuccNonzero}{1}}
  \proofstep{1}{\Pred(\Succ(n))=n}{\FormulaRefAuto{PeanoPred}{2,3}}
\end{tabproof}

\FormulaThmDeltaK[Nichtnullzahlen werden durch ihren Vorgaenger erzeugt]{%
  n\in\mathbb{N}\dsep n\neq0\vdash n=\Succ(\Pred(n))%
}{PeanoPredNonzeroValue}{%
  \DeltaRow{Mengen}{n\dsep y}%
}
\begin{tabproof}
  \proofstep{1}{n\in\mathbb{N}}{\rA}
  \proofstep{2}{n\neq0}{\rA}
  \proofstep{1,2}{\exists y\in\mathbb{N}\,\bigl(n=\Succ(y)\bigr)}{%
    \FormulaRefAuto{PeanoNonzeroIsSucc}{1,2}}
  \proofstep{1,2}{n=\Succ(\Pred(n))}{\FormulaRefAuto{PeanoPred}{1,2,3}}
\end{tabproof}

\FormulaThmDeltaK[Vorgaenger rekonstruiert Nichtnullzahl]{%
  k\in\mathbb{N}\dsep k\neq0\dsep \Pred(k)=n\vdash k=\Succ(n)%
}{PeanoPredRecoverSucc}{%
  \DeltaRow{Mengen}{k\dsep n}%
}
\begin{tabproof}
  \proofstep{1}{k\in\mathbb{N}}{\rA}
  \proofstep{2}{k\neq0}{\rA}
  \proofstep{3}{\Pred(k)=n}{\rA}
  \proofstep{1,2}{k=\Succ(\Pred(k))}{\FormulaRefAuto{PeanoPredNonzeroValue}{1,2}}
  \proofstep{1,2,3}{k=\Succ(n)}{\rIE{3,4}}
\end{tabproof}

\section{Dedekindscher Rekursionssatz}

\paragraph{Beweisskizze}

\begingroup
\footnotesize
\[
\begin{array}{@{}c@{}}
\boxed{
\begin{array}{c}
m_0\in M,\qquad f\colon M\to M
\\[-.1ex]
\text{gesucht: eine eindeutig bestimmte Abbildung }
F\colon\mathbb{N}\to M
\\[-.1ex]
F(0)=m_0,\qquad F(\Succ(n))=f(F(n))
\end{array}}
\\[1.1ex]
\Downarrow
\\[-.2ex]
\boxed{
\begin{array}{c}
\AdmRec{m_0}{f}(H)
\\[-.1ex]
H\subseteq\mathbb{N}\times M,\qquad
(0,m_0)\in H
\\[-.1ex]
(n,a)\in H\Longrightarrow(\Succ(n),f(a))\in H
\end{array}}
\\[1.1ex]
\Downarrow
\\[-.2ex]
\boxed{
\begin{array}{c}
\RecCore{m_0}{f}
=
\displaystyle\bigcap\{H\subseteq\mathbb{N}\times M
\mid \AdmRec{m_0}{f}(H)\}
\\[-.1ex]
\text{kleinste rekursionsadmissible Teilmenge von }
\mathbb{N}\times M
\end{array}}
\\[1.1ex]
\Downarrow
\\[-.2ex]
\boxed{
\begin{array}{c}
\text{Nullstufe: }(0,a)\in\RecCore{m_0}{f}
\Longrightarrow a=m_0
\\[-.1ex]
\text{Nachfolgerstufe: }(n,a)\in\RecCore{m_0}{f}
\Longrightarrow
(\Succ(n),f(a))\in\RecCore{m_0}{f}
\end{array}}
\\[1.1ex]
\Downarrow
\\[-.2ex]
\boxed{
\begin{array}{c}
\forall n\in\mathbb{N}\,\exists!a\in M\,
\bigl((n,a)\in\RecCore{m_0}{f}\bigr)
\\[-.1ex]
\text{Existenz durch Induktion, Eindeutigkeit durch Induktion}
\end{array}}
\\[1.1ex]
\Downarrow
\\[-.2ex]
\boxed{
\begin{array}{c}
\RecFun{m_0}{f}\colon\mathbb{N}\to M,
\qquad
\RecFun{m_0}{f}(n)
=
\iota a\,\bigl((n,a)\in\RecCore{m_0}{f}\bigr)
\\[-.1ex]
\RecFun{m_0}{f}(0)=m_0,\qquad
\RecFun{m_0}{f}(\Succ(n))=f(\RecFun{m_0}{f}(n))
\end{array}}
\\[1.1ex]
\Downarrow
\\[-.2ex]
\boxed{
\begin{array}{c}
G\colon\mathbb{N}\to M,\quad
G(0)=m_0,\quad
G(\Succ(n))=f(G(n))
\\[-.1ex]
\Longrightarrow\quad
G=\RecFun{m_0}{f}
\end{array}}
\end{array}
\]
\endgroup

Der Beweis ersetzt die direkte Definition einer Folge durch eine
Schnittkonstruktion. Zuerst werden diejenigen Teilmengen von
\(\mathbb{N}\times M\) betrachtet, die den Anfangswert enthalten und unter dem
Rekursionsschritt abgeschlossen sind. Ihr Durchschnitt ist der Rekursionskern.
Die Fixierungslemmata zeigen, dass der Kern in jeder Stufe höchstens einen
Wert enthält; die Induktion liefert zugleich die Existenz eines solchen Werts.
Damit ist der Kern der Graph einer Funktion. Die Rekursionsgleichungen folgen
dann aus der Konstruktion des Kerns, und die Eindeutigkeit einer beliebigen
anderen rekursiv definierten Abbildung wird erneut per Induktion bewiesen.

\paragraph{Aufbau des Abschnitts}

Der Abschnitt trennt die Hauptlinie des Beweises von den technischen
Fixierungsargumenten. Zuerst werden rekursionsadmissible Teilmengen und der
Rekursionskern konstruiert. Danach zeigen die Fixierungslemmata, dass jede
Stufe des Kerns höchstens einen Wert enthält; die Induktion liefert die
Existenz eines Werts in jeder Stufe. Aus eindeutiger Existenz wird der Kern als
Funktion gelesen. Die Rekursionsgleichungen und die Eindeutigkeit beliebiger
rekursiver Abbildungen ergeben schließlich den Dedekindschen Rekursionssatz.
Die danach eingeführte Rekursionsabbildung ist nur noch das benannte Objekt,
das durch diesen Satz eindeutig bestimmt ist.

\subsubsection{Hauptkonstruktion: Rekursionsadmissible Mengen und Rekursionskern}

\paragraph{Axiome der Rekursionsadmissibilität}

\FormulaAxiomDelta[Teilmenge des Produktraums]{%
  H\subseteq \mathbb{N}\times M
}{%
  \DeltaRow{Mengen}{H\dsep M}%
}

\FormulaAxiomDelta[Startpaar]{%
  m_0\in M \vdash (0,m_0)\in H
}{%
  \DeltaRow{Mengen}{H\dsep M\dsep m_0}%
}

\FormulaAxiomDelta[Rekursionsschritt]{%
  f\colon M\to M\dsep n\in\mathbb{N}\dsep a\in M\dsep (n,a)\in H
  \vdash (\Succ(n),f(a))\in H
}{%
  \DeltaRow{Mengen}{H\dsep M\dsep n\dsep a}%
  \DeltaRow{Funktionen}{f}%
}

\paragraph{Definition der Rekursionsadmissibilität}

\FormulaDefDeltaR[Begriff der rekursionsadmissiblen Teilmenge]{%
  \AdmRec{m_0}{f}(H)
}{%
  \AdmRec{m_0}{f}(H)
}{%
  \DeltaRow{Mengen}{H\dsep M\dsep n\dsep a\dsep m_0}%
  \DeltaRow{Funktionen}{f}%
  \DeltaRow{\textbf{Axiome}}{}%

  \DeltaRow{\makecell[l]{Teilmenge des\\Produktraums}}
           {H\subseteq \mathbb{N}\times M}
           [\FormulaRefAuto{H\subseteq \mathbb{N}\times M}[ax]]

  \DeltaRow{Startpaar}
           {m_0\in M \vdash (0,m_0)\in H}
           [\FormulaRefAuto{m_0\in M \vdash (0,m_0)\in H}[ax]]

  \DeltaRow{\makecell[l]{Rekursions-\\schritt}}
           {%
             \begin{aligned}[t]
               &f\colon M\to M\dsep n\in\mathbb{N}\dsep{}\\
               &a\in M\dsep (n,a)\in H{}\\
               &\vdash (\Succ(n),f(a))\in H
             \end{aligned}
           }
           [\FormulaRefAuto{f\colon M\to M\dsep n\in\mathbb{N}\dsep a\in M\dsep (n,a)\in H \vdash (\Succ(n),f(a))\in H}[ax]]

  \DeltaRow{\textbf{Neue Symbole}}{}%
  \DeltaPrem{\makecell[l]{Rekursionsadmissible\\Teilmenge}}{ }%
}

\paragraph{Nichtleere Admissibilität und Hilfskonstruktionen}

Der volle Produktraum zeigt zunächst, dass es überhaupt rekursionsadmissible
Teilmengen gibt. Die beiden Fixierungsmengen sind technische Werkzeuge: Sie
werden später benutzt, um falsche Werte in der Nullstufe und in einer
Nachfolgerstufe aus dem Rekursionskern auszuschließen.

\subparagraph{Der volle Produktraum}

\FormulaThmDelta[Der volle Produktraum ist rekursionsadmissibel]{%
  m_0\in M \dsep f\colon M\to M \vdash \AdmRec{m_0}{f}(\mathbb{N}\times M)
}{%
  \DeltaDecl{Mengen}{M\dsep n\dsep a\dsep m_0}%
  \DeltaDecl{Funktionen}{f}%
}
\begin{tabproofsplitwide}
  \proofpartwideind[Teilmenge des Produktraums]{%
    \mathbb{N}\times M \subseteq \mathbb{N}\times M}
    \proofstepwidestar{\mathbb{N}\times M \subseteq \mathbb{N}\times M}{%
      \FormulaRefAuto{A \subseteq A}}
  \closeproofpartwideind

  \proofpartwideindR[Startpaar]{%
    m_0\in M \vdash (0,m_0)\in \mathbb{N}\times M
  }{%
    m_0\in M \vdash (0,m_0)\in \mathbb{N}\times M
  }
    \proofstepwidestar[1]{m_0\in M}{\rA}
    \proofstepwidestar{0\in\mathbb{N}}{%
      \FormulaRefAuto{PeanoZero}}
    \proofstepwidestar[1]{(0,m_0)\in\mathbb{N}\times M}{%
      \FormulaRefAuto{a\in A \dsep b\in B \vdash (a,b)\in A\times B}{2,1}}
  \closeproofpartwideindR

  \proofpartwideindR[Rekursionsschritt]{%
    f\colon M\to M\dsep n\in\mathbb{N}\dsep a\in M\dsep
    (n,a)\in \mathbb{N}\times M \vdash
    (\Succ(n),f(a))\in \mathbb{N}\times M
  }{%
    f\colon M\to M\dsep n\in\mathbb{N}\dsep a\in M\dsep
    (n,a)\in \mathbb{N}\times M \vdash
    (\Succ(n),f(a))\in \mathbb{N}\times M
  }
    \proofstepwidestar[1]{f\colon M\to M}{\rA}
    \proofstepwidestar[2]{n\in\mathbb{N}}{\rA}
    \proofstepwidestar[3]{a\in M}{\rA}
    \proofstepwidestar[4]{(n,a)\in \mathbb{N}\times M}{\rA}

    \proofstepwidestar[2]{\Succ(n)\in\mathbb{N}}{%
      \FormulaRefAuto{PeanoSuccClosure}{2}}

    \proofstepwidestar[1,3]{f(a)\in M}{%
      \FormulaRefAuto{F\colon A\to B\dsep x\in A\vdash F(x)\in B}{1,3}}

    \proofstepwidestar[1,2,3,4]{(\Succ(n),f(a))\in\mathbb{N}\times M}{%
      \FormulaRefAuto{a\in A \dsep b\in B \vdash (a,b)\in A\times B}{5,6}}
  \closeproofpartwideindR

  \proofpartwide{Schluss}
    \proofstepwidestar[1]{m_0\in M}{\rA}
    \proofstepwidestar[2]{f\colon M\to M}{\rA}

    \proofstepwidestar[1,2]{\AdmRec{m_0}{f}(\mathbb{N}\times M)}{%
      \makecell[l]{%
        \FormulaRefAuto{\AdmRec{m_0}{f}(H)}(%
        \\ \quad \FormulaRefAuto{\mathbb{N}\times M \subseteq \mathbb{N}\times M},%
        \\ \quad \FormulaRefAuto{m_0\in M \vdash (0,m_0)\in \mathbb{N}\times M}{1},%
        \\ \quad \FormulaRefAuto{f\colon M\to M\dsep n\in\mathbb{N}\dsep a\in M\dsep (n,a)\in \mathbb{N}\times M \vdash (\Succ(n),f(a))\in \mathbb{N}\times M}{2}%
        \\ )%
      }%
    }
  \closeproofpartwide
\end{tabproofsplitwide}

\subparagraph{Technisches Lemma: Fixierungsmenge der Nullstufe}

\FormulaDefDeltaR[Fixierungsmenge der Nullstufe]{%
  \begin{aligned}[t]
    &m_0\in M \vdash{}\\
    &\ZeroFix{m_0}\coloneq
    \{(n,a)\in \mathbb{N}\times M \mid n=0 \rightarrow a=m_0\}
  \end{aligned}
}{%
  m_0\in M \vdash
  \ZeroFix{m_0}\coloneq
  \{(n,a)\in \mathbb{N}\times M \mid n=0 \rightarrow a=m_0\} }{%
  \DeltaDecl{Mengen}{M\dsep n\dsep a\dsep m_0}%
}

\FormulaThmDelta[Fixierungsmenge der Nullstufe ist rekursionsadmissibel]{%
  m_0\in M\dsep f\colon M\to M \vdash \AdmRec{m_0}{f}(\ZeroFix{m_0}) }{%
  \DeltaDecl{Mengen}{M\dsep n\dsep a\dsep m_0}%
  \DeltaDecl{Funktionen}{f}%
}
\begin{tabproofsplitwide}
  \proofpartwideindR[Teilmenge des Produktraums]{%
    m_0\in M \vdash \ZeroFix{m_0}\subseteq \mathbb{N}\times M
  }{%
    m_0\in M \vdash \ZeroFix{m_0}\subseteq \mathbb{N}\times M
  }
    \proofstepwidestar[1]{m_0\in M}{\rA}
    \proofstepwidestar[]{%
      \{(n,a)\in \mathbb{N}\times M \mid n=0 \rightarrow a=m_0\}
      \subseteq \mathbb{N}\times M
}{%
      \FormulaRefAuto{\{ x \in A \mid P(x) \} \subseteq A}}
    \proofstepwidestar[1]{\ZeroFix{m_0}\subseteq \mathbb{N}\times M}{%
      \FormulaRefAuto{m_0\in M \vdash \ZeroFix{m_0}\coloneq \{(n,a)\in \mathbb{N}\times M \mid n=0 \rightarrow a=m_0\}}{1,2}}
  \closeproofpartwideindR

  \proofpartwideindR[Startpaar]{%
    m_0\in M \vdash (0,m_0)\in \ZeroFix{m_0}
  }{%
    m_0\in M \vdash (0,m_0)\in \ZeroFix{m_0}
  }
    \proofstepwidestar[1]{m_0\in M}{\rA}
    \proofstepwidestar{0\in\mathbb{N}}{%
      \FormulaRefAuto{PeanoZero}}
    \proofstepwidestar[1]{(0,m_0)\in \mathbb{N}\times M}{%
      \FormulaRefAuto{a\in A \dsep b\in B \vdash (a,b)\in A\times B}{2,1}}
    \proofstepwidestar[]{m_0=m_0}{\rII}
    \proofstepwidestar[]{%
      0=0 \rightarrow m_0=m_0
    }{%
      \FormulaRefAuto{Q\vdash P \rightarrow Q}{4}}
    \proofstepwidestar[1]{%
      (0,m_0)\in
      \{(n,a)\in \mathbb{N}\times M \mid n=0 \rightarrow a=m_0\}
    }{%
      \FormulaRefAuto{x \in A\dsep P(x)\vdash x \in \{x \in A \mid P(x)\}}{3,5}}
    \proofstepwidestar[1]{(0,m_0)\in \ZeroFix{m_0}}{%
      \FormulaRefAuto{m_0\in M \vdash \ZeroFix{m_0}\coloneq \{(n,a)\in \mathbb{N}\times M \mid n=0 \rightarrow a=m_0\}}{1,6}}
  \closeproofpartwideindR

  \proofpartwideindR[Rekursionsschritt]{%
    m_0\in M\dsep f\colon M\to M\dsep n\in\mathbb{N}\dsep a\in M\dsep
    (n,a)\in \ZeroFix{m_0} \vdash
    (\Succ(n),f(a))\in \ZeroFix{m_0}
  }{%
    m_0\in M\dsep f\colon M\to M\dsep n\in\mathbb{N}\dsep a\in M\dsep
    (n,a)\in \ZeroFix{m_0} \vdash
    (\Succ(n),f(a))\in \ZeroFix{m_0}
  }
    \proofstepwidestar[1]{m_0\in M}{\rA}
    \proofstepwidestar[2]{f\colon M\to M}{\rA}
    \proofstepwidestar[3]{n\in\mathbb{N}}{\rA}
    \proofstepwidestar[4]{a\in M}{\rA}
    \proofstepwidestar[5]{(n,a)\in \ZeroFix{m_0}}{\rA}

    \proofstepwidestar[3]{\Succ(n)\in\mathbb{N}}{%
      \FormulaRefAuto{PeanoSuccClosure}{3}}

    \proofstepwidestar[2,4]{f(a)\in M}{%
      \FormulaRefAuto{F\colon A\to B\dsep x\in A\vdash F(x)\in B}{2,4}}

    \proofstepwidestar[2,3,4]{(\Succ(n),f(a))\in \mathbb{N}\times M}{%
      \FormulaRefAuto{a\in A \dsep b\in B \vdash (a,b)\in A\times B}{6,7}}

    \proofstepwidestar[3]{\Succ(n)\neq0}{%
      \FormulaRefAuto{PeanoSuccNonzero}{3}}

    \proofstepwidestar[3]{\Succ(n)=0 \rightarrow f(a)=m_0}{%
      \FormulaRefAuto{\neg P \vdash P \rightarrow Q}{9}}

    \proofstepwidestar[1,2,3,4,5]{%
      (\Succ(n),f(a))\in
      \{(n,a)\in \mathbb{N}\times M \mid n=0 \rightarrow a=m_0\}
    }{%
      \FormulaRefAuto{x \in A\dsep P(x)\vdash x \in \{x \in A \mid P(x)\}}{8,10}}

    \proofstepwidestar[1,2,3,4,5]{(\Succ(n),f(a))\in \ZeroFix{m_0}}{%
      \FormulaRefAuto{m_0\in M \vdash \ZeroFix{m_0}\coloneq \{(n,a)\in \mathbb{N}\times M \mid n=0 \rightarrow a=m_0\}}{1,11}}
  \closeproofpartwideindR

  \proofpartwide{Schluss}
    \proofstepwidestar[1]{m_0\in M}{\rA}
    \proofstepwidestar[2]{f\colon M\to M}{\rA}

    \proofstepwidestar[1,2]{\AdmRec{m_0}{f}(\ZeroFix{m_0})}{%
      \makecell[l]{%
        \FormulaRefAuto{\AdmRec{m_0}{f}(H)}(%
        \\ \quad \FormulaRefAuto{m_0\in M \vdash \ZeroFix{m_0}\subseteq \mathbb{N}\times M}{1},%
        \\ \quad \FormulaRefAuto{m_0\in M \vdash (0,m_0)\in \ZeroFix{m_0}}{1},%
        \\ \quad \FormulaRefAuto{m_0\in M\dsep f\colon M\to M\dsep n\in\mathbb{N}\dsep a\in M\dsep (n,a)\in \ZeroFix{m_0} \vdash (\Succ(n),f(a))\in \ZeroFix{m_0}}{1,2}%
        \\ )%
      }%
    }
  \closeproofpartwide
\end{tabproofsplitwide}

\subparagraph{Schnittkonstruktion des Rekursionskerns}

Der Rekursionskern ist die kleinste rekursionsadmissible Teilmenge von
\(\mathbb{N}\times M\): Ein Paar liegt genau dann im Kern, wenn es in jeder
rekursionsadmissiblen Teilmenge liegt.

\FormulaDefDelta[Rekursionskern]{%
  m_0\in M \dsep f\colon M\to M \vdash
  \RecCore{m_0}{f}\coloneq
  \bigcap_{\AdmRec{m_0}{f}(H)} H
}{%
  \DeltaDecl{Mengen}{H\dsep M\dsep m_0}%
  \DeltaDecl{Funktionen}{f}%
}

\FormulaThmDeltaR[Charakterisierung des Rekursionskerns]{%
  \begin{aligned}
    &m_0\in M\dsep f\colon M\to M \vdash{}\\
    &(n,a)\in \RecCore{m_0}{f}\leftrightarrow{}\\
    &\forall H\,\bigl(\AdmRec{m_0}{f}(H)\rightarrow (n,a)\in H\bigr)
  \end{aligned}
}{%
  m_0\in M\dsep f\colon M\to M \vdash
  (n,a)\in \RecCore{m_0}{f}\leftrightarrow
  \forall H\,\bigl(\AdmRec{m_0}{f}(H)\rightarrow (n,a)\in H\bigr)
}{%
  \DeltaDecl{Mengen}{H\dsep M\dsep n\dsep a\dsep m_0}%
  \DeltaDecl{Funktionen}{f}%
}
\begin{tabproofsplitwide}
  \proofpartwideindR[Hinrichtung]{%
    \begin{aligned}[t]
      &m_0\in M\dsep f\colon M\to M \vdash{}\\
      &(n,a)\in \RecCore{m_0}{f}\rightarrow{}\\
      &\forall H\,\bigl(\AdmRec{m_0}{f}(H)\rightarrow (n,a)\in H\bigr)
    \end{aligned}
  }{%
    m_0\in M\dsep f\colon M\to M\dsep (n,a)\in \RecCore{m_0}{f}
    \vdash \forall H\,\bigl(\AdmRec{m_0}{f}(H)\rightarrow (n,a)\in H\bigr)
  }
    \proofstepwidestar[1]{m_0\in M}{\rA}
    \proofstepwidestar[2]{f\colon M\to M}{\rA}
    \proofstepwidestar[3]{(n,a)\in \RecCore{m_0}{f}}{\rA}

    \proofstepwidestar[1,2]{\AdmRec{m_0}{f}(\mathbb{N}\times M)}{%
      \FormulaRefAuto{m_0\in M\dsep f\colon M\to M\vdash \AdmRec{m_0}{f}(\mathbb{N}\times M)}{1,2}}
    \proofstepwidestar[1,2]{\exists H\,\AdmRec{m_0}{f}(H)}{\rEI{4}}

    \proofstepwidestar[1,2,3]{(n,a)\in \bigcap_{\AdmRec{m_0}{f}(H)} H}{%
      \rIE{\FormulaRefAuto{m_0\in M\dsep f\colon M\to M\vdash
        \RecCore{m_0}{f}\coloneq \bigcap_{\AdmRec{m_0}{f}(H)} H}{1,2},3}}

    \proofstepwidestar[1,2,3]{%
      \forall H\,\bigl(\AdmRec{m_0}{f}(H)\rightarrow (n,a)\in H\bigr)
    }{%
      \FormulaRefAuto{%
        \exists A(P(A)) \vdash
        x \in \bigcap_{P(B)} B \leftrightarrow
        \forall C\, (P(C) \rightarrow x \in C)
      }{5,6}}

    \proofstepwidestar[1,2]{%
    \begin{aligned}[t]
      &(n,a)\in \RecCore{m_0}{f}\rightarrow{}\\
      &\forall H\,\bigl(\AdmRec{m_0}{f}(H)\rightarrow (n,a)\in H\bigr)
    \end{aligned}
    }{\rRI{3,7}}
  \closeproofpartwideindR

  \proofpartwideindR[Rückrichtung]{%
    \begin{aligned}[t]
      &m_0\in M\dsep f\colon M\to M \vdash{}\\
      &\forall H\,\bigl(\AdmRec{m_0}{f}(H)\rightarrow (n,a)\in H\bigr)\rightarrow{}\\
      &(n,a)\in \RecCore{m_0}{f}
    \end{aligned}
  }{%
    m_0\in M\dsep f\colon M\to M\dsep
    \forall H\,\bigl(\AdmRec{m_0}{f}(H)\rightarrow (n,a)\in H\bigr)
    \vdash (n,a)\in \RecCore{m_0}{f}
  }
    \proofstepwidestar[1]{m_0\in M}{\rA}
    \proofstepwidestar[2]{f\colon M\to M}{\rA}
    \proofstepwidestar[3]{%
      \forall H\,\bigl(\AdmRec{m_0}{f}(H)\rightarrow (n,a)\in H\bigr)
    }{\rA}

    \proofstepwidestar[1,2]{\AdmRec{m_0}{f}(\mathbb{N}\times M)}{%
      \FormulaRefAuto{m_0\in M\dsep f\colon M\to M\vdash \AdmRec{m_0}{f}(\mathbb{N}\times M)}{1,2}}
    \proofstepwidestar[1,2]{\exists H\,\AdmRec{m_0}{f}(H)}{\rEI{4}}

    \proofstepwidestar[1,2,3]{(n,a)\in \bigcap_{\AdmRec{m_0}{f}(H)} H}{%
      \FormulaRefAuto{%
        \exists A(P(A)) \vdash
        x \in \bigcap_{P(B)} B \leftrightarrow
        \forall C\, (P(C) \rightarrow x \in C)
      }{5,3}}

    \proofstepwidestar[1,2,3]{(n,a)\in \RecCore{m_0}{f}}{%
      \rIE{\FormulaRefAuto{m_0\in M\dsep f\colon M\to M\vdash
        \RecCore{m_0}{f}\coloneq \bigcap_{\AdmRec{m_0}{f}(H)} H}{1,2},6}}

    \proofstepwidestar[1,2]{%
    \begin{aligned}[t]
      &\forall H\,\bigl(\AdmRec{m_0}{f}(H)\rightarrow (n,a)\in H\bigr)\rightarrow{}\\
      &(n,a)\in \RecCore{m_0}{f}
    \end{aligned}
    }{\rRI{3,7}}
  \closeproofpartwideindR

\proofpartwide{Schluss}
  \proofstepwidestar[1]{m_0\in M}{\rA}
  \proofstepwidestar[2]{f\colon M\to M}{\rA}

  \proofstepwidestar[1,2]{%
    \begin{aligned}[t]
      &(n,a)\in \RecCore{m_0}{f}\rightarrow{}\\
      &\forall H\,\bigl(\AdmRec{m_0}{f}(H)\rightarrow (n,a)\in H\bigr)
    \end{aligned}
  }{%
    \FormulaRefAuto{%
      m_0\in M\dsep f\colon M\to M\dsep (n,a)\in \RecCore{m_0}{f}
      \vdash \forall H\,\bigl(\AdmRec{m_0}{f}(H)\rightarrow (n,a)\in H\bigr)
    }{1,2}}

  \proofstepwidestar[1,2]{%
    \begin{aligned}[t]
      &\forall H\,\bigl(\AdmRec{m_0}{f}(H)\rightarrow (n,a)\in H\bigr)\rightarrow{}\\
      &(n,a)\in \RecCore{m_0}{f}
    \end{aligned}
  }{%
    \FormulaRefAuto{%
      m_0\in M\dsep f\colon M\to M\dsep
      \forall H\,\bigl(\AdmRec{m_0}{f}(H)\rightarrow (n,a)\in H\bigr)
      \vdash (n,a)\in \RecCore{m_0}{f}
    }{1,2}}

  \proofstepwidestar[1,2]{%
    \begin{aligned}[t]
      &(n,a)\in \RecCore{m_0}{f}\leftrightarrow{}\\
      &\forall H\,\bigl(\AdmRec{m_0}{f}(H)\rightarrow (n,a)\in H\bigr)
    \end{aligned}
  }{\rLRI{3,4}}
\closeproofpartwide
\end{tabproofsplitwide}

\FormulaThmDelta[Minimalität des Rekursionskerns]{%
  m_0\in M\dsep f\colon M\to M\dsep \AdmRec{m_0}{f}(H)
  \vdash \RecCore{m_0}{f}\subseteq H
}{%
  \DeltaDecl{Mengen}{H\dsep M\dsep m_0}%
  \DeltaDecl{Funktionen}{f}%
}
\begin{tabproof}
  \proofstep{1}{m_0\in M}{\rA}
  \proofstep{2}{f\colon M\to M}{\rA}
  \proofstep{3}{\AdmRec{m_0}{f}(H)}{\rA}
  \proofstep{3}{\bigcap_{\AdmRec{m_0}{f}(X)} X \subseteq H}{%
    \FormulaRefAuto{P(C)\vdash \bigcap_{P(A)} A \subseteq C}{3}}
  \proofstep{1,2,3}{\RecCore{m_0}{f}\subseteq H}{%
    \rIE{\FormulaRefAuto{%
      m_0\in M \dsep f\colon M\to M \vdash
      \RecCore{m_0}{f}\coloneq
      \bigcap_{\AdmRec{m_0}{f}(X)} X
    }{1,2},4}}
\end{tabproof}

\FormulaThmDelta[Der Rekursionskern liegt im Produktraum]{%
  m_0\in M\dsep f\colon M\to M \vdash \RecCore{m_0}{f}\subseteq \mathbb{N}\times M
}{%
  \DeltaDecl{Mengen}{M\dsep m_0}%
  \DeltaDecl{Funktionen}{f}%
}
\begin{tabproof}
  \proofstep{1}{m_0\in M}{\rA}
  \proofstep{2}{f\colon M\to M}{\rA}
  \proofstep{1,2}{\AdmRec{m_0}{f}(\mathbb{N}\times M)}{%
    \FormulaRefAuto{m_0\in M \dsep f\colon M\to M \vdash \AdmRec{m_0}{f}(\mathbb{N}\times M)}{1,2}}
  \proofstep{1,2}{\RecCore{m_0}{f}\subseteq \mathbb{N}\times M}{%
    \FormulaRefAuto{m_0\in M\dsep f\colon M\to M\dsep \AdmRec{m_0}{f}(H)
    \vdash \RecCore{m_0}{f}\subseteq H}{1,2,3}}
\end{tabproof}

\FormulaThmDelta[Das Startpaar liegt im Rekursionskern]{%
  m_0\in M\dsep f\colon M\to M \vdash (0,m_0)\in \RecCore{m_0}{f}
}{%
  \DeltaDecl{Mengen}{H\dsep M\dsep m_0}%
  \DeltaDecl{Funktionen}{f}%
}
\begin{tabproof}
  \proofstep{1}{m_0\in M}{\rA}
  \proofstep{2}{f\colon M\to M}{\rA}

  \proofstep{3}{\AdmRec{m_0}{f}(H)}{\rA}
  \proofstep{1,3}{(0,m_0)\in H}{%
    \FormulaRefAuto{\AdmRec{m_0}{f}(H)}{3}}

  \proofstep{}{%
    \forall H\,\bigl(\AdmRec{m_0}{f}(H)\rightarrow (0,m_0)\in H\bigr)
  }{\rUI{\rRI{3,4}}}

  \proofstep{1,2}{(0,m_0)\in \RecCore{m_0}{f}}{%
    \FormulaRefAuto{%
      m_0\in M\dsep f\colon M\to M \vdash
      (n,a)\in \RecCore{m_0}{f}\leftrightarrow
      \forall H\,\bigl(\AdmRec{m_0}{f}(H)\rightarrow (n,a)\in H\bigr)
    }{1,2,5}}
\end{tabproof}

\FormulaThmDeltaR[Abschluss des Rekursionskerns]{%
  \begin{aligned}
    &m_0\in M\dsep f\colon M\to M\dsep n\in\mathbb{N}\dsep{}\\
    &a\in M\dsep (n,a)\in \RecCore{m_0}{f}
      \vdash (\Succ(n),f(a))\in \RecCore{m_0}{f}
  \end{aligned}
}{%
  m_0\in M\dsep f\colon M\to M\dsep n\in\mathbb{N}\dsep
  a\in M\dsep (n,a)\in \RecCore{m_0}{f}
  \vdash (\Succ(n),f(a))\in \RecCore{m_0}{f}
}{%
  \DeltaDecl{Mengen}{H\dsep M\dsep n\dsep a\dsep m_0}%
  \DeltaDecl{Funktionen}{f}%
}
\begin{tabproof}
  \proofstep{1}{m_0\in M}{\rA}
  \proofstep{2}{f\colon M\to M}{\rA}
  \proofstep{3}{n\in\mathbb{N}}{\rA}
  \proofstep{4}{a\in M}{\rA}
  \proofstep{5}{(n,a)\in \RecCore{m_0}{f}}{\rA}

  \proofstep{1,2}{%
    \begin{aligned}
      &(n,a)\in \RecCore{m_0}{f}\leftrightarrow{}\\
      &\forall H\,\bigl(\AdmRec{m_0}{f}(H)\rightarrow (n,a)\in H\bigr)
    \end{aligned}
  }{%
    \FormulaRefAuto{%
      m_0\in M\dsep f\colon M\to M \vdash
      (n,a)\in \RecCore{m_0}{f}\leftrightarrow
      \forall H\,\bigl(\AdmRec{m_0}{f}(H)\rightarrow (n,a)\in H\bigr)
    }{1,2}}

  \proofstep{1,2,5}{%
    \begin{aligned}
      &\forall H\,\bigl(\AdmRec{m_0}{f}(H)\rightarrow{}\\
      &\qquad (n,a)\in H\bigr)
    \end{aligned}
  }{%
    \FormulaRefAuto{P \leftrightarrow Q, P \vdash Q}{6,5}}

  \proofstep{8}{\AdmRec{m_0}{f}(H)}{\rA}
  \proofstep{1,2,5,8}{(n,a)\in H}{\rRE{\rUE{7},8}}

  \proofstep{2,3,4,8,9}{(\Succ(n),f(a))\in H}{%
    \FormulaRefAuto{f\colon M\to M\dsep n\in\mathbb{N}\dsep a\in M\dsep (n,a)\in H \vdash (\Succ(n),f(a))\in H}[ax]{2,3,4,9}}

  \proofstep{1,2,3,4,5}{%
    \begin{aligned}
      &\forall H\,\bigl(\AdmRec{m_0}{f}(H)\rightarrow{}\\
      &\qquad (\Succ(n),f(a))\in H\bigr)
    \end{aligned}
  }{\rUI{\rRI{8,10}}}

  \proofstep{1,2}{%
    \begin{aligned}
      &(\Succ(n),f(a))\in \RecCore{m_0}{f}\leftrightarrow{}\\
      &\forall H\,\bigl(\AdmRec{m_0}{f}(H)\rightarrow (\Succ(n),f(a))\in H\bigr)
    \end{aligned}
  }{%
    \FormulaRefAuto{%
      m_0\in M\dsep f\colon M\to M \vdash
      (n,a)\in \RecCore{m_0}{f}\leftrightarrow
      \forall H\,\bigl(\AdmRec{m_0}{f}(H)\rightarrow (n,a)\in H\bigr)
    }{1,2}}

  \proofstep{1,2,3,4,5}{(\Succ(n),f(a))\in \RecCore{m_0}{f}}{%
    \FormulaRefAuto{P \leftrightarrow Q, Q \vdash P}{12,11}}
\end{tabproof}

\FormulaThmDelta[Der Rekursionskern ist rekursionsadmissibel]{%
  m_0\in M\dsep f\colon M\to M \vdash \AdmRec{m_0}{f}(\RecCore{m_0}{f})
}{%
  \DeltaDecl{Mengen}{H\dsep M\dsep n\dsep a\dsep m_0}%
  \DeltaDecl{Funktionen}{f}%
}
\begin{tabproofwide}
  \proofstepwidestar[1]{m_0\in M}{\rA}
  \proofstepwidestar[2]{f\colon M\to M}{\rA}

  \proofstepwidestar[1,2]{\AdmRec{m_0}{f}(\RecCore{m_0}{f})}{%
    \makecell[l]{%
      \FormulaRefAuto{\AdmRec{m_0}{f}(H)}(%
      \\ \quad \FormulaRefAuto{m_0\in M\dsep f\colon M\to M \vdash \RecCore{m_0}{f}\subseteq \mathbb{N}\times M}{1,2},%
      \\ \quad \FormulaRefAuto{m_0\in M\dsep f\colon M\to M \vdash (0,m_0)\in \RecCore{m_0}{f}}{1,2},%
      \\ \quad \FormulaRefAuto{%
        m_0\in M\dsep f\colon M\to M\dsep n\in\mathbb{N}\dsep
        a\in M\dsep (n,a)\in \RecCore{m_0}{f}
        \vdash (\Succ(n),f(a))\in \RecCore{m_0}{f}
      }{1,2}%
      \\ )%
    }%
  }
\end{tabproofwide}

\FormulaThmDeltaK[Rekursionskern ist kleinster rekursionsadmissibler Graphkandidat]{%
  \begin{aligned}[t]
    &m_0\in M\dsep f\colon M\to M \vdash{}\\
    &\RecCore{m_0}{f}\subseteq \mathbb{N}\times M \land{}\\
    &(0,m_0)\in \RecCore{m_0}{f}\land{}\\
    &\AdmRec{m_0}{f}(\RecCore{m_0}{f})\land{}\\
    &\forall H\,\bigl(\AdmRec{m_0}{f}(H)\rightarrow
      \RecCore{m_0}{f}\subseteq H\bigr)
  \end{aligned}
}{RecCoreLeastAdmissible}{%
  \DeltaDecl{Mengen}{H\dsep M\dsep m_0}%
  \DeltaDecl{Funktionen}{f}%
}
\begin{tabproofwide}
  \proofstepwidestar[1]{m_0\in M}{\rA}
  \proofstepwidestar[2]{f\colon M\to M}{\rA}

  \proofstepwidestar[1,2]{\RecCore{m_0}{f}\subseteq \mathbb{N}\times M}{%
    \FormulaRefAuto{m_0\in M\dsep f\colon M\to M \vdash \RecCore{m_0}{f}\subseteq \mathbb{N}\times M}{1,2}}
  \proofstepwidestar[1,2]{(0,m_0)\in \RecCore{m_0}{f}}{%
    \FormulaRefAuto{m_0\in M\dsep f\colon M\to M \vdash (0,m_0)\in \RecCore{m_0}{f}}{1,2}}
  \proofstepwidestar[1,2]{\AdmRec{m_0}{f}(\RecCore{m_0}{f})}{%
    \FormulaRefAuto{m_0\in M\dsep f\colon M\to M \vdash \AdmRec{m_0}{f}(\RecCore{m_0}{f})}{1,2}}

  \proofstepwidestar[5]{\AdmRec{m_0}{f}(H)}{\rA}
  \proofstepwidestar[1,2,5]{\RecCore{m_0}{f}\subseteq H}{%
    \FormulaRefAuto{m_0\in M\dsep f\colon M\to M\dsep \AdmRec{m_0}{f}(H)
    \vdash \RecCore{m_0}{f}\subseteq H}{1,2,5}}
  \proofstepwidestar[1,2]{%
    \forall H\,\bigl(\AdmRec{m_0}{f}(H)\rightarrow
    \RecCore{m_0}{f}\subseteq H\bigr)
  }{\rUI{\rRI{5,6}}}

  \proofstepwidestar[1,2]{%
    \begin{aligned}[t]
      &\RecCore{m_0}{f}\subseteq \mathbb{N}\times M\land{}\\
      &(0,m_0)\in \RecCore{m_0}{f}
    \end{aligned}
  }{\rAI{3,4}}
  \proofstepwidestar[1,2]{%
    \begin{aligned}[t]
      &\RecCore{m_0}{f}\subseteq \mathbb{N}\times M\land{}\\
      &(0,m_0)\in \RecCore{m_0}{f}\land{}\\
      &\AdmRec{m_0}{f}(\RecCore{m_0}{f})
    \end{aligned}
  }{\rAI{8,5}}
  \proofstepwidestar[1,2]{%
    \begin{aligned}[t]
      &\RecCore{m_0}{f}\subseteq \mathbb{N}\times M\land{}\\
      &(0,m_0)\in \RecCore{m_0}{f}\land{}\\
      &\AdmRec{m_0}{f}(\RecCore{m_0}{f})\land{}\\
      &\forall H\,\bigl(\AdmRec{m_0}{f}(H)\rightarrow
        \RecCore{m_0}{f}\subseteq H\bigr)
    \end{aligned}
  }{\rAI{9,7}}
\end{tabproofwide}

\subparagraph{Fixierung der Nullstufe des Kerns}

\FormulaThmDelta[Die Nullstufe des Rekursionskerns ist festgelegt]{%
  m_0\in M\dsep f\colon M\to M\dsep
  (0,a)\in \RecCore{m_0}{f}\vdash a=m_0
}{%
  \DeltaDecl{Mengen}{M\dsep a\dsep m_0}%
  \DeltaDecl{Funktionen}{f}%
}
\begin{tabproofwide}
  \proofstepwidestar[1]{m_0\in M}{\rA}
  \proofstepwidestar[2]{f\colon M\to M}{\rA}
  \proofstepwidestar[3]{(0,a)\in \RecCore{m_0}{f}}{\rA}

  \proofstepwidestar[1,2]{\AdmRec{m_0}{f}(\ZeroFix{m_0})}{%
    \FormulaRefAuto{m_0\in M\dsep f\colon M\to M \vdash \AdmRec{m_0}{f}(\ZeroFix{m_0})}{1,2}}

  \proofstepwidestar[1,2]{\RecCore{m_0}{f}\subseteq \ZeroFix{m_0}}{%
    \FormulaRefAuto{m_0\in M\dsep f\colon M\to M\dsep \AdmRec{m_0}{f}(H)\vdash \RecCore{m_0}{f}\subseteq H}{1,2,4}}

  \proofstepwidestar[1,2,3]{(0,a)\in \ZeroFix{m_0}}{%
    \FormulaRefAuto{A\subseteq B\dsep x\in A\vdash x\in B}{5,3}}

  \proofstepwidestar[1]{%
    \begin{aligned}[t]
      &\ZeroFix{m_0}\coloneq{}\\
      &\{(n,a)\in \mathbb{N}\times M \mid{}\\
      &\qquad n=0 \rightarrow a=m_0\}
    \end{aligned}
  }{%
    \FormulaRefAuto{m_0\in M \vdash \ZeroFix{m_0}\coloneq \{(n,a)\in \mathbb{N}\times M \mid n=0 \rightarrow a=m_0\}}{1}}

  \proofstepwidestar[1,3]{%
    (0,a)\in
    \{(n,a)\in \mathbb{N}\times M \mid n=0 \rightarrow a=m_0\}
  }{%
    \rIE{7,6}}

  \proofstepwidestar[1,3]{0=0 \rightarrow a=m_0}{%
    \FormulaRefAuto{x \in \{x \in A \mid P(x)\}\vdash P(x)}{8}}

  \proofstepwidestar{0=0}{\rII}
  \proofstepwidestar[1,3]{a=m_0}{\rRE{10,9}}
\end{tabproofwide}

\subparagraph{Technisches Lemma: Fixierungsmenge einer Nachfolgerstufe}

\FormulaDefDeltaR[Fixierungsmenge der Nachfolgerstufe]{%
  \begin{aligned}[t]
    &m_0\in M\dsep f\colon M\to M\dsep x\in\mathbb{N}\dsep u\in M \vdash{}\\
    &\Phi_{x,u}\coloneq
    \{(y,d)\in \RecCore{m_0}{f}\mid y=\Succ(x)\rightarrow d=f(u)\}
  \end{aligned}
}{%
  m_0\in M\dsep f\colon M\to M\dsep x\in\mathbb{N}\dsep u\in M \vdash
  \Phi_{x,u}\coloneq
  \{(y,d)\in \RecCore{m_0}{f}\mid y=\Succ(x)\rightarrow d=f(u)\}
}{%
  \DeltaDecl{Mengen}{M\dsep x\dsep y\dsep d\dsep u\dsep m_0}%
  \DeltaDecl{Funktionen}{f}%
}

\FormulaThmDeltaR[Fixierungsmenge der Nachfolgerstufe ist rekursionsadmissibel]{%
  \begin{aligned}[t]
    &m_0\in M\dsep f\colon M\to M\dsep x\in\mathbb{N}\dsep u\in M\dsep{}\\
    &(x,u)\in \RecCore{m_0}{f}\dsep{}\\
    &\forall c\in M\,\bigl((x,c)\in \RecCore{m_0}{f}\rightarrow c=u\bigr)\vdash{}\\
    &\AdmRec{m_0}{f}(\Phi_{x,u})
  \end{aligned}
}{%
  m_0\in M\dsep f\colon M\to M\dsep x\in\mathbb{N}\dsep u\in M\dsep
  (x,u)\in \RecCore{m_0}{f}\dsep
  \forall c\in M\,\bigl((x,c)\in \RecCore{m_0}{f}\rightarrow c=u\bigr)
  \vdash \AdmRec{m_0}{f}(\Phi_{x,u})
}{%
  \DeltaDecl{Mengen}{M\dsep x\dsep y\dsep d\dsep c\dsep u\dsep m_0}%
  \DeltaDecl{Funktionen}{f}%
}
\begin{tabproofsplitwide}
  \proofpartwideindR[Elimination aus \(\Phi_{x,u}\): Kernzugehörigkeit]{%
    \begin{aligned}[t]
      &m_0\in M\dsep f\colon M\to M\dsep x\in\mathbb{N}\dsep u\in M\dsep{}\\
      &(y,d)\in \Phi_{x,u}\vdash (y,d)\in \RecCore{m_0}{f}
    \end{aligned}
  }{%
    m_0\in M\dsep f\colon M\to M\dsep x\in\mathbb{N}\dsep u\in M\dsep
    (y,d)\in \Phi_{x,u}\vdash (y,d)\in \RecCore{m_0}{f}
  }
    \proofstepwidestar[1]{m_0\in M}{\rA}
    \proofstepwidestar[2]{f\colon M\to M}{\rA}
    \proofstepwidestar[3]{x\in\mathbb{N}}{\rA}
    \proofstepwidestar[4]{u\in M}{\rA}
    \proofstepwidestar[5]{(y,d)\in \Phi_{x,u}}{\rA}
    \proofstepwidestar[1,2,3,4]{%
      \begin{aligned}[t]
        &\Phi_{x,u}\coloneq{}\\
        &\{(y,d)\in \RecCore{m_0}{f}\mid{}\\
        &y=\Succ(x)\rightarrow d=f(u)\}
      \end{aligned}
    }{%
      \FormulaRefAuto{m_0\in M\dsep f\colon M\to M\dsep x\in\mathbb{N}\dsep u\in M \vdash \Phi_{x,u}\coloneq \{(y,d)\in \RecCore{m_0}{f}\mid y=\Succ(x)\rightarrow d=f(u)\}}{1,2,3,4}}
    \proofstepwidestar[1,2,3,4,5]{%
      \begin{aligned}[t]
        &(y,d)\in{}\\
        &\{(y,d)\in \RecCore{m_0}{f}\mid{}\\
        &y=\Succ(x)\rightarrow d=f(u)\}
      \end{aligned}
    }{\rIE{6,5}}
    \proofstepwidestar[1,2,3,4,5]{(y,d)\in \RecCore{m_0}{f}}{%
      \FormulaRefAuto{x \in \{x \in A \mid P(x)\}\vdash x\in A}{7}}
  \closeproofpartwideindR

  \proofpartwideindR[Einführung in \(\Phi_{x,u}\)]{%
    \begin{aligned}[t]
      &m_0\in M\dsep f\colon M\to M\dsep x\in\mathbb{N}\dsep u\in M\dsep{}\\
      &(z,e)\in \RecCore{m_0}{f}\dsep
      z=\Succ(x)\rightarrow e=f(u)\vdash{}\\
      &(z,e)\in \Phi_{x,u}
    \end{aligned}
  }{%
    m_0\in M\dsep f\colon M\to M\dsep x\in\mathbb{N}\dsep u\in M\dsep
    (z,e)\in \RecCore{m_0}{f}\dsep
    z=\Succ(x)\rightarrow e=f(u)\vdash (z,e)\in \Phi_{x,u}
  }
    \proofstepwidestar[1]{m_0\in M}{\rA}
    \proofstepwidestar[2]{f\colon M\to M}{\rA}
    \proofstepwidestar[3]{x\in\mathbb{N}}{\rA}
    \proofstepwidestar[4]{u\in M}{\rA}
    \proofstepwidestar[5]{(z,e)\in \RecCore{m_0}{f}}{\rA}
    \proofstepwidestar[6]{z=\Succ(x)\rightarrow e=f(u)}{\rA}
    \proofstepwidestar[5,6]{%
      \begin{aligned}[t]
        &(z,e)\in{}\\
        &\{(y,d)\in \RecCore{m_0}{f}\mid{}\\
        &y=\Succ(x)\rightarrow d=f(u)\}
      \end{aligned}
    }{%
      \FormulaRefAuto{x \in A\dsep P(x)\vdash x \in \{x \in A \mid P(x)\}}{5,6}}
    \proofstepwidestar[1,2,3,4]{%
      \begin{aligned}[t]
        &\Phi_{x,u}\coloneq{}\\
        &\{(y,d)\in \RecCore{m_0}{f}\mid{}\\
        &y=\Succ(x)\rightarrow d=f(u)\}
      \end{aligned}
    }{%
      \FormulaRefAuto{m_0\in M\dsep f\colon M\to M\dsep x\in\mathbb{N}\dsep u\in M \vdash \Phi_{x,u}\coloneq \{(y,d)\in \RecCore{m_0}{f}\mid y=\Succ(x)\rightarrow d=f(u)\}}{1,2,3,4}}
    \proofstepwidestar[1,2,3,4,5,6]{(z,e)\in \Phi_{x,u}}{%
      \rIE{8,7}}
  \closeproofpartwideindR

  \proofpartwideindR[Teilmenge des Produktraums]{%
    \begin{aligned}[t]
      &m_0\in M\dsep f\colon M\to M\dsep x\in\mathbb{N}\dsep u\in M \vdash{}\\
      &\Phi_{x,u}\subseteq \mathbb{N}\times M
    \end{aligned}
  }{%
    m_0\in M\dsep f\colon M\to M\dsep x\in\mathbb{N}\dsep u\in M \vdash
    \Phi_{x,u}\subseteq \mathbb{N}\times M
  }
    \proofstepwidestar[1]{m_0\in M}{\rA}
    \proofstepwidestar[2]{f\colon M\to M}{\rA}
    \proofstepwidestar[3]{x\in\mathbb{N}}{\rA}
    \proofstepwidestar[4]{u\in M}{\rA}
    \proofstepwidestar[1,2,3,4]{%
      \begin{aligned}[t]
        &\Phi_{x,u}\coloneq{}\\
        &\{(y,d)\in \RecCore{m_0}{f}\mid{}\\
        &y=\Succ(x)\rightarrow d=f(u)\}
      \end{aligned}
    }{%
      \FormulaRefAuto{m_0\in M\dsep f\colon M\to M\dsep x\in\mathbb{N}\dsep u\in M \vdash \Phi_{x,u}\coloneq \{(y,d)\in \RecCore{m_0}{f}\mid y=\Succ(x)\rightarrow d=f(u)\}}{1,2,3,4}}
    \proofstepwidestar{%
      \begin{aligned}[t]
        &\{(y,d)\in \RecCore{m_0}{f}\mid{}\\
        &y=\Succ(x)\rightarrow d=f(u)\}\\
        &\subseteq \RecCore{m_0}{f}
      \end{aligned}
    }{%
      \FormulaRefAuto{\{ x \in A \mid P(x) \} \subseteq A}}
    \proofstepwidestar[1,2,3,4]{%
      \Phi_{x,u}\subseteq \RecCore{m_0}{f}
    }{\rIE{5,6}}
    \proofstepwidestar[1,2]{\RecCore{m_0}{f}\subseteq \mathbb{N}\times M}{%
      \FormulaRefAuto{m_0\in M\dsep f\colon M\to M \vdash \RecCore{m_0}{f}\subseteq \mathbb{N}\times M}{1,2}}
    \proofstepwidestar[1,2,3,4]{\Phi_{x,u}\subseteq \mathbb{N}\times M}{%
      \FormulaRefAuto{A \subseteq B, B \subseteq C \vdash A \subseteq C}{7,8}}
  \closeproofpartwideindR

  \proofpartwideindR[Startpaar]{%
    \begin{aligned}[t]
      &m_0\in M\dsep f\colon M\to M\dsep x\in\mathbb{N}\dsep u\in M \vdash{}\\
      &(0,m_0)\in \Phi_{x,u}
    \end{aligned}
  }{%
    m_0\in M\dsep f\colon M\to M\dsep x\in\mathbb{N}\dsep u\in M \vdash
    (0,m_0)\in \Phi_{x,u}
  }
    \proofstepwidestar[1]{m_0\in M}{\rA}
    \proofstepwidestar[2]{f\colon M\to M}{\rA}
    \proofstepwidestar[3]{x\in\mathbb{N}}{\rA}
    \proofstepwidestar[4]{u\in M}{\rA}
    \proofstepwidestar[3]{\Succ(x)\neq0}{%
      \FormulaRefAuto{PeanoSuccNonzero}{3}}
    \proofstepwidestar[6]{0=\Succ(x)}{\rA}
    \proofstepwidestar[6]{\Succ(x)=0}{%
      \FormulaRefAuto{a=b\vdash b=a}{6}}
    \proofstepwidestar[3]{0=\Succ(x)\rightarrow \Succ(x)=0}{%
      \rRI{6,7}}
    \proofstepwidestar[3]{\neg(0=\Succ(x))}{%
      \FormulaRefAuto{P \rightarrow Q, \neg Q \vdash \neg P}{8,5}}
    \proofstepwidestar[3]{0=\Succ(x)\rightarrow m_0=f(u)}{%
      \FormulaRefAuto{\neg P \vdash P \rightarrow Q}{9}}
    \proofstepwidestar[1,2]{(0,m_0)\in \RecCore{m_0}{f}}{%
      \FormulaRefAuto{m_0\in M\dsep f\colon M\to M \vdash (0,m_0)\in \RecCore{m_0}{f}}{1,2}}
    \proofstepwidestar[1,2,3,4]{%
      (0,m_0)\in \Phi_{x,u}
    }{%
      \FormulaRefAuto{m_0\in M\dsep f\colon M\to M\dsep x\in\mathbb{N}\dsep u\in M\dsep (z,e)\in \RecCore{m_0}{f}\dsep z=\Succ(x)\rightarrow e=f(u)\vdash (z,e)\in \Phi_{x,u}}{1,2,3,4,11,10}}
  \closeproofpartwideindR

  \proofpartwideindR[Nachfolgerbarriere]{%
    \begin{aligned}[t]
      &m_0\in M\dsep f\colon M\to M\dsep x\in\mathbb{N}\dsep u\in M\dsep{}\\
      &\forall c\in M\,\bigl((x,c)\in \RecCore{m_0}{f}\rightarrow c=u\bigr)\dsep{}\\
      &y\in\mathbb{N}\dsep d\in M\dsep (y,d)\in \Phi_{x,u}\vdash{}\\
      &\Succ(y)=\Succ(x)\rightarrow f(d)=f(u)
    \end{aligned}
  }{%
    m_0\in M\dsep f\colon M\to M\dsep x\in\mathbb{N}\dsep u\in M\dsep
    \forall c\in M\,\bigl((x,c)\in \RecCore{m_0}{f}\rightarrow c=u\bigr)\dsep
    y\in\mathbb{N}\dsep d\in M\dsep (y,d)\in \Phi_{x,u}
    \vdash \Succ(y)=\Succ(x)\rightarrow f(d)=f(u)
  }
    \proofstepwidestar[1]{m_0\in M}{\rA}
    \proofstepwidestar[2]{f\colon M\to M}{\rA}
    \proofstepwidestar[3]{x\in\mathbb{N}}{\rA}
    \proofstepwidestar[4]{u\in M}{\rA}
    \proofstepwidestar[5]{%
      \begin{aligned}[t]
        &\forall c\in M\,\bigl(
        (x,c)\in \RecCore{m_0}{f}\\
        &\rightarrow c=u
        \bigr)
      \end{aligned}
    }{\rA}
    \proofstepwidestar[6]{y\in\mathbb{N}}{\rA}
    \proofstepwidestar[7]{d\in M}{\rA}
    \proofstepwidestar[8]{(y,d)\in \Phi_{x,u}}{\rA}
    \proofstepwidestar[1,2,3,4,8]{(y,d)\in \RecCore{m_0}{f}}{%
      \FormulaRefAuto{m_0\in M\dsep f\colon M\to M\dsep x\in\mathbb{N}\dsep u\in M\dsep (y,d)\in \Phi_{x,u}\vdash (y,d)\in \RecCore{m_0}{f}}{1,2,3,4,8}}
    \proofstepwidestar[10]{\Succ(y)=\Succ(x)}{\rA}
    \proofstepwidestar[10]{\Succ(x)=\Succ(y)}{%
      \FormulaRefAuto{a=b\vdash b=a}{10}}
    \proofstepwidestar[3,6,11]{x=y}{%
      \FormulaRefAuto{PeanoSuccInjective}{3,6,11}}
    \proofstepwidestar[3,6,7,10]{(x,d)=(y,d)}{%
      \FormulaRefAuto{a=c \dsep b=d\vdash (a,b)=(c,d)}{12,\rII}}
    \proofstepwidestar[1,2,3,4,6,7,8,10]{(x,d)\in \RecCore{m_0}{f}}{%
      \FormulaRefAuto{a=b \dsep b\in A \vdash a\in A}{13,9}}
    \proofstepwidestar[5,7]{(x,d)\in \RecCore{m_0}{f}\rightarrow d=u}{%
      \rRE{\rUE{5},7}}
    \proofstepwidestar[1,2,3,4,5,6,7,8,10]{d=u}{\rRE{15,14}}
    \proofstepwidestar[1,2,3,4,5,6,7,8,10]{f(d)=f(u)}{%
      \FormulaRefAuto{x\in A\dsep y\in A\dsep x=y\vdash F(x)=F(y)}{7,4,16}}
    \proofstepwidestar[1,2,3,4,5,6,7,8]{\Succ(y)=\Succ(x)\rightarrow f(d)=f(u)}{%
      \rRI{10,17}}
  \closeproofpartwideindR

  \proofpartwideindR[Rekursionsschritt]{%
    \begin{aligned}[t]
      &m_0\in M\dsep f\colon M\to M\dsep x\in\mathbb{N}\dsep u\in M\dsep{}\\
      &\forall c\in M\,\bigl((x,c)\in \RecCore{m_0}{f}\rightarrow c=u\bigr)\dsep{}\\
      &y\in\mathbb{N}\dsep d\in M\dsep (y,d)\in \Phi_{x,u}\vdash{}\\
      &(\Succ(y),f(d))\in \Phi_{x,u}
    \end{aligned}
  }{%
    m_0\in M\dsep f\colon M\to M\dsep x\in\mathbb{N}\dsep u\in M\dsep
    \forall c\in M\,\bigl((x,c)\in \RecCore{m_0}{f}\rightarrow c=u\bigr)\dsep
    y\in\mathbb{N}\dsep d\in M\dsep (y,d)\in \Phi_{x,u}
    \vdash (\Succ(y),f(d))\in \Phi_{x,u}
  }
    \proofstepwidestar[1]{m_0\in M}{\rA}
    \proofstepwidestar[2]{f\colon M\to M}{\rA}
    \proofstepwidestar[3]{x\in\mathbb{N}}{\rA}
    \proofstepwidestar[4]{u\in M}{\rA}
    \proofstepwidestar[5]{%
      \begin{aligned}[t]
        &\forall c\in M\,\bigl(
        (x,c)\in \RecCore{m_0}{f}\\
        &\rightarrow c=u
        \bigr)
      \end{aligned}
    }{\rA}
    \proofstepwidestar[6]{y\in\mathbb{N}}{\rA}
    \proofstepwidestar[7]{d\in M}{\rA}
    \proofstepwidestar[8]{(y,d)\in \Phi_{x,u}}{\rA}
    \proofstepwidestar[1,2,3,4,8]{(y,d)\in \RecCore{m_0}{f}}{%
      \FormulaRefAuto{m_0\in M\dsep f\colon M\to M\dsep x\in\mathbb{N}\dsep u\in M\dsep (y,d)\in \Phi_{x,u}\vdash (y,d)\in \RecCore{m_0}{f}}{1,2,3,4,8}}
    \proofstepwidestar[1,2,6,7,9]{(\Succ(y),f(d))\in \RecCore{m_0}{f}}{%
      \FormulaRefAuto{m_0\in M\dsep f\colon M\to M\dsep n\in\mathbb{N}\dsep a\in M\dsep (n,a)\in \RecCore{m_0}{f}\vdash (\Succ(n),f(a))\in \RecCore{m_0}{f}}{1,2,6,7,9}}
    \proofstepwidestar[1,2,3,4,5,6,7,8]{%
      \Succ(y)=\Succ(x)\rightarrow f(d)=f(u)
    }{%
      \FormulaRefAuto{m_0\in M\dsep f\colon M\to M\dsep x\in\mathbb{N}\dsep u\in M\dsep \forall c\in M\,\bigl((x,c)\in \RecCore{m_0}{f}\rightarrow c=u\bigr)\dsep y\in\mathbb{N}\dsep d\in M\dsep (y,d)\in \Phi_{x,u}\vdash \Succ(y)=\Succ(x)\rightarrow f(d)=f(u)}{1,2,3,4,5,6,7,8}}
    \proofstepwidestar[1,2,3,4,5,6,7,8]{%
      (\Succ(y),f(d))\in \Phi_{x,u}
    }{%
      \FormulaRefAuto{m_0\in M\dsep f\colon M\to M\dsep x\in\mathbb{N}\dsep u\in M\dsep (z,e)\in \RecCore{m_0}{f}\dsep z=\Succ(x)\rightarrow e=f(u)\vdash (z,e)\in \Phi_{x,u}}{1,2,3,4,10,11}}
  \closeproofpartwideindR

  \proofpartwide{Schluss}
    \proofstepwidestar[1]{m_0\in M}{\rA}
    \proofstepwidestar[2]{f\colon M\to M}{\rA}
    \proofstepwidestar[3]{x\in\mathbb{N}}{\rA}
    \proofstepwidestar[4]{u\in M}{\rA}
    \proofstepwidestar[5]{(x,u)\in \RecCore{m_0}{f}}{\rA}
    \proofstepwidestar[6]{%
      \begin{aligned}[t]
        &\forall c\in M\,\bigl(
        (x,c)\in \RecCore{m_0}{f}\\
        &\rightarrow c=u
        \bigr)
      \end{aligned}
    }{\rA}
    \proofstepwidestar[1,2,3,4,5,6]{\AdmRec{m_0}{f}(\Phi_{x,u})}{%
      \makecell[l]{%
        \FormulaRefAuto{\AdmRec{m_0}{f}(H)}(%
        \\ \quad \FormulaRefAuto{m_0\in M\dsep f\colon M\to M\dsep x\in\mathbb{N}\dsep u\in M \vdash \Phi_{x,u}\subseteq \mathbb{N}\times M}{1,2,3,4},%
        \\ \quad \FormulaRefAuto{m_0\in M\dsep f\colon M\to M\dsep x\in\mathbb{N}\dsep u\in M \vdash (0,m_0)\in \Phi_{x,u}}{1,2,3,4},%
        \\ \quad \FormulaRefAuto{m_0\in M\dsep f\colon M\to M\dsep x\in\mathbb{N}\dsep u\in M\dsep \forall c\in M\,\bigl((x,c)\in \RecCore{m_0}{f}\rightarrow c=u\bigr)\dsep y\in\mathbb{N}\dsep d\in M\dsep (y,d)\in \Phi_{x,u}\vdash (\Succ(y),f(d))\in \Phi_{x,u}}{1,2,3,4,6}%
        \\ )%
      }%
    }
  \closeproofpartwide
\end{tabproofsplitwide}

\FormulaThmDeltaR[Die Nachfolgerstufe des Rekursionskerns ist durch die Vorstufe festgelegt]{%
  \begin{aligned}[t]
    &m_0\in M\dsep f\colon M\to M\dsep n\in\mathbb{N}\dsep a\in M\dsep{}\\
    &(n,a)\in \RecCore{m_0}{f}\dsep{}\\
    &\forall c\in M\,\bigl((n,c)\in \RecCore{m_0}{f}\rightarrow c=a\bigr)\dsep{}\\
    &(\Succ(n),b)\in \RecCore{m_0}{f}\vdash b=f(a)
  \end{aligned}
}{%
  m_0\in M\dsep f\colon M\to M\dsep n\in\mathbb{N}\dsep a\in M\dsep
  (n,a)\in \RecCore{m_0}{f}\dsep
  \forall c\in M\,\bigl((n,c)\in \RecCore{m_0}{f}\rightarrow c=a\bigr)\dsep
  (\Succ(n),b)\in \RecCore{m_0}{f}\vdash b=f(a)
}{%
  \DeltaDecl{Mengen}{M\dsep n\dsep a\dsep b\dsep c\dsep m_0}%
  \DeltaDecl{Funktionen}{f}%
}
\begin{tabproofwide}
  \proofstepwidestar[1]{m_0\in M}{\rA}
  \proofstepwidestar[2]{f\colon M\to M}{\rA}
  \proofstepwidestar[3]{n\in\mathbb{N}}{\rA}
  \proofstepwidestar[4]{a\in M}{\rA}
  \proofstepwidestar[5]{(n,a)\in \RecCore{m_0}{f}}{\rA}
  \proofstepwidestar[6]{%
    \forall c\in M\,\bigl((n,c)\in \RecCore{m_0}{f}\rightarrow c=a\bigr)
  }{\rA}
  \proofstepwidestar[7]{(\Succ(n),b)\in \RecCore{m_0}{f}}{\rA}

  \proofstepwidestar[1,2,3,4,5,6]{\AdmRec{m_0}{f}(\Phi_{n,a})}{%
    \FormulaRefAuto{m_0\in M\dsep f\colon M\to M\dsep x\in\mathbb{N}\dsep u\in M\dsep (x,u)\in \RecCore{m_0}{f}\dsep \forall c\in M\,\bigl((x,c)\in \RecCore{m_0}{f}\rightarrow c=u\bigr)\vdash \AdmRec{m_0}{f}(\Phi_{x,u})}{1,2,3,4,5,6}}

  \proofstepwidestar[1,2]{\RecCore{m_0}{f}\subseteq \Phi_{n,a}}{%
    \FormulaRefAuto{m_0\in M\dsep f\colon M\to M\dsep \AdmRec{m_0}{f}(H)\vdash \RecCore{m_0}{f}\subseteq H}{1,2,8}}

  \proofstepwidestar[1,2,3,4,5,6,7]{(\Succ(n),b)\in \Phi_{n,a}}{%
    \FormulaRefAuto{A\subseteq B\dsep x\in A\vdash x\in B}{9,7}}

  \proofstepwidestar[1,2,3,4]{%
    \begin{aligned}[t]
      &\Phi_{n,a}\coloneq{}\\
      &\{(y,d)\in \RecCore{m_0}{f}\mid{}\\
      &\qquad y=\Succ(n)\rightarrow d=f(a)\}
    \end{aligned}
  }{%
    \FormulaRefAuto{m_0\in M\dsep f\colon M\to M\dsep x\in\mathbb{N}\dsep u\in M \vdash \Phi_{x,u}\coloneq \{(y,d)\in \RecCore{m_0}{f}\mid y=\Succ(x)\rightarrow d=f(u)\}}{1,2,3,4}}

  \proofstepwidestar[1,2,3,4,5,6,7]{%
    \begin{aligned}[t]
      &(\Succ(n),b)\in{}\\
      &\{(y,d)\in \RecCore{m_0}{f}\mid{}\\
      &\qquad y=\Succ(n)\rightarrow d=f(a)\}
    \end{aligned}
  }{%
    \rIE{11,10}}

  \proofstepwidestar[1,2,3,4,5,6,7]{%
    \Succ(n)=\Succ(n)\rightarrow b=f(a)
  }{%
    \FormulaRefAuto{x \in \{x \in A \mid P(x)\}\vdash P(x)}{12}}

  \proofstepwidestar{\Succ(n)=\Succ(n)}{\rII}

  \proofstepwidestar[1,2,3,4,5,6,7]{b=f(a)}{\rRE{14,13}}
\end{tabproofwide}

\FormulaThmDeltaK[Fixierungslemma für Rekursionskernstufen]{%
  \begin{aligned}[t]
    &m_0\in M\dsep f\colon M\to M\dsep x\in\mathbb{N}\dsep u\in M\dsep{}\\
    &(x,u)\in \RecCore{m_0}{f}\dsep
    \forall c\in M\,\bigl((x,c)\in \RecCore{m_0}{f}\rightarrow c=u\bigr)
    \vdash{}\\
    &\bigl((0,a)\in \RecCore{m_0}{f}\rightarrow a=m_0\bigr)\land{}\\
    &\bigl((\Succ(x),b)\in \RecCore{m_0}{f}\rightarrow b=f(u)\bigr)
  \end{aligned}
}{RecCoreStageBarrier}{%
  \DeltaDecl{Mengen}{M\dsep x\dsep u\dsep a\dsep b\dsep c\dsep m_0}%
  \DeltaDecl{Funktionen}{f}%
}
\begin{tabproofwide}
  \proofstepwidestar[1]{m_0\in M}{\rA}
  \proofstepwidestar[2]{f\colon M\to M}{\rA}
  \proofstepwidestar[3]{x\in\mathbb{N}}{\rA}
  \proofstepwidestar[4]{u\in M}{\rA}
  \proofstepwidestar[5]{(x,u)\in \RecCore{m_0}{f}}{\rA}
  \proofstepwidestar[6]{%
    \forall c\in M\,\bigl((x,c)\in \RecCore{m_0}{f}\rightarrow c=u\bigr)
  }{\rA}

  \proofstepwidestar[7]{(0,a)\in \RecCore{m_0}{f}}{\rA}
  \proofstepwidestar[1,2,7]{a=m_0}{%
    \FormulaRefAuto{m_0\in M\dsep f\colon M\to M\dsep
    (0,a)\in \RecCore{m_0}{f}\vdash a=m_0}{1,2,7}}
  \proofstepwidestar[1,2]{(0,a)\in \RecCore{m_0}{f}\rightarrow a=m_0}{%
    \rRI{7,8}}

  \proofstepwidestar[10]{(\Succ(x),b)\in \RecCore{m_0}{f}}{\rA}
  \proofstepwidestar[1,2,3,4,5,6,10]{b=f(u)}{%
    \FormulaRefAuto{m_0\in M\dsep f\colon M\to M\dsep n\in\mathbb{N}\dsep a\in M\dsep
    (n,a)\in \RecCore{m_0}{f}\dsep
    \forall c\in M\,\bigl((n,c)\in \RecCore{m_0}{f}\rightarrow c=a\bigr)\dsep
    (\Succ(n),b)\in \RecCore{m_0}{f}\vdash b=f(a)}{1,2,3,4,5,6,10}}
  \proofstepwidestar[1,2,3,4,5,6]{%
    (\Succ(x),b)\in \RecCore{m_0}{f}\rightarrow b=f(u)
  }{\rRI{10,11}}

  \proofstepwidestar[1,2,3,4,5,6]{%
    \bigl((0,a)\in \RecCore{m_0}{f}\rightarrow a=m_0\bigr)\land
    \bigl((\Succ(x),b)\in \RecCore{m_0}{f}\rightarrow b=f(u)\bigr)
  }{\rAI{9,12}}
\end{tabproofwide}

\subsubsection{Das Stufenlemma}

Dieser Block ist der eigentliche Punkt der Konstruktion. Statt Existenz und
Eindeutigkeit jeder Stufe getrennt zu beweisen, wird direkt über das Prädikat
\[
  U(n)\;:\Longleftrightarrow\;
  \exists!a\in M\,\bigl((n,a)\in\RecCore{m_0}{f}\bigr)
\]
induktiv argumentiert.

\FormulaThmDeltaK[Stufenlemma des Rekursionskerns]{%
  m_0\in M\dsep f\colon M\to M\dsep n\in\mathbb{N}
  \vdash \exists! a\in M\,\bigl((n,a)\in \RecCore{m_0}{f}\bigr)
}{RecCoreStageUnique}{%
  \DeltaDecl{Mengen}{M\dsep n\dsep x\dsep a\dsep b\dsep c\dsep u\dsep m_0}%
  \DeltaDecl{Funktionen}{f}%
}
\begin{tabproofsplitwide}
  \proofpartwideindR[Induktionsanfang]{%
    m_0\in M\dsep f\colon M\to M
    \vdash \exists! a\in M\,\bigl((0,a)\in \RecCore{m_0}{f}\bigr)
  }{%
    m_0\in M\dsep f\colon M\to M
    \vdash \exists! a\in M\,\bigl((0,a)\in \RecCore{m_0}{f}\bigr)
  }
    \proofstepwidestar[1]{m_0\in M}{\rA}
    \proofstepwidestar[2]{f\colon M\to M}{\rA}
    \proofstepwidestar[1,2]{(0,m_0)\in \RecCore{m_0}{f}}{%
      \FormulaRefAuto{m_0\in M\dsep f\colon M\to M \vdash (0,m_0)\in \RecCore{m_0}{f}}{1,2}}
    \proofstepwidestar[1,2]{%
      \exists a\in M\,\bigl((0,a)\in \RecCore{m_0}{f}\bigr)
    }{\rEI{\rAI{1,3}}}

    \proofstepwidestar[5]{u\in M}{\rA}
    \proofstepwidestar[6]{c\in M}{\rA}
    \proofstepwidestar[7]{(0,u)\in \RecCore{m_0}{f}}{\rA}
    \proofstepwidestar[8]{(0,c)\in \RecCore{m_0}{f}}{\rA}
    \proofstepwidestar[1,2,7]{u=m_0}{%
      \FormulaRefAuto{m_0\in M\dsep f\colon M\to M\dsep
      (0,a)\in \RecCore{m_0}{f}\vdash a=m_0}{1,2,7}}
    \proofstepwidestar[1,2,8]{c=m_0}{%
      \FormulaRefAuto{m_0\in M\dsep f\colon M\to M\dsep
      (0,a)\in \RecCore{m_0}{f}\vdash a=m_0}{1,2,8}}
    \proofstepwidestar[1,2,7,8]{u=c}{%
      \FormulaRefAuto{a = b,\, c = b \vdash a = c}{9,10}}

    \proofstepwidestar[1,2]{%
      \exists! a\in M\,\bigl((0,a)\in \RecCore{m_0}{f}\bigr)
    }{\UEI{4,5,6,7,8,11}}
  \closeproofpartwideindR

  \proofpartwideindR[Induktionsschritt]{%
    \begin{aligned}[t]
      &m_0\in M\dsep f\colon M\to M\dsep x\in\mathbb{N}\dsep{}\\
      &\exists! a\in M\,\bigl((x,a)\in \RecCore{m_0}{f}\bigr)
      \vdash{}\\
      &\exists! b\in M\,\bigl((\Succ(x),b)\in \RecCore{m_0}{f}\bigr)
    \end{aligned}
  }{%
    m_0\in M\dsep f\colon M\to M\dsep x\in\mathbb{N}\dsep
    \exists! a\in M\,\bigl((x,a)\in \RecCore{m_0}{f}\bigr)
    \vdash
    \exists! b\in M\,\bigl((\Succ(x),b)\in \RecCore{m_0}{f}\bigr)
  }
    \proofstepwidestar[1]{m_0\in M}{\rA}
    \proofstepwidestar[2]{f\colon M\to M}{\rA}
    \proofstepwidestar[3]{x\in\mathbb{N}}{\rA}
    \proofstepwidestar[4]{%
      \exists! a\in M\,\bigl((x,a)\in \RecCore{m_0}{f}\bigr)
    }{\rA}

    \proofstepwidestar[5]{u\in M}{\rA}
    \proofstepwidestar[6]{(x,u)\in \RecCore{m_0}{f}}{\rA}
    \proofstepwidestar[1,2,3,5,6]{(\Succ(x),f(u))\in \RecCore{m_0}{f}}{%
      \FormulaRefAuto{m_0\in M\dsep f\colon M\to M\dsep n\in\mathbb{N}\dsep a\in M\dsep
      (n,a)\in \RecCore{m_0}{f}
      \vdash (\Succ(n),f(a))\in \RecCore{m_0}{f}}{1,2,3,5,6}}
    \proofstepwidestar[1,2,3,5,6]{%
      \exists b\in M\,\bigl((\Succ(x),b)\in \RecCore{m_0}{f}\bigr)
    }{\rEI{\rAI{%
      \FormulaRefAuto{F\colon A\to B\dsep x\in A\vdash F(x)\in B}{2,5},7}}}

    \proofstepwidestar[4,5,6]{%
      \forall c\in M\,\bigl((x,c)\in \RecCore{m_0}{f}\rightarrow c=u\bigr)
    }{\UEE{4,5,6}}

    \proofstepwidestar[10]{b\in M}{\rA}
    \proofstepwidestar[11]{c\in M}{\rA}
    \proofstepwidestar[12]{(\Succ(x),b)\in \RecCore{m_0}{f}}{\rA}
    \proofstepwidestar[13]{(\Succ(x),c)\in \RecCore{m_0}{f}}{\rA}
    \proofstepwidestar[1,2,3,5,6,9,12]{b=f(u)}{%
      \FormulaRefAuto{m_0\in M\dsep f\colon M\to M\dsep n\in\mathbb{N}\dsep a\in M\dsep
      (n,a)\in \RecCore{m_0}{f}\dsep
      \forall c\in M\,\bigl((n,c)\in \RecCore{m_0}{f}\rightarrow c=a\bigr)\dsep
      (\Succ(n),b)\in \RecCore{m_0}{f}\vdash b=f(a)}{1,2,3,5,6,9,12}}
    \proofstepwidestar[1,2,3,5,6,9,13]{c=f(u)}{%
      \FormulaRefAuto{m_0\in M\dsep f\colon M\to M\dsep n\in\mathbb{N}\dsep a\in M\dsep
      (n,a)\in \RecCore{m_0}{f}\dsep
      \forall c\in M\,\bigl((n,c)\in \RecCore{m_0}{f}\rightarrow c=a\bigr)\dsep
      (\Succ(n),b)\in \RecCore{m_0}{f}\vdash b=f(a)}{1,2,3,5,6,9,13}}
    \proofstepwidestar[1,2,3,5,6,9,12,13]{b=c}{%
      \FormulaRefAuto{a = b,\, c = b \vdash a = c}{14,15}}

    \proofstepwidestar[1,2,3,4,5,6]{%
      \exists! b\in M\,\bigl((\Succ(x),b)\in \RecCore{m_0}{f}\bigr)
    }{\UEI{8,10,11,12,13,16}}
    \proofstepwidestar[1,2,3,4]{%
      \exists! b\in M\,\bigl((\Succ(x),b)\in \RecCore{m_0}{f}\bigr)
    }{\UEE{4,5,6,17}}
  \closeproofpartwideindR

  \proofpartwide{Schluss}
    \proofstepwidestar[1]{m_0\in M}{\rA}
    \proofstepwidestar[2]{f\colon M\to M}{\rA}
    \proofstepwidestar[3]{n\in\mathbb{N}}{\rA}

    \proofstepwidestar[1,2,3]{%
      \exists! a\in M\,\bigl((n,a)\in \RecCore{m_0}{f}\bigr)
    }{%
      \makecell[l]{%
        \FormulaRefAuto{PeanoInductionRule}(
        \\ \quad \FormulaRefAuto{m_0\in M\dsep f\colon M\to M
          \vdash \exists! a\in M\,\bigl((0,a)\in \RecCore{m_0}{f}\bigr)}{1,2},
        \\ \quad \FormulaRefAuto{m_0\in M\dsep f\colon M\to M\dsep x\in\mathbb{N}\dsep
          \exists! a\in M\,\bigl((x,a)\in \RecCore{m_0}{f}\bigr)
          \vdash
          \exists! b\in M\,\bigl((\Succ(x),b)\in \RecCore{m_0}{f}\bigr)}{1,2},
        \\ \quad 3
        \\ )
      }%
    }
  \closeproofpartwide
\end{tabproofsplitwide}

\iffalse

\paragraph{Existenz}

\FormulaThmDelta[Existenz eines Werts in jeder Stufe]{%
  m_0\in M\dsep f\colon M\to M\dsep n\in\mathbb{N}
  \vdash \exists a\in M\,\bigl((n,a)\in \RecCore{m_0}{f}\bigr)
}{%
  \DeltaDecl{Mengen}{M\dsep n\dsep x\dsep a\dsep b\dsep m_0}%
  \DeltaDecl{Funktionen}{f}%
}
\begin{tabproofsplitwide}
  \proofpartwideindR[Induktionsanfang]{%
    m_0\in M\dsep f\colon M\to M
    \vdash \exists a\in M\,\bigl((0,a)\in \RecCore{m_0}{f}\bigr)
  }{%
    m_0\in M\dsep f\colon M\to M
    \vdash \exists a\in M\,\bigl((0,a)\in \RecCore{m_0}{f}\bigr)
  }
    \proofstepwidestar[1]{m_0\in M}{\rA}
    \proofstepwidestar[2]{f\colon M\to M}{\rA}
    \proofstepwidestar[1,2]{(0,m_0)\in \RecCore{m_0}{f}}{%
      \FormulaRefAuto{m_0\in M\dsep f\colon M\to M \vdash (0,m_0)\in \RecCore{m_0}{f}}{1,2}}
    \proofstepwidestar[1,2]{%
      \exists a\in M\,\bigl((0,a)\in \RecCore{m_0}{f}\bigr)
    }{%
      \rEI{\rAI{1,3}}}
  \closeproofpartwideindR

    \proofpartwideindR[Induktionsschritt]{%
      \begin{aligned}[t]
        &m_0\in M\dsep f\colon M\to M\dsep{}\\
        &x\in\mathbb{N}\dsep
          \exists a\in M\,\bigl((x,a)\in \RecCore{m_0}{f}\bigr)\vdash{}\\
        &\exists b\in M\,\bigl((\Succ(x),b)\in \RecCore{m_0}{f}\bigr)
      \end{aligned}
    }{%
      m_0\in M\dsep f\colon M\to M\dsep
      x\in\mathbb{N}\dsep
      \exists a\in M\,\bigl((x,a)\in \RecCore{m_0}{f}\bigr)\vdash
      \exists b\in M\,\bigl((\Succ(x),b)\in \RecCore{m_0}{f}\bigr)
    }
    \proofstepwidestar[1]{m_0\in M}{\rA}
    \proofstepwidestar[2]{f\colon M\to M}{\rA}
    \proofstepwidestar[3]{x\in\mathbb{N}}{\rA}
    \proofstepwidestar[4]{%
      \exists a\in M\,\bigl((x,a)\in \RecCore{m_0}{f}\bigr)
    }{\rA}

    \proofstepwidestar[5]{a\in M}{\rA}
    \proofstepwidestar[6]{(x,a)\in \RecCore{m_0}{f}}{\rA}

    \proofstepwidestar[2,5]{f(a)\in M}{%
      \FormulaRefAuto{F\colon A\to B\dsep x\in A\vdash F(x)\in B}{2,5}}

    \proofstepwidestar[1,2,3,5,6]{%
      (\Succ(x),f(a))\in \RecCore{m_0}{f}
    }{%
      \FormulaRefAuto{m_0\in M\dsep f\colon M\to M\dsep n\in\mathbb{N}\dsep a\in M\dsep (n,a)\in \RecCore{m_0}{f}\vdash (\Succ(n),f(a))\in \RecCore{m_0}{f}}{1,2,3,5,6}}

    \proofstepwidestar[1,2,3,5,6]{%
      \exists b\in M\,\bigl((\Succ(x),b)\in \RecCore{m_0}{f}\bigr)
    }{%
      \rEI{\rAI{7,8}}}

    \proofstepwidestar[1,2,3,4]{%
      \exists b\in M\,\bigl((\Succ(x),b)\in \RecCore{m_0}{f}\bigr)
    }{%
      \rEE{4,5,6,9}}
  \closeproofpartwideindR

  \proofpartwide{Schluss}
    \proofstepwidestar[1]{m_0\in M}{\rA}
    \proofstepwidestar[2]{f\colon M\to M}{\rA}
    \proofstepwidestar[3]{n\in\mathbb{N}}{\rA}

    \proofstepwidestar[1,2,3]{%
      \exists a\in M\,\bigl((n,a)\in \RecCore{m_0}{f}\bigr)
    }{%
      \makecell[l]{%
        \FormulaRefAuto{PeanoInductionRule}(%
        \\ \quad \FormulaRefAuto{m_0\in M\dsep f\colon M\to M \vdash \exists a\in M\,\bigl((0,a)\in \RecCore{m_0}{f}\bigr)}{1,2},%
        \\ \quad \FormulaRefAuto{m_0\in M\dsep f\colon M\to M\dsep
      x\in\mathbb{N}\dsep
      \exists a\in M\,\bigl((x,a)\in \RecCore{m_0}{f}\bigr)\vdash
      \exists b\in M\,\bigl((\Succ(x),b)\in \RecCore{m_0}{f}\bigr)}{1,2},%
        \\ \quad 3%
        \\ )%
      }%
    }
  \closeproofpartwide
\end{tabproofsplitwide}

\paragraph{Eindeutigkeit}

\FormulaThmDeltaR[Eindeutigkeit in der Nullstufe]{%
  \begin{aligned}[t]
    &m_0\in M\dsep f\colon M\to M\dsep{}\\
    &(0,u)\in \RecCore{m_0}{f}\dsep
      (0,c)\in \RecCore{m_0}{f}\vdash{}\\
    &u=c
  \end{aligned}
}{%
  m_0\in M\dsep f\colon M\to M\dsep
  (0,u)\in \RecCore{m_0}{f}\dsep
  (0,c)\in \RecCore{m_0}{f}\vdash u=c
}{%
  \DeltaDecl{Mengen}{M\dsep u\dsep c\dsep m_0}%
  \DeltaDecl{Funktionen}{f}%
}
\begin{tabproof}
  \proofstep{1}{m_0\in M}{\rA}
  \proofstep{2}{f\colon M\to M}{\rA}
  \proofstep{3}{(0,u)\in \RecCore{m_0}{f}}{\rA}
  \proofstep{4}{(0,c)\in \RecCore{m_0}{f}}{\rA}

  \proofstep{1,2,3}{u=m_0}{%
    \FormulaRefAuto{m_0\in M\dsep f\colon M\to M\dsep (0,a)\in \RecCore{m_0}{f}\vdash a=m_0}{1,2,3}}

  \proofstep{1,2,4}{c=m_0}{%
    \FormulaRefAuto{m_0\in M\dsep f\colon M\to M\dsep (0,a)\in \RecCore{m_0}{f}\vdash a=m_0}{1,2,4}}

  \proofstep{1,2,3,4}{u=c}{%
    \FormulaRefAuto{a = b,\, c = b \vdash a = c}{5,6}}
\end{tabproof}


\FormulaThmDeltaR[Eindeutigkeit in der Nachfolgerstufe bei eindeutiger Vorstufe]{%
  \begin{aligned}[t]
    &m_0\in M\dsep f\colon M\to M\dsep x\in\mathbb{N}\dsep{}\\
    &\forall u\in M\,\bigl((x,u)\in \RecCore{m_0}{f}\rightarrow{}\\
    &\qquad \forall c\in M\,\bigl((x,c)\in \RecCore{m_0}{f}\rightarrow u=c\bigr)\bigr)\dsep{}\\
    &(\Succ(x),d)\in \RecCore{m_0}{f}\dsep
      (\Succ(x),v)\in \RecCore{m_0}{f}\vdash{}\\
    &d=v
  \end{aligned}
}{%
  m_0\in M\dsep f\colon M\to M\dsep x\in\mathbb{N}\dsep
  \forall u\in M\,\bigl((x,u)\in \RecCore{m_0}{f}\rightarrow
  \forall c\in M\,\bigl((x,c)\in \RecCore{m_0}{f}\rightarrow u=c\bigr)\bigr)\dsep
  (\Succ(x),d)\in \RecCore{m_0}{f}\dsep
  (\Succ(x),v)\in \RecCore{m_0}{f}\vdash d=v
}{%
  \DeltaDecl{Mengen}{M\dsep x\dsep d\dsep v\dsep u\dsep c\dsep m_0}%
  \DeltaDecl{Funktionen}{f}%
}
\begin{tabproof}
  \proofstep{1}{m_0\in M}{\rA}
  \proofstep{2}{f\colon M\to M}{\rA}
  \proofstep{3}{x\in\mathbb{N}}{\rA}
  \proofstep{4}{%
    \begin{aligned}[t]
      &\forall u\in M\,\bigl(
      (x,u)\in \RecCore{m_0}{f}\rightarrow{}\\
      &\forall c\in M\,\bigl(
      (x,c)\in \RecCore{m_0}{f}\rightarrow u=c
      \bigr)\bigr)
    \end{aligned}
  }{\rA}
  \proofstep{5}{(\Succ(x),d)\in \RecCore{m_0}{f}}{\rA}
  \proofstep{6}{(\Succ(x),v)\in \RecCore{m_0}{f}}{\rA}

  \proofstep{1,2,3}{%
    \exists u\in M\,\bigl((x,u)\in \RecCore{m_0}{f}\bigr)
  }{%
    \FormulaRefAuto{m_0\in M\dsep f\colon M\to M\dsep n\in\mathbb{N}\vdash \exists a\in M\,\bigl((n,a)\in \RecCore{m_0}{f}\bigr)}{1,2,3}}

  \proofstep{8}{u\in M}{\rA}
  \proofstep{9}{(x,u)\in \RecCore{m_0}{f}}{\rA}

  \proofstep{4,8}{%
    \begin{aligned}[t]
      &(x,u)\in \RecCore{m_0}{f}\rightarrow{}\\
      &\forall c\in M\,\bigl(
      (x,c)\in \RecCore{m_0}{f}\rightarrow u=c
      \bigr)
    \end{aligned}
  }{%
    \rRE{\rUE{4},8}}

  \proofstep{4,8,9}{%
    \begin{aligned}[t]
      &\forall c\in M\,\bigl(
      (x,c)\in \RecCore{m_0}{f}\rightarrow{}\\
      &u=c
      \bigr)
    \end{aligned}
  }{%
    \rRE{9,10}}

  \proofstep{1,2,3,8,9,11,5}{d=f(u)}{%
    \FormulaRefAuto{m_0\in M\dsep f\colon M\to M\dsep n\in\mathbb{N}\dsep a\in M\dsep (n,a)\in \RecCore{m_0}{f}\dsep \forall c\in M\,\bigl((n,c)\in \RecCore{m_0}{f}\rightarrow c=a\bigr)\dsep (\Succ(n),b)\in \RecCore{m_0}{f}\vdash b=f(a)}{1,2,3,8,9,11,5}}

  \proofstep{1,2,3,8,9,11,6}{v=f(u)}{%
    \FormulaRefAuto{m_0\in M\dsep f\colon M\to M\dsep n\in\mathbb{N}\dsep a\in M\dsep (n,a)\in \RecCore{m_0}{f}\dsep \forall c\in M\,\bigl((n,c)\in \RecCore{m_0}{f}\rightarrow c=a\bigr)\dsep (\Succ(n),b)\in \RecCore{m_0}{f}\vdash b=f(a)}{1,2,3,8,9,11,6}}

  \proofstep{1,2,3,4,5,6,8,9}{d=v}{%
    \FormulaRefAuto{a = b,\, c = b \vdash a = c}{12,13}}

  \proofstep{1,2,3,4,5,6}{d=v}{%
    \rEE{7,8,9,14}}
\end{tabproof}

\FormulaThmDeltaR[Eindeutigkeit eines Werts in jeder Stufe]{%
  \begin{aligned}[t]
    &m_0\in M\dsep f\colon M\to M\dsep n\in\mathbb{N}\dsep{}\\
    &a\in M\dsep b\in M\dsep{}\\
    &(n,a)\in \RecCore{m_0}{f}\dsep (n,b)\in \RecCore{m_0}{f}\vdash a=b
  \end{aligned}
}{%
  m_0\in M\dsep f\colon M\to M\dsep n\in\mathbb{N}\dsep
  a\in M\dsep b\in M\dsep
  (n,a)\in \RecCore{m_0}{f}\dsep (n,b)\in \RecCore{m_0}{f}\vdash a=b
}{%
  \DeltaDecl{Mengen}{M\dsep n\dsep x\dsep a\dsep b\dsep c\dsep d\dsep u\dsep m_0}%
  \DeltaDecl{Funktionen}{f}%
}
\begin{tabproofsplitwide}
  \proofpartwideindR[Induktionsanfang]{%
    \begin{aligned}[t]
      &m_0\in M\dsep f\colon M\to M \vdash{}\\
      &\forall u\in M\,\bigl((0,u)\in \RecCore{m_0}{f}\rightarrow{}\\
      &\qquad \forall c\in M\,\bigl((0,c)\in \RecCore{m_0}{f}\rightarrow u=c\bigr)\bigr)
    \end{aligned}
  }{%
    m_0\in M\dsep f\colon M\to M \vdash
    \forall u\in M\,\bigl((0,u)\in \RecCore{m_0}{f}\rightarrow
    \forall c\in M\,\bigl((0,c)\in \RecCore{m_0}{f}\rightarrow u=c\bigr)\bigr)
  }
    \proofstepwidestar[1]{m_0\in M}{\rA}
    \proofstepwidestar[2]{f\colon M\to M}{\rA}
    \proofstepwidestar[3]{u\in M}{\rA}
    \proofstepwidestar[4]{(0,u)\in \RecCore{m_0}{f}}{\rA}
    \proofstepwidestar[5]{c\in M}{\rA}
    \proofstepwidestar[6]{(0,c)\in \RecCore{m_0}{f}}{\rA}

    \proofstepwidestar[1,2,4,6]{u=c}{%
      \FormulaRefAuto{m_0\in M\dsep f\colon M\to M\dsep (0,u)\in \RecCore{m_0}{f}\dsep (0,c)\in \RecCore{m_0}{f}\vdash u=c}{1,2,4,6}}

    \proofstepwidestar[1,2,4,5]{%
      (0,c)\in \RecCore{m_0}{f}\rightarrow u=c
    }{%
      \rRI{6,7}}

    \proofstepwidestar[1,2,4]{%
      \forall c\in M\,\bigl((0,c)\in \RecCore{m_0}{f}\rightarrow u=c\bigr)
    }{%
      \rUI{\rRI{5,8}}}

    \proofstepwidestar[1,2,3]{%
      \begin{aligned}[t]
        &(0,u)\in \RecCore{m_0}{f}\rightarrow{}\\
        &\forall c\in M\,\bigl((0,c)\in \RecCore{m_0}{f}\rightarrow u=c\bigr)
      \end{aligned}
    }{%
      \rRI{4,9}}

    \proofstepwidestar[1,2]{%
      \begin{aligned}[t]
        &\forall u\in M\,\bigl((0,u)\in \RecCore{m_0}{f}\rightarrow{}\\
        &\qquad \forall c\in M\,\bigl((0,c)\in \RecCore{m_0}{f}\rightarrow u=c\bigr)\bigr)
      \end{aligned}
    }{%
      \rUI{\rRI{3,10}}}
  \closeproofpartwideindR

  \proofpartwideindR[Induktionsschritt]{%
  \begin{aligned}[t]
    &m_0\in M\dsep f\colon M\to M\dsep x\in\mathbb{N}\dsep{}\\
    &\forall u\in M\,\bigl((x,u)\in \RecCore{m_0}{f}\rightarrow{}\\
    &\qquad \forall c\in M\,\bigl((x,c)\in \RecCore{m_0}{f}\rightarrow u=c\bigr)\bigr)\vdash{}\\
    &\forall a\in M\,\bigl((\Succ(x),a)\in \RecCore{m_0}{f}\rightarrow{}\\
    &\qquad \forall c\in M\,\bigl((\Succ(x),c)\in \RecCore{m_0}{f}\rightarrow a=c\bigr)\bigr)
  \end{aligned}
}{%
  m_0\in M\dsep f\colon M\to M\dsep x\in\mathbb{N}\dsep
  \forall u\in M\,\bigl((x,u)\in \RecCore{m_0}{f}\rightarrow
  \forall c\in M\,\bigl((x,c)\in \RecCore{m_0}{f}\rightarrow u=c\bigr)\bigr)\vdash
  \forall a\in M\,\bigl((\Succ(x),a)\in \RecCore{m_0}{f}\rightarrow
  \forall c\in M\,\bigl((\Succ(x),c)\in \RecCore{m_0}{f}\rightarrow a=c\bigr)\bigr)
}
    \proofstepwidestar[1]{m_0\in M}{\rA}
    \proofstepwidestar[2]{f\colon M\to M}{\rA}
    \proofstepwidestar[3]{x\in\mathbb{N}}{\rA}
    \proofstepwidestar[4]{%
      \begin{aligned}[t]
        &\forall u\in M\,\bigl((x,u)\in \RecCore{m_0}{f}\rightarrow{}\\
        &\qquad \forall c\in M\,\bigl((x,c)\in \RecCore{m_0}{f}\rightarrow u=c\bigr)\bigr)
      \end{aligned}
    }{\rA}

    \proofstepwidestar[5]{a\in M}{\rA}
    \proofstepwidestar[6]{(\Succ(x),a)\in \RecCore{m_0}{f}}{\rA}
    \proofstepwidestar[7]{c\in M}{\rA}
    \proofstepwidestar[8]{(\Succ(x),c)\in \RecCore{m_0}{f}}{\rA}

    \proofstepwidestar[1,2,3,4,6,8]{a=c}{%
      \FormulaRefAuto{%
        m_0\in M\dsep f\colon M\to M\dsep x\in\mathbb{N}\dsep
        \forall u\in M\,\bigl((x,u)\in \RecCore{m_0}{f}\rightarrow
        \forall c\in M\,\bigl((x,c)\in \RecCore{m_0}{f}\rightarrow u=c\bigr)\bigr)\dsep
        (\Succ(x),d)\in \RecCore{m_0}{f}\dsep
        (\Succ(x),v)\in \RecCore{m_0}{f}\vdash d=v
      }{1,2,3,4,6,8}}

    \proofstepwidestar[1,2,3,5,6,7]{%
      (\Succ(x),c)\in \RecCore{m_0}{f}\rightarrow a=c
    }{%
      \rRI{8,9}}

    \proofstepwidestar[1,2,3,5,6]{%
      \forall c\in M\,\bigl((\Succ(x),c)\in \RecCore{m_0}{f}\rightarrow a=c\bigr)
    }{%
      \rUI{\rRI{7,10}}}

    \proofstepwidestar[1,2,3,5]{%
      \begin{aligned}[t]
        &(\Succ(x),a)\in \RecCore{m_0}{f}\rightarrow{}\\
        &\forall c\in M\,\bigl((\Succ(x),c)\in \RecCore{m_0}{f}\rightarrow a=c\bigr)
      \end{aligned}
    }{%
      \rRI{6,11}}

    \proofstepwidestar[1,2,3]{%
      \begin{aligned}[t]
        &\forall a\in M\,\bigl(
        (\Succ(x),a)\in \RecCore{m_0}{f}\rightarrow{}\\
        &\forall c\in M\,\bigl(
        (\Succ(x),c)\in \RecCore{m_0}{f}\rightarrow{}\\
        &a=c
        \bigr)\bigr)
      \end{aligned}
    }{%
      \rUI{\rRI{5,12}}}

    \proofstepwidestar[1,2,3,4]{%
      \begin{aligned}[t]
        &\forall a\in M\,\bigl(
        (\Succ(x),a)\in \RecCore{m_0}{f}\rightarrow{}\\
        &\forall c\in M\,\bigl(
        (\Succ(x),c)\in \RecCore{m_0}{f}\rightarrow{}\\
        &a=c
        \bigr)\bigr)
      \end{aligned}
    }{%
      \rEE{9,10,11,13}}
  \closeproofpartwideindR

  \proofpartwide{Schluss}
    \proofstepwidestar[1]{m_0\in M}{\rA}
    \proofstepwidestar[2]{f\colon M\to M}{\rA}
    \proofstepwidestar[3]{n\in\mathbb{N}}{\rA}
    \proofstepwidestar[4]{a\in M}{\rA}
    \proofstepwidestar[5]{b\in M}{\rA}
    \proofstepwidestar[6]{(n,a)\in \RecCore{m_0}{f}}{\rA}
    \proofstepwidestar[7]{(n,b)\in \RecCore{m_0}{f}}{\rA}

    \proofstepwidestar[1,2,3]{%
      \begin{aligned}[t]
        &\forall a\in M\,\bigl((n,a)\in \RecCore{m_0}{f}\rightarrow{}\\
        &\qquad \forall c\in M\,\bigl((n,c)\in \RecCore{m_0}{f}\rightarrow a=c\bigr)\bigr)
      \end{aligned}
    }{%
      \makecell[l]{%
        \FormulaRefAuto{PeanoInductionRule}(%
        \\ \quad \FormulaRefAuto{m_0\in M\dsep f\colon M\to M \vdash \forall u\in M\,\bigl((0,u)\in \RecCore{m_0}{f}\rightarrow \forall c\in M\,\bigl((0,c)\in \RecCore{m_0}{f}\rightarrow u=c\bigr)\bigr)}{1,2},%
        \\ \quad \FormulaRefAuto{  m_0\in M\dsep f\colon M\to M\dsep x\in\mathbb{N}\dsep
  \forall u\in M\,\bigl((x,u)\in \RecCore{m_0}{f}\rightarrow
  \forall c\in M\,\bigl((x,c)\in \RecCore{m_0}{f}\rightarrow u=c\bigr)\bigr)\vdash
  \forall a\in M\,\bigl((\Succ(x),a)\in \RecCore{m_0}{f}\rightarrow
  \forall c\in M\,\bigl((\Succ(x),c)\in \RecCore{m_0}{f}\rightarrow a=c\bigr)\bigr)}{1,2},%
        \\ \quad 3%
        \\ )%
      }%
    }

    \proofstepwidestar[1,2,3,4]{%
      \begin{aligned}[t]
        &(n,a)\in \RecCore{m_0}{f}\rightarrow{}\\
        &\forall c\in M\,\bigl((n,c)\in \RecCore{m_0}{f}\rightarrow a=c\bigr)
      \end{aligned}
    }{%
      \rRE{\rUE{8},4}}

    \proofstepwidestar[1,2,3,4,6]{%
      \forall c\in M\,\bigl((n,c)\in \RecCore{m_0}{f}\rightarrow a=c\bigr)
    }{%
      \rRE{6,9}}

    \proofstepwidestar[1,2,3,4,5]{%
      (n,b)\in \RecCore{m_0}{f}\rightarrow a=b
    }{%
      \rRE{\rUE{10},5}}

    \proofstepwidestar[1,2,3,4,5,6,7]{a=b}{%
      \rRE{7,11}}
  \closeproofpartwide
\end{tabproofsplitwide}

\paragraph{Eindeutige Existenz}

\FormulaThmDelta[Eindeutige Existenz eines Werts in jeder Stufe]{%
  m_0\in M\dsep f\colon M\to M\dsep n\in\mathbb{N}
  \vdash \exists! a\in M\,\bigl((n,a)\in \RecCore{m_0}{f}\bigr)
}{%
  \DeltaDecl{Mengen}{M\dsep n\dsep a\dsep b\dsep m_0}%
  \DeltaDecl{Funktionen}{f}%
}
\begin{tabproof}
  \proofstep{1}{m_0\in M}{\rA}
  \proofstep{2}{f\colon M\to M}{\rA}
  \proofstep{3}{n\in\mathbb{N}}{\rA}
  \proofstep{1,2,3}{\exists a\in M\,\bigl((n,a)\in \RecCore{m_0}{f}\bigr)}{%
    \FormulaRefAuto{m_0\in M\dsep f\colon M\to M\dsep n\in\mathbb{N}\vdash \exists a\in M\,\bigl((n,a)\in \RecCore{m_0}{f}\bigr)}{1,2,3}}
  \proofstep{5}{a\in M}{\rA}
  \proofstep{6}{b\in M}{\rA}
  \proofstep{7}{(n,a)\in \RecCore{m_0}{f}}{\rA}
  \proofstep{8}{(n,b)\in \RecCore{m_0}{f}}{\rA}
  \proofstep{1,2,3,5,6,7,8}{a=b}{%
    \FormulaRefAuto{m_0\in M\dsep f\colon M\to M\dsep n\in\mathbb{N}\dsep a\in M\dsep b\in M\dsep (n,a)\in \RecCore{m_0}{f}\dsep (n,b)\in \RecCore{m_0}{f}\vdash a=b}{1,2,3,5,6,7,8}}
  \proofstep{1,2,3}{\exists! a\in M\,\bigl((n,a)\in \RecCore{m_0}{f}\bigr)}{%
    \UEI{4,5,6,7,8,9}}
\end{tabproof}

\fi

\subsubsection{Der Rekursionskern als Funktion}

Nachdem jede Stufe genau einen Wert besitzt, kann der Rekursionskern selbst als
Graph einer Funktion \(\mathbb{N}\to M\) gelesen werden.

\paragraph{Anwendung auf den Rekursionskern}

\FormulaThmDelta[Der Rekursionskern ist eine Abbildung]{%
  m_0\in M\dsep f\colon M\to M \vdash \RecCore{m_0}{f}\colon \mathbb{N}\to M
}{%
  \DeltaDecl{Mengen}{M\dsep x\dsep y\dsep z\dsep m_0}%
  \DeltaDecl{Funktionen}{f}%
}
\begin{tabproof}
  \proofstep{1}{m_0\in M}{\rA}
  \proofstep{2}{f\colon M\to M}{\rA}
  \proofstep{1,2}{\RecCore{m_0}{f}\subseteq \mathbb{N}\times M}{%
    \FormulaRefAuto{m_0\in M\dsep f\colon M\to M \vdash \RecCore{m_0}{f}\subseteq \mathbb{N}\times M}{1,2}}
  \proofstep{4}{n\in\mathbb{N}}{\rA}
  \proofstep{1,2,4}{%
    \exists! a\in M\,\bigl((n,a)\in \RecCore{m_0}{f}\bigr)
  }{\FormulaRefAuto{RecCoreStageUnique}{1,2,4}}
  \proofstep{1,2}{%
    \forall n\in\mathbb{N}\,
    \exists! a\in M\,\bigl((n,a)\in \RecCore{m_0}{f}\bigr)
  }{\rUI{\rRI{4,5}}}
  \proofstep{1,2}{\RecCore{m_0}{f}\colon \mathbb{N}\to M}{%
    \FormulaRefAuto{UniqueValuedGraphFunction}{3,6}}
\end{tabproof}

\paragraph{Rekursionsgleichungen}

\FormulaThmDelta[Anfangswert des Rekursionskerns]{%
  m_0\in M\dsep f\colon M\to M \vdash \RecCore{m_0}{f}(0)=m_0
}{%
  \DeltaDecl{Mengen}{M\dsep m_0}%
  \DeltaDecl{Funktionen}{f}%
}
\begin{tabproof}
  \proofstep{1}{m_0\in M}{\rA}
  \proofstep{2}{f\colon M\to M}{\rA}

  \proofstep{1,2}{\RecCore{m_0}{f}\colon \mathbb{N}\to M}{%
    \FormulaRefAuto{m_0\in M\dsep f\colon M\to M \vdash \RecCore{m_0}{f}\colon \mathbb{N}\to M}{1,2}}

  \proofstep{1,2}{(0,m_0)\in \RecCore{m_0}{f}}{%
    \FormulaRefAuto{m_0\in M\dsep f\colon M\to M \vdash (0,m_0)\in \RecCore{m_0}{f}}{1,2}}

  \proofstep{1,2}{\RecCore{m_0}{f}(0)=m_0}{%
    \FormulaRefAuto{F\colon A\to B\dsep (x,y)\in F\vdash F(x)=y}{3,4}}
\end{tabproof}

\FormulaThmDeltaR[Rekursionsgleichung des Rekursionskerns]{%
  \begin{aligned}
    &m_0\in M\dsep f\colon M\to M\dsep n\in\mathbb{N}\\
    &\vdash \RecCore{m_0}{f}(\Succ(n))=f(\RecCore{m_0}{f}(n))
  \end{aligned}
}{%
  m_0\in M\dsep f\colon M\to M\dsep n\in\mathbb{N}\vdash \RecCore{m_0}{f}(\Succ(n))=f(\RecCore{m_0}{f}(n))
}{%
  \DeltaDecl{Mengen}{M\dsep n\dsep m_0}%
  \DeltaDecl{Funktionen}{f}%
}
\begin{tabproof}
  \proofstep{1}{m_0\in M}{\rA}
  \proofstep{2}{f\colon M\to M}{\rA}
  \proofstep{3}{n\in\mathbb{N}}{\rA}

  \proofstep{1,2}{\RecCore{m_0}{f}\colon \mathbb{N}\to M}{%
    \FormulaRefAuto{m_0\in M\dsep f\colon M\to M \vdash \RecCore{m_0}{f}\colon \mathbb{N}\to M}{1,2}}

  \proofstep{1,2,3}{(n,\RecCore{m_0}{f}(n))\in \RecCore{m_0}{f}}{%
    \FormulaRefAuto{x\in A\vdash (x,F(x))\in F}{4,3}}

  \proofstep{1,2,3}{\RecCore{m_0}{f}(n)\in M}{%
    \FormulaRefAuto{F\colon A\to B\dsep x\in A\vdash F(x)\in B}{4,3}}

  \proofstep{1,2,3}{(\Succ(n),f(\RecCore{m_0}{f}(n)))\in \RecCore{m_0}{f}}{%
    \FormulaRefAuto{m_0\in M\dsep f\colon M\to M\dsep n\in\mathbb{N}\dsep a\in M\dsep (n,a)\in \RecCore{m_0}{f}\vdash (\Succ(n),f(a))\in \RecCore{m_0}{f}}{1,2,3,6,5}}

  \proofstep{1,2,3}{\RecCore{m_0}{f}(\Succ(n))=f(\RecCore{m_0}{f}(n))}{%
    \FormulaRefAuto{(x,y)\in F\vdash F(x)=y}{4,7}}
\end{tabproof}

\paragraph{Eindeutigkeit rekursiv definierter Abbildungen}

Die folgende Induktion trennt die spätere Eindeutigkeitsaussage vom
Existenzbeweis: Jede Abbildung mit demselben Anfangswert und derselben
Rekursionsgleichung stimmt punktweise mit jeder anderen solchen Abbildung
überein.

\FormulaThmDeltaR[Eindeutigkeit rekursiv definierter Abbildungen]{%
  \begin{aligned}
    &G\colon \mathbb{N}\to M\dsep H\colon \mathbb{N}\to M\dsep
    G(0)=m_0\dsep H(0)=m_0\dsep{}\\
    &\forall n\in\mathbb{N}\,G(\Succ(n))=f(G(n))\dsep
    \forall n\in\mathbb{N}\,H(\Succ(n))=f(H(n))
    \vdash G=H
  \end{aligned}
}{%
  G\colon \mathbb{N}\to M\dsep H\colon \mathbb{N}\to M\dsep
  G(0)=m_0\dsep H(0)=m_0\dsep
  \forall n\in\mathbb{N}\,G(\Succ(n))=f(G(n))\dsep
  \forall n\in\mathbb{N}\,H(\Succ(n))=f(H(n))
  \vdash G=H
}{%
  \DeltaDecl{Mengen}{M\dsep n\dsep x\dsep m_0}%
  \DeltaDecl{Funktionen}{f\dsep G\dsep H}%
}
\begin{tabproofsplitwide}
  \proofpartwideindR[Induktionsanfang]{%
    \begin{aligned}[t]
      &G\colon \mathbb{N}\to M\dsep H\colon \mathbb{N}\to M\dsep{}\\
      &G(0)=m_0\dsep H(0)=m_0\dsep{}\\
      &\forall n\in\mathbb{N}\,G(\Succ(n))=f(G(n))\dsep{}\\
      &\forall n\in\mathbb{N}\,H(\Succ(n))=f(H(n))
      \vdash G(0)=H(0)
    \end{aligned}
  }{%
    G\colon \mathbb{N}\to M\dsep H\colon \mathbb{N}\to M\dsep
    G(0)=m_0\dsep H(0)=m_0\dsep
    \forall n\in\mathbb{N}\,G(\Succ(n))=f(G(n))\dsep
    \forall n\in\mathbb{N}\,H(\Succ(n))=f(H(n))
    \vdash G(0)=H(0)
  }
    \proofstepwidestar[1]{G\colon \mathbb{N}\to M}{\rA}
    \proofstepwidestar[2]{H\colon \mathbb{N}\to M}{\rA}
    \proofstepwidestar[3]{G(0)=m_0}{\rA}
    \proofstepwidestar[4]{H(0)=m_0}{\rA}
    \proofstepwidestar[5]{\forall n\in\mathbb{N}\,G(\Succ(n))=f(G(n))}{\rA}
    \proofstepwidestar[6]{\forall n\in\mathbb{N}\,H(\Succ(n))=f(H(n))}{\rA}

    \proofstepwidestar[3,4]{G(0)=H(0)}{%
      \FormulaRefAuto{a = b,\, c = b \vdash a = c}{3,4}}
  \closeproofpartwideindR

  \proofpartwideindR[Induktionsschritt]{%
    \begin{aligned}[t]
      &G\colon \mathbb{N}\to M\dsep H\colon \mathbb{N}\to M\dsep{}\\
      &G(0)=m_0\dsep H(0)=m_0\dsep{}\\
      &\forall n\in\mathbb{N}\,G(\Succ(n))=f(G(n))\dsep{}\\
      &\forall n\in\mathbb{N}\,H(\Succ(n))=f(H(n))\dsep{}\\
      &x\in\mathbb{N}\dsep G(x)=H(x)\vdash G(\Succ(x))=H(\Succ(x))
    \end{aligned}
  }{%
    G\colon \mathbb{N}\to M\dsep H\colon \mathbb{N}\to M\dsep
    G(0)=m_0\dsep H(0)=m_0\dsep
    \forall n\in\mathbb{N}\,G(\Succ(n))=f(G(n))\dsep
    \forall n\in\mathbb{N}\,H(\Succ(n))=f(H(n))\dsep
    x\in\mathbb{N}\dsep G(x)=H(x)\vdash G(\Succ(x))=H(\Succ(x))
  }
    \proofstepwidestar[1]{G\colon \mathbb{N}\to M}{\rA}
    \proofstepwidestar[2]{H\colon \mathbb{N}\to M}{\rA}
    \proofstepwidestar[3]{G(0)=m_0}{\rA}
    \proofstepwidestar[4]{H(0)=m_0}{\rA}
    \proofstepwidestar[5]{\forall n\in\mathbb{N}\,G(\Succ(n))=f(G(n))}{\rA}
    \proofstepwidestar[6]{\forall n\in\mathbb{N}\,H(\Succ(n))=f(H(n))}{\rA}
    \proofstepwidestar[7]{x\in\mathbb{N}}{\rA}
    \proofstepwidestar[8]{G(x)=H(x)}{\rA}

    \proofstepwidestar[1,7]{G(x)\in M}{%
      \FormulaRefAuto{F\colon A\to B\dsep x\in A\vdash F(x)\in B}{1,7}}

    \proofstepwidestar[2,7]{H(x)\in M}{%
      \FormulaRefAuto{F\colon A\to B\dsep x\in A\vdash F(x)\in B}{2,7}}

    \proofstepwidestar[5]{x\in\mathbb{N}\rightarrow G(\Succ(x))=f(G(x))}{\rUE{5}}
    \proofstepwidestar[5,7]{G(\Succ(x))=f(G(x))}{\rRE{11,7}}

    \proofstepwidestar[6]{x\in\mathbb{N}\rightarrow H(\Succ(x))=f(H(x))}{\rUE{6}}
    \proofstepwidestar[6,7]{H(\Succ(x))=f(H(x))}{\rRE{13,7}}

    \proofstepwidestar[1,2,7,8]{f(G(x))=f(H(x))}{%
      \FormulaRefAuto{x\in A\dsep y\in A\dsep x=y\vdash F(x)=F(y)}{9,10,8}}

    \proofstepwidestar[1,2,6,7,8]{H(\Succ(x))=f(G(x))}{%
      \FormulaRefAuto{a = b,\, c = b \vdash a = c}{14,15}}

    \proofstepwidestar[1,2,6,7,8]{f(G(x))=H(\Succ(x))}{%
      \FormulaRefAuto{a=b\vdash b=a}{16}}

    \proofstepwidestar[1,2,5,6,7,8]{G(\Succ(x))=H(\Succ(x))}{%
      \FormulaRefAuto{a = b,\, c = b \vdash a = c}{12,17}}
  \closeproofpartwideindR

  \proofpartwide{Schluss}
    \proofstepwidestar[1]{G\colon \mathbb{N}\to M}{\rA}
    \proofstepwidestar[2]{H\colon \mathbb{N}\to M}{\rA}
    \proofstepwidestar[3]{G(0)=m_0}{\rA}
    \proofstepwidestar[4]{H(0)=m_0}{\rA}
    \proofstepwidestar[5]{\forall n\in\mathbb{N}\,G(\Succ(n))=f(G(n))}{\rA}
    \proofstepwidestar[6]{\forall n\in\mathbb{N}\,H(\Succ(n))=f(H(n))}{\rA}

    \proofstepwidestar[7]{n\in\mathbb{N}}{\rA}

    \proofstepwidestar[1,2,3,4,5,6,7]{G(n)=H(n)}{%
      \makecell[l]{%
        \FormulaRefAuto{PeanoInductionRule}(%
        \\ \quad \FormulaRefAuto{%
          G\colon \mathbb{N}\to M\dsep H\colon \mathbb{N}\to M\dsep
          G(0)=m_0\dsep H(0)=m_0\dsep
          \forall n\in\mathbb{N}\,G(\Succ(n))=f(G(n))\dsep
          \forall n\in\mathbb{N}\,H(\Succ(n))=f(H(n))
          \vdash G(0)=H(0)
        }{1,2,3,4,5,6},%
        \\ \quad \FormulaRefAuto{%
          G\colon \mathbb{N}\to M\dsep H\colon \mathbb{N}\to M\dsep
          G(0)=m_0\dsep H(0)=m_0\dsep
          \forall n\in\mathbb{N}\,G(\Succ(n))=f(G(n))\dsep
          \forall n\in\mathbb{N}\,H(\Succ(n))=f(H(n))\dsep
          x\in\mathbb{N}\dsep G(x)=H(x)\vdash G(\Succ(x))=H(\Succ(x))
        }{1,2,3,4,5,6},%
        \\ \quad 7%
        \\ )%
      }%
    }

    \proofstepwidestar[1,2,3,4,5,6]{\forall n\in\mathbb{N}\,\bigl(G(n)=H(n)\bigr)}{%
      \rUI{\rRI{7,8}}}

    \proofstepwidestar[1,2,3,4,5,6]{G=H}{%
      \FormulaRefAuto{\forall x\in A (F(x)=G(x)) \eqvdash F=G}{9}}
  \closeproofpartwide
\end{tabproofsplitwide}

\subsubsection{Der Dedekindsche Rekursionssatz}

Der Hauptsatz ist nun nur noch die Zusammenführung der beiden vorherigen
Blöcke. Die Existenz kommt vom Rekursionskern als Funktion, die Eindeutigkeit
von der punktweisen Induktion über rekursiv definierte Abbildungen.

\FormulaThmDeltaR[Existenz einer rekursiv definierten Abbildung]{%
  \begin{aligned}
    &m_0\in M\dsep f\colon M\to M\\
    \vdash{}&\exists G\colon \mathbb{N}\to M\,\Bigl(
      G(0)=m_0 \land{}\\
      &\qquad\qquad\qquad\forall n\in\mathbb{N}\,
      G(\Succ(n))=f(G(n))
    \Bigr)
  \end{aligned}
}{%
  m_0\in M\dsep f\colon M\to M\vdash \exists G\colon \mathbb{N}\to M\,\Bigl(G(0)=m_0 \land \forall n\in\mathbb{N}\,G(\Succ(n))=f(G(n))\Bigr)
}{%
  \DeltaDecl{Mengen}{M\dsep n\dsep m_0}%
  \DeltaDecl{Funktionen}{f\dsep G}%
}
\begin{tabproof}
  \proofstep{1}{m_0\in M}{\rA}
  \proofstep{2}{f\colon M\to M}{\rA}

  \proofstep{1,2}{\RecCore{m_0}{f}\colon \mathbb{N}\to M}{%
    \FormulaRefAuto{m_0\in M\dsep f\colon M\to M \vdash \RecCore{m_0}{f}\colon \mathbb{N}\to M}{1,2}}

  \proofstep{1,2}{\RecCore{m_0}{f}(0)=m_0}{%
    \FormulaRefAuto{m_0\in M\dsep f\colon M\to M \vdash \RecCore{m_0}{f}(0)=m_0}{1,2}}

  \proofstep{5}{n\in\mathbb{N}}{\rA}

  \proofstep{1,2,5}{%
    \begin{aligned}[t]
      &\RecCore{m_0}{f}(\Succ(n))={}\\
      &f(\RecCore{m_0}{f}(n))
    \end{aligned}
  }{%
    \FormulaRefAuto{m_0\in M\dsep f\colon M\to M\dsep n\in\mathbb{N}\vdash \RecCore{m_0}{f}(\Succ(n))=f(\RecCore{m_0}{f}(n))}{1,2,5}}

  \proofstep{1,2}{%
    \begin{aligned}[t]
      &\forall n\in\mathbb{N}\,\bigl({}\\
      &\qquad \RecCore{m_0}{f}(\Succ(n))={}\\
      &\qquad f(\RecCore{m_0}{f}(n))\bigr)
    \end{aligned}
  }{%
    \rUI{\rRI{5,6}}}

  \proofstep{1,2}{%
    \begin{aligned}[t]
      &\RecCore{m_0}{f}(0)=m_0 \land{}\\
      &\forall n\in\mathbb{N}\,\bigl({}\\
      &\qquad \RecCore{m_0}{f}(\Succ(n))={}\\
      &\qquad f(\RecCore{m_0}{f}(n))\bigr)
    \end{aligned}
  }{\rAI{4,7}}

  \proofstep{1,2}{%
    \begin{aligned}[t]
      &\exists G\colon \mathbb{N}\to M\,\Bigl({}\\
      &\qquad G(0)=m_0 \land{}\\
      &\qquad \forall n\in\mathbb{N}\,\bigl({}\\
      &\qquad\qquad G(\Succ(n))=f(G(n))\bigr)\Bigr)
    \end{aligned}
  }{%
    \rEI{\rAI{3,8}}}
\end{tabproof}



\FormulaThmDeltaR[Dedekindscher Rekursionssatz]{%
  \begin{aligned}
    &m_0\in M\dsep f\colon M\to M\\
    \vdash{}&\exists! G\colon \mathbb{N}\to M\,\Bigl(
      G(0)=m_0 \land{}\\
      &\qquad\qquad\qquad\forall n\in\mathbb{N}\,
      G(\Succ(n))=f(G(n))
    \Bigr)
  \end{aligned}
}{%
  m_0\in M\dsep f\colon M\to M\vdash \exists! G\colon \mathbb{N}\to M\,\Bigl(G(0)=m_0 \land \forall n\in\mathbb{N}\,G(\Succ(n))=f(G(n))\Bigr)
}{%
  \DeltaDecl{Mengen}{M\dsep m_0}%
  \DeltaDecl{Funktionen}{f\dsep G\dsep H}%
}
\begin{tabproof}
  \proofstep{1}{m_0\in M}{\rA}
  \proofstep{2}{f\colon M\to M}{\rA}

  \proofstep{1,2}{%
    \begin{aligned}[t]
      &\exists G\colon \mathbb{N}\to M\,\Bigl({}\\
      &\qquad G(0)=m_0 \land{}\\
      &\qquad \forall n\in\mathbb{N}\,\bigl({}\\
      &\qquad\qquad G(\Succ(n))=f(G(n))\bigr)\Bigr)
    \end{aligned}
  }{%
    \FormulaRefAuto{m_0\in M\dsep f\colon M\to M\vdash \exists G\colon \mathbb{N}\to M\,\Bigl(G(0)=m_0 \land \forall n\in\mathbb{N}\,G(\Succ(n))=f(G(n))\Bigr)}{1,2}}

  \proofstep{4}{G\colon \mathbb{N}\to M}{\rA}

  \proofstep{5}{%
    \begin{aligned}[t]
      &G(0)=m_0 \land{}\\
      &\forall n\in\mathbb{N}\,\bigl({}\\
      &\qquad G(\Succ(n))=f(G(n))\bigr)
    \end{aligned}
  }{\rA}

  \proofstep{6}{H\colon \mathbb{N}\to M}{\rA}

  \proofstep{7}{%
    \begin{aligned}[t]
      &H(0)=m_0 \land{}\\
      &\forall n\in\mathbb{N}\,\bigl({}\\
      &\qquad H(\Succ(n))=f(H(n))\bigr)
    \end{aligned}
  }{\rA}

  \proofstep{5}{G(0)=m_0}{\rAEa{5}}

  \proofstep{5}{%
    \begin{aligned}[t]
      &\forall n\in\mathbb{N}\,\bigl({}\\
      &\qquad G(\Succ(n))=f(G(n))\bigr)
    \end{aligned}
  }{\rAEb{5}}

  \proofstep{7}{H(0)=m_0}{\rAEa{7}}

  \proofstep{7}{%
    \begin{aligned}[t]
      &\forall n\in\mathbb{N}\,\bigl({}\\
      &\qquad H(\Succ(n))=f(H(n))\bigr)
    \end{aligned}
  }{\rAEb{7}}

  \proofstep{4,5,6,7}{G=H}{%
    \FormulaRefAuto{G\colon \mathbb{N}\to M\dsep H\colon \mathbb{N}\to M\dsep G(0)=m_0\dsep H(0)=m_0\dsep \forall n\in\mathbb{N}\,G(\Succ(n))=f(G(n))\dsep \forall n\in\mathbb{N}\,H(\Succ(n))=f(H(n))\vdash G=H}{4,6,8,10,9,11}}

  \proofstep{1,2}{%
    \begin{aligned}[t]
      &\exists! G\colon \mathbb{N}\to M\,\Bigl({}\\
      &\qquad G(0)=m_0 \land{}\\
      &\qquad \forall n\in\mathbb{N}\,\bigl({}\\
      &\qquad\qquad G(\Succ(n))=f(G(n))\bigr)\Bigr)
    \end{aligned}
  }{%
    \UEI{3,4,5,6,7,12}}
\end{tabproof}

\subsubsection{Die Rekursionsabbildung als benanntes Objekt}

Nach dem Stufenlemma ist der Rekursionskern bereits der Graph einer eindeutig
bestimmten Abbildung. Die Rekursionsabbildung wird deshalb nicht noch einmal
als neuer \(\iota\)-Zeuge eingeführt, sondern direkt als dieser Kern benannt.
Die folgenden Sätze übertragen die bereits bewiesenen Rekursionsgleichungen
auf dieses Symbol und sammeln Zusatzlemmata, die später für Einbettungen
benötigt werden.

\paragraph{Definition der Rekursionsabbildung}

\FormulaDefDeltaK[Definition der Rekursionsabbildung]{%
  m_0\in M\dsep f\colon M\to M \vdash
  \RecFun{m_0}{f}\coloneqq \RecCore{m_0}{f}%
}{RecFunDef}{%
  \DeltaDecl{Mengen}{M\dsep m_0}%
  \DeltaDecl{Funktionen}{f}%
}

\paragraph{Grundlegende Eigenschaften}

\FormulaThmDelta[Die Rekursionsabbildung ist eine Abbildung]{%
  m_0\in M\dsep f\colon M\to M \vdash
  \RecFun{m_0}{f}\colon \mathbb{N}\to M
}{%
  \DeltaDecl{Mengen}{M\dsep m_0}%
  \DeltaDecl{Funktionen}{f}%
}
\begin{tabproof}
  \proofstep{1}{m_0\in M}{\rA}
  \proofstep{2}{f\colon M\to M}{\rA}

  \proofstep{1,2}{\RecFun{m_0}{f}=\RecCore{m_0}{f}}{%
    \FormulaRefAuto{RecFunDef}{1,2}}
  \proofstep{1,2}{\RecCore{m_0}{f}\colon \mathbb{N}\to M}{%
    \FormulaRefAuto{m_0\in M\dsep f\colon M\to M \vdash
    \RecCore{m_0}{f}\colon \mathbb{N}\to M}{1,2}}
  \proofstep{1,2}{\RecFun{m_0}{f}\colon \mathbb{N}\to M}{%
    \rIE{3,4}}
\end{tabproof}

\FormulaThmDelta[Werte der Rekursionsabbildung liegen in $M$]{%
  m_0\in M\dsep f\colon M\to M\dsep n\in\mathbb{N}\vdash
  \RecFun{m_0}{f}(n)\in M%
}{%
  \DeltaDecl{Mengen}{M\dsep m_0\dsep n}%
  \DeltaDecl{Funktionen}{f}%
}
\begin{tabproof}
  \proofstep{1}{m_0\in M}{\rA}
  \proofstep{2}{f\colon M\to M}{\rA}
  \proofstep{3}{n\in\mathbb{N}}{\rA}
  \proofstep{1,2}{\RecFun{m_0}{f}\colon \mathbb{N}\to M}{%
    \FormulaRefAuto{m_0\in M\dsep f\colon M\to M \vdash
    \RecFun{m_0}{f}\colon \mathbb{N}\to M}{1,2}}
  \proofstep{1,2,3}{\RecFun{m_0}{f}(n)\in M}{%
    \FormulaRefAuto{F\colon A\to B\dsep a\in A\vdash F(a)\in B}{4,3}}
\end{tabproof}

\FormulaThmDelta[Anfangswert der Rekursionsabbildung]{%
  m_0\in M\dsep f\colon M\to M \vdash
  \RecFun{m_0}{f}(0)=m_0
}{%
  \DeltaDecl{Mengen}{M\dsep m_0}%
  \DeltaDecl{Funktionen}{f}%
}
\begin{tabproof}
  \proofstep{1}{m_0\in M}{\rA}
  \proofstep{2}{f\colon M\to M}{\rA}

  \proofstep{1,2}{\RecFun{m_0}{f}=\RecCore{m_0}{f}}{%
    \FormulaRefAuto{RecFunDef}{1,2}}
  \proofstep{1,2}{\RecCore{m_0}{f}(0)=m_0}{%
    \FormulaRefAuto{m_0\in M\dsep f\colon M\to M \vdash
    \RecCore{m_0}{f}(0)=m_0}{1,2}}
  \proofstep{1,2}{\RecFun{m_0}{f}(0)=m_0}{\rIE{3,4}}
\end{tabproof}

\FormulaThmDeltaR[Rekursionsgleichung der Rekursionsabbildung]{%
  \begin{aligned}[t]
    &m_0\in M\dsep f\colon M\to M\dsep n\in\mathbb{N}\\
    &\vdash \RecFun{m_0}{f}(\Succ(n))=f(\RecFun{m_0}{f}(n))
  \end{aligned}
}{%
  m_0\in M\dsep f\colon M\to M\dsep n\in\mathbb{N}\vdash
  \RecFun{m_0}{f}(\Succ(n))=f(\RecFun{m_0}{f}(n))
}{%
  \DeltaDecl{Mengen}{M\dsep n\dsep m_0}%
  \DeltaDecl{Funktionen}{f}%
}
\begin{tabproof}
  \proofstep{1}{m_0\in M}{\rA}
  \proofstep{2}{f\colon M\to M}{\rA}
  \proofstep{3}{n\in\mathbb{N}}{\rA}

  \proofstep{1,2}{\RecFun{m_0}{f}=\RecCore{m_0}{f}}{%
    \FormulaRefAuto{RecFunDef}{1,2}}
  \proofstep{1,2,3}{%
    \RecCore{m_0}{f}(\Succ(n))=f(\RecCore{m_0}{f}(n))
  }{%
    \FormulaRefAuto{m_0\in M\dsep f\colon M\to M\dsep n\in\mathbb{N}\vdash
    \RecCore{m_0}{f}(\Succ(n))=f(\RecCore{m_0}{f}(n))}{1,2,3}}
  \proofstep{1,2,3}{%
    \RecFun{m_0}{f}(\Succ(n))=f(\RecFun{m_0}{f}(n))
  }{\rIE{4,5}}
\end{tabproof}

\FormulaThmDeltaR[Rekursionsgleichung der Rekursionsabbildung für Nichtnullstellen]{%
  \begin{aligned}[t]
    &m_0\in M\dsep f\colon M\to M\dsep n\in\mathbb{N}\dsep n\neq0\\
    &\vdash \RecFun{m_0}{f}(n)=f(\RecFun{m_0}{f}(\Pred(n)))
  \end{aligned}
}{%
  m_0\in M\dsep f\colon M\to M\dsep n\in\mathbb{N}\dsep n\neq0\vdash
  \RecFun{m_0}{f}(n)=f(\RecFun{m_0}{f}(\Pred(n)))
}{%
  \DeltaDecl{Mengen}{M\dsep n\dsep k\dsep m_0}%
  \DeltaDecl{Funktionen}{f}%
}
\begin{tabproofwide}
  \proofstepwidestar[1]{m_0\in M}{\rA}
  \proofstepwidestar[2]{f\colon M\to M}{\rA}
  \proofstepwidestar[3]{n\in\mathbb{N}}{\rA}
  \proofstepwidestar[4]{n\neq0}{\rA}

  \proofstepwidestar[3,4]{%
    \exists k\in\mathbb{N}\,\bigl(n=\Succ(k)\bigr)
  }{%
    \FormulaRefAuto{PeanoNonzeroIsSucc}{3,4}}

  \proofstepwidestar[6]{k\in\mathbb{N}}{\rA}
  \proofstepwidestar[7]{n=\Succ(k)}{\rA}

  \proofstepwidestar[6]{\Pred(\Succ(k))=k}{%
    \FormulaRefAuto{PeanoPredSucc}{6}}

  \proofstepwidestar[3,6,7]{\Pred(n)=k}{\rIE{7,8}}
  \proofstepwidestar[3,6,7]{k=\Pred(n)}{%
    \FormulaRefAuto{a=b\vdash b=a}{9}}

  \proofstepwide[1,2,3,4,6,7]{\RecFun{m_0}{f}(n)}{=}{\RecFun{m_0}{f}(\Succ(k))}{\rIE{7}}

  \proofstepwide[1,2,3,4,6,7]{}{=}{f(\RecFun{m_0}{f}(k))}{%
    \FormulaRefAuto{%
      m_0\in M\dsep f\colon M\to M\dsep n\in\mathbb{N}\vdash
      \RecFun{m_0}{f}(\Succ(n))=f(\RecFun{m_0}{f}(n))
    }{1,2,6}}

  \proofstepwide[1,2,3,4,6,7]{}{=}{f(\RecFun{m_0}{f}(\Pred(n)))}{\rIE{10}}

  \proofstepwide[1,2,3,4,6,7]{\RecFun{m_0}{f}(n)}{=}{f(\RecFun{m_0}{f}(\Pred(n)))}{\rChain{11,12,13}}

  \proofstepwidestar[1,2,3,4]{%
    \RecFun{m_0}{f}(n)=f(\RecFun{m_0}{f}(\Pred(n)))
  }{%
    \rEE{5,6,7,14}}
\end{tabproofwide}

\FormulaThmDeltaR{%
  \begin{aligned}[t]
    &m_0\in M\dsep f\colon M\to M\dsep n\in\mathbb{N}\dsep n\neq0\\
    &\vdash f(\RecFun{m_0}{f}(\Pred(n)))=\RecFun{m_0}{f}(n)
  \end{aligned}
}{%
  m_0\in M\dsep f\colon M\to M\dsep n\in\mathbb{N}\dsep n\neq0\vdash
  f(\RecFun{m_0}{f}(\Pred(n)))=\RecFun{m_0}{f}(n)
}{%
  \DeltaDecl{Mengen}{M\dsep n\dsep m_0}%
  \DeltaDecl{Funktionen}{f}%
}
\begin{tabproofwide}
  \proofstepwidestar[1]{m_0\in M}{\rA}
  \proofstepwidestar[2]{f\colon M\to M}{\rA}
  \proofstepwidestar[3]{n\in\mathbb{N}}{\rA}
  \proofstepwidestar[4]{n\neq0}{\rA}
  \proofstepwidestar[1,2,3,4]{%
    \RecFun{m_0}{f}(n)=f(\RecFun{m_0}{f}(\Pred(n)))
  }{%
    \FormulaRefAuto{%
      m_0\in M\dsep f\colon M\to M\dsep n\in\mathbb{N}\dsep n\neq0\vdash
      \RecFun{m_0}{f}(n)=f(\RecFun{m_0}{f}(\Pred(n)))
    }{1,2,3,4}}
  \proofstepwidestar[1,2,3,4]{%
    f(\RecFun{m_0}{f}(\Pred(n)))=\RecFun{m_0}{f}(n)
  }{%
    \FormulaRefAuto{a=b\vdash b=a}{5}}
\end{tabproofwide}
\paragraph{Zusammenhang mit dem Rekursionskern}

Per Definition ist die Rekursionsabbildung der Rekursionskern:
\(\RecFun{m_0}{f}=\RecCore{m_0}{f}\). Entsprechende Beweise verwenden die
Definition \(\FormulaRefAuto{RecFunDef}\) direkt.


\paragraph{Hilfssätze zur Injektivität der Rekursionsabbildung}

\FormulaThmDeltaR[Nullfall des Induktionsschritts für die Rekursionsabbildung]{%
  \begin{aligned}[t]
    &m_0\in M\dsep f\colon M\inj M\dsep{}\\
    &\forall x\in M\,\bigl(f(x)\neq m_0\bigr)\dsep{}\\
    &u\in\mathbb{N}\dsep k\in\mathbb{N}\dsep{}\\
    &\RecFun{m_0}{f}(k)=\RecFun{m_0}{f}(\Succ(u))\vdash k\neq0
  \end{aligned}
}{%
  m_0\in M\dsep f\colon M\inj M\dsep
  \forall x\in M\,\bigl(f(x)\neq m_0\bigr)\dsep
  u\in\mathbb{N}\dsep k\in\mathbb{N}\dsep
  \RecFun{m_0}{f}(k)=\RecFun{m_0}{f}(\Succ(u))\vdash k\neq0
}{%
  \DeltaDecl{Mengen}{M\dsep u\dsep k\dsep x\dsep m_0}%
  \DeltaDecl{Funktionen}{f}%
}
\begin{tabproofwide}
  \proofstepwidestar[1]{m_0\in M}{\rA}
  \proofstepwidestar[2]{f\colon M\inj M}{\rA}
  \proofstepwidestar[2]{f\colon M\to M}{\FormulaRefAuto{Injektive Funktion}{2}}
  \proofstepwidestar[4]{\forall x\in M\,\bigl(f(x)\neq m_0\bigr)}{\rA}
  \proofstepwidestar[5]{u\in\mathbb{N}}{\rA}
  \proofstepwidestar[6]{k\in\mathbb{N}}{\rA}
  \proofstepwidestar[7]{\RecFun{m_0}{f}(k)=\RecFun{m_0}{f}(\Succ(u))}{\rA}

  \proofstepwidestar[1,2,5]{\RecFun{m_0}{f}(u)\in M}{%
    \FormulaRefAuto{m_0\in M\dsep f\colon M\to M\dsep n\in\mathbb{N}\vdash
    \RecFun{m_0}{f}(n)\in M}{1,3,5}}

  \proofstepwidestar[1,2,4,5]{f(\RecFun{m_0}{f}(u))\neq m_0}{%
    \rRE{\rUE{4},8}}

  \proofstepwidestar[10]{k=0}{\rA}

  \proofstepwide[1,2,5]{f(\RecFun{m_0}{f}(u))}{=}{\RecFun{m_0}{f}(\Succ(u))}{%
    \FormulaRefAuto{a=b\vdash b=a}{%
      \FormulaRefAuto{%
        m_0\in M\dsep f\colon M\to M\dsep n\in\mathbb{N}\vdash
        \RecFun{m_0}{f}(\Succ(n))=f(\RecFun{m_0}{f}(n))
      }{1,3,5}}}
  \proofstepwide[7]{}{=}{\RecFun{m_0}{f}(k)}{%
    \FormulaRefAuto{a=b\vdash b=a}{7}}
  \proofstepwide[10]{}{=}{\RecFun{m_0}{f}(0)}{\rIE{10}}
  \proofstepwide[1,2]{}{=}{m_0}{%
    \FormulaRefAuto{m_0\in M\dsep f\colon M\to M\vdash
      \RecFun{m_0}{f}(0)=m_0}{1,3}}
  \proofstepwide[1,2,5,7,10]{f(\RecFun{m_0}{f}(u))}{=}{m_0}{\rChain{11,12,13,14}}

  \proofstepwidestar[1,2,4,5,7,10]{\bot}{\rBI{9,15}}

  \proofstepwidestar[1,2,4,5,7]{k\neq0}{\rCI{10,16}}
\end{tabproofwide}


\FormulaThmDeltaR[Hilfssatz zum Induktionsanfang der Eindeutigkeit der Rekursionsabbildung]{%
  \begin{aligned}[t]
    &m_0\in M\dsep f\colon M\inj M\dsep{}\\
    &\forall x\in M\,\bigl(f(x)\neq m_0\bigr)\dsep{}\\
    &k\in\mathbb{N}\dsep
    \RecFun{m_0}{f}(k)=\RecFun{m_0}{f}(0)\vdash k=0
  \end{aligned}
}{%
  m_0\in M\dsep f\colon M\inj M\dsep
  \forall x\in M\,\bigl(f(x)\neq m_0\bigr)\dsep
  k\in\mathbb{N}\dsep
  \RecFun{m_0}{f}(k)=\RecFun{m_0}{f}(0)\vdash k=0
}{%
  \DeltaDecl{Mengen}{M\dsep k\dsep x\dsep m_0}%
  \DeltaDecl{Funktionen}{f}%
}
\begin{tabproofwide}
  \proofstepwidestar[1]{m_0\in M}{\rA}
  \proofstepwidestar[2]{f\colon M\inj M}{\rA}
  \proofstepwidestar[2]{f\colon M\to M}{\FormulaRefAuto{Injektive Funktion}{2}}
  \proofstepwidestar[4]{\forall x\in M\,\bigl(f(x)\neq m_0\bigr)}{\rA}
  \proofstepwidestar[5]{k\in\mathbb{N}}{\rA}
  \proofstepwidestar[6]{\RecFun{m_0}{f}(k)=\RecFun{m_0}{f}(0)}{\rA}

  \proofstepwidestar[7]{k\neq0}{\rA}

  \proofstepwidestar[5]{\Pred(k)\in\mathbb{N}}{%
    \FormulaRefAuto{PeanoPredNatural}{5}}

  \proofstepwidestar[1,2,5]{\RecFun{m_0}{f}(\Pred(k))\in M}{%
    \FormulaRefAuto{m_0\in M\dsep f\colon M\to M\dsep n\in\mathbb{N}\vdash
    \RecFun{m_0}{f}(n)\in M}{1,3,8}}

  \proofstepwidestar[1,2,4,5]{f(\RecFun{m_0}{f}(\Pred(k)))\neq m_0}{%
    \rRE{\rUE{4},9}}

  \proofstepwide[1,2,5,7]{f(\RecFun{m_0}{f}(\Pred(k)))}{=}{\RecFun{m_0}{f}(k)}{%
    \FormulaRefAuto{%
      m_0\in M\dsep f\colon M\to M\dsep n\in\mathbb{N}\dsep n\neq0\vdash
      f(\RecFun{m_0}{f}(\Pred(n)))=\RecFun{m_0}{f}(n)
    }{1,3,5,7}}
  \proofstepwide[6]{}{=}{\RecFun{m_0}{f}(0)}{\rA}
  \proofstepwide[1,2]{}{=}{m_0}{%
    \FormulaRefAuto{m_0\in M\dsep f\colon M\to M \vdash
    \RecFun{m_0}{f}(0)=m_0}{1,3}}
  \proofstepwide[1,2,5,6,7]{f(\RecFun{m_0}{f}(\Pred(k)))}{=}{m_0}{\rChain{11,12,13}}

  \proofstepwidestar[1,2,4,5,6,7]{\bot}{\rBI{10,14}}

  \proofstepwidestar[1,2,4,5,6]{k=0}{\rCE{7,15}}
\end{tabproofwide}


\FormulaThmDeltaR[Hilfssatz zum Induktionsschritt der Eindeutigkeit der Rekursionsabbildung]{%
  \begin{aligned}[t]
    &m_0\in M\dsep f\colon M\inj M\dsep{}\\
    &\forall x\in M\,\bigl(f(x)\neq m_0\bigr)\dsep{}\\
    &u\in\mathbb{N}\dsep{}\\
    &\forall x\in\mathbb{N}\,\Bigl(
      \RecFun{m_0}{f}(x)=\RecFun{m_0}{f}(u)\rightarrow x=u
    \Bigr)\dsep{}\\
    &k\in\mathbb{N}\dsep
    \RecFun{m_0}{f}(k)=\RecFun{m_0}{f}(\Succ(u))\vdash k=\Succ(u)
  \end{aligned}
}{%
  m_0\in M\dsep f\colon M\inj M\dsep
  \forall x\in M\,\bigl(f(x)\neq m_0\bigr)\dsep
  u\in\mathbb{N}\dsep
  \forall x\in\mathbb{N}\,\Bigl(
    \RecFun{m_0}{f}(x)=\RecFun{m_0}{f}(u)\rightarrow x=u
  \Bigr)\dsep
  k\in\mathbb{N}\dsep
  \RecFun{m_0}{f}(k)=\RecFun{m_0}{f}(\Succ(u))\vdash k=\Succ(u)
}{%
  \DeltaDecl{Mengen}{M\dsep u\dsep k\dsep x\dsep m_0}%
  \DeltaDecl{Funktionen}{f}%
}
\begin{tabproofwide}
  \proofstepwidestar[1]{m_0\in M}{\rA}
  \proofstepwidestar[2]{f\colon M\inj M}{\rA}
  \proofstepwidestar[2]{f\colon M\to M}{\FormulaRefAuto{Injektive Funktion}{2}}
  \proofstepwidestar[4]{\forall x\in M\,\bigl(f(x)\neq m_0\bigr)}{\rA}
  \proofstepwidestar[5]{u\in\mathbb{N}}{\rA}
    \proofstepwidestar[6]{%
      \begin{aligned}[t]
        &\forall x\in\mathbb{N}\,\Bigl(
          \RecFun{m_0}{f}(x)=\RecFun{m_0}{f}(u)\\
        &\rightarrow x=u
        \Bigr)
      \end{aligned}
    }{\rA}
  \proofstepwidestar[7]{k\in\mathbb{N}}{\rA}
  \proofstepwidestar[8]{\RecFun{m_0}{f}(k)=\RecFun{m_0}{f}(\Succ(u))}{\rA}

  \proofstepwidestar[1,2,4,5,7,8]{k\neq0}{%
    \FormulaRefAuto{%
      m_0\in M\dsep f\colon M\inj M\dsep
      \forall x\in M\,\bigl(f(x)\neq m_0\bigr)\dsep
      u\in\mathbb{N}\dsep k\in\mathbb{N}\dsep
      \RecFun{m_0}{f}(k)=\RecFun{m_0}{f}(\Succ(u))\vdash k\neq0
    }{1,2,4,5,7,8}}

  \proofstepwidestar[7]{\Pred(k)\in\mathbb{N}}{%
    \FormulaRefAuto{PeanoPredNatural}{7}}

  \proofstepwide[1,2,4,5,7,8]{f(\RecFun{m_0}{f}(\Pred(k)))}{=}{\RecFun{m_0}{f}(k)}{%
    \FormulaRefAuto{%
      m_0\in M\dsep f\colon M\to M\dsep n\in\mathbb{N}\dsep n\neq0\vdash
      f(\RecFun{m_0}{f}(\Pred(n)))=\RecFun{m_0}{f}(n)
    }{1,3,7,9}}
  \proofstepwide[8]{}{=}{\RecFun{m_0}{f}(\Succ(u))}{\FormulaRefAuto{P\vdash P}{8}}
  \proofstepwide[1,2,5]{}{=}{f(\RecFun{m_0}{f}(u))}{%
    \FormulaRefAuto{%
      m_0\in M\dsep f\colon M\to M\dsep n\in\mathbb{N}\vdash
      \RecFun{m_0}{f}(\Succ(n))=f(\RecFun{m_0}{f}(n))
    }{1,3,5}}
  \proofstepwide[1,2,4,5,7,8]{f(\RecFun{m_0}{f}(\Pred(k)))}{=}{f(\RecFun{m_0}{f}(u))}{\rChain{11,12,13}}

  \proofstepwidestar[1,2,4,5,7,8]{%
    \RecFun{m_0}{f}(\Pred(k))=\RecFun{m_0}{f}(u)
  }{%
    \FormulaRefAuto{Injektive Funktion}{%
      2,
      \FormulaRefAuto{m_0\in M\dsep f\colon M\to M\dsep n\in\mathbb{N}\vdash
        \RecFun{m_0}{f}(n)\in M}{1,3,10},
      \FormulaRefAuto{m_0\in M\dsep f\colon M\to M\dsep n\in\mathbb{N}\vdash
        \RecFun{m_0}{f}(n)\in M}{1,3,5},14}}

    \proofstepwidestar[6,10]{%
      \begin{aligned}[t]
        &\RecFun{m_0}{f}(\Pred(k))=\RecFun{m_0}{f}(u)\\
        &\rightarrow \Pred(k)=u
      \end{aligned}
    }{\rRE{\rUE{6},10}}

  \proofstepwidestar[1,2,4,5,6,7,8]{\Pred(k)=u}{\rRE{15,16}}

  \proofstepwidestar[1,2,4,5,6,7,8]{k=\Succ(u)}{%
    \FormulaRefAuto{PeanoPredRecoverSucc}{7,9,17}}
\end{tabproofwide}


\FormulaThmDeltaR[Eindeutigkeit der Werte der Rekursionsabbildung]{%
  \begin{aligned}[t]
    &m_0\in M\dsep f\colon M\inj M\dsep{}\\
    &\forall x\in M\,\bigl(f(x)\neq m_0\bigr)\dsep{}\\
    &n\in\mathbb{N}\dsep k\in\mathbb{N}\dsep{}\\
    &\RecFun{m_0}{f}(k)=\RecFun{m_0}{f}(n)\vdash k=n
  \end{aligned}
}{%
  m_0\in M\dsep f\colon M\inj M\dsep
  \forall x\in M\,\bigl(f(x)\neq m_0\bigr)\dsep
  n\in\mathbb{N}\dsep k\in\mathbb{N}\dsep
  \RecFun{m_0}{f}(k)=\RecFun{m_0}{f}(n)\vdash k=n
}{%
  \DeltaDecl{Mengen}{M\dsep n\dsep k\dsep x\dsep u\dsep m_0}%
  \DeltaDecl{Funktionen}{f}%
}
\begin{tabproofsplitwide}
  \proofpartwideindR[Induktionsanfang]{%
    \begin{aligned}[t]
      &m_0\in M\dsep f\colon M\inj M\dsep{}\\
      &\forall x\in M\,\bigl(f(x)\neq m_0\bigr)\vdash{}\\
      &\forall k\in\mathbb{N}\,\Bigl(
        \RecFun{m_0}{f}(k)=\RecFun{m_0}{f}(0)\rightarrow
        k=0
      \Bigr)
    \end{aligned}
  }{%
    m_0\in M\dsep f\colon M\inj M\dsep
    \forall x\in M\,\bigl(f(x)\neq m_0\bigr)\vdash
    \forall k\in\mathbb{N}\,\Bigl(
      \RecFun{m_0}{f}(k)=\RecFun{m_0}{f}(0)\rightarrow
      k=0
    \Bigr)
  }
    \proofstepwidestar[1]{m_0\in M}{\rA}
    \proofstepwidestar[2]{f\colon M\inj M}{\rA}
    \proofstepwidestar[3]{\forall x\in M\,\bigl(f(x)\neq m_0\bigr)}{\rA}

    \proofstepwidestar[4]{k\in\mathbb{N}}{\rA}
    \proofstepwidestar[5]{\RecFun{m_0}{f}(k)=\RecFun{m_0}{f}(0)}{\rA}

    \proofstepwidestar[1,2,3,4,5]{k=0}{%
      \FormulaRefAuto{%
        m_0\in M\dsep f\colon M\inj M\dsep
        \forall x\in M\,\bigl(f(x)\neq m_0\bigr)\dsep
        k\in\mathbb{N}\dsep
        \RecFun{m_0}{f}(k)=\RecFun{m_0}{f}(0)\vdash k=0
      }{1,2,3,4,5}}

    \proofstepwidestar[1,2,3,4]{%
      \RecFun{m_0}{f}(k)=\RecFun{m_0}{f}(0)\rightarrow
      k=0
    }{\rRI{5,6}}

    \proofstepwidestar[1,2,3]{%
      \forall k\in\mathbb{N}\,\Bigl(
        \RecFun{m_0}{f}(k)=\RecFun{m_0}{f}(0)\rightarrow
        k=0
      \Bigr)
    }{\rUI{\rRI{4,7}}}
  \closeproofpartwideindR

\proofpartwideindR[Induktionsschritt]{%
  \begin{aligned}[t]
    &m_0\in M\dsep f\colon M\inj M\dsep{}\\
    &\forall x\in M\,\bigl(f(x)\neq m_0\bigr)\dsep{}\\
    &u\in\mathbb{N}\dsep{}\\
    &\forall x\in\mathbb{N}\,\Bigl(
      \RecFun{m_0}{f}(x)=\RecFun{m_0}{f}(u)\\
    &\rightarrow x=u
    \Bigr)\vdash{}\\
    &\forall k\in\mathbb{N}\,\Bigl(
      \RecFun{m_0}{f}(k)=\RecFun{m_0}{f}(\Succ(u))\\
    &\rightarrow k=\Succ(u)
    \Bigr)
  \end{aligned}
}{%
  m_0\in M\dsep f\colon M\inj M\dsep
  \forall x\in M\,\bigl(f(x)\neq m_0\bigr)\dsep
  u\in\mathbb{N}\dsep
  \forall x\in\mathbb{N}\,\Bigl(
    \RecFun{m_0}{f}(x)=\RecFun{m_0}{f}(u)\rightarrow x=u
  \Bigr)\vdash
  \forall k\in\mathbb{N}\,\Bigl(
    \RecFun{m_0}{f}(k)=\RecFun{m_0}{f}(\Succ(u))\rightarrow
    k=\Succ(u)
  \Bigr)
}
  \proofstepwidestar[1]{m_0\in M}{\rA}
  \proofstepwidestar[2]{f\colon M\inj M}{\rA}
  \proofstepwidestar[3]{\forall x\in M\,\bigl(f(x)\neq m_0\bigr)}{\rA}
  \proofstepwidestar[4]{u\in\mathbb{N}}{\rA}
  \proofstepwidestar[5]{%
    \begin{aligned}[t]
      &\forall x\in\mathbb{N}\,\Bigl(
        \RecFun{m_0}{f}(x)=\RecFun{m_0}{f}(u)\\
      &\rightarrow x=u
      \Bigr)
    \end{aligned}
  }{\rA}

  \proofstepwidestar[6]{k\in\mathbb{N}}{\rA}
  \proofstepwidestar[7]{%
    \RecFun{m_0}{f}(k)=\RecFun{m_0}{f}(\Succ(u))
  }{\rA}

  \proofstepwidestar[1,2,3,4,5,6,7]{k=\Succ(u)}{%
    \FormulaRefAuto{%
      m_0\in M\dsep f\colon M\inj M\dsep
      \forall x\in M\,\bigl(f(x)\neq m_0\bigr)\dsep
      u\in\mathbb{N}\dsep
      \forall x\in\mathbb{N}\,\Bigl(
        \RecFun{m_0}{f}(x)=\RecFun{m_0}{f}(u)\rightarrow x=u
      \Bigr)\dsep
      k\in\mathbb{N}\dsep
      \RecFun{m_0}{f}(k)=\RecFun{m_0}{f}(\Succ(u))\vdash k=\Succ(u)
    }{1,2,3,4,5,6,7}}

  \proofstepwidestar[1,2,3,4,5,6]{%
    \begin{aligned}[t]
      &\RecFun{m_0}{f}(k)=\RecFun{m_0}{f}(\Succ(u))\\
      &\rightarrow k=\Succ(u)
    \end{aligned}
  }{\rRI{7,8}}

  \proofstepwidestar[1,2,3,4,5]{%
    \begin{aligned}[t]
      &\forall k\in\mathbb{N}\,\Bigl(
        \RecFun{m_0}{f}(k)=\RecFun{m_0}{f}(\Succ(u))\\
      &\rightarrow k=\Succ(u)
      \Bigr)
    \end{aligned}
  }{\rUI{\rRI{6,9}}}
\closeproofpartwideindR

\proofpartwide{Schluss}
  \proofstepwidestar[1]{m_0\in M}{\rA}
  \proofstepwidestar[2]{f\colon M\inj M}{\rA}
  \proofstepwidestar[3]{\forall x\in M\,\bigl(f(x)\neq m_0\bigr)}{\rA}
  \proofstepwidestar[4]{n\in\mathbb{N}}{\rA}
  \proofstepwidestar[5]{k\in\mathbb{N}}{\rA}
  \proofstepwidestar[6]{\RecFun{m_0}{f}(k)=\RecFun{m_0}{f}(n)}{\rA}

  \proofstepwidestar[1,2,3,4]{%
    \begin{aligned}[t]
      &\forall k\in\mathbb{N}\,\Bigl(
        \RecFun{m_0}{f}(k)=\RecFun{m_0}{f}(n)\\
      &\rightarrow k=n
      \Bigr)
    \end{aligned}
  }{%
    \makecell[l]{%
      \FormulaRefAuto{PeanoInductionRule}(%
      \\ \quad \FormulaRefAuto{%
        m_0\in M\dsep f\colon M\inj M\dsep
        \forall x\in M\,\bigl(f(x)\neq m_0\bigr)\vdash
        \forall k\in\mathbb{N}\,\Bigl(
          \RecFun{m_0}{f}(k)=\RecFun{m_0}{f}(0)\rightarrow
          k=0
        \Bigr)
      }{1,2,3},%
      \\ \quad \FormulaRefAuto{%
        m_0\in M\dsep f\colon M\inj M\dsep
        \forall x\in M\,\bigl(f(x)\neq m_0\bigr)\dsep
        u\in\mathbb{N}\dsep
        \forall x\in\mathbb{N}\,\Bigl(
          \RecFun{m_0}{f}(x)=\RecFun{m_0}{f}(u)\rightarrow x=u
        \Bigr)\vdash
        \forall k\in\mathbb{N}\,\Bigl(
          \RecFun{m_0}{f}(k)=\RecFun{m_0}{f}(\Succ(u))\rightarrow
          k=\Succ(u)
        \Bigr)
      }{1,2,3},%
      \\ \quad 4%
      \\ )%
    }%
  }

  \proofstepwidestar[1,2,3,4,5]{%
    \begin{aligned}[t]
      &\RecFun{m_0}{f}(k)=\RecFun{m_0}{f}(n)\\
      &\rightarrow k=n
    \end{aligned}
  }{\rRE{\rUE{7},5}}

  \proofstepwidestar[1,2,3,4,5,6]{k=n}{\rRE{6,8}}
\closeproofpartwide
\end{tabproofsplitwide}

\FormulaThmDelta[Die Rekursionsabbildung ist injektiv]{%
  m_0\in M\dsep f\colon M\inj M\dsep
  \forall x\in M\,\bigl(f(x)\neq m_0\bigr)\vdash
  \RecFun{m_0}{f}\colon \mathbb{N}\inj M
}{%
  \DeltaDecl{Mengen}{M\dsep x\dsep y\dsep z\dsep n\dsep k\dsep m_0}%
  \DeltaDecl{Funktionen}{f}%
}
\begin{tabproof}
  \proofstep{1}{m_0\in M}{\rA}
  \proofstep{2}{f\colon M\inj M}{\rA}
  \proofstep{3}{\forall x\in M\,\bigl(f(x)\neq m_0\bigr)}{\rA}

  \proofstep{1,2}{\RecFun{m_0}{f}\colon \mathbb{N}\to M}{%
    \FormulaRefAuto{%
      m_0\in M\dsep f\colon M\to M \vdash
      \RecFun{m_0}{f}\colon \mathbb{N}\to M
    }{1,2}}

  \proofstep{1,2,3}{%
    \begin{aligned}[t]
      \forall x\in\mathbb{N}\,\forall y\in\mathbb{N}\,\bigl(
      &\RecFun{m_0}{f}(x)=\RecFun{m_0}{f}(y)\\
      &\rightarrow x=y
      \bigr)
    \end{aligned}
  }{%
    \FormulaRefAuto{%
      m_0\in M\dsep f\colon M\inj M\dsep
      \forall z\in M\,\bigl(f(z)\neq m_0\bigr)\dsep
      n\in\mathbb{N}\dsep k\in\mathbb{N}\dsep
      \RecFun{m_0}{f}(k)=\RecFun{m_0}{f}(n)\vdash k=n
    }{1,2,3}}

  \proofstep{1,2,3}{\RecFun{m_0}{f}\colon \mathbb{N}\inj M}{%
    \FormulaRefAuto{Injektive Funktion}{4,5}}
\end{tabproof}
\subsection{Anwendung: Folgenverschiebung entlang einer Einbettung}

Im folgenden Konstruktionsschritt bezeichnet \(\tSuccShift{H}\) nur die
punktweise Funktionsvorschrift; die eigentliche Folgenverschiebung ist
\(\SuccShift{H}\).

\subsubsection{Funktionsvorschrift und Typisierung}

\paragraph{Definition der Vorschrift}

\FormulaDefDeltaR[Funktionsvorschrift der Folgenverschiebung]{%
  \begin{aligned}[t]
    &H\colon \mathbb{N}\inj M\dsep x\in M \vdash{}\\
    &\tSuccShift{H}(x)\coloneqq
      \IfThenElse{x\in H[\mathbb{N}]}
        {H\bigl(\Succ(H^{-1}(x))\bigr)}
        {x}
  \end{aligned}
}{%
  H\colon \mathbb{N}\inj M\dsep x\in M \vdash
  \tSuccShift{H}(x)\coloneqq
  \IfThenElse{x\in H[\mathbb{N}]}
    {H\bigl(\Succ(H^{-1}(x))\bigr)}
    {x}
}{%
  \DeltaDecl{Mengen}{M\dsep x}%
  \DeltaDecl{Funktionen}{H}%
  \DeltaDecl{Funktionensymbol}{\tSuccShift{H}}%
}

\paragraph{Auswertung der Vorschrift}

\FormulaThmDelta[Funktionsvorschrift der Folgenverschiebung im Bildfall]{%
  H\colon \mathbb{N}\inj M\dsep x\in M\dsep x\in H[\mathbb{N}]
  \vdash \tSuccShift{H}(x)=H\bigl(\Succ(H^{-1}(x))\bigr)
}{%
  \DeltaDecl{Mengen}{M\dsep x}%
  \DeltaDecl{Funktionen}{H}%
  \DeltaDecl{Funktionensymbol}{\tSuccShift{H}}%
}
\begin{tabproof}
  \proofstep{1}{H\colon \mathbb{N}\inj M}{\rA}
  \proofstep{2}{x\in M}{\rA}
  \proofstep{3}{x\in H[\mathbb{N}]}{\rA}
  \proofstep{1,2}{%
    \begin{aligned}[t]
      \tSuccShift{H}(x)
      &=\IfThenElse{x\in H[\mathbb{N}]}
        {H\bigl(\Succ(H^{-1}(x))\bigr)}{x}
    \end{aligned}
  }{%
    \FormulaRefAuto{%
      H\colon \mathbb{N}\inj M\dsep x\in M \vdash
      \tSuccShift{H}(x)\coloneqq
      \IfThenElse{x\in H[\mathbb{N}]}
        {H\bigl(\Succ(H^{-1}(x))\bigr)}{x}}{1,2}}
  \proofstep{3}{%
    \begin{aligned}[t]
      &\IfThenElse{x\in H[\mathbb{N}]}
        {H\bigl(\Succ(H^{-1}(x))\bigr)}{x}\\
      &\qquad=H\bigl(\Succ(H^{-1}(x))\bigr)
    \end{aligned}
  }{\FormulaRefAuto{P \vdash \IfThenElse{P}{a}{b}=a}{3}}
  \proofstep{1,2,3}{%
    \begin{aligned}[t]
      \tSuccShift{H}(x)
      &=H\bigl(\Succ(H^{-1}(x))\bigr)
    \end{aligned}
  }{\FormulaRefAuto{a = b,\, b = c \vdash a = c}{4,5}}
\end{tabproof}

\FormulaThmDelta[Funktionsvorschrift der Folgenverschiebung auf Folgengliedern]{%
  H\colon \mathbb{N}\inj M\dsep n\in\mathbb{N}\vdash
  \tSuccShift{H}(H(n))=H(\Succ(n))
}{%
  \DeltaDecl{Mengen}{M\dsep n}%
  \DeltaDecl{Funktionen}{H}%
  \DeltaDecl{Funktionensymbol}{\tSuccShift{H}}%
}
\begin{tabproof}
  \proofstep{1}{H\colon \mathbb{N}\inj M}{\rA}
  \proofstep{2}{n\in\mathbb{N}}{\rA}
  \proofstep{1,2}{H(n)\in M}{%
    \FormulaRefAuto{F\colon A\inj B\dsep x\in A\vdash F(x)\in B}{1,2}}
  \proofstep{1,2}{H(n)\in H[\mathbb{N}]}{%
    \FormulaRefAuto{C\subseteq A\dsep F\colon A\inj B\dsep x\in C\vdash F(x)\in F[C]}{%
      \FormulaRefAuto{A\subseteq A},1,2}}
  \proofstep{1,2}{%
    \tSuccShift{H}(H(n))=H\bigl(\Succ(H^{-1}(H(n)))\bigr)
  }{%
    \FormulaRefAuto{%
      H\colon \mathbb{N}\inj M\dsep x\in M\dsep x\in H[\mathbb{N}]
      \vdash \tSuccShift{H}(x)=H\bigl(\Succ(H^{-1}(x))\bigr)}{1,3,4}}
  \proofstep{1,2}{H^{-1}(H(n))=n}{%
    \FormulaRefAuto{x\in A \vdash F^{-1}(F(x)) = x}{%
      \FormulaRefAuto{F\colon A\inj B \vdash F\colon A\bij F[A]}{1},2}}
  \proofstep{1,2}{H\bigl(\Succ(H^{-1}(H(n)))\bigr)=H(\Succ(n))}{%
    \rIE{6,\rII}}
  \proofstep{1,2}{\tSuccShift{H}(H(n))=H(\Succ(n))}{%
    \FormulaRefAuto{a = b,\, b = c \vdash a = c}{5,7}}
\end{tabproof}

\FormulaThmDelta[Funktionsvorschrift der Folgenverschiebung im Nicht-Bildfall]{%
  H\colon \mathbb{N}\inj M\dsep x\in M\dsep x\notin H[\mathbb{N}]
  \vdash \tSuccShift{H}(x)=x
}{%
  \DeltaDecl{Mengen}{M\dsep x}%
  \DeltaDecl{Funktionen}{H}%
  \DeltaDecl{Funktionensymbol}{\tSuccShift{H}}%
}
\begin{tabproof}
  \proofstep{1}{H\colon \mathbb{N}\inj M}{\rA}
  \proofstep{2}{x\in M}{\rA}
  \proofstep{3}{x\notin H[\mathbb{N}]}{\rA}
  \proofstep{1,2,3}{\tSuccShift{H}(x)=x}{%
    \FormulaRefAuto{a = b,\, b = c \vdash a = c}{%
      \FormulaRefAuto{%
        H\colon \mathbb{N}\inj M\dsep x\in M \vdash
        \tSuccShift{H}(x)\coloneqq
        \IfThenElse{x\in H[\mathbb{N}]}
          {H\bigl(\Succ(H^{-1}(x))\bigr)}{x}}{1,2},
      \FormulaRefAuto{\neg P \vdash \IfThenElse{P}{a}{b}=b}{3}}}
\end{tabproof}

\paragraph{Typisierung der Vorschrift}

\FormulaThmDelta[Typisierung der Funktionsvorschrift der Folgenverschiebung]{%
  H\colon \mathbb{N}\inj M\dsep x\in M\vdash \tSuccShift{H}(x)\in M
}{%
  \DeltaDecl{Mengen}{M\dsep x}%
  \DeltaDecl{Funktionen}{H}%
  \DeltaDecl{Funktionensymbol}{\tSuccShift{H}}%
}
\begin{tabproof}
  \proofstep{1}{H\colon \mathbb{N}\inj M}{\rA}
  \proofstep{2}{x\in M}{\rA}
  \proofstep{}{x\in H[\mathbb{N}]\lor x\notin H[\mathbb{N}]}{%
    \FormulaRefAuto{P\lor \neg P}}

  \proofstep{4}{x\in H[\mathbb{N}]}{\rA}
  \proofstep{1,4}{H^{-1}(x)\in\mathbb{N}}{%
    \FormulaRefAuto{H\colon \mathbb{N}\inj M\dsep x\in H[\mathbb{N}]
      \vdash H^{-1}(x)\in\mathbb{N}}{1,4}}
  \proofstep{1,4}{\Succ(H^{-1}(x))\in\mathbb{N}}{%
    \FormulaRefAuto{PeanoSuccClosure}{5}}
  \proofstep{1}{H\colon \mathbb{N}\to M}{%
    \FormulaRefAuto{Injektive Funktion}{1}}
  \proofstep{1,4}{H(\Succ(H^{-1}(x)))\in M}{%
    \FormulaRefAuto{F\colon A\to B\dsep x\in A\vdash F(x)\in B}{7,6}}
  \proofstep{1,2,4}{%
    \tSuccShift{H}(x)=H\bigl(\Succ(H^{-1}(x))\bigr)
  }{%
    \FormulaRefAuto{%
      H\colon \mathbb{N}\inj M\dsep x\in M\dsep x\in H[\mathbb{N}]
      \vdash \tSuccShift{H}(x)=H\bigl(\Succ(H^{-1}(x))\bigr)}{1,2,4}}
  \proofstep{1,2,4}{\tSuccShift{H}(x)\in M}{\rIE{9,8}}

  \proofstep{11}{x\notin H[\mathbb{N}]}{\rA}
  \proofstep{1,2,11}{\tSuccShift{H}(x)=x}{%
    \FormulaRefAuto{%
      H\colon \mathbb{N}\inj M\dsep x\in M\dsep x\notin H[\mathbb{N}]
      \vdash \tSuccShift{H}(x)=x}{1,2,11}}
  \proofstep{1,2,11}{\tSuccShift{H}(x)\in M}{\rIE{12,2}}

  \proofstep{1,2}{\tSuccShift{H}(x)\in M}{\rOE{3,4,10,11,13}}
\end{tabproof}

\subsubsection{Existenz und Definition der Folgenverschiebung}

\paragraph{Existenz als Funktion}

\FormulaThmDelta[Folgenverschiebung als Funktion]{%
  H\colon \mathbb{N}\inj M\vdash
  \exists! S\colon M\to M\,\forall x\in M\,S(x)=\tSuccShift{H}(x)
}{%
  \DeltaDecl{Mengen}{M\dsep x}%
  \DeltaDecl{Funktionen}{H\dsep S}%
  \DeltaDecl{Funktionensymbol}{\tSuccShift{H}}%
}
\begin{tabproof}
  \proofstep{1}{H\colon \mathbb{N}\inj M}{\rA}
  \proofstep{1}{%
    \exists! S\colon M\to M\,\forall x\in M\,S(x)=\tSuccShift{H}(x)
  }{%
    \FormulaRefAuto{%
      x\in A\vdash t(x)\coloneq y\dsep x\in A\vdash t(x)\in B
      \exists! F\colon A\to B\,\forall x\in A\,F(x) = t(x)}{%
      \FormulaRefAuto{H\colon \mathbb{N}\inj M\dsep x\in M\vdash
        \tSuccShift{H}(x)\in M}{1}}}
\end{tabproof}

\paragraph{Definition und Selbstabbildung}

\FormulaDefDeltaR[Folgenverschiebung entlang einer Einbettung]{%
  \begin{aligned}[t]
    &H\colon \mathbb{N}\inj M \vdash{}\\
    &\SuccShift{H}\coloneqq
    \iota S\Bigl(
      S\colon M\to M \land{}\\
    &\qquad \forall x\in M\,S(x)=\tSuccShift{H}(x)
    \Bigr)
  \end{aligned}
}{%
  H\colon \mathbb{N}\inj M \vdash
  \SuccShift{H}\coloneqq
  \iota S\Bigl(
    S\colon M\to M \land
    \forall x\in M\,S(x)=\tSuccShift{H}(x)
  \Bigr)
}{%
  \DeltaDecl{Mengen}{M\dsep x}%
  \DeltaDecl{Funktionen}{H\dsep S}%
  \DeltaDecl{Funktionensymbol}{\tSuccShift{H}\dsep \SuccShift{H}}%
}

\FormulaThmDelta[Die Folgenverschiebung ist eine Selbstabbildung]{%
  H\colon \mathbb{N}\inj M \vdash \SuccShift{H}\colon M\to M
}{%
  \DeltaDecl{Mengen}{M}%
  \DeltaDecl{Funktionen}{H}%
  \DeltaDecl{Funktionensymbol}{\SuccShift{H}}%
}
\begin{tabproof}
  \proofstep{1}{H\colon \mathbb{N}\inj M}{\rA}
  \proofstep{1}{\SuccShift{H}\colon M\to M}{%
    \rAEa{\FormulaRefAuto{%
      H\colon \mathbb{N}\inj M \vdash
      \SuccShift{H}\coloneqq
      \iota S\Bigl(
        S\colon M\to M \land
        \forall x\in M\,S(x)=\tSuccShift{H}(x)
      \Bigr)}{1}}}
\end{tabproof}

\subsubsection{Grundgleichungen der Folgenverschiebung}

\paragraph{Zusammenhang mit der Vorschrift}

\FormulaThmDelta[Grundgleichung zur Funktionsvorschrift]{%
  H\colon \mathbb{N}\inj M\dsep x\in M\vdash
  \SuccShift{H}(x)=\tSuccShift{H}(x)
}{%
  \DeltaDecl{Mengen}{M\dsep x}%
  \DeltaDecl{Funktionen}{H}%
  \DeltaDecl{Funktionensymbol}{\tSuccShift{H}\dsep \SuccShift{H}}%
}
\begin{tabproof}
  \proofstep{1}{H\colon \mathbb{N}\inj M}{\rA}
  \proofstep{2}{x\in M}{\rA}
  \proofstep{1}{%
    \forall x\in M\,\SuccShift{H}(x)=\tSuccShift{H}(x)
  }{%
    \rAEb{\FormulaRefAuto{%
      H\colon \mathbb{N}\inj M \vdash
      \SuccShift{H}\coloneqq
      \iota S\Bigl(
        S\colon M\to M \land
        \forall x\in M\,S(x)=\tSuccShift{H}(x)
      \Bigr)}{1}}}
  \proofstep{1,2}{\SuccShift{H}(x)=\tSuccShift{H}(x)}{%
    \rRE{\rUE{3},2}}
\end{tabproof}

\paragraph{Auswertung auf Bild und Komplement}

\FormulaThmDelta[Grundgleichung der Folgenverschiebung]{%
  H\colon \mathbb{N}\inj M\dsep n\in\mathbb{N}\vdash
  \SuccShift{H}(H(n))=H(\Succ(n))
}{%
  \DeltaDecl{Mengen}{M\dsep n}%
  \DeltaDecl{Funktionen}{H}%
  \DeltaDecl{Funktionensymbol}{\SuccShift{H}}%
}
\begin{tabproof}
  \proofstep{1}{H\colon \mathbb{N}\inj M}{\rA}
  \proofstep{2}{n\in\mathbb{N}}{\rA}
  \proofstep{1}{H\colon \mathbb{N}\to M}{%
    \FormulaRefAuto{Injektive Funktion}{1}}
  \proofstep{1,2}{H(n)\in M}{%
    \FormulaRefAuto{F\colon A\to B\dsep x\in A\vdash F(x)\in B}{3,2}}
  \proofstep{1,2}{\SuccShift{H}(H(n))=\tSuccShift{H}(H(n))}{%
    \FormulaRefAuto{%
      H\colon \mathbb{N}\inj M\dsep x\in M\vdash
      \SuccShift{H}(x)=\tSuccShift{H}(x)}{1,4}}
  \proofstep{1,2}{\tSuccShift{H}(H(n))=H(\Succ(n))}{%
    \FormulaRefAuto{%
      H\colon \mathbb{N}\inj M\dsep n\in\mathbb{N}\vdash
      \tSuccShift{H}(H(n))=H(\Succ(n))}{1,2}}
  \proofstep{1,2}{\SuccShift{H}(H(n))=H(\Succ(n))}{%
    \FormulaRefAuto{a = b,\, b = c \vdash a = c}{5,6}}
\end{tabproof}

\FormulaThmDelta[Fixpunkte der Folgenverschiebung außerhalb des Bildes]{%
  H\colon \mathbb{N}\inj M\dsep x\in M\dsep x\notin H[\mathbb{N}]
  \vdash \SuccShift{H}(x)=x
}{%
  \DeltaDecl{Mengen}{M\dsep x}%
  \DeltaDecl{Funktionen}{H}%
  \DeltaDecl{Funktionensymbol}{\SuccShift{H}}%
}
\begin{tabproof}
  \proofstep{1}{H\colon \mathbb{N}\inj M}{\rA}
  \proofstep{2}{x\in M}{\rA}
  \proofstep{3}{x\notin H[\mathbb{N}]}{\rA}
  \proofstep{1,2}{\SuccShift{H}(x)=\tSuccShift{H}(x)}{%
    \FormulaRefAuto{%
      H\colon \mathbb{N}\inj M\dsep x\in M\vdash
      \SuccShift{H}(x)=\tSuccShift{H}(x)}{1,2}}
  \proofstep{1,2,3}{\tSuccShift{H}(x)=x}{%
    \FormulaRefAuto{%
      H\colon \mathbb{N}\inj M\dsep x\in M\dsep x\notin H[\mathbb{N}]
      \vdash \tSuccShift{H}(x)=x}{1,2,3}}
  \proofstep{1,2,3}{\SuccShift{H}(x)=x}{%
    \FormulaRefAuto{a = b,\, b = c \vdash a = c}{4,5}}
\end{tabproof}

\subsubsection{Injektivität der Folgenverschiebung}

\paragraph{Punktweise Injektivität}

\FormulaThmDelta[Punktweise Injektivität der Folgenverschiebung]{%
  H\colon \mathbb{N}\inj M\dsep x\in M\dsep y\in M
  \dsep \SuccShift{H}(x)=\SuccShift{H}(y)\vdash x=y
}{%
  \DeltaDecl{Mengen}{M\dsep x\dsep y\dsep n\dsep k}%
  \DeltaDecl{Funktionen}{H}%
  \DeltaDecl{Funktionensymbol}{\SuccShift{H}}%
}
\begin{tabproofsplitwide}
  \proofpartwideindR[Außen-Bild-Widerspruch]{%
    \begin{aligned}[t]
      &H\colon \mathbb{N}\inj M\dsep x\in M\dsep x\notin H[\mathbb{N}]
      \dsep n\in\mathbb{N}\dsep{}\\
      &\SuccShift{H}(x)=\SuccShift{H}(H(n))\vdash \bot
    \end{aligned}
  }{%
    H\colon \mathbb{N}\inj M\dsep x\in M\dsep x\notin H[\mathbb{N}]
    \dsep n\in\mathbb{N}\dsep
    \SuccShift{H}(x)=\SuccShift{H}(H(n))\vdash \bot
  }
    \proofstepwidestar[1]{H\colon \mathbb{N}\inj M}{\rA}
    \proofstepwidestar[2]{x\in M}{\rA}
    \proofstepwidestar[3]{x\notin H[\mathbb{N}]}{\rA}
    \proofstepwidestar[4]{n\in\mathbb{N}}{\rA}
    \proofstepwidestar[5]{\SuccShift{H}(x)=\SuccShift{H}(H(n))}{\rA}
    \proofstepwidestar[1,2,3]{\SuccShift{H}(x)=x}{%
      \FormulaRefAuto{%
        H\colon \mathbb{N}\inj M\dsep x\in M\dsep x\notin H[\mathbb{N}]
        \vdash \SuccShift{H}(x)=x}{1,2,3}}
    \proofstepwidestar[1,4]{\SuccShift{H}(H(n))=H(\Succ(n))}{%
      \FormulaRefAuto{%
        H\colon \mathbb{N}\inj M\dsep n\in\mathbb{N}\vdash
        \SuccShift{H}(H(n))=H(\Succ(n))}{1,4}}
    \proofstepwidestar[1,2,3,4,5]{x=H(\Succ(n))}{%
      \FormulaRefAuto{a = b\dsep a = c\dsep b=d\vdash c = d}{5,6,7}}
    \proofstepwidestar[1,4]{H(\Succ(n))\in H[\mathbb{N}]}{%
      \FormulaRefAuto{%
        H\colon \mathbb{N}\inj M\dsep n\in\mathbb{N}
        \vdash H(\Succ(n))\in H[\mathbb{N}]}{1,4}}
    \proofstepwidestar[1,2,3,4,5]{x\in H[\mathbb{N}]}{\rIE{8,9}}
    \proofstepwidestar[1,2,3,4,5]{\bot}{\rBI{10,3}}
  \closeproofpartwideindR

  \proofpartwideindR[Bild-Außen-Widerspruch]{%
    \begin{aligned}[t]
      &H\colon \mathbb{N}\inj M\dsep n\in\mathbb{N}\dsep x\in M
      \dsep x\notin H[\mathbb{N}]\dsep{}\\
      &\SuccShift{H}(H(n))=\SuccShift{H}(x)\vdash \bot
    \end{aligned}
  }{%
    H\colon \mathbb{N}\inj M\dsep n\in\mathbb{N}\dsep x\in M
    \dsep x\notin H[\mathbb{N}]\dsep
    \SuccShift{H}(H(n))=\SuccShift{H}(x)\vdash \bot
  }
    \proofstepwidestar[1]{H\colon \mathbb{N}\inj M}{\rA}
    \proofstepwidestar[2]{n\in\mathbb{N}}{\rA}
    \proofstepwidestar[3]{x\in M}{\rA}
    \proofstepwidestar[4]{x\notin H[\mathbb{N}]}{\rA}
    \proofstepwidestar[5]{\SuccShift{H}(H(n))=\SuccShift{H}(x)}{\rA}
    \proofstepwidestar[5]{\SuccShift{H}(x)=\SuccShift{H}(H(n))}{%
      \FormulaRefAuto{a=b\vdash b=a}{5}}
    \proofstepwidestar[1,2,3,4,5]{\bot}{%
      \FormulaRefAuto{%
        H\colon \mathbb{N}\inj M\dsep x\in M\dsep x\notin H[\mathbb{N}]
        \dsep n\in\mathbb{N}\dsep
        \SuccShift{H}(x)=\SuccShift{H}(H(n))\vdash \bot}{1,3,4,2,6}}
  \closeproofpartwideindR

  \proofpartwideindR[Bild-Bild-Injektivität]{%
    \begin{aligned}[t]
      &H\colon \mathbb{N}\inj M\dsep n\in\mathbb{N}\dsep k\in\mathbb{N}
      \dsep{}\\
      &\SuccShift{H}(H(n))=\SuccShift{H}(H(k))
      \vdash H(n)=H(k)
    \end{aligned}
  }{%
    H\colon \mathbb{N}\inj M\dsep n\in\mathbb{N}\dsep k\in\mathbb{N}\dsep
    \SuccShift{H}(H(n))=\SuccShift{H}(H(k))\vdash H(n)=H(k)
  }
    \proofstepwidestar[1]{H\colon \mathbb{N}\inj M}{\rA}
    \proofstepwidestar[2]{n\in\mathbb{N}}{\rA}
    \proofstepwidestar[3]{k\in\mathbb{N}}{\rA}
    \proofstepwidestar[4]{\SuccShift{H}(H(n))=\SuccShift{H}(H(k))}{\rA}
    \proofstepwidestar[1,2]{\SuccShift{H}(H(n))=H(\Succ(n))}{%
      \FormulaRefAuto{%
        H\colon \mathbb{N}\inj M\dsep n\in\mathbb{N}\vdash
        \SuccShift{H}(H(n))=H(\Succ(n))}{1,2}}
    \proofstepwidestar[1,3]{\SuccShift{H}(H(k))=H(\Succ(k))}{%
      \FormulaRefAuto{%
        H\colon \mathbb{N}\inj M\dsep n\in\mathbb{N}\vdash
        \SuccShift{H}(H(n))=H(\Succ(n))}{1,3}}
    \proofstepwidestar[1,2,3,4]{H(\Succ(n))=H(\Succ(k))}{%
      \FormulaRefAuto{a = b\dsep a = c\dsep b=d\vdash c = d}{4,5,6}}
    \proofstepwidestar[2]{\Succ(n)\in\mathbb{N}}{%
      \FormulaRefAuto{PeanoSuccClosure}{2}}
    \proofstepwidestar[3]{\Succ(k)\in\mathbb{N}}{%
      \FormulaRefAuto{PeanoSuccClosure}{3}}
    \proofstepwidestar[1,2,3,4]{\Succ(n)=\Succ(k)}{%
      \FormulaRefAuto{F\colon A\inj B\dsep x\in A\dsep y\in A\dsep F(x)=F(y)\vdash x=y}{1,8,9,7}}
    \proofstepwidestar[1,2,3,4]{n=k}{%
      \FormulaRefAuto{%
        x\in\mathbb{N}\dsep y\in\mathbb{N}\dsep \Succ(x)=\Succ(y)\vdash x=y}{2,3,10}}
    \proofstepwidestar[1,2,3,4]{H(n)=H(k)}{\rIE{11,\rII}}
  \closeproofpartwideindR

  \proofpartwideindR[Bild-Bild-Fall]{%
    \begin{aligned}[t]
      &H\colon \mathbb{N}\inj M\dsep x\in M\dsep y\in M
      \dsep x\in H[\mathbb{N}]\dsep y\in H[\mathbb{N}]\dsep{}\\
      &\SuccShift{H}(x)=\SuccShift{H}(y)\vdash x=y
    \end{aligned}
  }{%
    H\colon \mathbb{N}\inj M\dsep x\in M\dsep y\in M
    \dsep x\in H[\mathbb{N}]\dsep y\in H[\mathbb{N}]\dsep
    \SuccShift{H}(x)=\SuccShift{H}(y)\vdash x=y
  }
    \proofstepwidestar[1]{H\colon \mathbb{N}\inj M}{\rA}
    \proofstepwidestar[2]{x\in M}{\rA}
    \proofstepwidestar[3]{y\in M}{\rA}
    \proofstepwidestar[4]{x\in H[\mathbb{N}]}{\rA}
    \proofstepwidestar[5]{y\in H[\mathbb{N}]}{\rA}
    \proofstepwidestar[6]{\SuccShift{H}(x)=\SuccShift{H}(y)}{\rA}
    \proofstepwidestar[1]{H\colon \mathbb{N}\to M}{%
      \FormulaRefAuto{Injektive Funktion}{1}}
    \proofstepwidestar[1,4]{\exists n\in\mathbb{N}\,\bigl(x=H(n)\bigr)}{%
      \FormulaRefAuto{F\colon A\to B\dsep y\in F[A]\vdash \exists x\in A\,y=F(x)}{7,4}}
    \proofstepwidestar[9]{n\in\mathbb{N}}{\rA}
    \proofstepwidestar[10]{x=H(n)}{\rA}
    \proofstepwidestar[1,5]{\exists k\in\mathbb{N}\,\bigl(y=H(k)\bigr)}{%
      \FormulaRefAuto{F\colon A\to B\dsep y\in F[A]\vdash \exists x\in A\,y=F(x)}{7,5}}
    \proofstepwidestar[12]{k\in\mathbb{N}}{\rA}
    \proofstepwidestar[13]{y=H(k)}{\rA}
    \proofstepwidestar[6,10,13]{%
      \begin{aligned}[t]
        \SuccShift{H}(H(n))
        &=\SuccShift{H}(H(k))
      \end{aligned}
    }{%
      \rIE{10,\rIE{13,6}}}
    \proofstepwidestar[1,9,12,14]{H(n)=H(k)}{%
      \FormulaRefAuto{%
        H\colon \mathbb{N}\inj M\dsep n\in\mathbb{N}\dsep k\in\mathbb{N}\dsep
        \SuccShift{H}(H(n))=\SuccShift{H}(H(k))
        \vdash H(n)=H(k)}{1,9,12,14}}
    \proofstepwidestar[1,6,9,10,12,13]{x=y}{%
      \FormulaRefAuto{a = b\dsep c = a\dsep d = b \vdash c = d}{15,10,13}}
    \proofstepwidestar[1,5,6,9,10]{x=y}{\rEE{11,12,13,16}}
    \proofstepwidestar[1,4,5,6]{x=y}{\rEE{8,9,10,17}}
  \closeproofpartwideindR

  \proofpartwideindR[Bild-Außen-Fall]{%
    \begin{aligned}[t]
      &H\colon \mathbb{N}\inj M\dsep x\in M\dsep y\in M
      \dsep x\in H[\mathbb{N}]\dsep y\notin H[\mathbb{N}]\dsep{}\\
      &\SuccShift{H}(x)=\SuccShift{H}(y)\vdash x=y
    \end{aligned}
  }{%
    H\colon \mathbb{N}\inj M\dsep x\in M\dsep y\in M
    \dsep x\in H[\mathbb{N}]\dsep y\notin H[\mathbb{N}]\dsep
    \SuccShift{H}(x)=\SuccShift{H}(y)\vdash x=y
  }
    \proofstepwidestar[1]{H\colon \mathbb{N}\inj M}{\rA}
    \proofstepwidestar[2]{x\in M}{\rA}
    \proofstepwidestar[3]{y\in M}{\rA}
    \proofstepwidestar[4]{x\in H[\mathbb{N}]}{\rA}
    \proofstepwidestar[5]{y\notin H[\mathbb{N}]}{\rA}
    \proofstepwidestar[6]{\SuccShift{H}(x)=\SuccShift{H}(y)}{\rA}
    \proofstepwidestar[1]{H\colon \mathbb{N}\to M}{%
      \FormulaRefAuto{Injektive Funktion}{1}}
    \proofstepwidestar[1,4]{\exists n\in\mathbb{N}\,\bigl(x=H(n)\bigr)}{%
      \FormulaRefAuto{F\colon A\to B\dsep y\in F[A]\vdash \exists x\in A\,y=F(x)}{7,4}}
    \proofstepwidestar[9]{n\in\mathbb{N}}{\rA}
    \proofstepwidestar[10]{x=H(n)}{\rA}
    \proofstepwidestar[6,10]{\SuccShift{H}(H(n))=\SuccShift{H}(y)}{%
      \rIE{10,6}}
    \proofstepwidestar[1,3,5,9,11]{\bot}{%
      \FormulaRefAuto{%
        H\colon \mathbb{N}\inj M\dsep n\in\mathbb{N}\dsep x\in M
        \dsep x\notin H[\mathbb{N}]\dsep
        \SuccShift{H}(H(n))=\SuccShift{H}(x)\vdash \bot}{1,9,3,5,11}}
    \proofstepwidestar[1,3,5,6,9,10]{x=y}{\rCE{12}}
    \proofstepwidestar[1,3,4,5,6]{x=y}{\rEE{8,9,10,13}}
  \closeproofpartwideindR

  \proofpartwideindR[Außen-Bild-Fall]{%
    \begin{aligned}[t]
      &H\colon \mathbb{N}\inj M\dsep x\in M\dsep y\in M
      \dsep x\notin H[\mathbb{N}]\dsep y\in H[\mathbb{N}]\dsep{}\\
      &\SuccShift{H}(x)=\SuccShift{H}(y)\vdash x=y
    \end{aligned}
  }{%
    H\colon \mathbb{N}\inj M\dsep x\in M\dsep y\in M
    \dsep x\notin H[\mathbb{N}]\dsep y\in H[\mathbb{N}]\dsep
    \SuccShift{H}(x)=\SuccShift{H}(y)\vdash x=y
  }
    \proofstepwidestar[1]{H\colon \mathbb{N}\inj M}{\rA}
    \proofstepwidestar[2]{x\in M}{\rA}
    \proofstepwidestar[3]{y\in M}{\rA}
    \proofstepwidestar[4]{x\notin H[\mathbb{N}]}{\rA}
    \proofstepwidestar[5]{y\in H[\mathbb{N}]}{\rA}
    \proofstepwidestar[6]{\SuccShift{H}(x)=\SuccShift{H}(y)}{\rA}
    \proofstepwidestar[1]{H\colon \mathbb{N}\to M}{%
      \FormulaRefAuto{Injektive Funktion}{1}}
    \proofstepwidestar[1,5]{\exists n\in\mathbb{N}\,\bigl(y=H(n)\bigr)}{%
      \FormulaRefAuto{F\colon A\to B\dsep y\in F[A]\vdash \exists x\in A\,y=F(x)}{7,5}}
    \proofstepwidestar[9]{n\in\mathbb{N}}{\rA}
    \proofstepwidestar[10]{y=H(n)}{\rA}
    \proofstepwidestar[6,10]{\SuccShift{H}(x)=\SuccShift{H}(H(n))}{%
      \rIE{10,6}}
    \proofstepwidestar[1,2,4,9,11]{\bot}{%
      \FormulaRefAuto{%
        H\colon \mathbb{N}\inj M\dsep x\in M\dsep x\notin H[\mathbb{N}]
        \dsep n\in\mathbb{N}\dsep
        \SuccShift{H}(x)=\SuccShift{H}(H(n))\vdash \bot}{1,2,4,9,11}}
    \proofstepwidestar[1,2,4,6,9,10]{x=y}{\rCE{12}}
    \proofstepwidestar[1,2,4,5,6]{x=y}{\rEE{8,9,10,13}}
  \closeproofpartwideindR

  \proofpartwideindR[Außen-Außen-Fall]{%
    \begin{aligned}[t]
      &H\colon \mathbb{N}\inj M\dsep x\in M\dsep y\in M
      \dsep x\notin H[\mathbb{N}]\dsep y\notin H[\mathbb{N}]\dsep{}\\
      &\SuccShift{H}(x)=\SuccShift{H}(y)\vdash x=y
    \end{aligned}
  }{%
    H\colon \mathbb{N}\inj M\dsep x\in M\dsep y\in M
    \dsep x\notin H[\mathbb{N}]\dsep y\notin H[\mathbb{N}]\dsep
    \SuccShift{H}(x)=\SuccShift{H}(y)\vdash x=y
  }
    \proofstepwidestar[1]{H\colon \mathbb{N}\inj M}{\rA}
    \proofstepwidestar[2]{x\in M}{\rA}
    \proofstepwidestar[3]{y\in M}{\rA}
    \proofstepwidestar[4]{x\notin H[\mathbb{N}]}{\rA}
    \proofstepwidestar[5]{y\notin H[\mathbb{N}]}{\rA}
    \proofstepwidestar[6]{\SuccShift{H}(x)=\SuccShift{H}(y)}{\rA}
    \proofstepwidestar[1,2,4]{\SuccShift{H}(x)=x}{%
      \FormulaRefAuto{%
        H\colon \mathbb{N}\inj M\dsep x\in M\dsep x\notin H[\mathbb{N}]
        \vdash \SuccShift{H}(x)=x}{1,2,4}}
    \proofstepwidestar[1,3,5]{\SuccShift{H}(y)=y}{%
      \FormulaRefAuto{%
        H\colon \mathbb{N}\inj M\dsep x\in M\dsep x\notin H[\mathbb{N}]
        \vdash \SuccShift{H}(x)=x}{1,3,5}}
    \proofstepwidestar[1,2,3,4,5,6]{x=y}{%
      \FormulaRefAuto{a = b\dsep a = c\dsep b=d\vdash c = d}{6,7,8}}
  \closeproofpartwideindR

  \proofpartwideindR[Schlussfall $x$ im Bild]{%
    H\colon \mathbb{N}\inj M\dsep x\in M\dsep y\in M\dsep
    \SuccShift{H}(x)=\SuccShift{H}(y)\dsep x\in H[\mathbb{N}]
    \vdash x=y
  }{%
    H\colon \mathbb{N}\inj M\dsep x\in M\dsep y\in M\dsep
    \SuccShift{H}(x)=\SuccShift{H}(y)\dsep x\in H[\mathbb{N}]
    \vdash x=y
  }
    \proofstepwidestar[1]{H\colon \mathbb{N}\inj M}{\rA}
    \proofstepwidestar[2]{x\in M}{\rA}
    \proofstepwidestar[3]{y\in M}{\rA}
    \proofstepwidestar[4]{\SuccShift{H}(x)=\SuccShift{H}(y)}{\rA}
    \proofstepwidestar[5]{x\in H[\mathbb{N}]}{\rA}
    \proofstepwidestar{y\in H[\mathbb{N}]\lor y\notin H[\mathbb{N}]}{%
      \FormulaRefAuto{P\lor \neg P}}
    \proofstepwidestar[7]{y\in H[\mathbb{N}]}{\rA}
    \proofstepwidestar[1,2,3,4,5,7]{x=y}{%
      \FormulaRefAuto{%
        H\colon \mathbb{N}\inj M\dsep x\in M\dsep y\in M
        \dsep x\in H[\mathbb{N}]\dsep y\in H[\mathbb{N}]\dsep
        \SuccShift{H}(x)=\SuccShift{H}(y)\vdash x=y}{1,2,3,5,7,4}}
    \proofstepwidestar[9]{y\notin H[\mathbb{N}]}{\rA}
    \proofstepwidestar[1,2,3,4,5,9]{x=y}{%
      \FormulaRefAuto{%
        H\colon \mathbb{N}\inj M\dsep x\in M\dsep y\in M
        \dsep x\in H[\mathbb{N}]\dsep y\notin H[\mathbb{N}]\dsep
        \SuccShift{H}(x)=\SuccShift{H}(y)\vdash x=y}{1,2,3,5,9,4}}
    \proofstepwidestar[1,2,3,4,5]{x=y}{\rOE{6,7,8,9,10}}
  \closeproofpartwideindR

  \proofpartwideindR[Schlussfall $x$ außerhalb des Bildes]{%
    H\colon \mathbb{N}\inj M\dsep x\in M\dsep y\in M\dsep
    \SuccShift{H}(x)=\SuccShift{H}(y)\dsep x\notin H[\mathbb{N}]
    \vdash x=y
  }{%
    H\colon \mathbb{N}\inj M\dsep x\in M\dsep y\in M\dsep
    \SuccShift{H}(x)=\SuccShift{H}(y)\dsep x\notin H[\mathbb{N}]
    \vdash x=y
  }
    \proofstepwidestar[1]{H\colon \mathbb{N}\inj M}{\rA}
    \proofstepwidestar[2]{x\in M}{\rA}
    \proofstepwidestar[3]{y\in M}{\rA}
    \proofstepwidestar[4]{\SuccShift{H}(x)=\SuccShift{H}(y)}{\rA}
    \proofstepwidestar[5]{x\notin H[\mathbb{N}]}{\rA}
    \proofstepwidestar{y\in H[\mathbb{N}]\lor y\notin H[\mathbb{N}]}{%
      \FormulaRefAuto{P\lor \neg P}}
    \proofstepwidestar[7]{y\in H[\mathbb{N}]}{\rA}
    \proofstepwidestar[1,2,3,4,5,7]{x=y}{%
      \FormulaRefAuto{%
        H\colon \mathbb{N}\inj M\dsep x\in M\dsep y\in M
        \dsep x\notin H[\mathbb{N}]\dsep y\in H[\mathbb{N}]\dsep
        \SuccShift{H}(x)=\SuccShift{H}(y)\vdash x=y}{1,2,3,5,7,4}}
    \proofstepwidestar[9]{y\notin H[\mathbb{N}]}{\rA}
    \proofstepwidestar[1,2,3,4,5,9]{x=y}{%
      \FormulaRefAuto{%
        H\colon \mathbb{N}\inj M\dsep x\in M\dsep y\in M
        \dsep x\notin H[\mathbb{N}]\dsep y\notin H[\mathbb{N}]\dsep
        \SuccShift{H}(x)=\SuccShift{H}(y)\vdash x=y}{1,2,3,5,9,4}}
    \proofstepwidestar[1,2,3,4,5]{x=y}{\rOE{6,7,8,9,10}}
  \closeproofpartwideindR

  \proofpartwide{Schluss}
    \proofstepwidestar[1]{H\colon \mathbb{N}\inj M}{\rA}
    \proofstepwidestar[2]{x\in M}{\rA}
    \proofstepwidestar[3]{y\in M}{\rA}
    \proofstepwidestar[4]{\SuccShift{H}(x)=\SuccShift{H}(y)}{\rA}
    \proofstepwidestar{x\in H[\mathbb{N}]\lor x\notin H[\mathbb{N}]}{%
      \FormulaRefAuto{P\lor \neg P}}
    \proofstepwidestar[6]{x\in H[\mathbb{N}]}{\rA}
    \proofstepwidestar[1,2,3,4,6]{x=y}{%
      \FormulaRefAuto{%
        H\colon \mathbb{N}\inj M\dsep x\in M\dsep y\in M\dsep
        \SuccShift{H}(x)=\SuccShift{H}(y)\dsep x\in H[\mathbb{N}]
        \vdash x=y}{1,2,3,4,6}}
    \proofstepwidestar[8]{x\notin H[\mathbb{N}]}{\rA}
    \proofstepwidestar[1,2,3,4,8]{x=y}{%
      \FormulaRefAuto{%
        H\colon \mathbb{N}\inj M\dsep x\in M\dsep y\in M\dsep
        \SuccShift{H}(x)=\SuccShift{H}(y)\dsep x\notin H[\mathbb{N}]
        \vdash x=y}{1,2,3,4,8}}
    \proofstepwidestar[1,2,3,4]{x=y}{\rOE{5,6,7,8,9}}
  \closeproofpartwide
\end{tabproofsplitwide}

\paragraph{Injektivität als Funktion}

\FormulaThmDelta[Injektivität der Folgenverschiebung]{%
  H\colon \mathbb{N}\inj M \vdash \SuccShift{H}\colon M\inj M
}{%
  \DeltaDecl{Mengen}{M\dsep x\dsep y}%
  \DeltaDecl{Funktionen}{H}%
  \DeltaDecl{Funktionensymbol}{\SuccShift{H}}%
}
\begin{tabproof}
  \proofstep{1}{H\colon \mathbb{N}\inj M}{\rA}
  \proofstep{1}{\SuccShift{H}\colon M\to M}{%
    \FormulaRefAuto{H\colon \mathbb{N}\inj M \vdash \SuccShift{H}\colon M\to M}{1}}
  \proofstep{3}{x\in M}{\rA}
  \proofstep{4}{y\in M}{\rA}
  \proofstep{5}{\SuccShift{H}(x)=\SuccShift{H}(y)}{\rA}
  \proofstep{1,3,4,5}{x=y}{%
    \FormulaRefAuto{%
      H\colon \mathbb{N}\inj M\dsep x\in M\dsep y\in M
      \dsep \SuccShift{H}(x)=\SuccShift{H}(y)\vdash x=y}{1,3,4,5}}
  \proofstep{1}{\SuccShift{H}\colon M\inj M}{%
    \FormulaRefAuto{Injektive Funktion}{2,6}}
\end{tabproof}

\subsubsection{Nichtsurjektivität der Folgenverschiebung}

\paragraph{Kein Urbild des Anfangspunktes}

\FormulaThmDelta[Kein Urbild des Anfangspunktes]{%
  H\colon \mathbb{N}\inj M\dsep x\in M\dsep
  \SuccShift{H}(x)=H(0)\vdash \bot
}{%
  \DeltaDecl{Mengen}{M\dsep x\dsep n}%
  \DeltaDecl{Funktionen}{H}%
  \DeltaDecl{Funktionensymbol}{\SuccShift{H}}%
}
\begin{tabproofsplitwide}
  \proofpartwideindR[Bildfall]{%
    H\colon \mathbb{N}\inj M\dsep n\in\mathbb{N}
    \dsep \SuccShift{H}(H(n))=H(0)\vdash \bot
  }{%
    H\colon \mathbb{N}\inj M\dsep n\in\mathbb{N}
    \dsep \SuccShift{H}(H(n))=H(0)\vdash \bot
  }
    \proofstepwidestar[1]{H\colon \mathbb{N}\inj M}{\rA}
    \proofstepwidestar[2]{n\in\mathbb{N}}{\rA}
    \proofstepwidestar[3]{\SuccShift{H}(H(n))=H(0)}{\rA}
    \proofstepwidestar[1,2]{\SuccShift{H}(H(n))=H(\Succ(n))}{%
      \FormulaRefAuto{%
        H\colon \mathbb{N}\inj M\dsep n\in\mathbb{N}\vdash
        \SuccShift{H}(H(n))=H(\Succ(n))}{1,2}}
    \proofstepwidestar[1,2,3]{H(\Succ(n))=H(0)}{%
      \FormulaRefAuto{a = b,\, a = c \vdash b = c}{4,3}}
    \proofstepwidestar[1,2]{H(\Succ(n))\neq H(0)}{%
      \FormulaRefAuto{%
        H\colon \mathbb{N}\inj M\dsep n\in\mathbb{N}
        \vdash H(\Succ(n))\neq H(0)}{1,2}}
    \proofstepwidestar[1,2,3]{\bot}{\rBI{5,6}}
  \closeproofpartwideindR

  \proofpartwideindR[Außenfall]{%
    H\colon \mathbb{N}\inj M\dsep x\in M\dsep x\notin H[\mathbb{N}]
    \dsep \SuccShift{H}(x)=H(0)\vdash \bot
  }{%
    H\colon \mathbb{N}\inj M\dsep x\in M\dsep x\notin H[\mathbb{N}]
    \dsep \SuccShift{H}(x)=H(0)\vdash \bot
  }
    \proofstepwidestar[1]{H\colon \mathbb{N}\inj M}{\rA}
    \proofstepwidestar[2]{x\in M}{\rA}
    \proofstepwidestar[3]{x\notin H[\mathbb{N}]}{\rA}
    \proofstepwidestar[4]{\SuccShift{H}(x)=H(0)}{\rA}
    \proofstepwidestar[1,2,3]{\SuccShift{H}(x)=x}{%
      \FormulaRefAuto{%
        H\colon \mathbb{N}\inj M\dsep x\in M\dsep x\notin H[\mathbb{N}]
        \vdash \SuccShift{H}(x)=x}{1,2,3}}
    \proofstepwidestar[1,2,3,4]{x=H(0)}{%
      \FormulaRefAuto{a = b,\, a = c \vdash b = c}{5,4}}
    \proofstepwidestar[1]{H(0)\in H[\mathbb{N}]}{%
      \FormulaRefAuto{H\colon \mathbb{N}\inj M \vdash H(0)\in H[\mathbb{N}]}{1}}
    \proofstepwidestar[1,2,3,4]{x\in H[\mathbb{N}]}{\rIE{6,7}}
    \proofstepwidestar[1,2,3,4]{\bot}{\rBI{8,3}}
  \closeproofpartwideindR

  \proofpartwide{Schluss}
  \proofstepwidestar[1]{H\colon \mathbb{N}\inj M}{\rA}
  \proofstepwidestar[2]{x\in M}{\rA}
  \proofstepwidestar[3]{\SuccShift{H}(x)=H(0)}{\rA}
  \proofstepwidestar{x\in H[\mathbb{N}]\lor x\notin H[\mathbb{N}]}{%
    \FormulaRefAuto{P\lor \neg P}}

  \proofstepwidestar[5]{x\in H[\mathbb{N}]}{\rA}
  \proofstepwidestar[1]{H\colon \mathbb{N}\to M}{%
    \FormulaRefAuto{Injektive Funktion}{1}}
  \proofstepwidestar[1,5]{\exists n\in\mathbb{N}\,\bigl(x=H(n)\bigr)}{%
    \FormulaRefAuto{F\colon A\to B\dsep y\in F[A]\vdash \exists x\in A\,y=F(x)}{6,5}}
  \proofstepwidestar[8]{n\in\mathbb{N}}{\rA}
  \proofstepwidestar[9]{x=H(n)}{\rA}
  \proofstepwidestar[3,9]{\SuccShift{H}(H(n))=H(0)}{\rIE{9,3}}
  \proofstepwidestar[1,8,10]{\bot}{%
    \FormulaRefAuto{%
      H\colon \mathbb{N}\inj M\dsep n\in\mathbb{N}
      \dsep \SuccShift{H}(H(n))=H(0)\vdash \bot}{1,8,10}}
  \proofstepwidestar[1,3,5]{\bot}{\rEE{7,8,9,11}}

  \proofstepwidestar[13]{x\notin H[\mathbb{N}]}{\rA}
  \proofstepwidestar[1,2,3,13]{\bot}{%
    \FormulaRefAuto{%
      H\colon \mathbb{N}\inj M\dsep x\in M\dsep x\notin H[\mathbb{N}]
      \dsep \SuccShift{H}(x)=H(0)\vdash \bot}{1,2,13,3}}

  \proofstepwidestar[1,2,3]{\bot}{\rOE{4,5,12,13,14}}
  \closeproofpartwide
\end{tabproofsplitwide}

\paragraph{Folgerung für Surjektivität}

\FormulaThmDelta[Die Folgenverschiebung lässt den Anfangspunkt aus]{%
  H\colon \mathbb{N}\inj M \vdash \neg(\SuccShift{H}\colon M\sur M)
}{%
  \DeltaDecl{Mengen}{M\dsep x}%
  \DeltaDecl{Funktionen}{H}%
  \DeltaDecl{Funktionensymbol}{\SuccShift{H}}%
}
\begin{tabproof}
  \proofstep{1}{H\colon \mathbb{N}\inj M}{\rA}
  \proofstep{1}{H\colon \mathbb{N}\to M}{%
    \FormulaRefAuto{Injektive Funktion}{1}}
  \proofstep{1}{\SuccShift{H}\colon M\to M}{%
    \FormulaRefAuto{H\colon \mathbb{N}\inj M \vdash \SuccShift{H}\colon M\to M}{1}}
  \proofstep{}{0\in\mathbb{N}}{%
    \FormulaRefAuto{PeanoZero}}
  \proofstep{1}{H(0)\in M}{%
    \FormulaRefAuto{F\colon A\to B\dsep x\in A\vdash F(x)\in B}{2,4}}
  \proofstep{6}{\SuccShift{H}\colon M\sur M}{\rA}
  \proofstep{1,6}{\exists x\in M\,\SuccShift{H}(x)=H(0)}{%
    \FormulaRefAuto{Surjektivität}{6,5}}
  \proofstep{8}{x\in M}{\rA}
  \proofstep{9}{\SuccShift{H}(x)=H(0)}{\rA}
  \proofstep{1,8,9}{\bot}{%
    \FormulaRefAuto{%
      H\colon \mathbb{N}\inj M\dsep x\in M\dsep
      \SuccShift{H}(x)=H(0)\vdash \bot}{1,8,9}}
  \proofstep{1,6}{\bot}{\rEE{7,8,9,10}}
  \proofstep{1}{\neg(\SuccShift{H}\colon M\sur M)}{\rCI{6,11}}
\end{tabproof}

\section{Addition}

Für festes \(a\in\mathbb{N}\) liefert der Dedekindsche Rekursionssatz genau
eine Abbildung \(G_a\colon\mathbb{N}\to\mathbb{N}\) mit
\[
  G_a(0)=a,\qquad
  G_a(\Succ(b))=\Succ(G_a(b)).
\]
In der eingeführten Schreibweise ist diese Abbildung
\(\RecFun{a}{\Succ}\). Ihr Wert am zweiten Argument ist die Summe.

\FormulaDefDeltaK[Addition natürlicher Zahlen]{%
  a\in\mathbb{N}\dsep b\in\mathbb{N}\vdash
  a+b\coloneqq \RecFun{a}{\Succ}(b)%
}{Addition auf \(\mathbb{N}\)}{%
  \DeltaRow{Mengen}{a\dsep b}%
  \DeltaRow{Funktionensymbole}{\Succ\dsep \RecFun{a}{\Succ}}%
  \DeltaRow{Neue Symbole}{+}%
}

\FormulaThmDelta[Abgeschlossenheit der Addition]{%
  a\in\mathbb{N}\dsep b\in\mathbb{N}\vdash a+b\in\mathbb{N}%
}{%
  \DeltaRow{Mengen}{a\dsep b}%
}
\begin{tabproof}
  \proofstep{1}{a\in\mathbb{N}}{\rA}
  \proofstep{2}{b\in\mathbb{N}}{\rA}
  \proofstep{}{\Succ\colon\mathbb{N}\to\mathbb{N}}{%
    \FormulaRefAuto{PeanoSuccFunction}[ax]}
  \proofstep{1,2}{\RecFun{a}{\Succ}(b)\in\mathbb{N}}{%
    \FormulaRefAuto{m_0\in M\dsep f\colon M\to M\dsep n\in\mathbb{N}\vdash
      \RecFun{m_0}{f}(n)\in M}{1,3,2}}
  \proofstep{1,2}{a+b=\RecFun{a}{\Succ}(b)}{%
    \FormulaRefAuto{a\in\mathbb{N}\dsep b\in\mathbb{N}\vdash
      a+b\coloneqq \RecFun{a}{\Succ}(b)}{1,2}}
  \proofstep{1,2}{a+b\in\mathbb{N}}{\rIE{5,4}}
\end{tabproof}

\FormulaThmDelta[Nullregel der Addition rechts]{%
  a\in\mathbb{N}\vdash a+0=a%
}{%
  \DeltaRow{Mengen}{a}%
}
\begin{tabproof}
  \proofstep{1}{a\in\mathbb{N}}{\rA}
  \proofstep{}{0\in\mathbb{N}}{%
    \FormulaRefAuto{PeanoZero}}
  \proofstep{}{\Succ\colon\mathbb{N}\to\mathbb{N}}{%
    \FormulaRefAuto{PeanoSuccFunction}[ax]}
  \proofstep{1}{a+0=\RecFun{a}{\Succ}(0)}{%
    \FormulaRefAuto{a\in\mathbb{N}\dsep b\in\mathbb{N}\vdash
      a+b\coloneqq \RecFun{a}{\Succ}(b)}{1,2}}
  \proofstep{1}{\RecFun{a}{\Succ}(0)=a}{%
    \FormulaRefAuto{m_0\in M\dsep f\colon M\to M \vdash
      \RecFun{m_0}{f}(0)=m_0}{1,3}}
  \proofstep{1}{a+0=a}{%
    \FormulaRefAuto{a = b,\, b = c \vdash a = c}{4,5}}
\end{tabproof}

\FormulaThmDeltaR[Nachfolgerregel der Addition rechts]{%
  \begin{aligned}[t]
    &a\in\mathbb{N}\dsep b\in\mathbb{N}\vdash{}\\
    &a+\Succ(b)=\Succ(a+b)
  \end{aligned}
}{%
  a\in\mathbb{N}\dsep b\in\mathbb{N}\vdash a+\Succ(b)=\Succ(a+b)
}{%
  \DeltaRow{Mengen}{a\dsep b}%
  \DeltaRow{Funktionensymbole}{\Succ\dsep \RecFun{a}{\Succ}}%
}
\begin{tabproof}
  \proofstep{1}{a\in\mathbb{N}}{\rA}
  \proofstep{2}{b\in\mathbb{N}}{\rA}
  \proofstep{}{\Succ\colon\mathbb{N}\to\mathbb{N}}{%
    \FormulaRefAuto{PeanoSuccFunction}[ax]}
  \proofstep{2}{\Succ(b)\in\mathbb{N}}{%
    \FormulaRefAuto{PeanoSuccClosure}{2}}
  \proofstep{1,2}{a+\Succ(b)=\RecFun{a}{\Succ}(\Succ(b))}{%
    \FormulaRefAuto{a\in\mathbb{N}\dsep b\in\mathbb{N}\vdash
      a+b\coloneqq \RecFun{a}{\Succ}(b)}{1,4}}
  \proofstep{1,2}{\RecFun{a}{\Succ}(\Succ(b))=
    \Succ(\RecFun{a}{\Succ}(b))}{%
    \FormulaRefAuto{m_0\in M\dsep f\colon M\to M\dsep n\in\mathbb{N}\vdash
      \RecFun{m_0}{f}(\Succ(n))=f(\RecFun{m_0}{f}(n))}{1,3,2}}
  \proofstep{1,2}{a+\Succ(b)=\Succ(\RecFun{a}{\Succ}(b))}{%
    \FormulaRefAuto{a = b,\, b = c \vdash a = c}{5,6}}
  \proofstep{1,2}{a+b=\RecFun{a}{\Succ}(b)}{%
    \FormulaRefAuto{a\in\mathbb{N}\dsep b\in\mathbb{N}\vdash
      a+b\coloneqq \RecFun{a}{\Succ}(b)}{1,2}}
  \proofstep{1,2}{\Succ(a+b)=\Succ(\RecFun{a}{\Succ}(b))}{\rIE{8}}
  \proofstep{1,2}{\Succ(\RecFun{a}{\Succ}(b))=\Succ(a+b)}{%
    \FormulaRefAuto{a=b\vdash b=a}{9}}
  \proofstep{1,2}{a+\Succ(b)=\Succ(a+b)}{%
    \FormulaRefAuto{a = b,\, b = c \vdash a = c}{7,10}}
\end{tabproof}

\FormulaThmDeltaK[Nullregel der Addition links]{%
  n\in\mathbb{N}\vdash 0+n=n%
}{PeanoAddLeftZero}{%
  \DeltaRow{Mengen}{n}%
}
\begin{tabproofsplitwide}
  \proofpartwideindR[Induktionsanfang]{0+0=0}{0+0=0}
    \proofstepwidestar[]{0\in\mathbb{N}}{%
      \FormulaRefAuto{PeanoZero}[ax]}
    \proofstepwidestar[]{0+0=0}{%
      \FormulaRefAuto{a\in\mathbb{N}\vdash a+0=a}{1}}
  \closeproofpartwideindR

  \proofpartwideindR[Induktionsschritt]{%
    n\in\mathbb{N}\dsep 0+n=n\vdash 0+\Succ(n)=\Succ(n)
  }{%
    n\in\mathbb{N}\dsep 0+n=n\vdash 0+\Succ(n)=\Succ(n)
  }
    \proofstepwidestar[1]{n\in\mathbb{N}}{\rA}
    \proofstepwidestar[2]{0+n=n}{\rA}
    \proofstepwidestar[]{0\in\mathbb{N}}{%
      \FormulaRefAuto{PeanoZero}[ax]}
    \proofstepwidestar[1]{0+\Succ(n)=\Succ(0+n)}{%
      \FormulaRefAuto{a\in\mathbb{N}\dsep b\in\mathbb{N}\vdash
        a+\Succ(b)=\Succ(a+b)}{3,1}}
    \proofstepwidestar[2]{\Succ(0+n)=\Succ(n)}{\rIE{2}}
    \proofstepwidestar[1,2]{0+\Succ(n)=\Succ(n)}{%
      \FormulaRefAuto{a = b,\, b = c \vdash a = c}{4,5}}
  \closeproofpartwideindR

  \proofpartwide{Induktionsschluss}
    \proofstepwidestar[1]{n\in\mathbb{N}}{\rA}
    \proofstepwidestar[1]{0+n=n}{%
      \FormulaRefAuto{PeanoInductionRule}{IA,IS,1}}
  \closeproofpartwide
\end{tabproofsplitwide}

\FormulaThmDeltaK[Nachfolgerabstand \(1\)]{%
  n\in\mathbb{N}\vdash n+\Succ(0)=\Succ(n)%
}{PeanoAddSuccZero}{%
  \DeltaRow{Mengen}{n}%
}
\begin{tabproof}
  \proofstep{1}{n\in\mathbb{N}}{\rA}
  \proofstep{}{0\in\mathbb{N}}{%
    \FormulaRefAuto{PeanoZero}[ax]}
  \proofstep{1}{n+\Succ(0)=\Succ(n+0)}{%
    \FormulaRefAuto{a\in\mathbb{N}\dsep b\in\mathbb{N}\vdash
      a+\Succ(b)=\Succ(a+b)}{1,2}}
  \proofstep{1}{n+0=n}{%
    \FormulaRefAuto{a\in\mathbb{N}\vdash a+0=a}{1}}
  \proofstep{1}{\Succ(n+0)=\Succ(n)}{\rIE{4}}
  \proofstep{1}{n+\Succ(0)=\Succ(n)}{%
    \FormulaRefAuto{a = b,\, b = c \vdash a = c}{3,5}}
\end{tabproof}

\FormulaThmDeltaK[Nachfolgerregel der Addition links]{%
  a\in\mathbb{N}\dsep b\in\mathbb{N}\vdash
  \Succ(a)+b=\Succ(a+b)%
}{PeanoAddSuccLeft}{%
  \DeltaRow{Mengen}{a\dsep b}%
}
\begin{tabproofsplitwide}
  \proofpartwideindR[Induktionsanfang]{%
    a\in\mathbb{N}\vdash \Succ(a)+0=\Succ(a+0)
  }{%
    a\in\mathbb{N}\vdash \Succ(a)+0=\Succ(a+0)
  }
    \proofstepwidestar[1]{a\in\mathbb{N}}{\rA}
    \proofstepwidestar[1]{\Succ(a)\in\mathbb{N}}{%
      \FormulaRefAuto{PeanoSuccClosure}{1}}
    \proofstepwidestar[1]{\Succ(a)+0=\Succ(a)}{%
      \FormulaRefAuto{a\in\mathbb{N}\vdash a+0=a}{2}}
    \proofstepwidestar[1]{a+0=a}{%
      \FormulaRefAuto{a\in\mathbb{N}\vdash a+0=a}{1}}
    \proofstepwidestar[1]{\Succ(a+0)=\Succ(a)}{\rIE{4}}
    \proofstepwidestar[1]{\Succ(a)=\Succ(a+0)}{%
      \FormulaRefAuto{a=b\vdash b=a}{5}}
    \proofstepwidestar[1]{\Succ(a)+0=\Succ(a+0)}{%
      \FormulaRefAuto{a = b,\, b = c \vdash a = c}{3,6}}
  \closeproofpartwideindR

  \proofpartwideindR[Induktionsschritt]{%
    \begin{aligned}[t]
      &b\in\mathbb{N}\dsep
        \bigl(a\in\mathbb{N}\vdash \Succ(a)+b=\Succ(a+b)\bigr)
        \vdash{}\\
      &a\in\mathbb{N}\vdash
        \Succ(a)+\Succ(b)=\Succ(a+\Succ(b))
    \end{aligned}
  }{%
    b\in\mathbb{N}\dsep
    \bigl(a\in\mathbb{N}\vdash \Succ(a)+b=\Succ(a+b)\bigr)
    \vdash
    a\in\mathbb{N}\vdash \Succ(a)+\Succ(b)=\Succ(a+\Succ(b))
  }
    \proofstepwidestar[1]{b\in\mathbb{N}}{\rA}
    \proofstepwidestar[2]{a\in\mathbb{N}\vdash \Succ(a)+b=\Succ(a+b)}{\rA}
    \proofstepwidestar[3]{a\in\mathbb{N}}{\rA}
    \proofstepwidestar[3]{\Succ(a)\in\mathbb{N}}{%
      \FormulaRefAuto{PeanoSuccClosure}{3}}
    \proofstepwidestar[1,3]{\Succ(a)+\Succ(b)=\Succ(\Succ(a)+b)}{%
      \FormulaRefAuto{a\in\mathbb{N}\dsep b\in\mathbb{N}\vdash
        a+\Succ(b)=\Succ(a+b)}{4,1}}
    \proofstepwidestar[2,3]{\Succ(a)+b=\Succ(a+b)}{\rRE{2,3}}
    \proofstepwidestar[2,3]{\Succ(\Succ(a)+b)=\Succ(\Succ(a+b))}{\rIE{6}}
    \proofstepwidestar[1,2,3]{\Succ(a)+\Succ(b)=\Succ(\Succ(a+b))}{%
      \FormulaRefAuto{a = b,\, b = c \vdash a = c}{5,7}}
    \proofstepwidestar[1,3]{a+\Succ(b)=\Succ(a+b)}{%
      \FormulaRefAuto{a\in\mathbb{N}\dsep b\in\mathbb{N}\vdash
        a+\Succ(b)=\Succ(a+b)}{3,1}}
    \proofstepwidestar[1,3]{\Succ(a+\Succ(b))=\Succ(\Succ(a+b))}{\rIE{9}}
    \proofstepwidestar[1,3]{\Succ(\Succ(a+b))=\Succ(a+\Succ(b))}{%
      \FormulaRefAuto{a=b\vdash b=a}{10}}
    \proofstepwidestar[1,2,3]{\Succ(a)+\Succ(b)=\Succ(a+\Succ(b))}{%
      \FormulaRefAuto{a = b,\, b = c \vdash a = c}{8,11}}
    \proofstepwidestar[1,2]{%
      a\in\mathbb{N}\vdash \Succ(a)+\Succ(b)=\Succ(a+\Succ(b))}{%
      \rRI{3,12}}
  \closeproofpartwideindR

  \proofpartwide{Induktionsschluss}
    \proofstepwidestar[1]{a\in\mathbb{N}}{\rA}
    \proofstepwidestar[2]{b\in\mathbb{N}}{\rA}
    \proofstepwidestar[1,2]{\Succ(a)+b=\Succ(a+b)}{%
      \FormulaRefAuto{PeanoInductionRule}{IA,IS,2,1}}
  \closeproofpartwide
\end{tabproofsplitwide}

\FormulaThmDeltaK[Assoziativität der Addition]{%
  a\in\mathbb{N}\dsep b\in\mathbb{N}\dsep c\in\mathbb{N}\vdash
  (a+b)+c=a+(b+c)%
}{PeanoAddAssociative}{%
  \DeltaRow{Mengen}{a\dsep b\dsep c}%
}
\begin{tabproofsplitwide}
  \proofpartwideindR[Induktionsanfang]{%
    a\in\mathbb{N}\dsep b\in\mathbb{N}\vdash (a+b)+0=a+(b+0)
  }{%
    a\in\mathbb{N}\dsep b\in\mathbb{N}\vdash (a+b)+0=a+(b+0)
  }
    \proofstepwidestar[1]{a\in\mathbb{N}}{\rA}
    \proofstepwidestar[2]{b\in\mathbb{N}}{\rA}
    \proofstepwidestar[1,2]{a+b\in\mathbb{N}}{%
      \FormulaRefAuto{a\in\mathbb{N}\dsep b\in\mathbb{N}\vdash a+b\in\mathbb{N}}{1,2}}
    \proofstepwidestar[1,2]{(a+b)+0=a+b}{%
      \FormulaRefAuto{a\in\mathbb{N}\vdash a+0=a}{3}}
    \proofstepwidestar[2]{b+0=b}{%
      \FormulaRefAuto{a\in\mathbb{N}\vdash a+0=a}{2}}
    \proofstepwidestar[2]{a+(b+0)=a+b}{\rIE{5}}
    \proofstepwidestar[2]{a+b=a+(b+0)}{%
      \FormulaRefAuto{a=b\vdash b=a}{6}}
    \proofstepwidestar[1,2]{(a+b)+0=a+(b+0)}{%
      \FormulaRefAuto{a = b,\, b = c \vdash a = c}{4,7}}
  \closeproofpartwideindR

  \proofpartwideindR[Induktionsschritt]{%
    \begin{aligned}[t]
      &c\in\mathbb{N}\dsep
        \bigl(a\in\mathbb{N}\dsep b\in\mathbb{N}\vdash
        (a+b)+c=a+(b+c)\bigr)\vdash{}\\
      &a\in\mathbb{N}\dsep b\in\mathbb{N}\vdash
        (a+b)+\Succ(c)=a+(b+\Succ(c))
    \end{aligned}
  }{%
    c\in\mathbb{N}\dsep
    \bigl(a\in\mathbb{N}\dsep b\in\mathbb{N}\vdash
    (a+b)+c=a+(b+c)\bigr)
    \vdash
    a\in\mathbb{N}\dsep b\in\mathbb{N}\vdash
    (a+b)+\Succ(c)=a+(b+\Succ(c))
  }
    \proofstepwidestar[1]{c\in\mathbb{N}}{\rA}
    \proofstepwidestar[2]{%
      a\in\mathbb{N}\dsep b\in\mathbb{N}\vdash
      (a+b)+c=a+(b+c)}{\rA}
    \proofstepwidestar[3]{a\in\mathbb{N}}{\rA}
    \proofstepwidestar[4]{b\in\mathbb{N}}{\rA}
    \proofstepwidestar[3,4]{a+b\in\mathbb{N}}{%
      \FormulaRefAuto{a\in\mathbb{N}\dsep b\in\mathbb{N}\vdash a+b\in\mathbb{N}}{3,4}}
    \proofstepwidestar[3,4,1]{(a+b)+\Succ(c)=\Succ((a+b)+c)}{%
      \FormulaRefAuto{a\in\mathbb{N}\dsep b\in\mathbb{N}\vdash
        a+\Succ(b)=\Succ(a+b)}{5,1}}
    \proofstepwidestar[2,3,4]{(a+b)+c=a+(b+c)}{\rRE{2,3,4}}
    \proofstepwidestar[2,3,4]{\Succ((a+b)+c)=\Succ(a+(b+c))}{\rIE{7}}
    \proofstepwidestar[1,2,3,4]{(a+b)+\Succ(c)=\Succ(a+(b+c))}{%
      \FormulaRefAuto{a = b,\, b = c \vdash a = c}{6,8}}
    \proofstepwidestar[4,1]{b+\Succ(c)=\Succ(b+c)}{%
      \FormulaRefAuto{a\in\mathbb{N}\dsep b\in\mathbb{N}\vdash
        a+\Succ(b)=\Succ(a+b)}{4,1}}
    \proofstepwidestar[4,1]{b+c\in\mathbb{N}}{%
      \FormulaRefAuto{a\in\mathbb{N}\dsep b\in\mathbb{N}\vdash a+b\in\mathbb{N}}{4,1}}
    \proofstepwidestar[3,4,1]{a+(b+\Succ(c))=a+\Succ(b+c)}{\rIE{10}}
    \proofstepwidestar[3,4,1]{a+\Succ(b+c)=\Succ(a+(b+c))}{%
      \FormulaRefAuto{a\in\mathbb{N}\dsep b\in\mathbb{N}\vdash
        a+\Succ(b)=\Succ(a+b)}{3,11}}
    \proofstepwidestar[3,4,1]{a+(b+\Succ(c))=\Succ(a+(b+c))}{%
      \FormulaRefAuto{a = b,\, b = c \vdash a = c}{12,13}}
    \proofstepwidestar[3,4,1]{\Succ(a+(b+c))=a+(b+\Succ(c))}{%
      \FormulaRefAuto{a=b\vdash b=a}{14}}
    \proofstepwidestar[1,2,3,4]{(a+b)+\Succ(c)=a+(b+\Succ(c))}{%
      \FormulaRefAuto{a = b,\, b = c \vdash a = c}{9,15}}
    \proofstepwidestar[1,2]{%
      a\in\mathbb{N}\dsep b\in\mathbb{N}\vdash
      (a+b)+\Succ(c)=a+(b+\Succ(c))}{%
      \rRI{3,4,16}}
  \closeproofpartwideindR

  \proofpartwide{Induktionsschluss}
    \proofstepwidestar[1]{a\in\mathbb{N}}{\rA}
    \proofstepwidestar[2]{b\in\mathbb{N}}{\rA}
    \proofstepwidestar[3]{c\in\mathbb{N}}{\rA}
    \proofstepwidestar[1,2,3]{(a+b)+c=a+(b+c)}{%
      \FormulaRefAuto{PeanoInductionRule}{IA,IS,3,1,2}}
  \closeproofpartwide
\end{tabproofsplitwide}

\FormulaDefDeltaK[Zahlzeichen \(1\)]{%
  1\coloneqq\Succ(0)%
}{PeanoOne}{%
  \DeltaRow{Mengen}{0\dsep 1}%
  \DeltaRow{Funktionensymbol}{\Succ}%
  \DeltaRow{Definitionsprinzip}{Nachfolger von 0}%
}

\FormulaThmDeltaK[Das Zahlzeichen \(1\) ist natürlich]{%
  1\in\mathbb{N}%
}{PeanoOneInN}{%
  \DeltaRow{Mengen}{1}%
}
\begin{tabproof}
  \proofstep{}{1=\Succ(0)}{\FormulaRefAuto{PeanoOne}}
  \proofstep{}{0\in\mathbb{N}}{\FormulaRefAuto{PeanoZero}[ax]}
  \proofstep{}{\Succ(0)\in\mathbb{N}}{\FormulaRefAuto{PeanoSuccClosure}{2}}
  \proofstep{}{1\in\mathbb{N}}{\rIE{1,3}}
\end{tabproof}

\FormulaThmDeltaK[Nachfolger als Addition von \(1\)]{%
  n\in\mathbb{N}\vdash n+1=\Succ(n)%
}{PeanoAddOneSucc}{%
  \DeltaRow{Mengen}{n}%
}
\begin{tabproof}
  \proofstep{1}{n\in\mathbb{N}}{\rA}
  \proofstep{}{1=\Succ(0)}{\FormulaRefAuto{PeanoOne}}
  \proofstep{1}{n+1=n+\Succ(0)}{\rIE{2}}
  \proofstep{1}{n+\Succ(0)=\Succ(n)}{\FormulaRefAuto{PeanoAddSuccZero}{1}}
  \proofstep{1}{n+1=\Succ(n)}{%
    \FormulaRefAuto{a = b,\, b = c \vdash a = c}{3,4}}
\end{tabproof}

\section{Natürliche Ordnung}

Die Ordnung auf \(\mathbb{N}\) wird in diesem Band aus der Addition gewonnen.
Sie ist daher nicht die Mitgliedschaftsrelation der von-Neumann-Konstruktion.
Mit dem Zahlzeichen \(1\) schreiben wir Nachfolgerabstände bevorzugt als
\(n+1\); formal wird dies durch \(n+1=\Succ(n)\) abgesichert.

\FormulaDefDeltaK[Ordnungsrelation der natürlichen Zahlen]{%
  \leq_{\mathbb{N}}\coloneqq
  \{(m,n)\in \mathbb{N}\times\mathbb{N}\mid
    \exists k\in\mathbb{N}\,(m+k=n)\}%
}{Natürliche Ordnung}{%
  \DeltaRow{Mengen}{k\dsep m\dsep n}%
  \DeltaRow{Relation}{\leq_{\mathbb{N}}}%
  \DeltaRow{Definitionsprinzip}{Addition}%
}

\FormulaDefDeltaK[Notation der natürlichen Ordnung]{%
  m\in\mathbb{N}\dsep n\in\mathbb{N}\vdash
  m\leq n \coloneqq \exists k\in\mathbb{N}\,(m+k=n)%
}{Notation \(m\leq n\)}{%
  \DeltaRow{Mengen}{k\dsep m\dsep n}%
  \DeltaRow{Relation}{\leq}%
}

\FormulaThmDelta[Additive Charakterisierung der natürlichen Ordnung]{%
  m\in\mathbb{N}\dsep n\in\mathbb{N}\vdash
  m\leq n \leftrightarrow \exists k\in\mathbb{N}\,(m+k=n)%
}{%
  \DeltaRow{Mengen}{k\dsep m\dsep n}%
}
\begin{tabproof}
  \proofstep{1}{m\in\mathbb{N}}{\rA}
  \proofstep{2}{n\in\mathbb{N}}{\rA}
  \proofstep{1,2}{m\leq n \leftrightarrow \exists k\in\mathbb{N}\,(m+k=n)}{%
    \FormulaRefAuto{m\in\mathbb{N}\dsep n\in\mathbb{N}\vdash
      m\leq n \coloneqq \exists k\in\mathbb{N}\,(m+k=n)}{1,2}}
\end{tabproof}

\FormulaDefDeltaK[Strikte Ordnungsrelation der natürlichen Zahlen]{%
  <_{\mathbb{N}}\coloneqq
  \{(m,n)\in \mathbb{N}\times\mathbb{N}\mid
    \exists k\in\mathbb{N}\,(m+\Succ(k)=n)\}%
}{Strikte natürliche Ordnung}{%
  \DeltaRow{Mengen}{k\dsep m\dsep n}%
  \DeltaRow{Relation}{<_{\mathbb{N}}}%
  \DeltaRow{Definitionsprinzip}{Addition und Nachfolger}%
}

\FormulaDefDeltaK[Notation der strikten natürlichen Ordnung]{%
  m\in\mathbb{N}\dsep n\in\mathbb{N}\vdash
  m<n \coloneqq \exists k\in\mathbb{N}\,(m+\Succ(k)=n)%
}{Notation \(m<n\)}{%
  \DeltaRow{Mengen}{k\dsep m\dsep n}%
  \DeltaRow{Relation}{<}%
}

\FormulaThmDelta[Additive Charakterisierung der strikten natürlichen Ordnung]{%
  m\in\mathbb{N}\dsep n\in\mathbb{N}\vdash
  m<n \leftrightarrow \exists k\in\mathbb{N}\,(m+\Succ(k)=n)%
}{%
  \DeltaRow{Mengen}{k\dsep m\dsep n}%
}
\begin{tabproof}
  \proofstep{1}{m\in\mathbb{N}}{\rA}
  \proofstep{2}{n\in\mathbb{N}}{\rA}
  \proofstep{1,2}{m<n \leftrightarrow \exists k\in\mathbb{N}\,(m+\Succ(k)=n)}{%
    \FormulaRefAuto{m\in\mathbb{N}\dsep n\in\mathbb{N}\vdash
      m<n \coloneqq \exists k\in\mathbb{N}\,(m+\Succ(k)=n)}{1,2}}
\end{tabproof}

\FormulaThmDeltaK[Additive Charakterisierung der strikten natürlichen Ordnung mit \(+1\)]{%
  m\in\mathbb{N}\dsep n\in\mathbb{N}\vdash
  m<n \leftrightarrow \exists k\in\mathbb{N}\,(m+(k+1)=n)%
}{PeanoStrictAddOneDistance}{%
  \DeltaRow{Mengen}{k\dsep m\dsep n}%
}
\begin{tabproofsplitwide}
  \proofpartwideindR[Hinrichtung]{%
    m\in\mathbb{N}\dsep n\in\mathbb{N}\dsep m<n
    \vdash \exists k\in\mathbb{N}\,(m+(k+1)=n)
  }{%
    m\in\mathbb{N}\dsep n\in\mathbb{N}\dsep m<n
    \vdash \exists k\in\mathbb{N}\,(m+(k+1)=n)
  }
    \proofstepwidestar[1]{m\in\mathbb{N}}{\rA}
    \proofstepwidestar[2]{n\in\mathbb{N}}{\rA}
    \proofstepwidestar[3]{m<n}{\rA}
    \proofstepwidestar[1,2,3]{\exists k\in\mathbb{N}\,(m+\Succ(k)=n)}{%
      \FormulaRefAuto{m\in\mathbb{N}\dsep n\in\mathbb{N}\vdash
        m<n \leftrightarrow \exists k\in\mathbb{N}\,(m+\Succ(k)=n)}{1,2,3}}
    \proofstepwidestar[5]{k\in\mathbb{N}\land m+\Succ(k)=n}{\rA}
    \proofstepwidestar[5]{k\in\mathbb{N}}{\rAEa{5}}
    \proofstepwidestar[5]{m+\Succ(k)=n}{\rAEb{5}}
    \proofstepwidestar[5]{k+1=\Succ(k)}{\FormulaRefAuto{PeanoAddOneSucc}{6}}
    \proofstepwidestar[5]{m+(k+1)=m+\Succ(k)}{\rIE{8}}
    \proofstepwidestar[5]{m+(k+1)=n}{%
      \FormulaRefAuto{a = b,\, b = c \vdash a = c}{9,7}}
    \proofstepwidestar[5]{\exists k\in\mathbb{N}\,(m+(k+1)=n)}{\rEI{\rAI{6,10}}}
    \proofstepwidestar[1,2,3]{\exists k\in\mathbb{N}\,(m+(k+1)=n)}{\rEE{4,5,11}}
  \closeproofpartwideindR

  \proofpartwideindR[Rückrichtung]{%
    m\in\mathbb{N}\dsep n\in\mathbb{N}\dsep
    \exists k\in\mathbb{N}\,(m+(k+1)=n)\vdash m<n
  }{%
    m\in\mathbb{N}\dsep n\in\mathbb{N}\dsep
    \exists k\in\mathbb{N}\,(m+(k+1)=n)\vdash m<n
  }
    \proofstepwidestar[1]{m\in\mathbb{N}}{\rA}
    \proofstepwidestar[2]{n\in\mathbb{N}}{\rA}
    \proofstepwidestar[3]{\exists k\in\mathbb{N}\,(m+(k+1)=n)}{\rA}
    \proofstepwidestar[4]{k\in\mathbb{N}\land m+(k+1)=n}{\rA}
    \proofstepwidestar[4]{k\in\mathbb{N}}{\rAEa{4}}
    \proofstepwidestar[4]{m+(k+1)=n}{\rAEb{4}}
    \proofstepwidestar[4]{k+1=\Succ(k)}{\FormulaRefAuto{PeanoAddOneSucc}{5}}
    \proofstepwidestar[4]{m+(k+1)=m+\Succ(k)}{\rIE{7}}
    \proofstepwidestar[4]{m+\Succ(k)=m+(k+1)}{\FormulaRefAuto{a=b\vdash b=a}{8}}
    \proofstepwidestar[4]{m+\Succ(k)=n}{%
      \FormulaRefAuto{a = b,\, b = c \vdash a = c}{9,6}}
    \proofstepwidestar[4]{\exists k\in\mathbb{N}\,(m+\Succ(k)=n)}{\rEI{\rAI{5,10}}}
    \proofstepwidestar[1,2,4]{m<n}{%
      \FormulaRefAuto{m\in\mathbb{N}\dsep n\in\mathbb{N}\vdash
        m<n \leftrightarrow \exists k\in\mathbb{N}\,(m+\Succ(k)=n)}{1,2,11}}
    \proofstepwidestar[1,2,3]{m<n}{\rEE{3,4,12}}
  \closeproofpartwideindR

  \proofpartwide{Schluss}
    \proofstepwidestar[1]{m\in\mathbb{N}}{\rA}
    \proofstepwidestar[2]{n\in\mathbb{N}}{\rA}
    \proofstepwidestar[1,2]{%
      m<n\rightarrow \exists k\in\mathbb{N}\,(m+(k+1)=n)}{\rRI{3,4}}
    \proofstepwidestar[1,2]{%
      \exists k\in\mathbb{N}\,(m+(k+1)=n)\rightarrow m<n}{\rRI{5,6}}
    \proofstepwidestar[1,2]{%
      m<n \leftrightarrow \exists k\in\mathbb{N}\,(m+(k+1)=n)}{\rLRI{7,8}}
  \closeproofpartwide
\end{tabproofsplitwide}

\FormulaThmDeltaK[Reflexivität der natürlichen Ordnung]{%
  n\in\mathbb{N}\vdash n\leq n%
}{PeanoLeqReflexive}{%
  \DeltaRow{Mengen}{n}%
}
\begin{tabproof}
  \proofstep{1}{n\in\mathbb{N}}{\rA}
  \proofstep{}{0\in\mathbb{N}}{%
    \FormulaRefAuto{PeanoZero}[ax]}
  \proofstep{1}{n+0=n}{%
    \FormulaRefAuto{a\in\mathbb{N}\vdash a+0=a}{1}}
  \proofstep{1}{\exists k\in\mathbb{N}\,(n+k=n)}{\rEI{\rAI{2,3}}}
  \proofstep{1}{n\leq n}{%
    \FormulaRefAuto{m\in\mathbb{N}\dsep n\in\mathbb{N}\vdash
      m\leq n \leftrightarrow \exists k\in\mathbb{N}\,(m+k=n)}{1,1,4}}
\end{tabproof}

\FormulaThmDelta[Element ist kleiner oder gleich seinem Nachfolger]{%
  n\in\mathbb{N}\vdash n\leq \Succ(n)%
}{%
  \DeltaRow{Mengen}{n}%
}
\begin{tabproof}
  \proofstep{1}{n\in\mathbb{N}}{\rA}
  \proofstep{1}{\Succ(n)\in\mathbb{N}}{%
    \FormulaRefAuto{PeanoSuccClosure}{1}}
  \proofstep{}{0\in\mathbb{N}}{%
    \FormulaRefAuto{PeanoZero}[ax]}
  \proofstep{}{\Succ(0)\in\mathbb{N}}{%
    \FormulaRefAuto{PeanoSuccClosure}{3}}
  \proofstep{1}{n+\Succ(0)=\Succ(n)}{%
    \FormulaRefAuto{PeanoAddSuccZero}{1}}
  \proofstep{1}{\exists k\in\mathbb{N}\,(n+k=\Succ(n))}{\rEI{\rAI{4,5}}}
  \proofstep{1}{n\leq \Succ(n)}{%
    \FormulaRefAuto{m\in\mathbb{N}\dsep n\in\mathbb{N}\vdash
      m\leq n \leftrightarrow \exists k\in\mathbb{N}\,(m+k=n)}{1,2,6}}
\end{tabproof}

\FormulaThmDeltaK[Element ist kleiner oder gleich \(n+1\)]{%
  n\in\mathbb{N}\vdash n\leq n+1%
}{PeanoLeqAddOne}{%
  \DeltaRow{Mengen}{n}%
}
\begin{tabproof}
  \proofstep{1}{n\in\mathbb{N}}{\rA}
  \proofstep{1}{n\leq \Succ(n)}{%
    \FormulaRefAuto{n\in\mathbb{N}\vdash n\leq \Succ(n)}{1}}
  \proofstep{1}{n+1=\Succ(n)}{\FormulaRefAuto{PeanoAddOneSucc}{1}}
  \proofstep{1}{\Succ(n)=n+1}{\FormulaRefAuto{a=b\vdash b=a}{3}}
  \proofstep{1}{n\leq n+1}{\rIE{4,2}}
\end{tabproof}

\FormulaThmDelta[Element ist strikt kleiner als sein Nachfolger]{%
  n\in\mathbb{N}\vdash n<\Succ(n)%
}{%
  \DeltaRow{Mengen}{n}%
}
\begin{tabproof}
  \proofstep{1}{n\in\mathbb{N}}{\rA}
  \proofstep{1}{\Succ(n)\in\mathbb{N}}{%
    \FormulaRefAuto{PeanoSuccClosure}{1}}
  \proofstep{}{0\in\mathbb{N}}{%
    \FormulaRefAuto{PeanoZero}[ax]}
  \proofstep{1}{n+\Succ(0)=\Succ(n)}{%
    \FormulaRefAuto{PeanoAddSuccZero}{1}}
  \proofstep{1}{\exists k\in\mathbb{N}\,(n+\Succ(k)=\Succ(n))}{\rEI{\rAI{3,4}}}
  \proofstep{1}{n<\Succ(n)}{%
    \FormulaRefAuto{m\in\mathbb{N}\dsep n\in\mathbb{N}\vdash
      m<n \leftrightarrow \exists k\in\mathbb{N}\,(m+\Succ(k)=n)}{1,2,5}}
\end{tabproof}

\FormulaThmDeltaK[Element ist strikt kleiner als \(n+1\)]{%
  n\in\mathbb{N}\vdash n<n+1%
}{PeanoLtAddOne}{%
  \DeltaRow{Mengen}{n}%
}
\begin{tabproof}
  \proofstep{1}{n\in\mathbb{N}}{\rA}
  \proofstep{1}{n<\Succ(n)}{%
    \FormulaRefAuto{n\in\mathbb{N}\vdash n<\Succ(n)}{1}}
  \proofstep{1}{n+1=\Succ(n)}{\FormulaRefAuto{PeanoAddOneSucc}{1}}
  \proofstep{1}{\Succ(n)=n+1}{\FormulaRefAuto{a=b\vdash b=a}{3}}
  \proofstep{1}{n<n+1}{\rIE{4,2}}
\end{tabproof}

\FormulaThmDeltaK[Strikte Ordnung impliziert natürliche Ordnung]{%
  m\in\mathbb{N}\dsep n\in\mathbb{N}\dsep m<n\vdash m\leq n%
}{PeanoStrictImpliesLeq}{%
  \DeltaRow{Mengen}{m\dsep n}%
}
\begin{tabproof}
  \proofstep{1}{m\in\mathbb{N}}{\rA}
  \proofstep{2}{n\in\mathbb{N}}{\rA}
  \proofstep{3}{m<n}{\rA}
  \proofstep{1,2,3}{\exists k\in\mathbb{N}\,(m+\Succ(k)=n)}{%
    \FormulaRefAuto{m\in\mathbb{N}\dsep n\in\mathbb{N}\vdash
      m<n \leftrightarrow \exists k\in\mathbb{N}\,(m+\Succ(k)=n)}{1,2,3}}
  \proofstep{5}{k\in\mathbb{N}\land m+\Succ(k)=n}{\rA}
  \proofstep{5}{k\in\mathbb{N}}{\rAEa{5}}
  \proofstep{5}{m+\Succ(k)=n}{\rAEb{5}}
  \proofstep{5}{\Succ(k)\in\mathbb{N}}{%
    \FormulaRefAuto{PeanoSuccClosure}{6}}
  \proofstep{5}{\exists j\in\mathbb{N}\,(m+j=n)}{\rEI{\rAI{8,7}}}
  \proofstep{1,2,3}{\exists j\in\mathbb{N}\,(m+j=n)}{\rEE{4,5,9}}
  \proofstep{1,2,3}{m\leq n}{%
    \FormulaRefAuto{m\in\mathbb{N}\dsep n\in\mathbb{N}\vdash
      m\leq n \leftrightarrow \exists k\in\mathbb{N}\,(m+k=n)}{1,2,10}}
\end{tabproof}

\FormulaThmDeltaK[Abstandsdarstellung über den Nachfolger]{%
  m\in\mathbb{N}\dsep n\in\mathbb{N}\vdash
  \Succ(m)\leq n \leftrightarrow
  \exists k\in\mathbb{N}\,(m+\Succ(k)=n)%
}{PeanoSuccLeqDistance}{%
  \DeltaRow{Mengen}{k\dsep m\dsep n}%
}
\begin{tabproofsplitwide}
  \proofpartwideindR[Hinrichtung]{%
    m\in\mathbb{N}\dsep n\in\mathbb{N}\dsep \Succ(m)\leq n
    \vdash \exists k\in\mathbb{N}\,(m+\Succ(k)=n)
  }{%
    m\in\mathbb{N}\dsep n\in\mathbb{N}\dsep \Succ(m)\leq n
    \vdash \exists k\in\mathbb{N}\,(m+\Succ(k)=n)
  }
    \proofstepwidestar[1]{m\in\mathbb{N}}{\rA}
    \proofstepwidestar[2]{n\in\mathbb{N}}{\rA}
    \proofstepwidestar[3]{\Succ(m)\leq n}{\rA}
    \proofstepwidestar[1]{\Succ(m)\in\mathbb{N}}{%
      \FormulaRefAuto{PeanoSuccClosure}{1}}
    \proofstepwidestar[1,2,3]{%
      \exists k\in\mathbb{N}\,(\Succ(m)+k=n)}{%
      \FormulaRefAuto{m\in\mathbb{N}\dsep n\in\mathbb{N}\vdash
        m\leq n \leftrightarrow \exists k\in\mathbb{N}\,(m+k=n)}{4,2,3}}
    \proofstepwidestar[6]{k\in\mathbb{N}\land \Succ(m)+k=n}{\rA}
    \proofstepwidestar[6]{k\in\mathbb{N}}{\rAEa{6}}
    \proofstepwidestar[6]{\Succ(m)+k=n}{\rAEb{6}}
    \proofstepwidestar[1,6]{\Succ(m)+k=\Succ(m+k)}{%
      \FormulaRefAuto{PeanoAddSuccLeft}{1,7}}
    \proofstepwidestar[1,6]{m+\Succ(k)=\Succ(m+k)}{%
      \FormulaRefAuto{a\in\mathbb{N}\dsep b\in\mathbb{N}\vdash
        a+\Succ(b)=\Succ(a+b)}{1,7}}
    \proofstepwidestar[1,6]{\Succ(m+k)=n}{%
      \FormulaRefAuto{a=b\dsep c=b\vdash a=c}{8,9}}
    \proofstepwidestar[1,6]{m+\Succ(k)=n}{%
      \FormulaRefAuto{a = b,\, b = c \vdash a = c}{10,11}}
    \proofstepwidestar[1,6]{%
      \exists k\in\mathbb{N}\,(m+\Succ(k)=n)}{\rEI{\rAI{7,12}}}
    \proofstepwidestar[1,2,3]{%
      \exists k\in\mathbb{N}\,(m+\Succ(k)=n)}{\rEE{5,6,13}}
  \closeproofpartwideindR

  \proofpartwideindR[Rückrichtung]{%
    m\in\mathbb{N}\dsep n\in\mathbb{N}\dsep
    \exists k\in\mathbb{N}\,(m+\Succ(k)=n)
    \vdash \Succ(m)\leq n
  }{%
    m\in\mathbb{N}\dsep n\in\mathbb{N}\dsep
    \exists k\in\mathbb{N}\,(m+\Succ(k)=n)
    \vdash \Succ(m)\leq n
  }
    \proofstepwidestar[1]{m\in\mathbb{N}}{\rA}
    \proofstepwidestar[2]{n\in\mathbb{N}}{\rA}
    \proofstepwidestar[3]{\exists k\in\mathbb{N}\,(m+\Succ(k)=n)}{\rA}
    \proofstepwidestar[1]{\Succ(m)\in\mathbb{N}}{%
      \FormulaRefAuto{PeanoSuccClosure}{1}}
    \proofstepwidestar[5]{k\in\mathbb{N}\land m+\Succ(k)=n}{\rA}
    \proofstepwidestar[5]{k\in\mathbb{N}}{\rAEa{5}}
    \proofstepwidestar[5]{m+\Succ(k)=n}{\rAEb{5}}
    \proofstepwidestar[1,5]{\Succ(m)+k=\Succ(m+k)}{%
      \FormulaRefAuto{PeanoAddSuccLeft}{1,6}}
    \proofstepwidestar[1,5]{m+\Succ(k)=\Succ(m+k)}{%
      \FormulaRefAuto{a\in\mathbb{N}\dsep b\in\mathbb{N}\vdash
        a+\Succ(b)=\Succ(a+b)}{1,6}}
    \proofstepwidestar[1,5]{\Succ(m+k)=n}{%
      \FormulaRefAuto{a=b\dsep a=c\vdash b=c}{9,7}}
    \proofstepwidestar[1,5]{\Succ(m)+k=n}{%
      \FormulaRefAuto{a = b,\, b = c \vdash a = c}{8,10}}
    \proofstepwidestar[1,5]{%
      \exists j\in\mathbb{N}\,(\Succ(m)+j=n)}{\rEI{\rAI{6,11}}}
    \proofstepwidestar[1,2,5]{\Succ(m)\leq n}{%
      \FormulaRefAuto{m\in\mathbb{N}\dsep n\in\mathbb{N}\vdash
        m\leq n \leftrightarrow \exists k\in\mathbb{N}\,(m+k=n)}{4,2,12}}
    \proofstepwidestar[1,2,3]{\Succ(m)\leq n}{\rEE{3,5,13}}
  \closeproofpartwideindR

  \proofpartwide{Schluss}
    \proofstepwidestar[1]{m\in\mathbb{N}}{\rA}
    \proofstepwidestar[2]{n\in\mathbb{N}}{\rA}
    \proofstepwidestar[1,2]{%
      \Succ(m)\leq n \rightarrow
      \exists k\in\mathbb{N}\,(m+\Succ(k)=n)}{\rRI{3,4}}
    \proofstepwidestar[1,2]{%
      \exists k\in\mathbb{N}\,(m+\Succ(k)=n)
      \rightarrow \Succ(m)\leq n}{\rRI{5,6}}
    \proofstepwidestar[1,2]{%
      \Succ(m)\leq n \leftrightarrow
      \exists k\in\mathbb{N}\,(m+\Succ(k)=n)}{\rLRI{7,8}}
  \closeproofpartwide
\end{tabproofsplitwide}

\FormulaThmDelta[Nachfolgercharakterisierung der strikten natürlichen Ordnung]{%
  m\in\mathbb{N}\dsep n\in\mathbb{N}\vdash
  m<n \leftrightarrow \Succ(m)\leq n%
}{%
  \DeltaRow{Mengen}{m\dsep n}%
}
\begin{tabproof}
  \proofstep{1}{m\in\mathbb{N}}{\rA}
  \proofstep{2}{n\in\mathbb{N}}{\rA}
  \proofstep{1,2}{%
    m<n \leftrightarrow \exists k\in\mathbb{N}\,(m+\Succ(k)=n)}{%
    \FormulaRefAuto{m\in\mathbb{N}\dsep n\in\mathbb{N}\vdash
      m<n \leftrightarrow \exists k\in\mathbb{N}\,(m+\Succ(k)=n)}{1,2}}
  \proofstep{1,2}{%
    \Succ(m)\leq n \leftrightarrow
    \exists k\in\mathbb{N}\,(m+\Succ(k)=n)}{%
    \FormulaRefAuto{PeanoSuccLeqDistance}{1,2}}
  \proofstep{1,2}{m<n \leftrightarrow \Succ(m)\leq n}{%
    \FormulaRefAuto{P \leftrightarrow Q, R \leftrightarrow Q \vdash P \leftrightarrow R}{3,4}}
\end{tabproof}

\FormulaThmDeltaK[\(1\)-Charakterisierung der strikten natürlichen Ordnung]{%
  m\in\mathbb{N}\dsep n\in\mathbb{N}\vdash
  m<n \leftrightarrow m+1\leq n%
}{PeanoLtAddOneLeq}{%
  \DeltaRow{Mengen}{m\dsep n}%
}
\begin{tabproof}
  \proofstep{1}{m\in\mathbb{N}}{\rA}
  \proofstep{2}{n\in\mathbb{N}}{\rA}
  \proofstep{1,2}{m<n \leftrightarrow \Succ(m)\leq n}{%
    \FormulaRefAuto{m\in\mathbb{N}\dsep n\in\mathbb{N}\vdash
      m<n \leftrightarrow \Succ(m)\leq n}{1,2}}
  \proofstep{1}{m+1=\Succ(m)}{\FormulaRefAuto{PeanoAddOneSucc}{1}}
  \proofstep{1}{\Succ(m)=m+1}{\FormulaRefAuto{a=b\vdash b=a}{4}}
  \proofstep{1,2}{\Succ(m)\leq n\leftrightarrow m+1\leq n}{\rIE{5}}
  \proofstep{1,2}{m<n \leftrightarrow m+1\leq n}{%
    \FormulaRefAuto{P\leftrightarrow Q\dsep Q\leftrightarrow R\vdash P\leftrightarrow R}{3,6}}
\end{tabproof}

\FormulaThmDelta[Einführung in die natürliche Ordnungsrelation]{%
  m\in\mathbb{N}\dsep n\in\mathbb{N}\dsep m\leq n
  \vdash (m,n)\in \leq_{\mathbb{N}}%
}{%
  \DeltaRow{Mengen}{m\dsep n}%
  \DeltaRow{Relation}{\leq_{\mathbb{N}}}%
}
\begin{tabproof}
  \proofstep{1}{m\in\mathbb{N}}{\rA}
  \proofstep{2}{n\in\mathbb{N}}{\rA}
  \proofstep{3}{m\leq n}{\rA}
  \proofstep{1,2,3}{\exists k\in\mathbb{N}\,(m+k=n)}{%
    \FormulaRefAuto{m\in\mathbb{N}\dsep n\in\mathbb{N}\vdash
      m\leq n \leftrightarrow \exists k\in\mathbb{N}\,(m+k=n)}{1,2,3}}
  \proofstep{1,2}{(m,n)\in \mathbb{N}\times\mathbb{N}}{%
    \FormulaRefAuto{a\in A\dsep b\in B\vdash (a,b)\in A\times B}{1,2}}
  \proofstep{1,2,3}{%
    (m,n)\in
    \{(u,v)\in \mathbb{N}\times\mathbb{N}\mid
      \exists k\in\mathbb{N}\,(u+k=v)\}}{%
    \FormulaRefAuto{x\in A \dsep P(x) \vdash x\in \{\,x\in A \mid P(x)\,\}}{5,4}}
  \proofstep{1,2,3}{(m,n)\in \leq_{\mathbb{N}}}{%
    \rIE{\FormulaRefAuto{\leq_{\mathbb{N}}\coloneqq
      \{(m,n)\in \mathbb{N}\times\mathbb{N}\mid
        \exists k\in\mathbb{N}\,(m+k=n)\}},6}}
\end{tabproof}

\FormulaThmDelta[Elimination der natürlichen Ordnungsrelation]{%
  (m,n)\in \leq_{\mathbb{N}}\vdash
  m\in\mathbb{N}\land n\in\mathbb{N}\land m\leq n%
}{%
  \DeltaRow{Mengen}{m\dsep n}%
  \DeltaRow{Relation}{\leq_{\mathbb{N}}}%
}
\begin{tabproof}
  \proofstep{1}{(m,n)\in \leq_{\mathbb{N}}}{\rA}
  \proofstep{}{\leq_{\mathbb{N}}=
    \{(u,v)\in \mathbb{N}\times\mathbb{N}\mid
      \exists k\in\mathbb{N}\,(u+k=v)\}}{%
    \FormulaRefAuto{\leq_{\mathbb{N}}\coloneqq
      \{(m,n)\in \mathbb{N}\times\mathbb{N}\mid
        \exists k\in\mathbb{N}\,(m+k=n)\}}}
  \proofstep{1}{%
    (m,n)\in
    \{(u,v)\in \mathbb{N}\times\mathbb{N}\mid
      \exists k\in\mathbb{N}\,(u+k=v)\}}{\rIE{2,1}}
  \proofstep{1}{%
    (m,n)\in \mathbb{N}\times\mathbb{N}\land
    \exists k\in\mathbb{N}\,(m+k=n)}{%
    \FormulaRefAuto{x\in\{u\in A\mid P(u)\}\eqvdash x\in A\land P(x)}{3}}
  \proofstep{1}{(m,n)\in \mathbb{N}\times\mathbb{N}}{\rAEa{4}}
  \proofstep{1}{\exists k\in\mathbb{N}\,(m+k=n)}{\rAEb{4}}
  \proofstep{1}{m\in\mathbb{N}\land n\in\mathbb{N}}{%
    \FormulaRefAuto{(a,b)\in A\times B \eqvdash a\in A \land b\in B}{5}}
  \proofstep{1}{m\in\mathbb{N}}{\rAEa{7}}
  \proofstep{1}{n\in\mathbb{N}}{\rAEb{7}}
  \proofstep{1}{m\leq n}{%
    \FormulaRefAuto{m\in\mathbb{N}\dsep n\in\mathbb{N}\vdash
      m\leq n \leftrightarrow \exists k\in\mathbb{N}\,(m+k=n)}{8,9,6}}
  \proofstep{1}{m\in\mathbb{N}\land n\in\mathbb{N}\land m\leq n}{%
    \FormulaRefAuto{P\dsep Q\dsep R\vdash P\land Q\land R}{8,9,10}}
\end{tabproof}

\FormulaThmDeltaK[Fallunterscheidung der natürlichen Ordnung]{%
  m\in\mathbb{N}\dsep n\in\mathbb{N}\dsep m\leq n
  \vdash m<n\lor m=n%
}{PeanoLeqSplit}{%
  \DeltaRow{Mengen}{k\dsep m\dsep n\dsep y}%
}
\begin{tabproofsplitwide}
  \proofpartwideindR[Abstandszeuge]{%
    m\in\mathbb{N}\dsep n\in\mathbb{N}\dsep m\leq n
    \vdash \exists k\in\mathbb{N}\,(m+k=n)
  }{%
    m\in\mathbb{N}\dsep n\in\mathbb{N}\dsep m\leq n
    \vdash \exists k\in\mathbb{N}\,(m+k=n)
  }
    \proofstepwidestar[1]{m\in\mathbb{N}}{\rA}
    \proofstepwidestar[2]{n\in\mathbb{N}}{\rA}
    \proofstepwidestar[3]{m\leq n}{\rA}
    \proofstepwidestar[1,2,3]{\exists k\in\mathbb{N}\,(m+k=n)}{%
      \FormulaRefAuto{m\in\mathbb{N}\dsep n\in\mathbb{N}\vdash
        m\leq n \leftrightarrow \exists k\in\mathbb{N}\,(m+k=n)}{1,2,3}}
  \closeproofpartwideindR

  \proofpartwideindR[Nullabstand]{%
    m\in\mathbb{N}\dsep k\in\mathbb{N}\dsep m+k=n\dsep k=0
    \vdash m=n
  }{%
    m\in\mathbb{N}\dsep k\in\mathbb{N}\dsep m+k=n\dsep k=0
    \vdash m=n
  }
    \proofstepwidestar[1]{m\in\mathbb{N}}{\rA}
    \proofstepwidestar[2]{k\in\mathbb{N}}{\rA}
    \proofstepwidestar[3]{m+k=n}{\rA}
    \proofstepwidestar[4]{k=0}{\rA}
    \proofstepwidestar[1,4]{m+k=m}{%
      \rIE{4,\FormulaRefAuto{a\in\mathbb{N}\vdash a+0=a}{1}}}
    \proofstepwidestar[1,3,4]{m=n}{%
      \FormulaRefAuto{a=b\dsep a=c\vdash b=c}{5,3}}
  \closeproofpartwideindR

  \proofpartwideindR[Positiver Abstand]{%
    m\in\mathbb{N}\dsep n\in\mathbb{N}\dsep k\in\mathbb{N}\dsep
    m+k=n\dsep k\neq 0 \vdash m<n
  }{%
    m\in\mathbb{N}\dsep n\in\mathbb{N}\dsep k\in\mathbb{N}\dsep
    m+k=n\dsep k\neq 0 \vdash m<n
  }
    \proofstepwidestar[1]{m\in\mathbb{N}}{\rA}
    \proofstepwidestar[2]{n\in\mathbb{N}}{\rA}
    \proofstepwidestar[3]{k\in\mathbb{N}}{\rA}
    \proofstepwidestar[4]{m+k=n}{\rA}
    \proofstepwidestar[5]{k\neq 0}{\rA}
    \proofstepwidestar[3,5]{\exists y\in\mathbb{N}\,(k=\Succ(y))}{%
      \FormulaRefAuto{PeanoNonzeroIsSucc}{3,5}}
    \proofstepwidestar[7]{y\in\mathbb{N}\land k=\Succ(y)}{\rA}
    \proofstepwidestar[7]{y\in\mathbb{N}}{\rAEa{7}}
    \proofstepwidestar[7]{k=\Succ(y)}{\rAEb{7}}
    \proofstepwidestar[4,7]{m+\Succ(y)=n}{\rIE{9,4}}
    \proofstepwidestar[4,7]{\exists y\in\mathbb{N}\,(m+\Succ(y)=n)}{%
      \rEI{\rAI{8,10}}}
    \proofstepwidestar[1,2,4,7]{m<n}{%
      \FormulaRefAuto{m\in\mathbb{N}\dsep n\in\mathbb{N}\vdash
        m<n \leftrightarrow \exists k\in\mathbb{N}\,(m+\Succ(k)=n)}{1,2,11}}
    \proofstepwidestar[1,2,3,4,5]{m<n}{\rEE{6,7,12}}
  \closeproofpartwideindR

  \proofpartwide{Schluss}
    \proofstepwidestar[1]{m\in\mathbb{N}}{\rA}
    \proofstepwidestar[2]{n\in\mathbb{N}}{\rA}
    \proofstepwidestar[3]{m\leq n}{\rA}
    \proofstepwidestar[1,2,3]{\exists k\in\mathbb{N}\,(m+k=n)}{%
      \FormulaRefAuto{m\in\mathbb{N}\dsep n\in\mathbb{N}\dsep m\leq n
        \vdash \exists k\in\mathbb{N}\,(m+k=n)}{1,2,3}}
    \proofstepwidestar[5]{k\in\mathbb{N}\land m+k=n}{\rA}
    \proofstepwidestar[5]{k\in\mathbb{N}}{\rAEa{5}}
    \proofstepwidestar[5]{m+k=n}{\rAEb{5}}
    \proofstepwidestar[]{k=0\lor k\neq 0}{\FormulaRefAuto{P\lor\neg P}}
    \proofstepwidestar[9]{k=0}{\rA}
    \proofstepwidestar[1,5,9]{m=n}{%
      \FormulaRefAuto{m\in\mathbb{N}\dsep k\in\mathbb{N}\dsep m+k=n\dsep k=0
        \vdash m=n}{1,6,7,9}}
    \proofstepwidestar[1,5,9]{m<n\lor m=n}{\rOIb{10}}
    \proofstepwidestar[12]{k\neq 0}{\rA}
    \proofstepwidestar[1,2,5,12]{m<n}{%
      \FormulaRefAuto{m\in\mathbb{N}\dsep n\in\mathbb{N}\dsep k\in\mathbb{N}\dsep
        m+k=n\dsep k\neq 0 \vdash m<n}{1,2,6,7,12}}
    \proofstepwidestar[1,2,5,12]{m<n\lor m=n}{\rOIa{13}}
    \proofstepwidestar[1,2,5]{m<n\lor m=n}{\rOE{8,9,11,12,14}}
    \proofstepwidestar[1,2,3]{m<n\lor m=n}{\rEE{4,5,15}}
  \closeproofpartwide
\end{tabproofsplitwide}

\FormulaThmDeltaK[Positive Additionsabstände sind nicht stationär]{%
  m\in\mathbb{N}\dsep k\in\mathbb{N}\vdash m+\Succ(k)\neq m%
}{PeanoAddPositiveDistance}{%
  \DeltaRow{Mengen}{k\dsep m}%
}
\begin{tabproofsplitwide}
  \proofpartwideindR[Induktionsanfang]{%
    k\in\mathbb{N}\vdash 0+\Succ(k)\neq 0
  }{%
    k\in\mathbb{N}\vdash 0+\Succ(k)\neq 0
  }
    \proofstepwidestar[1]{k\in\mathbb{N}}{\rA}
    \proofstepwidestar[1]{\Succ(k)\in\mathbb{N}}{%
      \FormulaRefAuto{PeanoSuccClosure}{1}}
    \proofstepwidestar[1]{0+\Succ(k)=\Succ(k)}{%
      \FormulaRefAuto{PeanoAddLeftZero}{2}}
    \proofstepwidestar[1]{\Succ(k)\neq 0}{%
      \FormulaRefAuto{PeanoSuccNonzero}{1}}
    \proofstepwidestar[1]{0+\Succ(k)\neq 0}{\rIE{3,4}}
  \closeproofpartwideindR

  \proofpartwideindR[Induktionsschritt]{%
    \begin{aligned}[t]
      &m\in\mathbb{N}\dsep
        \bigl(k\in\mathbb{N}\vdash m+\Succ(k)\neq m\bigr)
        \vdash{}\\
      &k\in\mathbb{N}\vdash \Succ(m)+\Succ(k)\neq \Succ(m)
    \end{aligned}
  }{%
    m\in\mathbb{N}\dsep
    \bigl(k\in\mathbb{N}\vdash m+\Succ(k)\neq m\bigr)
    \vdash
    k\in\mathbb{N}\vdash \Succ(m)+\Succ(k)\neq \Succ(m)
  }
    \proofstepwidestar[1]{m\in\mathbb{N}}{\rA}
    \proofstepwidestar[2]{k\in\mathbb{N}\vdash m+\Succ(k)\neq m}{\rA}
    \proofstepwidestar[3]{k\in\mathbb{N}}{\rA}
    \proofstepwidestar[4]{\Succ(m)+\Succ(k)=\Succ(m)}{\rA}
    \proofstepwidestar[1]{\Succ(m)\in\mathbb{N}}{%
      \FormulaRefAuto{PeanoSuccClosure}{1}}
    \proofstepwidestar[1,3]{\Succ(m)+\Succ(k)=\Succ(m+\Succ(k))}{%
      \FormulaRefAuto{PeanoAddSuccLeft}{1,%
        \FormulaRefAuto{PeanoSuccClosure}{3}}}
    \proofstepwidestar[1,3,4]{\Succ(m+\Succ(k))=\Succ(m)}{%
      \FormulaRefAuto{a=b\dsep a=c\vdash b=c}{6,4}}
    \proofstepwidestar[1,3,4]{m+\Succ(k)=m}{%
      \FormulaRefAuto{PeanoSuccInjective}{%
        \FormulaRefAuto{a\in\mathbb{N}\dsep b\in\mathbb{N}\vdash a+b\in\mathbb{N}}{1,%
          \FormulaRefAuto{PeanoSuccClosure}{3}},1,7}}
    \proofstepwidestar[2,3]{m+\Succ(k)\neq m}{\rRE{2,3}}
    \proofstepwidestar[1,2,3,4]{\bot}{\rBI{8,9}}
    \proofstepwidestar[1,2,3]{\Succ(m)+\Succ(k)\neq \Succ(m)}{\rCI{4,10}}
    \proofstepwidestar[1,2]{%
      k\in\mathbb{N}\vdash \Succ(m)+\Succ(k)\neq \Succ(m)}{\rRI{3,11}}
  \closeproofpartwideindR

  \proofpartwide{Induktionsschluss}
    \proofstepwidestar[1]{m\in\mathbb{N}}{\rA}
    \proofstepwidestar[2]{k\in\mathbb{N}}{\rA}
    \proofstepwidestar[1,2]{m+\Succ(k)\neq m}{%
      \FormulaRefAuto{PeanoInductionRule}{IA,IS,1,2}}
  \closeproofpartwide
\end{tabproofsplitwide}

\FormulaThmDeltaK[Zyklenfreiheit der natürlichen Ordnung]{%
  m\in\mathbb{N}\dsep n\in\mathbb{N}\dsep m<n\dsep n\leq m\vdash \bot%
}{PeanoOrderNoCycle}{%
  \DeltaRow{Mengen}{m\dsep n}%
}
\begin{tabproofsplitwide}
  \proofpartwideindR[Strikter Abstandszeuge]{%
    m\in\mathbb{N}\dsep n\in\mathbb{N}\dsep m<n
    \vdash \exists k\in\mathbb{N}\,(m+\Succ(k)=n)
  }{%
    m\in\mathbb{N}\dsep n\in\mathbb{N}\dsep m<n
    \vdash \exists k\in\mathbb{N}\,(m+\Succ(k)=n)
  }
    \proofstepwidestar[1]{m\in\mathbb{N}}{\rA}
    \proofstepwidestar[2]{n\in\mathbb{N}}{\rA}
    \proofstepwidestar[3]{m<n}{\rA}
    \proofstepwidestar[1,2,3]{\exists k\in\mathbb{N}\,(m+\Succ(k)=n)}{%
      \FormulaRefAuto{m\in\mathbb{N}\dsep n\in\mathbb{N}\vdash
        m<n \leftrightarrow \exists k\in\mathbb{N}\,(m+\Succ(k)=n)}{1,2,3}}
  \closeproofpartwideindR

  \proofpartwideindR[Rückabstandszeuge]{%
    m\in\mathbb{N}\dsep n\in\mathbb{N}\dsep n\leq m
    \vdash \exists l\in\mathbb{N}\,(n+l=m)
  }{%
    m\in\mathbb{N}\dsep n\in\mathbb{N}\dsep n\leq m
    \vdash \exists l\in\mathbb{N}\,(n+l=m)
  }
    \proofstepwidestar[1]{m\in\mathbb{N}}{\rA}
    \proofstepwidestar[2]{n\in\mathbb{N}}{\rA}
    \proofstepwidestar[3]{n\leq m}{\rA}
    \proofstepwidestar[1,2,3]{\exists l\in\mathbb{N}\,(n+l=m)}{%
      \FormulaRefAuto{m\in\mathbb{N}\dsep n\in\mathbb{N}\vdash
        m\leq n \leftrightarrow \exists k\in\mathbb{N}\,(m+k=n)}{2,1,3}}
  \closeproofpartwideindR

  \proofpartwideindR[Positive Schleife]{%
    m\in\mathbb{N}\dsep k\in\mathbb{N}\dsep l\in\mathbb{N}\dsep
    m+\Succ(k)=n\dsep n+l=m \vdash \bot
  }{%
    m\in\mathbb{N}\dsep k\in\mathbb{N}\dsep l\in\mathbb{N}\dsep
    m+\Succ(k)=n\dsep n+l=m \vdash \bot
  }
    \proofstepwidestar[1]{m\in\mathbb{N}}{\rA}
    \proofstepwidestar[2]{k\in\mathbb{N}}{\rA}
    \proofstepwidestar[3]{l\in\mathbb{N}}{\rA}
    \proofstepwidestar[4]{m+\Succ(k)=n}{\rA}
    \proofstepwidestar[5]{n+l=m}{\rA}
    \proofstepwidestar[2]{\Succ(k)\in\mathbb{N}}{%
      \FormulaRefAuto{PeanoSuccClosure}{2}}
    \proofstepwidestar[1,2]{m+\Succ(k)\in\mathbb{N}}{%
      \FormulaRefAuto{a\in\mathbb{N}\dsep b\in\mathbb{N}\vdash a+b\in\mathbb{N}}{1,6}}
    \proofstepwidestar[4]{(m+\Succ(k))+l=n+l}{\rIE{4}}
    \proofstepwidestar[4,5]{(m+\Succ(k))+l=m}{%
      \FormulaRefAuto{a = b,\, b = c \vdash a = c}{8,5}}
    \proofstepwidestar[1,2,3]{(m+\Succ(k))+l=m+(\Succ(k)+l)}{%
      \FormulaRefAuto{PeanoAddAssociative}{1,6,3}}
    \proofstepwidestar[1,2,3,4,5]{m+(\Succ(k)+l)=m}{%
      \FormulaRefAuto{a=b\dsep a=c\vdash b=c}{10,9}}
    \proofstepwidestar[2,3]{\Succ(k)+l=\Succ(k+l)}{%
      \FormulaRefAuto{PeanoAddSuccLeft}{2,3}}
    \proofstepwidestar[2,3]{k+l\in\mathbb{N}}{%
      \FormulaRefAuto{a\in\mathbb{N}\dsep b\in\mathbb{N}\vdash a+b\in\mathbb{N}}{2,3}}
    \proofstepwidestar[1,2,3,4,5]{m+\Succ(k+l)=m}{\rIE{12,11}}
    \proofstepwidestar[1,2,3]{m+\Succ(k+l)\neq m}{%
      \FormulaRefAuto{PeanoAddPositiveDistance}{1,13}}
    \proofstepwidestar[1,2,3,4,5]{\bot}{\rBI{14,15}}
  \closeproofpartwideindR

  \proofpartwide{Schluss}
    \proofstepwidestar[1]{m\in\mathbb{N}}{\rA}
    \proofstepwidestar[2]{n\in\mathbb{N}}{\rA}
    \proofstepwidestar[3]{m<n}{\rA}
    \proofstepwidestar[4]{n\leq m}{\rA}
    \proofstepwidestar[1,2,3]{\exists k\in\mathbb{N}\,(m+\Succ(k)=n)}{%
      \FormulaRefAuto{m\in\mathbb{N}\dsep n\in\mathbb{N}\dsep m<n
        \vdash \exists k\in\mathbb{N}\,(m+\Succ(k)=n)}{1,2,3}}
    \proofstepwidestar[1,2,4]{\exists l\in\mathbb{N}\,(n+l=m)}{%
      \FormulaRefAuto{m\in\mathbb{N}\dsep n\in\mathbb{N}\dsep n\leq m
        \vdash \exists l\in\mathbb{N}\,(n+l=m)}{1,2,4}}
    \proofstepwidestar[7]{k\in\mathbb{N}\land m+\Succ(k)=n}{\rA}
    \proofstepwidestar[7]{k\in\mathbb{N}}{\rAEa{7}}
    \proofstepwidestar[7]{m+\Succ(k)=n}{\rAEb{7}}
    \proofstepwidestar[10]{l\in\mathbb{N}\land n+l=m}{\rA}
    \proofstepwidestar[10]{l\in\mathbb{N}}{\rAEa{10}}
    \proofstepwidestar[10]{n+l=m}{\rAEb{10}}
    \proofstepwidestar[1,7,10]{\bot}{%
      \FormulaRefAuto{m\in\mathbb{N}\dsep k\in\mathbb{N}\dsep l\in\mathbb{N}\dsep
        m+\Succ(k)=n\dsep n+l=m \vdash \bot}{1,8,11,9,12}}
    \proofstepwidestar[1,2,3,4,7]{\bot}{\rEE{6,10,13}}
    \proofstepwidestar[1,2,3,4]{\bot}{\rEE{5,7,14}}
  \closeproofpartwide
\end{tabproofsplitwide}

\FormulaThmDeltaK[Vergleichbarkeit der natürlichen Ordnung]{%
  m\in\mathbb{N}\dsep n\in\mathbb{N}\vdash m\leq n\lor n\leq m%
}{PeanoOrderComparison}{%
  \DeltaRow{Mengen}{m\dsep n}%
}
\begin{tabproofsplitwide}
  \proofpartwideindR[Induktionsanfang]{%
    m\in\mathbb{N}\vdash m\leq 0\lor 0\leq m
  }{%
    m\in\mathbb{N}\vdash m\leq 0\lor 0\leq m
  }[PeanoOrderComparisonBase]
    \proofstepwidestar[1]{m\in\mathbb{N}}{\rA}
    \proofstepwidestar[]{0\in\mathbb{N}}{%
      \FormulaRefAuto{PeanoZero}[ax]}
    \proofstepwidestar[1]{0+m=m}{%
      \FormulaRefAuto{PeanoAddLeftZero}{1}}
    \proofstepwidestar[1]{\exists k\in\mathbb{N}\,(0+k=m)}{\rEI{\rAI{1,3}}}
    \proofstepwidestar[1]{0\leq m}{%
      \FormulaRefAuto{m\in\mathbb{N}\dsep n\in\mathbb{N}\vdash
        m\leq n \leftrightarrow \exists k\in\mathbb{N}\,(m+k=n)}{2,1,4}}
    \proofstepwidestar[1]{m\leq 0\lor 0\leq m}{\rOIb{5}}
  \closeproofpartwideindR

  \proofpartwideindR[Linker Fortschritt]{%
    m\in\mathbb{N}\dsep n\in\mathbb{N}\dsep m\leq n
    \vdash m\leq \Succ(n)
  }{%
    m\in\mathbb{N}\dsep n\in\mathbb{N}\dsep m\leq n
    \vdash m\leq \Succ(n)
  }
    \proofstepwidestar[1]{m\in\mathbb{N}}{\rA}
    \proofstepwidestar[2]{n\in\mathbb{N}}{\rA}
    \proofstepwidestar[3]{m\leq n}{\rA}
    \proofstepwidestar[1,2,3]{\exists k\in\mathbb{N}\,(m+k=n)}{%
      \FormulaRefAuto{m\in\mathbb{N}\dsep n\in\mathbb{N}\vdash
        m\leq n \leftrightarrow \exists k\in\mathbb{N}\,(m+k=n)}{1,2,3}}
    \proofstepwidestar[5]{k\in\mathbb{N}\land m+k=n}{\rA}
    \proofstepwidestar[5]{k\in\mathbb{N}}{\rAEa{5}}
    \proofstepwidestar[5]{m+k=n}{\rAEb{5}}
    \proofstepwidestar[1,5]{m+\Succ(k)=\Succ(m+k)}{%
      \FormulaRefAuto{a\in\mathbb{N}\dsep b\in\mathbb{N}\vdash
        a+\Succ(b)=\Succ(a+b)}{1,6}}
    \proofstepwidestar[5]{\Succ(m+k)=\Succ(n)}{\rIE{7}}
    \proofstepwidestar[1,5]{m+\Succ(k)=\Succ(n)}{%
      \FormulaRefAuto{a = b,\, b = c \vdash a = c}{8,9}}
    \proofstepwidestar[1,5]{\exists k\in\mathbb{N}\,(m+\Succ(k)=\Succ(n))}{%
      \rEI{\rAI{6,10}}}
    \proofstepwidestar[2]{\Succ(n)\in\mathbb{N}}{%
      \FormulaRefAuto{PeanoSuccClosure}{2}}
    \proofstepwidestar[1,2,5]{m<\Succ(n)}{%
      \FormulaRefAuto{m\in\mathbb{N}\dsep n\in\mathbb{N}\vdash
        m<n \leftrightarrow \exists k\in\mathbb{N}\,(m+\Succ(k)=n)}{1,12,11}}
    \proofstepwidestar[1,2,5]{m\leq \Succ(n)}{%
      \FormulaRefAuto{m\in\mathbb{N}\dsep n\in\mathbb{N}\dsep m<n\vdash m\leq n}{1,12,13}}
    \proofstepwidestar[1,2,3]{m\leq \Succ(n)}{\rEE{4,5,14}}
  \closeproofpartwideindR

  \proofpartwideindR[Rechter Fortschritt]{%
    n\in\mathbb{N}\dsep m\in\mathbb{N}\dsep n\leq m
    \vdash m\leq \Succ(n)\lor \Succ(n)\leq m
  }{%
    n\in\mathbb{N}\dsep m\in\mathbb{N}\dsep n\leq m
    \vdash m\leq \Succ(n)\lor \Succ(n)\leq m
  }
    \proofstepwidestar[1]{n\in\mathbb{N}}{\rA}
    \proofstepwidestar[2]{m\in\mathbb{N}}{\rA}
    \proofstepwidestar[3]{n\leq m}{\rA}
    \proofstepwidestar[1,2,3]{n<m\lor n=m}{%
      \FormulaRefAuto{PeanoLeqSplit}{1,2,3}}
    \proofstepwidestar[5]{n<m}{\rA}
    \proofstepwidestar[1,2,5]{\Succ(n)\leq m}{%
      \FormulaRefAuto{m\in\mathbb{N}\dsep n\in\mathbb{N}\vdash
        m<n \leftrightarrow \Succ(m)\leq n}{1,2,5}}
    \proofstepwidestar[1,2,5]{m\leq \Succ(n)\lor \Succ(n)\leq m}{\rOIb{6}}
    \proofstepwidestar[8]{n=m}{\rA}
    \proofstepwidestar[1]{n\leq \Succ(n)}{%
      \FormulaRefAuto{n\in\mathbb{N}\vdash n\leq \Succ(n)}{1}}
    \proofstepwidestar[1,8]{m\leq \Succ(n)}{\rIE{8,9}}
    \proofstepwidestar[1,8]{m\leq \Succ(n)\lor \Succ(n)\leq m}{\rOIa{10}}
    \proofstepwidestar[1,2,3]{m\leq \Succ(n)\lor \Succ(n)\leq m}{%
      \rOE{4,5,7,8,11}}
  \closeproofpartwideindR

  \proofpartwideindR[Induktionsschritt]{%
    n\in\mathbb{N}\dsep
    \bigl(m\in\mathbb{N}\vdash m\leq n\lor n\leq m\bigr)
    \vdash m\in\mathbb{N}\vdash m\leq \Succ(n)\lor \Succ(n)\leq m
  }{%
    n\in\mathbb{N}\dsep
    \bigl(m\in\mathbb{N}\vdash m\leq n\lor n\leq m\bigr)
    \vdash m\in\mathbb{N}\vdash m\leq \Succ(n)\lor \Succ(n)\leq m
  }[PeanoOrderComparisonStep]
    \proofstepwidestar[1]{n\in\mathbb{N}}{\rA}
    \proofstepwidestar[2]{m\in\mathbb{N}\vdash m\leq n\lor n\leq m}{\rA}
    \proofstepwidestar[3]{m\in\mathbb{N}}{\rA}
    \proofstepwidestar[2,3]{m\leq n\lor n\leq m}{\rRE{2,3}}
    \proofstepwidestar[5]{m\leq n}{\rA}
    \proofstepwidestar[1,3,5]{m\leq \Succ(n)}{%
      \FormulaRefAuto{m\in\mathbb{N}\dsep n\in\mathbb{N}\dsep m\leq n
        \vdash m\leq \Succ(n)}{3,1,5}}
    \proofstepwidestar[1,3,5]{m\leq \Succ(n)\lor \Succ(n)\leq m}{\rOIa{6}}
    \proofstepwidestar[8]{n\leq m}{\rA}
    \proofstepwidestar[1,3,8]{m\leq \Succ(n)\lor \Succ(n)\leq m}{%
      \FormulaRefAuto{n\in\mathbb{N}\dsep m\in\mathbb{N}\dsep n\leq m
        \vdash m\leq \Succ(n)\lor \Succ(n)\leq m}{1,3,8}}
    \proofstepwidestar[1,2,3]{m\leq \Succ(n)\lor \Succ(n)\leq m}{%
      \rOE{4,5,7,8,9}}
    \proofstepwidestar[1,2]{m\in\mathbb{N}\vdash
      m\leq \Succ(n)\lor \Succ(n)\leq m}{\rRI{3,10}}
  \closeproofpartwideindR

  \proofpartwide{Induktionsschluss}
    \proofstepwidestar[1]{m\in\mathbb{N}}{\rA}
    \proofstepwidestar[2]{n\in\mathbb{N}}{\rA}
    \proofstepwidestar[1,2]{m\leq n\lor n\leq m}{%
      \FormulaRefAuto{PeanoInductionRule}{%
        PeanoOrderComparisonBase,PeanoOrderComparisonStep,2,1}}
  \closeproofpartwide
\end{tabproofsplitwide}

\FormulaThmDeltaK[Trichotomie der natürlichen Ordnung]{%
  m\in\mathbb{N}\dsep n\in\mathbb{N}\vdash m<n\lor n<m\lor m=n%
}{PeanoOrderTrichotomy}{%
  \DeltaRow{Mengen}{m\dsep n}%
}
\begin{tabproofsplitwide}
  \proofpartwideindR[Linker Vergleich]{%
    m\in\mathbb{N}\dsep n\in\mathbb{N}\dsep m\leq n
    \vdash m<n\lor n<m\lor m=n
  }{%
    m\in\mathbb{N}\dsep n\in\mathbb{N}\dsep m\leq n
    \vdash m<n\lor n<m\lor m=n
  }
    \proofstepwidestar[1]{m\in\mathbb{N}}{\rA}
    \proofstepwidestar[2]{n\in\mathbb{N}}{\rA}
    \proofstepwidestar[3]{m\leq n}{\rA}
    \proofstepwidestar[1,2,3]{m<n\lor m=n}{%
      \FormulaRefAuto{PeanoLeqSplit}{1,2,3}}
    \proofstepwidestar[5]{m<n}{\rA}
    \proofstepwidestar[5]{m<n\lor n<m\lor m=n}{\rOIa{5}}
    \proofstepwidestar[7]{m=n}{\rA}
    \proofstepwidestar[7]{n<m\lor m=n}{\rOIb{7}}
    \proofstepwidestar[7]{m<n\lor n<m\lor m=n}{\rOIb{8}}
    \proofstepwidestar[1,2,3]{m<n\lor n<m\lor m=n}{\rOE{4,5,6,7,9}}
  \closeproofpartwideindR

  \proofpartwideindR[Rechter Vergleich]{%
    m\in\mathbb{N}\dsep n\in\mathbb{N}\dsep n\leq m
    \vdash m<n\lor n<m\lor m=n
  }{%
    m\in\mathbb{N}\dsep n\in\mathbb{N}\dsep n\leq m
    \vdash m<n\lor n<m\lor m=n
  }
    \proofstepwidestar[1]{m\in\mathbb{N}}{\rA}
    \proofstepwidestar[2]{n\in\mathbb{N}}{\rA}
    \proofstepwidestar[3]{n\leq m}{\rA}
    \proofstepwidestar[1,2,3]{n<m\lor n=m}{%
      \FormulaRefAuto{PeanoLeqSplit}{2,1,3}}
    \proofstepwidestar[5]{n<m}{\rA}
    \proofstepwidestar[5]{n<m\lor m=n}{\rOIa{5}}
    \proofstepwidestar[5]{m<n\lor n<m\lor m=n}{\rOIb{6}}
    \proofstepwidestar[8]{n=m}{\rA}
    \proofstepwidestar[8]{m=n}{\FormulaRefAuto{a=b\vdash b=a}{8}}
    \proofstepwidestar[8]{n<m\lor m=n}{\rOIb{9}}
    \proofstepwidestar[8]{m<n\lor n<m\lor m=n}{\rOIb{10}}
    \proofstepwidestar[1,2,3]{m<n\lor n<m\lor m=n}{\rOE{4,5,7,8,11}}
  \closeproofpartwideindR

  \proofpartwide{Schluss}
    \proofstepwidestar[1]{m\in\mathbb{N}}{\rA}
    \proofstepwidestar[2]{n\in\mathbb{N}}{\rA}
    \proofstepwidestar[1,2]{m\leq n\lor n\leq m}{%
      \FormulaRefAuto{PeanoOrderComparison}{1,2}}
    \proofstepwidestar[4]{m\leq n}{\rA}
    \proofstepwidestar[1,2,4]{m<n\lor n<m\lor m=n}{%
      \FormulaRefAuto{m\in\mathbb{N}\dsep n\in\mathbb{N}\dsep m\leq n
        \vdash m<n\lor n<m\lor m=n}{1,2,4}}
    \proofstepwidestar[6]{n\leq m}{\rA}
    \proofstepwidestar[1,2,6]{m<n\lor n<m\lor m=n}{%
      \FormulaRefAuto{m\in\mathbb{N}\dsep n\in\mathbb{N}\dsep n\leq m
        \vdash m<n\lor n<m\lor m=n}{1,2,6}}
    \proofstepwidestar[1,2]{m<n\lor n<m\lor m=n}{\rOE{3,4,5,6,7}}
  \closeproofpartwide
\end{tabproofsplitwide}

\FormulaThmDelta[Transitivität der natürlichen Ordnung]{%
  k\in\mathbb{N}\dsep m\in\mathbb{N}\dsep n\in\mathbb{N}\dsep
  k\leq m\dsep m\leq n\vdash k\leq n%
}{%
  \DeltaRow{Mengen}{k\dsep m\dsep n}%
}
\begin{tabproofsplitwide}
  \proofpartwideindR[Erster Abstand]{%
    k\in\mathbb{N}\dsep m\in\mathbb{N}\dsep k\leq m
    \vdash \exists u\in\mathbb{N}\,(k+u=m)
  }{%
    k\in\mathbb{N}\dsep m\in\mathbb{N}\dsep k\leq m
    \vdash \exists u\in\mathbb{N}\,(k+u=m)
  }
    \proofstepwidestar[1]{k\in\mathbb{N}}{\rA}
    \proofstepwidestar[2]{m\in\mathbb{N}}{\rA}
    \proofstepwidestar[3]{k\leq m}{\rA}
    \proofstepwidestar[1,2,3]{\exists u\in\mathbb{N}\,(k+u=m)}{%
      \FormulaRefAuto{m\in\mathbb{N}\dsep n\in\mathbb{N}\vdash
        m\leq n \leftrightarrow \exists k\in\mathbb{N}\,(m+k=n)}{1,2,3}}
  \closeproofpartwideindR

  \proofpartwideindR[Zweiter Abstand]{%
    m\in\mathbb{N}\dsep n\in\mathbb{N}\dsep m\leq n
    \vdash \exists v\in\mathbb{N}\,(m+v=n)
  }{%
    m\in\mathbb{N}\dsep n\in\mathbb{N}\dsep m\leq n
    \vdash \exists v\in\mathbb{N}\,(m+v=n)
  }
    \proofstepwidestar[1]{m\in\mathbb{N}}{\rA}
    \proofstepwidestar[2]{n\in\mathbb{N}}{\rA}
    \proofstepwidestar[3]{m\leq n}{\rA}
    \proofstepwidestar[1,2,3]{\exists v\in\mathbb{N}\,(m+v=n)}{%
      \FormulaRefAuto{m\in\mathbb{N}\dsep n\in\mathbb{N}\vdash
        m\leq n \leftrightarrow \exists k\in\mathbb{N}\,(m+k=n)}{1,2,3}}
  \closeproofpartwideindR

  \proofpartwideindR[Addition der Abstände]{%
    k\in\mathbb{N}\dsep n\in\mathbb{N}\dsep u\in\mathbb{N}\dsep v\in\mathbb{N}
    \dsep k+u=m\dsep m+v=n \vdash k\leq n
  }{%
    k\in\mathbb{N}\dsep n\in\mathbb{N}\dsep u\in\mathbb{N}\dsep v\in\mathbb{N}
    \dsep k+u=m\dsep m+v=n \vdash k\leq n
  }
    \proofstepwidestar[1]{k\in\mathbb{N}}{\rA}
    \proofstepwidestar[2]{n\in\mathbb{N}}{\rA}
    \proofstepwidestar[3]{u\in\mathbb{N}}{\rA}
    \proofstepwidestar[4]{v\in\mathbb{N}}{\rA}
    \proofstepwidestar[5]{k+u=m}{\rA}
    \proofstepwidestar[6]{m+v=n}{\rA}
    \proofstepwidestar[3,4]{u+v\in\mathbb{N}}{%
      \FormulaRefAuto{a\in\mathbb{N}\dsep b\in\mathbb{N}\vdash a+b\in\mathbb{N}}{3,4}}
    \proofstepwidestar[1,3,4]{(k+u)+v=k+(u+v)}{%
      \FormulaRefAuto{PeanoAddAssociative}{1,3,4}}
    \proofstepwidestar[5]{(k+u)+v=m+v}{\rIE{5}}
    \proofstepwidestar[5,6]{(k+u)+v=n}{%
      \FormulaRefAuto{a = b,\, b = c \vdash a = c}{9,6}}
    \proofstepwidestar[1,3,4,5,6]{k+(u+v)=n}{%
      \FormulaRefAuto{a=b\dsep a=c\vdash b=c}{8,10}}
    \proofstepwidestar[1,3,4,5,6]{\exists w\in\mathbb{N}\,(k+w=n)}{%
      \rEI{\rAI{7,11}}}
    \proofstepwidestar[1,2,3,4,5,6]{k\leq n}{%
      \FormulaRefAuto{m\in\mathbb{N}\dsep n\in\mathbb{N}\vdash
        m\leq n \leftrightarrow \exists k\in\mathbb{N}\,(m+k=n)}{1,2,12}}
  \closeproofpartwideindR

  \proofpartwide{Schluss}
    \proofstepwidestar[1]{k\in\mathbb{N}}{\rA}
    \proofstepwidestar[2]{m\in\mathbb{N}}{\rA}
    \proofstepwidestar[3]{n\in\mathbb{N}}{\rA}
    \proofstepwidestar[4]{k\leq m}{\rA}
    \proofstepwidestar[5]{m\leq n}{\rA}
    \proofstepwidestar[1,2,4]{\exists u\in\mathbb{N}\,(k+u=m)}{%
      \FormulaRefAuto{k\in\mathbb{N}\dsep m\in\mathbb{N}\dsep k\leq m
        \vdash \exists u\in\mathbb{N}\,(k+u=m)}{1,2,4}}
    \proofstepwidestar[2,3,5]{\exists v\in\mathbb{N}\,(m+v=n)}{%
      \FormulaRefAuto{m\in\mathbb{N}\dsep n\in\mathbb{N}\dsep m\leq n
        \vdash \exists v\in\mathbb{N}\,(m+v=n)}{2,3,5}}
    \proofstepwidestar[8]{u\in\mathbb{N}\land k+u=m}{\rA}
    \proofstepwidestar[8]{u\in\mathbb{N}}{\rAEa{8}}
    \proofstepwidestar[8]{k+u=m}{\rAEb{8}}
    \proofstepwidestar[11]{v\in\mathbb{N}\land m+v=n}{\rA}
    \proofstepwidestar[11]{v\in\mathbb{N}}{\rAEa{11}}
    \proofstepwidestar[11]{m+v=n}{\rAEb{11}}
    \proofstepwidestar[1,3,8,11]{k\leq n}{%
      \FormulaRefAuto{k\in\mathbb{N}\dsep n\in\mathbb{N}\dsep u\in\mathbb{N}\dsep v\in\mathbb{N}
        \dsep k+u=m\dsep m+v=n \vdash k\leq n}{1,3,9,12,10,13}}
    \proofstepwidestar[1,2,3,4,5,8]{k\leq n}{\rEE{7,11,14}}
    \proofstepwidestar[1,2,3,4,5]{k\leq n}{\rEE{6,8,15}}
  \closeproofpartwide
\end{tabproofsplitwide}

\FormulaThmDelta[Antisymmetrie der natürlichen Ordnung]{%
  m\in\mathbb{N}\dsep n\in\mathbb{N}\dsep m\leq n\dsep n\leq m\vdash m=n%
}{%
  \DeltaRow{Mengen}{m\dsep n}%
}
\begin{tabproofsplitwide}
  \proofpartwideindR[Strikter Fall ausgeschlossen]{%
    m\in\mathbb{N}\dsep n\in\mathbb{N}\dsep n\leq m\dsep m<n
    \vdash \bot
  }{%
    m\in\mathbb{N}\dsep n\in\mathbb{N}\dsep n\leq m\dsep m<n
    \vdash \bot
  }
    \proofstepwidestar[1]{m\in\mathbb{N}}{\rA}
    \proofstepwidestar[2]{n\in\mathbb{N}}{\rA}
    \proofstepwidestar[3]{n\leq m}{\rA}
    \proofstepwidestar[4]{m<n}{\rA}
    \proofstepwidestar[1,2,3,4]{\bot}{%
      \FormulaRefAuto{PeanoOrderNoCycle}{1,2,4,3}}
  \closeproofpartwideindR

  \proofpartwide{Schluss}
    \proofstepwidestar[1]{m\in\mathbb{N}}{\rA}
    \proofstepwidestar[2]{n\in\mathbb{N}}{\rA}
    \proofstepwidestar[3]{m\leq n}{\rA}
    \proofstepwidestar[4]{n\leq m}{\rA}
    \proofstepwidestar[1,2,3]{m<n\lor m=n}{%
      \FormulaRefAuto{PeanoLeqSplit}{1,2,3}}
    \proofstepwidestar[6]{m<n}{\rA}
    \proofstepwidestar[1,2,4,6]{\bot}{%
      \FormulaRefAuto{m\in\mathbb{N}\dsep n\in\mathbb{N}\dsep n\leq m\dsep m<n
        \vdash \bot}{1,2,4,6}}
    \proofstepwidestar[1,2,3,4,6]{m=n}{\rCE{6,7}}
    \proofstepwidestar[9]{m=n}{\rA}
    \proofstepwidestar[1,2,3,4]{m=n}{\rOE{5,6,8,9,9}}
  \closeproofpartwide
\end{tabproofsplitwide}

\FormulaThmDelta[Totale Ordnung der natürlichen Zahlen]{%
  \TotOrd{\leq_{\mathbb{N}},\mathbb{N}}%
}{%
  \DeltaRow{Relation}{\leq_{\mathbb{N}}}%
  \DeltaRow{Trägermenge}{\mathbb{N}}%
}
\begin{tabproofsplitwide}
  \proofpartwideindR[Reflexivität]{%
    x\in\mathbb{N}\vdash (x,x)\in \leq_{\mathbb{N}}
  }{%
    x\in\mathbb{N}\vdash (x,x)\in \leq_{\mathbb{N}}
  }
    \proofstepwidestar[1]{x\in\mathbb{N}}{\rA}
    \proofstepwidestar[1]{x\leq x}{%
      \FormulaRefAuto{n\in\mathbb{N}\vdash n\leq n}{1}}
    \proofstepwidestar[1]{(x,x)\in \leq_{\mathbb{N}}}{%
      \FormulaRefAuto{m\in\mathbb{N}\dsep n\in\mathbb{N}\dsep m\leq n
        \vdash (m,n)\in \leq_{\mathbb{N}}}{1,1,2}}
  \closeproofpartwideindR

  \proofpartwideindR[Transitivität]{%
    (x,y)\in \leq_{\mathbb{N}}\dsep (y,z)\in \leq_{\mathbb{N}}
    \vdash (x,z)\in \leq_{\mathbb{N}}
  }{%
    (x,y)\in \leq_{\mathbb{N}}\dsep (y,z)\in \leq_{\mathbb{N}}
    \vdash (x,z)\in \leq_{\mathbb{N}}
  }
    \proofstepwidestar[1]{(x,y)\in \leq_{\mathbb{N}}}{\rA}
    \proofstepwidestar[2]{(y,z)\in \leq_{\mathbb{N}}}{\rA}
    \proofstepwidestar[1]{x\in\mathbb{N}\land y\in\mathbb{N}\land x\leq y}{%
      \FormulaRefAuto{(m,n)\in \leq_{\mathbb{N}}\vdash
        m\in\mathbb{N}\land n\in\mathbb{N}\land m\leq n}{1}}
    \proofstepwidestar[2]{y\in\mathbb{N}\land z\in\mathbb{N}\land y\leq z}{%
      \FormulaRefAuto{(m,n)\in \leq_{\mathbb{N}}\vdash
        m\in\mathbb{N}\land n\in\mathbb{N}\land m\leq n}{2}}
    \proofstepwidestar[1]{x\in\mathbb{N}}{\rAEa{\rAEa{3}}}
    \proofstepwidestar[1]{y\in\mathbb{N}}{\FormulaRefAuto{P\land Q\land R\vdash Q}{3}}
    \proofstepwidestar[2]{z\in\mathbb{N}}{\FormulaRefAuto{P\land Q\land R\vdash Q}{4}}
    \proofstepwidestar[1]{x\leq y}{\FormulaRefAuto{P\land Q\land R\vdash R}{3}}
    \proofstepwidestar[2]{y\leq z}{\FormulaRefAuto{P\land Q\land R\vdash R}{4}}
    \proofstepwidestar[1,2]{x\leq z}{%
      \FormulaRefAuto{k\in\mathbb{N}\dsep m\in\mathbb{N}\dsep n\in\mathbb{N}\dsep
        k\leq m\dsep m\leq n\vdash k\leq n}{5,6,7,8,9}}
    \proofstepwidestar[1,2]{(x,z)\in \leq_{\mathbb{N}}}{%
      \FormulaRefAuto{m\in\mathbb{N}\dsep n\in\mathbb{N}\dsep m\leq n
        \vdash (m,n)\in \leq_{\mathbb{N}}}{5,7,10}}
  \closeproofpartwideindR

  \proofpartwideindR[Antisymmetrie]{%
    (x,y)\in \leq_{\mathbb{N}}\dsep (y,x)\in \leq_{\mathbb{N}}
    \vdash x=y
  }{%
    (x,y)\in \leq_{\mathbb{N}}\dsep (y,x)\in \leq_{\mathbb{N}}
    \vdash x=y
  }
    \proofstepwidestar[1]{(x,y)\in \leq_{\mathbb{N}}}{\rA}
    \proofstepwidestar[2]{(y,x)\in \leq_{\mathbb{N}}}{\rA}
    \proofstepwidestar[1]{x\in\mathbb{N}\land y\in\mathbb{N}\land x\leq y}{%
      \FormulaRefAuto{(m,n)\in \leq_{\mathbb{N}}\vdash
        m\in\mathbb{N}\land n\in\mathbb{N}\land m\leq n}{1}}
    \proofstepwidestar[2]{y\in\mathbb{N}\land x\in\mathbb{N}\land y\leq x}{%
      \FormulaRefAuto{(m,n)\in \leq_{\mathbb{N}}\vdash
        m\in\mathbb{N}\land n\in\mathbb{N}\land m\leq n}{2}}
    \proofstepwidestar[1]{x\in\mathbb{N}}{\rAEa{\rAEa{3}}}
    \proofstepwidestar[1]{y\in\mathbb{N}}{\FormulaRefAuto{P\land Q\land R\vdash Q}{3}}
    \proofstepwidestar[1]{x\leq y}{\FormulaRefAuto{P\land Q\land R\vdash R}{3}}
    \proofstepwidestar[2]{y\leq x}{\FormulaRefAuto{P\land Q\land R\vdash R}{4}}
    \proofstepwidestar[1,2]{x=y}{%
      \FormulaRefAuto{m\in\mathbb{N}\dsep n\in\mathbb{N}\dsep
        m\leq n\dsep n\leq m\vdash m=n}{5,6,7,8}}
  \closeproofpartwideindR

  \proofpartwideindR[Partielle Ordnung]{%
    \PartOrd{\leq_{\mathbb{N}},\mathbb{N}}
  }{%
    \PartOrd{\leq_{\mathbb{N}},\mathbb{N}}
  }
    \proofstepwidestar[]{\leq_{\mathbb{N}}\subseteq \mathbb{N}\times\mathbb{N}}{%
      \FormulaRefAuto{\leq_{\mathbb{N}}\coloneqq
        \{(m,n)\in \mathbb{N}\times\mathbb{N}\mid
          \exists k\in\mathbb{N}\,(m+k=n)\}}}
    \proofstepwidestar[]{x\in\mathbb{N}\vdash (x,x)\in \leq_{\mathbb{N}}}{%
      \FormulaRefAuto{x\in\mathbb{N}\vdash (x,x)\in \leq_{\mathbb{N}}}{}}
    \proofstepwidestar[]{(x,y)\in \leq_{\mathbb{N}}\dsep (y,z)\in \leq_{\mathbb{N}}
      \vdash (x,z)\in \leq_{\mathbb{N}}}{%
      \FormulaRefAuto{(x,y)\in \leq_{\mathbb{N}}\dsep (y,z)\in \leq_{\mathbb{N}}
        \vdash (x,z)\in \leq_{\mathbb{N}}}{}}
    \proofstepwidestar[]{(x,y)\in \leq_{\mathbb{N}}\dsep (y,x)\in \leq_{\mathbb{N}}
      \vdash x=y}{%
      \FormulaRefAuto{(x,y)\in \leq_{\mathbb{N}}\dsep (y,x)\in \leq_{\mathbb{N}}
        \vdash x=y}{}}
    \proofstepwidestar[]{\PartOrd{\leq_{\mathbb{N}},\mathbb{N}}}{%
      \FormulaRefAuto{Partielle Ordnung}[def]{1,2,3,4}}
  \closeproofpartwideindR

  \proofpartwideindR[Totalitätsbedingung]{%
    x\in\mathbb{N}\dsep y\in\mathbb{N}\vdash
    (x,y)\in \leq_{\mathbb{N}}\lor (y,x)\in \leq_{\mathbb{N}}
  }{%
    x\in\mathbb{N}\dsep y\in\mathbb{N}\vdash
    (x,y)\in \leq_{\mathbb{N}}\lor (y,x)\in \leq_{\mathbb{N}}
  }
    \proofstepwidestar[1]{x\in\mathbb{N}}{\rA}
    \proofstepwidestar[2]{y\in\mathbb{N}}{\rA}
    \proofstepwidestar[1,2]{x\leq y\lor y\leq x}{%
      \FormulaRefAuto{m\in\mathbb{N}\dsep n\in\mathbb{N}\vdash m\leq n\lor n\leq m}{1,2}}
    \proofstepwidestar[4]{x\leq y}{\rA}
    \proofstepwidestar[1,2,4]{(x,y)\in \leq_{\mathbb{N}}}{%
      \FormulaRefAuto{m\in\mathbb{N}\dsep n\in\mathbb{N}\dsep m\leq n
        \vdash (m,n)\in \leq_{\mathbb{N}}}{1,2,4}}
    \proofstepwidestar[1,2,4]{(x,y)\in \leq_{\mathbb{N}}\lor (y,x)\in \leq_{\mathbb{N}}}{%
      \rOIa{5}}
    \proofstepwidestar[7]{y\leq x}{\rA}
    \proofstepwidestar[1,2,7]{(y,x)\in \leq_{\mathbb{N}}}{%
      \FormulaRefAuto{m\in\mathbb{N}\dsep n\in\mathbb{N}\dsep m\leq n
        \vdash (m,n)\in \leq_{\mathbb{N}}}{2,1,7}}
    \proofstepwidestar[1,2,7]{(x,y)\in \leq_{\mathbb{N}}\lor (y,x)\in \leq_{\mathbb{N}}}{%
      \rOIb{8}}
    \proofstepwidestar[1,2]{(x,y)\in \leq_{\mathbb{N}}\lor (y,x)\in \leq_{\mathbb{N}}}{%
      \rOE{3,4,6,7,9}}
  \closeproofpartwideindR

  \proofpartwide{Schluss}
    \proofstepwidestar[]{\PartOrd{\leq_{\mathbb{N}},\mathbb{N}}}{%
      \FormulaRefAuto{\PartOrd{\leq_{\mathbb{N}},\mathbb{N}}}{}}
    \proofstepwidestar[]{x\in\mathbb{N}\dsep y\in\mathbb{N}\vdash
      (x,y)\in \leq_{\mathbb{N}}\lor (y,x)\in \leq_{\mathbb{N}}}{%
      \FormulaRefAuto{x\in\mathbb{N}\dsep y\in\mathbb{N}\vdash
        (x,y)\in \leq_{\mathbb{N}}\lor (y,x)\in \leq_{\mathbb{N}}}{}}
    \proofstepwidestar[]{\TotOrd{\leq_{\mathbb{N}},\mathbb{N}}}{%
      \FormulaRefAuto{Totale Ordnung}[def]{1,2}}
  \closeproofpartwide
\end{tabproofsplitwide}

\subsection{Linkskürzung und injektive Peano-Translationen}

Für die spätere Auslassung eines beliebigen Urbildpunktes benötigen wir eine
injektive Folge, die bei diesem Punkt beginnt.  Sie wird durch die
Peano-Translation $k\mapsto a+k$ geliefert.  Die dazu nötige Linkskürzung
wird hier ausschließlich aus der additiven Ordnung hergeleitet.

\FormulaThmDeltaK[Strikte Monotonie der Addition im rechten Argument]{%
  a\in\mathbb{N}\dsep b\in\mathbb{N}\dsep c\in\mathbb{N}\dsep b<c
  \vdash a+b<a+c%
}{PeanoAddLeftStrictMonotone}{%
  \DeltaRow{Mengen}{a\dsep b\dsep c\dsep k}%
}
\begin{tabproof}
  \proofstep{1}{a\in\mathbb{N}}{\rA}
  \proofstep{2}{b\in\mathbb{N}}{\rA}
  \proofstep{3}{c\in\mathbb{N}}{\rA}
  \proofstep{4}{b<c}{\rA}
  \proofstep{2,3,4}{\exists k\in\mathbb{N}\,(b+\Succ(k)=c)}{%
    \FormulaRefAuto{m\in\mathbb{N}\dsep n\in\mathbb{N}\vdash
      m<n\leftrightarrow\exists k\in\mathbb{N}\,(m+\Succ(k)=n)}{2,3,4}}
  \proofstep{6}{k\in\mathbb{N}\land b+\Succ(k)=c}{\rA}
  \proofstep{6}{k\in\mathbb{N}}{\rAEa{6}}
  \proofstep{6}{b+\Succ(k)=c}{\rAEb{6}}
  \proofstep{6}{\Succ(k)\in\mathbb{N}}{\FormulaRefAuto{PeanoSuccClosure}{7}}
  \proofstep{1,2,6}{a+(b+\Succ(k))=(a+b)+\Succ(k)}{%
    \FormulaRefAuto{a=b\vdash b=a}{%
      \FormulaRefAuto{PeanoAddAssociative}{1,2,9}}}
  \proofstep{1,6}{a+(b+\Succ(k))=a+c}{\rIE{8,\rII}}
  \proofstep{1,2,6}{(a+b)+\Succ(k)=a+c}{%
    \FormulaRefAuto{a=b\dsep a=c\vdash b=c}{10,11}}
  \proofstep{1,2,6}{a+b\in\mathbb{N}}{%
    \FormulaRefAuto{a\in\mathbb{N}\dsep b\in\mathbb{N}\vdash a+b\in\mathbb{N}}{1,2}}
  \proofstep{1,3,6}{a+c\in\mathbb{N}}{%
    \FormulaRefAuto{a\in\mathbb{N}\dsep b\in\mathbb{N}\vdash a+b\in\mathbb{N}}{1,3}}
  \proofstep{1,2,6}{\exists j\in\mathbb{N}\,((a+b)+\Succ(j)=a+c)}{%
    \rEI{\rAI{7,12}}}
  \proofstep{1,2,3,6}{a+b<a+c}{%
    \FormulaRefAuto{m\in\mathbb{N}\dsep n\in\mathbb{N}\vdash
      m<n\leftrightarrow\exists k\in\mathbb{N}\,(m+\Succ(k)=n)}{13,14,15}}
  \proofstep{1,2,3,4}{a+b<a+c}{\rEE{5,6,16}}
\end{tabproof}

\FormulaThmDeltaK[Strikte Ordnung schließt Gleichheit aus]{%
  m\in\mathbb{N}\dsep n\in\mathbb{N}\dsep m<n\dsep m=n\vdash\bot%
}{PeanoStrictEqualityImpossible}{%
  \DeltaRow{Mengen}{m\dsep n}%
}
\begin{tabproof}
  \proofstep{1}{m\in\mathbb{N}}{\rA}
  \proofstep{2}{n\in\mathbb{N}}{\rA}
  \proofstep{3}{m<n}{\rA}
  \proofstep{4}{m=n}{\rA}
  \proofstep{4}{n=m}{\FormulaRefAuto{a=b\vdash b=a}{4}}
  \proofstep{2}{n\leq n}{\FormulaRefAuto{PeanoLeqReflexive}{2}}
  \proofstep{2,4}{n\leq m}{\rIE{5,6}}
  \proofstep{1,2,3,4}{\bot}{%
    \FormulaRefAuto{PeanoOrderNoCycle}{1,2,3,7}}
\end{tabproof}

\FormulaThmDeltaK[Linkskürzung der Addition]{%
  a\in\mathbb{N}\dsep b\in\mathbb{N}\dsep c\in\mathbb{N}\dsep
  a+b=a+c\vdash b=c%
}{PeanoAddLeftCancellation}{%
  \DeltaRow{Mengen}{a\dsep b\dsep c}%
}
\begin{tabproof}
  \proofstep{1}{a\in\mathbb{N}}{\rA}
  \proofstep{2}{b\in\mathbb{N}}{\rA}
  \proofstep{3}{c\in\mathbb{N}}{\rA}
  \proofstep{4}{a+b=a+c}{\rA}
  \proofstep{1,2}{a+b\in\mathbb{N}}{%
    \FormulaRefAuto{a\in\mathbb{N}\dsep b\in\mathbb{N}\vdash a+b\in\mathbb{N}}{1,2}}
  \proofstep{1,3}{a+c\in\mathbb{N}}{%
    \FormulaRefAuto{a\in\mathbb{N}\dsep b\in\mathbb{N}\vdash a+b\in\mathbb{N}}{1,3}}
  \proofstep{2,3}{b<c\lor c<b\lor b=c}{%
    \FormulaRefAuto{m\in\mathbb{N}\dsep n\in\mathbb{N}\vdash
      m<n\lor n<m\lor m=n}{2,3}}
  \proofstep{8}{b<c}{\rA}
  \proofstep{1,2,3,8}{a+b<a+c}{\FormulaRefAuto{PeanoAddLeftStrictMonotone}{1,2,3,8}}
  \proofstep{1,2,3,4,8}{\bot}{\FormulaRefAuto{PeanoStrictEqualityImpossible}{5,6,9,4}}
  \proofstep{1,2,3,4,8}{b=c}{\rCE{10}}
  \proofstep{12}{c<b}{\rA}
  \proofstep{1,2,3,12}{a+c<a+b}{\FormulaRefAuto{PeanoAddLeftStrictMonotone}{1,3,2,12}}
  \proofstep{4}{a+c=a+b}{\FormulaRefAuto{a=b\vdash b=a}{4}}
  \proofstep{1,2,3,4,12}{\bot}{\FormulaRefAuto{PeanoStrictEqualityImpossible}{6,5,13,14}}
  \proofstep{1,2,3,4,12}{b=c}{\rCE{15}}
  \proofstep{17}{b=c}{\rA}
  \proofstep{1,2,3,4}{b=c}{\rOE{7,8,11,12,16,17,17}}
\end{tabproof}

\FormulaThmDeltaK[Injektivität der Peano-Translation]{%
  a\in\mathbb{N}\vdash \RecFun{a}{\Succ}\colon\mathbb{N}\inj\mathbb{N}%
}{PeanoTranslateInjection}{%
  \DeltaRow{Mengen}{a\dsep x\dsep y}%
  \DeltaRow{Funktionensymbol}{\RecFun{a}{\Succ}}%
}
\begin{tabproof}
  \proofstep{1}{a\in\mathbb{N}}{\rA}
  \proofstep{}{\Succ\colon\mathbb{N}\to\mathbb{N}}{\FormulaRefAuto{PeanoSuccFunction}[ax]}
  \proofstep{1}{\RecFun{a}{\Succ}\colon\mathbb{N}\to\mathbb{N}}{%
    \FormulaRefAuto{m_0\in M\dsep f\colon M\to M\vdash
      \RecFun{m_0}{f}\colon\mathbb{N}\to M}{1,2}}
  \proofstep{4}{x\in\mathbb{N}}{\rA}
  \proofstep{5}{y\in\mathbb{N}}{\rA}
  \proofstep{6}{\RecFun{a}{\Succ}(x)=\RecFun{a}{\Succ}(y)}{\rA}
  \proofstep{1,4}{a+x=\RecFun{a}{\Succ}(x)}{%
    \FormulaRefAuto{a\in\mathbb{N}\dsep b\in\mathbb{N}\vdash
      a+b\coloneqq\RecFun{a}{\Succ}(b)}{1,4}}
  \proofstep{1,5}{a+y=\RecFun{a}{\Succ}(y)}{%
    \FormulaRefAuto{a\in\mathbb{N}\dsep b\in\mathbb{N}\vdash
      a+b\coloneqq\RecFun{a}{\Succ}(b)}{1,5}}
  \proofstep{1,5}{\RecFun{a}{\Succ}(y)=a+y}{\FormulaRefAuto{a=b\vdash b=a}{8}}
  \proofstep{1,4,5,6}{a+x=\RecFun{a}{\Succ}(y)}{%
    \FormulaRefAuto{a=b\dsep b=c\vdash a=c}{7,6}}
  \proofstep{1,4,5,6}{a+x=a+y}{%
    \FormulaRefAuto{a=b\dsep b=c\vdash a=c}{10,9}}
  \proofstep{1,4,5,6}{x=y}{\FormulaRefAuto{PeanoAddLeftCancellation}{1,4,5,11}}
  \proofstep{1,4,5}{\RecFun{a}{\Succ}(x)=\RecFun{a}{\Succ}(y)\rightarrow x=y}{\rRI{6,12}}
  \proofstep{1,4}{y\in\mathbb{N}\rightarrow
    (\RecFun{a}{\Succ}(x)=\RecFun{a}{\Succ}(y)\rightarrow x=y)}{\rRI{5,13}}
  \proofstep{1,4}{\forall y\in\mathbb{N}\,
    (\RecFun{a}{\Succ}(x)=\RecFun{a}{\Succ}(y)\rightarrow x=y)}{\rUI{14}}
  \proofstep{1}{x\in\mathbb{N}\rightarrow\forall y\in\mathbb{N}\,
    (\RecFun{a}{\Succ}(x)=\RecFun{a}{\Succ}(y)\rightarrow x=y)}{\rRI{4,15}}
  \proofstep{1}{\forall x\in\mathbb{N}\,\forall y\in\mathbb{N}\,
    (\RecFun{a}{\Succ}(x)=\RecFun{a}{\Succ}(y)\rightarrow x=y)}{\rUI{16}}
  \proofstep{1}{\RecFun{a}{\Succ}\colon\mathbb{N}\inj\mathbb{N}}{%
    \FormulaRefAuto{Injektive Funktion}{3,17}}
\end{tabproof}

\section{Peano-Anfangsabschnitte}

Die endlichen Anfangsbereiche werden in diesem Band nicht mit natuerlichen
Zahlen als Mengen identifiziert. Fuer eine natuerliche Zahl \(n\) bezeichnet
\(\NatSeg{n}\) den durch die Peano-Ordnung bestimmten Anfangsabschnitt.

\FormulaDefDeltaK[Peano-Anfangsabschnitt]{%
  n\in\mathbb{N}\vdash
  \NatSeg{n}\coloneqq \{i\in\mathbb{N}\mid i<n\}%
}{PeanoNatSeg}{%
  \DeltaRow{Mengen}{i\dsep n}%
  \DeltaRow{Definitionsprinzip}{Peano-Ordnung}%
  \DeltaRow{Neue Symbole}{\NatSeg{n}}%
}

\FormulaThmDeltaK[Charakterisierung des Peano-Anfangsabschnitts]{%
  n\in\mathbb{N}\dsep i\in\mathbb{N}\vdash
  i\in\NatSeg{n}\leftrightarrow i<n%
}{PeanoNatSegMembership}{%
  \DeltaRow{Mengen}{i\dsep n}%
}
\begin{tabproof}
  \proofstep{1}{n\in\mathbb{N}}{\rA}
  \proofstep{2}{i\in\mathbb{N}}{\rA}
  \proofstep{1}{\NatSeg{n}=\{i\in\mathbb{N}\mid i<n\}}{%
    \FormulaRefAuto{PeanoNatSeg}{1}}
  \proofstep{1,2}{i\in\NatSeg{n}\leftrightarrow i<n}{%
    \FormulaRefAuto{x\in\{u\in A\mid P(u)\}\eqvdash x\in A\land P(x)}{3,2}}
\end{tabproof}

\FormulaThmDeltaK[Anfangsabschnitte sind Teilmengen von \(\mathbb{N}\)]{%
  n\in\mathbb{N}\vdash \NatSeg{n}\subseteq\mathbb{N}%
}{PeanoNatSegSubsetN}{%
  \DeltaRow{Mengen}{n}%
}
\begin{tabproof}
  \proofstep{1}{n\in\mathbb{N}}{\rA}
  \proofstep{1}{\NatSeg{n}=\{i\in\mathbb{N}\mid i<n\}}{\FormulaRefAuto{PeanoNatSeg}{1}}
  \proofstep{1}{\NatSeg{n}\subseteq\mathbb{N}}{%
    \FormulaRefAuto{A=\{x\in B\mid P(x)\}\vdash A\subseteq B}{2}}
\end{tabproof}

\FormulaThmDeltaK[Null ist kleiner oder gleich jeder natuerlichen Zahl]{%
  n\in\mathbb{N}\vdash 0\leq n%
}{PeanoZeroLeq}{%
  \DeltaRow{Mengen}{n\dsep k}%
}
\begin{tabproof}
  \proofstep{1}{n\in\mathbb{N}}{\rA}
  \proofstep{}{0\in\mathbb{N}}{\FormulaRefAuto{PeanoZero}[ax]}
  \proofstep{1}{0+n=n}{\FormulaRefAuto{PeanoAddLeftZero}{1}}
  \proofstep{1}{\exists k\in\mathbb{N}\,(0+k=n)}{\rEI{\rAI{1,3}}}
  \proofstep{1}{0\leq n}{%
    \FormulaRefAuto{m\in\mathbb{N}\dsep n\in\mathbb{N}\vdash
      m\leq n\leftrightarrow\exists k\in\mathbb{N}\,(m+k=n)}{2,1,4}}
\end{tabproof}

\FormulaThmDeltaK[Keine natuerliche Zahl ist kleiner als Null]{%
  n\in\mathbb{N}\vdash \neg(n<0)%
}{PeanoNoLtZero}{%
  \DeltaRow{Mengen}{n}%
}
\begin{tabproof}
  \proofstep{1}{n\in\mathbb{N}}{\rA}
  \proofstep{}{0\in\mathbb{N}}{\FormulaRefAuto{PeanoZero}[ax]}
  \proofstep{1}{0\leq n}{\FormulaRefAuto{PeanoZeroLeq}{1}}
  \proofstep{4}{n<0}{\rA}
  \proofstep{1,4}{\bot}{%
    \FormulaRefAuto{PeanoOrderNoCycle}{1,2,4,3}}
  \proofstep{1}{\neg(n<0)}{\rCI{4,5}}
\end{tabproof}

\FormulaThmDeltaK[Null-Anfangsabschnitt]{%
  \NatSeg{0}=\varnothing%
}{PeanoNatSegZero}{%
  \DeltaRow{Mengen}{}%
}
\begin{tabproof}
  \proofstep{}{0\in\mathbb{N}}{\FormulaRefAuto{PeanoZero}[ax]}
  \proofstep{}{\NatSeg{0}=\{i\in\mathbb{N}\mid i<0\}}{%
    \FormulaRefAuto{PeanoNatSeg}{1}}
  \proofstep{3}{i\in\mathbb{N}}{\rA}
  \proofstep{3}{\neg(i<0)}{\FormulaRefAuto{PeanoNoLtZero}{3}}
  \proofstep{}{\forall i\in\mathbb{N}\,\neg(i<0)}{%
    \rUI{\rRI{3,4}}}
  \proofstep{}{\NatSeg{0}=\varnothing}{%
    \FormulaRefAuto{EmptySeparationByNegPredicate}{2,5}}
\end{tabproof}

\FormulaThmDeltaK[Nichtstrikte Ordnung als strikte Nachfolgerordnung]{%
  m\in\mathbb{N}\dsep n\in\mathbb{N}\vdash
  m\leq n\leftrightarrow m<\Succ(n)%
}{PeanoLeqSuccStrict}{%
  \DeltaRow{Mengen}{k\dsep m\dsep n}%
  \DeltaRow{Relationen}{\leq_{\mathbb{N}}\dsep <_{\mathbb{N}}}%
}
\begin{tabproofsplitwide}
  \proofpartwideindR[Hinrichtung]{%
    m\in\mathbb{N}\dsep n\in\mathbb{N}\dsep m\leq n
    \vdash m<\Succ(n)%
  }{%
    m\in\mathbb{N}\dsep n\in\mathbb{N}\dsep m\leq n
    \vdash m<\Succ(n)%
  }[PeanoLeqSuccStrictForward]
    \proofstepwidestar[1]{m\in\mathbb{N}}{\rA}
    \proofstepwidestar[2]{n\in\mathbb{N}}{\rA}
    \proofstepwidestar[3]{m\leq n}{\rA}
    \proofstepwidestar[2]{\Succ(n)\in\mathbb{N}}{%
      \FormulaRefAuto{PeanoSuccClosure}{2}}
    \proofstepwidestar[1,2,3]{\exists k\in\mathbb{N}\,(m+k=n)}{%
      \FormulaRefAuto{m\in\mathbb{N}\dsep n\in\mathbb{N}\vdash
        m\leq n\leftrightarrow\exists k\in\mathbb{N}\,(m+k=n)}{1,2,3}}
    \proofstepwidestar[5]{k\in\mathbb{N}\land m+k=n}{\rA}
    \proofstepwidestar[5]{k\in\mathbb{N}}{\rAEa{6}}
    \proofstepwidestar[5]{m+k=n}{\rAEb{6}}
    \proofstepwidestar[1,5]{m+\Succ(k)=\Succ(m+k)}{%
      \FormulaRefAuto{a\in\mathbb{N}\dsep b\in\mathbb{N}\vdash
        a+\Succ(b)=\Succ(a+b)}{1,7}}
    \proofstepwidestar[1,5]{m+\Succ(k)=\Succ(n)}{\rIE{8,9}}
    \proofstepwidestar[1,5]{%
      \exists k\in\mathbb{N}\,(m+\Succ(k)=\Succ(n))}{%
      \rEI{\rAI{7,10}}}
    \proofstepwidestar[1,2,3,5]{m<\Succ(n)}{%
      \FormulaRefAuto{m\in\mathbb{N}\dsep n\in\mathbb{N}\vdash
        m<n\leftrightarrow\exists k\in\mathbb{N}\,(m+\Succ(k)=n)}{1,4,11}}
    \proofstepwidestar[1,2,3]{m<\Succ(n)}{\rEE{5,6,12}}
  \closeproofpartwideindR

  \proofpartwideindR[Rueckrichtung]{%
    m\in\mathbb{N}\dsep n\in\mathbb{N}\dsep m<\Succ(n)
    \vdash m\leq n%
  }{%
    m\in\mathbb{N}\dsep n\in\mathbb{N}\dsep m<\Succ(n)
    \vdash m\leq n%
  }[PeanoLeqSuccStrictBackward]
    \proofstepwidestar[1]{m\in\mathbb{N}}{\rA}
    \proofstepwidestar[2]{n\in\mathbb{N}}{\rA}
    \proofstepwidestar[3]{m<\Succ(n)}{\rA}
    \proofstepwidestar[2]{\Succ(n)\in\mathbb{N}}{%
      \FormulaRefAuto{PeanoSuccClosure}{2}}
    \proofstepwidestar[1,2,3]{%
      \exists k\in\mathbb{N}\,(m+\Succ(k)=\Succ(n))}{%
      \FormulaRefAuto{m\in\mathbb{N}\dsep n\in\mathbb{N}\vdash
        m<n\leftrightarrow\exists k\in\mathbb{N}\,(m+\Succ(k)=n)}{1,4,3}}
    \proofstepwidestar[5]{%
      k\in\mathbb{N}\land m+\Succ(k)=\Succ(n)}{\rA}
    \proofstepwidestar[5]{k\in\mathbb{N}}{\rAEa{6}}
    \proofstepwidestar[5]{m+\Succ(k)=\Succ(n)}{\rAEb{6}}
    \proofstepwidestar[1,5]{m+k\in\mathbb{N}}{%
      \FormulaRefAuto{a\in\mathbb{N}\dsep b\in\mathbb{N}
        \vdash a+b\in\mathbb{N}}{1,7}}
    \proofstepwidestar[1,5]{m+\Succ(k)=\Succ(m+k)}{%
      \FormulaRefAuto{a\in\mathbb{N}\dsep b\in\mathbb{N}\vdash
        a+\Succ(b)=\Succ(a+b)}{1,7}}
    \proofstepwidestar[1,5]{\Succ(m+k)=\Succ(n)}{%
      \FormulaRefAuto{a=b\dsep a=c\vdash b=c}{10,8}}
    \proofstepwidestar[1,2,5]{m+k=n}{%
      \FormulaRefAuto{PeanoSuccInjective}[ax]{9,2,11}}
    \proofstepwidestar[1,2,5]{\exists k\in\mathbb{N}\,(m+k=n)}{%
      \rEI{\rAI{7,12}}}
    \proofstepwidestar[1,2,5]{m\leq n}{%
      \FormulaRefAuto{m\in\mathbb{N}\dsep n\in\mathbb{N}\vdash
        m\leq n\leftrightarrow\exists k\in\mathbb{N}\,(m+k=n)}{1,2,13}}
    \proofstepwidestar[1,2,3]{m\leq n}{\rEE{5,6,14}}
  \closeproofpartwideindR

  \proofpartwide{Schluss}
    \proofstepwidestar[1]{m\in\mathbb{N}}{\rA}
    \proofstepwidestar[2]{n\in\mathbb{N}}{\rA}
    \proofstepwidestar[3]{m\leq n}{\rA}
    \proofstepwidestar[1,2,3]{m<\Succ(n)}{%
      \FormulaRefAuto{PeanoLeqSuccStrictForward}{1,2,3}}
    \proofstepwidestar[1,2]{m\leq n\rightarrow m<\Succ(n)}{\rRI{3,4}}
    \proofstepwidestar[6]{m<\Succ(n)}{\rA}
    \proofstepwidestar[1,2,6]{m\leq n}{%
      \FormulaRefAuto{PeanoLeqSuccStrictBackward}{1,2,6}}
    \proofstepwidestar[1,2]{m<\Succ(n)\rightarrow m\leq n}{\rRI{6,7}}
    \proofstepwidestar[1,2]{m\leq n\leftrightarrow m<\Succ(n)}{%
      \rLRI{5,8}}
  \closeproofpartwide
\end{tabproofsplitwide}

\FormulaThmDeltaK[Nachfolgerfall der strikten Ordnung]{%
  n\in\mathbb{N}\dsep i\in\mathbb{N}\vdash
  i<\Succ(n)\leftrightarrow i<n\lor i=n%
}{PeanoStrictSuccSplit}{%
  \DeltaRow{Mengen}{i\dsep n}%
  \DeltaRow{Relationen}{<_{\mathbb{N}}}%
}
\begin{tabproof}
  \proofstep{1}{n\in\mathbb{N}}{\rA}
  \proofstep{2}{i\in\mathbb{N}}{\rA}
  \proofstep{3}{i<\Succ(n)}{\rA}
  \proofstep{1,2,3}{i\leq n}{%
    \FormulaRefAuto{PeanoLeqSuccStrict}{2,1,3}}
  \proofstep{1,2,3}{i<n\lor i=n}{%
    \FormulaRefAuto{PeanoLeqSplit}{2,1,4}}
  \proofstep{1,2}{i<\Succ(n)\rightarrow i<n\lor i=n}{\rRI{3,5}}
  \proofstep{7}{i<n\lor i=n}{\rA}
  \proofstep{8}{i<n}{\rA}
  \proofstep{1,2,8}{i\leq n}{%
    \FormulaRefAuto{PeanoStrictImpliesLeq}{2,1,8}}
  \proofstep{1,2,8}{i<\Succ(n)}{%
    \FormulaRefAuto{PeanoLeqSuccStrict}{2,1,9}}
  \proofstep{11}{i=n}{\rA}
  \proofstep{1}{n\leq n}{\FormulaRefAuto{PeanoLeqReflexive}{1}}
  \proofstep{1,11}{i\leq n}{\rIE{11,12}}
  \proofstep{1,2,11}{i<\Succ(n)}{%
    \FormulaRefAuto{PeanoLeqSuccStrict}{2,1,13}}
  \proofstep{1,2,7}{i<\Succ(n)}{\rOE{7,8,10,11,14}}
  \proofstep{1,2}{(i<n\lor i=n)\rightarrow i<\Succ(n)}{\rRI{7,15}}
  \proofstep{1,2}{i<\Succ(n)\leftrightarrow i<n\lor i=n}{\rLRI{6,16}}
\end{tabproof}

\FormulaThmDeltaK[Nachfolger-Anfangsabschnitt]{%
  n\in\mathbb{N}\vdash
  \NatSeg{\Succ(n)}=\NatSeg{n}\cup\{n\}%
}{PeanoNatSegSucc}{%
  \DeltaRow{Mengen}{n}%
}
\begin{tabproof}
  \proofstep{1}{n\in\mathbb{N}}{\rA}
  \proofstep{1}{\Succ(n)\in\mathbb{N}}{\FormulaRefAuto{PeanoSuccClosure}{1}}
  \proofstep{1}{\NatSeg{\Succ(n)}=\{i\in\mathbb{N}\mid i<\Succ(n)\}}{%
    \FormulaRefAuto{PeanoNatSeg}{2}}
  \proofstep{1}{\NatSeg{n}=\{i\in\mathbb{N}\mid i<n\}}{%
    \FormulaRefAuto{PeanoNatSeg}{1}}
  \proofstep{1}{\NatSeg{\Succ(n)}=\NatSeg{n}\cup\{n\}}{%
    \FormulaRefAuto{n\in\mathbb{N}\dsep i\in\mathbb{N}\vdash i<\Succ(n)\leftrightarrow i<n\lor i=n}{1}}
\end{tabproof}
\FormulaThmDeltaK[Elemente eines Nachfolger-Anfangsabschnitts]{%
  n\in\mathbb{N}\dsep x\in\NatSeg{\Succ(n)}\vdash
  x\in\NatSeg{n}\lor x=n%
}{PeanoNatSegSuccElim}{%
  \DeltaRow{Mengen}{n\dsep x}%
}
\begin{tabproof}
  \proofstep{1}{n\in\mathbb{N}}{\rA}
  \proofstep{2}{x\in\NatSeg{\Succ(n)}}{\rA}
  \proofstep{1}{\NatSeg{\Succ(n)}=\NatSeg{n}\cup\{n\}}{%
    \FormulaRefAuto{PeanoNatSegSucc}{1}}
  \proofstep{1,2}{x\in\NatSeg{n}\cup\{n\}}{\rIE{3,2}}
  \proofstep{1,2}{x\in\NatSeg{n}\lor x=n}{%
    \FormulaRefAuto{x\in A\cup\{a\} \vdash x\in A\lor x=a}{4}}
\end{tabproof}

\FormulaThmDeltaK[Einbettung des Anfangsabschnitts in den Nachfolger]{%
  n\in\mathbb{N}\vdash \NatSeg{n}\subseteq\NatSeg{\Succ(n)}%
}{PeanoNatSegSubsetSucc}{%
  \DeltaRow{Mengen}{n}%
}
\begin{tabproof}
  \proofstep{1}{n\in\mathbb{N}}{\rA}
  \proofstep{1}{\NatSeg{\Succ(n)}=\NatSeg{n}\cup\{n\}}{%
    \FormulaRefAuto{PeanoNatSegSucc}{1}}
  \proofstep{}{\NatSeg{n}\subseteq\NatSeg{n}\cup\{n\}}{%
    \FormulaRefAuto{A\subseteq A\cup B}}
  \proofstep{1}{\NatSeg{n}\subseteq\NatSeg{\Succ(n)}}{\rIE{2,3}}
\end{tabproof}

\FormulaThmDeltaK[Strikte Ordnung als Mitgliedschaft im Peano-Anfangsabschnitt]{%
  m\in\mathbb{N}\dsep n\in\mathbb{N}\vdash
  m<n \leftrightarrow m\in \NatSeg{n}%
}{PeanoNatSegStrictMembership}{%
  \DeltaRow{Mengen}{m\dsep n}%
  \DeltaRow{Relationen}{<_{\mathbb{N}}}%
}
\begin{tabproof}
  \proofstep{1}{m\in\mathbb{N}}{\rA}
  \proofstep{2}{n\in\mathbb{N}}{\rA}
  \proofstep{2,1}{m\in\NatSeg{n}\leftrightarrow m<n}{%
    \FormulaRefAuto{PeanoNatSegMembership}{2,1}}
  \proofstep{1,2}{m<n \leftrightarrow m\in\NatSeg{n}}{%
    \FormulaRefAuto{P \leftrightarrow Q \vdash Q \leftrightarrow P}{3}}
\end{tabproof}
\FormulaThmDeltaK[Kein Endpunkt im eigenen Anfangsabschnitt]{%
  n\in\mathbb{N}\vdash n\notin \NatSeg{n}%
}{PeanoNatSegEndpointNotIn}{%
  \DeltaRow{Mengen}{n}%
}
\begin{tabproof}
  \proofstep{1}{n\in\mathbb{N}}{\rA}
  \proofstep{2}{n\in\NatSeg{n}}{\rA}
  \proofstep{1,2}{n<n}{%
    \FormulaRefAuto{PeanoNatSegStrictMembership}{1,1,2}}
  \proofstep{1}{n\leq n}{%
    \FormulaRefAuto{PeanoLeqReflexive}{1}}
  \proofstep{1,2}{\bot}{%
    \FormulaRefAuto{PeanoOrderNoCycle}{1,1,3,4}}
  \proofstep{1}{n\notin\NatSeg{n}}{\rCI{2,5}}
\end{tabproof}

\FormulaThmDeltaK[Endpunkt im Nachfolger-Anfangsabschnitt]{%
  n\in\mathbb{N}\vdash n\in \NatSeg{\Succ(n)}%
}{PeanoNatSegEndpointInSucc}{%
  \DeltaRow{Mengen}{n}%
}
\begin{tabproof}
  \proofstep{1}{n\in\mathbb{N}}{\rA}
  \proofstep{1}{n<\Succ(n)}{%
    \FormulaRefAuto{n\in\mathbb{N}\vdash n<\Succ(n)}{1}}
  \proofstep{1}{\Succ(n)\in\mathbb{N}}{%
    \FormulaRefAuto{PeanoSuccClosure}{1}}
  \proofstep{1}{n\in\NatSeg{\Succ(n)}}{%
    \FormulaRefAuto{PeanoNatSegStrictMembership}{1,3,2}}
\end{tabproof}

\FormulaThmDeltaK[Nichtstrikte Ordnung als Nachfolger-Anfangsabschnitt]{%
  m\in\mathbb{N}\dsep n\in\mathbb{N}\vdash
  m\leq n \leftrightarrow m\in \NatSeg{\Succ(n)}%
}{PeanoNatSegLeqSuccMembership}{%
  \DeltaRow{Mengen}{m\dsep n}%
  \DeltaRow{Relationen}{\leq_{\mathbb{N}}}%
}
\begin{tabproof}
  \proofstep{1}{m\in\mathbb{N}}{\rA}
  \proofstep{2}{n\in\mathbb{N}}{\rA}
  \proofstep{2}{\Succ(n)\in\mathbb{N}}{%
    \FormulaRefAuto{PeanoSuccClosure}{2}}
  \proofstep{1,2}{m\leq n\leftrightarrow m<\Succ(n)}{%
    \FormulaRefAuto{PeanoLeqSuccStrict}{1,2}}
  \proofstep{1,2}{m<\Succ(n)\leftrightarrow m\in\NatSeg{\Succ(n)}}{%
    \FormulaRefAuto{PeanoNatSegStrictMembership}{1,3}}
  \proofstep{1,2}{m\leq n \leftrightarrow m\in \NatSeg{\Succ(n)}}{%
    \FormulaRefAuto{P\leftrightarrow Q\dsep Q\leftrightarrow R\vdash P\leftrightarrow R}{4,5}}
\end{tabproof}

\FormulaThmDeltaK[Ordnung liefert Anfangsabschnittsinklusion]{%
  m\in\mathbb{N}\dsep n\in\mathbb{N}\dsep m\leq n
  \vdash \NatSeg{m}\subseteq\NatSeg{n}%
}{PeanoNatSegLeqSubsetForward}{%
  \DeltaRow{Mengen}{m\dsep n}%
  \DeltaRow{Relationen}{\leq_{\mathbb{N}}}%
}
\begin{tabproof}
  \proofstep{1}{m\in\mathbb{N}}{\rA}
  \proofstep{2}{n\in\mathbb{N}}{\rA}
  \proofstep{3}{m\leq n}{\rA}
  \proofstep{4}{x\in\NatSeg{m}}{\rA}
  \proofstep{1}{\NatSeg{m}\subseteq\mathbb{N}}{%
    \FormulaRefAuto{PeanoNatSegSubsetN}{1}}
  \proofstep{1,4}{x\in\mathbb{N}}{%
    \FormulaRefAuto{A\subseteq B\dsep x\in A\vdash x\in B}{5,4}}
  \proofstep{1,4}{x<m}{\FormulaRefAuto{PeanoNatSegStrictMembership}{6,1,4}}
  \proofstep{1,4}{\Succ(x)\leq m}{%
    \FormulaRefAuto{m\in\mathbb{N}\dsep n\in\mathbb{N}\vdash
      m<n\leftrightarrow\Succ(m)\leq n}{6,1,7}}
  \proofstep{1,4}{\Succ(x)\in\mathbb{N}}{\FormulaRefAuto{PeanoSuccClosure}{6}}
  \proofstep{1,2,3,4}{\Succ(x)\leq n}{%
    \FormulaRefAuto{k\in\mathbb{N}\dsep m\in\mathbb{N}\dsep n\in\mathbb{N}
      \dsep k\leq m\dsep m\leq n\vdash k\leq n}{9,1,2,8,3}}
  \proofstep{1,2,3,4}{x<n}{%
    \FormulaRefAuto{m\in\mathbb{N}\dsep n\in\mathbb{N}\vdash
      m<n\leftrightarrow\Succ(m)\leq n}{6,2,10}}
  \proofstep{1,2,3,4}{x\in\NatSeg{n}}{%
    \FormulaRefAuto{PeanoNatSegStrictMembership}{6,2,11}}
  \proofstep{1,2,3}{\NatSeg{m}\subseteq\NatSeg{n}}{%
    \FormulaRefAuto{A\subseteq B\coloneq\forall x\,(x\in A\rightarrow x\in B)}{%
      \rUI{\rRI{4,12}}}}
\end{tabproof}

\FormulaThmDeltaK[Anfangsabschnittsinklusion liefert Ordnung]{%
  m\in\mathbb{N}\dsep n\in\mathbb{N}\dsep
  \NatSeg{m}\subseteq\NatSeg{n}\vdash m\leq n%
}{PeanoNatSegSubsetLeqBackward}{%
  \DeltaRow{Mengen}{m\dsep n}%
  \DeltaRow{Relationen}{\leq_{\mathbb{N}}}%
}
\begin{tabproofsplitwide}
  \proofpartwideindR[Der Gegenvergleich erzwingt Gleichheit]{%
    m\in\mathbb{N}\dsep n\in\mathbb{N}\dsep
    \NatSeg{m}\subseteq\NatSeg{n}\dsep n\leq m\vdash m\leq n
  }{%
    m\in\mathbb{N}\dsep n\in\mathbb{N}\dsep
    \NatSeg{m}\subseteq\NatSeg{n}\dsep n\leq m\vdash m\leq n
  }
    \proofstepwidestar[1]{m\in\mathbb{N}}{\rA}
    \proofstepwidestar[2]{n\in\mathbb{N}}{\rA}
    \proofstepwidestar[3]{\NatSeg{m}\subseteq\NatSeg{n}}{\rA}
    \proofstepwidestar[4]{n\leq m}{\rA}
    \proofstepwidestar[1,2,4]{n<m\lor n=m}{%
      \FormulaRefAuto{PeanoLeqSplit}{2,1,4}}
    \proofstepwidestar[6]{n<m}{\rA}
    \proofstepwidestar[1,2,6]{n\in\NatSeg{m}}{%
      \FormulaRefAuto{PeanoNatSegStrictMembership}{2,1,6}}
    \proofstepwidestar[1,2,3,6]{n\in\NatSeg{n}}{%
      \FormulaRefAuto{A\subseteq B\dsep x\in A\vdash x\in B}{3,7}}
    \proofstepwidestar[2]{n\notin\NatSeg{n}}{%
      \FormulaRefAuto{PeanoNatSegEndpointNotIn}{2}}
    \proofstepwidestar[1,2,3,6]{\bot}{\rBI{8,9}}
    \proofstepwidestar[1,2,3,6]{m\leq n}{\rCE{10}}
    \proofstepwidestar[12]{n=m}{\rA}
    \proofstepwidestar[12]{m=n}{\FormulaRefAuto{a=b\vdash b=a}{12}}
    \proofstepwidestar[1]{m\leq m}{\FormulaRefAuto{PeanoLeqReflexive}{1}}
    \proofstepwidestar[1,12]{m\leq n}{\rIE{13,14}}
    \proofstepwidestar[1,2,3,4]{m\leq n}{\rOE{5,6,11,12,15}}
  \closeproofpartwideindR

  \proofpartwide{Schluss}
    \proofstepwidestar[1]{m\in\mathbb{N}}{\rA}
    \proofstepwidestar[2]{n\in\mathbb{N}}{\rA}
    \proofstepwidestar[3]{\NatSeg{m}\subseteq\NatSeg{n}}{\rA}
    \proofstepwidestar[1,2]{m\leq n\lor n\leq m}{%
      \FormulaRefAuto{PeanoOrderComparison}{1,2}}
    \proofstepwidestar[5]{m\leq n}{\rA}
    \proofstepwidestar[6]{n\leq m}{\rA}
    \proofstepwidestar[1,2,3,6]{m\leq n}{%
      \FormulaRefAuto{m\in\mathbb{N}\dsep n\in\mathbb{N}\dsep
        \NatSeg{m}\subseteq\NatSeg{n}\dsep n\leq m\vdash m\leq n}{1,2,3,6}}
    \proofstepwidestar[1,2,3]{m\leq n}{\rOE{4,5,5,6,7}}
  \closeproofpartwide
\end{tabproofsplitwide}

\FormulaThmDeltaK[Ordnung und Anfangsabschnittsinklusion]{%
  m\in\mathbb{N}\dsep n\in\mathbb{N}\vdash
  m\leq n \leftrightarrow \NatSeg{m}\subseteq \NatSeg{n}%
}{PeanoNatSegLeqSubset}{%
  \DeltaRow{Mengen}{m\dsep n}%
  \DeltaRow{Relationen}{\leq_{\mathbb{N}}}%
}
\begin{tabproof}
  \proofstep{1}{m\in\mathbb{N}}{\rA}
  \proofstep{2}{n\in\mathbb{N}}{\rA}
  \proofstep{3}{m\leq n}{\rA}
  \proofstep{1,2,3}{\NatSeg{m}\subseteq\NatSeg{n}}{%
    \FormulaRefAuto{PeanoNatSegLeqSubsetForward}{1,2,3}}
  \proofstep{1,2}{m\leq n\rightarrow\NatSeg{m}\subseteq\NatSeg{n}}{\rRI{3,4}}
  \proofstep{6}{\NatSeg{m}\subseteq\NatSeg{n}}{\rA}
  \proofstep{1,2,6}{m\leq n}{%
    \FormulaRefAuto{PeanoNatSegSubsetLeqBackward}{1,2,6}}
  \proofstep{1,2}{\NatSeg{m}\subseteq\NatSeg{n}\rightarrow m\leq n}{\rRI{6,7}}
  \proofstep{1,2}{m\leq n\leftrightarrow\NatSeg{m}\subseteq\NatSeg{n}}{\rLRI{5,8}}
\end{tabproof}

\FormulaThmDeltaK[Strikte Ordnung liefert Anfangsabschnittsinklusion]{%
  m\in\mathbb{N}\dsep n\in\mathbb{N}\dsep m\in \NatSeg{n}
  \vdash \NatSeg{m}\subseteq \NatSeg{n}%
}{PeanoNatSegMembershipSubset}{%
  \DeltaRow{Mengen}{m\dsep n}%
}
\begin{tabproof}
  \proofstep{1}{m\in\mathbb{N}}{\rA}
  \proofstep{2}{n\in\mathbb{N}}{\rA}
  \proofstep{3}{m\in\NatSeg{n}}{\rA}
  \proofstep{1,2,3}{m<n}{%
    \FormulaRefAuto{PeanoNatSegStrictMembership}{1,2,3}}
  \proofstep{1,2,3}{m\leq n}{%
    \FormulaRefAuto{m\in\mathbb{N}\dsep n\in\mathbb{N}\dsep m<n\vdash m\leq n}{1,2,4}}
  \proofstep{1,2,3}{\NatSeg{m}\subseteq\NatSeg{n}}{%
    \FormulaRefAuto{PeanoNatSegLeqSubset}{1,2,5}}
\end{tabproof}

\FormulaThmDeltaK[Teilmengen eines Nachfolger-Anfangsabschnitts ohne Endpunkt]{%
  n\in\mathbb{N}\dsep B\subseteq \NatSeg{\Succ(n)}\dsep n\notin B
  \vdash B\subseteq \NatSeg{n}%
}{PeanoNatSegSubsetSuccWithoutEndpoint}{%
  \DeltaRow{Mengen}{B\dsep n}%
}
\begin{tabproof}
  \proofstep{1}{n\in\mathbb{N}}{\rA}
  \proofstep{2}{B\subseteq \NatSeg{\Succ(n)}}{\rA}
  \proofstep{3}{n\notin B}{\rA}
  \proofstep{4}{x\in B}{\rA}
  \proofstep{2,4}{x\in\NatSeg{\Succ(n)}}{%
    \FormulaRefAuto{A\subseteq B,\, x\in A\vdash x\in B}{2,4}}
  \proofstep{1,2,4}{x\in\NatSeg{n}\lor x=n}{%
    \FormulaRefAuto{PeanoNatSegSuccElim}{1,5}}
  \proofstep{7}{x\in\NatSeg{n}}{\rA}
  \proofstep{8}{x=n}{\rA}
  \proofstep{3,4,8}{\bot}{\rBI{\rIE{8,4},3}}
  \proofstep{3,4,8}{x\in\NatSeg{n}}{\FormulaRefAuto{P,\,\neg P \vdash Q}{4,9}}
  \proofstep{1,2,3,4}{x\in\NatSeg{n}}{\rOE{6,7,7,8,10}}
  \proofstep{1,2,3}{B\subseteq\NatSeg{n}}{%
    \FormulaRefAuto{A \subseteq B \coloneq \forall x\,(x\in A \rightarrow x\in B)}{\rUI{\rRI{4,11}}}}
\end{tabproof}

\FormulaThmDeltaK[Teilmengen eines Nachfolger-Anfangsabschnitts mit Endpunkt]{%
  n\in\mathbb{N}\dsep B\subseteq \NatSeg{\Succ(n)}\dsep n\in B
  \vdash B=(B\cap \NatSeg{n})\cup\{n\}%
}{PeanoNatSegSubsetSuccWithEndpoint}{%
  \DeltaRow{Mengen}{B\dsep n}%
}
\begin{tabproofsplitwide}
  \proofpartwideindR[Hinrichtung]{%
    n\in\mathbb{N}\dsep B\subseteq\NatSeg{\Succ(n)}\dsep n\in B
    \vdash B\subseteq(B\cap\NatSeg{n})\cup\{n\}
  }{%
    n\in\mathbb{N}\dsep B\subseteq\NatSeg{\Succ(n)}\dsep n\in B
    \vdash B\subseteq(B\cap\NatSeg{n})\cup\{n\}
  }
    \proofstepwidestar[1]{n\in\mathbb{N}}{\rA}
    \proofstepwidestar[2]{B\subseteq\NatSeg{\Succ(n)}}{\rA}
    \proofstepwidestar[3]{n\in B}{\rA}
    \proofstepwidestar[4]{x\in B}{\rA}
    \proofstepwidestar[2,4]{x\in\NatSeg{\Succ(n)}}{%
      \FormulaRefAuto{A\subseteq B\dsep x\in A\vdash x\in B}{2,4}}
    \proofstepwidestar[1,2,4]{x\in\NatSeg{n}\lor x=n}{%
      \FormulaRefAuto{PeanoNatSegSuccElim}{1,5}}
    \proofstepwidestar[7]{x\in\NatSeg{n}}{\rA}
    \proofstepwidestar[4,7]{x\in B\cap\NatSeg{n}}{%
      \FormulaRefAuto{x\in A\cap B\eqvdash x\in A\land x\in B}{\rAI{4,7}}}
    \proofstepwidestar[4,7]{x\in(B\cap\NatSeg{n})\cup\{n\}}{%
      \FormulaRefAuto{x\in A\vdash x\in A\cup B}{8}}
    \proofstepwidestar[10]{x=n}{\rA}
    \proofstepwidestar[10]{x\in\{n\}}{%
      \FormulaRefAuto{x\in\{a\}\eqvdash x=a}{10}}
    \proofstepwidestar[10]{x\in(B\cap\NatSeg{n})\cup\{n\}}{%
      \FormulaRefAuto{x\in B\vdash x\in A\cup B}{11}}
    \proofstepwidestar[1,2,4]{x\in(B\cap\NatSeg{n})\cup\{n\}}{%
      \rOE{6,7,9,10,12}}
    \proofstepwidestar[1,2,3]{B\subseteq(B\cap\NatSeg{n})\cup\{n\}}{%
      \FormulaRefAuto{A\subseteq B\coloneq\forall x\,(x\in A\rightarrow x\in B)}{%
        \rUI{\rRI{4,13}}}}
  \closeproofpartwideindR

  \proofpartwideindR[Rückrichtung]{%
    n\in\mathbb{N}\dsep B\subseteq\NatSeg{\Succ(n)}\dsep n\in B
    \vdash(B\cap\NatSeg{n})\cup\{n\}\subseteq B
  }{%
    n\in\mathbb{N}\dsep B\subseteq\NatSeg{\Succ(n)}\dsep n\in B
    \vdash(B\cap\NatSeg{n})\cup\{n\}\subseteq B
  }
    \proofstepwidestar[1]{n\in\mathbb{N}}{\rA}
    \proofstepwidestar[2]{B\subseteq\NatSeg{\Succ(n)}}{\rA}
    \proofstepwidestar[3]{n\in B}{\rA}
    \proofstepwidestar[4]{x\in(B\cap\NatSeg{n})\cup\{n\}}{\rA}
    \proofstepwidestar[4]{x\in B\cap\NatSeg{n}\lor x\in\{n\}}{%
      \FormulaRefAuto{x\in A\cup B\eqvdash x\in A\lor x\in B}{4}}
    \proofstepwidestar[6]{x\in B\cap\NatSeg{n}}{\rA}
    \proofstepwidestar[6]{x\in B}{\FormulaRefAuto{x\in A\cap B\vdash x\in A}{6}}
    \proofstepwidestar[8]{x\in\{n\}}{\rA}
    \proofstepwidestar[8]{x=n}{\FormulaRefAuto{x\in\{a\}\eqvdash x=a}{8}}
    \proofstepwidestar[3,8]{x\in B}{\rIE{9,3}}
    \proofstepwidestar[3,4]{x\in B}{\rOE{5,6,7,8,10}}
    \proofstepwidestar[1,2,3]{(B\cap\NatSeg{n})\cup\{n\}\subseteq B}{%
      \FormulaRefAuto{A\subseteq B\coloneq\forall x\,(x\in A\rightarrow x\in B)}{%
        \rUI{\rRI{4,11}}}}
  \closeproofpartwideindR

  \proofpartwide{Schluss}
    \proofstepwidestar[1]{n\in\mathbb{N}}{\rA}
    \proofstepwidestar[2]{B\subseteq\NatSeg{\Succ(n)}}{\rA}
    \proofstepwidestar[3]{n\in B}{\rA}
    \proofstepwidestar[1,2,3]{B\subseteq(B\cap\NatSeg{n})\cup\{n\}}{%
      \FormulaRefAuto{n\in\mathbb{N}\dsep B\subseteq\NatSeg{\Succ(n)}\dsep n\in B
        \vdash B\subseteq(B\cap\NatSeg{n})\cup\{n\}}{1,2,3}}
    \proofstepwidestar[1,2,3]{(B\cap\NatSeg{n})\cup\{n\}\subseteq B}{%
      \FormulaRefAuto{n\in\mathbb{N}\dsep B\subseteq\NatSeg{\Succ(n)}\dsep n\in B
        \vdash(B\cap\NatSeg{n})\cup\{n\}\subseteq B}{1,2,3}}
    \proofstepwidestar[1,2,3]{B=(B\cap\NatSeg{n})\cup\{n\}}{%
      \FormulaRefAuto{A\subseteq B\dsep B\subseteq A\vdash A=B}{4,5}}
  \closeproofpartwide
\end{tabproofsplitwide}

\subsection{Auslassungsabbildung und Injektionen in Anfangsabschnitte}

Die allgemeine Folgenverschiebung aus dem Rekursionsabschnitt wird nun auf
die Translation $k\mapsto x_0+k$ angewandt.  Sie spielt hier die Rolle der
Peano-Auslassungsabbildung: Sie ist injektiv und lässt ihren Anfangspunkt
$x_0$ aus.

\FormulaThmDeltaK[Die verschobene Translation lässt ihren Anfangspunkt aus]{%
  x_0\in\mathbb{N}\dsep x\in\mathbb{N}\vdash
  \SuccShift{\RecFun{x_0}{\Succ}}(x)\neq x_0%
}{PeanoTranslateShiftOmitStart}{%
  \DeltaRow{Mengen}{x_0\dsep x}%
  \DeltaRow{Funktionensymbole}{\RecFun{x_0}{\Succ}\dsep
    \SuccShift{\RecFun{x_0}{\Succ}}}%
}
\begin{tabproof}
  \proofstep{1}{x_0\in\mathbb{N}}{\rA}
  \proofstep{2}{x\in\mathbb{N}}{\rA}
  \proofstep{1}{\RecFun{x_0}{\Succ}\colon\mathbb{N}\inj\mathbb{N}}{%
    \FormulaRefAuto{PeanoTranslateInjection}{1}}
  \proofstep{}{\Succ\colon\mathbb{N}\to\mathbb{N}}{\FormulaRefAuto{PeanoSuccFunction}[ax]}
  \proofstep{1}{\RecFun{x_0}{\Succ}(0)=x_0}{%
    \FormulaRefAuto{m_0\in M\dsep f\colon M\to M\vdash
      \RecFun{m_0}{f}(0)=m_0}{1,4}}
  \proofstep{6}{\SuccShift{\RecFun{x_0}{\Succ}}(x)=x_0}{\rA}
  \proofstep{1}{x_0=\RecFun{x_0}{\Succ}(0)}{\FormulaRefAuto{a=b\vdash b=a}{5}}
  \proofstep{1,2,6}{%
    \SuccShift{\RecFun{x_0}{\Succ}}(x)=\RecFun{x_0}{\Succ}(0)}{%
    \FormulaRefAuto{a=b\dsep b=c\vdash a=c}{6,7}}
  \proofstep{1,2,6}{\bot}{%
    \FormulaRefAuto{H\colon\mathbb{N}\inj M\dsep x\in M\dsep
      \SuccShift{H}(x)=H(0)\vdash\bot}{3,2,8}}
  \proofstep{1,2}{\SuccShift{\RecFun{x_0}{\Succ}}(x)\neq x_0}{\rCI{6,9}}
\end{tabproof}

\FormulaThmDeltaK[Die Auslassungskomposition trifft den Endpunkt nicht]{%
  \begin{aligned}[t]
    &n\in\mathbb{N}\dsep F\colon\mathbb{N}\inj\NatSeg{\Succ(n)}
      \dsep x_0\in\mathbb{N}\dsep F(x_0)=n\dsep x\in\mathbb{N}\vdash{}\\
    &\bigl(F\circ\SuccShift{\RecFun{x_0}{\Succ}}\bigr)(x)\neq n
  \end{aligned}%
}{PeanoOmissionCompositeAvoidsEndpoint}{%
  \DeltaRow{Mengen}{n\dsep x_0\dsep x}%
  \DeltaRow{Funktionen}{F}%
  \DeltaRow{Funktionensymbole}{\RecFun{x_0}{\Succ}\dsep
    \SuccShift{\RecFun{x_0}{\Succ}}}%
}
\begin{tabproofwide}
  \proofstepwidestar[1]{n\in\mathbb{N}}{\rA}
  \proofstepwidestar[2]{F\colon\mathbb{N}\inj\NatSeg{\Succ(n)}}{\rA}
  \proofstepwidestar[3]{x_0\in\mathbb{N}}{\rA}
  \proofstepwidestar[4]{F(x_0)=n}{\rA}
  \proofstepwidestar[5]{x\in\mathbb{N}}{\rA}
  \proofstepwidestar[3]{\RecFun{x_0}{\Succ}\colon\mathbb{N}\inj\mathbb{N}}{%
    \FormulaRefAuto{PeanoTranslateInjection}{3}}
  \proofstepwidestar[3]{\SuccShift{\RecFun{x_0}{\Succ}}\colon
    \mathbb{N}\inj\mathbb{N}}{%
    \FormulaRefAuto{H\colon\mathbb{N}\inj M\vdash\SuccShift{H}\colon M\inj M}{6}}
  \proofstepwidestar[3,5]{\SuccShift{\RecFun{x_0}{\Succ}}(x)\in\mathbb{N}}{%
    \FormulaRefAuto{F\colon A\inj B\dsep x\in A\vdash F(x)\in B}{7,5}}
  \proofstepwidestar[3,5]{\SuccShift{\RecFun{x_0}{\Succ}}(x)\neq x_0}{%
    \FormulaRefAuto{PeanoTranslateShiftOmitStart}{3,5}}
  \proofstepwidestar[10]{%
    \bigl(F\circ\SuccShift{\RecFun{x_0}{\Succ}}\bigr)(x)=n}{\rA}
  \proofstepwidestar[5]{%
    \bigl(F\circ\SuccShift{\RecFun{x_0}{\Succ}}\bigr)(x)
      =F\bigl(\SuccShift{\RecFun{x_0}{\Succ}}(x)\bigr)}{%
    \FormulaRefAuto{x\in A\vdash(G\circ F)(x)=G(F(x))}{5}}
  \proofstepwidestar[5,10]{F\bigl(\SuccShift{\RecFun{x_0}{\Succ}}(x)\bigr)=n}{%
    \FormulaRefAuto{a=b\dsep a=c\vdash b=c}{11,10}}
  \proofstepwidestar[4]{n=F(x_0)}{\FormulaRefAuto{a=b\vdash b=a}{4}}
  \proofstepwidestar[4,5,10]{%
    F\bigl(\SuccShift{\RecFun{x_0}{\Succ}}(x)\bigr)=F(x_0)}{%
    \FormulaRefAuto{a=b\dsep b=c\vdash a=c}{12,13}}
  \proofstepwidestar[2,3,4,5,10]{\SuccShift{\RecFun{x_0}{\Succ}}(x)=x_0}{%
    \FormulaRefAuto{F\colon A\inj B\dsep x\in A\dsep y\in A\dsep
      F(x)=F(y)\vdash x=y}{2,8,3,14}}
  \proofstepwidestar[2,3,4,5,10]{\bot}{\rBI{15,9}}
  \proofstepwidestar[2,3,4,5]{%
    \bigl(F\circ\SuccShift{\RecFun{x_0}{\Succ}}\bigr)(x)\neq n}{\rCI{10,16}}
\end{tabproofwide}

\FormulaThmDeltaK[Werte der Auslassungskomposition liegen im kleineren Abschnitt]{%
  \begin{aligned}[t]
    &n\in\mathbb{N}\dsep F\colon\mathbb{N}\inj\NatSeg{\Succ(n)}
      \dsep x_0\in\mathbb{N}\dsep F(x_0)=n\dsep x\in\mathbb{N}\vdash{}\\
    &\bigl(F\circ\SuccShift{\RecFun{x_0}{\Succ}}\bigr)(x)\in\NatSeg{n}
  \end{aligned}%
}{PeanoOmissionCompositeRange}{%
  \DeltaRow{Mengen}{n\dsep x_0\dsep x}%
  \DeltaRow{Funktionen}{F}%
}
\begin{tabproofwide}
  \proofstepwidestar[1]{n\in\mathbb{N}}{\rA}
  \proofstepwidestar[2]{F\colon\mathbb{N}\inj\NatSeg{\Succ(n)}}{\rA}
  \proofstepwidestar[3]{x_0\in\mathbb{N}}{\rA}
  \proofstepwidestar[4]{F(x_0)=n}{\rA}
  \proofstepwidestar[5]{x\in\mathbb{N}}{\rA}
  \proofstepwidestar[3]{\RecFun{x_0}{\Succ}\colon\mathbb{N}\inj\mathbb{N}}{%
    \FormulaRefAuto{PeanoTranslateInjection}{3}}
  \proofstepwidestar[3]{\SuccShift{\RecFun{x_0}{\Succ}}\colon
    \mathbb{N}\inj\mathbb{N}}{%
    \FormulaRefAuto{H\colon\mathbb{N}\inj M\vdash\SuccShift{H}\colon M\inj M}{6}}
  \proofstepwidestar[2,3]{F\circ\SuccShift{\RecFun{x_0}{\Succ}}\colon
    \mathbb{N}\inj\NatSeg{\Succ(n)}}{%
    \FormulaRefAuto{F\colon A\inj B\dsep G\colon B\inj C\vdash
      G\circ F\colon A\inj C}{7,2}}
  \proofstepwidestar[2,3,5]{%
    \bigl(F\circ\SuccShift{\RecFun{x_0}{\Succ}}\bigr)(x)
      \in\NatSeg{\Succ(n)}}{%
    \FormulaRefAuto{F\colon A\inj B\dsep x\in A\vdash F(x)\in B}{8,5}}
  \proofstepwidestar[1,2,3,5]{%
    \bigl(F\circ\SuccShift{\RecFun{x_0}{\Succ}}\bigr)(x)\in\NatSeg{n}
    \lor \bigl(F\circ\SuccShift{\RecFun{x_0}{\Succ}}\bigr)(x)=n}{%
    \FormulaRefAuto{PeanoNatSegSuccElim}{1,9}}
  \proofstepwidestar[1,2,3,4,5]{%
    \bigl(F\circ\SuccShift{\RecFun{x_0}{\Succ}}\bigr)(x)\neq n}{%
    \FormulaRefAuto{PeanoOmissionCompositeAvoidsEndpoint}{1,2,3,4,5}}
  \proofstepwidestar[1,2,3,4,5]{%
    \bigl(F\circ\SuccShift{\RecFun{x_0}{\Succ}}\bigr)(x)\in\NatSeg{n}}{%
    \FormulaRefAuto{P\lor Q\dsep\neg Q\vdash P}{10,11}}
\end{tabproofwide}

\FormulaThmDeltaK[Auslassen eines Endpunkturbildes]{%
  \begin{aligned}[t]
    &n\in\mathbb{N}\dsep F\colon\mathbb{N}\inj\NatSeg{\Succ(n)}
      \dsep x_0\in\mathbb{N}\dsep F(x_0)=n\vdash{}\\
    &F\circ\SuccShift{\RecFun{x_0}{\Succ}}\colon
      \mathbb{N}\inj\NatSeg{n}
  \end{aligned}%
}{PeanoNatSegOmitEndpointPreimage}{%
  \DeltaRow{Mengen}{n\dsep x_0\dsep x}%
  \DeltaRow{Funktionen}{F}%
}
\begin{tabproofwide}
  \proofstepwidestar[1]{n\in\mathbb{N}}{\rA}
  \proofstepwidestar[2]{F\colon\mathbb{N}\inj\NatSeg{\Succ(n)}}{\rA}
  \proofstepwidestar[3]{x_0\in\mathbb{N}}{\rA}
  \proofstepwidestar[4]{F(x_0)=n}{\rA}
  \proofstepwidestar[3]{\RecFun{x_0}{\Succ}\colon\mathbb{N}\inj\mathbb{N}}{%
    \FormulaRefAuto{PeanoTranslateInjection}{3}}
  \proofstepwidestar[3]{\SuccShift{\RecFun{x_0}{\Succ}}\colon
    \mathbb{N}\inj\mathbb{N}}{%
    \FormulaRefAuto{H\colon\mathbb{N}\inj M\vdash\SuccShift{H}\colon M\inj M}{5}}
  \proofstepwidestar[2,3]{F\circ\SuccShift{\RecFun{x_0}{\Succ}}\colon
    \mathbb{N}\inj\NatSeg{\Succ(n)}}{%
    \FormulaRefAuto{F\colon A\inj B\dsep G\colon B\inj C\vdash
      G\circ F\colon A\inj C}{6,2}}
  \proofstepwidestar[8]{x\in\mathbb{N}}{\rA}
  \proofstepwidestar[1,2,3,4,8]{%
    \bigl(F\circ\SuccShift{\RecFun{x_0}{\Succ}}\bigr)(x)\in\NatSeg{n}}{%
    \FormulaRefAuto{PeanoOmissionCompositeRange}{1,2,3,4,8}}
  \proofstepwidestar[1,2,3,4]{\forall x\in\mathbb{N}\,
    \bigl(F\circ\SuccShift{\RecFun{x_0}{\Succ}}\bigr)(x)\in\NatSeg{n}}{%
    \rUI{\rRI{8,9}}}
  \proofstepwidestar[1,2,3,4]{F\circ\SuccShift{\RecFun{x_0}{\Succ}}\colon
    \mathbb{N}\inj\NatSeg{n}}{%
    \FormulaRefAuto{F\colon A\inj B\dsep\forall x\in A\,F(x)\in C
      \vdash F\colon A\inj C}{7,10}}
\end{tabproofwide}

\FormulaThmDeltaK[Injektion ohne Endpunktbild reduziert sich auf den Vorgängerabschnitt]{%
  n\in\mathbb{N}\dsep F\colon\mathbb{N}\inj\NatSeg{\Succ(n)}
  \dsep n\notin F[\mathbb{N}]\vdash F\colon\mathbb{N}\inj\NatSeg{n}%
}{PeanoNatSegAvoidEndpointReduction}{%
  \DeltaRow{Mengen}{n\dsep x}%
  \DeltaRow{Funktionen}{F}%
}
\begin{tabproof}
  \proofstep{1}{n\in\mathbb{N}}{\rA}
  \proofstep{2}{F\colon\mathbb{N}\inj\NatSeg{\Succ(n)}}{\rA}
  \proofstep{3}{n\notin F[\mathbb{N}]}{\rA}
  \proofstep{4}{x\in\mathbb{N}}{\rA}
  \proofstep{2,4}{F(x)\in\NatSeg{\Succ(n)}}{%
    \FormulaRefAuto{F\colon A\inj B\dsep x\in A\vdash F(x)\in B}{2,4}}
  \proofstep{1,2,4}{F(x)\in\NatSeg{n}\lor F(x)=n}{%
    \FormulaRefAuto{PeanoNatSegSuccElim}{1,5}}
  \proofstep{7}{F(x)=n}{\rA}
  \proofstep{2,4}{F(x)\in F[\mathbb{N}]}{%
    \FormulaRefAuto{F\colon A\inj B\dsep x\in A\vdash F(x)\in F[A]}{2,4}}
  \proofstep{2,4,7}{n\in F[\mathbb{N}]}{\rIE{7,8}}
  \proofstep{2,3,4,7}{\bot}{\rBI{9,3}}
  \proofstep{2,3,4}{F(x)\neq n}{\rCI{7,10}}
  \proofstep{1,2,3,4}{F(x)\in\NatSeg{n}}{%
    \FormulaRefAuto{P\lor Q\dsep\neg Q\vdash P}{6,11}}
  \proofstep{1,2,3}{\forall x\in\mathbb{N}\,F(x)\in\NatSeg{n}}{%
    \rUI{\rRI{4,12}}}
  \proofstep{1,2,3}{F\colon\mathbb{N}\inj\NatSeg{n}}{%
    \FormulaRefAuto{F\colon A\inj B\dsep\forall x\in A\,F(x)\in C
      \vdash F\colon A\inj C}{2,13}}
\end{tabproof}

\FormulaThmDeltaK[Endpunktbild liefert eine Injektion in den Vorgängerabschnitt]{%
  n\in\mathbb{N}\dsep F\colon\mathbb{N}\inj\NatSeg{\Succ(n)}
  \dsep n\in F[\mathbb{N}]\vdash
  \exists G\colon\mathbb{N}\inj\NatSeg{n}%
}{PeanoNatSegEndpointImageReduction}{%
  \DeltaRow{Mengen}{n\dsep x_0}%
  \DeltaRow{Funktionen}{F\dsep G}%
}
\begin{tabproof}
  \proofstep{1}{n\in\mathbb{N}}{\rA}
  \proofstep{2}{F\colon\mathbb{N}\inj\NatSeg{\Succ(n)}}{\rA}
  \proofstep{3}{n\in F[\mathbb{N}]}{\rA}
  \proofstep{2}{F\colon\mathbb{N}\to\NatSeg{\Succ(n)}}{%
    \FormulaRefAuto{Injektive Funktion}{2}}
  \proofstep{2,3}{\exists x_0\in\mathbb{N}\,n=F(x_0)}{%
    \FormulaRefAuto{F\colon A\to B\dsep y\in F[A]\vdash
      \exists x\in A\,y=F(x)}{4,3}}
  \proofstep{6}{x_0\in\mathbb{N}\land n=F(x_0)}{\rA}
  \proofstep{6}{x_0\in\mathbb{N}}{\rAEa{6}}
  \proofstep{6}{n=F(x_0)}{\rAEb{6}}
  \proofstep{6}{F(x_0)=n}{\FormulaRefAuto{a=b\vdash b=a}{8}}
  \proofstep{1,2,6}{F\circ\SuccShift{\RecFun{x_0}{\Succ}}\colon
    \mathbb{N}\inj\NatSeg{n}}{%
    \FormulaRefAuto{PeanoNatSegOmitEndpointPreimage}{1,2,7,9}}
  \proofstep{1,2,6}{\exists G\colon\mathbb{N}\inj\NatSeg{n}}{\rEI{10}}
  \proofstep{1,2,3}{\exists G\colon\mathbb{N}\inj\NatSeg{n}}{\rEE{5,6,11}}
\end{tabproof}

\FormulaThmDeltaK[Keine Injektion von \(\mathbb{N}\) in einen Anfangsabschnitt]{%
  n\in\mathbb{N}\vdash \neg\exists F\colon\mathbb{N}\inj\NatSeg{n}%
}{PeanoNatSegNoNatInjection}{%
  \DeltaRow{Mengen}{n}%
  \DeltaRow{Funktionen}{F\dsep G}%
}
\begin{tabproofsplitwide}
  \proofpartwideindR[Induktionsanfang]{%
    \vdash\neg\exists F\colon\mathbb{N}\inj\NatSeg{0}
  }{%
    \vdash\neg\exists F\colon\mathbb{N}\inj\NatSeg{0}
  }[PeanoNatSegNoNatInjectionBase]
    \proofstepwidestar[1]{\exists F\colon\mathbb{N}\inj\NatSeg{0}}{\rA}
    \proofstepwidestar[2]{F\colon\mathbb{N}\inj\NatSeg{0}}{\rA}
    \proofstepwidestar[]{0\in\mathbb{N}}{\FormulaRefAuto{PeanoZero}[ax]}
    \proofstepwidestar[2]{F(0)\in\NatSeg{0}}{%
      \FormulaRefAuto{F\colon A\inj B\dsep x\in A\vdash F(x)\in B}{2,3}}
    \proofstepwidestar[]{\NatSeg{0}=\varnothing}{\FormulaRefAuto{PeanoNatSegZero}}
    \proofstepwidestar[2]{F(0)\in\varnothing}{\rIE{5,4}}
    \proofstepwidestar[]{F(0)\notin\varnothing}{\FormulaRefAuto{x\not\in\varnothing}}
    \proofstepwidestar[2]{\bot}{\rBI{6,7}}
    \proofstepwidestar[1]{\bot}{\rEE{1,2,8}}
    \proofstepwidestar[]{\neg\exists F\colon\mathbb{N}\inj\NatSeg{0}}{\rCI{1,9}}
  \closeproofpartwideindR

  \proofpartwideindR[Induktionsschritt]{%
    n\in\mathbb{N}\dsep\neg\exists G\colon\mathbb{N}\inj\NatSeg{n}
    \vdash\neg\exists F\colon\mathbb{N}\inj\NatSeg{\Succ(n)}
  }{%
    n\in\mathbb{N}\dsep\neg\exists G\colon\mathbb{N}\inj\NatSeg{n}
    \vdash\neg\exists F\colon\mathbb{N}\inj\NatSeg{\Succ(n)}
  }[PeanoNatSegNoNatInjectionStep]
    \proofstepwidestar[1]{n\in\mathbb{N}}{\rA}
    \proofstepwidestar[2]{\neg\exists G\colon\mathbb{N}\inj\NatSeg{n}}{\rA}
    \proofstepwidestar[3]{\exists F\colon\mathbb{N}\inj\NatSeg{\Succ(n)}}{\rA}
    \proofstepwidestar[4]{F\colon\mathbb{N}\inj\NatSeg{\Succ(n)}}{\rA}
    \proofstepwidestar[]{n\in F[\mathbb{N}]\lor n\notin F[\mathbb{N}]}{%
      \FormulaRefAuto{P\lor\neg P}}
    \proofstepwidestar[6]{n\in F[\mathbb{N}]}{\rA}
    \proofstepwidestar[1,4,6]{\exists G\colon\mathbb{N}\inj\NatSeg{n}}{%
      \FormulaRefAuto{PeanoNatSegEndpointImageReduction}{1,4,6}}
    \proofstepwidestar[1,2,4,6]{\bot}{\rBI{7,2}}
    \proofstepwidestar[9]{n\notin F[\mathbb{N}]}{\rA}
    \proofstepwidestar[1,4,9]{F\colon\mathbb{N}\inj\NatSeg{n}}{%
      \FormulaRefAuto{PeanoNatSegAvoidEndpointReduction}{1,4,9}}
    \proofstepwidestar[1,4,9]{\exists G\colon\mathbb{N}\inj\NatSeg{n}}{\rEI{10}}
    \proofstepwidestar[1,2,4,9]{\bot}{\rBI{11,2}}
    \proofstepwidestar[1,2,4]{\bot}{\rOE{5,6,8,9,12}}
    \proofstepwidestar[1,2,3]{\bot}{\rEE{3,4,13}}
    \proofstepwidestar[1,2]{\neg\exists F\colon\mathbb{N}\inj\NatSeg{\Succ(n)}}{%
      \rCI{3,14}}
  \closeproofpartwideindR

  \proofpartwide{Induktionsschluss}
    \proofstepwidestar[1]{n\in\mathbb{N}}{\rA}
    \proofstepwidestar[1]{\neg\exists F\colon\mathbb{N}\inj\NatSeg{n}}{%
      \FormulaRefAuto{PeanoInductionRule}{%
        PeanoNatSegNoNatInjectionBase,PeanoNatSegNoNatInjectionStep,1}}
  \closeproofpartwide
\end{tabproofsplitwide}


\FormulaThmDeltaK[Schrittweise Inklusionen induzieren Monotonie]{%
  \forall u\in\mathbb{N}\,\bigl(F(u)\subseteq F(\Succ(u))\bigr)\dsep
  k\in\mathbb{N}\dsep n\in\mathbb{N}\dsep k<n
  \vdash F(\Succ(k))\subseteq F(n)%
}{B04PeanoStrictMonotoneSucc}{%
  \DeltaRow{Mengen}{k\dsep n\dsep u}%
  \DeltaRow{Funktionen}{F}%
}
\begin{tabproofsplitwide}
  \proofpartwideindR[Induktionsanfang]{%
    \forall u\in\mathbb{N}\,\bigl(F(u)\subseteq F(\Succ(u))\bigr)
    \vdash
    \forall k\in\mathbb{N}\,\bigl(k<0\rightarrow
      F(\Succ(k))\subseteq F(0)\bigr)%
  }{%
    \forall u\in\mathbb{N}\,\bigl(F(u)\subseteq F(\Succ(u))\bigr)
    \vdash
    \forall k\in\mathbb{N}\,\bigl(k<0\rightarrow
      F(\Succ(k))\subseteq F(0)\bigr)%
  }[B04PeanoStrictMonotoneSuccBase]
    \proofstepwidestar[1]{%
      \forall u\in\mathbb{N}\,\bigl(F(u)\subseteq F(\Succ(u))\bigr)}{\rA}
    \proofstepwidestar[2]{k\in\mathbb{N}}{\rA}
    \proofstepwidestar[3]{k<0}{\rA}
    \proofstepwidestar[2]{\neg(k<0)}{%
      \FormulaRefAuto{PeanoNoLtZero}[thm]{2}}
    \proofstepwidestar[2,3]{F(\Succ(k))\subseteq F(0)}{%
      \FormulaRefAuto{P,\,\neg P\vdash Q}{3,4}}
    \proofstepwidestar[2]{%
      k<0\rightarrow F(\Succ(k))\subseteq F(0)}{\rRI{3,5}}
    \proofstepwidestar[]{%
      \forall k\in\mathbb{N}\,\bigl(k<0\rightarrow
        F(\Succ(k))\subseteq F(0)\bigr)}{\rUI{\rRI{2,6}}}
  \closeproofpartwideindR

  \proofpartwideindR[Induktionsschritt]{%
    \begin{aligned}[t]
      &\forall v\in\mathbb{N}\,\bigl(F(v)\subseteq F(\Succ(v))\bigr)
       \dsep u\in\mathbb{N}\dsep{}\\
      &\forall k\in\mathbb{N}\,\bigl(k<u\rightarrow
        F(\Succ(k))\subseteq F(u)\bigr)
      \vdash{}\\
      &\forall k\in\mathbb{N}\,\bigl(k<\Succ(u)\rightarrow
        F(\Succ(k))\subseteq F(\Succ(u))\bigr)
    \end{aligned}%
  }{%
    \forall v\in\mathbb{N}\,\bigl(F(v)\subseteq F(\Succ(v))\bigr)
    \dsep u\in\mathbb{N}\dsep
    \forall k\in\mathbb{N}\,\bigl(k<u\rightarrow
      F(\Succ(k))\subseteq F(u)\bigr)
    \vdash
    \forall k\in\mathbb{N}\,\bigl(k<\Succ(u)\rightarrow
      F(\Succ(k))\subseteq F(\Succ(u))\bigr)%
  }[B04PeanoStrictMonotoneSuccStep]
    \proofstepwidestar[1]{%
      \forall v\in\mathbb{N}\,\bigl(F(v)\subseteq F(\Succ(v))\bigr)}{\rA}
    \proofstepwidestar[2]{u\in\mathbb{N}}{\rA}
    \proofstepwidestar[3]{%
      \forall k\in\mathbb{N}\,\bigl(k<u\rightarrow
        F(\Succ(k))\subseteq F(u)\bigr)}{\rA}
    \proofstepwidestar[4]{k\in\mathbb{N}}{\rA}
    \proofstepwidestar[5]{k<\Succ(u)}{\rA}
    \proofstepwidestar[2,4,5]{k<u\lor k=u}{%
      \FormulaRefAuto{PeanoStrictSuccSplit}[thm]{4,2,5}}
    \proofstepwidestar[7]{k<u}{\rA}
    \proofstepwidestar[3,4,7]{F(\Succ(k))\subseteq F(u)}{%
      \rRE{7,\rRE{4,\rUE{3}}}}
    \proofstepwidestar[1,2]{F(u)\subseteq F(\Succ(u))}{%
      \rRE{2,\rUE{1}}}
    \proofstepwidestar[1,2,3,4,7]{%
      F(\Succ(k))\subseteq F(\Succ(u))}{%
      \FormulaRefAuto{A\subseteq B\dsep B\subseteq C\vdash
        A\subseteq C}{8,9}}
    \proofstepwidestar[11]{k=u}{\rA}
    \proofstepwidestar[11]{\Succ(k)=\Succ(u)}{\rIE{11}}
    \proofstepwidestar[11]{F(\Succ(k))=F(\Succ(u))}{\rIE{12}}
    \proofstepwidestar[11]{F(\Succ(k))\subseteq F(\Succ(u))}{%
      \FormulaRefAuto{A=B\vdash A\subseteq B}{13}}
    \proofstepwidestar[1,2,3,4,5]{%
      F(\Succ(k))\subseteq F(\Succ(u))}{\rOE{6,7,10,11,14}}
    \proofstepwidestar[1,2,3,4]{%
      k<\Succ(u)\rightarrow F(\Succ(k))\subseteq F(\Succ(u))}{\rRI{5,15}}
    \proofstepwidestar[1,2,3]{%
      \forall k\in\mathbb{N}\,\bigl(k<\Succ(u)\rightarrow
        F(\Succ(k))\subseteq F(\Succ(u))\bigr)}{\rUI{\rRI{4,16}}}
  \closeproofpartwideindR

  \proofpartwide{Induktionsschluss}
    \proofstepwidestar[1]{%
      \forall u\in\mathbb{N}\,\bigl(F(u)\subseteq F(\Succ(u))\bigr)}{\rA}
    \proofstepwidestar[2]{k\in\mathbb{N}}{\rA}
    \proofstepwidestar[3]{n\in\mathbb{N}}{\rA}
    \proofstepwidestar[4]{k<n}{\rA}
    \proofstepwidestar[1,3]{%
      \forall k\in\mathbb{N}\,\bigl(k<n\rightarrow
        F(\Succ(k))\subseteq F(n)\bigr)}{%
      \FormulaRefAuto{PeanoInductionRule}[ax]{IA,IS,3}}
    \proofstepwidestar[1,2,3]{%
      k<n\rightarrow F(\Succ(k))\subseteq F(n)}{\rRE{2,\rUE{5}}}
    \proofstepwidestar[1,2,3,4]{F(\Succ(k))\subseteq F(n)}{\rRE{6,4}}
  \closeproofpartwide
\end{tabproofsplitwide}

\FormulaThmDeltaK[Strikte Punkttrennung liefert Injektivität]{%
  F\colon\mathbb{N}\to M\dsep
  \forall k,n\in\mathbb{N}\,\bigl(k<n\rightarrow F(k)\neq F(n)\bigr)
  \vdash F\colon\mathbb{N}\inj M%
}{B04PeanoStrictSeparationInjection}{%
  \DeltaRow{Mengen}{M\dsep k\dsep n}%
  \DeltaRow{Funktionen}{F}%
}
\begin{tabproof}
  \proofstep{1}{F\colon\mathbb{N}\to M}{\rA}
  \proofstep{2}{%
    \forall k,n\in\mathbb{N}\,\bigl(k<n\rightarrow F(k)\neq F(n)\bigr)}{\rA}
  \proofstep{3}{k\in\mathbb{N}}{\rA}
  \proofstep{4}{n\in\mathbb{N}}{\rA}
  \proofstep{5}{F(k)=F(n)}{\rA}
  \proofstep{3,4}{k<n\lor n<k\lor k=n}{%
    \FormulaRefAuto{PeanoOrderTrichotomy}[thm]{3,4}}
  \proofstep{7}{k<n}{\rA}
  \proofstep{2,3,4,7}{F(k)\neq F(n)}{%
    \rRE{7,\rRE{4,\rRE{3,\rUE{2}}}}}
  \proofstep{2,3,4,5,7}{k=n}{%
    \FormulaRefAuto{P,\,\neg P\vdash Q}{5,8}}
  \proofstep{10}{n<k}{\rA}
  \proofstep{2,3,4,10}{F(n)\neq F(k)}{%
    \rRE{10,\rRE{3,\rRE{4,\rUE{2}}}}}
  \proofstep{5}{F(n)=F(k)}{\FormulaRefAuto{a=b\vdash b=a}{5}}
  \proofstep{2,3,4,5,10}{k=n}{%
    \FormulaRefAuto{P,\,\neg P\vdash Q}{12,11}}
  \proofstep{14}{k=n}{\rA}
  \proofstep{14}{k=n}{\FormulaRefAuto{P\vdash P}{14}}
  \proofstep{2,3,4,5}{k=n}{%
    \FormulaRefAuto{P_1\lor P_2\lor P_3\dsep
      [P_1]\vdots R\dsep [P_2]\vdots R\dsep [P_3]\vdots R\vdash R}{%
      6,7,9,10,13,14,15}}
  \proofstep{1,2}{%
    \forall k,n\in\mathbb{N}\,\bigl(F(k)=F(n)\rightarrow k=n\bigr)}{%
    \rUI{\rRI{3,\rRI{4,\rRI{5,16}}}}}
  \proofstep{1,2}{F\colon\mathbb{N}\inj M}{%
    \FormulaRefAuto{Injektive Funktion}{1,17}}
\end{tabproof}



\section{Multiplikation}

Die Multiplikation wird im zweiten Argument rekursiv definiert. Für ein
festes \(a\in\mathbb{N}\) benutzen wir als Schrittabbildung die Addition von
\(a\) auf der rechten Seite.

\FormulaDefDeltaK[Funktionsvorschrift des additiven Schritts]{%
  a\in\mathbb{N}\dsep x\in\mathbb{N}\vdash
  \tAddStep{a}(x)\coloneqq x+a%
}{Funktionsvorschrift des additiven Schritts}{%
  \DeltaRow{Mengen}{a\dsep x}%
  \DeltaRow{Funktionensymbol}{\tAddStep{a}}%
}

\FormulaThmDelta[Typisierung des additiven Schrittterms]{%
  a\in\mathbb{N}\dsep x\in\mathbb{N}\vdash \tAddStep{a}(x)\in\mathbb{N}%
}{%
  \DeltaRow{Mengen}{a\dsep x}%
  \DeltaRow{Funktionensymbol}{\tAddStep{a}}%
}
\begin{tabproof}
  \proofstep{1}{a\in\mathbb{N}}{\rA}
  \proofstep{2}{x\in\mathbb{N}}{\rA}
  \proofstep{1,2}{x+a\in\mathbb{N}}{%
    \FormulaRefAuto{a\in\mathbb{N}\dsep b\in\mathbb{N}\vdash a+b\in\mathbb{N}}{2,1}}
  \proofstep{1,2}{\tAddStep{a}(x)=x+a}{%
    \FormulaRefAuto{a\in\mathbb{N}\dsep x\in\mathbb{N}\vdash
      \tAddStep{a}(x)\coloneqq x+a}{1,2}}
  \proofstep{1,2}{\tAddStep{a}(x)\in\mathbb{N}}{\rIE{4,3}}
\end{tabproof}

\FormulaThmDelta[Existenz des additiven Schritts]{%
  a\in\mathbb{N}\vdash
  \exists! F\colon\mathbb{N}\to\mathbb{N}\,
  \forall x\in\mathbb{N}\,F(x)=\tAddStep{a}(x)%
}{%
  \DeltaRow{Mengen}{a\dsep x}%
  \DeltaRow{Funktionensymbol}{\tAddStep{a}}%
}
\begin{tabproof}
  \proofstep{1}{a\in\mathbb{N}}{\rA}
  \proofstep{2}{x\in\mathbb{N}}{\rA}
  \proofstep{1,2}{\tAddStep{a}(x)=x+a}{%
    \FormulaRefAuto{a\in\mathbb{N}\dsep x\in\mathbb{N}\vdash
      \tAddStep{a}(x)\coloneqq x+a}{1,2}}
  \proofstep{1}{\forall x\in\mathbb{N}\,\tAddStep{a}(x)=x+a}{%
    \rUI{\rRI{2,3}}}
  \proofstep{1,2}{\tAddStep{a}(x)\in\mathbb{N}}{%
    \FormulaRefAuto{a\in\mathbb{N}\dsep x\in\mathbb{N}\vdash
      \tAddStep{a}(x)\in\mathbb{N}}{1,2}}
  \proofstep{1}{\forall x\in\mathbb{N}\,\tAddStep{a}(x)\in\mathbb{N}}{%
    \rUI{\rRI{2,5}}}
  \proofstep{1}{%
    \exists! F\colon\mathbb{N}\to\mathbb{N}\,
    \forall x\in\mathbb{N}\,F(x)=\tAddStep{a}(x)}{%
    \FormulaRefAuto{%
      x\in A\vdash t(x)\coloneq y\dsep x\in A\vdash t(x)\in B
      \exists! F\colon A\to B\,\forall x\in A\,F(x)=t(x)}{4,6}}
\end{tabproof}

\FormulaDefDeltaK[Additiver Schritt zu \(a\)]{%
  a\in\mathbb{N}\vdash
  \AddStep{a}\coloneqq \iota F\Bigl(
    F\colon\mathbb{N}\to\mathbb{N}\land
    \forall x\in\mathbb{N}\,F(x)=\tAddStep{a}(x)
  \Bigr)%
}{Additiver Schritt}{%
  \DeltaRow{Mengen}{a\dsep x}%
  \DeltaRow{Funktionensymbole}{\AddStep{a}\dsep \tAddStep{a}\dsep F}%
}

\FormulaThmDelta[Additiver Schritt ist eine Abbildung]{%
  a\in\mathbb{N}\vdash \AddStep{a}\colon\mathbb{N}\to\mathbb{N}%
}{%
  \DeltaRow{Mengen}{a}%
  \DeltaRow{Funktionensymbol}{\AddStep{a}}%
}
\begin{tabproof}
  \proofstep{1}{a\in\mathbb{N}}{\rA}
  \proofstep{1}{\AddStep{a}\colon\mathbb{N}\to\mathbb{N}}{%
    \rAEa{\FormulaRefAuto{a\in\mathbb{N}\vdash
      \AddStep{a}\coloneqq \iota F\Bigl(
        F\colon\mathbb{N}\to\mathbb{N}\land
        \forall x\in\mathbb{N}\,F(x)=\tAddStep{a}(x)
      \Bigr)}{1}}}
\end{tabproof}

\FormulaThmDelta[Wert des additiven Schritts]{%
  a\in\mathbb{N}\dsep x\in\mathbb{N}\vdash \AddStep{a}(x)=x+a%
}{%
  \DeltaRow{Mengen}{a\dsep x}%
  \DeltaRow{Funktionensymbole}{\AddStep{a}\dsep \tAddStep{a}}%
}
\begin{tabproof}
  \proofstep{1}{a\in\mathbb{N}}{\rA}
  \proofstep{2}{x\in\mathbb{N}}{\rA}
  \proofstep{1}{\forall x\in\mathbb{N}\,\AddStep{a}(x)=\tAddStep{a}(x)}{%
    \rAEb{\FormulaRefAuto{a\in\mathbb{N}\vdash
      \AddStep{a}\coloneqq \iota F\Bigl(
        F\colon\mathbb{N}\to\mathbb{N}\land
        \forall x\in\mathbb{N}\,F(x)=\tAddStep{a}(x)
      \Bigr)}{1}}}
  \proofstep{1,2}{\AddStep{a}(x)=\tAddStep{a}(x)}{\rRE{2,\rUE{3}}}
  \proofstep{1,2}{\tAddStep{a}(x)=x+a}{%
    \FormulaRefAuto{a\in\mathbb{N}\dsep x\in\mathbb{N}\vdash
      \tAddStep{a}(x)\coloneqq x+a}{1,2}}
  \proofstep{1,2}{\AddStep{a}(x)=x+a}{%
    \FormulaRefAuto{a = b,\, b = c \vdash a = c}{4,5}}
\end{tabproof}

\FormulaDefDeltaK[Multiplikation natürlicher Zahlen]{%
  a\in\mathbb{N}\dsep b\in\mathbb{N}\vdash
  a\cdot b\coloneqq \RecFun{0}{\AddStep{a}}(b)%
}{Multiplikation auf \(\mathbb{N}\)}{%
  \DeltaRow{Mengen}{a\dsep b}%
  \DeltaRow{Funktionensymbole}{\AddStep{a}\dsep \RecFun{0}{\AddStep{a}}}%
  \DeltaRow{Neue Symbole}{\cdot}%
}

\FormulaThmDelta[Abgeschlossenheit der Multiplikation]{%
  a\in\mathbb{N}\dsep b\in\mathbb{N}\vdash a\cdot b\in\mathbb{N}%
}{%
  \DeltaRow{Mengen}{a\dsep b}%
}
\begin{tabproof}
  \proofstep{1}{a\in\mathbb{N}}{\rA}
  \proofstep{2}{b\in\mathbb{N}}{\rA}
  \proofstep{}{0\in\mathbb{N}}{%
    \FormulaRefAuto{PeanoZero}}
  \proofstep{1}{\AddStep{a}\colon\mathbb{N}\to\mathbb{N}}{%
    \FormulaRefAuto{a\in\mathbb{N}\vdash \AddStep{a}\colon\mathbb{N}\to\mathbb{N}}{1}}
  \proofstep{1,2}{\RecFun{0}{\AddStep{a}}(b)\in\mathbb{N}}{%
    \FormulaRefAuto{m_0\in M\dsep f\colon M\to M\dsep n\in\mathbb{N}\vdash
      \RecFun{m_0}{f}(n)\in M}{3,4,2}}
  \proofstep{1,2}{a\cdot b=\RecFun{0}{\AddStep{a}}(b)}{%
    \FormulaRefAuto{a\in\mathbb{N}\dsep b\in\mathbb{N}\vdash
      a\cdot b\coloneqq \RecFun{0}{\AddStep{a}}(b)}{1,2}}
  \proofstep{1,2}{a\cdot b\in\mathbb{N}}{\rIE{6,5}}
\end{tabproof}

\FormulaThmDelta[Nullregel der Multiplikation rechts]{%
  a\in\mathbb{N}\vdash a\cdot0=0%
}{%
  \DeltaRow{Mengen}{a}%
}
\begin{tabproof}
  \proofstep{1}{a\in\mathbb{N}}{\rA}
  \proofstep{}{0\in\mathbb{N}}{%
    \FormulaRefAuto{PeanoZero}}
  \proofstep{1}{\AddStep{a}\colon\mathbb{N}\to\mathbb{N}}{%
    \FormulaRefAuto{a\in\mathbb{N}\vdash \AddStep{a}\colon\mathbb{N}\to\mathbb{N}}{1}}
  \proofstep{1}{a\cdot0=\RecFun{0}{\AddStep{a}}(0)}{%
    \FormulaRefAuto{a\in\mathbb{N}\dsep b\in\mathbb{N}\vdash
      a\cdot b\coloneqq \RecFun{0}{\AddStep{a}}(b)}{1,2}}
  \proofstep{1}{\RecFun{0}{\AddStep{a}}(0)=0}{%
    \FormulaRefAuto{m_0\in M\dsep f\colon M\to M \vdash
      \RecFun{m_0}{f}(0)=m_0}{2,3}}
  \proofstep{1}{a\cdot0=0}{%
    \FormulaRefAuto{a = b,\, b = c \vdash a = c}{4,5}}
\end{tabproof}

\FormulaThmDeltaR[Nachfolgerregel der Multiplikation rechts]{%
  \begin{aligned}[t]
    &a\in\mathbb{N}\dsep b\in\mathbb{N}\vdash{}\\
    &a\cdot\Succ(b)=(a\cdot b)+a
  \end{aligned}
}{%
  a\in\mathbb{N}\dsep b\in\mathbb{N}\vdash a\cdot\Succ(b)=(a\cdot b)+a
}{%
  \DeltaRow{Mengen}{a\dsep b}%
  \DeltaRow{Funktionensymbole}{\AddStep{a}\dsep \RecFun{0}{\AddStep{a}}}%
}
\begin{tabproof}
  \proofstep{1}{a\in\mathbb{N}}{\rA}
  \proofstep{2}{b\in\mathbb{N}}{\rA}
  \proofstep{}{0\in\mathbb{N}}{%
    \FormulaRefAuto{PeanoZero}}
  \proofstep{1}{\AddStep{a}\colon\mathbb{N}\to\mathbb{N}}{%
    \FormulaRefAuto{a\in\mathbb{N}\vdash \AddStep{a}\colon\mathbb{N}\to\mathbb{N}}{1}}
  \proofstep{2}{\Succ(b)\in\mathbb{N}}{%
    \FormulaRefAuto{PeanoSuccClosure}{2}}
  \proofstep{1,2}{a\cdot\Succ(b)=
    \RecFun{0}{\AddStep{a}}(\Succ(b))}{%
    \FormulaRefAuto{a\in\mathbb{N}\dsep b\in\mathbb{N}\vdash
      a\cdot b\coloneqq \RecFun{0}{\AddStep{a}}(b)}{1,5}}
  \proofstep{1,2}{%
    \RecFun{0}{\AddStep{a}}(\Succ(b))=
    \AddStep{a}(\RecFun{0}{\AddStep{a}}(b))}{%
    \FormulaRefAuto{m_0\in M\dsep f\colon M\to M\dsep n\in\mathbb{N}\vdash
      \RecFun{m_0}{f}(\Succ(n))=f(\RecFun{m_0}{f}(n))}{3,4,2}}
  \proofstep{1,2}{a\cdot\Succ(b)=
    \AddStep{a}(\RecFun{0}{\AddStep{a}}(b))}{%
    \FormulaRefAuto{a = b,\, b = c \vdash a = c}{6,7}}
  \proofstep{1,2}{\RecFun{0}{\AddStep{a}}(b)\in\mathbb{N}}{%
    \FormulaRefAuto{m_0\in M\dsep f\colon M\to M\dsep n\in\mathbb{N}\vdash
      \RecFun{m_0}{f}(n)\in M}{3,4,2}}
  \proofstep{1,2}{%
    \AddStep{a}(\RecFun{0}{\AddStep{a}}(b))=
    \RecFun{0}{\AddStep{a}}(b)+a}{%
    \FormulaRefAuto{a\in\mathbb{N}\dsep x\in\mathbb{N}\vdash
      \AddStep{a}(x)=x+a}{1,9}}
  \proofstep{1,2}{a\cdot\Succ(b)=
    \RecFun{0}{\AddStep{a}}(b)+a}{%
    \FormulaRefAuto{a = b,\, b = c \vdash a = c}{8,10}}
  \proofstep{1,2}{a\cdot b=\RecFun{0}{\AddStep{a}}(b)}{%
    \FormulaRefAuto{a\in\mathbb{N}\dsep b\in\mathbb{N}\vdash
      a\cdot b\coloneqq \RecFun{0}{\AddStep{a}}(b)}{1,2}}
  \proofstep{1,2}{(a\cdot b)+a=
    \RecFun{0}{\AddStep{a}}(b)+a}{\rIE{12}}
  \proofstep{1,2}{\RecFun{0}{\AddStep{a}}(b)+a=
    (a\cdot b)+a}{%
    \FormulaRefAuto{a=b\vdash b=a}{13}}
  \proofstep{1,2}{a\cdot\Succ(b)=(a\cdot b)+a}{%
    \FormulaRefAuto{a = b,\, b = c \vdash a = c}{11,14}}
\end{tabproof}

Damit sind Addition und Multiplikation als totale Operationen auf
\(\mathbb{N}\) eingeführt. Ihre vertrauten Rechenregeln, etwa Assoziativität,
Kommutativität und Distributivität, sind nun Aussagen über diese
rekursiv definierten Abbildungen und können anschließend per Induktion
bewiesen werden.

\section{Summen- und Produktzeichen}

Die Zeichen \(\sum\) und \(\prod\) werden ebenfalls nicht als neue
Rechenprinzipien eingeführt. Sie sind Abkürzungen für Werte von
Dedekindschen Rekursionsabbildungen. Da der Rekursionssatz nur eine feste
Schrittabbildung \(S\colon M\to M\) verwendet, tragen wir den laufenden Index
als Teil des Rekursionszustands mit.

\subsection{Summen}

\subsubsection{Summenschritt}

Sei \(u\colon\mathbb{N}\to\mathbb{N}\), also eine Folge natürlicher Zahlen.
Der Summenschritt ist diejenige Abbildung auf
\(\mathbb{N}\times\mathbb{N}\), die aus einem Zustand \((k,x)\) den Zustand
\((\Succ(k),x + u(k))\) bildet.

\FormulaDefDeltaK[Funktionsvorschrift des Summenschritts]{%
  u\colon\mathbb{N}\to\mathbb{N}\dsep
  k\in\mathbb{N}\dsep x\in\mathbb{N}\vdash
  \tSumStep{u}(k,x)\coloneqq(\Succ(k),x + u(k))%
}{Funktionsvorschrift des Summenschritts}{%
  \DeltaRow{Mengen}{u\dsep k\dsep x}%
  \DeltaRow{Funktionensymbol}{\tSumStep{u}}%
}

\FormulaThmDelta[Typisierung des Summenschrittterms]{%
  u\colon\mathbb{N}\to\mathbb{N}\dsep
  k\in\mathbb{N}\dsep x\in\mathbb{N}
  \vdash \tSumStep{u}(k,x)\in\mathbb{N}\times\mathbb{N}%
}{%
  \DeltaRow{Mengen}{u\dsep k\dsep x}%
  \DeltaRow{Funktionensymbol}{\tSumStep{u}}%
}
\begin{tabproof}
  \proofstep{1}{u\colon\mathbb{N}\to\mathbb{N}}{\rA}
  \proofstep{2}{k\in\mathbb{N}}{\rA}
  \proofstep{3}{x\in\mathbb{N}}{\rA}
  \proofstep{2}{\Succ(k)\in\mathbb{N}}{%
    \FormulaRefAuto{PeanoSuccClosure}{2}}
  \proofstep{1,2}{u(k)\in\mathbb{N}}{%
    \FormulaRefAuto{F\colon A\to B\dsep x\in A\vdash F(x)\in B}{1,2}}
  \proofstep{1,2,3}{x + u(k)\in\mathbb{N}}{%
    \FormulaRefAuto{a\in\mathbb{N}\dsep b\in\mathbb{N}\vdash a+b\in\mathbb{N}}{3,5}}
  \proofstep{1,2,3}{\tSumStep{u}(k,x)=(\Succ(k),x + u(k))}{%
    \FormulaRefAuto{u\colon\mathbb{N}\to\mathbb{N}\dsep
      k\in\mathbb{N}\dsep x\in\mathbb{N}\vdash
      \tSumStep{u}(k,x)\coloneqq(\Succ(k),x + u(k))}{1,2,3}}
  \proofstep{1,2,3}{(\Succ(k),x + u(k))\in\mathbb{N}\times\mathbb{N}}{%
    \FormulaRefAuto{a\in A\dsep b\in B\vdash (a,b)\in A\times B}{4,6}}
  \proofstep{1,2,3}{\tSumStep{u}(k,x)\in\mathbb{N}\times\mathbb{N}}{\rIE{7,8}}
\end{tabproof}

\FormulaThmDelta[Existenz des Summenschritts]{%
  u\colon\mathbb{N}\to\mathbb{N}\vdash
  \exists! F\colon\mathbb{N}\times\mathbb{N}\to\mathbb{N}\times\mathbb{N}\,
  \forall k\in\mathbb{N}\,\forall x\in\mathbb{N}\,
  F(k,x)=\tSumStep{u}(k,x)%
}{%
  \DeltaRow{Mengen}{u\dsep k\dsep x}%
  \DeltaRow{Funktionensymbol}{\tSumStep{u}}%
}
\begin{tabproof}
  \proofstep{1}{u\colon\mathbb{N}\to\mathbb{N}}{\rA}
  \proofstep{2}{k\in\mathbb{N}}{\rA}
  \proofstep{3}{x\in\mathbb{N}}{\rA}
  \proofstep{1,2,3}{\tSumStep{u}(k,x)=(\Succ(k),x + u(k))}{%
    \FormulaRefAuto{u\colon\mathbb{N}\to\mathbb{N}\dsep
      k\in\mathbb{N}\dsep x\in\mathbb{N}\vdash
      \tSumStep{u}(k,x)\coloneqq(\Succ(k),x + u(k))}{1,2,3}}
  \proofstep{1,2}{%
    \forall x\in\mathbb{N}\,\tSumStep{u}(k,x)=(\Succ(k),x + u(k))}{%
    \rUI{\rRI{3,4}}}
  \proofstep{1}{%
    \forall k\in\mathbb{N}\,\forall x\in\mathbb{N}\,
    \tSumStep{u}(k,x)=(\Succ(k),x + u(k))}{%
    \rUI{\rRI{2,5}}}
  \proofstep{1,2,3}{\tSumStep{u}(k,x)\in\mathbb{N}\times\mathbb{N}}{%
    \FormulaRefAuto{u\colon\mathbb{N}\to\mathbb{N}\dsep
      k\in\mathbb{N}\dsep x\in\mathbb{N}
      \vdash \tSumStep{u}(k,x)\in\mathbb{N}\times\mathbb{N}}{1,2,3}}
  \proofstep{1,2}{%
    \forall x\in\mathbb{N}\,
    \tSumStep{u}(k,x)\in\mathbb{N}\times\mathbb{N}}{%
    \rUI{\rRI{3,7}}}
  \proofstep{1}{%
    \forall k\in\mathbb{N}\,\forall x\in\mathbb{N}\,
    \tSumStep{u}(k,x)\in\mathbb{N}\times\mathbb{N}}{%
    \rUI{\rRI{2,8}}}
  \proofstep{1}{%
    \exists! F\colon\mathbb{N}\times\mathbb{N}\to\mathbb{N}\times\mathbb{N}\,
    \forall k\in\mathbb{N}\,\forall x\in\mathbb{N}\,
    F(k,x)=\tSumStep{u}(k,x)}{%
    \FormulaRefAuto{%
      x\in A\vdash t(x)\coloneq y\dsep x\in A\vdash t(x)\in B
      \exists! F\colon A\to B\,\forall x\in A\,F(x)=t(x)}{6,9}}
\end{tabproof}

\FormulaDefDeltaR[Summenschritt einer Folge natürlicher Zahlen]{%
  \begin{aligned}[t]
    &u\colon\mathbb{N}\to\mathbb{N}\vdash{}\\
    &\SumStep{u}\coloneqq \iota F\Bigl(
      F\colon\mathbb{N}\times\mathbb{N}\to\mathbb{N}\times\mathbb{N}\land{}\\
    &\qquad \forall k\in\mathbb{N}\,\forall x\in\mathbb{N}\,
      F(k,x)=\tSumStep{u}(k,x)
    \Bigr)
  \end{aligned}
}{%
  u\colon\mathbb{N}\to\mathbb{N}\vdash
  \SumStep{u}\coloneqq \iota F\Bigl(
    F\colon\mathbb{N}\times\mathbb{N}\to\mathbb{N}\times\mathbb{N}\land
    \forall k\in\mathbb{N}\,\forall x\in\mathbb{N}\,
    F(k,x)=\tSumStep{u}(k,x)
  \Bigr)
}{%
  \DeltaRow{Mengen}{u\dsep k\dsep x}%
  \DeltaRow{Funktionensymbole}{\SumStep{u}\dsep \tSumStep{u}\dsep F}%
}

\FormulaThmDelta[Summenschritt ist eine Abbildung]{%
  u\colon\mathbb{N}\to\mathbb{N}\vdash
  \SumStep{u}\colon\mathbb{N}\times\mathbb{N}\to\mathbb{N}\times\mathbb{N}%
}{%
  \DeltaRow{Mengen}{u}%
  \DeltaRow{Funktionensymbol}{\SumStep{u}}%
}
\begin{tabproof}
  \proofstep{1}{u\colon\mathbb{N}\to\mathbb{N}}{\rA}
  \proofstep{1}{\SumStep{u}\colon\mathbb{N}\times\mathbb{N}\to\mathbb{N}\times\mathbb{N}}{%
    \rAEa{\FormulaRefAuto{u\colon\mathbb{N}\to\mathbb{N}\vdash
      \SumStep{u}\coloneqq \iota F\Bigl(
        F\colon\mathbb{N}\times\mathbb{N}\to\mathbb{N}\times\mathbb{N}\land
        \forall k\in\mathbb{N}\,\forall x\in\mathbb{N}\,
        F(k,x)=\tSumStep{u}(k,x)
      \Bigr)}{1}}}
\end{tabproof}

\FormulaThmDelta[Summenschrittwert liegt im Zustandsraum]{%
  u\colon\mathbb{N}\to\mathbb{N}\dsep
  z\in\mathbb{N}\times\mathbb{N}\vdash
  \SumStep{u}(z)\in\mathbb{N}\times\mathbb{N}%
}{%
  \DeltaRow{Mengen}{u\dsep z}%
  \DeltaRow{Funktionensymbol}{\SumStep{u}}%
}
\begin{tabproofwide}
  \proofstepwidestar[1]{u\colon\mathbb{N}\to\mathbb{N}}{\rA}
  \proofstepwidestar[2]{z\in\mathbb{N}\times\mathbb{N}}{\rA}
  \proofstepwidestar[1]{%
    \SumStep{u}\colon\mathbb{N}\times\mathbb{N}\to\mathbb{N}\times\mathbb{N}}{%
    \FormulaRefAuto{u\colon\mathbb{N}\to\mathbb{N}\vdash
      \SumStep{u}\colon\mathbb{N}\times\mathbb{N}\to\mathbb{N}\times\mathbb{N}}{1}}
  \proofstepwidestar[1,2]{\SumStep{u}(z)\in\mathbb{N}\times\mathbb{N}}{%
    \FormulaRefAuto{F\colon A\to B\dsep x\in A\vdash F(x)\in B}{3,2}}
\end{tabproofwide}

\FormulaThmDelta[Wert des Summenschritts]{%
  u\colon\mathbb{N}\to\mathbb{N}\dsep
  k\in\mathbb{N}\dsep x\in\mathbb{N}\vdash
  \SumStep{u}(k,x)=(\Succ(k),x + u(k))%
}{%
  \DeltaRow{Mengen}{u\dsep k\dsep x}%
  \DeltaRow{Funktionensymbole}{\SumStep{u}\dsep \tSumStep{u}}%
}
\begin{tabproof}
  \proofstep{1}{u\colon\mathbb{N}\to\mathbb{N}}{\rA}
  \proofstep{2}{k\in\mathbb{N}}{\rA}
  \proofstep{3}{x\in\mathbb{N}}{\rA}
  \proofstep{1}{%
    \forall k\in\mathbb{N}\,\forall x\in\mathbb{N}\,
    \SumStep{u}(k,x)=\tSumStep{u}(k,x)}{%
    \rAEb{\FormulaRefAuto{u\colon\mathbb{N}\to\mathbb{N}\vdash
      \SumStep{u}\coloneqq \iota F\Bigl(
        F\colon\mathbb{N}\times\mathbb{N}\to\mathbb{N}\times\mathbb{N}\land
        \forall k\in\mathbb{N}\,\forall x\in\mathbb{N}\,
        F(k,x)=\tSumStep{u}(k,x)
      \Bigr)}{1}}}
  \proofstep{1,2,3}{\SumStep{u}(k,x)=\tSumStep{u}(k,x)}{\rRE{3,\rRE{2,\rUE{4}}}}
  \proofstep{1,2,3}{\tSumStep{u}(k,x)=(\Succ(k),x + u(k))}{%
    \FormulaRefAuto{u\colon\mathbb{N}\to\mathbb{N}\dsep
      k\in\mathbb{N}\dsep x\in\mathbb{N}\vdash
      \tSumStep{u}(k,x)\coloneqq(\Succ(k),x + u(k))}{1,2,3}}
  \proofstep{1,2,3}{\SumStep{u}(k,x)=(\Succ(k),x + u(k))}{%
    \FormulaRefAuto{a = b,\, b = c \vdash a = c}{5,6}}
\end{tabproof}

Den von \(\SumStep{u}\) erzeugten Rekursionszustand fassen wir als eigenes
definiertes Symbol zusammen.

\subsubsection{Summenzustand}

\FormulaDefDeltaK[Summenzustand]{%
  u\colon\mathbb{N}\to\mathbb{N}\dsep m\in\mathbb{N}\vdash
  R_u(m)\coloneqq\RecFun{(0,0)}{\SumStep{u}}(m)%
}{Summenzustand}{%
  \DeltaRow{Mengen}{u\dsep m}%
  \DeltaRow{Funktionensymbole}{\SumStep{u}\dsep
    \RecFun{(0,0)}{\SumStep{u}}\dsep R_u}%
}

\FormulaThmDelta[Typisierung des Summenzustands]{%
  u\colon\mathbb{N}\to\mathbb{N}\dsep m\in\mathbb{N}\vdash
  R_u(m)\in\mathbb{N}\times\mathbb{N}%
}{%
  \DeltaRow{Mengen}{u\dsep m}%
  \DeltaRow{Funktionensymbole}{\SumStep{u}\dsep
    \RecFun{(0,0)}{\SumStep{u}}\dsep R_u}%
}
\begin{tabproofwide}
  \proofstepwidestar[1]{u\colon\mathbb{N}\to\mathbb{N}}{\rA}
  \proofstepwidestar[2]{m\in\mathbb{N}}{\rA}
  \proofstepwidestar[]{0\in\mathbb{N}}{\FormulaRefAuto{PeanoZero}}
  \proofstepwidestar[]{(0,0)\in\mathbb{N}\times\mathbb{N}}{%
    \FormulaRefAuto{a\in A\dsep b\in B\vdash (a,b)\in A\times B}{3,3}}
  \proofstepwidestar[1]{\SumStep{u}\colon\mathbb{N}\times\mathbb{N}\to
    \mathbb{N}\times\mathbb{N}}{%
    \FormulaRefAuto{u\colon\mathbb{N}\to\mathbb{N}\vdash
      \SumStep{u}\colon\mathbb{N}\times\mathbb{N}\to\mathbb{N}\times\mathbb{N}}{1}}
  \proofstepwidestar[1,2]{%
    \RecFun{(0,0)}{\SumStep{u}}(m)\in\mathbb{N}\times\mathbb{N}}{%
    \FormulaRefAuto{m_0\in M\dsep f\colon M\to M\dsep n\in\mathbb{N}\vdash
      \RecFun{m_0}{f}(n)\in M}{4,5,2}}
  \proofstepwidestar[1,2]{%
    R_u(m)=\RecFun{(0,0)}{\SumStep{u}}(m)}{%
    \FormulaRefAuto{u\colon\mathbb{N}\to\mathbb{N}\dsep m\in\mathbb{N}\vdash
      R_u(m)\coloneqq\RecFun{(0,0)}{\SumStep{u}}(m)}{1,2}}
  \proofstepwidestar[1,2]{R_u(m)\in\mathbb{N}\times\mathbb{N}}{\rIE{7}}
\end{tabproofwide}

\FormulaThmDelta[Nullregel des Summenzustands]{%
  u\colon\mathbb{N}\to\mathbb{N}\vdash
  R_u(0)=(0,0)%
}{%
  \DeltaRow{Mengen}{u}%
  \DeltaRow{Funktionensymbole}{\SumStep{u}\dsep
    \RecFun{(0,0)}{\SumStep{u}}\dsep R_u}%
}
\begin{tabproofwide}
  \proofstepwidestar[1]{u\colon\mathbb{N}\to\mathbb{N}}{\rA}
  \proofstepwidestar[]{0\in\mathbb{N}}{\FormulaRefAuto{PeanoZero}}
  \proofstepwidestar[]{(0,0)\in\mathbb{N}\times\mathbb{N}}{%
    \FormulaRefAuto{a\in A\dsep b\in B\vdash (a,b)\in A\times B}{2,2}}
  \proofstepwidestar[1]{\SumStep{u}\colon\mathbb{N}\times\mathbb{N}\to
    \mathbb{N}\times\mathbb{N}}{%
    \FormulaRefAuto{u\colon\mathbb{N}\to\mathbb{N}\vdash
      \SumStep{u}\colon\mathbb{N}\times\mathbb{N}\to\mathbb{N}\times\mathbb{N}}{1}}
  \proofstepwidestar[1]{%
    R_u(0)=\RecFun{(0,0)}{\SumStep{u}}(0)}{%
    \FormulaRefAuto{u\colon\mathbb{N}\to\mathbb{N}\dsep m\in\mathbb{N}\vdash
      R_u(m)\coloneqq\RecFun{(0,0)}{\SumStep{u}}(m)}{1,2}}
  \proofstepwidestar[1]{%
    \RecFun{(0,0)}{\SumStep{u}}(0)
    =(0,0)}{%
    \FormulaRefAuto{m_0\in M\dsep f\colon M\to M \vdash
      \RecFun{m_0}{f}(0)=m_0}{3,4}}
  \proofstepwidestar[1]{R_u(0)=(0,0)}{%
    \FormulaRefAuto{a = b,\, b = c \vdash a = c}{5,6}}
\end{tabproofwide}

\FormulaThmDelta[Nachfolgerregel des Summenzustands]{%
  u\colon\mathbb{N}\to\mathbb{N}\dsep m\in\mathbb{N}\vdash
  R_u(\Succ(m))=\SumStep{u}(R_u(m))%
}{%
  \DeltaRow{Mengen}{u\dsep m}%
  \DeltaRow{Funktionensymbole}{\SumStep{u}\dsep
    \RecFun{(0,0)}{\SumStep{u}}\dsep R_u}%
}
\begin{tabproofwide}
  \proofstepwidestar[1]{u\colon\mathbb{N}\to\mathbb{N}}{\rA}
  \proofstepwidestar[2]{m\in\mathbb{N}}{\rA}
  \proofstepwidestar[]{0\in\mathbb{N}}{\FormulaRefAuto{PeanoZero}}
  \proofstepwidestar[]{(0,0)\in\mathbb{N}\times\mathbb{N}}{%
    \FormulaRefAuto{a\in A\dsep b\in B\vdash (a,b)\in A\times B}{3,3}}
  \proofstepwidestar[1]{\SumStep{u}\colon\mathbb{N}\times\mathbb{N}\to
    \mathbb{N}\times\mathbb{N}}{%
    \FormulaRefAuto{u\colon\mathbb{N}\to\mathbb{N}\vdash
      \SumStep{u}\colon\mathbb{N}\times\mathbb{N}\to\mathbb{N}\times\mathbb{N}}{1}}
  \proofstepwidestar[2]{\Succ(m)\in\mathbb{N}}{%
    \FormulaRefAuto{PeanoSuccClosure}{2}}
  \proofstepwidestar[1,2]{%
    R_u(\Succ(m))=\RecFun{(0,0)}{\SumStep{u}}(\Succ(m))}{%
    \FormulaRefAuto{u\colon\mathbb{N}\to\mathbb{N}\dsep m\in\mathbb{N}\vdash
      R_u(m)\coloneqq\RecFun{(0,0)}{\SumStep{u}}(m)}{1,6}}
  \proofstepwidestar[1,2]{%
    \RecFun{(0,0)}{\SumStep{u}}(\Succ(m))
    =\SumStep{u}\bigl(\RecFun{(0,0)}{\SumStep{u}}(m)\bigr)}{%
    \FormulaRefAuto{m_0\in M\dsep f\colon M\to M\dsep n\in\mathbb{N}\vdash
      \RecFun{m_0}{f}(\Succ(n))=f(\RecFun{m_0}{f}(n))}{4,5,2}}
  \proofstepwidestar[1,2]{%
    R_u(m)=\RecFun{(0,0)}{\SumStep{u}}(m)}{%
    \FormulaRefAuto{u\colon\mathbb{N}\to\mathbb{N}\dsep m\in\mathbb{N}\vdash
      R_u(m)\coloneqq\RecFun{(0,0)}{\SumStep{u}}(m)}{1,2}}
  \proofstepwidestar[1,2]{%
    \RecFun{(0,0)}{\SumStep{u}}(m)=R_u(m)}{%
    \FormulaRefAuto{a=b\vdash b=a}{9}}
  \proofstepwidestar[1,2]{%
    \SumStep{u}\bigl(\RecFun{(0,0)}{\SumStep{u}}(m)\bigr)
    =\SumStep{u}(R_u(m))}{\rIE{10}}
  \proofstepwidestar[1,2]{%
    R_u(\Succ(m))=
    \SumStep{u}\bigl(\RecFun{(0,0)}{\SumStep{u}}(m)\bigr)}{%
    \FormulaRefAuto{a = b,\, b = c \vdash a = c}{7,8}}
  \proofstepwidestar[1,2]{R_u(\Succ(m))=\SumStep{u}(R_u(m))}{%
    \FormulaRefAuto{a = b,\, b = c \vdash a = c}{12,11}}
\end{tabproofwide}

Die endliche Summe über den Anfangsabschnitt der Länge \(n\) ist die zweite
Komponente der durch diesen Schritt erzeugten Rekursionsfolge. Der erste
Zustandswert ist der nächste zu lesende Index, der zweite Zustandswert ist
die bisher aufaddierte Summe.

\subsubsection{Summenzeichen}

\FormulaDefDeltaK[Summenzeichen über Anfangsabschnitten]{%
  u\colon\mathbb{N}\to\mathbb{N}\dsep n\in\mathbb{N}\vdash
  \sum_{i<n}u(i)\coloneqq
  \pi_2\bigl(R_u(n)\bigr)%
}{Summenzeichen}{%
  \DeltaRow{Mengen}{u\dsep n\dsep i}%
  \DeltaRow{Funktionensymbole}{\SumStep{u}\dsep R_u}%
  \DeltaRow{Neue Symbole}{\sum}%
}

Für spätere Rekursionsrechnungen trennen wir zunächst die beiden Projektionen
des Summenschritts ab.
Dabei verwenden wir wie bei den Projektionsfunktionen die Schreibweise
\(F(a,b)\) als Kurzform für \(F((a,b))\), wenn \(F\) auf einem kartesischen
Produkt definiert ist.

\FormulaThmDelta[Wert des Summenschritts am Zustand]{%
  u\colon\mathbb{N}\to\mathbb{N}\dsep z\in\mathbb{N}\times\mathbb{N}\vdash
  \SumStep{u}(z)=
  \bigl(\Succ(\pi_1(z)),\pi_2(z) + u(\pi_1(z))\bigr)%
}{%
  \DeltaRow{Mengen}{u\dsep z}%
  \DeltaRow{Funktionensymbol}{\SumStep{u}}%
}
\begin{tabproofwide}
  \proofstepwidestar[1]{u\colon\mathbb{N}\to\mathbb{N}}{\rA}
  \proofstepwidestar[2]{z\in\mathbb{N}\times\mathbb{N}}{\rA}
  \proofstepwidestar[2]{z=\bigl(\pi_1(z),\pi_2(z)\bigr)}{%
    \FormulaRefAuto{z\in A\times B \vdash z=(\pi_1(z),\pi_2(z))}{2}}
  \proofstepwidestar[1,2]{\SumStep{u}(z)=\SumStep{u}(z)}{\rII}
  \proofstepwidestar[1,2]{\SumStep{u}(z)=
    \SumStep{u}\bigl(\pi_1(z),\pi_2(z)\bigr)}{\rIE{3,4}}
  \proofstepwidestar[2]{\pi_1(z)\in\mathbb{N}}{%
    \FormulaRefAuto{z\in A\times B\vdash \pi_1(z)\in A}{2}}
  \proofstepwidestar[2]{\pi_2(z)\in\mathbb{N}}{%
    \FormulaRefAuto{z\in A\times B\vdash \pi_2(z)\in B}{2}}
  \proofstepwidestar[1,2]{%
    \SumStep{u}\bigl(\pi_1(z),\pi_2(z)\bigr)=
    \bigl(\Succ(\pi_1(z)),\pi_2(z) + u(\pi_1(z))\bigr)}{%
    \FormulaRefAuto{u\colon\mathbb{N}\to\mathbb{N}\dsep
      k\in\mathbb{N}\dsep x\in\mathbb{N}\vdash
      \SumStep{u}(k,x)=(\Succ(k),x + u(k))}{1,6,7}}
  \proofstepwidestar[1,2]{%
    \SumStep{u}(z)=
    \bigl(\Succ(\pi_1(z)),\pi_2(z) + u(\pi_1(z))\bigr)}{%
    \FormulaRefAuto{a = b,\, b = c \vdash a = c}{5,8}}
\end{tabproofwide}

\FormulaThmDelta[Erste Projektion des Summenschritts]{%
  u\colon\mathbb{N}\to\mathbb{N}\dsep z\in\mathbb{N}\times\mathbb{N}\vdash
  \pi_1(\SumStep{u}(z))=\Succ(\pi_1(z))%
}{%
  \DeltaRow{Mengen}{u\dsep z}%
  \DeltaRow{Funktionensymbol}{\SumStep{u}}%
}
\begin{tabproofwide}
  \proofstepwidestar[1]{u\colon\mathbb{N}\to\mathbb{N}}{\rA}
  \proofstepwidestar[2]{z\in\mathbb{N}\times\mathbb{N}}{\rA}
  \proofstepwidestar[1,2]{%
    \SumStep{u}(z)=
    \bigl(\Succ(\pi_1(z)),\pi_2(z) + u(\pi_1(z))\bigr)}{%
    \FormulaRefAuto{u\colon\mathbb{N}\to\mathbb{N}\dsep
      z\in\mathbb{N}\times\mathbb{N}\vdash
      \SumStep{u}(z)=
      \bigl(\Succ(\pi_1(z)),\pi_2(z) + u(\pi_1(z))\bigr)}{1,2}}
  \proofstepwidestar[1,2]{\SumStep{u}(z)\in\mathbb{N}\times\mathbb{N}}{%
    \FormulaRefAuto{u\colon\mathbb{N}\to\mathbb{N}\dsep
      z\in\mathbb{N}\times\mathbb{N}\vdash
      \SumStep{u}(z)\in\mathbb{N}\times\mathbb{N}}{1,2}}
  \proofstepwidestar[1,2]{\pi_1(\SumStep{u}(z))=\Succ(\pi_1(z))}{%
    \FormulaRefAuto{a\in A\times B\dsep a=(x,y)\vdash \pi_1(a)=x}{4,3}}
\end{tabproofwide}

\FormulaThmDelta[Zweite Projektion des Summenschritts]{%
  u\colon\mathbb{N}\to\mathbb{N}\dsep z\in\mathbb{N}\times\mathbb{N}\vdash
  \pi_2(\SumStep{u}(z))=\pi_2(z) + u(\pi_1(z))%
}{%
  \DeltaRow{Mengen}{u\dsep z}%
  \DeltaRow{Funktionensymbol}{\SumStep{u}}%
}
\begin{tabproofwide}
  \proofstepwidestar[1]{u\colon\mathbb{N}\to\mathbb{N}}{\rA}
  \proofstepwidestar[2]{z\in\mathbb{N}\times\mathbb{N}}{\rA}
  \proofstepwidestar[1,2]{%
    \SumStep{u}(z)=
    \bigl(\Succ(\pi_1(z)),\pi_2(z) + u(\pi_1(z))\bigr)}{%
    \FormulaRefAuto{u\colon\mathbb{N}\to\mathbb{N}\dsep
      z\in\mathbb{N}\times\mathbb{N}\vdash
      \SumStep{u}(z)=
      \bigl(\Succ(\pi_1(z)),\pi_2(z) + u(\pi_1(z))\bigr)}{1,2}}
  \proofstepwidestar[1,2]{\SumStep{u}(z)\in\mathbb{N}\times\mathbb{N}}{%
    \FormulaRefAuto{u\colon\mathbb{N}\to\mathbb{N}\dsep
      z\in\mathbb{N}\times\mathbb{N}\vdash
      \SumStep{u}(z)\in\mathbb{N}\times\mathbb{N}}{1,2}}
  \proofstepwidestar[1,2]{\pi_2(\SumStep{u}(z))=\pi_2(z) + u(\pi_1(z))}{%
    \FormulaRefAuto{a\in A\times B\dsep a=(x,y)\vdash \pi_2(a)=y}{4,3}}
\end{tabproofwide}

Die erste Komponente des Rekursionszustands zählt die bereits gelesenen
Glieder. Durch den Projektionssatz zum Summenschritt ist der Induktionsschritt
jetzt nur noch die Zustandsrekursion und ein Ersetzen mit der
Induktionsannahme.

\FormulaThmDeltaR[Indexkomponente des Summenzustands]{%
  \begin{aligned}[t]
    &u\colon\mathbb{N}\to\mathbb{N}\dsep n\in\mathbb{N}\vdash{}\\
    &\pi_1(R_u(n))=n
  \end{aligned}
}{%
  u\colon\mathbb{N}\to\mathbb{N}\dsep n\in\mathbb{N}\vdash
  \pi_1(R_u(n))=n
}{%
  \DeltaRow{Mengen}{u\dsep n\dsep m}%
  \DeltaRow{Funktionensymbole}{\SumStep{u}\dsep R_u}%
}
\begin{tabproofsplitwide}
  \proofpartwideindR[Induktionsanfang]{%
    u\colon\mathbb{N}\to\mathbb{N}\vdash \pi_1(R_u(0))=0
  }{%
    u\colon\mathbb{N}\to\mathbb{N}\vdash \pi_1(R_u(0))=0
  }
    \proofstepwidestar[1]{u\colon\mathbb{N}\to\mathbb{N}}{\rA}
    \proofstepwidestar[]{0\in\mathbb{N}}{%
      \FormulaRefAuto{PeanoZero}}
    \proofstepwidestar[]{(0,0)\in\mathbb{N}\times\mathbb{N}}{%
      \FormulaRefAuto{a\in A\dsep b\in B\vdash (a,b)\in A\times B}{2,2}}
    \proofstepwidestar[1]{R_u(0)=(0,0)}{%
      \FormulaRefAuto{u\colon\mathbb{N}\to\mathbb{N}\vdash
        R_u(0)=(0,0)}{1}}
    \proofstepwidestar[1]{\pi_1(R_u(0))=\pi_1(0,0)}{\rIE{4}}
    \proofstepwidestar[]{\pi_1(0,0)=0}{%
      \FormulaRefAuto{(x,y)\in A\times B \vdash \pi_1(x,y)=x}{3}}
    \proofstepwidestar[1]{\pi_1(R_u(0))=0}{%
      \FormulaRefAuto{a = b,\, b = c \vdash a = c}{5,6}}
  \closeproofpartwideindR

  \proofpartwideindR[Induktionsschritt]{%
    \begin{aligned}[t]
      &u\colon\mathbb{N}\to\mathbb{N}\dsep m\in\mathbb{N}\dsep
        \pi_1(R_u(m))=m\vdash{}\\
      &\pi_1(R_u(\Succ(m)))=\Succ(m)
    \end{aligned}
  }{%
    u\colon\mathbb{N}\to\mathbb{N}\dsep m\in\mathbb{N}\dsep
    \pi_1(R_u(m))=m\vdash \pi_1(R_u(\Succ(m)))=\Succ(m)
  }
    \proofstepwidestar[1]{u\colon\mathbb{N}\to\mathbb{N}}{\rA}
    \proofstepwidestar[2]{m\in\mathbb{N}}{\rA}
    \proofstepwidestar[3]{\pi_1(R_u(m))=m}{\rA}
    \proofstepwidestar[1,2]{R_u(m)\in\mathbb{N}\times\mathbb{N}}{%
      \FormulaRefAuto{u\colon\mathbb{N}\to\mathbb{N}\dsep m\in\mathbb{N}\vdash
        R_u(m)\in\mathbb{N}\times\mathbb{N}}{1,2}}
    \proofstepwidestar[1,2]{R_u(\Succ(m))=\SumStep{u}(R_u(m))}{%
    \FormulaRefAuto{u\colon\mathbb{N}\to\mathbb{N}\dsep m\in\mathbb{N}\vdash
      R_u(\Succ(m))=\SumStep{u}(R_u(m))}{1,2}}
    \proofstepwidestar[1,2]{%
      \pi_1(R_u(\Succ(m)))=\pi_1(\SumStep{u}(R_u(m)))}{\rIE{5}}
    \proofstepwidestar[1,2]{\pi_1(\SumStep{u}(R_u(m)))=\Succ(\pi_1(R_u(m)))}{%
      \FormulaRefAuto{u\colon\mathbb{N}\to\mathbb{N}\dsep z\in\mathbb{N}\times\mathbb{N}\vdash
        \pi_1(\SumStep{u}(z))=\Succ(\pi_1(z))}{1,4}}
    \proofstepwidestar[1,2]{%
      \pi_1(R_u(\Succ(m)))=\Succ(\pi_1(R_u(m)))}{%
      \FormulaRefAuto{a = b,\, b = c \vdash a = c}{6,7}}
    \proofstepwidestar[3]{\Succ(\pi_1(R_u(m)))=\Succ(m)}{\rIE{3}}
    \proofstepwidestar[1,2,3]{\pi_1(R_u(\Succ(m)))=\Succ(m)}{%
      \FormulaRefAuto{a = b,\, b = c \vdash a = c}{8,9}}
  \closeproofpartwideindR

  \proofpartwide{Induktionsschluss}
    \proofstepwidestar[1]{u\colon\mathbb{N}\to\mathbb{N}}{\rA}
    \proofstepwidestar[2]{n\in\mathbb{N}}{\rA}
    \proofstepwidestar[1,2]{\pi_1(R_u(n))=n}{%
      \FormulaRefAuto{PeanoInductionRule}{IA,IS,2}}
  \closeproofpartwide
\end{tabproofsplitwide}

Die zweite Komponente des Summenzustands erfüllt damit unmittelbar die
Rekursionsgleichung des Akkumulators.

\FormulaThmDeltaR[Nachfolgerregel der Summenakkumulation]{%
  \begin{aligned}[t]
    &u\colon\mathbb{N}\to\mathbb{N}\dsep n\in\mathbb{N}\vdash{}\\
    &\pi_2(R_u(\Succ(n)))=\pi_2(R_u(n)) + u(n)
  \end{aligned}
}{%
  u\colon\mathbb{N}\to\mathbb{N}\dsep n\in\mathbb{N}\vdash
  \pi_2(R_u(\Succ(n)))=\pi_2(R_u(n)) + u(n)
}{%
  \DeltaRow{Mengen}{u\dsep n\dsep m}%
  \DeltaRow{Funktionensymbole}{\SumStep{u}\dsep R_u}%
}
\begin{tabproofwide}
  \proofstepwidestar[1]{u\colon\mathbb{N}\to\mathbb{N}}{\rA}
  \proofstepwidestar[2]{n\in\mathbb{N}}{\rA}
  \proofstepwidestar[1,2]{R_u(n)\in\mathbb{N}\times\mathbb{N}}{%
    \FormulaRefAuto{u\colon\mathbb{N}\to\mathbb{N}\dsep m\in\mathbb{N}\vdash
      R_u(m)\in\mathbb{N}\times\mathbb{N}}{1,2}}
  \proofstepwidestar[1,2]{R_u(\Succ(n))=\SumStep{u}(R_u(n))}{%
    \FormulaRefAuto{u\colon\mathbb{N}\to\mathbb{N}\dsep m\in\mathbb{N}\vdash
      R_u(\Succ(m))=\SumStep{u}(R_u(m))}{1,2}}
  \proofstepwidestar[1,2]{%
    \pi_2\bigl(R_u(\Succ(n))\bigr)=
    \pi_2(\SumStep{u}(R_u(n)))}{\rIE{4}}
  \proofstepwidestar[1,2]{%
    \pi_2(\SumStep{u}(R_u(n)))=\pi_2(R_u(n)) + u(\pi_1(R_u(n)))}{%
    \FormulaRefAuto{u\colon\mathbb{N}\to\mathbb{N}\dsep z\in\mathbb{N}\times\mathbb{N}\vdash
      \pi_2(\SumStep{u}(z))=\pi_2(z) + u(\pi_1(z))}{1,3}}
  \proofstepwidestar[1,2]{%
    \pi_2\bigl(R_u(\Succ(n))\bigr)=\pi_2(R_u(n)) + u(\pi_1(R_u(n)))}{%
    \FormulaRefAuto{a = b,\, b = c \vdash a = c}{5,6}}
  \proofstepwidestar[1,2]{\pi_1(R_u(n))=n}{%
    \FormulaRefAuto{u\colon\mathbb{N}\to\mathbb{N}\dsep n\in\mathbb{N}\vdash
      \pi_1(R_u(n))=n}{1,2}}
  \proofstepwidestar[1,2]{u(\pi_1(R_u(n)))=u(n)}{\rIE{8}}
  \proofstepwidestar[1,2]{%
    \pi_2(R_u(n)) + u(\pi_1(R_u(n)))=\pi_2(R_u(n)) + u(n)}{\rIE{9}}
  \proofstepwidestar[1,2]{%
    \pi_2\bigl(R_u(\Succ(n))\bigr)=\pi_2(R_u(n)) + u(n)}{%
    \FormulaRefAuto{a = b,\, b = c \vdash a = c}{7,10}}
\end{tabproofwide}

\FormulaThmDelta[Nullregel des Summenzeichens]{%
  u\colon\mathbb{N}\to\mathbb{N}\vdash
  \sum_{i<0}u(i)=0%
}{%
  \DeltaRow{Mengen}{u\dsep i}%
  \DeltaRow{Funktionensymbole}{\SumStep{u}\dsep R_u}%
}
\begin{tabproof}
  \proofstep{1}{u\colon\mathbb{N}\to\mathbb{N}}{\rA}
  \proofstep{}{0\in\mathbb{N}}{\FormulaRefAuto{PeanoZero}}
  \proofstep{}{(0,0)\in\mathbb{N}\times\mathbb{N}}{%
    \FormulaRefAuto{a\in A\dsep b\in B\vdash (a,b)\in A\times B}{2,2}}
  \proofstep{1}{%
    \sum_{i<0}u(i)=
    \pi_2\bigl(R_u(0)\bigr)}{%
    \FormulaRefAuto{u\colon\mathbb{N}\to\mathbb{N}\dsep n\in\mathbb{N}\vdash
      \sum_{i<n}u(i)\coloneqq
      \pi_2\bigl(R_u(n)\bigr)}{1,2}}
  \proofstep{1}{R_u(0)=(0,0)}{%
    \FormulaRefAuto{u\colon\mathbb{N}\to\mathbb{N}\vdash
      R_u(0)=(0,0)}{1}}
  \proofstep{1}{\pi_2\bigl(R_u(0)\bigr)=\pi_2(0,0)}{\rIE{5}}
  \proofstep{}{\pi_2(0,0)=0}{%
    \FormulaRefAuto{(x,y)\in A\times B \vdash \pi_2(x,y)=y}{3}}
  \proofstep{1}{\pi_2\bigl(R_u(0)\bigr)=0}{%
    \FormulaRefAuto{a = b,\, b = c \vdash a = c}{6,7}}
  \proofstep{1}{\sum_{i<0}u(i)=0}{%
    \FormulaRefAuto{a = b,\, b = c \vdash a = c}{4,8}}
\end{tabproof}

\FormulaThmDeltaR[Nachfolgerregel des Summenzeichens]{%
  \begin{aligned}[t]
    &u\colon\mathbb{N}\to\mathbb{N}\dsep n\in\mathbb{N}\vdash{}\\
    &\sum_{i<\Succ(n)}u(i)=\left(\sum_{i<n}u(i)\right) + u(n)
  \end{aligned}
}{%
  u\colon\mathbb{N}\to\mathbb{N}\dsep n\in\mathbb{N}\vdash
  \sum_{i<\Succ(n)}u(i)=\left(\sum_{i<n}u(i)\right) + u(n)
}{%
  \DeltaRow{Mengen}{u\dsep n\dsep i}%
  \DeltaRow{Funktionensymbole}{\SumStep{u}\dsep R_u}%
}
\begin{tabproofwide}
  \proofstepwidestar[1]{u\colon\mathbb{N}\to\mathbb{N}}{\rA}
  \proofstepwidestar[2]{n\in\mathbb{N}}{\rA}
  \proofstepwidestar[2]{\Succ(n)\in\mathbb{N}}{%
    \FormulaRefAuto{PeanoSuccClosure}{2}}
  \proofstepwidestar[1,2]{%
    \sum_{i<\Succ(n)}u(i)=\pi_2\bigl(R_u(\Succ(n))\bigr)}{%
    \FormulaRefAuto{u\colon\mathbb{N}\to\mathbb{N}\dsep n\in\mathbb{N}\vdash
      \sum_{i<n}u(i)\coloneqq
      \pi_2\bigl(R_u(n)\bigr)}{1,3}}
  \proofstepwidestar[1,2]{%
    \pi_2\bigl(R_u(\Succ(n))\bigr)=\pi_2(R_u(n)) + u(n)}{%
    \FormulaRefAuto{u\colon\mathbb{N}\to\mathbb{N}\dsep n\in\mathbb{N}\vdash
      \pi_2(R_u(\Succ(n)))=\pi_2(R_u(n)) + u(n)}{1,2}}
  \proofstepwidestar[1,2]{%
    \sum_{i<\Succ(n)}u(i)=\pi_2(R_u(n)) + u(n)}{%
    \FormulaRefAuto{a = b,\, b = c \vdash a = c}{4,5}}
  \proofstepwidestar[1,2]{\sum_{i<n}u(i)=\pi_2(R_u(n))}{%
    \FormulaRefAuto{u\colon\mathbb{N}\to\mathbb{N}\dsep n\in\mathbb{N}\vdash
      \sum_{i<n}u(i)\coloneqq
      \pi_2\bigl(R_u(n)\bigr)}{1,2}}
  \proofstepwidestar[1,2]{\pi_2(R_u(n))=\sum_{i<n}u(i)}{%
    \FormulaRefAuto{a=b\vdash b=a}{7}}
  \proofstepwidestar[1,2]{%
    \pi_2(R_u(n)) + u(n)=\left(\sum_{i<n}u(i)\right) + u(n)}{\rIE{8}}
  \proofstepwidestar[1,2]{%
    \pi_2\bigl(R_u(\Succ(n))\bigr)=\left(\sum_{i<n}u(i)\right) + u(n)}{%
    \FormulaRefAuto{a = b,\, b = c \vdash a = c}{6,9}}
  \proofstepwidestar[1,2]{%
    \sum_{i<\Succ(n)}u(i)=\left(\sum_{i<n}u(i)\right) + u(n)}{%
    \FormulaRefAuto{a = b,\, b = c \vdash a = c}{4,10}}
\end{tabproofwide}

Die übliche Schreibweise mit oberer Grenze ist eine abgeleitete Notation.

\subsubsection{Schreibweise mit oberer Grenze}

\FormulaDefDeltaK[Summenzeichen mit oberer Grenze]{%
  u\colon\mathbb{N}\to\mathbb{N}\dsep n\in\mathbb{N}\vdash
  \sum_{i=0}^{n}u(i)\coloneqq
  \sum_{i<\Succ(n)}u(i)%
}{Summenzeichen mit oberer Grenze}{%
  \DeltaRow{Mengen}{u\dsep n\dsep i}%
  \DeltaRow{Funktionensymbole}{\SumStep{u}\dsep R_u}%
  \DeltaRow{Neue Schreibweise}{\sum_{i=0}^{n}u(i)}%
}

Analog wird das Produktzeichen über eine zweite Zustandsrekursion definiert.
Der Anfangswert des Produktakkumulators ist der erste Nachfolger der Null,
also \(\Succ(0)\). Die spätere Einseigenschaft dieses Elements ist
eine Rechenregel der Multiplikation; für die Definition genügt hier, dass
\(\Succ(0)\in\mathbb{N}\) gilt.

\subsection{Produkte}

\subsubsection{Produktschritt}

\FormulaDefDeltaK[Funktionsvorschrift des Produktschritts]{%
  u\colon\mathbb{N}\to\mathbb{N}\dsep
  k\in\mathbb{N}\dsep x\in\mathbb{N}\vdash
  \tProdStep{u}(k,x)\coloneqq(\Succ(k),x \cdot u(k))%
}{Funktionsvorschrift des Produktschritts}{%
  \DeltaRow{Mengen}{u\dsep k\dsep x}%
  \DeltaRow{Funktionensymbol}{\tProdStep{u}}%
}

\FormulaThmDelta[Typisierung des Produktschrittterms]{%
  u\colon\mathbb{N}\to\mathbb{N}\dsep
  k\in\mathbb{N}\dsep x\in\mathbb{N}
  \vdash \tProdStep{u}(k,x)\in\mathbb{N}\times\mathbb{N}%
}{%
  \DeltaRow{Mengen}{u\dsep k\dsep x}%
  \DeltaRow{Funktionensymbol}{\tProdStep{u}}%
}
\begin{tabproof}
  \proofstep{1}{u\colon\mathbb{N}\to\mathbb{N}}{\rA}
  \proofstep{2}{k\in\mathbb{N}}{\rA}
  \proofstep{3}{x\in\mathbb{N}}{\rA}
  \proofstep{2}{\Succ(k)\in\mathbb{N}}{%
    \FormulaRefAuto{PeanoSuccClosure}{2}}
  \proofstep{1,2}{u(k)\in\mathbb{N}}{%
    \FormulaRefAuto{F\colon A\to B\dsep x\in A\vdash F(x)\in B}{1,2}}
  \proofstep{1,2,3}{x \cdot u(k)\in\mathbb{N}}{%
    \FormulaRefAuto{a\in\mathbb{N}\dsep b\in\mathbb{N}\vdash a\cdot b\in\mathbb{N}}{3,5}}
  \proofstep{1,2,3}{\tProdStep{u}(k,x)=(\Succ(k),x \cdot u(k))}{%
    \FormulaRefAuto{u\colon\mathbb{N}\to\mathbb{N}\dsep
      k\in\mathbb{N}\dsep x\in\mathbb{N}\vdash
      \tProdStep{u}(k,x)\coloneqq(\Succ(k),x \cdot u(k))}{1,2,3}}
  \proofstep{1,2,3}{(\Succ(k),x \cdot u(k))\in\mathbb{N}\times\mathbb{N}}{%
    \FormulaRefAuto{a\in A\dsep b\in B\vdash (a,b)\in A\times B}{4,6}}
  \proofstep{1,2,3}{\tProdStep{u}(k,x)\in\mathbb{N}\times\mathbb{N}}{\rIE{7,8}}
\end{tabproof}

\FormulaThmDelta[Existenz des Produktschritts]{%
  u\colon\mathbb{N}\to\mathbb{N}\vdash
  \exists! F\colon\mathbb{N}\times\mathbb{N}\to\mathbb{N}\times\mathbb{N}\,
  \forall k\in\mathbb{N}\,\forall x\in\mathbb{N}\,
  F(k,x)=\tProdStep{u}(k,x)%
}{%
  \DeltaRow{Mengen}{u\dsep k\dsep x}%
  \DeltaRow{Funktionensymbol}{\tProdStep{u}}%
}
\begin{tabproof}
  \proofstep{1}{u\colon\mathbb{N}\to\mathbb{N}}{\rA}
  \proofstep{2}{k\in\mathbb{N}}{\rA}
  \proofstep{3}{x\in\mathbb{N}}{\rA}
  \proofstep{1,2,3}{\tProdStep{u}(k,x)=(\Succ(k),x \cdot u(k))}{%
    \FormulaRefAuto{u\colon\mathbb{N}\to\mathbb{N}\dsep
      k\in\mathbb{N}\dsep x\in\mathbb{N}\vdash
      \tProdStep{u}(k,x)\coloneqq(\Succ(k),x \cdot u(k))}{1,2,3}}
  \proofstep{1,2}{%
    \forall x\in\mathbb{N}\,\tProdStep{u}(k,x)=(\Succ(k),x \cdot u(k))}{%
    \rUI{\rRI{3,4}}}
  \proofstep{1}{%
    \forall k\in\mathbb{N}\,\forall x\in\mathbb{N}\,
    \tProdStep{u}(k,x)=(\Succ(k),x \cdot u(k))}{%
    \rUI{\rRI{2,5}}}
  \proofstep{1,2,3}{\tProdStep{u}(k,x)\in\mathbb{N}\times\mathbb{N}}{%
    \FormulaRefAuto{u\colon\mathbb{N}\to\mathbb{N}\dsep
      k\in\mathbb{N}\dsep x\in\mathbb{N}
      \vdash \tProdStep{u}(k,x)\in\mathbb{N}\times\mathbb{N}}{1,2,3}}
  \proofstep{1,2}{%
    \forall x\in\mathbb{N}\,
    \tProdStep{u}(k,x)\in\mathbb{N}\times\mathbb{N}}{%
    \rUI{\rRI{3,7}}}
  \proofstep{1}{%
    \forall k\in\mathbb{N}\,\forall x\in\mathbb{N}\,
    \tProdStep{u}(k,x)\in\mathbb{N}\times\mathbb{N}}{%
    \rUI{\rRI{2,8}}}
  \proofstep{1}{%
    \exists! F\colon\mathbb{N}\times\mathbb{N}\to\mathbb{N}\times\mathbb{N}\,
    \forall k\in\mathbb{N}\,\forall x\in\mathbb{N}\,
    F(k,x)=\tProdStep{u}(k,x)}{%
    \FormulaRefAuto{%
      x\in A\vdash t(x)\coloneq y\dsep x\in A\vdash t(x)\in B
      \exists! F\colon A\to B\,\forall x\in A\,F(x)=t(x)}{6,9}}
\end{tabproof}

\FormulaDefDeltaR[Produktschritt einer Folge natürlicher Zahlen]{%
  \begin{aligned}[t]
    &u\colon\mathbb{N}\to\mathbb{N}\vdash{}\\
    &\ProdStep{u}\coloneqq \iota F\Bigl(
      F\colon\mathbb{N}\times\mathbb{N}\to\mathbb{N}\times\mathbb{N}\land{}\\
    &\qquad \forall k\in\mathbb{N}\,\forall x\in\mathbb{N}\,
      F(k,x)=\tProdStep{u}(k,x)
    \Bigr)
  \end{aligned}
}{%
  u\colon\mathbb{N}\to\mathbb{N}\vdash
  \ProdStep{u}\coloneqq \iota F\Bigl(
    F\colon\mathbb{N}\times\mathbb{N}\to\mathbb{N}\times\mathbb{N}\land
    \forall k\in\mathbb{N}\,\forall x\in\mathbb{N}\,
    F(k,x)=\tProdStep{u}(k,x)
  \Bigr)
}{%
  \DeltaRow{Mengen}{u\dsep k\dsep x}%
  \DeltaRow{Funktionensymbole}{\ProdStep{u}\dsep \tProdStep{u}\dsep F}%
}

\FormulaThmDelta[Produktschritt ist eine Abbildung]{%
  u\colon\mathbb{N}\to\mathbb{N}\vdash
  \ProdStep{u}\colon\mathbb{N}\times\mathbb{N}\to\mathbb{N}\times\mathbb{N}%
}{%
  \DeltaRow{Mengen}{u}%
  \DeltaRow{Funktionensymbol}{\ProdStep{u}}%
}
\begin{tabproof}
  \proofstep{1}{u\colon\mathbb{N}\to\mathbb{N}}{\rA}
  \proofstep{1}{\ProdStep{u}\colon\mathbb{N}\times\mathbb{N}\to\mathbb{N}\times\mathbb{N}}{%
    \rAEa{\FormulaRefAuto{u\colon\mathbb{N}\to\mathbb{N}\vdash
      \ProdStep{u}\coloneqq \iota F\Bigl(
        F\colon\mathbb{N}\times\mathbb{N}\to\mathbb{N}\times\mathbb{N}\land
        \forall k\in\mathbb{N}\,\forall x\in\mathbb{N}\,
        F(k,x)=\tProdStep{u}(k,x)
      \Bigr)}{1}}}
\end{tabproof}

\FormulaThmDelta[Produktschrittwert liegt im Zustandsraum]{%
  u\colon\mathbb{N}\to\mathbb{N}\dsep
  z\in\mathbb{N}\times\mathbb{N}\vdash
  \ProdStep{u}(z)\in\mathbb{N}\times\mathbb{N}%
}{%
  \DeltaRow{Mengen}{u\dsep z}%
  \DeltaRow{Funktionensymbol}{\ProdStep{u}}%
}
\begin{tabproofwide}
  \proofstepwidestar[1]{u\colon\mathbb{N}\to\mathbb{N}}{\rA}
  \proofstepwidestar[2]{z\in\mathbb{N}\times\mathbb{N}}{\rA}
  \proofstepwidestar[1]{%
    \ProdStep{u}\colon\mathbb{N}\times\mathbb{N}\to\mathbb{N}\times\mathbb{N}}{%
    \FormulaRefAuto{u\colon\mathbb{N}\to\mathbb{N}\vdash
      \ProdStep{u}\colon\mathbb{N}\times\mathbb{N}\to\mathbb{N}\times\mathbb{N}}{1}}
  \proofstepwidestar[1,2]{\ProdStep{u}(z)\in\mathbb{N}\times\mathbb{N}}{%
    \FormulaRefAuto{F\colon A\to B\dsep x\in A\vdash F(x)\in B}{3,2}}
\end{tabproofwide}

\FormulaThmDelta[Wert des Produktschritts]{%
  u\colon\mathbb{N}\to\mathbb{N}\dsep
  k\in\mathbb{N}\dsep x\in\mathbb{N}\vdash
  \ProdStep{u}(k,x)=(\Succ(k),x \cdot u(k))%
}{%
  \DeltaRow{Mengen}{u\dsep k\dsep x}%
  \DeltaRow{Funktionensymbole}{\ProdStep{u}\dsep \tProdStep{u}}%
}
\begin{tabproof}
  \proofstep{1}{u\colon\mathbb{N}\to\mathbb{N}}{\rA}
  \proofstep{2}{k\in\mathbb{N}}{\rA}
  \proofstep{3}{x\in\mathbb{N}}{\rA}
  \proofstep{1}{%
    \forall k\in\mathbb{N}\,\forall x\in\mathbb{N}\,
    \ProdStep{u}(k,x)=\tProdStep{u}(k,x)}{%
    \rAEb{\FormulaRefAuto{u\colon\mathbb{N}\to\mathbb{N}\vdash
      \ProdStep{u}\coloneqq \iota F\Bigl(
        F\colon\mathbb{N}\times\mathbb{N}\to\mathbb{N}\times\mathbb{N}\land
        \forall k\in\mathbb{N}\,\forall x\in\mathbb{N}\,
        F(k,x)=\tProdStep{u}(k,x)
      \Bigr)}{1}}}
  \proofstep{1,2,3}{\ProdStep{u}(k,x)=\tProdStep{u}(k,x)}{\rRE{3,\rRE{2,\rUE{4}}}}
  \proofstep{1,2,3}{\tProdStep{u}(k,x)=(\Succ(k),x \cdot u(k))}{%
    \FormulaRefAuto{u\colon\mathbb{N}\to\mathbb{N}\dsep
      k\in\mathbb{N}\dsep x\in\mathbb{N}\vdash
      \tProdStep{u}(k,x)\coloneqq(\Succ(k),x \cdot u(k))}{1,2,3}}
  \proofstep{1,2,3}{\ProdStep{u}(k,x)=(\Succ(k),x \cdot u(k))}{%
    \FormulaRefAuto{a = b,\, b = c \vdash a = c}{5,6}}
\end{tabproof}

\subsubsection{Produktzustand}

\FormulaDefDeltaK[Produktzustand]{%
  u\colon\mathbb{N}\to\mathbb{N}\dsep m\in\mathbb{N}\vdash
  P_u(m)\coloneqq
  \RecFun{(0,\Succ(0))}{\ProdStep{u}}(m)%
}{Produktzustand}{%
  \DeltaRow{Mengen}{u\dsep m}%
  \DeltaRow{Funktionensymbole}{\ProdStep{u}\dsep
    \RecFun{(0,\Succ(0))}{\ProdStep{u}}\dsep P_u}%
}

\FormulaThmDelta[Typisierung des Produktzustands]{%
  u\colon\mathbb{N}\to\mathbb{N}\dsep m\in\mathbb{N}\vdash
  P_u(m)\in\mathbb{N}\times\mathbb{N}%
}{%
  \DeltaRow{Mengen}{u\dsep m}%
  \DeltaRow{Funktionensymbole}{\ProdStep{u}\dsep
    \RecFun{(0,\Succ(0))}{\ProdStep{u}}\dsep P_u}%
}
\begin{tabproofwide}
  \proofstepwidestar[1]{u\colon\mathbb{N}\to\mathbb{N}}{\rA}
  \proofstepwidestar[2]{m\in\mathbb{N}}{\rA}
  \proofstepwidestar[]{0\in\mathbb{N}}{\FormulaRefAuto{PeanoZero}}
  \proofstepwidestar[]{\Succ(0)\in\mathbb{N}}{%
    \FormulaRefAuto{PeanoSuccClosure}{3}}
  \proofstepwidestar[]{(0,\Succ(0))\in\mathbb{N}\times\mathbb{N}}{%
    \FormulaRefAuto{a\in A\dsep b\in B\vdash (a,b)\in A\times B}{3,4}}
  \proofstepwidestar[1]{\ProdStep{u}\colon\mathbb{N}\times\mathbb{N}\to
    \mathbb{N}\times\mathbb{N}}{%
    \FormulaRefAuto{u\colon\mathbb{N}\to\mathbb{N}\vdash
      \ProdStep{u}\colon\mathbb{N}\times\mathbb{N}\to\mathbb{N}\times\mathbb{N}}{1}}
  \proofstepwidestar[1,2]{%
    \RecFun{(0,\Succ(0))}{\ProdStep{u}}(m)
    \in\mathbb{N}\times\mathbb{N}}{%
    \FormulaRefAuto{m_0\in M\dsep f\colon M\to M\dsep n\in\mathbb{N}\vdash
      \RecFun{m_0}{f}(n)\in M}{5,6,2}}
  \proofstepwidestar[1,2]{%
    P_u(m)=\RecFun{(0,\Succ(0))}{\ProdStep{u}}(m)}{%
    \FormulaRefAuto{u\colon\mathbb{N}\to\mathbb{N}\dsep m\in\mathbb{N}\vdash
      P_u(m)\coloneqq
      \RecFun{(0,\Succ(0))}{\ProdStep{u}}(m)}{1,2}}
  \proofstepwidestar[1,2]{P_u(m)\in\mathbb{N}\times\mathbb{N}}{\rIE{8}}
\end{tabproofwide}

\FormulaThmDelta[Nullregel des Produktzustands]{%
  u\colon\mathbb{N}\to\mathbb{N}\vdash
  P_u(0)=(0,\Succ(0))%
}{%
  \DeltaRow{Mengen}{u}%
  \DeltaRow{Funktionensymbole}{\ProdStep{u}\dsep
    \RecFun{(0,\Succ(0))}{\ProdStep{u}}\dsep P_u}%
}
\begin{tabproofwide}
  \proofstepwidestar[1]{u\colon\mathbb{N}\to\mathbb{N}}{\rA}
  \proofstepwidestar[]{0\in\mathbb{N}}{\FormulaRefAuto{PeanoZero}}
  \proofstepwidestar[]{\Succ(0)\in\mathbb{N}}{%
    \FormulaRefAuto{PeanoSuccClosure}{2}}
  \proofstepwidestar[]{(0,\Succ(0))\in\mathbb{N}\times\mathbb{N}}{%
    \FormulaRefAuto{a\in A\dsep b\in B\vdash (a,b)\in A\times B}{2,3}}
  \proofstepwidestar[1]{\ProdStep{u}\colon\mathbb{N}\times\mathbb{N}\to
    \mathbb{N}\times\mathbb{N}}{%
    \FormulaRefAuto{u\colon\mathbb{N}\to\mathbb{N}\vdash
      \ProdStep{u}\colon\mathbb{N}\times\mathbb{N}\to\mathbb{N}\times\mathbb{N}}{1}}
  \proofstepwidestar[1]{%
    P_u(0)=
    \RecFun{(0,\Succ(0))}{\ProdStep{u}}(0)}{%
    \FormulaRefAuto{u\colon\mathbb{N}\to\mathbb{N}\dsep m\in\mathbb{N}\vdash
      P_u(m)\coloneqq
      \RecFun{(0,\Succ(0))}{\ProdStep{u}}(m)}{1,2}}
  \proofstepwidestar[1]{%
    \RecFun{(0,\Succ(0))}{\ProdStep{u}}(0)
    =(0,\Succ(0))}{%
    \FormulaRefAuto{m_0\in M\dsep f\colon M\to M \vdash
      \RecFun{m_0}{f}(0)=m_0}{4,5}}
  \proofstepwidestar[1]{P_u(0)=(0,\Succ(0))}{%
    \FormulaRefAuto{a = b,\, b = c \vdash a = c}{6,7}}
\end{tabproofwide}

\FormulaThmDelta[Nachfolgerregel des Produktzustands]{%
  u\colon\mathbb{N}\to\mathbb{N}\dsep m\in\mathbb{N}\vdash
  P_u(\Succ(m))=\ProdStep{u}(P_u(m))%
}{%
  \DeltaRow{Mengen}{u\dsep m}%
  \DeltaRow{Funktionensymbole}{\ProdStep{u}\dsep
    \RecFun{(0,\Succ(0))}{\ProdStep{u}}\dsep P_u}%
}
\begin{tabproofwide}
  \proofstepwidestar[1]{u\colon\mathbb{N}\to\mathbb{N}}{\rA}
  \proofstepwidestar[2]{m\in\mathbb{N}}{\rA}
  \proofstepwidestar[]{0\in\mathbb{N}}{\FormulaRefAuto{PeanoZero}}
  \proofstepwidestar[]{\Succ(0)\in\mathbb{N}}{%
    \FormulaRefAuto{PeanoSuccClosure}{3}}
  \proofstepwidestar[]{(0,\Succ(0))\in\mathbb{N}\times\mathbb{N}}{%
    \FormulaRefAuto{a\in A\dsep b\in B\vdash (a,b)\in A\times B}{3,4}}
  \proofstepwidestar[1]{\ProdStep{u}\colon\mathbb{N}\times\mathbb{N}\to
    \mathbb{N}\times\mathbb{N}}{%
    \FormulaRefAuto{u\colon\mathbb{N}\to\mathbb{N}\vdash
      \ProdStep{u}\colon\mathbb{N}\times\mathbb{N}\to\mathbb{N}\times\mathbb{N}}{1}}
  \proofstepwidestar[2]{\Succ(m)\in\mathbb{N}}{%
    \FormulaRefAuto{PeanoSuccClosure}{2}}
  \proofstepwidestar[1,2]{%
    P_u(\Succ(m))=
    \RecFun{(0,\Succ(0))}{\ProdStep{u}}(\Succ(m))}{%
    \FormulaRefAuto{u\colon\mathbb{N}\to\mathbb{N}\dsep m\in\mathbb{N}\vdash
      P_u(m)\coloneqq
      \RecFun{(0,\Succ(0))}{\ProdStep{u}}(m)}{1,7}}
  \proofstepwidestar[1,2]{%
    \RecFun{(0,\Succ(0))}{\ProdStep{u}}(\Succ(m))
    =\ProdStep{u}\bigl(\RecFun{(0,\Succ(0))}{\ProdStep{u}}(m)\bigr)}{%
    \FormulaRefAuto{m_0\in M\dsep f\colon M\to M\dsep n\in\mathbb{N}\vdash
      \RecFun{m_0}{f}(\Succ(n))=f(\RecFun{m_0}{f}(n))}{5,6,2}}
  \proofstepwidestar[1,2]{%
    P_u(m)=\RecFun{(0,\Succ(0))}{\ProdStep{u}}(m)}{%
    \FormulaRefAuto{u\colon\mathbb{N}\to\mathbb{N}\dsep m\in\mathbb{N}\vdash
      P_u(m)\coloneqq
      \RecFun{(0,\Succ(0))}{\ProdStep{u}}(m)}{1,2}}
  \proofstepwidestar[1,2]{%
    \RecFun{(0,\Succ(0))}{\ProdStep{u}}(m)=P_u(m)}{%
    \FormulaRefAuto{a=b\vdash b=a}{10}}
  \proofstepwidestar[1,2]{%
    \ProdStep{u}\bigl(\RecFun{(0,\Succ(0))}{\ProdStep{u}}(m)\bigr)
    =\ProdStep{u}(P_u(m))}{\rIE{11}}
  \proofstepwidestar[1,2]{%
    P_u(\Succ(m))=
    \ProdStep{u}\bigl(\RecFun{(0,\Succ(0))}{\ProdStep{u}}(m)\bigr)}{%
    \FormulaRefAuto{a = b,\, b = c \vdash a = c}{8,9}}
  \proofstepwidestar[1,2]{P_u(\Succ(m))=\ProdStep{u}(P_u(m))}{%
    \FormulaRefAuto{a = b,\, b = c \vdash a = c}{13,12}}
\end{tabproofwide}

\subsubsection{Produktzeichen}

\FormulaDefDeltaK[Produktzeichen über Anfangsabschnitten]{%
  u\colon\mathbb{N}\to\mathbb{N}\dsep n\in\mathbb{N}\vdash
  \prod_{i<n}u(i)\coloneqq
  \pi_2\bigl(P_u(n)\bigr)%
}{Produktzeichen}{%
  \DeltaRow{Mengen}{u\dsep n\dsep i}%
  \DeltaRow{Funktionensymbole}{\ProdStep{u}\dsep P_u}%
  \DeltaRow{Neue Symbole}{\prod}%
}

\FormulaThmDelta[Abgeschlossenheit des Produktzeichens]{%
  u\colon\mathbb{N}\to\mathbb{N}\dsep n\in\mathbb{N}\vdash
  \prod_{i<n}u(i)\in\mathbb{N}%
}{%
  \DeltaRow{Mengen}{u\dsep n\dsep i}%
  \DeltaRow{Funktionensymbole}{\ProdStep{u}\dsep P_u}%
}
\begin{tabproof}
  \proofstep{1}{u\colon\mathbb{N}\to\mathbb{N}}{\rA}
  \proofstep{2}{n\in\mathbb{N}}{\rA}
  \proofstep{1,2}{P_u(n)\in\mathbb{N}\times\mathbb{N}}{%
    \FormulaRefAuto{u\colon\mathbb{N}\to\mathbb{N}\dsep m\in\mathbb{N}\vdash
      P_u(m)\in\mathbb{N}\times\mathbb{N}}{1,2}}
  \proofstep{1,2}{\pi_2\bigl(P_u(n)\bigr)\in\mathbb{N}}{%
    \FormulaRefAuto{z\in A\times B \vdash \pi_2(z)\in B}{3}}
  \proofstep{1,2}{%
    \prod_{i<n}u(i)=\pi_2\bigl(P_u(n)\bigr)}{%
    \FormulaRefAuto{u\colon\mathbb{N}\to\mathbb{N}\dsep n\in\mathbb{N}\vdash
      \prod_{i<n}u(i)\coloneqq
      \pi_2\bigl(P_u(n)\bigr)}{1,2}}
  \proofstep{1,2}{\prod_{i<n}u(i)\in\mathbb{N}}{\rIE{5,4}}
\end{tabproof}

\FormulaThmDelta[Wert des Produktschritts am Zustand]{%
  u\colon\mathbb{N}\to\mathbb{N}\dsep z\in\mathbb{N}\times\mathbb{N}\vdash
  \ProdStep{u}(z)=
  \bigl(\Succ(\pi_1(z)),\pi_2(z)\cdot u(\pi_1(z))\bigr)%
}{%
  \DeltaRow{Mengen}{u\dsep z}%
  \DeltaRow{Funktionensymbol}{\ProdStep{u}}%
}
\begin{tabproofwide}
  \proofstepwidestar[1]{u\colon\mathbb{N}\to\mathbb{N}}{\rA}
  \proofstepwidestar[2]{z\in\mathbb{N}\times\mathbb{N}}{\rA}
  \proofstepwidestar[2]{z=\bigl(\pi_1(z),\pi_2(z)\bigr)}{%
    \FormulaRefAuto{z\in A\times B \vdash z=(\pi_1(z),\pi_2(z))}{2}}
  \proofstepwidestar[1,2]{\ProdStep{u}(z)=\ProdStep{u}(z)}{\rII}
  \proofstepwidestar[1,2]{\ProdStep{u}(z)=
    \ProdStep{u}\bigl(\pi_1(z),\pi_2(z)\bigr)}{\rIE{3,4}}
  \proofstepwidestar[2]{\pi_1(z)\in\mathbb{N}}{%
    \FormulaRefAuto{z\in A\times B\vdash \pi_1(z)\in A}{2}}
  \proofstepwidestar[2]{\pi_2(z)\in\mathbb{N}}{%
    \FormulaRefAuto{z\in A\times B\vdash \pi_2(z)\in B}{2}}
  \proofstepwidestar[1,2]{%
    \ProdStep{u}\bigl(\pi_1(z),\pi_2(z)\bigr)=
    \bigl(\Succ(\pi_1(z)),\pi_2(z)\cdot u(\pi_1(z))\bigr)}{%
    \FormulaRefAuto{u\colon\mathbb{N}\to\mathbb{N}\dsep
      k\in\mathbb{N}\dsep x\in\mathbb{N}\vdash
      \ProdStep{u}(k,x)=(\Succ(k),x \cdot u(k))}{1,6,7}}
  \proofstepwidestar[1,2]{%
    \ProdStep{u}(z)=
    \bigl(\Succ(\pi_1(z)),\pi_2(z)\cdot u(\pi_1(z))\bigr)}{%
    \FormulaRefAuto{a = b,\, b = c \vdash a = c}{5,8}}
\end{tabproofwide}

\FormulaThmDelta[Erste Projektion des Produktschritts]{%
  u\colon\mathbb{N}\to\mathbb{N}\dsep z\in\mathbb{N}\times\mathbb{N}\vdash
  \pi_1(\ProdStep{u}(z))=\Succ(\pi_1(z))%
}{%
  \DeltaRow{Mengen}{u\dsep z}%
  \DeltaRow{Funktionensymbol}{\ProdStep{u}}%
}
\begin{tabproofwide}
  \proofstepwidestar[1]{u\colon\mathbb{N}\to\mathbb{N}}{\rA}
  \proofstepwidestar[2]{z\in\mathbb{N}\times\mathbb{N}}{\rA}
  \proofstepwidestar[1,2]{%
    \ProdStep{u}(z)=
    \bigl(\Succ(\pi_1(z)),\pi_2(z)\cdot u(\pi_1(z))\bigr)}{%
    \FormulaRefAuto{u\colon\mathbb{N}\to\mathbb{N}\dsep
      z\in\mathbb{N}\times\mathbb{N}\vdash
      \ProdStep{u}(z)=
      \bigl(\Succ(\pi_1(z)),\pi_2(z)\cdot u(\pi_1(z))\bigr)}{1,2}}
  \proofstepwidestar[1,2]{\ProdStep{u}(z)\in\mathbb{N}\times\mathbb{N}}{%
    \FormulaRefAuto{u\colon\mathbb{N}\to\mathbb{N}\dsep
      z\in\mathbb{N}\times\mathbb{N}\vdash
      \ProdStep{u}(z)\in\mathbb{N}\times\mathbb{N}}{1,2}}
  \proofstepwidestar[1,2]{\pi_1(\ProdStep{u}(z))=\Succ(\pi_1(z))}{%
    \FormulaRefAuto{a\in A\times B\dsep a=(x,y)\vdash \pi_1(a)=x}{4,3}}
\end{tabproofwide}

\FormulaThmDelta[Zweite Projektion des Produktschritts]{%
  u\colon\mathbb{N}\to\mathbb{N}\dsep z\in\mathbb{N}\times\mathbb{N}\vdash
  \pi_2(\ProdStep{u}(z))=\pi_2(z)\cdot u(\pi_1(z))%
}{%
  \DeltaRow{Mengen}{u\dsep z}%
  \DeltaRow{Funktionensymbol}{\ProdStep{u}}%
}
\begin{tabproofwide}
  \proofstepwidestar[1]{u\colon\mathbb{N}\to\mathbb{N}}{\rA}
  \proofstepwidestar[2]{z\in\mathbb{N}\times\mathbb{N}}{\rA}
  \proofstepwidestar[1,2]{%
    \ProdStep{u}(z)=
    \bigl(\Succ(\pi_1(z)),\pi_2(z)\cdot u(\pi_1(z))\bigr)}{%
    \FormulaRefAuto{u\colon\mathbb{N}\to\mathbb{N}\dsep
      z\in\mathbb{N}\times\mathbb{N}\vdash
      \ProdStep{u}(z)=
      \bigl(\Succ(\pi_1(z)),\pi_2(z)\cdot u(\pi_1(z))\bigr)}{1,2}}
  \proofstepwidestar[1,2]{\ProdStep{u}(z)\in\mathbb{N}\times\mathbb{N}}{%
    \FormulaRefAuto{u\colon\mathbb{N}\to\mathbb{N}\dsep
      z\in\mathbb{N}\times\mathbb{N}\vdash
      \ProdStep{u}(z)\in\mathbb{N}\times\mathbb{N}}{1,2}}
  \proofstepwidestar[1,2]{\pi_2(\ProdStep{u}(z))=\pi_2(z)\cdot u(\pi_1(z))}{%
    \FormulaRefAuto{a\in A\times B\dsep a=(x,y)\vdash \pi_2(a)=y}{4,3}}
\end{tabproofwide}

\FormulaThmDeltaR[Indexkomponente des Produktzustands]{%
  \begin{aligned}[t]
    &u\colon\mathbb{N}\to\mathbb{N}\dsep n\in\mathbb{N}\vdash{}\\
    &\pi_1(P_u(n))=n
  \end{aligned}
}{%
  u\colon\mathbb{N}\to\mathbb{N}\dsep n\in\mathbb{N}\vdash
  \pi_1(P_u(n))=n
}{%
  \DeltaRow{Mengen}{u\dsep n\dsep m}%
  \DeltaRow{Funktionensymbole}{\ProdStep{u}\dsep P_u}%
}
\begin{tabproofsplitwide}
  \proofpartwideindR[Induktionsanfang]{%
    u\colon\mathbb{N}\to\mathbb{N}\vdash \pi_1(P_u(0))=0
  }{%
    u\colon\mathbb{N}\to\mathbb{N}\vdash \pi_1(P_u(0))=0
  }
    \proofstepwidestar[1]{u\colon\mathbb{N}\to\mathbb{N}}{\rA}
    \proofstepwidestar[]{0\in\mathbb{N}}{%
      \FormulaRefAuto{PeanoZero}}
    \proofstepwidestar[]{\Succ(0)\in\mathbb{N}}{%
      \FormulaRefAuto{PeanoSuccClosure}{2}}
    \proofstepwidestar[]{%
      (0,\Succ(0))\in\mathbb{N}\times\mathbb{N}}{%
      \FormulaRefAuto{a\in A\dsep b\in B\vdash (a,b)\in A\times B}{2,3}}
    \proofstepwidestar[1]{P_u(0)=(0,\Succ(0))}{%
      \FormulaRefAuto{u\colon\mathbb{N}\to\mathbb{N}\vdash
        P_u(0)=(0,\Succ(0))}{1}}
    \proofstepwidestar[1]{%
      \pi_1(P_u(0))=\pi_1(0,\Succ(0))}{\rIE{5}}
    \proofstepwidestar[]{\pi_1(0,\Succ(0))=0}{%
      \FormulaRefAuto{(x,y)\in A\times B \vdash \pi_1(x,y)=x}{4}}
    \proofstepwidestar[1]{\pi_1(P_u(0))=0}{%
      \FormulaRefAuto{a = b,\, b = c \vdash a = c}{6,7}}
  \closeproofpartwideindR

  \proofpartwideindR[Induktionsschritt]{%
    \begin{aligned}[t]
      &u\colon\mathbb{N}\to\mathbb{N}\dsep m\in\mathbb{N}\dsep
        \pi_1(P_u(m))=m\vdash{}\\
      &\pi_1(P_u(\Succ(m)))=\Succ(m)
    \end{aligned}
  }{%
    u\colon\mathbb{N}\to\mathbb{N}\dsep m\in\mathbb{N}\dsep
    \pi_1(P_u(m))=m\vdash \pi_1(P_u(\Succ(m)))=\Succ(m)
  }
    \proofstepwidestar[1]{u\colon\mathbb{N}\to\mathbb{N}}{\rA}
    \proofstepwidestar[2]{m\in\mathbb{N}}{\rA}
    \proofstepwidestar[3]{\pi_1(P_u(m))=m}{\rA}
    \proofstepwidestar[1,2]{P_u(m)\in\mathbb{N}\times\mathbb{N}}{%
      \FormulaRefAuto{u\colon\mathbb{N}\to\mathbb{N}\dsep m\in\mathbb{N}\vdash
        P_u(m)\in\mathbb{N}\times\mathbb{N}}{1,2}}
    \proofstepwidestar[1,2]{P_u(\Succ(m))=\ProdStep{u}(P_u(m))}{%
      \FormulaRefAuto{u\colon\mathbb{N}\to\mathbb{N}\dsep m\in\mathbb{N}\vdash
        P_u(\Succ(m))=\ProdStep{u}(P_u(m))}{1,2}}
    \proofstepwidestar[1,2]{%
      \pi_1(P_u(\Succ(m)))=\pi_1(\ProdStep{u}(P_u(m)))}{\rIE{5}}
    \proofstepwidestar[1,2]{\pi_1(\ProdStep{u}(P_u(m)))=\Succ(\pi_1(P_u(m)))}{%
      \FormulaRefAuto{u\colon\mathbb{N}\to\mathbb{N}\dsep z\in\mathbb{N}\times\mathbb{N}\vdash
        \pi_1(\ProdStep{u}(z))=\Succ(\pi_1(z))}{1,4}}
    \proofstepwidestar[1,2]{\pi_1(P_u(\Succ(m)))=\Succ(\pi_1(P_u(m)))}{%
      \FormulaRefAuto{a = b,\, b = c \vdash a = c}{6,7}}
    \proofstepwidestar[3]{\Succ(\pi_1(P_u(m)))=\Succ(m)}{\rIE{3}}
    \proofstepwidestar[1,2,3]{\pi_1(P_u(\Succ(m)))=\Succ(m)}{%
      \FormulaRefAuto{a = b,\, b = c \vdash a = c}{8,9}}
  \closeproofpartwideindR

  \proofpartwide{Induktionsschluss}
    \proofstepwidestar[1]{u\colon\mathbb{N}\to\mathbb{N}}{\rA}
    \proofstepwidestar[2]{n\in\mathbb{N}}{\rA}
    \proofstepwidestar[1,2]{\pi_1(P_u(n))=n}{%
      \FormulaRefAuto{PeanoInductionRule}{IA,IS,2}}
  \closeproofpartwide
\end{tabproofsplitwide}

\FormulaThmDeltaR[Nachfolgerregel der Produktakkumulation]{%
  \begin{aligned}[t]
    &u\colon\mathbb{N}\to\mathbb{N}\dsep n\in\mathbb{N}\vdash{}\\
    &\pi_2(P_u(\Succ(n)))=\pi_2(P_u(n))\cdot u(n)
  \end{aligned}
}{%
  u\colon\mathbb{N}\to\mathbb{N}\dsep n\in\mathbb{N}\vdash
  \pi_2(P_u(\Succ(n)))=\pi_2(P_u(n))\cdot u(n)
}{%
  \DeltaRow{Mengen}{u\dsep n\dsep m}%
  \DeltaRow{Funktionensymbole}{\ProdStep{u}\dsep P_u}%
}
\begin{tabproofwide}
  \proofstepwidestar[1]{u\colon\mathbb{N}\to\mathbb{N}}{\rA}
  \proofstepwidestar[2]{n\in\mathbb{N}}{\rA}
  \proofstepwidestar[1,2]{P_u(n)\in\mathbb{N}\times\mathbb{N}}{%
    \FormulaRefAuto{u\colon\mathbb{N}\to\mathbb{N}\dsep m\in\mathbb{N}\vdash
      P_u(m)\in\mathbb{N}\times\mathbb{N}}{1,2}}
  \proofstepwidestar[1,2]{P_u(\Succ(n))=\ProdStep{u}(P_u(n))}{%
    \FormulaRefAuto{u\colon\mathbb{N}\to\mathbb{N}\dsep m\in\mathbb{N}\vdash
      P_u(\Succ(m))=\ProdStep{u}(P_u(m))}{1,2}}
  \proofstepwidestar[1,2]{%
    \pi_2\bigl(P_u(\Succ(n))\bigr)=
    \pi_2(\ProdStep{u}(P_u(n)))}{\rIE{4}}
  \proofstepwidestar[1,2]{%
    \pi_2(\ProdStep{u}(P_u(n)))=\pi_2(P_u(n))\cdot u(\pi_1(P_u(n)))}{%
    \FormulaRefAuto{u\colon\mathbb{N}\to\mathbb{N}\dsep z\in\mathbb{N}\times\mathbb{N}\vdash
      \pi_2(\ProdStep{u}(z))=\pi_2(z)\cdot u(\pi_1(z))}{1,3}}
  \proofstepwidestar[1,2]{%
    \pi_2\bigl(P_u(\Succ(n))\bigr)=\pi_2(P_u(n))\cdot u(\pi_1(P_u(n)))}{%
    \FormulaRefAuto{a = b,\, b = c \vdash a = c}{5,6}}
  \proofstepwidestar[1,2]{\pi_1(P_u(n))=n}{%
    \FormulaRefAuto{u\colon\mathbb{N}\to\mathbb{N}\dsep n\in\mathbb{N}\vdash
      \pi_1(P_u(n))=n}{1,2}}
  \proofstepwidestar[1,2]{u(\pi_1(P_u(n)))=u(n)}{\rIE{8}}
  \proofstepwidestar[1,2]{%
    \pi_2(P_u(n))\cdot u(\pi_1(P_u(n)))=\pi_2(P_u(n))\cdot u(n)}{\rIE{9}}
  \proofstepwidestar[1,2]{%
    \pi_2\bigl(P_u(\Succ(n))\bigr)=\pi_2(P_u(n))\cdot u(n)}{%
    \FormulaRefAuto{a = b,\, b = c \vdash a = c}{7,10}}
\end{tabproofwide}

\FormulaThmDelta[Nullregel des Produktzeichens]{%
  u\colon\mathbb{N}\to\mathbb{N}\vdash
  \prod_{i<0}u(i)=\Succ(0)%
}{%
  \DeltaRow{Mengen}{u\dsep i}%
  \DeltaRow{Funktionensymbole}{\ProdStep{u}\dsep P_u}%
}
\begin{tabproofwide}
  \proofstepwidestar[1]{u\colon\mathbb{N}\to\mathbb{N}}{\rA}
  \proofstepwidestar[]{0\in\mathbb{N}}{\FormulaRefAuto{PeanoZero}}
  \proofstepwidestar[]{\Succ(0)\in\mathbb{N}}{%
    \FormulaRefAuto{PeanoSuccClosure}{2}}
  \proofstepwidestar[]{(0,\Succ(0))\in\mathbb{N}\times\mathbb{N}}{%
    \FormulaRefAuto{a\in A\dsep b\in B\vdash (a,b)\in A\times B}{2,3}}
  \proofstepwidestar[1]{%
    \prod_{i<0}u(i)=
    \pi_2\bigl(P_u(0)\bigr)}{%
    \FormulaRefAuto{u\colon\mathbb{N}\to\mathbb{N}\dsep n\in\mathbb{N}\vdash
      \prod_{i<n}u(i)\coloneqq
      \pi_2\bigl(P_u(n)\bigr)}{1,2}}
  \proofstepwidestar[1]{P_u(0)=(0,\Succ(0))}{%
    \FormulaRefAuto{u\colon\mathbb{N}\to\mathbb{N}\vdash
      P_u(0)=(0,\Succ(0))}{1}}
  \proofstepwidestar[1]{%
    \pi_2\bigl(P_u(0)\bigr)=\pi_2(0,\Succ(0))}{\rIE{6}}
  \proofstepwidestar[]{\pi_2(0,\Succ(0))=\Succ(0)}{%
    \FormulaRefAuto{(x,y)\in A\times B \vdash \pi_2(x,y)=y}{4}}
  \proofstepwidestar[1]{\pi_2\bigl(P_u(0)\bigr)=\Succ(0)}{%
    \FormulaRefAuto{a = b,\, b = c \vdash a = c}{7,8}}
  \proofstepwidestar[1]{\prod_{i<0}u(i)=\Succ(0)}{%
    \FormulaRefAuto{a = b,\, b = c \vdash a = c}{5,9}}
\end{tabproofwide}

\FormulaThmDeltaR[Nachfolgerregel des Produktzeichens]{%
  \begin{aligned}[t]
    &u\colon\mathbb{N}\to\mathbb{N}\dsep n\in\mathbb{N}\vdash{}\\
    &\prod_{i<\Succ(n)}u(i)=\left(\prod_{i<n}u(i)\right)\cdot u(n)
  \end{aligned}
}{%
  u\colon\mathbb{N}\to\mathbb{N}\dsep n\in\mathbb{N}\vdash
  \prod_{i<\Succ(n)}u(i)=\left(\prod_{i<n}u(i)\right)\cdot u(n)
}{%
  \DeltaRow{Mengen}{u\dsep n\dsep i}%
  \DeltaRow{Funktionensymbole}{\ProdStep{u}\dsep P_u}%
}
\begin{tabproofwide}
  \proofstepwidestar[1]{u\colon\mathbb{N}\to\mathbb{N}}{\rA}
  \proofstepwidestar[2]{n\in\mathbb{N}}{\rA}
  \proofstepwidestar[2]{\Succ(n)\in\mathbb{N}}{%
    \FormulaRefAuto{PeanoSuccClosure}{2}}
  \proofstepwidestar[1,2]{%
    \prod_{i<\Succ(n)}u(i)=\pi_2\bigl(P_u(\Succ(n))\bigr)}{%
    \FormulaRefAuto{u\colon\mathbb{N}\to\mathbb{N}\dsep n\in\mathbb{N}\vdash
      \prod_{i<n}u(i)\coloneqq
      \pi_2\bigl(P_u(n)\bigr)}{1,3}}
  \proofstepwidestar[1,2]{%
    \pi_2\bigl(P_u(\Succ(n))\bigr)=\pi_2(P_u(n))\cdot u(n)}{%
    \FormulaRefAuto{u\colon\mathbb{N}\to\mathbb{N}\dsep n\in\mathbb{N}\vdash
      \pi_2(P_u(\Succ(n)))=\pi_2(P_u(n))\cdot u(n)}{1,2}}
  \proofstepwidestar[1,2]{%
    \prod_{i<\Succ(n)}u(i)=\pi_2(P_u(n))\cdot u(n)}{%
    \FormulaRefAuto{a = b,\, b = c \vdash a = c}{4,5}}
  \proofstepwidestar[1,2]{\prod_{i<n}u(i)=\pi_2(P_u(n))}{%
    \FormulaRefAuto{u\colon\mathbb{N}\to\mathbb{N}\dsep n\in\mathbb{N}\vdash
      \prod_{i<n}u(i)\coloneqq
      \pi_2\bigl(P_u(n)\bigr)}{1,2}}
  \proofstepwidestar[1,2]{\pi_2(P_u(n))=\prod_{i<n}u(i)}{%
    \FormulaRefAuto{a=b\vdash b=a}{7}}
  \proofstepwidestar[1,2]{%
    \pi_2(P_u(n))\cdot u(n)=\left(\prod_{i<n}u(i)\right)\cdot u(n)}{\rIE{8}}
  \proofstepwidestar[1,2]{%
    \pi_2\bigl(P_u(\Succ(n))\bigr)=\left(\prod_{i<n}u(i)\right)\cdot u(n)}{%
    \FormulaRefAuto{a = b,\, b = c \vdash a = c}{6,9}}
  \proofstepwidestar[1,2]{%
    \prod_{i<\Succ(n)}u(i)=\left(\prod_{i<n}u(i)\right)\cdot u(n)}{%
    \FormulaRefAuto{a = b,\, b = c \vdash a = c}{4,10}}
\end{tabproofwide}

Für die Schreibweise mit oberer Grenze setzen wir entsprechend:

\subsubsection{Schreibweise mit oberer Grenze}

\FormulaDefDeltaK[Produktzeichen mit oberer Grenze]{%
  u\colon\mathbb{N}\to\mathbb{N}\dsep n\in\mathbb{N}\vdash
  \prod_{i=0}^{n}u(i)\coloneqq
  \prod_{i<\Succ(n)}u(i)%
}{Produktzeichen mit oberer Grenze}{%
  \DeltaRow{Mengen}{u\dsep n\dsep i}%
  \DeltaRow{Funktionensymbole}{\ProdStep{u}\dsep P_u}%
  \DeltaRow{Neue Schreibweise}{\prod_{i=0}^{n}u(i)}%
}

Damit sind Summen- und Produktzeichen auf Anfangsabschnitten natürlicher
Folgen vollständig auf die Rekursionsfunktion \(\RecFun{m_0}{S}\)
zurückgeführt. Spätere Rechenregeln für diese Zeichen sind daher keine
zusätzlichen Axiome, sondern Induktionssätze über die hier definierten
Rekursionsabbildungen.

\section{Potenzen natürlicher Zahlen}

Potenzen natürlicher Zahlen werden nach der Multiplikation als endliche
Produkte eingeführt. Damit gehören sie nicht zur Definition der
Multiplikation selbst, sondern zu den abgeleiteten Notationen, die auf
Anfangsabschnitten natürlicher Folgen beruhen. Als konstante Folge mit Wert
\(b\) verwenden wir die bereits eingeführte konstante Funktion
\(\ConstFun{\mathbb{N}}{\mathbb{N}}{b}\colon\mathbb{N}\to\mathbb{N}\).

\FormulaDefDeltaK[Potenz natürlicher Zahlen]{%
  b\in\mathbb{N}\dsep n\in\mathbb{N}\vdash
  b^n\coloneqq \prod_{i<n}\ConstFun{\mathbb{N}}{\mathbb{N}}{b}(i)%
}{Potenz natürlicher Zahlen}{%
  \DeltaRow{Mengen}{b\dsep n\dsep i}%
  \DeltaRow{Funktionensymbol}{\ConstFun{\mathbb{N}}{\mathbb{N}}{b}}%
  \DeltaRow{Neue Symbole}{(\,\cdot\,)^{(\,\cdot\,)}}%
}

\FormulaThmDelta[Abgeschlossenheit der Potenz natürlicher Zahlen]{%
  b\in\mathbb{N}\dsep n\in\mathbb{N}\vdash b^n\in\mathbb{N}%
}{%
  \DeltaRow{Mengen}{b\dsep n\dsep i}%
  \DeltaRow{Funktionensymbol}{\ConstFun{\mathbb{N}}{\mathbb{N}}{b}}%
}
\begin{tabproof}
  \proofstep{1}{b\in\mathbb{N}}{\rA}
  \proofstep{2}{n\in\mathbb{N}}{\rA}
  \proofstep{1}{%
    \ConstFun{\mathbb{N}}{\mathbb{N}}{b}\colon\mathbb{N}\to\mathbb{N}}{%
    \FormulaRefAuto{b\in B\vdash \ConstFun{A}{B}{b}\colon A\to B}{1}}
  \proofstep{1,2}{%
    b^n=\prod_{i<n}\ConstFun{\mathbb{N}}{\mathbb{N}}{b}(i)}{%
    \FormulaRefAuto{b\in\mathbb{N}\dsep n\in\mathbb{N}\vdash
      b^n\coloneqq \prod_{i<n}\ConstFun{\mathbb{N}}{\mathbb{N}}{b}(i)}{1,2}}
  \proofstep{1,2}{%
    \prod_{i<n}\ConstFun{\mathbb{N}}{\mathbb{N}}{b}(i)\in\mathbb{N}}{%
    \FormulaRefAuto{u\colon\mathbb{N}\to\mathbb{N}\dsep n\in\mathbb{N}\vdash
      \prod_{i<n}u(i)\in\mathbb{N}}{3,2}}
  \proofstep{1,2}{b^n\in\mathbb{N}}{\rIE{4,5}}
\end{tabproof}

Damit sind die Rekursionsgleichungen für Potenzen nur die Produktgleichungen
für konstante Folgen. Wir formulieren sie als eigene Sätze, weil sie die
Rekursionsform der Potenznotation festhalten.

\FormulaThmDelta[Nullregel der Potenz natürlicher Zahlen]{%
  b\in\mathbb{N}\vdash b^{0}=\Succ(0)%
}{%
  \DeltaRow{Mengen}{b\dsep i}%
  \DeltaRow{Funktionensymbol}{\ConstFun{\mathbb{N}}{\mathbb{N}}{b}}%
}
\begin{tabproofwide}
  \proofstepwidestar[1]{b\in\mathbb{N}}{\rA}
  \proofstepwidestar[]{0\in\mathbb{N}}{%
    \FormulaRefAuto{PeanoZero}}
  \proofstepwidestar[1]{%
    \ConstFun{\mathbb{N}}{\mathbb{N}}{b}\colon\mathbb{N}\to\mathbb{N}}{%
    \FormulaRefAuto{b\in B\vdash \ConstFun{A}{B}{b}\colon A\to B}{1}}
  \proofstepwidestar[1]{%
    b^{0}=
    \prod_{i<0}\ConstFun{\mathbb{N}}{\mathbb{N}}{b}(i)}{%
    \FormulaRefAuto{b\in\mathbb{N}\dsep n\in\mathbb{N}\vdash
      b^n\coloneqq \prod_{i<n}\ConstFun{\mathbb{N}}{\mathbb{N}}{b}(i)}{1,2}}
  \proofstepwidestar[1]{%
    \prod_{i<0}\ConstFun{\mathbb{N}}{\mathbb{N}}{b}(i)=
    \Succ(0)}{%
    \FormulaRefAuto{u\colon\mathbb{N}\to\mathbb{N}\vdash
      \prod_{i<0}u(i)=\Succ(0)}{3}}
  \proofstepwidestar[1]{b^{0}=\Succ(0)}{%
    \FormulaRefAuto{a = b,\, b = c \vdash a = c}{4,5}}
\end{tabproofwide}

\FormulaThmDeltaR[Nachfolgerregel der Potenz natürlicher Zahlen]{%
  \begin{aligned}[t]
    &b\in\mathbb{N}\dsep n\in\mathbb{N}\vdash{}\\
    &b^{\Succ(n)}=b^n\cdot b
  \end{aligned}
}{%
  b\in\mathbb{N}\dsep n\in\mathbb{N}\vdash b^{\Succ(n)}=b^n\cdot b%
}{%
  \DeltaRow{Mengen}{b\dsep n\dsep i\dsep m}%
  \DeltaRow{Funktionensymbol}{\ConstFun{\mathbb{N}}{\mathbb{N}}{b}}%
}
\begin{tabproofwide}
  \proofstepwidestar[1]{b\in\mathbb{N}}{\rA}
  \proofstepwidestar[2]{n\in\mathbb{N}}{\rA}
  \proofstepwidestar[1]{%
    \ConstFun{\mathbb{N}}{\mathbb{N}}{b}\colon\mathbb{N}\to\mathbb{N}}{%
    \FormulaRefAuto{b\in B\vdash \ConstFun{A}{B}{b}\colon A\to B}{1}}
  \proofstepwidestar[2]{\Succ(n)\in\mathbb{N}}{%
    \FormulaRefAuto{PeanoSuccClosure}{2}}
  \proofstepwidestar[1,2]{%
    b^{\Succ(n)}=
    \prod_{i<\Succ(n)}\ConstFun{\mathbb{N}}{\mathbb{N}}{b}(i)}{%
    \FormulaRefAuto{b\in\mathbb{N}\dsep n\in\mathbb{N}\vdash
      b^n\coloneqq \prod_{i<n}\ConstFun{\mathbb{N}}{\mathbb{N}}{b}(i)}{1,4}}
  \proofstepwidestar[1,2]{%
    \prod_{i<\Succ(n)}\ConstFun{\mathbb{N}}{\mathbb{N}}{b}(i)
    =
    \left(\prod_{i<n}\ConstFun{\mathbb{N}}{\mathbb{N}}{b}(i)\right)\cdot
    \ConstFun{\mathbb{N}}{\mathbb{N}}{b}(n)}{%
    \FormulaRefAuto{u\colon\mathbb{N}\to\mathbb{N}\dsep n\in\mathbb{N}\vdash
      \prod_{i<\Succ(n)}u(i)=\left(\prod_{i<n}u(i)\right)\cdot u(n)}{3,2}}
  \proofstepwidestar[1,2]{%
    \ConstFun{\mathbb{N}}{\mathbb{N}}{b}(n)=b}{%
    \FormulaRefAuto{b\in B\dsep x\in A\vdash \ConstFun{A}{B}{b}(x)=b}{1,2}}
  \proofstepwidestar[1,2]{%
    \left(\prod_{i<n}\ConstFun{\mathbb{N}}{\mathbb{N}}{b}(i)\right)\cdot
    \ConstFun{\mathbb{N}}{\mathbb{N}}{b}(n)
    =
    \left(\prod_{i<n}\ConstFun{\mathbb{N}}{\mathbb{N}}{b}(i)\right)\cdot b}{\rIE{7}}
  \proofstepwidestar[1,2]{%
    \prod_{i<\Succ(n)}\ConstFun{\mathbb{N}}{\mathbb{N}}{b}(i)
    =
    \left(\prod_{i<n}\ConstFun{\mathbb{N}}{\mathbb{N}}{b}(i)\right)\cdot b}{%
    \FormulaRefAuto{a = b,\, b = c \vdash a = c}{6,8}}
  \proofstepwidestar[1,2]{%
    b^n=\prod_{i<n}\ConstFun{\mathbb{N}}{\mathbb{N}}{b}(i)}{%
    \FormulaRefAuto{b\in\mathbb{N}\dsep n\in\mathbb{N}\vdash
      b^n\coloneqq \prod_{i<n}\ConstFun{\mathbb{N}}{\mathbb{N}}{b}(i)}{1,2}}
  \proofstepwidestar[1,2]{%
    \prod_{i<n}\ConstFun{\mathbb{N}}{\mathbb{N}}{b}(i)=b^n}{%
    \FormulaRefAuto{a=b\vdash b=a}{10}}
  \proofstepwidestar[1,2]{%
    \left(\prod_{i<n}\ConstFun{\mathbb{N}}{\mathbb{N}}{b}(i)\right)\cdot b
    =b^n\cdot b}{\rIE{11}}
  \proofstepwidestar[1,2]{%
    \prod_{i<\Succ(n)}\ConstFun{\mathbb{N}}{\mathbb{N}}{b}(i)=b^n\cdot b}{%
    \FormulaRefAuto{a = b,\, b = c \vdash a = c}{9,12}}
  \proofstepwidestar[1,2]{b^{\Succ(n)}=b^n\cdot b}{%
    \FormulaRefAuto{a = b,\, b = c \vdash a = c}{5,13}}
\end{tabproofwide}

\chapter{\(b\)-adische Stellenwertsysteme für natürliche Zahlen}

Die zuvor eingeführten endlichen Summen und Potenzen erlauben nun eine
einheitliche Definition von Stellenwertdarstellungen natürlicher Zahlen. Eine
endliche Ziffernfolge ist dabei nicht selbst die natürliche Zahl, sondern ein
Zahlzeichen, dessen Wert durch eine Auswertungsfunktion bestimmt wird. Damit
die späteren Beweise eindeutig auf die hier eingeführten Begriffe verweisen
können, werden Basisbedingung, Ziffernmenge, Ziffernfolge, Auswertung und
Zahlzeichen als eigene Symbole festgelegt.

\section{Zahlzeichen, Basen und Ziffern}

Das Zahlzeichen \(1\) wurde bereits vor der additiven Ordnung eingeführt.
Die übrigen vertrauten kleinen Zahlzeichen werden hier als Abkürzungen für
die entsprechenden natürlichen Zahlen ergänzt; jedes folgende Zahlzeichen
bezeichnet den Nachfolger des vorherigen.

\FormulaDefDeltaK[Natürliche Zahlzeichen von zwei bis zehn]{%
  \begin{gathered}
    2\coloneqq\Succ(1),\qquad
    3\coloneqq\Succ(2),\\
    4\coloneqq\Succ(3),\qquad
    5\coloneqq\Succ(4),\qquad
    6\coloneqq\Succ(5),\qquad
    7\coloneqq\Succ(6),\\
    8\coloneqq\Succ(7),\qquad
    9\coloneqq\Succ(8),\qquad
    10\coloneqq\Succ(9)
  \end{gathered}
}{Natürliche Zahlzeichen von zwei bis zehn}{%
  \DeltaRow{Mengen}{1\dsep 2\dsep 3\dsep 4\dsep 5\dsep 6\dsep 7\dsep 8\dsep 9\dsep 10}%
  \DeltaRow{Funktionensymbol}{\Succ}%
}

Eine Basis eines Stellenwertsystems ist eine natürliche Zahl, die mindestens
\(2\) ist. Diese Bedingung wird als eigenes Prädikat notiert.

\FormulaDefDeltaK[Basis eines Stellenwertsystems]{%
  \BaseOK{b}\coloneqq b\in\mathbb{N}\land 2\leq b%
}{Basis eines Stellenwertsystems}{%
  \DeltaRow{Mengen}{b}%
  \DeltaRow{Neue Symbole}{\BaseOK{b}}%
}

\FormulaThmDelta[Basis ist natürliche Zahl]{%
  \BaseOK{b}\vdash b\in\mathbb{N}%
}{%
  \DeltaRow{Mengen}{b}%
  \DeltaRow{Symbol}{\BaseOK{b}}%
}
\begin{tabproof}
  \proofstep{1}{\BaseOK{b}}{\rA}
  \proofstep{1}{b\in\mathbb{N}\land 2\leq b}{%
    \FormulaRefAuto{\BaseOK{b}\coloneqq b\in\mathbb{N}\land 2\leq b}{1}}
  \proofstep{1}{b\in\mathbb{N}}{\rAEa{2}}
\end{tabproof}

Die zulässigen Ziffern zur Basis \(b\) sind genau die natürlichen Zahlen, die
echt kleiner als \(b\) sind.

\FormulaDefDeltaK[Ziffernmenge zur Basis \(b\)]{%
  \BaseOK{b}\vdash
  \DigitSet{b}\coloneqq \{d\in\mathbb{N}\mid d<b\}%
}{Ziffernmenge zur Basis \(b\)}{%
  \DeltaRow{Mengen}{b\dsep d}%
  \DeltaRow{Neue Symbole}{\DigitSet{b}}%
}

\section{Ziffernfolgen und Stellenwertauswertung}

Eine Ziffernfolge zur Basis \(b\) wird als Funktion
\(a\colon\mathbb{N}\to\mathbb{N}\) behandelt. Das Symbol
\(\DigitSeq{b}{a}{\ell}\) bedeutet, dass die Anfangswerte
\(a(0),\ldots,a(\ell)\) zulässige Ziffern zur Basis \(b\) sind. Der letzte
Parameter \(\ell\) bezeichnet also die höchste Stelle, nicht die Länge.

\FormulaDefDeltaR[Ziffernfolge zur Basis \(b\)]{%
  \begin{aligned}[t]
    &\BaseOK{b}\dsep \ell\in\mathbb{N}\vdash{}\\
    &\DigitSeq{b}{a}{\ell}\coloneqq
      \bigl(a\colon\mathbb{N}\to\mathbb{N}\land
      \forall i<\Succ(\ell)\,a(i)\in\DigitSet{b}\bigr)
  \end{aligned}
}{%
  \BaseOK{b}\dsep \ell\in\mathbb{N}\vdash
  \DigitSeq{b}{a}{\ell}\coloneqq
  \bigl(a\colon\mathbb{N}\to\mathbb{N}\land
  \forall i<\Succ(\ell)\,a(i)\in\DigitSet{b}\bigr)
}{%
  \DeltaRow{Mengen}{a\dsep b\dsep \ell\dsep i}%
  \DeltaRow{Neue Symbole}{\DigitSeq{b}{a}{\ell}}%
  \DeltaRow{Funktionensymbol}{\Succ}%
}

Das Summenzeichen wurde oben für Folgen \(u\colon\mathbb{N}\to\mathbb{N}\)
eingeführt. Damit die Stellenwertsumme genau in diese Form fällt, trennen wir
wie bei Summen- und Produktschritten zunächst die punktweise Vorschrift von
der durch sie bestimmten Funktion.

\FormulaDefDeltaK[Funktionsvorschrift der Stellenwertfolge]{%
  \BaseOK{b}\dsep a\colon\mathbb{N}\to\mathbb{N}\dsep i\in\mathbb{N}\vdash
  \tPlaceTermSeq{b}{a}(i)\coloneqq a(i)\cdot b^i%
}{Funktionsvorschrift der Stellenwertfolge}{%
  \DeltaRow{Mengen}{a\dsep b\dsep i}%
  \DeltaRow{Funktionensymbol}{\tPlaceTermSeq{b}{a}}%
}

\FormulaThmDelta[Typisierung des Stellenwertterms]{%
  \BaseOK{b}\dsep a\colon\mathbb{N}\to\mathbb{N}\dsep i\in\mathbb{N}
  \vdash \tPlaceTermSeq{b}{a}(i)\in\mathbb{N}%
}{%
  \DeltaRow{Mengen}{a\dsep b\dsep i}%
  \DeltaRow{Funktionensymbol}{\tPlaceTermSeq{b}{a}}%
}
\begin{tabproof}
  \proofstep{1}{\BaseOK{b}}{\rA}
  \proofstep{2}{a\colon\mathbb{N}\to\mathbb{N}}{\rA}
  \proofstep{3}{i\in\mathbb{N}}{\rA}
  \proofstep{1}{b\in\mathbb{N}}{%
    \FormulaRefAuto{\BaseOK{b}\vdash b\in\mathbb{N}}{1}}
  \proofstep{2,3}{a(i)\in\mathbb{N}}{%
    \FormulaRefAuto{F\colon A\to B\dsep x\in A\vdash F(x)\in B}{2,3}}
  \proofstep{1,3}{b^i\in\mathbb{N}}{%
    \FormulaRefAuto{b\in\mathbb{N}\dsep n\in\mathbb{N}\vdash
      b^n\in\mathbb{N}}{4,3}}
  \proofstep{1,2,3}{a(i)\cdot b^i\in\mathbb{N}}{%
    \FormulaRefAuto{a\in\mathbb{N}\dsep b\in\mathbb{N}\vdash
      a\cdot b\in\mathbb{N}}{5,6}}
  \proofstep{1,2,3}{\tPlaceTermSeq{b}{a}(i)=a(i)\cdot b^i}{%
    \FormulaRefAuto{\BaseOK{b}\dsep a\colon\mathbb{N}\to\mathbb{N}\dsep
      i\in\mathbb{N}\vdash
      \tPlaceTermSeq{b}{a}(i)\coloneqq a(i)\cdot b^i}{1,2,3}}
  \proofstep{1,2,3}{\tPlaceTermSeq{b}{a}(i)\in\mathbb{N}}{\rIE{8,7}}
\end{tabproof}

\FormulaThmDelta[Existenz der Stellenwertfolge zur Basis \(b\)]{%
  \BaseOK{b}\dsep a\colon\mathbb{N}\to\mathbb{N}\vdash
  \exists! w\colon\mathbb{N}\to\mathbb{N}\,
  \forall i\in\mathbb{N}\,w(i)=\tPlaceTermSeq{b}{a}(i)%
}{%
  \DeltaRow{Mengen}{a\dsep b\dsep i}%
  \DeltaRow{Funktionensymbol}{\tPlaceTermSeq{b}{a}}%
}
\begin{tabproof}
  \proofstep{1}{\BaseOK{b}}{\rA}
  \proofstep{2}{a\colon\mathbb{N}\to\mathbb{N}}{\rA}
  \proofstep{3}{i\in\mathbb{N}}{\rA}
  \proofstep{1,2,3}{\tPlaceTermSeq{b}{a}(i)=a(i)\cdot b^i}{%
    \FormulaRefAuto{\BaseOK{b}\dsep a\colon\mathbb{N}\to\mathbb{N}\dsep
      i\in\mathbb{N}\vdash
      \tPlaceTermSeq{b}{a}(i)\coloneqq a(i)\cdot b^i}{1,2,3}}
  \proofstep{1,2}{%
    \forall i\in\mathbb{N}\,\tPlaceTermSeq{b}{a}(i)=a(i)\cdot b^i}{%
    \rUI{\rRI{3,4}}}
  \proofstep{1,2,3}{\tPlaceTermSeq{b}{a}(i)\in\mathbb{N}}{%
    \FormulaRefAuto{\BaseOK{b}\dsep a\colon\mathbb{N}\to\mathbb{N}\dsep
      i\in\mathbb{N}\vdash
      \tPlaceTermSeq{b}{a}(i)\in\mathbb{N}}{1,2,3}}
  \proofstep{1,2}{%
    \forall i\in\mathbb{N}\,\tPlaceTermSeq{b}{a}(i)\in\mathbb{N}}{%
    \rUI{\rRI{3,6}}}
  \proofstep{1,2}{%
    \exists! w\colon\mathbb{N}\to\mathbb{N}\,
    \forall i\in\mathbb{N}\,w(i)=\tPlaceTermSeq{b}{a}(i)}{%
    \FormulaRefAuto{%
      x\in A\vdash t(x)\coloneq y\dsep x\in A\vdash t(x)\in B
      \exists! F\colon A\to B\,\forall x\in A\,F(x)=t(x)}{5,7}}
\end{tabproof}

Die Stellenwertfolge \(\PlaceTermSeq{b}{a}\) ist die durch diese Vorschrift
eindeutig bestimmte Folge. Ihr \(i\)-tes Glied ist der Beitrag der \(i\)-ten
Ziffer zur Basis \(b\).

\FormulaDefDeltaR[Stellenwertfolge zur Basis \(b\)]{%
  \begin{aligned}[t]
    &\BaseOK{b}\dsep a\colon\mathbb{N}\to\mathbb{N}\vdash{}\\
    &\PlaceTermSeq{b}{a}\coloneqq \iota w\Bigl(
      w\colon\mathbb{N}\to\mathbb{N}\land{}\\
    &\qquad \forall i\in\mathbb{N}\,
      w(i)=\tPlaceTermSeq{b}{a}(i)
    \Bigr)
  \end{aligned}
}{%
  \BaseOK{b}\dsep a\colon\mathbb{N}\to\mathbb{N}\vdash
  \PlaceTermSeq{b}{a}\coloneqq \iota w\Bigl(
    w\colon\mathbb{N}\to\mathbb{N}\land
    \forall i\in\mathbb{N}\,w(i)=\tPlaceTermSeq{b}{a}(i)
  \Bigr)
}{%
  \DeltaRow{Mengen}{b\dsep i}%
  \DeltaRow{Funktionen}{a\dsep w}%
  \DeltaRow{Funktionensymbole}{\tPlaceTermSeq{b}{a}\dsep \PlaceTermSeq{b}{a}}%
}

\FormulaThmDelta[Stellenwertfolge ist eine Folge natürlicher Zahlen]{%
  \BaseOK{b}\dsep a\colon\mathbb{N}\to\mathbb{N}\vdash
  \PlaceTermSeq{b}{a}\colon\mathbb{N}\to\mathbb{N}%
}{%
  \DeltaRow{Mengen}{a\dsep b}%
  \DeltaRow{Funktionensymbol}{\PlaceTermSeq{b}{a}}%
}
\begin{tabproof}
  \proofstep{1}{\BaseOK{b}}{\rA}
  \proofstep{2}{a\colon\mathbb{N}\to\mathbb{N}}{\rA}
  \proofstep{1,2}{\PlaceTermSeq{b}{a}\colon\mathbb{N}\to\mathbb{N}}{%
    \rAEa{\FormulaRefAuto{\BaseOK{b}\dsep a\colon\mathbb{N}\to\mathbb{N}\vdash
      \PlaceTermSeq{b}{a}\coloneqq \iota w\Bigl(
        w\colon\mathbb{N}\to\mathbb{N}\land
        \forall i\in\mathbb{N}\,w(i)=\tPlaceTermSeq{b}{a}(i)
      \Bigr)}{1,2}}}
\end{tabproof}

\FormulaThmDelta[Grundgleichung der Stellenwertfolge]{%
  \BaseOK{b}\dsep a\colon\mathbb{N}\to\mathbb{N}\dsep i\in\mathbb{N}\vdash
  \PlaceTermSeq{b}{a}(i)=a(i)\cdot b^i%
}{%
  \DeltaRow{Mengen}{a\dsep b\dsep i}%
  \DeltaRow{Funktionensymbole}{\tPlaceTermSeq{b}{a}\dsep \PlaceTermSeq{b}{a}}%
}
\begin{tabproof}
  \proofstep{1}{\BaseOK{b}}{\rA}
  \proofstep{2}{a\colon\mathbb{N}\to\mathbb{N}}{\rA}
  \proofstep{3}{i\in\mathbb{N}}{\rA}
  \proofstep{1,2}{%
    \forall i\in\mathbb{N}\,
    \PlaceTermSeq{b}{a}(i)=\tPlaceTermSeq{b}{a}(i)}{%
    \rAEb{\FormulaRefAuto{\BaseOK{b}\dsep a\colon\mathbb{N}\to\mathbb{N}\vdash
      \PlaceTermSeq{b}{a}\coloneqq \iota w\Bigl(
        w\colon\mathbb{N}\to\mathbb{N}\land
        \forall i\in\mathbb{N}\,w(i)=\tPlaceTermSeq{b}{a}(i)
      \Bigr)}{1,2}}}
  \proofstep{1,2,3}{\PlaceTermSeq{b}{a}(i)=\tPlaceTermSeq{b}{a}(i)}{%
    \rRE{3,\rUE{4}}}
  \proofstep{1,2,3}{\tPlaceTermSeq{b}{a}(i)=a(i)\cdot b^i}{%
    \FormulaRefAuto{\BaseOK{b}\dsep a\colon\mathbb{N}\to\mathbb{N}\dsep
      i\in\mathbb{N}\vdash
      \tPlaceTermSeq{b}{a}(i)\coloneqq a(i)\cdot b^i}{1,2,3}}
  \proofstep{1,2,3}{\PlaceTermSeq{b}{a}(i)=a(i)\cdot b^i}{%
    \FormulaRefAuto{a = b,\, b = c \vdash a = c}{5,6}}
\end{tabproof}

Hat eine Ziffernfolge die betrachtete Länge \(n\), so wird ihr Wert durch die
endliche Summe über die Stellenwertfolge bestimmt. Die Schreibweise
\(\sum_{i<n}a(i)\cdot b^i\) ist damit die ausgeschriebene Form von
\(\sum_{i<n}\PlaceTermSeq{b}{a}(i)\).

\FormulaDefDeltaR[\(b\)-adische Stellenwertauswertung]{%
  \begin{aligned}[t]
    &\BaseOK{b}\dsep n\in\mathbb{N}\dsep
      a\colon\mathbb{N}\to\mathbb{N}\dsep
      \forall i<n\,a(i)\in\DigitSet{b}\vdash{}\\
    &\PlaceVal{b}{a}{n}\coloneqq
      \sum_{i<n} \PlaceTermSeq{b}{a}(i)
  \end{aligned}
}{%
  \BaseOK{b}\dsep n\in\mathbb{N}\dsep
  a\colon\mathbb{N}\to\mathbb{N}\dsep
  \forall i<n\,a(i)\in\DigitSet{b}\vdash
  \PlaceVal{b}{a}{n}\coloneqq
  \sum_{i<n} \PlaceTermSeq{b}{a}(i)
}{%
  \DeltaRow{Mengen}{a\dsep b\dsep n\dsep i}%
  \DeltaRow{Funktionensymbole}{\PlaceTermSeq{b}{a}\dsep \PlaceVal{b}{a}{n}}%
}

In der üblichen Schreibweise werden die Ziffern von links nach rechts
notiert, während die Summe von der niedrigsten Stelle \(i=0\) ausgeht. Für
eine Ziffernfolge \(\DigitSeq{b}{a}{\ell}\) wird das zugehörige
\(b\)-adische Zahlzeichen deshalb mit dem Basisindex \(b\) notiert.

\FormulaDefDeltaR[\(b\)-adisches Zahlzeichen]{%
  \begin{aligned}[t]
    &\BaseOK{b}\dsep \ell\in\mathbb{N}\dsep
      \DigitSeq{b}{a}{\ell}\vdash{}\\
    &\BaseNum{b}{a}{\ell}\coloneqq
      \PlaceVal{b}{a}{\Succ(\ell)}
  \end{aligned}
}{%
  \BaseOK{b}\dsep \ell\in\mathbb{N}\dsep
  \DigitSeq{b}{a}{\ell}\vdash
  \BaseNum{b}{a}{\ell}\coloneqq
  \PlaceVal{b}{a}{\Succ(\ell)}
}{%
  \DeltaRow{Mengen}{a\dsep b\dsep \ell}%
  \DeltaRow{Funktionensymbol}{\PlaceVal{b}{a}{\Succ(\ell)}}%
  \DeltaRow{Neue Schreibweise}{\BaseNum{b}{a}{\ell}}%
}

Die rechte Ziffer \(a_0\) ist also die Einerstelle, \(a_1\) die \(b\)-Stelle,
\(a_2\) die \(b^2\)-Stelle und so weiter. In formalen Beweisen steht
\(\BaseNum{b}{a}{\ell}\) immer für die durch \(a\), \(b\) und \(\ell\)
festgelegte Auswertung, nicht für eine unabhängige Zeichenkette.

\section{Normalisierte Darstellungen}

Eine Darstellung ist normalisiert, wenn vor der höchsten Stelle keine
führende Null steht. Bei einer Darstellung mit nur einer Stelle ist auch die
Ziffer \(0\) normalisiert.

\FormulaDefDeltaR[Normalisierte Ziffernfolge]{%
  \begin{aligned}[t]
    &\BaseOK{b}\dsep \ell\in\mathbb{N}\vdash{}\\
    &\NormDigitSeq{b}{a}{\ell}\coloneqq
      \bigl(\DigitSeq{b}{a}{\ell}\land
      (\ell=0\lor a(\ell)\neq 0)\bigr)
  \end{aligned}
}{%
  \BaseOK{b}\dsep \ell\in\mathbb{N}\vdash
  \NormDigitSeq{b}{a}{\ell}\coloneqq
  \bigl(\DigitSeq{b}{a}{\ell}\land
  (\ell=0\lor a(\ell)\neq 0)\bigr)
}{%
  \DeltaRow{Mengen}{a\dsep b\dsep \ell}%
  \DeltaRow{Neue Symbole}{\NormDigitSeq{b}{a}{\ell}}%
}

Führende Nullen ändern den Wert der Stellenwertsumme nicht; sie gehören
deshalb nicht zu normalisierten Darstellungen mit mehr als einer Stelle.

\section{Dualsystem}

Das Dualsystem ist das \(b\)-adische Stellenwertsystem zur Basis \(2\). Seine
Ziffernmenge ist
\[
  \BinDigitSet
  =
  \{d\in\mathbb{N}\mid d<2\}
  =
  \{0,1\}.
\]
Die beiden Ziffern \(0\) und \(1\) heißen im Dualsystem auch \emph{Bits}. Für
binäre Ziffernfolgen verwenden wir das basisindizierte Symbol
\(\BinDigitSeq{a}{\ell}\).

\FormulaDefDeltaK[Binäre Ziffernfolge]{%
  \ell\in\mathbb{N}\vdash
  \BinDigitSeq{a}{\ell}\coloneqq \DigitSeq{2}{a}{\ell}%
}{Binäre Ziffernfolge}{%
  \DeltaRow{Mengen}{a\dsep \ell}%
  \DeltaRow{Neue Symbole}{\BinDigitSeq{a}{\ell}}%
}

\FormulaDefDeltaK[Binäres Zahlzeichen]{%
  \ell\in\mathbb{N}\dsep \BinDigitSeq{a}{\ell}\vdash
  \BinNum{a}{\ell}\coloneqq \PlaceVal{2}{a}{\Succ(\ell)}%
}{Binäres Zahlzeichen}{%
  \DeltaRow{Mengen}{a\dsep \ell}%
  \DeltaRow{Neue Schreibweise}{\BinNum{a}{\ell}}%
}

Damit gilt für jede binäre Ziffernfolge \(\BinDigitSeq{a}{\ell}\):
\[
  \BinNum{a}{\ell}
  =
  \sum_{i<\Succ(\ell)} \PlaceTermSeq{2}{a}(i)
  =
  \sum_{i<\Succ(\ell)} a(i)\cdot 2^i.
\]
Zum Beispiel sind
\[
  \BinDigitString{0}=0,\qquad
  \BinDigitString{1}=1,\qquad
  \BinDigitString{10}=2,\qquad
  \BinDigitString{101}=1\cdot 2^2+0\cdot 2^1+1\cdot 2^0.
\]

\section{Dezimalsystem}

Das Dezimalsystem ist das \(b\)-adische Stellenwertsystem zur Basis \(10\).
Da dieses System im Folgenden das Standardsystem für die gewöhnliche
Schreibweise natürlicher Zahlen ist, erhalten seine Ziffernmenge,
Ziffernfolgen und Zahlzeichen Symbole ohne Basisindex \(10\).

\FormulaDefDeltaK[Dezimale Ziffernmenge]{%
  \DecDigitSet\coloneqq \DigitSet{10}%
}{Dezimale Ziffernmenge}{%
  \DeltaRow{Mengen}{10}%
  \DeltaRow{Neue Symbole}{\DecDigitSet}%
}

Damit ist
\[
  \DecDigitSet
  =
  \{d\in\mathbb{N}\mid d<10\}
  =
  \{0,1,2,3,4,5,6,7,8,9\}.
\]
Die Ziffern des Dezimalsystems sind also genau die zehn Zeichen
\[
  0,\ 1,\ 2,\ 3,\ 4,\ 5,\ 6,\ 7,\ 8,\ 9.
\]

\FormulaDefDeltaK[Dezimale Ziffernfolge]{%
  \ell\in\mathbb{N}\vdash
  \DecDigitSeq{a}{\ell}\coloneqq \DigitSeq{10}{a}{\ell}%
}{Dezimale Ziffernfolge}{%
  \DeltaRow{Mengen}{a\dsep \ell}%
  \DeltaRow{Neue Symbole}{\DecDigitSeq{a}{\ell}}%
}

\FormulaDefDeltaK[Dezimale Stellenwertauswertung]{%
  n\in\mathbb{N}\dsep a\colon\mathbb{N}\to\mathbb{N}\dsep
  \forall i<n\,a(i)\in\DecDigitSet\vdash
  \DecPlaceVal{a}{n}\coloneqq \PlaceVal{10}{a}{n}%
}{Dezimale Stellenwertauswertung}{%
  \DeltaRow{Mengen}{a\dsep n\dsep i}%
  \DeltaRow{Funktionensymbol}{\DecPlaceVal{a}{n}}%
}

\FormulaDefDeltaK[Dezimales Zahlzeichen]{%
  \ell\in\mathbb{N}\dsep \DecDigitSeq{a}{\ell}\vdash
  \DecNum{a}{\ell}\coloneqq \DecPlaceVal{a}{\Succ(\ell)}%
}{Dezimales Zahlzeichen}{%
  \DeltaRow{Mengen}{a\dsep \ell}%
  \DeltaRow{Neue Schreibweise}{\DecNum{a}{\ell}}%
}

Für jede dezimale Ziffernfolge \(\DecDigitSeq{a}{\ell}\) gilt folglich
\[
  \DecNum{a}{\ell}
  =
  \sum_{i<\Succ(\ell)} \PlaceTermSeq{10}{a}(i)
  =
  \sum_{i<\Succ(\ell)} a(i)\cdot 10^i.
\]
Das sichtbare Zahlzeichen trägt im Dezimalsystem weder Klammern noch
Basisindex; es steht direkt für die dezimale Stellenwertauswertung.



\end{document}
